% BEGIN CONTINUATION_NIGHT36_RECEIVING

\subsection{Checkpoint36 source results with continuation review pending}

The following source statements and their complete supplied arguments are
retained here while their continuation checks are completed. Their source
records report proofs by GPT-6 Astra (OpenAI), ChatGPT Pro mode, effort Pro,
S6\_06\_BLIND, 16--18 November 2026. Each receiving status records the
current scope below; importing the source does not add a proof check.
The separate event ledger records the algebraic substeps already checked
with new code and every later mathematical comparison.


\subsection{The full holomorphic endpoint group of the generic marked E6 source}
\begin{theorem}[Source result P3K.109 under continuation review]\label{night:P3K.109}
\StatementRecord{night:P3K.109}{open}{Historical events: GPT-6 Astra (OpenAI), ChatGPT Pro mode, effort Pro, S6\_06\_BLIND, 16--18 November 2026: targeted proof check. Also checked in this writing. Continuation mathematical and source checks: GPT-6 Astra (OpenAI), ChatGPT Pro mode, effort Pro, S6\_10\_PAPERS continuation, 19 November 2026: computation or formal certificate; targeted proof check; used and rechecked in this run. Continuation metadata events: GPT-6 Astra (OpenAI), ChatGPT Pro mode, effort Pro, S6\_10\_PAPERS continuation, 19 November 2026: used and rechecked in this run; explicitly scoped metadata reconciliation; metadata reconciliation only; no new mathematical check. GPT-6 Astra (OpenAI), ChatGPT Pro mode, effort Pro, S6\_10\_PAPERS continuation, 19 November 2026: used and rechecked in this run; source-row intake; metadata reconciliation only; no new mathematical check. GPT-6 Astra (OpenAI), ChatGPT Pro mode, effort Pro, S6\_10\_PAPERS continuation, 19 November 2026: used and rechecked in this run; receiving source placement; metadata reconciliation only; no new mathematical check. \emph{Pre-continuation recorded proof scope.} Historical event: GPT-6 Astra (OpenAI), ChatGPT Pro mode, effort Pro, S6\_06\_BLIND, 16--18 November 2026: targeted proof check, recorded producing-run re-derivation. Also checked in this writing by GPT-6 Astra (OpenAI), ChatGPT Pro mode, effort Pro, S6\_10\_PAPERS continuation, 19 November 2026: computation or formal certificate for the explicitly listed algebraic substeps in the event ledger; full continuation rederivation and renewed credit comparison remain pending.}{\cref{night:P3K.94,night:P3K.96,night:P3F.35,night:P3F.46}; the original source, point line, roots, characters, reference, selected fan and scalar range stated here}
Set
\[
 P_r=\frac{324r^3-405r^2+144r-16}{2187},\qquad
 Q_r=\frac{2(9r-2)(972r^3-891r^2+288r-32)}{531441},
\]
and let \(r\) lie in the explicit source locus
\[
 \mathcal U=\{r\in\mathbb C:r(4r-1)(3r-1)P_rQ_r\ne0\}.
\]
Retain the original source
\[
 E_r:\ y^2=x^3+rx^2+tx+t^2,\qquad G=(0,t),
\]
one chosen root \(t_c\) of
\(27t^2+(4-18r)t+r^2(4r-1)=0\), the selected height
\(q=t-t_c\), its actual branch marking, and any prescribed
lattice-basis fan. Let the point line and suspension be
\(M_a=\mathcal O(aG-O)\) and translation by \(-bG\), with
either of the complete cap progressions
\[
 (a,b;m)=(6,1;3l+2),\quad l\ge0,
 \qquad
 (a,b;m)=(2,3;3l),\quad l\ge1.
\]
Use either signed primitive norm \(n=\pm1\), the prescribed finite
character \(\zeta_m^d\), the full annular reference of P3K.96,
the ordinary attachments at the other finite source fibres, and a
fixed nonzero scalar in the strict contraction range. Let \(X\)
be the resulting completed complex threefold.

Multiplication on the genuine point line, carried through the full
root, reference and toric maps, gives the exact group isomorphism
\[
       \mathbb C^*_{\rm physical}
          \xrightarrow{\ \cong\ }\operatorname{Aut}_{\rm hol}(X).
\]
Every holomorphic automorphism preserves the original base pointwise
and is isotopic to the identity through this physical action. For
two constructed holomorphic endpoint identifications of the same
cut family, their ratio is precisely such a physical scalar; the
inverse-endpoint returns are therefore isotopic by the displayed
completed maps.

The proof includes all possible descended base maps on this explicit
source locus, both cap progressions, both norm signs and every
retained finite character. It assigns no new source or scalar return
and no shorter endpoint outside the constructed family. A concrete
object is
\[
 r=\tfrac12,\quad t_c=(5+i\sqrt2)/54,\quad
 (a,b;m,n,d)=(6,1;2,1,0),\quad T_0=T(e_1,e_2),
\]
with the complete P3K.94--96 reference.
\end{theorem}
\begin{proof}[Supplied source argument; continuation review in progress]

% Expanded verified cp36 input: notes/part3_e6_fixed_source_endpoints_proof.tex
\par

\medskip\noindent\textbf{EF.1. The actual source locus, line, cap and first object}

\par\nobreak

Put
\[
 E_r:\quad y^2=x^3+rx^2+tx+t^2,\qquad G=(0,t),\qquad
 s=\frac14-r.
 \tag{EF.1}
\]
Write
\[
 \begin{split}
 S_r&=(18r-4)/27,\quad C_r=r^2(4r-1)/27,\\
 P_r&=(324r^3-405r^2+144r-16)/2187,\\
 Q_r&=2(9r-2)(972r^3-891r^2+288r-32)/531441.
 \end{split}                                                    \tag{EF.2}
\]
The source locus in this proof is the explicit open set
\[
 \mathcal U=\{r\in\mathbb C:
 r(4r-1)(3r-1)P_rQ_r\ne0\}.
 \tag{EF.3}
\]
Choose one of the two roots of
\(27t^2+(4-18r)t+r^2(4r-1)=0\), call it \(t_c\), and retain
that marked selected I1, its actual branch coordinate and its prescribed
lattice-basis fan. Write \(q=t-t_c\). The other I1 and the I2 at zero
have their ordinary source attachments. There is one finite cap at the
original infinity, with the entire reference map of K96.

The two data in the scope are
\[
 (a,b;m)=(6,1;3l+2),\ l\ge0,
 \qquad (a,b;m)=(2,3;3l),\ l\ge1,
 \qquad n\in\{1,-1\}.
 \tag{EF.4}
\]
The point line is the actual \(M_a=\mathcal O(aG-O)\); the source
suspension is translation by \(-bG\). The character \(\zeta_m^d\)
and the whole annular reference cochain are retained. The scalar is a
fixed nonzero value in the contraction domain of the construction. The
source, root and line maps in EX.1-\/-19 of K96 construct these actual
models. The integral maps of F46 and the common blowup, two-ruling,
Gysin and van Kampen maps of F35 give the sphere on each prescribed fan.
Equivalently, one can transport the explicit K96 sphere along a compact
path of marked simple selected nodes to the chosen point of EF.3: all
EX charts are holomorphic along that path, and a common smaller scalar
gives its strict compact norm shell. Changing that scalar within the
contraction domain gives the same smooth sphere. Thus the completed
threefold used here has \(H^2=H^4=0\), with its actual source fibration.

A fully specified first object is
\[
 r=\tfrac12,\quad t_c=(5+i\sqrt2)/54,\quad
 (a,b;m,n,d)=(6,1;2,1,0),\quad T_0=T(e_1,e_2),
 \tag{EF.5}
\]
with the fixed small scalar and the K94/K96 reference lift. Here
\(P_r=-19/8748\) and \(Q_r=215/2125764\), so it belongs to EF.3.

The conclusion proved below is the exact group equality
\[
       \operatorname{Aut}_{\rm hol}(X)=\mathbb C^*_{\rm physical}.
 \tag{EF.6}
\]
The right-hand map is multiplication on the actual point line, followed
through all root, annular and Cox arrows. In particular it is injective
and every endpoint automorphism is joined to the identity by that action.

\par

\medskip\noindent\textbf{EF.2. Recover the base and fix it by the actual critical values}

\par\nobreak

Let \(\pi:X\to\mathbb P^1_t\) be the constructed fibration. A smooth
fibre is a compact complex two-torus. Give it any flat Kähler form
\(\omega\). A horizontal effective divisor on \(X\) would restrict
to a nonzero effective divisor \(D_T\) on a generic fibre. Since its
line class restricts from \(H^2(X)=0\),
\[
 0=\int_T c_1(\mathcal O(D_T))\wedge\omega
   =\int_{D_T}\omega>0.
 \tag{EF.7}
\]
The divisor fundamental-class identity includes every multiplicity, so
this excludes every horizontal polar divisor of a meromorphic function.
The function is constant on each generic compact fibre, by removal of
its remaining codimension-two indeterminacy there. It therefore descends
on the regular base. At a generic smooth point of a reduced special-fibre
component choose a transverse disc. After taking the extending root of
its nonzero base unit, its exact base map is \(t-t_0=u^e\).
The pulled-back meromorphic function has a Laurent expansion with finite
principal part. Constancy on the generic fibres makes it invariant under
\(u\mapsto\zeta_eu\), so every occurring exponent is a multiple of
\(e\). This gives its meromorphic base extension at \(t_0\). Thus
\(\pi^*\) identifies \(\mathbb C(X)\) with \(\mathbb C(t)\).
Every biholomorphism consequently has the exact equation
\(\pi F=h\pi\) for a Möbius map \(h\), first on the regular locus
and then on the whole space.

The direct source discriminant is
\[
 -t^2\{27t^2+(4-18r)t+r^2(4r-1)\},
 \qquad \operatorname{disc}_t\{\cdots\}=16(1-3r)^3.
 \tag{EF.8}
\]
The local actual source maps at the two simple nodes are \(xy=q_i\)
after the nonzero analytic Morse factor; their invariant differential
has an annular logarithmic period. A positive node meridian therefore
has a nontrivial unipotent elliptic block. The I2 has its two-node
unipotent block. The selected toric characters retain the same base
function \(q=t-t_c\), and introduce no other base critical value.
The ordinary source/line quotients are submersions at every other finite
value. Thus the finite critical set is exactly \(\{0,t_+,t_-\}\).

At infinity use the actual minimal good cubic source root
\[
 q_\infty=z^3,\quad x_0=z^4X,\quad y_0=z^6Y,\qquad
 Y^2=X^3+r z^2X^2+zX+1.
 \tag{EF.9}
\]
Its deck is \((z,X,Y)\mapsto(\zeta_3z,\zeta_3^2X,Y)\).
On the central curve the differential has eigenvalue \(\zeta_3^2\),
so the elliptic monodromy polynomial is \(x^2+x+1\). The full cap
annulus of EF.7 below conjugates the original periods with their
functions, and retains this elliptic quotient. Its full rank-four
polynomial is \((x-1)^2(x^2+x+1)\). Infinity is therefore the unique
critical value with these non-unipotent eigenvalues. The base map fixes
infinity and is affine on \(\mathbb C_t\).

It permutes the finite set in EF.8. Its centroid is \(S_r/3\).
In the centered coordinate \(u=t-S_r/3\), the monic polynomial of
that set is exactly \(u^3+P_ru+Q_r\). An affine permutation fixes
the centroid and has the form \(u\mapsto c u\). Comparing its two
nonzero coefficients gives \(c^2=c^3=1\), hence \(c=1\).
This proves \(h=1\) on the stated source locus, without presupposing
that the automorphism preserves the elliptic source or its marked point.

\par

\medskip\noindent\textbf{EF.3. The actual point-theta period module}

\par\nobreak

On a simply connected regular base chart choose oriented elliptic
periods, set \(\tau=\omega_2/\omega_1\), and let \(g\) be the
chosen Abel integral of \(G\), divided by \(\omega_1\). Put
\(Q=e^{2\pi i\tau}\), \(P=e^{2\pi ig}\), and
\(Z=e^{2\pi iz}\). Write
\[
 f_R(Z)=\frac{\Theta(Z/R;Q)}{\Theta(Z;Q)},\qquad
 V_{\rm pt}=f_{P^a}(Z)v_M.
 \tag{EF.10}
\]
This is the coefficient of the actual point line in its constant-factor
frame. In additive coordinates \((z,w)\), with
\(V_{\rm pt}=e^{2\pi iw}\), the four literal periods, ordered as
\((\gamma,e_1,e_2,\delta)\), are
\[
 \gamma=(-bg,\ell),\qquad e_1=(1,0),\qquad
 e_2=(\tau,ag),\qquad \delta=(0,1).
 \tag{EF.11}
\]
In this frame the complete suspension coefficient is constant on each
smooth elliptic fibre: the ratio with the translation lift having the
same constant factors is an everywhere nonzero holomorphic function on
that compact curve. This proves the constant-factor assertion while
retaining its actual base value. Here \(\ell\) is the chosen logarithm
of the \textbf{entire} multiplier
of \(B_\lambda=\lambda B_{\rm sel}\) in the point frame. It includes
\(q/t^3\), the full Miller coefficient, the original scalar and every
Cartier unit. It is never replaced by a leading coefficient.

For comparison with the earlier tensor coordinates retain the exact
typed identity
\[
 V_{\rm tensor}=f_P(Z)^a f_a(p)^{-1}v_M
       =C_a(t)V_{\rm pt},\qquad
 C_a(t)=\frac{f_P(Z)^a}{f_a(p)f_{P^a}(Z)}.
 \tag{EF.12}
\]
The full Miller divisor identity makes this ratio independent of the
source point. Its complete base value and its continuation are retained
in \(\ell\). EF.12 is used in both directions at each original
Cartier chart, including coincident sections.

There are two exact changes in EF.11. First, an Abel representative
change \(g'=g+\mu+\nu\tau\), with integral \(\mu,\nu\), gives
\(P'^a=P^aQ^{a\nu}\). The identity
\(\Theta(Q^{-k}z)=(-z)^kQ^{-k(k+1)/2}\Theta(z)\) gives the full arrow
\[
 V_{\rm pt}'=(-Z/P^a)^{a\nu}
            Q^{-a\nu(a\nu+1)/2}V_{\rm pt}.
 \tag{EF.13}
\]
Its inverse uses the reciprocal at the same source point. Pulling the
new periods through that inverse gives
\[
 e_1'=e_1-a\nu\delta,\qquad
 e_2'=e_2+a\mu\delta,\qquad
 \gamma'=\gamma-b\mu e_1-b\nu e_2+c\delta,
 \tag{EF.14}
\]
where the complete integer is
\[
 c=\ell'-\ell+ab\nu(2g+\mu+\nu\tau).
 \tag{EF.15}
\]
This is an integer because these are periods of the same actual closed
torus map. Formula EF.15 specifies it using the entire logarithms; it
does not set it to zero. All signs and the source-period reindex are
fixed by the inverse of EF.13.

Second, for a change of oriented elliptic basis
\(A=\left(\begin{smallmatrix}p&q_0\\r_0&s_0\end{smallmatrix}\right)
\in SL_2(\mathbb Z)\), put \(\lambda_A=p+r_0\tau\). The directed
source and line map is
\[
 z'=z/\lambda_A,\quad \tau'=(q_0+s_0\tau)/\lambda_A,
 \quad g'=g/\lambda_A,\quad
 w'=w-\frac{a r_0g}{\lambda_A}z+c_A.
 \tag{EF.16}
\]
The full constant is defined by the actual point-divisor arrow:
\[
 e^{2\pi i c_A}=
 \frac{f_{e^{2\pi iag'}}(e^{2\pi iz'};e^{2\pi i\tau'})}
      {f_{P^a}(Z;Q)}
 \exp\!\left(2\pi i\frac{a r_0g}{\lambda_A}z\right).
 \tag{EF.17}
\]
The numerator and denominator have the same transported point divisor.
Their period factors cancel by EF.16, so the right side is a nonzero
constant on the compact elliptic fibre. This proves the displayed
Cartier arrow, and specifies its complete value by evaluation at any
ordinary point. Its logarithm branch is carried into \(\ell\).
The two old vectors \(p e_1+r_0e_2\), \(q_0e_1+s_0e_2\) go
exactly to \((1,0),(\tau',ag')\). The suspension line coefficient
changes by \(ab r_0g^2/\lambda_A\), before its retained integral
logarithm adjustment. Thus this change introduces no unrecorded mixed
period term.

Combining these two actual arrows gives the monodromy on the full
integral lattice:
\[
 M_\alpha=
 \begin{pmatrix}
 1&0&0\\
 b v_\alpha&A_\alpha&0\\
 c_\alpha&a v_\alpha^t J A_\alpha&1
 \end{pmatrix},\qquad
 J=\begin{pmatrix}0&-1\\1&0\end{pmatrix}.
 \tag{EF.18}
\]
The first and last block sizes are one. In EF.14,
\(v=(-\mu,-\nu)^t\). In the complete atlas the cocycle and its
central cochain obey the literal product equations
\[
 v_{\alpha\beta}=v_\alpha+A_\alpha v_\beta,\qquad
 c_{\alpha\beta}=c_\alpha+c_\beta+
             ab v_\alpha^t J A_\alpha v_\beta.
 \tag{EF.19}
\]
They are the multiplication equations of the displayed integral
matrices. The line scalar loop \(\delta\) is primitive, as is the
quotient homomorphism \(\psi(\gamma)=1\),
\(\psi(e_1)=\psi(e_2)=\psi(\delta)=0\): the nonzero-line cover
has fundamental fibre lattice generated by \(e_1,e_2,\delta\),
and adjoining the actual suspension adds its primitive deck generator.
EF.7 below retains these generators on the full cap annulus.

\par

\medskip\noindent\textbf{EF.4. The Abel class is nonzero and the full centralizer is computed}

\par\nobreak

The cocycle \(v_\alpha\) of EF.18 is nonzero in rational cohomology.
Here is a direct proof on the actual source. A rational coboundary would,
after multiplication by an integer and subtraction of that integral
period cochain, give a single-valued holomorphic Lie logarithm of
\(NG\) on the regular base. Multiply that integer further by six.
Then the section lies in the identity component at every special fibre.

At each I1 and at I2 the actual Tate identity coordinate of \(NG\)
is an extending nonzero unit. Its chosen analytic logarithm gives a
local extending Lie section. The difference from the purported global
logarithm is a monodromy-invariant period. For an I1 or I2 the invariant
period is the annular vanishing-cycle period; integrating the invariant
differential in \(xy=q_i\) shows that it is a holomorphic unit times
\(2\pi i\). It therefore extends in the same Lie frame. There is
no logarithmic second period in an invariant vector.

At infinity the actual formulas are particularly explicit. The chords
give
\[
 \begin{split}
 2G&=(s,-t-s/2),\\
 x(3G)&=4t^2/s^2+2t/s,\\
 y(3G)&=8t^3/s^3+6t^2/s^2+t/s-t.
 \end{split}                                                    \tag{EF.20}
\]
In EF.9 the root formal logarithm of \(3G'\) is
\(-s z^2/2+O(z^5)\). The differential in the original minimal infinity chart is
\(\omega_0=dx_0/(2y_0)=z^{-2}\omega_{\rm good}\); the affine
differential is instead \(\omega_t=z\omega_{\rm good}\). The
root logarithm transforms by \(\zeta_3^2\) under the actual deck,
so its quotient by \(z^2\) is invariant and descends holomorphically
in \(q_\infty=z^3\). Consequently division
of that root logarithm by \(z^2\) is the extending original Lie
logarithm. Integer multiples retain this extension. The difference
between the local and purported global logarithms is an invariant
elliptic period, and EF.9 has no nonzero invariant period. Thus this
difference is zero, including its full branch.

The resulting logarithm is a global holomorphic section of the actual
Lie line. In the original coordinates
\(x_0=q_\infty^2x,y_0=q_\infty^3y\), its differential frames obey
\(\omega_0=q_\infty^{-1}\omega_t\), so its dual frames obey
\(W_0=q_\infty W_t\). The Lie line is therefore \(\mathcal O(-1)\).
In coefficients, an entire function \(u_t(t)\) must satisfy
\(u_0(q_\infty)=q_\infty^{-1}u_t(q_\infty^{-1})\) holomorphic
at zero. Hence \(u_t(t)=O(t^{-1})\) at infinity, and is zero.
This would give \(NG=O\).

The explicit original central coordinate disproves that conclusion.
In the IV* Weierstrass chart EF.20 gives
\[
 x_0(3G)=4/s^2,\qquad y_0(3G)=8/s^3,\qquad
                 -x_0(3G)/y_0(3G)=-s/2\ne0.
 \tag{EF.21}
\]
The parameter \(-x_0/y_0\) is the additive coordinate on the smooth
identity component of the central cusp; direct substitution
\((x_0,y_0)=(u^2,u^3)\) gives
\(dx_0/(2y_0)=d(-1/u)\). Its integer multiples are nonzero.
Thus \(G\) is not torsion and the rational Abel cocycle is nonzero.

The elliptic monodromy centralizer is \(\mathbb Q I\). Indeed the
infinity matrix has irreducible polynomial \(x^2+x+1\), so its
rational centralizer is the quadratic field it generates. A nontrivial
I1 unipotent cannot belong to this field: its nonzero nilpotent part
would be a zero divisor. A non-scalar element of that quadratic field
generates it, so no non-scalar element commutes with the I1 matrix.

We now solve the full rank-four problem, including every
\(c_\alpha\). Modulo \(\delta\), an invariant vector with
\(\gamma\) coefficient \(z\) would give
\((A_\alpha-I)u+bzv_\alpha=0\). The nonzero Abel class forces
\(z=0\), and the absence of elliptic invariants forces \(u=0\).
Thus the invariant line of the full module is exactly
\(\mathbb Q\delta\). The dual computation using
\(v_\alpha^tJA_\alpha\) gives the invariant covector line exactly
\(\mathbb Q\psi\). A commuting matrix must preserve both and has
the form
\[
 X=\begin{pmatrix}
 s_0&0&0\\ u&A&0\\ c&w&t_0
 \end{pmatrix}.
 \tag{EF.22}
\]
Its elliptic quotient is \(A=eI\). The middle-left commutation
equation is
\((I-A_\alpha)u=b(s_0-e)v_\alpha\), hence \(s_0=e,u=0\).
The lower-middle equation is
\(w(A_\alpha-I)=a(e-t_0)v_\alpha^tJA_\alpha\).
Writing \(w=z^tJ\) and using \(A_\alpha^tJA_\alpha=J\) turns
this into the same Abel coboundary equation. Therefore
\(t_0=e,w=0\). The remaining equations allow precisely
\[
                 X=eI+c\delta\psi.
 \tag{EF.23}
\]
They include the full central cochain, because this matrix commutes
with every EF.18 for its actual value of \(c_\alpha\). Integral
unimodularity gives \(e=\pm1\), \(c\in\mathbb Z\).

A fibre biholomorphism is affine: its lifted derivative is periodic,
so its entries are holomorphic functions on the compact torus and are
constant. Its real linear map is complex linear. Subtracting
\(eI\) from EF.23 leaves the real rank-one map \(c\delta\psi\)
when \(c\ne0\); a complex-linear nonzero map has even real rank.
Consequently every actual holomorphic endpoint has linear part exactly
\(+I\) or \(-I\).

\par

\medskip\noindent\textbf{EF.5. The completed fixed-curve obstruction excludes the negative part}

\par\nobreak

Take the actual cap-completed automorphism with linear part \(-I\)
whose possibility remains in EF.4. On each regular two-torus it has
the affine form \(x\mapsto-x+a(t)\). Multiplication by two on its
integral lattice has degree \(2^4=16\), so the actual fixed set over
the regular base has exactly sixteen points in each fibre. Its vertical
derivative is \(-2I\); the analytic implicit-function theorem makes
this fixed set a finite étale curve there.

The global equalizer of the two actual holomorphic maps \(F\) and
identity is a closed analytic subset of the compact completed \(X\).
Choose an irreducible component containing one of these regular fixed
points. Its local dimension is one, hence its entire irreducible
component is a compact analytic curve. This is the proper analytic
closure of that horizontal branch, with all source, toric and cap
limits retained. Its map to the base has positive degree. Equivalently,
its intersection with a generic smooth torus fibre is a positive sum
of the transverse complex intersection multiplicities.

But \(H^4(X;\mathbb Z)=0\), so the Poincaré dual of this curve is
zero and the same generic intersection is zero. One may also pair its
normalization with the exact class \(\pi^*[\omega_{\mathbb P^1}]\)
in \(H^2(X)=0\): its positive degree gives a positive integral and
Stokes gives zero. This contradiction concerns the actual global
completed map. It does not use a local inverse whose cap extension was
only postulated. Thus every holomorphic endpoint has linear part \(+I\).

\par

\medskip\noindent\textbf{EF.6. The identity-linear endpoint factors through the actual scalar}

\par\nobreak

On the regular fibres an identity-linear map is an actual translation.
It preserves the primitive \(\psi\) and lifts locally to the cyclic
nonzero-line cover. It commutes with its generator \(B_\lambda\).
Its lifted affine map is equivariant for the full physical \(\mathbb
C^*\), and its source projection is an elliptic translation.

At an ordinary critical disc, this cover is the actual quotient of
\(M_a^\times\) over the resolved source by the invertible contraction
\(B_\lambda\). The base disc is simply connected and its total
fundamental group is obtained from the generic fibre by its vanishing
cycles. Those cycles lie in the kernel of \(\psi\), as the displayed
cover shows. Thus the preserved fibre homomorphism extends and the map
lifts holomorphically across that disc. Equivariance and source
projection extend from the dense regular locus. Evaluation on the
chosen nonzero vector over \(O\) gives the actual extending source
section there.

At infinity take the literal fibre product with \(z^m=q_\infty\)
and its normalization. EF.7 identifies it with the root quotient by
\(B_\lambda\), because the kernel of the full cap generator has
\(A_\lambda^m=B_\lambda^n\), \(n=\pm1\). The canonical lift of
the base-preserving map fixes \(z\). Its nonzero-line lift has a
commutator \(B_\lambda^k\) with the actual \(A_\lambda\).
Iteration \(m\) times, and its already proved commutation with
\(B_\lambda^n\), give \(B_\lambda^{mk}=1\). Strict contraction
gives infinite order, hence \(k=0\). This retains the entire finite
character and both signed norms. Source evaluation therefore extends
on the genuine root and on its source graph, not merely on a generic
elliptic fibre.

At the selected I1 choose a source section in an open toric stratum.
In a chart with the actual equation \(q=z_0z_1z_2\), its image under
the actual endpoint is a section over \(q\). The vanishing orders of
these three holomorphic coordinates are nonnegative and sum to one;
the source and line torus characters consequently have integral finite
orders. Write the complete periods there as
\[
 h:(Z,V,q)\mapsto(QZ,\alpha V,q),\qquad
 B:(Z,V,q)\mapsto(gZ,q\eta V,q),\quad Q=\kappa q,
 \tag{EF.24}
\]
where \(\kappa,\alpha,g,\eta\) are the actual nonzero holomorphic
source/Cartier factors. Reindex by the literal \(h^rB^k\) to remove
the two integral orders; retain its full unit factors. Evaluation of
the endpoint and its inverse then gives extending nonzero units
\(K(q),C(q)\). On every fan character their exact map is
\[
 Z^uV^vq^h\longmapsto K(q)^uC(q)^vZ^uV^vq^h.
 \tag{EF.25}
\]
The inverses use the inverse units at the actual inverse point. This
proves source-section extension across the selected I1 and the full
bounded line arrow on every prescribed cone; the integral \(B\)
reindex will be used below, not deleted.

Away from the selected point the local lifts differ by integral
suspension powers. The orbifold base is \(\mathbb P^1\setminus
\{t_c\}\) with the retained order-\(m\) cap; its fundamental group
is the finite cyclic cap group. A homomorphism to \(\mathbb Z\) is
zero. The integer cocycle of lifts is therefore removed by its actual
deck powers. The resulting source section is global on the minimal
elliptic surface.

For clarity the marked source generator is identified integrally.
The explicit rational chart
\[
 U=y-t-x/2,\quad V=y+t+x/2,\quad UV=x^2(x-s)
 \tag{EF.26}
\]
shows the source is rational. Its resolved elliptic surface has Euler
number \(8+2+1+1=12\), hence the unimodular rank-ten rational-surface
intersection lattice. The fibre and zero section give the unimodular
rank-two summand. The actual IV* and I2 component lattices are E6 and
A1, with determinant six. Restriction of a divisor to the generic
fibre identifies the quotient by these vertical classes with the
Mordell-\/-Weil group, of free rank one. The actual component and
intersection arrows can be read directly. At finite base values \(G\)
is in the affine chart and \(O\) is at infinity. At the A1 point of
EF.26 the tangent cone is \(UV=-s x^2\). The section \(G\) has direction
\([x:U:V]=[0:0:1]\), whereas the strict central fibre has equation
\(V-U-x=0\). Thus the lift of \(G\) meets the exceptional nonidentity
A1 component once, off its intersection with the identity component.
The two sections remain disjoint. At infinity, EF.9 has the distinct
fixed sections \((0,1)\) and \(O\) on the good central elliptic curve.
The actual quotient charts at each have weights \((2,1)\), resolved
by \((2/3,1/3),(1/3,2/3)\), so their section arms remain disjoint.
Label the nonidentity outer component met by \(G\) as \(C_0\), and
order the other E6 components by the edges
\((0,1),(1,2),(2,3),(3,4),(2,5)\). The arm through \(O\) has outer
component \(F-C_0-2C_1-3C_2-2C_3-C_4-2C_5\).
These are the literal component labels of F41 SC.5 and the integral
source matrix in J.4, carried by the displayed root charts.

The positive Cartan matrix and actual correction column are
\[
 A=\begin{pmatrix}
 2&-1&0&0&0&0\\-1&2&-1&0&0&0\\
 0&-1&2&-1&0&-1\\0&0&-1&2&-1&0\\
 0&0&0&-1&2&0\\0&0&-1&0&0&2
 \end{pmatrix},\quad
 A^{-1}e_0=(4/3,5/3,2,4/3,2/3,1)^t.
 \tag{EF.26a}
\]
The section intersection vector is exactly \(e_0\); at A1 it is
\((1)\), with Cartan matrix \((2)\). Hence the two local corrections
are \(e_0^tA^{-1}e_0=4/3\) and \(1/2\). The displayed disjoint
section charts prove \(G\cdot O=0\), so the height is
\(2-4/3-1/2=1/6\). The normal line of either smooth section is the
actual Lie line, so \(O^2=G^2=-1\). In the order
\((O,F,C_0,\ldots,C_5,A_1,G)\), put \(OF=GF=1\), the component block
\(-A\oplus(-2)\), \(G C_0=G A_1=1\), and all other unspecified
pairings zero. This full integral Gram matrix has determinant \(-1\),
as checked directly. The ten actual classes therefore give the whole
unimodular source lattice, and its quotient by the first nine is the
free group generated by \(G\). The height determinant gives the same
primitive conclusion explicitly. Writing \(\iota\) for its free index and
\(t_{\rm MW}\) for the torsion order, the integral determinant
equations are
\[
 1=6h(G_{\rm prim})/t_{\rm MW}^2,\qquad
 \iota^2h(G_{\rm prim})=1/6,\qquad
                         \iota^2t_{\rm MW}^2=1.
 \tag{EF.27}
\]
Both positive integers are one. Thus the actual section above is
\(jG\) for an integer \(j\), with no additional torsion section.

Keep the full Miller functions of K96. Their typed one-step map is
\[
 B_{a,1}(p)=f_2(p)^a\frac{f_a(p-G)}{f_a(p)},\quad
 U_G=B_{a,1}^{-1},\quad
 U_G(p,v)=(p+G,B_{a,1}(p+G)^{-1}v).
 \tag{EF.28}
\]
All inverses therefore evaluate at the actual target source point.
The chronological product \(B_{a,b}\) is \(U_G^{-b}\), and
\[
 B_{\rm sel}=(q/t^3)B_{a,b},\qquad
              B_{\rm sel}^{-1}=(t^3/q)U_G^b.
 \tag{EF.29}
\]
An endpoint lift over \(+jG\) is \(H(t)U_G^j\): the ratio is an
invertible function on each compact elliptic fibre and is a base
function. Its meromorphic extension at the finite singular fibres and
root is in the actual Cartier module charts. The selected argument
EF.24 gives its finite integer order there. Hence \(H\) is a rational
base function.

The actual generic-component valuations are
\[
\begin{array}{c|ccc}
 &t=0&z=0\ (q_\infty=z^m)&q=0\\ \hline
 U_G&-3/b&2m/b&0\\
 H&3j/b&-2mj/b&-k
\end{array}                                                     \tag{EF.30}
\]
Here are the exact full-arrow antecedents of these orders. For M2,
EX.1 gives \(B_{2,1}(p)=f_2(p)f_2(p-G)\). On the identity I2
component \(f_2\) has order zero; on its other component the actual
chart \(x=t\alpha,y=t\beta\) gives
\[
 f_2=t\,\frac{\beta-1-\alpha/2}{-s+t\alpha}.
 \tag{EF.30a}
\]
Both displayed generic coefficients are nonzero. The actual \(-G\)
map exchanges the components, so \(B_{2,1}\) has order one on each,
and its inverse has order \(-1=-3/b\).

For M6, the complete K94 EU.2-\/-3 formula is
\[
 \begin{gathered}
 U=y-t-x/2,\quad V=y+t+x/2,\quad
 f_2=U/(x-s),\\
 f_3=\frac{s xU}{sV-(s+2t)x},\qquad
 f_6=f_3^2 f_{2,3G},\qquad
 B_{6,1}=f_2^6\,\frac{f_6(p-G)}{f_6(p)},\\
 \left.\frac{B_{6,1}}{t^3}\right|_{C_{\rm id}}
             =\frac{8((1-s)w-s)}{s^3},\qquad
 x|_{C_{\rm id}}=s+w^2-w .
 \end{gathered}                                                  \tag{EF.30b}
\]
The last coefficient is not identically zero for any \(s\ne0\).
For completeness its full factors are
\[
 \begin{aligned}
 f_2(p)&\to x_0/w,&f_3(p)&\to x_0,\\
 f_3(p-G)/t&\to2w^3/[x_0(w-s)],&
 f_{2,3G}(p)&\to x_0(1-s)/[(1-s)w-s],\\
 f_{2,3G}(p-G)/t&\to
                  2(1-s)(w-s)^2/(s^3x_0).
 \end{aligned}                                                  \tag{EF.30c}
\]
Their literal product is EF.30b. The value \(s=1\), which remains
in EF.3, requires specialization of the full map before taking this
limit. Do that in the exact source chart
\[
 u=w-1,\quad v=w+h,\quad x=1+uv,\quad
 t=xh/2,\quad y=x(u+v)/2 .
 \tag{EF.30e}
\]
It satisfies the original equation with \(r=-3/4\). The complete
inverse translation is \(x'=vh,\ y'=t-(v-1/2)x'\). Set \(s=1\)
in the actual tangent at \(3G\) first. Its slope and doubled point are
\[
 \lambda_{3G}=\frac{24t^3+24t^2+3t-1}{2t(4t+3)},\quad
 x(6G)=
 \frac{(4t^2+2t-1)(16t^4+24t^3+16t^2+4t-1)}
      {4t^2(4t+3)^2}.
 \tag{EF.30f}
\]
Substitution of these full functions into the two tangent quotients,
and only then taking \(h\to0\), gives, with \(x_0=w^2-w+1\),
\[
 f_2\longrightarrow x_0/w,\qquad
 \frac{f_3(p-G)}{h f_3(p)}
       \longrightarrow\frac{w^3}{(w-1)x_0},\qquad
 \frac{f_{2,3G}(p-G)}{h f_{2,3G}(p)}
       \longrightarrow-\frac{(w-1)^2}{x_0}.
 \tag{EF.30g}
\]
Since \(t/h\to x_0/2\), their full product gives
\(B_{6,1}/t^3\to-8\) after this specialization. It agrees with EF.30b
but was computed before any generic-\(s\) leading-term cancellation.
The same exact doubling formula gives
\(\tau(6G)/t\to-6\), where \(\tau=-x/y\) at \(O\). On this actual
point/O chart keep the two Cartier equations
\(\tau-\tau(6G)(t)\) and \(\tau\), with both their quotient arrows.
They define the point-line module also at their coincidence. The
Miller identity carries these two modules in the Hom map. No
horizontal divisor is removed by evaluation at its generic point.
After the full order-three factor is removed, all vertical valuations
are zero and the Hom divisor is zero. Normality of the actual smooth
resolved source extends the map and its reciprocal at the remaining
points, including this coincidence.
The section \(6G\) and \(O\) meet the same I2 component, so \(M_6\)
has degree zero on both. The Hom divisor has zero component
intersections. The kernel of the full I2 matrix
\(\left(\begin{smallmatrix}-2&2\\2&-2\end{smallmatrix}\right)\)
is the primitive vector \((1,1)\). Thus the order three in EF.30b
is the order on both components, giving the first M6 entry.

For the root column the exact EX.4 and EX.13 scalings are
\[
 \begin{array}{c|c|c}
 (a,b;m)&B_{a,1,\rm old}&\text{generic central root coefficient}\\ \hline
 (2,3;3l)&z^{-2l}B'_{2,1}&
  B'_{2,1}=f(p)f(p-G'),\quad f=(Y-1)/X\\
 (6,1;3l+2)&z^{-6(l+1)}B'_6&
                         B'_6/z^2\longrightarrow4 .
 \end{array}                                                    \tag{EF.30d}
\]
The first coefficient is nonzero at a generic point of the smooth
root elliptic fibre. The second is the complete chord limit in K96
EX.13, in the line graph EX.10-\/-12. Hence the inverse orders are
\(2l=2m/3\) and \(6(l+1)-2=2m\), respectively. The remaining
root graph conditions are the literal two-section coefficients
EF.32 and the two correction sums EF.33 below, that is K96
EX.11-\/-14; their inverses and characters are carried into EF.35-\/-38.
They are not inferred from this generic order.

The resulting base order of \(H\) at infinity is \(-2j/b\).
The third column is exactly the integer \(B^k\) reindex of EF.24.
At a finite irreducible fibre the numerator and denominator
\(f_a(p-G)/f_a(p)\) have the same vertical order; each other
factor is generically nonzero. This is the component-DVR argument
of K94 EU.2 and K96 EX.1. There are no other fibre divisors.
All graph-correction
monomials are still part of the root modules; EF.30 reads their generic
component orders and does not discard the remaining extension
conditions. The degree of the actual rational divisor is zero:
\(3j/b-2j/b-k=0\), so \(j=bk\). Consequently
\[
 H=c(t^3/q)^k,\qquad H U_G^{bk}=c B_{\rm sel}^{-k},\quad
                                      c\in\mathbb C^*.
 \tag{EF.31}
\]
Since \(B_\lambda=\lambda B_{\rm sel}\) is the actual quotient
deck, EF.31 acts exactly by the physical scalar \(c\lambda^k\).
The equality holds first on the common dense chart and hence on the
entire completed endpoint; EF.7 shows explicitly that this scalar
commutes with all original root, character and annular arrows. This
also retains any additional extension restrictions on the meromorphic
representative in EF.30: the final map is the actual extending scalar,
with the same graph and finite characters.

\par

\medskip\noindent\textbf{EF.7. Both full cap and annular arrows, and the endpoint consequence}

\par\nobreak

Here are the root maps used above, in their original direction. For
\((a,b;m)=(6,1;3l+2)\) put \(k=2l+1\):
\[
 q_\infty=z^m,\quad x_0=z^{2k}X,\quad y_0=z^{3k}Y,\quad
 Y^2=X^3+r z^{2l+2}X^2+z^{l+2}X+z^2.
 \tag{EF.32}
\]
The deck is
\(D(z,X,Y)=(\zeta_mz,\zeta_m^{-2k}X,\zeta_m^{-3k}Y)\).
The three A2 charts of EX.6-\/-8 and their explicit chord inverses give
the actual source resolution and translation. The exact division
section is
\(R=2G'-[m]^{-1}((2m-1)G')\), with \(mR=G'\).
The graph of \(\tau'=z^k\tau_0\), including every ray
\((j,1)\), \(1\le j\le k\), retains separately the section
equations \(\tau'\) and \(\tau'-z^kL(s,z^m)\). Their actual
coefficient ratio is \(u/(u-L)\), or \(1/(1-vL)\) on the inverse
graph chart. It gives \(\beta^*N_6=\pi^*M_6\). The full multiplier
is \(B'_{\rm sel}=h z^{-2}B'_6\),
\(h=1-t_cz^m\), with central coefficient \(4h\).

For \((a,b;m)=(2,3;3l)\) the exact maps are
\[
 q_\infty=z^{3l},\quad x_0=z^{4l}X,\quad y_0=z^{6l}Y,
 \qquad Y^2=X^3+r z^{2l}X^2+z^lX+1.
 \tag{EF.33}
\]
The deck has weights \((1,-4l,-6l)\). Put \(u=z^l\); the
two point-line quotient rays are \((2/3,1/3),(1/3,2/3)\) in
\(\mathbb Z^2+\mathbb Z(2,1)/3\). Their full support difference
is \(\min(v_w,2l v_z)\), at both marked sections. With their
actual monomial generators this gives
\(\beta^*N_2=\pi^*M_2+E_P-E_O\). The source division is the
small section \(R=[m]^{-1}(3G')\). Its translation preserves the
ideals \((w,u^j),(u,w^j)\), \(j=1,2\), and their inverse charts,
so the full line correction is carried. The complete coefficient is
\(B'_{\rm sel}=hB'_{2,3}\); the literal three-orbit Miller product
has central value four. These are the EX.3-\/-5 maps, including the
nontrivial central degree-zero line \(N_2\).

For either branch the actual chord-over-vertical function
\[
 g_{P,R}(p)=\ell_{P,R}(p)/(x(p)-x(P+R))
 \tag{EF.34}
\]
has divisor \((P)+(R)-(P+R)-(O)\) and is the typed lift covering
\(p\mapsto p-R\). The normalized starting lift is
\(F=\tau(P)g_{P,R}\) for \(P=6G'\), and
\(F=\tau(R)g_{P,R}\) for \(P=2G'\). The actual near-O expansion
gives central coefficient one. The formal-coordinate weight and the
chord weight cancel, so it commutes with the full \(D\). The base
unit \(\beta=B'_{\rm sel}/F^m\) has central value four; its
reference-chosen analytic \(m\)-th root is \(D\)-invariant. Thus
\[
 \widetilde F=\beta^{1/m}F,\quad
 A_0=\zeta_m^d\widetilde F^{\,n}D,\quad
 B_\lambda=\lambda B'_{\rm sel},\quad
 A_\lambda=R_\lambda^n A_0,\quad R_\lambda^m=\lambda,
 \tag{EF.35}
\]
and literally \(A_\lambda^m=B_\lambda^n\),
\([A_\lambda,B_\lambda]=1\). These equations retain \(d\),
both signs of \(n\), every source argument of a chord coefficient,
and the selected-height function \(h\).

On the two fixed annular sectors set
\[
 s_\alpha=\log_\alpha q_\infty/(2\pi i),\quad
 z_\alpha=e^{2\pi i s_\alpha/m},\quad
 c_\lambda=n\log R_\lambda/(2\pi i)+d/m,\quad
 f_\alpha=\operatorname{Log}_\alpha\widetilde F.
\]
The full ordinary-to-root map is
\[
 \Phi_\alpha=\operatorname{Exp}
  (n s_\alpha f_\alpha+s_\alpha c_\lambda\delta+h_\alpha)
                                                     I_{z_\alpha}.
 \tag{EF.36}
\]
When \(s_\beta=s_\alpha+j\) and
\(f_\beta=D_*^jf_\alpha+\Pi_{\alpha\beta}\), its actual
reference cochain is
\[
 h_\beta-D_*^jh_\alpha
 =-n s_\alpha\Pi_{\alpha\beta}
       +k_{\alpha\beta}\operatorname{Log}_\beta B_\lambda
       +\Lambda_{\alpha\beta}.
 \tag{EF.37}
\]
The columns \(\Pi,\Lambda\) and integer \(k_{\alpha\beta}\)
are precisely the transported complete reference columns. Subtracting
the target-sheet logarithm \(jn f_\beta+jc_\lambda\delta\)
gives the literal squares
\[
 \Phi_\beta=A_\lambda^jB_\lambda^{k_{\alpha\beta}}\Phi_\alpha,
 \qquad \Phi_\alpha B_\lambda=B_\lambda\Phi_\alpha.
 \tag{EF.38}
\]
Their inverses use the same columns with inverse group elements. The
full source displacement is the signed \(-n s_\alpha\) Abel
integral of \(R\), and the scalar is
\(R_\lambda^{n s_\alpha}e^{2\pi i d s_\alpha/m}\).
This is the complete cap-period extension used in EF.3 and EF.6.

Physical multiplication by \(c\) commutes with EF.34-\/-38 and acts
on each cusp character by \(Z^uV^vq^h\mapsto c^vZ^uV^vq^h\).
It therefore extends on the same root and fan with the prescribed
reference. A path in \(\mathbb C^*\) from one to \(c\) is an
actual isotopy of these completed maps. Conversely EF.2-\/-6 show every
holomorphic automorphism is one of them. Its kernel is trivial: on a
regular source fibre an equality with a suspension power would give
an integral multiple of \(bG\) equal to zero as a global source
section, contradicted by EF.21; then the scalar itself is one.

For two actual holomorphic endpoint identifications of one and the
same cut family, their ratio is such an automorphism of the common
completed endpoint. The physical return uses the inverse endpoint;
the scalar path therefore supplies an actual isotopy between those
inverse-endpoint returns, with the specified source, fan and cap
markings carried. This conclusion concerns endpoints that have been
constructed. It assigns neither a new shorter scalar endpoint nor a
return to an unconstructed source loop.

\par

\medskip\noindent\textbf{EF.8. Exact verification and scope ledger}

\par\nobreak

The companion program
\path{code/part3/e6_fixed_source_endpoints.py} passes 1,401 exact
checks. It records the discriminant and affine-rigidity
polynomials, the literal 2G/3G chord identities and IV* additive limit,
the two point-theta shift directions, integral monodromy multiplication,
the complete centralizer block equations, both root coordinate/deck
weight identities, full cap norm exponents, and the rational divisor
factorization. These computations support the displayed maps. The
extension of the Lie logarithm, proper analytic fixed-curve closure,
cover-lifting argument and actual graph/annular extension are the
geometric proofs above, not conclusions drawn from a finite sample.

The original proof sources retained as inputs are K96 EX.1-\/-6, K97
DZ6.6, the marked-node source note, and the explicitly cited F35/F46
topology maps. K79 is a methodological comparison; none of its OS49
monodromy matrices, source sections or divisor orders is substituted
for EF.1 or EF.18. The full independent and root audits have now passed. Their exact source locks and the retained prior audit stages are recorded below.

\par\medskip\noindent\textbf{Final source, exact certificates and audits.}\par
All eight proof sections above are exported in full from
\path{notes/part3_e6_fixed_source_endpoints_working.md}.
The primary exact certificate passes 1401 checks. The root's complete
EF.1--8 read and independent certificate pass 259 checks, including
a separate specialization of the full Miller arrow at \(s=1\)
with the original \(x\) coordinate held fixed. The independent
period-module and centralizer proof passes 2126 checks and includes
116 exact integral elliptic basis changes. Its reviewed EF.3--4
text is byte-identical to the corresponding sections of this final
source, as the saved source bridge verifies.

The complete audit proofs and records are
\path{notes/part3_e6_fixed_source_endpoint_root_peer.md},
\path{notes/part3_e6_fixed_source_period_peer.md},
\path{records/root_audit_e6_fixed_source_endpoints.json}, and
\path{records/checks_audit_e6_fixed_source_period_module.json}.
The bridge is
\path{records/e6_fixed_source_period_peer_source_bridge.json}.
The earlier generic-leading-term stage, its original source and
certificate, and the subsequent literal \(s=1\) source specialization
are retained under
\path{records/attempts/e6_fixed_source_endpoints_before_s1_specialization/}.
The export manifest locks every source, proof, certificate and audit
at \path{records/k109_export_manifest.json}.

% End expanded verified cp36 input

\end{proof}


\subsection{All holomorphic scalar endpoints on the generic fixed E6 source}
\begin{theorem}[Source result P3K.110 under continuation review]\label{night:P3K.110}
\StatementRecord{night:P3K.110}{open}{Historical events: GPT-6 Astra (OpenAI), ChatGPT Pro mode, effort Pro, S6\_06\_BLIND, 16--18 November 2026: targeted proof check. Also checked in this writing. Continuation mathematical and source checks: GPT-6 Astra (OpenAI), ChatGPT Pro mode, effort Pro, S6\_10\_PAPERS continuation, 19 November 2026: computation or formal certificate; targeted proof check; used and rechecked in this run. Continuation metadata events: GPT-6 Astra (OpenAI), ChatGPT Pro mode, effort Pro, S6\_10\_PAPERS continuation, 19 November 2026: used and rechecked in this run; source-row intake; metadata reconciliation only; no new mathematical check. GPT-6 Astra (OpenAI), ChatGPT Pro mode, effort Pro, S6\_10\_PAPERS continuation, 19 November 2026: used and rechecked in this run; explicitly scoped metadata reconciliation; metadata reconciliation only; no new mathematical check. GPT-6 Astra (OpenAI), ChatGPT Pro mode, effort Pro, S6\_10\_PAPERS continuation, 19 November 2026: used and rechecked in this run; receiving source placement; metadata reconciliation only; no new mathematical check. \emph{Pre-continuation recorded proof scope.} Historical event: GPT-6 Astra (OpenAI), ChatGPT Pro mode, effort Pro, S6\_06\_BLIND, 16--18 November 2026: targeted proof check, recorded producing-run re-derivation. Also checked in this writing by GPT-6 Astra (OpenAI), ChatGPT Pro mode, effort Pro, S6\_10\_PAPERS continuation, 19 November 2026: computation or formal certificate for the explicitly listed algebraic substeps in the event ledger; full continuation rederivation and renewed credit comparison remain pending.}{\cref{night:P3K.109,night:P3K.96}; the original source, point line, roots, characters, reference, selected fan and scalar range stated here}
Retain an original marked source in the precise locus
\(\mathcal U\) of P3K.109, its selected height \(q=t-t_c\),
every prescribed lattice-basis fan, and either full cap progression
\[
 (a,b;m)=(2,3;3l),\quad l\ge1,
 \qquad (a,b;m)=(6,1;3l+2),\quad l\ge0.
\]
Use the actual point line \(M_a=\mathcal O(aG-O)\), suspension
by \(-bG\), either norm \(n=\pm1\), the prescribed finite
character \(\zeta_m^d\), and the complete reference of P3K.96.
On the scalar parameter cover retain
\[
 \lambda_\theta=e^{2\pi i\theta}\lambda_*,\qquad
 R_\theta=e^{2\pi i\theta/m}R_*,\qquad R_*^m=\lambda_*,
\]
with the full restored reference in SEF.3 below.

Every holomorphic endpoint has equal original scalar values at its
two ends, hence has integral scalar winding \(J\). The complete
groups of possible windings are
\[
 \begin{array}{c|c}
 (a,b;m)&\text{scalar winding group}\\\hline
 (2,3;3l)&l\mathbb Z\\
 (6,1;3l+2)&m\mathbb Z.
 \end{array}
\]
All these endpoints have linear part \(+I\) on the ordinary
torus fibres. For a fixed winding the endpoint ambiguity is exactly
the connected physical \(\mathbb C^*\) action of P3K.109.

For the good M2 cap, let
\[
 f_2(p)=\frac{y-t-x/2}{x-(1/4-r)},\qquad
 B_{2,1}(p)=f_2(p)f_2(p-G),\qquad U_G=B_{2,1}^{-1},
\]
where the inverse evaluates its coefficient at the actual translated
source point. Put \(S=tU_G\). On its genuine good root retain
the full map \(U_{G,\rm old}=z^{2l}U'_G\). The primitive
positive winding \(J=l\) has the following actual forward endpoint:
\[
 E_R^{\rm ord}=(R^{-l}S)^n,\qquad
 E_R^{\rm cap}=(R^{-l}U'_G)^n.
\]
Both signs use the actual inverse source/line arrow when required.
Their three chronological composites at
\(R,\zeta_3R,\zeta_3^2R\) are exactly
\[
 E_2^{\rm ord}E_1^{\rm ord}E_0^{\rm ord}
       =q^nB_\lambda^{-n},\qquad
 E_2^{\rm cap}E_1^{\rm cap}E_0^{\rm cap}
       =h^nB_\lambda^{-n},\quad h=q_\infty q.
\]
Thus their quotient maps give the canonical full forward
\(q^n/h^n\) closing. For the reduced M6 cap the primitive winding
\(J=m\) is precisely that canonical forward closing. Negative
windings use the chronological inverses. Every source, Cartier,
root, character, fan and annular multiplier is retained by the
proof below. The physical return uses the inverse endpoint;
no return class is assigned or divided by a covering degree here.
\end{theorem}
\begin{proof}[Supplied source argument; continuation review in progress]

% Expanded verified cp36 input: notes/part3_e6_fixed_source_scalar_endpoints_proof.tex
\par\noindent
The full constructive stage SEF.1--6 is retained first, including
its then-open local quadratic norm equation. The complete global
NE.1--4 proof follows and resolves that equation using the actual
selected lift and the retained whole annular deck. No mathematical
paragraph of either stage is omitted.

\par

\medskip\noindent\textbf{SEF.1. Fixed objects and the restored full-reference scalar family}

\par\nobreak

Retain the source locus EF.3, the chosen original selected I1 at \(t_c\),
the actual height \(q=t-t_c\), point line \(M_a=\mathcal O(aG-O)\),
all fan labels, and the complete reference of K96. Put
\[
 (a,b;m)=(2,3;3l),\quad l\ge1,\qquad n=\pm1,\quad d\in\mathbb Z/m.
 \tag{SEF.1}
\]
The source variable is \(t\). The scalar root is denoted \(R\);
it is not a source coordinate. Let \(\lambda_*=R_*^m\ne0\) be
the fixed small original scalar at the reference. On the complex
parameter cover set
\[
 \lambda_\theta=e^{2\pi i\theta}\lambda_*,\quad
 R_\theta=e^{2\pi i\theta/m}R_*,\quad
 B_\theta=e^{2\pi i\theta}B_*,\quad
 A_\theta=e^{2\pi i n\theta/m}A_* .
 \tag{SEF.2}
\]
These are the complete source/line operators, with
\(A_\theta^m=B_\theta^n\) and the prescribed finite character \(d\).

Use the actual reference sectors of EF.36-\/-38, with
\(s_\alpha=\log_\alpha q_\infty/(2\pi i)\). On their labelled
suspension lifts let \(J_\alpha\) be the retained integer representative;
on an overlap \(J_\beta-J_\alpha=j_{\alpha\beta}\) is exactly its
reference suspension reindex. The whole scalar family is
\[
 \Phi_{\theta,\alpha}
 =e^{2\pi i n\theta s_\alpha/m}\,
 B_\theta^{J_\alpha}
       \bigl(B_*^{-J_\alpha}\Phi_{*,\alpha}\bigr).
 \tag{SEF.3}
\]
This is the reference restoration used in SI3.5, with the present
order \(m\). It is a formula on the labelled reference groupoid:
the integer is retained when changing a sector or a suspension label.
Its source and line period columns \(\Pi,\Lambda\) are the original
functions, not their central values. Substitution in EF.38 gives
\[
 \Phi_{\theta,\beta}
  =A_\theta^{\,k}B_\theta^{j_{\alpha\beta}}\Phi_{\theta,\alpha},
 \qquad
 \Phi_{\theta,\alpha}B_\theta=B_\theta\Phi_{\theta,\alpha}.
 \tag{SEF.4}
\]
Here \(s_\beta=s_\alpha+k\); the source part of \(A_\theta\) and every
reference column are unchanged. The extra factors on the two sides
are respectively \(e^{2\pi i n\theta k/m}\) and
\(e^{2\pi i\theta j_{\alpha\beta}}\), exactly as SEF.2-\/-3 specify.
The inverse maps use the reciprocal entire factors and the same
integer labels. At an integral scalar winding \(J\),
\[
 B_{\theta+J}=B_\theta,\qquad
 \Phi_{\theta+J,\alpha}
       =q_\infty^{\,nJ/m}\Phi_{\theta,\alpha}
       =z^{nJ}\Phi_{\theta,\alpha}.
 \tag{SEF.5}
\]
The last equality is in the actual \(m\)-root coordinate, with its
chosen sector branch. No reference suspension factor remains
unrecorded: its factor is \(e^{2\pi iJJ_\alpha}=1\).

\par

\medskip\noindent\textbf{SEF.2. The complete finite-source map \(S=tU_G\)}

\par\nobreak

Keep the rational point-to-tensor map \(r_2=f_2^{-1}\), where
\[
 f_2(p)=\frac{y-t-x/2}{x-s},\quad s=\tfrac14-r,\qquad
 B_{2,1}(p)=f_2(p)f_2(p-G),\qquad U_G=B_{2,1}^{-1}.
 \tag{SEF.6}
\]
Thus
\[
 U_G(p,v_M)=\bigl(p+G,B_{2,1}(p+G)^{-1}v_M\bigr),\qquad
                         S=tU_G.
 \tag{SEF.7}
\]
This retains the actual target source argument of the inverse.
The generic divisor identity is exactly the one for the two point-line
modules; all numerator and denominator section generators remain
through their intersections.

There is a direct full I2 check. Put
\[
 u=w-1,\quad v=w+h,\quad x=s+uv,\quad
 t=xh/2,\quad y=x(u+v)/2.
 \tag{SEF.8}
\]
The exact inverse chord has
\(x(p-G)=vh,\ y(p-G)=t-(v-1/2)vh\). It gives, before taking a limit,
\[
                    \frac{B_{2,1}(p)}t=\frac{2v}{s-vh}.
 \tag{SEF.9}
\]
Its generic identity-component coefficient is \(2w/s\ne0\).
On the other component \(f_2=t(\beta-1-\alpha/2)/(-s+t\alpha)\);
the source map exchanges the two components. Thus \(B_{2,1}\)
has order one on both. The section \(2G\) lies on the identity
component, so the point line has zero degree on both components.
After division by \(t\) the full Hom divisor is zero. On the smooth
resolved source the actual map and its reciprocal therefore extend
through both components and their intersection, in the prescribed
point/O modules. At every other finite fibre the same Hom map has
zero vertical order; all horizontal factors cancel in these modules.
Consequently \(S\) is an actual invertible source/line map over the
entire finite base.

At the selected I1 the source displacement is \(+G\). In the actual
point-theta coordinates its full line coefficient is
\[
 C_S(t)=
 t\,B_{2,1}(p+G)^{-1}
       \frac{f_{P^2}(PZ;Q)}{f_{P^2}(Z;Q)}.
 \tag{SEF.10}
\]
Both the function and its reciprocal extend at the selected height.
The ratio with the constant-factor translation lift is an everywhere
nonzero function on the compact elliptic fibre, hence is constant
there; SEF.10 specifies its entire base value. The exact torus map is
\((Z,V)\mapsto(PZ,C_S V)\). Every character
\(Z^uV^vq^h\) is multiplied by \(P^uC_S^v\), and both factors are
extending nonzero units. This carries the same prescribed fan, all
source periods and both normal collars.

The chronological threefold product is literal:
\[
 B_{2,3}=U_G^{-3},\quad
 B_{\rm sel}=(q/t^3)B_{2,3},\quad
 S^3=t^3U_G^3=q B_{\rm sel}^{-1}.
 \tag{SEF.11}
\]
It uses the full Miller factors at their successive source points.

\par

\medskip\noindent\textbf{SEF.3. The complete good-root map and finite character square}

\par\nobreak

Use precisely the root and line of K96:
\[
 q_\infty=z^{3l},\quad
 x=z^{-2l}X,\quad y=z^{-3l}Y,\quad t=z^{-3l},\qquad
 Y^2=X^3+r z^{2l}X^2+z^lX+1.
 \tag{SEF.12}
\]
The point is \(G'=(0,1)\), the point line is \(N_2=\mathcal O(2G'-O)\),
and
\[
 D(z,X,Y)=(\zeta_m z,\zeta_m^{-4l}X,Y),\qquad
 \phi(p)=\frac{Y-1-z^lX/2}{X-sz^{2l}},\qquad
 B'_{2,1}(p)=\phi(p)\phi(p-G').
 \tag{SEF.13}
\]
The full coefficient identities in the rational point frames are
\[
 f_{2,\rm old}=z^{-l}\phi,\qquad
 B_{2,1,\rm old}=z^{-2l}B'_{2,1},\qquad
 U_{G,\rm old}=z^{2l}U'_G,\quad U'_G=(B'_{2,1})^{-1}.
 \tag{SEF.14}
\]
On the central elliptic curve, \(\phi=(Y-1)/X\),
\(\phi(p-G')=-(Y+1)/X\), and \(B'_{2,1}=-X\) in the
canonical meromorphic point frames. This is a nonzero map of the
actual Hom line, including the points \(O,G',2G'\); its apparent
zero or pole is exactly their retained Cartier generator. The
generic central order is zero, so the map and reciprocal extend
through the whole smooth root family. This proves that \(U'_G\)
is the actual invertible lift of \(+G'\).

The root/source point-line comparison retains
\(\beta^*N_2=\pi^*M_2+E_P-E_O\), with each exceptional coefficient
\(\min(v_w,2l v_z)\), as in EX.14. SEF.14 is the equality of
rational coefficient arrows on that common graph, with the original
and translated strict-section generators and their correction
monomials. It is not an identification of their Chern classes. On
the root side its extension is the Hom map just established; on the
ordinary side it is SEF.7-\/-9. Thus both graph projections carry this
same dense arrow, with its full point factors and inverse.

The deck character is computed before taking any central value.
The numerator of \(\phi\) has weight zero, its denominator has
weight \(2l\), and
\[
 \phi(Dp)=\zeta_m^l\phi(p),\quad
 B'_{2,1}(Dp)=\zeta_m^{2l}B'_{2,1}(p),\quad
                    U'_G D=\zeta_m^{-2l}D U'_G.
 \tag{SEF.15}
\]
These equations are in the canonical point frame; the two actual
Cartier generators supply the same equality in every regular
frame. The finite character \(\zeta_m^d\) is never removed.

The small source division section \(R_{\rm div}=[m]^{-1}(3G')\)
and its full normalized line lift \(\widetilde F\) are those of
EX.15-\/-16. Both \(U'_G\) and \(\widetilde F\) have constant
coefficients in a common constant-factor frame of \(N_2\), by the
compact-fibre Hom argument used in SEF.10. Their source
translations commute and their coefficients are retained, so
they commute as actual maps. The same holds for the analytic
norm root \(\beta^{1/m}\), a full base function. Consequently,
for \(A_R=R^n\zeta_m^d\widetilde F^{\,n}D\),
\[
 (U'_G)^n A_R
       =\zeta_m^{-2ln} A_R(U'_G)^n
       =\zeta_m^{ln} A_R(U'_G)^n .
 \tag{SEF.16}
\]
This is the exact finite square, using \(m=3l\), not a numerical
inference from the cubic source action.

\par

\medskip\noindent\textbf{SEF.4. The actual primitive good-root endpoint and its triple norm}

\par\nobreak

The positive scalar winding \(J=l\) changes \(R\) to \(\zeta_3R\).
Define the forward parameter endpoint on the two pieces by
\[
 E_R^{\rm ord}=(R^{-l}S)^n,\qquad
 E_R^{\rm cap}=(R^{-l}U'_G)^n .
 \tag{SEF.17}
\]
For \(n=-1\) these are the actual inverse source/line maps, with
their reciprocal factors at the inverse source point. On the
ordinary part the scalar at the two ends is the same, so SEF.11
gives commutation with the entire \(B_\lambda\). On the cap,
SEF.16 gives the directed equality
\[
                 E_R^{\rm cap} A_R=A_{\zeta_3R}E_R^{\rm cap}.
 \tag{SEF.18}
\]
It retains every character \(d\) and the same full small division
section. It does not replace the cap generator by its effective
order-three source action.

The actual annular comparison follows from SEF.5 and SEF.14:
\[
 \begin{split}
 \Phi_{\theta+l}E_R^{\rm ord}\Phi_\theta^{-1}
 &=z^{nl}R^{-ln}t^n z^{2ln}(U'_G)^n\\
 &=R^{-ln}(U'_G)^n=E_R^{\rm cap}.
 \end{split}                                                    \tag{SEF.19}
\]
Here \(t=z^{-3l}\) is the original base function. All \(\Pi,\Lambda\)
columns and the \(B_\theta^{J_\alpha}\) factors in SEF.3 occur
before this endpoint equality, where the latter are exactly one.
The source/line inverse satisfies the inverse square. The same
full units multiply the cusp characters in SEF.10. Thus SEF.17 is
a holomorphic endpoint of the completed sphere family with its
actual reference, on every prescribed fan.

The three chronological endpoints use respectively
\(R,\zeta_3R,\zeta_3^2R\). Their constant factors multiply to
\(R^{-3ln}=\lambda^{-n}\), since the extra exponent of \(\zeta_3\)
is \(-3ln\). Therefore
\[
 \begin{split}
 E_{\zeta_3^2R}^{\rm ord}
 E_{\zeta_3R}^{\rm ord}E_R^{\rm ord}
   &=\lambda^{-n}S^{3n}
     =q^n B_\lambda^{-n},\\
 E_{\zeta_3^2R}^{\rm cap}
 E_{\zeta_3R}^{\rm cap}E_R^{\rm cap}
   &=\lambda^{-n}(U'_G)^{3n}
     =h^n B_\lambda^{-n},\\
 &\hspace{12mm}h=1-t_cz^m=q_\infty q .
 \end{split}                                                    \tag{SEF.20}
\]
The root equality uses the full
\(B'_{\rm sel}=h B'_{2,3}=h(U'_G)^{-3}\).
The quotient maps are exactly the canonical forward endpoint
\(q^n\) and its cap map \(h^n\). The physical return uses their
inverse endpoints. This is an actual three-sheet family/closing
comparison; it assigns no division of a return class.

\par

\medskip\noindent\textbf{SEF.5. Full identity-linear winding restriction and the negative norm}

\par\nobreak

Every actual holomorphic endpoint for this fixed generic source
preserves the original base, by EF.2. Its linear matrix intertwines
the identical full monodromy modules EF.18; the scalar changes only
the full suspension logarithm. The algebra in EF.4 gives
\(eI+c\delta\psi\). The two complex-independent period vectors
\(e_1=(1,0)\), \(\delta=(0,1)\) force its complex linear map to be
\(eI\). Applying it to the actual suspension vectors gives
\[
                 e(\ell_{\rm target}-\ell_{\rm source})\in\mathbb Z .
 \tag{SEF.21}
\]
Their difference is the logarithm of the full original scalar ratio;
all source and Cartier factors are identical. Hence the endpoint
scalars themselves are equal. A scalar parameter path with an
endpoint therefore has integral winding \(J\). Its actual endpoint
linear part on this common ordinary torus is \(+I\) or \(-I\);
the integer in SEF.21 remains the chosen suspension-lift change.

For a \(+I\) endpoint the same cover/source argument as EF.6 gives
the actual ordinary lift \(H(t)U_G^j\), with \(j\in\mathbb Z\).
The full root endpoint differs by the scalar \(z^{nJ}\) of SEF.5.
Thus its exact component and selected orders are
\[
 \operatorname{ord}_0H=3j/b,\quad
 \operatorname{ord}_{q=0}H=-k,\quad
 m\operatorname{ord}_\infty H+2mj/b+nJ=0.
 \tag{SEF.22}
\]
There are no other divisors, by the full EF.30 point-line calculation.
The rational degree equation now gives
\[
                 j=bk+\frac{bnJ}{m},\qquad m\mid bJ .
 \tag{SEF.23}
\]
All orders are read in the actual source/root module frames.
For the good M2 branch this says \(l\mid J\).
SEF.17 constructs \(J=l\) and its chronological powers construct
every such integer. The ratio of two endpoints with the same winding
is an actual automorphism covered by EF.6, hence a physical scalar.

There is also an exact finite-cap restriction on the negative
linear part. Lift such an endpoint to the literal normalized
\(m\)-root fibre product and then to its nonzero-line cover, keeping
the base coordinate \(z\) fixed. It takes \(B_\lambda\) to
\(B_\lambda^{-1}\). The two lifts of the actual finite square differ
by an integer suspension deck:
\[
 F A_{\rm source}=B_\lambda^k A_{\rm target}F.
 \tag{SEF.24}
\]
Iteration \(m\) times, retaining both norms
\(A_{\rm source}^m=A_{\rm target}^m=B_\lambda^n\), gives
\[
 B_\lambda^{-n}F=B_\lambda^{mk+n}F,\qquad mk=-2n.
 \tag{SEF.25}
\]
The action is strictly contracting, so equality of powers has this
literal integer meaning. Since \(m=3l\ge3\) and \(n=\pm1\), no such
integer \(k\) exists. Hence the endpoints just constructed exhaust
the holomorphic scalar endpoints in this good-root scope: winding
group \(l\mathbb Z\), with the actual primitive endpoint SEF.17
and only the connected physical scalar ambiguity.

For the reduced M6 cap with \(m=3l+2\ge5\), the identical actual
EF.32-\/-38 root argument gives the identity-linear restriction
\(m\mid J\); its canonical endpoint \(q^n\), with cap \(h^n\), is
already constructed by the full scalar family. Equation SEF.25
excludes its negative linear endpoints as well. Thus its winding
group is exactly \(m\mathbb Z\), with the same physical ambiguity.

At this local-norm stage, for the quadratic reduced cap \(m=2\), the remaining exact negative
norm equation is \(k=-n\). It is compatible with that one local
integer norm, and this note does not substitute that compatibility
for an endpoint construction. The full source/Cartier/selected
square for this remaining negative endpoint is the next explicit
gate. Its identity-linear winding group is \(2\mathbb Z\), by the
same completed divisor proof. No other exceptional-source locus
or unconstructed parameter closing is assigned here.

\par

\medskip\noindent\textbf{SEF.6. Verification scope}

\par\nobreak

The companion exact calculation retains the finite-source identity
SEF.9, the two root Miller weights, the full \(D\) commutator, both
signed endpoint squares, the three successive root phases, and the
ordinary/cap triple identities with the full selected unit \(h\).
It also checks the integral divisor equations and finite norm
restriction. The full module, properness and reference arguments
are the actual maps above and EF/EX, not consequences of those
finite algebra checks. This constructive stage and its full annular completion below have now passed the complete independent root audit.

\par

\medskip\noindent\textbf{NE.1. Objects and the literal selected inversion}

\par\nobreak

Use exactly EF.1-\/-4, either full cap progression, \(m\ge2\),
\(n=\pm1\), the prescribed finite character and reference. A scalar
endpoint has integral scalar winding \(J\), by SEF.21; its ordinary
scalar is the same at both ends. Both completed endpoints therefore
have the identical ordinary and selected pieces. The base map is the
identity by the actual meromorphic-base and critical-triple proof EF.2.

On the full selected point-theta torus retain
\[
 h=(QZ,\alpha V,q),\qquad B_\lambda=(gZ,q\eta V,q),
 \qquad Q=\kappa q,
 \tag{NE.1}
\]
with every displayed coefficient an actual extending nonzero function.
There is a literal local inversion
\[
                  I(Z,V,q)=(Z^{-1},V^{-1},q).
 \tag{NE.2}
\]
It satisfies \(Ih=h^{-1}I\) and
\(IB_\lambda=B_\lambda^{-1}I\), with the full coefficient functions,
because it fixes the base coordinate. On the periodic fan of an
integral basis \((u,v)\), negation sends
\[
 \operatorname{conv}(p,p+u,p+u+v)
 \longmapsto
 \operatorname{conv}(p',p'+v,p'+u+v),\qquad p'=-p-u-v.
 \tag{NE.3}
\]
It sends the second triangle to the first in the same way. Thus it
carries precisely this same periodic fan.

Here are the complete local coordinate arrows, including the units.
For the standard basis write
\(A=\kappa g\), \(C=\alpha\eta\). In the first vertex
\((a,b,c)\), \(q=abc,Z=b,V=c\); in the second vertex
\((u_0,v_0,w_0)\), \(q=u_0v_0w_0,Z=u_0v_0,V=u_0w_0\).
The representative \(hB_\lambda I\) is
\[
 (a,b,c)\longmapsto(ACa,C^{-1}c,A^{-1}b),\qquad
 (u_0,v_0,w_0)\longmapsto((AC)^{-1}u_0,Aw_0,Cv_0).
 \tag{NE.4}
\]
Both products are exactly \(q\). Substitution gives
\(Z'=Aq/Z,V'=Cq/V\), and the two arrows are literal inverses.
All coefficients are evaluated at this unchanged \(q\). For a
general prescribed basis, apply its full integral torus character
change and its actual period basis; the two period valuations are
again \((1,0),(0,1)\), with their complete unit multipliers. NE.4
then uses their full products \(A,C\). This proves the analytic
extension of NE.2 on the actual selected quotient and all its charts.

\par

\medskip\noindent\textbf{NE.2. A single-valued ordinary cyclic-cover lift}

\par\nobreak

Let \(F\) be an endpoint with fibre linear part \(-I\). On every
ordinary smooth fibre it takes the primitive suspension homomorphism
\(\psi\) to \(-\psi\), so it lifts locally to the genuine nonzero-line
cover. Any two local lifts differ by a power of the actual
\(B_\lambda\), and each lift satisfies
\[
                   \widetilde F B_\lambda
                        =B_\lambda^{-1}\widetilde F.
 \tag{NE.5}
\]
At each ordinary singular source fibre these covers are the actual
nonzero Cartier-line covers of the resolved source. On its small
original base disc, restrict the actual Cartier line to the O-section
and choose a nonzero holomorphic line vector there. When the point
section meets O, use its retained numerator and denominator module
frames; their actual Cartier-unit transition gives that same nonzero
vector through the intersection. The resulting source/line section
is a disc in the quotient and kills the base meridian. Every loop can
be moved off the central fibre; after subtracting the loop on this
section it is a generic fibre loop, with exactly the actual
vanishing-cycle relations. Those vanishing cycles lie in the kernel
of \(\psi\). Thus the cover class is determined by the generic fibre
and is sent to \(-\psi\). The covering criterion, followed by the
local biholomorphic inverse, extends the lift there; the full
numerator and denominator Cartier modules remain present.

At the selected puncture compose \(F\) with the explicitly extended
NE.2 map. The composition has linear part \(+I\). Evaluate it on
an actual toric section over the \(q\)-disc. A lift of its image
near the central point lies in a vertex chart, where the three
coordinate orders are nonnegative and sum to one because their
product is exactly \(q\). The source and line torus characters
therefore have finite integral orders. The actual period powers
\(h^rB_\lambda^k\) remove these orders, with all units retained.
The coordinate arrow and its inverse are then
\[
 (Z,V,q)\longmapsto(K(q)Z,C_0(q)V,q),\qquad
 (Z,V,q)\longmapsto(K(q)^{-1}Z,C_0(q)^{-1}V,q),
 \tag{NE.6}
\]
for the actual extending units obtained by this evaluation. This is
the local EF.24-\/-25 calculation, now applied to \(FI\). It supplies
a single-valued lift on the entire punctured selected disc, rather
than only a choice on its universal cover. Compose that lift with
NE.2 and the retained integral period power to obtain a
single-valued negative lift of \(F\) there.

The local lift torsor over the finite base minus the selected point
has constant deck group \(\mathbb Z\langle B_\lambda\rangle\).
The ordinary singular discs just treated add no puncture to this
torsor. The fundamental group of the remaining base
\(\mathbb C_t\setminus\{t_c\}\) is generated by the selected
meridian. Its lift monodromy is zero by NE.6. Choose a lift at one
ordinary point and continue it; NE.5 and this zero monodromy give
one single-valued holomorphic map \(\widetilde F_{\rm ord}\) on
the full ordinary nonzero-line cover. In particular it is
single-valued around a large finite \(t\)-circle. This construction
retains, rather than sets to zero, every constant integer by which
the local lifts were joined.

\par

\medskip\noindent\textbf{NE.3. The actual root lift and complete annular deck}

\par\nobreak

Normalize the literal fibre product with \(z^m=q_\infty\). The
kernel of the full finite generator is \(\langle B_\lambda\rangle\),
since \(A_\lambda^m=B_\lambda^n\) and \(n=\pm1\). Thus the
normalized pullback is the genuine root quotient
\(Y=N_a^\times/\langle B_\lambda\rangle\) of EF.32-\/-38. The
canonical lift of the endpoint fixes \(z\).

The primitive class \(\psi\) on this root quotient is determined
by its restriction to a generic fibre. To see this with the actual
source and line, the zero section of the resolved root has a
nonzero holomorphic line vector on the whole small \(z\)-disc.
At the point/O intersection use precisely the Cartier frame
\(e=\sigma_a/[(\tau-\tau(aG'))/\tau]\); off that intersection
use the retained single-divisor frame. Their transition is the
actual Cartier unit. The pullback line on the disc has these
holomorphic nonzero frames, giving a source/line section of \(Y\).
Any loop in \(Y\) can be moved off the central fibre; its projected
base loop is killed by this section. The remaining loop is a
generic fibre loop, with the actual vanishing-cycle relations.
Consequently the endpoint takes the global cover class to
\(-\psi\). It lifts to a single-valued holomorphic
\(\widetilde F_{\rm cap}:N_{a,\rm source}^\times\to
N_{a,\rm target}^\times\), fixing \(z\) and obeying NE.5.

Now retain the entire reference, not only its finite character.
Follow one positive original \(q_\infty\)-circle through all
reference sectors in EF.38. Its complete arrow is
\[
                     A_\lambda B_\lambda^\nu,
                     \qquad \nu\in\mathbb Z.
 \tag{NE.7}
\]
Here \(\nu\) is the actual sum of the carried suspension integers
along that path. Every analytic source/line factor is already in
the displayed \(A_\lambda,B_\lambda\), and the full period columns
are their actual identities in the nonzero-line cover. Nothing in
this argument chooses \(\nu=0\). The commuting-square equations
EF.38 make its \(m\)-fold chronological product exactly
\[
 (A_\lambda B_\lambda^\nu)^m
       =B_\lambda^{N},\qquad N=n+m\nu.
 \tag{NE.8}
\]
This is one positive \(z\)-circle. The scalar endpoint family
SEF.3 retains the same integer \(\nu\) at both ends. At integral
winding \(J\), its extra full comparison is \(z^{nJ}\), which
is single-valued on this root circle. Its finite generator has the
same norm \(B_\lambda^n\), including the prescribed character.
Thus the source and target annular arrows both have precisely the
monodromy NE.8.

Choose the constant suspension deck so that, at one point of the
root annulus, the two actual endpoint lifts satisfy
\[
 \widetilde F_{\rm cap}
     =\Phi_{\rm target}\widetilde F_{\rm ord}
                         \Phi_{\rm source}^{-1}.
 \tag{NE.9}
\]
The equality holds throughout the connected cut annulus. On
continuing once around the positive \(z\)-circle, its right side
becomes, by NE.5 and NE.8,
\[
 B_\lambda^N\widetilde F_{\rm cap}B_\lambda^{-N}
                  =B_\lambda^{2N}\widetilde F_{\rm cap}.
 \tag{NE.10}
\]
The root lift in NE.9 and the ordinary lift constructed in
Section NE.2 are single-valued by their proved covering extensions. Hence
\(B_\lambda^{2N}=1\). Strict contraction gives its literal infinite
order, so \(N=0\). But \(N=n+m\nu\), \(m\ge2\) and
\(n=\pm1\); no integer \(\nu\) satisfies this equation.
This proves that a negative scalar endpoint does not exist on
the stated generic source with either full cap progression,
including the previously remaining quadratic reduced cap.

\par

\medskip\noindent\textbf{NE.4. The completed winding conclusion and its exact scope}

\par\nobreak

The positive winding calculation SEF.22-\/-23 and the constructed
good-root endpoint SEF.17 now give all fixed-source holomorphic
scalar endpoints in this explicit generic source locus:
\[
 \begin{array}{c|c|c}
 (a,b;m)&\text{scalar winding group}&\text{actual primitive forward endpoint}\\
 (2,3;3l),\ l\ge1&l\mathbb Z&(R^{-l}tU_G)^n\text{ on the ordinary piece}\\
 (6,1;3l+2),\ l\ge0&m\mathbb Z&q^n\text{ on the ordinary quotient}.
 \end{array}
 \tag{NE.11}
\]
The first cap arrow is \((R^{-l}U'_G)^n\), and its threefold
composite is the full \(q^n/h^n\) endpoint square of SEF.20.
The second cap arrow is \(h^n\) from the full original scalar
family. The ratio of any two endpoints with the same winding is
the physical scalar action proved in EF, carried through every
source, root, fan and reference chart. These statements concern
the actual fixed-source endpoint maps. No numerical return is
assigned, and no class is divided by a covering degree.


\par\medskip\noindent\textbf{Completed exact certificates and full audits.}\par
The full SEF source is
\path{notes/part3_e6_fixed_source_scalar_endpoints_working.md};
its exact certificate passes 6820 checks. The full annular source is
\path{notes/part3_e6_fixed_source_negative_endpoint_working.md};
its certificate passes 9269 checks, including 360 complete integral
basis fans and the q-dependent full-unit vertex Jacobian.

The independent root proof is
\path{notes/part3_e6_scalar_endpoints_root_peer.md}.
Its fresh certificate passes 529 checks and includes 232 basis fans.
It reads all ten source sections geometrically, including the actual
point-line extensions, complete reference restoration, finite
characters, both signs, single-valued ordinary and root lifts,
and the unrestricted annular integer \(\nu\).
The full audit is \path{records/root_audit_e6_scalar_endpoints.json}.
P3K.109 supplies the complete fixed-source endpoint theorem used in
the winding classification. The exact source, export, certificate,
audit and dependency locks are in
\path{records/k110_export_manifest.json}.

% End expanded verified cp36 input

\end{proof}


\subsection{The original M2 selected-I3 return in the actual fan \texorpdfstring{\(T_{-1}\)}{T-1}}
\begin{theorem}[Source result P3K.111 under continuation review]\label{night:P3K.111}
\StatementRecord{night:P3K.111}{open}{Historical events: GPT-6 Astra (OpenAI), ChatGPT Pro mode, effort Pro, S6\_06\_BLIND, 16--18 November 2026: targeted proof check. Also checked in this writing. Continuation mathematical and source checks: GPT-6 Astra (OpenAI), ChatGPT Pro mode, effort Pro, S6\_10\_PAPERS continuation, 19 November 2026: computation or formal certificate; targeted proof check; used and rechecked in this run. Continuation metadata events: GPT-6 Astra (OpenAI), ChatGPT Pro mode, effort Pro, S6\_10\_PAPERS continuation, 19 November 2026: used and rechecked in this run; explicitly scoped metadata reconciliation; metadata reconciliation only; no new mathematical check. GPT-6 Astra (OpenAI), ChatGPT Pro mode, effort Pro, S6\_10\_PAPERS continuation, 19 November 2026: used and rechecked in this run; receiving source placement; metadata reconciliation only; no new mathematical check. GPT-6 Astra (OpenAI), ChatGPT Pro mode, effort Pro, S6\_10\_PAPERS continuation, 19 November 2026: used and rechecked in this run; source-row intake; metadata reconciliation only; no new mathematical check. \emph{Pre-continuation recorded proof scope.} Historical event: GPT-6 Astra (OpenAI), ChatGPT Pro mode, effort Pro, S6\_06\_BLIND, 16--18 November 2026: targeted proof check, recorded producing-run re-derivation. Also checked in this writing by GPT-6 Astra (OpenAI), ChatGPT Pro mode, effort Pro, S6\_10\_PAPERS continuation, 19 November 2026: computation or formal certificate for the explicitly listed algebraic substeps in the event ledger; full continuation rederivation and renewed credit comparison remain pending.}{\cref{night:P3F.35,night:P3F.41,night:P3F.46,night:P3K.68,night:P3K.96}; the original source, point line, roots, characters, reference, selected fan and scalar range stated here}

Retain the original marked E6 source, doubled point section and physical line
\[
 y^2=x^3+r_{\rm src}x^2+tx+t^2,\qquad G=(0,t),\qquad
 s=\tfrac14-r_{\rm src},\qquad M_2=\mathcal O([2G]-[O]).
\]
The source step is the three chronological translations by \(-G\), hence
\(-3G\), with the original selected height \(q=t-t_c(s)\), \(t_c(0)=0\).
For \(T_G(p)=p-G\), retain the entire Miller section
\[
 \begin{gathered}
 f_2=\frac{y-t-x/2}{x-s},\qquad
 B_{2,3}(p)=f_2(p)f_2(T_G^3p)f_2(T_Gp)^2f_2(T_G^2p)^2,\\
 B_{\rm sel}=q\,t^{-3}B_{2,3}.
 \end{gathered}
\]
Use the actual good cubic cap \((m,n)=(3,1)\), its prescribed finite line
character and complete annular reference cochain. Choose \(\lambda_0\)
in the explicit strict contraction range in ASI2.1, choose
\(r_0^3=\lambda_0\), and a sufficiently small positive basepoint \(s_*\).
The auxiliary filling is \(\lambda=\lambda_0s^{12}\); the source return
keeps the original scalar \(\lambda_{\rm base}=(r_0s_*^4)^3\) fixed.
Traverse the positive original \(s\)-meridian with its stated source,
root, point-line and physical-circle closing maps.

The selected native star is the actual one with point-coordinate matrix
\[
 P_N=\begin{pmatrix}1&-1\\-1&2\end{pmatrix},
\]
labelled \(T_{-1}\) by the explicit triangle and period morphism RFB.4--7.
Its completed Jin physical characteristic pair, actual held-quaternion
native product, and oriented sphere class are
\[
 \begin{gathered}
 (k_M,\ell_M)=(71,-4),\qquad (\mu,\nu)=(-1,72),\\
 A_M=\frac{4k_M-\ell_M}{24}=12,\qquad
 [\Phi_{\rm source}]=[12]_{28}\in\Theta_7.
 \end{gathered}
\]
Here \(k_M=\langle c_1(M)^3,[M_f]\rangle\) and
\(\ell_M=\langle p_1(TM_f)c_1(M),[M_f]\rangle\), with the original
physical line and the K68 orientation. The native pair is obtained by
the complete physical-line and held-quaternion map in ASI2.7.
All full source, Cartier, normal, cap, finite-character, reference and
endpoint functions appearing in the construction are retained.

\end{theorem}
\begin{proof}[Supplied source argument; continuation review in progress]\mbox{}\par
\noindent\textbf{Complete original-source attachment and directed physical transfer.}\par\smallskip

The nineteen complete supporting proofs are reproduced in their locked
order below. SI2 and I2P give the original source, the entire Miller and
theta coefficients, the point/O graphs and their inverses. I2M, I2Z,
I2O, I2I, I2F, I2ZC and I2L give the exact period dictionary, all three
native parents, the 136442-cone common fan and its proper endpoint lifts,
with full source bands, semigroup maps and coefficient inverses.
I2CV and CFL then calculate the complete compact fixed-stratum pair.

RB retains the whole scalar reference body and its three vanishing discs.
NS, NC and MN determine the two extending native charts, the actual cap
maps and their inverses, and the separate collision and reference normal
bundles. The class expressions on the closed replacement in NC.4 and MN
are transferred to the actual cap-completed bundles by the displayed
five-dimensional bundle bordism. IL retains all twelve lower/reference
doubles, including the intrinsic cube divisor and its boundary connection.
RFB and CSN give the full finite reference-fan continuation and independent
native scalar maps. ASI2 carries these maps through the compact attachment,
parameter square and original outer transfer.

The signed assembly is explicit. In the retained positive-inner
orientation, the compact pair is \((2723,19124)\), the twelve lower
pairs sum to \((-2712,-19092)\), and the difference of the two own-cap
moving pairs is \((-10,-28)\). Their sum is \((1,4)\), with the actual
RB reference in fan \(T_1\). RFB constructs the whole Jin reference
change \(J_{r,-1}-J_{r,1}=(-72,0)\), whose positive-inner counterpart
is \((72,0)\). Subtracting that reference gives \((-71,4)\); the
explicit Jin interval reversal gives \((71,-4)\). The original
physical-line dual and held-quaternion map then give
\(k_M=\mu+\nu\), \(\ell_M=4\mu\), hence \((-1,72)\).

The supporting sources retain their original progress and local-scope
sentences as part of their recorded proofs. The completed seed assertion
is ASI2. In particular RB's original point-coordinate descriptions are
carried to the two extending native charts by NS and MN; their finite
whole-pair comparison is the explicit RFB construction. These are the
maps used in the final assembly. The numbered dependencies are K68 for
the physical/Hopf inverse clutch, K96 for the original M2 good-root
attachment, F35 for the primitive common-link/Gysin/van Kampen maps,
and F41/F46 for the reference sphere construction. Their actual
hypotheses and boundary arrows are checked in the following proofs.


% Expanded verified cp36 input: notes/part3_e6_selected_i3_m2_return_support_01.tex
\par\bigskip\noindent\textbf{P3K.111 supporting proof 1: Original source, M2 Miller/point periods, full cubic cap and parameter square}\par\nobreak
\par

\medskip\noindent\textbf{\(M_{2}\) selected-I3: complete source, point-line, good-root and period maps}

\par\nobreak

Working exact construction, 18 November 2026. The original source is
\[
 y^2=x^3+rx^2+tx+t^2,\quad G=(0,t),\quad s=1/4-r.
\]
We use the actual point line \(M_2=\mathcal O([2G]-[O])\), the
chronological translation by \(-3G\), the ordered selected node
\(t_c(s)\to0\), cap \((m,n)=(3,1)\), its prescribed finite character,
and a fixed original scalar \(\lambda_0=r_0^3\) in the common strict
contraction range. The positive original \(s\)-circle and the
unit-\(+1\) Tate marking specify the initial fan. This note constructs
the source/line/root/annular maps and exact period parents; it assigns
no characteristic pair or return.

\par

\medskip\noindent\textbf{SI2.1. Original selected height and the full node morphism}

\par\nobreak

Use the original marked-node cover near \(w=\pm1\):
\[
 v=w^2-1,\quad D=v-2,\quad r=(1-2v)/D^2,\quad
 (X,Y,T)=(D^2x,D^3y,D^3t).
\]
The inverse divides by the same nonzero functions. Its exact equation,
selected zero and source parameter are
\[
 Y^2=X^3+(1-2v)X^2+DTX+T^2,\quad q=T-v=D^3q_{\rm old},
\]
\[
 q_{\rm old}=t-t_c(s),\quad
 t_c(s(v))=v/D^3,\quad s=v(v+4)/(4D^2),\quad s'(0)=1/4.
 \tag{SI2.1}
\]
Thus each of these marked parameter discs carries one positive
original source meridian.

Put \(s_0=v+v^2/4\),
\(U=Y-T-DX/2,\ V_s=Y+T+DX/2\). The literal source identity is
\(UV_s=X^2(X-s_0)\). On its ordered source-resolution chart
\[
 U=Xz,\quad V_s=X\omega,\quad X=s_0+z\omega,\quad
 Y=X(z+\omega)/2,\quad T=X(\omega-z-D)/2.
 \tag{SI2.2}
\]
Its selected node is \(z=-v/2,\omega=v/2\). Set
\(a=z+v/2,\ b=\omega-v/2\) and retain the analytic square root
\[
 R_b=\sqrt{(b+2)^2+4v},\quad R_b(0)=2,\quad K_b=b+v+2+R_b.
\]
The full source map is
\[
 x_n=a-vb/K_b,\quad
 y_n=[bK_b-(2b+v)a]/4,\qquad x_ny_n=q.
 \tag{SI2.3}
\]
Its inverse is the unique small solution
\[
 b=\frac{4y_n+vx_n}{2(R_b(b,v)-x_n)},\qquad
 a=x_n+vb/K_b.
 \tag{SI2.4}
\]
The denominator is a unit. On a small fixed ball the first right
side has derivative of norm less than \(1/2\) in \(b\); successive
substitution gives its holomorphic inverse and the displayed
substitution verifies both compositions. The normal Jacobian is
\(\sqrt{1+v}\) with value one. This keeps the full selected height
and source parameter, including both source branches.

The original inverse Weierstrass/Tate map E81, with its unit-\(+1\)
root, supplies the actual functions
\[
 Q=P^2R,\quad P=t u_P,\quad R-P=s u_s,\quad
 u_P(0)=2,\quad u_s(0)=1,\quad t=P u_t.
\]
The complete discriminant comparison gives
\[
 Q=t^2q_{\rm old}u_Q,\quad u_Q(0)=8,\quad
 q_{\rm old}=R c_q,\quad c_q=u_P^2/u_Q,\quad c_q(0)=1/2.
 \tag{SI2.5}
\]
These are the original analytic source functions, not just their
orders. In particular \(R=0\) is exactly the selected-height divisor.

\par

\medskip\noindent\textbf{SI2.2. Six Miller factors and the genuine \(M_{2}\) point frame}

\par\nobreak

Let \(T_G(p)=p-G\). The literal point-to-tensor coefficient is
\[
 f_2=\frac{y-t-x/2}{x-s},\qquad r_2=f_2^{-1}.
\]
The full chronological suspension on the genuine point line is
\[
 B_{2,3}(p)=f_2(p)f_2(T_G^3p)
                 f_2(T_Gp)^2f_2(T_G^2p)^2,\qquad
 B_{\rm sel}^{(2,3)}=(q_{\rm old}/t^3)B_{2,3}.
 \tag{SI2.6}
\]
The source scaling in SI2.1 gives
\[
 f_{2,\rm old}=D^{-1}f_{2,\rm new},\quad
 r_{2,\rm old}=D r_{2,\rm new},\quad
 B_{2,3,\rm old}=D^{-6}B_{2,3,\rm new}.
\]
Therefore \(B_{\rm sel}^{(2,3)}=q B_{2,3,\rm new}/T^3\)
exactly. The canonical meromorphic point-section coefficients are
identified by the source isomorphism. Its balanced tensor coefficient
has the additional \(D\) factor written below; that factor is retained.

Write \(f_P(Z)=\Theta_Q(Z/P)/\Theta_Q(Z)\) and retain the full
degree-one coefficient
\[
 \eta=a_T\,\Theta'_Q(1)\,\Theta_Q(P^2)/\Theta_Q(P)^2.
\]
Here \(a_T\) is the actual Weierstrass differential scale, with \(a_T(0)=1\); the separate derivative has \(\Theta'_Q(1)(0)=-1\), hence \(\eta(0)=-1\). The derivative factor remains in every point/tensor arrow and its inverse. In particular the full sixth power is \(\eta^6=a_T^6\Theta'_Q(1)^6\Theta_Q(P^2)^6/\Theta_Q(P)^{12}\). The formulas below use the actual \(\eta\) fixed by the displayed Miller identity; this derivative factor is present in their full analytic coefficients, as well as in the first-power point-frame comparison. The original chord identity and exact telescope are
\[
 \eta=f_2(p)\frac{f_P(P^{-1}Z)}{f_P(Z)},\qquad
 f_P(Z)f_P(P^{-1}Z)=f_{P^2}(Z).
 \tag{SI2.7}
\]
They give the full module arrows
\[
 V_{2,T}=Z^{-1}f_P(Z)^2r_2(p)v_2,\quad
 V_{\rm reg}=ZV_{2,T}=\eta^{-1}V_{{\rm pt},2},\quad
 V_{{\rm pt},2}=f_{P^2}(Z)v_2.
 \tag{SI2.8}
\]
Under SI2.1, \(\eta_{\rm old}=D^{-1}\eta_{\rm new}\),
\(V_{{\rm pt},2,\rm old}=V_{{\rm pt},2,\rm new}\), and
\(V_{\rm reg,old}=D V_{\rm reg,new}\). This explicitly types the
point/tensor change and its inverse.

The six factors in SI2.6 telescope, with their order retained, to
\[
 B_{2,3}=\eta^6
       \frac{f_{P^2}(Z)}{f_{P^2}(P^{-3}Z)}.
 \tag{SI2.9}
\]
Consequently the complete source/line periods are
\[
 h:(Z,V_{2,T})\mapsto(P^2RZ,R^{-1}V_{2,T}),\quad
 B:(Z,V_{2,T})\mapsto(P^{-3}Z,\lambda q_{\rm old}K V_{2,T}),
\]
\[
 K=(P/t)^3\eta^6,\quad K(0)=8,
 \tag{SI2.10}
\]
or, in the actual point-related coefficient,
\[
 h:(V_{\rm reg},Z)\mapsto(P^2V_{\rm reg},P^2RZ),\quad
 B:(V_{\rm reg},Z)\mapsto
   (\lambda q_{\rm old}\eta^6t^{-3}V_{\rm reg},P^{-3}Z).
 \tag{SI2.11}
\]
All \(q_{\rm old}=Rc_q\), \(t=Pu_t\), and \(\eta\) factors remain.

At \(s=0\), set \(A_*=y-t-x/2,\ B_*=y+t+x/2\).
Then \(A_*B_*=x^3\) and the three chronological coefficients are
\(A_*/x,-B_*/x,2t/x\). Their product is \(-2t\), hence
\(B_{2,3}=4t^2\). The selected height is now exactly \(t\);
the genuine point Hom map is \(B_{\rm sel}^{(2,3)}=4\), covering
the identity \(T_G^3=1\). In the retained Tate lifts \(Q=P^3\),
so the actual central period is
\[
 hB:(Z,V_{\rm reg})\longmapsto(Z,4\lambda V_{\rm reg}).
 \tag{SI2.12}
\]
This is the full central coefficient, derived from the Miller
maps, not a chosen scalar normalization.

On the protected selected \(q_{\rm old}\)-disc over a nonzero
\(s\)-annulus, put
\[
 Q=q_{\rm old}\kappa,\quad \kappa=P^2/c_q,\quad
 \alpha=P^2,\quad g=P^{-3},\quad
 \eta_*=\lambda_{\rm base}\eta^6t^{-3}.
 \tag{SI2.13}
\]
Their exact parameter windings are \((2,2,-3,-3)\). The full
fixed-scalar winding matrix in the source/point coefficient marking is
\[
 \begin{pmatrix}2&-3\\2&-3\end{pmatrix}.
 \tag{SI2.14}
\]
No return or free coordinate is inferred from this matrix.

\par

\medskip\noindent\textbf{SI2.3. Point/O modules and both ordered markings}

\par\nobreak

For a retained source band \(j\), set
\[
 x=Q^jZ,\quad y=Q^{1-j}/Z,\quad xy=P^2R.
\]
The actual point graph is the blowup of \((x,P^2)\):
\[
 x=P^2a,\ ay=R;\qquad P^2=xc,\ y=Rc.
 \tag{SI2.15}
\]
The second chart is resolved by
\[
 P=uv,\ x=u,\ c=uv^2;\qquad P=uv,\ x=u^2v,\ c=v,
 \tag{SI2.16}
\]
with the literal overlap \(u_2=v_1^{-1},v_2=u_1v_1^2\).
The source functions \(s=(R-P)/u_s,\ t=Pu_t\) travel through these
proper graph maps. In the first chart \(ds\) has a nonzero \(dP\)
coefficient near the collision; in the two latter charts it has
a nonzero \(dR\) coefficient. Thus the actual source parameter is
preserved and the projection is a submersion on these small charts.

Put \(b=P^2,e=R\). The full theta product has
\[
 b f_b(x)=(b-x)E_{b,e}(x,y),
\quad
 E_{b,e}=
 \frac{\prod_{r\ge1}(1-Q^{r-1}ex)(1-Q^{r-1}by)}
 {(1-x)\prod_{r\ge1}(1-Q^rx)(1-Q^{r-1}y)}.
 \tag{SI2.17}
\]
On each chosen small node chart this product and its reciprocal
converge normally to units. In the actual point generators
\(e_a=(a-1)^{-1},e_c=(c-1)^{-1}\), writing
\(v_2=e_a A=e_c C\), the full arrows are
\[
 V_{\rm reg}=-\eta^{-1}P^{-2j}E_{b,e}A,\qquad
 V_{\rm reg}=\eta^{-1}P^{-2j}c^{-1}E_{b,e}C,\quad C=-cA.
 \tag{SI2.18}
\]
Their inverses retain every displayed factor. Near O use
\(v_2=(x-1)C_O\); then
\[
 V_{\rm reg}=\eta^{-1}P^{-2j-2}(x-P^2)E_OC_O,
 \tag{SI2.19}
\]
where
\[
 E_O=
 \frac{\prod_{r\ge1}(1-Q^rx/P^2)(1-Q^rP^2/x)}
 {\prod_{r\ge1}(1-Q^rx)(1-Q^r/x)} .
\]
In particular the additional \(P^{-2}\) at O remains in the
coefficient map. Under \(h\) use \(j'=j-1\); the local source and
module coefficients are unchanged. Under \(B\), for the actual
target band \(j'\),
\[
 x'=Q^{j'-j}P^{-3}x,\qquad
 v_2'=\lambda q_{\rm old}\eta^6t^{-3}P^{2(j'-j)}
                 \frac{f_b(x)}{f_b(x')}v_2.
 \tag{SI2.20}
\]
The reciprocal is evaluated at the inverse source point. This is
the full six-factor Hom by SI2.9.

The original unit-\(+1\) Tate source has
\(Z_T=Z_w\) near \(w=-1\) and \(Z_T=Z_w^{-1}\) near \(w=+1\).
The full theta inversion gives the point-line arrow
\[
 (Z,V_{{\rm pt},2})\mapsto(Z^{-1},P^2V_{{\rm pt},2}).
\]
The target point representative is \(P^{-1}\); its actual source
period is \((Q,P^{-2})\). Composing gives
\[
 (Z,V_{{\rm pt},2})\mapsto(Q/Z,V_{{\rm pt},2}).
 \tag{SI2.21}
\]
This is the literal reflected-fan character substitution, with both
point-module maps SI2.18--20 retained. The Weierstrass differential
changes sign with \(Z\)-inversion; the corresponding sign in
\(\Theta(P^{-2})/\Theta(P^{-1})^2\) leaves the full \(\eta\)
unchanged. The source/cap family itself is the original one.

\par

\medskip\noindent\textbf{SI2.4. The cubic good root and full annular reference}

\par\nobreak

Use \(q_\infty=z^3\) and the literal original-source arrow
\[
 t=z^{-3},\quad x=z^{-2}X,\quad y=z^{-3}Y,\qquad
 Y^2=X^3+r z^2X^2+zX+1 .
 \tag{SI2.22}
\]
The deck is \(D_3(z,X,Y)=(\zeta_3z,\zeta_3^{-4}X,Y)\).
The marked sections are \(G'=(0,1)\) and
\(2G'=(s z^2,-1-(s/2)z^3)\). The root source is a smooth proper
elliptic family near \(s=z=0\).
All six Miller factors give
\[
 f_{2,\rm old}=z^{-1}f'_2,\quad
 B_{2,3,\rm old}=z^{-6}B'_{2,3},\quad
 B'_{\rm sel}=h_\infty B'_{2,3},\quad
 h_\infty=1-t_c(s)z^3 .
 \tag{SI2.23}
\]
On \(z=0\), the three factors are
\((Y-1)/X,-(Y+1)/X,2/X\), of product \(-2\).
On \(s=0\), the selected height has \(t_c(0)=0\), so the complete
central Hom is \(B'_{\rm sel}=4\), on both parameter axes.

The original point-line correction is retained on the quotient
graph of weights \((w,z)=(2,1)\):
\[
 \beta^*\mathcal O(2G'-O)=\pi^*M_2+E_P-E_O .
 \tag{SI2.24}
\]
The quotient rays are \((2/3,1/3),(1/3,2/3)\); the correction is
the literal minimum \(\min(v_w,2v_z)\), with its monomial generator
at both marked sections. These are OG.6\textquotesingle s actual module maps.

Let \(R_0'\) be the unique small division \(3R_0'=3G'\). Its formal
coordinate is obtained from
\[
 \tau(3G')=z^2
 \frac{-2s(2+sz^3)}{8+6sz^3+(s^2-s^3)z^6};
\quad
 \tau(R_0')=s z^2\ell_R,\quad \ell_R(0,0)=-1/6.
 \tag{SI2.25}
\]
It and its inverse preserve the actual refinement ideals
\((w,z),(w,z^2),(z,w^2)\). Their canonical divisor sections carry
the complete correction SI2.24.
For \(P'=2G'\), the normalized chord is
\[
 F=\frac{\tau(R_0')\ell_{P',R_0'}}{X-X_{P'+R_0'}} .
 \tag{SI2.26}
\]
Writing \(X_{R_0'}=A/\tau^2,Y_{R_0'}=-A/\tau^3\), its numerator is
\[
 \tau(Y-Y_{P'})+
 \frac{A+\tau^3Y_{P'}}{A-\tau^2X_{P'}}(X-X_{P'}),
\]
whose denominator is a unit and whose exact Hom value is one at
\(\tau=0\). The horizontal divisor and the zero vertical orders
extend this Hom and its inverse jointly across both parameter axes.
The chord and \(\tau\) deck weights cancel, so \(F\) commutes with
\(D_3\).

Set \(\beta=B'_{\rm sel}/F^3\), with the full chronological product
understood. This is a nonzero base function of central value four.
Its prescribed invariant analytic cube root gives
\[
 \widetilde F=\beta^{1/3}F,\quad
 A_0=\zeta_3^d\widetilde F D_3,\quad
 A_\lambda=r_\lambda A_0,\quad
 B_\lambda=\lambda B'_{\rm sel},\quad r_\lambda^3=\lambda.
 \tag{SI2.27}
\]
Thus \(A_\lambda^3=B_\lambda\) and they commute, including
the full finite character and both point corrections.

On each original infinity sector let
\(\ell_\alpha=\log_\alpha q_\infty/(2\pi i)\),
\(z_\alpha=e^{2\pi i\ell_\alpha/3}\). The full original annular map is
\[
 \Phi_\alpha=
 \operatorname{Exp}\{\ell_\alpha\operatorname{Log}_\alpha\widetilde F+
 \ell_\alpha(\log r_\lambda/(2\pi i)+d/3)\delta+h_\alpha\}\,I_{z_\alpha}.
 \tag{SI2.28}
\]
Here \(I_{z_\alpha}\) includes SI2.24. The complete OG.12--13 sector
construction, with the displayed \(t_c(s)\) and \(M_{2}\) modules substituted,
gives
\(\Phi_\beta=A_\lambda^{j_0}B_\lambda^{j_{\alpha\beta}}\Phi_\alpha\)
and \(\Phi_\alpha B_\lambda=B_\lambda\Phi_\alpha\).
The reference-zero Cauchy operators use fixed annular contours;
the retained cochain contains the full source period and line lattice.
Its inverse reverses the sector, root and Lie arrows. Changing from
the distant selected node changes the normalized root by the actual
invariant cube root of
\((1-t_c(s)z^3)/(1-t_{\rm distant}(s)z^3)\).
This full unit and the original integer reference cochain remain
in SI2.28; neither is reset.

\par

\medskip\noindent\textbf{SI2.5. The admitted twelfth-power scalar and its complete phase square}

\par\nobreak

Choose
\[
 \lambda=\lambda_0s^{12},\quad r_\lambda=r_0s^4,\quad r_0^3=\lambda_0.
 \tag{SI2.29}
\]
For \(A=-\log|P|,\ C=-\log|R|,\ D_B=-\log|\lambda q_{\rm old}K|\),
the full period determinant is
\[
 (2A+C)D_B-3AC.
 \tag{SI2.30}
\]
The actual equation \(R-P=su_s\) gives
\(D_B\ge12\min(A,C)+C-C_0\), with bounded full-unit logarithms.
Since \(3AC/(2A+C)\le3\min(A,C)\), shrinking the bidisc gives a
strict positive determinant gap with every unit retained.

Use the independent root parameter \(r\), with
\(B=r^3B_0(s),A=rA_0(s)\), and the entire annular reference
\(F=F_0(s)+\rho\,\ell_\infty\delta\),
\(\rho=\log r/(2\pi i)\). Along SI2.29 the varying part is
\(4\tau\ell_\infty\delta\), \(\tau=\log s/(2\pi i)\).
The ordinary and cap displacements are
\[
 o_{\rm ord}=4\tau\frac{\log q_{\rm old}}{2\pi i}\delta,\qquad
 o_{\rm cap}=4\tau\frac{\log h_\infty}{2\pi i}\delta.
 \tag{SI2.31}
\]
The exact equality \(q_\infty q_{\rm old}=h_\infty\) gives
\(F+o_{\rm ord}-o_{\rm cap}=F_{\rm ref}\). The full forward square is
\[
 \Phi_{\rho+4}q_{\rm old}^{\,4}=h_\infty^4\Phi_\rho .
 \tag{SI2.32}
\]
Its endpoint-to-initial arrows are their full inverses.
With \(\beta_2=\eta^6t^{-3}\), the selected ordinary endpoint has
\[
 D_{\rm ord}^4=B^4R_4,\qquad
 R_4=(P^{12}Z,\lambda^{-4}\beta_2^{-4}V_{{\rm pt},2}).
 \tag{SI2.33}
\]
This retains the selected source translation and line factor at
the selected-height zero.

At a reference source point \(s_*\), the actual root basepoint is
\(r_{\rm base}=r_0s_*^4\). On the source cut put
\[
 G=\operatorname{Log}(P^{-3}),\quad
 J=\operatorname{Log}(\lambda_{\rm base}\beta_2),\qquad
 C_\varrho=(e^{\varrho G},
   e^{\varrho J+3\pi i\varrho(\varrho-1)}).
 \tag{SI2.34}
\]
The complete source-cut increments of both \(G\) and \(J\) are
\(-6\pi i\). With
\(\eta_\varrho=e^{6\pi i\varrho}\lambda_{\rm base}\beta_2\),
\[
 C_{\varrho+1}(g^{-1},\eta_\varrho^{-1})=C_\varrho,\quad
 C_0=1,\quad C_4=(g^4,\eta_*^4)=R_4^{-1}.
 \tag{SI2.35}
\]
For the native Cox lift
\((a_0,a_1,c_0,c_1)\mapsto(a_0,Aa_1,(AD)^{-1}c_0,Dc_1)\),
the source-cut mismatch is \((A,D)=(e^{-6\pi i\varrho},e^{-6\pi i\varrho})\).
Its literal quaternion map is
\[
 q_H\longmapsto e^{3\pi i\varrho}q_H e^{9\pi i\varrho}.
 \tag{SI2.36}
\]
At every integer seam its two half-phase signs are equal.
The base and orbit sphere rotate by the displayed left adjoint
map, which extends over the actual three-ball and carries the
principal physical line. The source-cut faces and their Spin
endpoints are therefore fully specified.

In the original DZ origin, substitution of the four root turns
in its actual quadratic product formula gives the retained factor
\[
 e^{3\pi i(16t^2+4t)}
   =e^{48\pi it(t-1)}e^{60\pi it}.
 \tag{SI2.37}
\]
The first logarithm has zero endpoint values; the second is the
physical loop of winding \(30\). It travels through both compared
boundary products, with their entire Cox and cap maps.
Thus the source/root parameter words are \((1,4)\), \((0,4)\)
and \((1,0)\), with the actual edge morphisms SI2.31--37 retained.
The open scalar path from \(r_{\rm base}\) to \(r_0\) is
\(r(a)=r_0\exp(4(1-a)\operatorname{Log}s_*)\), with its full lifted
log in SI2.28. The normal and compact central attachment has yet
to be evaluated, so these word identities assign no return.

\par

\medskip\noindent\textbf{SI2.6. Exact three endpoint period parents}

\par\nobreak

On \(R=pP\), put
\[
 c_{12}=\lambda_0c_q\eta^6/(u_s^{12}u_t^3),\quad
 c_{12}(0)=4\lambda_0,\quad f_0=(p-1)^{12}c_{12}.
 \tag{SI2.38}
\]
The original periods and the words \(K=hB,\ J=h\) are
\[
 h=(P^2V_{\rm reg},P^3pZ),\quad
 B=(P^{10}p f_0V_{\rm reg},P^{-3}Z),
\]
\[
 K=(P^{12}p f_0V_{\rm reg},pZ),\qquad
 J=(P^2V_{\rm reg},P^3pZ).
 \tag{SI2.39}
\]
In common characters \(X_c=V_{\rm reg}/Z,\ Z_c=Z\), take
\(L=K^2J,T=KJ\). Their entire arrows are
\[
 L=(P^{23}p^{-1}f_0^2X_c,P^3p^3Z_c),\qquad
 T=(P^{11}p^{-1}f_0X_c,P^3p^2Z_c).
 \tag{SI2.40}
\]
The generic period matrix is
\(\left(\begin{smallmatrix}23&11\\3&3\end{smallmatrix}\right)\)
with area \(36\). Both character changes and period words retain
their determinant and inverse; they do not replace a coefficient
by its valuation.

At zero use \(X_0=V_{\rm reg}^{-1},Y_0=Z\) and
\(L_0=B^{-1},T_0=hB^{-1}\):
\[
 L_0=(P^{10}p f_0X_0,P^3Y_0),\qquad
 T_0=(P^8p f_0X_0,P^6pY_0).
 \tag{SI2.41}
\]
Here \(f_0\) is a unit near \(p=0\). For \(r=\operatorname{ord}P\),
\(z=\operatorname{ord}p\), the full matrix and determinant are
\[
 M_0=\begin{pmatrix}10r+z&8r+z\\3r&6r+z\end{pmatrix},\qquad
 \det M_0=36r^2+13rz+z^2.
 \tag{SI2.42}
\]
The exact inspection map \(S=\left(\begin{smallmatrix}3&-3\\0&2\end{smallmatrix}\right)\)
identifies the original lattice with its index-six image lattice.
It gives
\[
 SM_0=\begin{pmatrix}21r+3z&6r\\6r&12r+2z\end{pmatrix}.
 \tag{SI2.43}
\]
The whole periodic support is the negative-square interpolation of
this symmetric integral form on that specified lattice. Its transverse
jumps are \(6r,15r+3z,6r+2z\). Pulling the exact dual cell back gives
\[
 (0,0),\ (-5r-z,0),\ (-10r-z,-3r),\
 (-13r-2z,-6r-z),\ (-8r-z,-6r-z),\ (-3r-z,-3r-z).
 \tag{SI2.44}
\]
This polygon spans two actual Tate source bands, separated by
\(y+3r+z=0\). Their maps are
\[
 (x_1,y_1)=(QZ,Z^{-1}),\qquad
 (x_2,y_2)=(Q^2Z,Q^{-1}Z^{-1}),
 \tag{SI2.45}
\]
with \(Q=P^3p\); all exponents are nonnegative in their respective
half-polygons. The coefficient arrow is SI2.18--20 with the actual
band index. The source-compatible regular four-fan is being
constructed from this periodic parent and these source cuts.

At one put \(v_1=p-1,\ X_1=Z^{-1},Y_1=V_{\rm reg}/Z\).
For \(L_1=h^{-1},T_1=hB\), the exact maps are
\[
 L_1=(P^3pX_1,PpY_1),\qquad
 T_1=(p^{-1}X_1,P^{12}v_1^{12}c_{12}Y_1).
 \tag{SI2.46}
\]
Their matrix is
\(\left(\begin{smallmatrix}3r&0\\r&12r+12z\end{smallmatrix}\right)\),
determinant \(36r^2+36rz\), with \(z=\operatorname{ord}v_1\).
The source and line \(p\)-factors are still present.

At infinity use \(u=P/R\) and
\[
 f_\infty=(1-u)^{12}c_{12}(Ru,R),\qquad f_0=u^{-12}f_\infty.
\]
The same common words are
\[
 L=(R^{23}f_\infty^2X_c,R^3Z_c),\quad
 T=(R^{11}f_\infty X_c,R^3uZ_c).
 \tag{SI2.47}
\]
Their matrix is
\(\left(\begin{smallmatrix}23r&11r\\3r&3r+z\end{smallmatrix}\right)\),
determinant \(36r^2+23rz\). The substitution \(P=Ru,p=u^{-1}\),
all reciprocal point generators and the selected-height factor remain
in this equality.

The actual original four-form, in the zero characters, is
\[
 \Omega=(2J_{ts}/a_T)\,P\,dP\wedge dp\wedge
                        d\log X_0\wedge d\log Y_0.
 \tag{SI2.48}
\]
Here \(J_{ts}=\det\partial(t,s)/\partial(P,R)\), with \(J_{ts}(0)=1/2\),
and \(a_T(0)=1\). Its full unit and exceptional \(P\) factor are
retained. Thus a primitive zero-end ray of base heights \((r,z)\)
has original canonical order \(2r+z-1\), with the actual chart
Jacobian sign still to be recorded on each regular cone.

The earlier sixth-power period attempt is preserved separately
with its exact constant-symmetrizer obstruction. It supplies no
completed boundary and no numerical input here.


% End expanded verified cp36 input


% Expanded verified cp36 input: notes/part3_e6_selected_i3_m2_return_support_02.tex
\par\bigskip\noindent\textbf{P3K.111 supporting proof 2: Original marked point/O graph and inverse}\par\nobreak
\par

\medskip\noindent\textbf{Selected-I3 \(M_{2}\) N12: the original point graph and complete Cartier coefficient}

\par\nobreak

This note retains SI2.1-\/-6, the original source parameter functions of
E81 and the locked common/end atlases I2F and its actual endpoint parents and lifts. It establishes
local source and line-module arrows. The finite marked boundary
attachment and compact fixed-normal ledger are separate calculations;
no characteristic pair is assigned here.

\par

\medskip\noindent\textbf{I2P.1. The literal point and source-node graph}

\par\nobreak

Write \(Q=P^{2} R\), \(b=P^{2}\) and \(e=R\), with every E81 analytic unit retained.
The actual Tate representative of 2G is b, and \(be=Q\). In a retained
source band j put \(x=Q^j Z\), \(y=Q^{1-j}/Z\), so \(xy=P^{2} R\).
The source period is always its written arrow; changing j is accompanied
by its \(P^{2j}\) line factor below.

The closure of the graph of \([x:P^{2}]\) is the blowup of the ideal
\((x,P^{2})\), with its strict-transform equations
\[
 x=P^2 a,\quad ay=R;\qquad
 P^2=x c,\quad y=Rc.                                    \tag{I2P.1}
\]
Both project to the original source by the displayed x,y,P,R, and
therefore preserve the original s,t through E81. On the common open
set \(c=a^{-1}\). The point section is \(a=1\) or \(c=1\). The first chart is
smooth with coordinates (P,a,y), \(R=ay\). The second is the product
of \(xc=P^{2}\) with the R coordinate. Its actual two-chart resolution is
\[
 P=uv,\quad x=u,\quad c=uv^2,\quad y=Ruv^2;
 \qquad
 P=uv,\quad x=u^2v,\quad c=v,\quad y=Rv.                 \tag{I2P.2}
\]
The two charts have the literal overlap \(u_2=v_1^{-1}\),
\(v_2=u_1v_1^{2}\). Their maps give the ordered P,P,R resolution of
\(xy=P^{2}R\). The projection is the composition of the graph blowup and
this explicit \(A_{1}\) resolution. These are proper morphisms over the
original parameter bidisc. The original zero section is retained
in its separate source chart x near 1 (or the adjacent-period chart).
On Q nonzero their Tate points cannot coincide: equality would give
\(P^{2}=Q^k\); the valuations for \(|P|,|R|<1\) exclude
this for every integer k. Indeed \(k\le 0\) would make the right side
have modulus at least 1, and \(k\ge 1\) gives \(1=P^{2k-2}R^k\), whose modulus
is strictly less than 1. Over the parameter axes the strict point
remains in \(a=1\) or \(c=1\), with \(x=P^{2}\) small, whereas the zero section is
in x near 1. These are disjoint source neighbourhoods on the graph.

\par

\medskip\noindent\textbf{I2P.2. Full theta products and the actual module generators}

\par\nobreak

The complete theta identity on a source-node neighbourhood is
\[
 b f_b(x)=(b-x)E_{b,e}(x,y),
\]
\[
 E_{b,e}=
 \frac{\prod_{r\ge1}(1-Q^{r-1}ex)(1-Q^{r-1}by)}
 {(1-x)\prod_{r\ge1}(1-Q^r x)(1-Q^{r-1}y)}.              \tag{I2P.3}
\]
This is obtained by factoring the \(r=0\) factor from
\(\Theta_Q(x/b)/\Theta_Q(x)\) and substituting \(Q=be=xy\). On the
chosen closed small node chart all factors in the displayed ratio
are units; the products converge normally and their reciprocals do
also. Its central value is 1. The actual Miller unit is
\(\eta=a_T \Theta'_Q(1) \Theta_Q(P^{2})/\Theta_Q(P)^{2}, \eta(0)=-1\),
with the original invariant differential \(a_T(0)=1\). This entire
function and its reciprocal remain in every following arrow.

Let \(v_{2}\) be the coefficient in the original canonical meromorphic
point-line frame. The exact SI2.2/band arrow is
\[
 V_{reg}=\eta^{-1}P^{-2j} f_b(x)v_2,
 \qquad v_2=\eta P^{2j}V_{reg}/f_b(x).                    \tag{I2P.4}
\]
This equality denotes a morphism of the specified Cartier modules.
In I2P.1 use \(e_a=1/(a-1)\), \(e_c=1/(c-1)\) as the local meromorphic
point-divisor generators, and write \(v_{2}=e_a A=e_c C\). Direct substitution
in I2P.3 gives the full frame arrows and their inverses:
\[
 V_{reg}=-\eta^{-1}P^{-2j}E_{b,e}A,
 \qquad A=-\eta P^{2j}E_{b,e}^{-1}V_{reg},               \tag{I2P.5}
\]
\[
 V_{reg}=\eta^{-1}P^{-2j}c^{-1}E_{b,e}C,
 \qquad C=\eta P^{2j}cE_{b,e}^{-1}V_{reg}.                \tag{I2P.6}
\]
The apparent c pole in I2P.6 has its stated Cartier type. It is
retained when pulling back to the two I2P.2 charts. On their point
neighbourhood c is a unit. On the two point charts \(e_c=-a e_a\),
so \(C=-c A\). Substitution verifies exactly the equality of I2P.5 and
I2P.6, including the minus sign, \(\eta\) and every P power.

Near O, let x be near 1 and use its actual generator \(e_O=x-1\) for
the negative point divisor. The full smooth-chart theta unit \(E_O\) is
\[
 f_b(x)=\frac{x-b}{b(x-1)}E_O,
\quad E_O=
 \frac{\prod_{r\ge1}(1-Q^r x/b)(1-Q^r b/x)}
 {\prod_{r\ge1}(1-Q^r x)(1-Q^r/x)}.
\]
Here \(Q/b=R\), so \(E_O\) and its reciprocal extend through the whole
small parameter bidisc. With \(v_{2}=e_O C_O\) the actual arrow is
\[
 V_{reg}=\eta^{-1}P^{-2j-2}(x-P^2)E_O C_O,
 \quad C_O=\eta P^{2j+2}(x-P^2)^{-1}E_O^{-1}V_{reg}.
                                                               \tag{I2P.7}
\]
The extra factor \(P^{2}\) at O is retained. It has not been folded into
\(\eta\) or an unspecified unit. The relation with the actual original
\(M_{2}\) module is SI2.2\textquotesingle s \(r_{2}=f_{2}^{-1}\) and SI2.2\textquotesingle s telescoping theta identity;
I2P.4 is exactly that relation before the graph substitutions. Thus
I2P.5-\/-7 carry both original marked divisors and the original
point-to-tensor coefficient, rather than choosing a new line from
its Chern class.

\par

\medskip\noindent\textbf{I2P.3. Directed source periods and suspension on every frame}

\par\nobreak

For the actual source period H, set \(j'=j-1\) after applying
\((Z,V_{\rm reg})\longmapsto (QZ,P^{2}V_{\rm reg})\). Then x,y and each coefficient \(A,C,C_O\)
in I2P.5-\/-7 are literally unchanged. This proves the full module
square on both source-node charts and on O; no unrecorded reindex
is available in this square.

Here \(q=t-t_c(s)=R c_q\) is the actual selected height,
\(R-P=s u_s\) and \(t=P u_t\). For the actual suspension B, write
\[
 d_B=\lambda_0s^{12}q\eta^6t^{-3},\quad
 (Z,V_{reg})\longmapsto(P^{-3}Z,d_B V_{reg}).             \tag{I2P.8}
\]
If the target retained band is j\textquotesingle, its full source coordinate is
\(x'=Q^{j'-j}P^{-3}x\). The meromorphic point coefficient transforms by
\[
 v_2'=d_B P^{2(j'-j)}\frac{f_b(x)}{f_b(x')}v_2.          \tag{I2P.9}
\]
The inverse has the reciprocal coefficient evaluated at the actual
inverse source point. Equations I2P.5-\/-7, on source and target, give
its local module coefficients by substitution, retaining both theta
units and their actual marked sections. At \(j'=j\) this is exactly
SI2.2\textquotesingle s six-factor chronological Hom: the original f2 theta identity
telescopes through all three source translations. For general j\textquotesingle{}
the equality \(f_b(Q^r z)=b^r f_b(z)\) supplies precisely the displayed
\(P^{2(j'-j)}\), so no new Hom scalar occurs.

On a common/end chart all x,y,P,R are the explicit monomials and
full units of I2F.8-\/-9, I2I.3, I2O.3 and the complete zero-parent source bands. Pulling I2P.1 back
means blowing up the actual ideal generated by those two functions,
not only a valuation vector. The source and line coefficient
morphisms above specify the marked finite graphs needed by that
pullback. Period-translated graph owners must be carried with the
full I2P.8-\/-9 maps when constructing the final compact normal ledger.
This last global assembly is not asserted by the present local note.

The source translation, full coefficient, point graph and inverse
Cartier maps above are fixed before any characteristic calculation.
The full \(\eta\) sixth power retains \(\Theta'_Q(1)^{6}\); the forward
point frame and its inverse retain \(\eta\) itself.


% End expanded verified cp36 input


% Expanded verified cp36 input: notes/part3_e6_selected_i3_m2_return_support_03.tex
\par\bigskip\noindent\textbf{P3K.111 supporting proof 3: Complete common period dictionary}\par\nobreak
\par

\medskip\noindent\textbf{Selected-I3 \(M_{2}\) N12: the complete common period dictionary}

\par\nobreak

The original SI2 source, point/O modules, finite character,
full annular reference and fixed original scalar are retained.
Use \(Q=P^{2}R\), \(R-P=s u_s\), \(t=P u_t\), \(q=R c_q\), \(p=R/P\) and

\[
\begin{gathered}
c_{12}=\lambda_0 c_q \eta^{6}/(u_s^{12} u_t^{3}), \\
f_{0}=(p-1)^{12} c_{12}.
\end{gathered}
\]

\par

\medskip\noindent\textbf{I2M.1. Common characters and full periods}

\par\nobreak

The genuine point coefficient has \(V_{\rm reg}=\eta^{-1} V_{{\rm pt},2}\). In the
ordered genuine characters \((V_{\rm reg},Z)\), the original periods are

\[
\begin{gathered}
K=hB=(P^{12} p f_{0} V_{\rm reg},pZ), \\
J=h=(P^{2} V_{\rm reg},P^{3}pZ).
\end{gathered}
\]

Set \((X_c,Z_c)=(V_{\rm reg}/Z,Z)\), \(L_c=K^{2}J\) and \(T_c=KJ\). Their exact maps are

\[
\begin{gathered}
L_c:(X_c,Z_c)\longmapsto (P^{23} p^{-1} f_{0}^{2} X_c,P^{3}p^{3} Z_c), \\
T_c:(X_c,Z_c)\longmapsto (P^{11} p^{-1} f_{0} X_c,P^{3}p^{2} Z_c).
\end{gathered}
\]

The inverse character is \(V_{\rm reg}=X_c Z_c\). The inverse period words
are \(K=L_c T_c^{-1}\) and \(J=T_c^{2} L_c^{-1}\). Both matrices have determinant 1.
Every entire coefficient is fixed by both period actions, so
the displayed compositions commute exactly, including their
units. The generic period matrix is \(\begin{pmatrix}23&11\\3&3\end{pmatrix}\), with
original area 36.

\par

\medskip\noindent\textbf{I2M.2. The zero and one parent maps in both directions}

\par\nobreak

At zero use \((X_{0},Y_{0})=(V_{\rm reg}^{-1},Z)\), \(L_{0}=K^{-1}J\) and \(T_{0}=K^{-1}J^{2}\).
The exact maps are

\[
\begin{gathered}
L_{0}:(X_{0},Y_{0})\longmapsto (P^{10} p f_{0} X_{0},P^{3}Y_{0}), \\
T_{0}:(X_{0},Y_{0})\longmapsto (P^{8} p f_{0} X_{0},P^{6} pY_{0}).
\end{gathered}
\]

The common character and period matrices are

\[
\begin{gathered}
A_{c0}=\begin{pmatrix}-1&-1\\0&1\end{pmatrix}, C_{c0}=\begin{pmatrix}-5&-3\\3&2\end{pmatrix}.
\end{gathered}
\]

They act by \((X_c,Z_c)=((X_{0}Y_{0})^{-1},Y_{0})\), \(L_c=L_{0}^{-5} T_{0}^{3}\) and
\(T_c=L_{0}^{-3} T_{0}^{2}\). Their exact inverse maps are
\((X_{0},Y_{0})=((X_cZ_c)^{-1},Z_c)\), \(L_{0}=L_c^{-2} T_c^{3}\), \(T_{0}=L_c^{-3} T_c^{5}\).
Both matrices have determinant −1. Direct composition of the
full functions in these displayed arrows gives I2M.1.

At one use \((X_{1},Y_{1})=(Z^{-1},V_{\rm reg}/Z)\), \(L_{1}=J^{-1}\) and \(T_{1}=K\).
Writing \(v=p-1\), the entire period functions are

\[
\begin{gathered}
L_{1}:(X_{1},Y_{1})\longmapsto (P^{3}p X_{1},PpY_{1}), \\
T_{1}:(X_{1},Y_{1})\longmapsto (p^{-1} X_{1},P^{12} v^{12} c_{12} Y_{1}).
\end{gathered}
\]

The exact common matrices are

\[
\begin{gathered}
A_{c1}=\begin{pmatrix}0&1\\-1&0\end{pmatrix}, C_{c1}=\begin{pmatrix}-1&-1\\2&1\end{pmatrix}.
\end{gathered}
\]

Thus \((X_c,Z_c)=(Y_{1},X_{1}^{-1})\), \(L_c=L_{1}^{-1} T_{1}^{2}\), \(T_c=L_{1}^{-1}T_{1}\).
Their inverses are \((X_{1},Y_{1})=(Z_c^{-1},X_c)\), \(L_{1}=L_c T_c^{-2}\),
\(T_{1}=L_c T_c^{-1}\). Both matrices have determinant 1. Substitution
of \(f_{0}=v^{12} c_{12}\) gives exactly the two full I2M.1 maps.
The actual generic matrix identities are

\[
\begin{gathered}
A_{c0} \begin{pmatrix}10&8\\3&6\end{pmatrix} C_{c0} \\
= A_{c1} \begin{pmatrix}3&0\\1&12\end{pmatrix} C_{c1} \\
= \begin{pmatrix}23&11\\3&3\end{pmatrix}.
\end{gathered}
\]

\par

\medskip\noindent\textbf{I2M.3. Infinity and every full translating coefficient}

\par\nobreak

At infinity the characters already equal \((X_c,Z_c)\). The actual
height is \(R=pP\), \(u=p^{-1}\) and \(f_{\rm infinity}=p^{-12} f_{0}\). Consequently

\[
\begin{gathered}
R^{23} f_{\rm infinity}^{2}=P^{23} p^{-1} f_{0}^{2}, \\
R^{11} f_{\rm infinity}=P^{11} p^{-1} f_{0}, \\
R^{3}=P^{3}p^{3}, R^{3}u=P^{3}p^{2}.
\end{gathered}
\]

These are exact full analytic equalities and the literal
arrows of I2I.1. For the translate \(L_c^i T_c^j\), the full unit
vectors in character order \((X_c,Z_c,P)\) are

\[
\begin{gathered}
u_p(i,j)=(-i-j,3i+2j,0), \\
u_f(i,j)=(2i+j,0,0).
\end{gathered}
\]

Their original valuation shift is \(((23i+11j)h,(3i+3j)h,0)\).
At zero p remains the second original parameter, so its full
orders remain coordinate orders. At one the exact substitution
\(f_{0}=v^{12} c_{12}\) adds \(12(2i+j)\) to the first character\textquotesingle s v order.
At infinity use both \(P=Ru\) and \(f_{0}=u^{-12} f_{\rm infinity}\). All inverse
arrows use the same full functions and opposite word. These
identities fix every period, source and unit map used by the
common toric overlay.


% End expanded verified cp36 input


% Expanded verified cp36 input: notes/part3_e6_selected_i3_m2_return_support_04.tex
\par\bigskip\noindent\textbf{P3K.111 supporting proof 4: Zero raw parent, lattice and full coefficients}\par\nobreak
\par

\medskip\noindent\textbf{Selected-I3 \(M_{2}\) source: the complete N12 zero endpoint}

\par\nobreak

This is the zero endpoint of SI2.1--6 for the actual line
\(M_2=\mathcal O([2G]-[O])\), translation by \(-3G\), original good
cap \((m,n)=(3,1)\), and coupled scalar
\(\lambda=s^{12}\lambda _0\), \(r_\lambda=s^4r_0\).
It constructs a finite regular endpoint with its two literal Tate source
bands, full period arrows, physical circle and original determinant.
The executed endpoint has 420 regular four-cones. The point/O modules
are carried on their actual graphs below. This note assigns no
characteristic pair or completed source return.

\par

\medskip\noindent\textbf{I2Z.1. Actual characters, units and the inspection lattice}

\par\nobreak

Use the SI2 characters and functions
\[
 X_0=V_{\rm reg}^{-1},\quad Y_0=Z,\quad p=R/P,\quad
 Q=P^3p,\quad
 f_0=(p-1)^{12}c_{12},\qquad
 c_{12}=\frac{\lambda _0c_q\eta^6}{u_s^{12}u_t^3}.
 \tag{I2Z.1}
\]
Here \(q_{\rm old}=Pc_qp\), \(s=P(p-1)/u_s\), \(t=Pu_t\), and
\[
 \eta=a_T\Theta'_Q(1)\frac{\Theta_Q(P^2)}{\Theta_Q(P)^2},
 \qquad \eta(0)=-1.
 \tag{I2Z.2}
\]
Thus the entire sixth power in (I2Z.1) includes
\(\Theta'_Q(1)^6\). All these functions are the original SI2 functions;
on a sufficiently small zero-end base disc \(p-1\), \(c_q,u_s,u_t,\eta\)
are nonzero. In particular \(f_0\) is an actual unit on this endpoint.
The original period words, without changing the physical line, are
\[
\begin{aligned}
 L_0=B^{-1}:&(X_0,Y_0)\longmapsto
       (P^{10}pf_0X_0,P^3Y_0),\\
 T_0=hB^{-1}:&(X_0,Y_0)\longmapsto
       (P^8pf_0X_0,P^6pY_0).
\end{aligned}
\tag{I2Z.3}
\]
The inverse arrows divide by the displayed factors at the same base
point. The base \(P,p\) is fixed by both periods, so these arrows commute,
including their units.

In the original valuation lattice
\(N=\mathbb Z^4_{(x,y,r,z)}
=\mathbb Z^4_{\operatorname{ord}(X_0,Y_0,P,p)}\), the period matrix is
\[
 M_0(r,z)=
 \begin{pmatrix}10r+z&8r+z\\3r&6r+z\end{pmatrix},
 \qquad \det M_0=36r^2+13rz+z^2.
 \tag{I2Z.4}
\]
The integer map
\[
 S=\begin{pmatrix}3&-3\\0&2\end{pmatrix},\qquad
 S\mathbb Z^2=3\mathbb Z\oplus2\mathbb Z
 \tag{I2Z.5}
\]
has index six. Its inverse is used only on that displayed image lattice.
It gives
\[
 SM_0=
 \begin{pmatrix}21r+3z&6r\\6r&12r+2z\end{pmatrix}.
 \tag{I2Z.6}
\]
Define \(\nu\) by affine interpolation of
\(\nu(n)=n(n-1)/2\) at every integer. Its slope on \((n,n+1)\)
is \(n\). The explicit convex potential
\[
 \Psi(\theta_1,\theta_2)=
 (15r+3z)\nu(\theta_1)+6r\nu(\theta_1+\theta_2)
       +(6r+2z)\nu(\theta_2)
 \tag{I2Z.6a}
\]
has the three stated positive slope jumps in the interior \(r,z>0\).
The identity
\(\nu(t+n)=\nu(t)+nt+n(n-1)/2\), valid on every interval
and at its endpoints, gives
\(\nabla\Psi(\theta+n)-\nabla\Psi(\theta)=SM_0n\).
The additive constant is exactly
\((15r+3z)n_1(n_1-1)/2+
6r(n_1+n_2)(n_1+n_2-1)/2+
(6r+2z)n_2(n_2-1)/2\).
At zero the subgradient is the sum of the three segments from zero
to \(-(15r+3z,0)\), \(-(6r,6r)\) and \(-(0,6r+2z)\).
All its vertices lie in the explicit lattice (I2Z.5).
Pulling this actual subgradient cell back by \(S^{-1}\), the
fundamental hexagon in the original lattice has vertices
\[
\begin{aligned}
 &(0,0),\ (-5r-z,0),\ (-10r-z,-3r),\\
 &(-13r-2z,-6r-z),\ (-8r-z,-6r-z),\
 (-3r-z,-3r-z).
\end{aligned}
\tag{I2Z.7}
\]
Direct substitution gives area \(\det M_0(r,z)\). Opposite edges
are paired by (I2Z.3), with the equations in I2Z.4 below. Thus the
index-six inspection map does not replace the original period lattice.

For comparison, the preserved exponent-six attempt had
\[
 M_r=\begin{pmatrix}4&2\\3&6\end{pmatrix},\quad
 M_z=\begin{pmatrix}1&1\\0&1\end{pmatrix},\quad
 M_z^{-1}M_r=\begin{pmatrix}1&-4\\3&6\end{pmatrix}.
 \tag{I2Z.8}
\]
Its characteristic polynomial is \(t^2-7t+18\), with discriminant
\(-23\). A constant real multiplier making both \(M_r,M_z\)
symmetric positive would make \(M_z^{-1}M_r\) self-adjoint for the
positive form supplied by \(M_z\), contradicting these nonreal roots.
This records the precise stopped constant-symmetrizer construction.
For the present exponent twelve the corresponding eigenvalues are
four and nine, and the actual map is (I2Z.5).

\par

\medskip\noindent\textbf{I2Z.2. The full raw endpoint and its invariant monomial maps}

\par\nobreak

At \(z=0,r=1\), retain the 72 triangles of the immutable upper file
\path{e6_selected_i3_m2_zero_upper.json}. Their vertices lie in the hexagon
\[
 (0,0),(-5,0),(-10,-3),(-13,-6),(-8,-6),(-3,-3),
 \tag{I2Z.9}
\]
and every triangle is primitive. Their union has area 36. This
triangulation includes the actual source line \(y+3r=0\).

Write the six lifted vertices as \(U_0,\ldots,U_5\), with
\((r,z)=(1,0)\), and set
\[
 A=(0,0,0,1),\ B_1=(-1,0,0,1),\
 B_2=(-2,-1,0,1),\ B_3=(-1,-1,0,1).
 \tag{I2Z.10}
\]
The following is the complete raw four-cone list, with each displayed
column set ordered to have positive determinant:

\begin{enumerate}
\def\labelenumi{\arabic{enumi}.}

\item
  \(A\) joined to each of the 72 upper triangles;
\item
  \((A,B_1,U_1,U_2)\);
\item
  \((A,B_1,(-10-k,-3-k,1,0),(-11-k,-4-k,1,0))\),
  \(k=0,1,2\);
\item
  \((A,B_1,U_3,B_2)\);
\item
  \((A,B_3,(-13+k,-6,1,0),(-12+k,-6,1,0))\),
  \(k=0,\ldots,4\);
\item
  \((A,B_3,U_3,B_2)\);
\item
  \((A,B_3,U_4,U_5)\).
\end{enumerate}

These 84 cones are obtained by coning the facets of the convex
\(r+z=1\) slice not containing \(A\). On each such facet the displayed
segments give its triangulation. Consequently the cones have disjoint
interiors and cover exactly the slice of (I2Z.7). The independent volume
identity is
\[
 \sum_\sigma\frac{|\det M_\sigma|}
 {6\prod_{\rho\subset\sigma}(r_\rho+z_\rho)}
 =\int_0^1\{36r^2+13r(1-r)+(1-r)^2\}\,dr
 =\frac{29}{2}.
 \tag{I2Z.11}
\]
Their determinants are one, two or three.

For a raw matrix \(M_\sigma\) of determinant \(d\), let \(u_1,\ldots,u_4\)
be its finite cover coordinates. The actual affine quotient functions
are the invariant monomials \(u^a\) with
\[
 a\in\mathbb Z_{\ge0}^4,\qquad aM_\sigma^{-1}\in\mathbb Z^4.
 \tag{I2Z.12}
\]
The saved list contains all such nonconstant monomials with
\(|a|\le d\). It generates the full invariant semigroup: an invariant
word of length greater than \(d\) in the character group of order \(d\)
has a nonempty zero-sum subword of length at most \(d\), by repetition
among its partial sums; remove it and continue. This supplies the
actual coordinate ring and finite quotient map.

For every refined cone with ray matrix \(N_\tau\), the recorded map is
\[
 u^a\longmapsto
 \prod_{\ell=1}^4 t_\ell^{(aM_\sigma^{-1}N_\tau)_\ell}.
 \tag{I2Z.13}
\]
Every exponent is an integer and nonnegative. On the dense torus its
inverse uses the same original four characters. This proves the
holomorphic morphism to the raw quotient, including its finite
isotropy, without adjoining fractional cover coordinates.

\par

\medskip\noindent\textbf{I2Z.3. Source-compatible resolution and the complete finite execution}

\par\nobreak

Insert, in this order, the three actual rays
\[
 (-9,-2,1,1),\quad(-7,-1,1,1),\quad(-6,-5,1,1).
 \tag{I2Z.14}
\]
The first two resolve the determinant-three raw cone; the third
resolves the determinant-two raw cone. The saved nonnegative
barycentric columns give each source and target monomial map through
(I2Z.13). All resulting cones have determinant one. All three inserted
rays have both base heights positive, so the upper 72 triangles remain
literally unchanged.

Now use the full source functional
\[
 \ell(x,y,r,z)=y+3r+z.
 \tag{I2Z.15}
\]
For every edge with opposite \(\ell\)-signs, the positive endpoint in
this execution has \(\ell=1\). Insert the sum of its two primitive
endpoints. In a regular chart this is the ordinary two-coordinate
blowup: on its two charts the original pair is \((a,ab)\) and
\((cd,d)\). Its inverse on the overlap divides by the displayed
nonzero coordinate. The negative endpoint value increases by one
along the resulting crossing edge, so the recorded finite sequence
ends with every cone on one side of (I2Z.15).

Complete the boundary period orbits of these rays in the already
subdivided fan. The order is material: an earlier attempt restarted
from the coarse regular fan and then inserted the same orbit members;
its executed cones still crossed (I2Z.15). That exact failed output is
preserved. The retained construction inserts into the actual
source-compatible fan, and every subsequent cone is contained in its
predecessor.

The resulting opposite period facets initially had different
diagonals. Their exact common subdivision is also retained. Normalize
each boundary ray by \(r+z\). On the \(L_0\) and \(L_0/T_0\) faces use
coordinates \((y/(r+z),r/(r+z))\); on the \(T_0\) faces use
\((x/(r+z),r/(r+z))\). Intersect all edges of a source triangulation
with all edges of the inverse-transported target triangulation using
the two-by-two rational determinant formula. Clear denominators and
divide by the greatest common divisor to obtain each actual primitive
intersection ray. Insert its full period orbit.

These insertions are performed through ordinary face blowups. In a
regular owner cone express the desired ray as an integral nonnegative
column. Insert the sum of the basis rays in its positive support,
then repeat in the child containing the desired ray. The sum of its
positive coefficients strictly decreases until that ray is present.
Every such insertion has the usual monomial blowup charts. The saved
execution has two face-overlay rounds: six new intersection rays
requiring twelve ordinary stars, then four requiring eight. The next
complete comparison has no mismatched period facet.

The full sequence has 77 inserted centres and 420 final regular
four-cones. Every final cone lies in its recorded raw parent and
still lies on one side of (I2Z.15). The exact parent-by-parent slice
volumes agree with their raw-parent volumes. The final incidence
table has 760 internal facets, 72 protected generic facets, two lower
facets and 86 period facets. These lists exhaust the boundary.

\par

\medskip\noindent\textbf{I2Z.4. Full boundary periods, source bands and point/O frames}

\par\nobreak

On valuation columns the complete two periods are
\[
 T_L=\begin{pmatrix}1&0&10&1\\0&1&3&0\\0&0&1&0\\0&0&0&1\end{pmatrix},
 \quad
 T_T=\begin{pmatrix}1&0&8&1\\0&1&6&1\\0&0&1&0\\0&0&0&1\end{pmatrix}.
 \tag{I2Z.16}
\]
Their three paired boundary types are
\[
\begin{array}{c|c|c}
\text{word}&\text{source equation}&\text{target equation}\\ \hline
L_0&x-y+7r+z=0&x-y=0\\
T_0&y+6r+z=0&y=0\\
L_0T_0^{-1}&3x-5y+15r+3z=0&3x-5y-6r-2z=0 .
\end{array}
\tag{I2Z.17}
\]
The full word \(L_0^aT_0^b\) adds the coefficient \(f_0^{a+b}\)
in the \(X_0\) character. If source and target ray matrices are
\(N_s,N_t\), the actual local coefficient arrow is
\[
 t'_i=f_0^{(N_t^{-1}(a+b,0,0,0)^t)_i}
       \prod_jt_j^{(N_t^{-1}T_L^aT_T^bN_s)_{ij}}.
 \tag{I2Z.18}
\]
The inverse uses \(N_s^{-1}T_T^{-b}T_L^{-a}N_t\) and
\(N_s^{-1}(-a-b,0,0,0)^t\). On the indicated localized faces all
exponents of vanishing coordinates are nonnegative integers. The
two matrices compose to identity, and the transported two unit
columns sum to zero. Thus every original \(p\), \(c_q\), \(u_s\),
\(u_t\), \(\lambda_0\) and \(\eta^6\) factor remains in the actual
forward and inverse maps.

For a final cone on \(\ell\ge0\), take \(j=1\); on \(\ell\le0\),
take \(j=2\). The literal source node coordinates are
\[
 x_j=Q^j Z,\qquad y_j=Q^{1-j}Z^{-1},\qquad x_jy_j=P^3p .
 \tag{I2Z.19}
\]
Their exponent rows on a ray \((x,y,r,z)\) are
\[
 y+3jr+jz,\qquad -y+3(1-j)r+(1-j)z.
 \tag{I2Z.20}
\]
They are nonnegative on every cone assigned to that band. Across
the band overlap \(x_2=Qx_1\), \(y_2=Q^{-1}y_1\), which is the actual
Tate node reindex. No full \(p\) power is suppressed.

Write \(f_b(x)=\Theta_Q(x/b)/\Theta_Q(x)\), \(b=P^2\).
The original \(M_{2}\) coefficient and the genuine point coefficient satisfy
\[
 V_{\rm reg}=\eta^{-1}P^{-2j}f_b(x_j)v_2,\qquad
 v_2=\frac{\eta P^{2j}}{X_0f_b(x_j)},\qquad
 V_{\rm pt}=\eta V_{\rm reg}.
 \tag{I2Z.21}
\]
For the actual point graph \(\operatorname{Bl}_{(x_j,P^2)}\), use
\(x_j=P^2a,\ ay_j=R\), or \(P^2=x_jc,\ y_j=Rc\).
The full SI2 theta product \(E\) gives the coefficient arrows
\[
 V_{\rm reg}=-\eta^{-1}P^{-2j}EA,\qquad
 V_{\rm reg}=\eta^{-1}P^{-2j}c^{-1}EC,\qquad C=-cA.
 \tag{I2Z.22}
\]
On the O chart \(v_2=(x_j-1)C_O\), the arrow is
\[
 V_{\rm reg}=\eta^{-1}P^{-2j-2}(x_j-P^2)E_O C_O .
 \tag{I2Z.23}
\]
The units \(E,E_O\), both graph charts, the section coincidence and
their inverse coefficient maps are those proved in SI2.3. Equations
(I2Z.19)--(I2Z.23) are the source and line morphisms carried by this
endpoint. Finite marked-boundary graph pullbacks use precisely these
ideals and module generators.

\par

\medskip\noindent\textbf{I2Z.5. Physical weights, original determinant and certificate scope}

\par\nobreak

The original physical circle has weight
\((-1,0,0,0)\) in the \((X_0,Y_0,P,p)\) characters. Therefore in
every final chart \(N_\tau\) its exact coordinate weights are
\[
 N_\tau^{-1}(-1,0,0,0)^t.
 \tag{I2Z.24}
\]
The nonzero weights, including every finite stabilizer they determine,
are retained in the output. Period arrows (I2Z.18) commute with this
circle because their coefficients are source-base functions.

With the original SI2 differential and
\(J_{ts}=\det\partial(t,s)/\partial(P,R)\), the original four-form is
\[
 \Omega_{\rm original}
 =\frac{2J_{ts}}{a_T}\,
 P\,dP\wedge dp\wedge d\log X_0\wedge d\log Y_0,\qquad
 J_{ts}(0)=\frac12,\quad a_T(0)=1 .
 \tag{I2Z.25}
\]
The character order and the wedge order in (I2Z.25) have positive
comparison. Each saved ray matrix has determinant \(+1\); hence a
ray \((x,y,r,z)\) has original canonical order \(2r+z-1\).
The unit \(2J_{ts}/a_T\), including its phase, is retained. This is
the original determinant map through the base blowup and the actual
source reindices.

The executable \path{e6_selected_i3_m2_zero_fan.py} proves 4,320 exact
identities and writes the full cone, invariant-semigroup, source,
point-module, unit, inverse, physical-weight and canonical-order
tables. Its input upper file is immutable. Its proof/source hashes
are locked on execution. Every stellar centre and every period-face
intersection is saved, so the finite construction is reproducible.
The two stopped source/period attempts are preserved under
\path{records/attempts/e6_selected_i3_m2_zero_first_source_orbit} and
\path{records/attempts/e6_selected_i3_m2_zero_before_period_face_overlay.}
The resulting endpoint is the zero parent for the subsequent common
generic fan lift; its global compact fixed-locus and outer transfer
are computed only after those full three-end identifications.


% End expanded verified cp36 input


% Expanded verified cp36 input: notes/part3_e6_selected_i3_m2_return_support_05.tex
\par\bigskip\noindent\textbf{P3K.111 supporting proof 5: One raw parent and full coefficients}\par\nobreak
\par

\medskip\noindent\textbf{Selected-I3 \(M_{2}\): the complete N12 one parent}

\par\nobreak

We retain the original source \(y^{2}=x^{3}+r_{\rm source} x^{2}+t x+t^{2}\), \(G=(0,t)\),
\(M_{2}=\mathcal O([2G]-[O])\), translation \(-3G\), \(s=1/4-r_{\rm source}\) and the selected
height \(q=t-t_c(s)\), \(t_c(0)=0\). SI2 fixes the actual coupled scalar
\(\lambda=\lambda_0 s^{12}\), its root \(r_{\rm sc}=r_{0} s^{4}\), and the full source,
finite character and four-root-turn reference square. This note
constructs the one parent in that original lattice.

\par

\medskip\noindent\textbf{I2O.1. Complete period functions and the original lattice}

\par\nobreak

Use the entire SI2 functions \(Q=P^{2}R\), \(R-P=s u_s\), \(t=P u_t\) and
\(q=R c_q\). Put \(p=R/P\), \(v=p-1\) and

\[
\begin{gathered}
c_{12}=\lambda_0 c_q \eta^{6}/(u_s^{12} u_t^{3}), \\
f_{0}=v^{12} c_{12}.
\end{gathered}
\]

The genuine-coordinate periods in order \((V_{\rm reg},Z)\) are

\[
\begin{gathered}
K=hB=(P^{12} p v^{12} c_{12} V_{\rm reg},p Z), \\
J=h=(P^{2} V_{\rm reg},P^{3}p Z).
\end{gathered}
\]

These retain their original point-line module \(V_{\rm reg}=\eta^{-1} V_{{\rm pt},2}\)
and the full theta identity with \(v_{2}\). Set \(X=Z^{-1},Y=V_{\rm reg}/Z\) and
\(L=J^{-1},T=K\). Their exact arrows are

\[
\begin{gathered}
L:(X,Y)\longmapsto (P^{3}p X,Pp Y), \\
T:(X,Y)\longmapsto (p^{-1} X,P^{12} v^{12} c_{12} Y).
\end{gathered}
\]

The character inverse is \(Z=X^{-1},V_{\rm reg}=Y/X\); the period inverse
is \(J=L^{-1},K=T\). Both matrices are integral and determinant 1.
Every inverse coefficient is the reciprocal of the displayed
full function. These functions are invariant under all periods.
For \(r=\operatorname{ord} P,z=\operatorname{ord} v\) the exact period matrix and determinant are

\[
\begin{gathered}
N=\begin{pmatrix}3r&0\\r&12r+12z\end{pmatrix}, \operatorname{det} N=36r^{2}+36rz.
\end{gathered}
\]

The inspection map \((x,y)\longmapsto (x,-x+3y)\) maps the original integer
lattice isomorphically onto its explicit index 3 image. It sends
N to \(\operatorname{diag}(3r,36r+36z)\). The rectangular quadratic cells for the
form \(3r \binom{n_{1}}{2}+(36r+36z)\binom{n_{2}}{2}\), pulled back through
this specified lattice map, have the original parallelogram

\[
\begin{gathered}
(0,0),(3r,r),(3r,13r+12z),(0,12r+12z).
\end{gathered}
\]

Its original area equals \(\operatorname{det} N\). Its opposite faces are paired
by exactly L and T, including the full p and \(c_{12}\) coefficients.
The finite-index image is retained throughout this construction.

\par

\medskip\noindent\textbf{I2O.2. Integral parent charts and proper regular subdivision}

\par\nobreak

At \(r=1,z=0\) triangulate the upper parallelogram by inserting its
integral points into every containing triangle or edge, starting
with the two diagonal triangles. The executed upper list has 72
primitive triangles, area 36, with matching period boundary edges.
At \(r=0\) use the thirteen primitive lower rays \(E_j=(0,j,0,1)\),
\(0\le j\le 12\). Join \(E_0\) to all upper triangles. For \(0\le j<12\) join
\(E_j\),\(E_{j+1}\) to the upper edge from \((3,13,1,0)\) to \((0,12,1,0)\).
These lower parents have determinant 3. The joins cover every
slice of the displayed parallelogram, and their common faces
are exactly the shared faces of these joins.

For each singular primitive matrix B, the code enumerates the
finite residue group \(B^{-1} \mathbb Z^4/\mathbb Z^4\) using its four inverse columns.
Its cardinality equals \(\operatorname{det} B\). Each nonzero representative gives
an integral point of the half-open fundamental parallelepiped.
Divide it by its gcd and insert the selected primitive ray in
every cone containing it. The positive barycentric coefficients
are strictly less than 1, so every changed determinant strictly
decreases. Every step retains its full parent and child list.
The process terminates in integral determinant-one cones.

For each raw parent the primitive dual rays and all lattice
points of their half-open parallelepiped generate its entire
dual semigroup: subtract the coefficient floors of an arbitrary
point. The certificate records this entire finite generating
list and all nonnegative integral exponent maps to the regular
charts. On their original dense torus the inverse is the exact
rational character matrix. The subdivisions preserve the whole
support; these affine maps therefore assemble to the proper
map from the regular parent to the original toric support.

\par

\medskip\noindent\textbf{I2O.3. Full boundary arrows, source modules and volume}

\par\nobreak

The complete periodic facets are

\[
\begin{gathered}
L: x=0 \longmapsto  x=3r, \\
T: 3y-x=0 \longmapsto  3y-x-36r-36z=0.
\end{gathered}
\]

Their original lattice shifts are \((3r,r)\) and \((0,12r+12z)\).
The full p columns are \((1,1)\) and \((-1,0)\), respectively; the \(c_{12}\)
columns are \((0,0)\) and \((0,1)\). The inverse maps use the opposite
shifts and reciprocal units. The executed full sets of final
three-faces agree under these exact maps. For each regular
source/target pair their matrix is \(B_{\rm target}^{-1} T_{\rm period} B_{\rm source}\)
and each full unit column is \(B_{\rm target}^{-1} u_{\rm period}\). Composing
with the actual inverse cancels both the matrix and each full
coefficient, giving the original quotient gluing.

Every source-band chart has the full functions

\[
\begin{gathered}
x_{\rm source}=QZ=pP^{3}/X, y_{\rm source}=Z^{-1}=X, \\
x_{\rm source} y_{\rm source}=Q.
\end{gathered}
\]

Their ray orders \(3r-x\) and x are nonnegative on every cone.
The full inverse Tate chart therefore extends to the original
Weierstrass source. The actual original point-line coefficient
and its inverse arrow are

\[
\begin{gathered}
v_{2}=\eta P^{2} Y/[X f_{P^{2}}(pP^{3}/X)], \\
Y=\eta^{-1} P^{-2} X f_{P^{2}}(pP^{3}/X) v_{2},
\end{gathered}
\]

as morphisms in the specified point/O Cartier modules. The
proper graph of \((x_{\rm source},P^{2})\) is carried through these exact
arrows, with its point and zero-section factors retained.

The physical circle has valuation vector \((0,1,0,0)\); its four
weights on each regular chart are the inverse matrix applied
to this vector. In the original source-volume coefficient the
change to \((X,Y)\) gives exactly

\[
\begin{gathered}
\Omega=-(2J_{\rm ts}/a_T) P \mathrm dP\wedge \mathrm dv\wedge \mathrm d\log X\wedge \mathrm d\log Y.
\end{gathered}
\]

The Jacobian \(J_{\rm ts}\) and invariant differential \(a_T\) remain their
entire nonzero units. Every original canonical ray order is
\(2r+z-1\). The coupled scalar contributes only source-base forms
to \(\mathrm d\log V\), whose wedge with \(\mathrm dP\wedge \mathrm dv\) vanishes. All source orders,
physical weights, inverse maps and canonical orders are saved
in the certificate. The characteristic pair awaits the common
attachment, complete point graph and compact normal ledger.

The \(\eta\) in every displayed point-module arrow is exactly
\(a_T \Theta'_Q(1) \Theta_Q(P^{2})/\Theta_Q(P)^{2}\), with \(\eta(0)=-1\)
and \(a_T(0)=1\). In \(c_{12}\) its full sixth power includes
\(\Theta'_Q(1)^{6}\). The source SI2 certificate retains this entire
identity and the reciprocal frame morphism.


% End expanded verified cp36 input


% Expanded verified cp36 input: notes/part3_e6_selected_i3_m2_return_support_06.tex
\par\bigskip\noindent\textbf{P3K.111 supporting proof 6: Infinity raw parent and full coefficients}\par\nobreak
\par

\medskip\noindent\textbf{Selected-I3 \(M_{2}\): the complete N12 infinity parent}

\par\nobreak

We retain the original source \(y^{2}=x^{3}+r_{\rm source} x^{2}+t x+t^{2}\), \(G=(0,t)\),
\(M_{2}=\mathcal O([2G]-[O])\), translation \(-3G\), \(s=1/4-r_{\rm source}\) and the selected
height \(q=t-t_c(s)\), \(t_c(0)=0\). The actual balanced and genuine point coefficients are
\(V_{2,T}=Z^{-1} f_P(Z)^{2} r_{2} v_{2}\) and
\(V_{\rm reg}=Z V_{2,T}=\eta^{-1} V_{{\rm pt},2}\), with the full OG/MP point/O module
map retained. For this actual coupled
filling choose \(\lambda=\lambda_0 s^{12}\) and root \(r_{\rm sc}=s^{4} r_{0}\),
\(r_{0}^{3}=\lambda_0\). The full four-root-turn reference square is carried
by the source proof. This note constructs its infinity source,
period and toric parent, without assigning a characteristic pair.

\par

\medskip\noindent\textbf{I2I.1. Entire functions and the original lattice}

\par\nobreak

Use \(Q=P^{2}R\), \(R-P=s u_s\), \(t=P u_t\), \(q=R c_q\), with the original
analytic units of E81/SI3. Put

\[
\begin{gathered}
c_{12}=\lambda_0 c_q \eta^{6}/(u_s^{12} u_t^{3}), \\
p=R/P, f_{0}=(p-1)^{12} c_{12}.
\end{gathered}
\]

The full genuine-coordinate periods, in order \((V_{\rm reg},Z)\), are

\[
\begin{gathered}
K=hB=(P^{12} p f_{0} V_{\rm reg},pZ), \\
J=h=(P^{2} V_{\rm reg},P^{3}pZ).
\end{gathered}
\]

They fix P,p and every source function. At infinity set \(P=Ru\),
\(p=u^{-1}\) and \(f=(1-u)^{12} c_{12}(P=Ru,R)\). The exact unit equation is
\(f_{0}=u^{-12} f\). Thus the actual periods are

\[
\begin{gathered}
K=(R^{12} u^{-1} f V_{\rm reg},u^{-1} Z), \\
J=(R^{2}u^{2} V_{\rm reg},R^{3}u^{2} Z).
\end{gathered}
\]

Every inverse uses these reciprocal full functions at the actual
inverse point. With \(X=V_{\rm reg}/Z\) and \(L=K^{2}J,T=KJ\), direct composition
gives

\[
\begin{gathered}
L:(X,Z)\longmapsto (R^{23} f^{2} X,R^{3} Z), \\
T:(X,Z)\longmapsto (R^{11} f X,R^{3}u Z).
\end{gathered}
\]

The inverse period maps are \(K=L T^{-1}\) and \(J=T^{2} L^{-1}\); the character
inverse is \(V_{\rm reg}=XZ\). The two integer matrices have determinant 1.
The valuation matrix for \(r=\operatorname{ord} R,z=\operatorname{ord} u\) is

\[
\begin{gathered}
N(r,z)=\begin{pmatrix}23r&11r\\3r&3r+z\end{pmatrix}, \\
\operatorname{det} N=36r^{2}+23rz.
\end{gathered}
\]

The inspection map \((X_{\rm val},Z_{\rm val})\longmapsto (3X_{\rm val},11Z_{\rm val})\) is an isomorphism
from the original integer lattice onto its stated index 33 image.
On that image the period matrix is the symmetric matrix
\(\begin{pmatrix}69r&33r\\33r&33r+11z\end{pmatrix}\). Its actual negative-square quadratic
on period labels is
\(69r \binom{n_{1}}{2}+(33r+11z)\binom{n_{2}}{2}+33r n_{1} n_{2}\).
Its affine interpolation has transverse jumps \(36r,33r,11z\).
The full change under a label shift is the bilinear form of
that matrix. Pulling its dual cells back by the specified
lattice isomorphism gives the original-lattice hexagon

\[
\begin{gathered}
(0,0), (-12r,0), (-23r,-3r), \\
(-23r,-3r-z), (-11r,-3r-z), (0,-z).
\end{gathered}
\]

Its opposite faces are glued by \(L,T,L/T\). Its area is exactly
\(\operatorname{det} N\). At \(z=0\) the generic parallelogram has area 36. This
constructs the complete periodic support in the original lattice.

\par

\medskip\noindent\textbf{I2I.2. Full finite parent and actual regular subdivision}

\par\nobreak

Triangulate the upper parallelogram at \(r=1,z=0\) by successively
inserting all its integral points into every containing triangle
and edge, starting with its two diagonal triangles. The exact
output has 72 determinant-one triangles and total area 36. The
boundary subdivisions agree under both original period sides.
At \(r=0\) the lower segment has endpoints \(A_{0}=(0,0,0,1)\) and
\(A_{1}=(0,-1,0,1)\). Join \(A_{0}\) to all upper triangles. Join \(A_{0},A_{1}\) to
each of the twelve consecutive horizontal unit edges from
\((-23,-3)\) to \((-11,-3)\), and to the remaining diagonal edge from
\((-11,-3)\) to \((0,0)\). These cones cover the hexagonal support on
each base slice \(r,z\ge 0\) and meet on their stated common faces.
The horizontal joins have determinant 1 and the diagonal join
has determinant 11.

For each singular cone with ordered primitive matrix B, enumerate
\(B^{-1} \mathbb Z^4/\mathbb Z^4\) by the residues of its four inverse-coordinate
columns, modulo its determinant. The resulting finite group has
exactly \(\operatorname{det} B\) elements. For each nonzero residue its vector in
the half-open fundamental parallelepiped is integral. Divide by
its gcd, choose a primitive ray, and subdivide every cone
containing it. A positive barycentric coefficient is strictly
less than 1, so every changed child\textquotesingle s determinant is a positive
integer smaller than its parent\textquotesingle s. The algorithm terminates.
Every step is the actual toric proper subdivision, with its
full integer character maps, and affects the complete incident
star. The certificate retains all selected rays and children.

For an original raw cone enumerate the primitive dual rays and
every integral vector in their half-open parallelepiped. They
generate the whole dual semigroup: subtract the nonnegative
coefficient floors of an arbitrary dual-cone lattice point.
The recorded monomials for every final regular cone are the
pairings of those generators with its primitive columns. They
are integral and nonnegative. Thus each chart has its full
holomorphic morphism to the original affine semigroup chart,
including the determinant-eleven raw parent. On the dense
original torus the inverse is the stated rational integer
character map. Every subdivision preserves the whole support,
so their assembled map is proper.

\par

\medskip\noindent\textbf{I2I.3. All original period faces, point modules and volume}

\par\nobreak

On the final four-dimensional fan the full face arrows are

\[
\begin{gathered}
L: x+23r=0 \longmapsto  x=0, \\
T: y+3r+z=0 \longmapsto  y=0, \\
L/T: 3x-11y+36r=0 \longmapsto  3x-11y-11z=0.
\end{gathered}
\]

Their lattice shifts are respectively \((23r,3r),(11r,3r+z)\),
\((12r,-z)\); their full unit columns for f are \((2,0),(1,0),(1,0)\).
The inverse arrows retain the opposite shifts and reciprocal
units. The executed boundary face sets are exactly paired by
these maps. At each regular chart their character matrix is
\(B_{\rm target}^{-1} T_{\rm period} B_{\rm source}\), with full unit column
\(B_{\rm target}^{-1} u_{\rm period}\). Its inverse has the opposite transported
column; the coefficient product cancels identically. This
specifies the full original toric quotient gluing.

In every chart the actual source-band functions are

\[
\begin{gathered}
x_{\rm source}=QZ=R^{3}u^{2} Z, y_{\rm source}=Z^{-1}, \\
x_{\rm source} y_{\rm source}=Q.
\end{gathered}
\]

Their ray orders are \(y+3r+2z\) and \(-y\). They are nonnegative on
all final cones; their sum is the full source order \(3r+2z\).
The full inverse Tate map therefore carries these source-node
charts to the original Weierstrass source. On this band the
actual original point coefficient is

\[
\begin{gathered}
v_{2}=\eta P^{2} XZ/f_{P^{2}}(QZ), P=Ru,
\end{gathered}
\]

with its reciprocal arrow in the stated Cartier point/O modules.
This is the full MP theta identity, including every \(\eta\), P and
source-band factor. Its subsequent point graph is the proper
graph of \((x_{\rm source},P^{2})\), to be attached through these very maps.

The physical circle has valuation vector \((1,0,0,0)\); each chart\textquotesingle s
four physical weights are its full inverse matrix applied to
that vector. The source Jacobian and invariant differential
remain their complete units. In the common coordinate volume,
\(\Omega=-(2J_{\rm ts}/a_T) P \mathrm dP\wedge \mathrm dp\wedge \mathrm d\log X\wedge \mathrm d\log Z\), substitute
\(P=Ru,p=u^{-1}\). Since \(P \mathrm dP\wedge \mathrm dp=-R \mathrm dR\wedge \mathrm du\), the full volume is

\[
\begin{gathered}
\Omega=(2J_{\rm ts}/a_T) R \mathrm dR\wedge \mathrm du\wedge \mathrm d\log X\wedge \mathrm d\log Z.
\end{gathered}
\]

Every ray therefore has original canonical order \(2r+z-1\), with
the actual nonzero coefficient retained. The change from the
coupled scalar\textquotesingle s line factor adds only source-base differentials
to \(\mathrm d\log V\), which wedge to 0 with the two displayed source forms.
The certificate checks every source exponent, inverse matrix,
physical weight and canonical order. No global characteristic
pair is assigned by this parent construction.

The executed certificate \path{code/part3_root/e6_selected_i3_m2_infinity_fan.py}
passes 598 exact checks. Its raw upper slice has 72 regular triangles;
its complete regular infinity parent has 149 four-cones after 18
recorded primitive stellar insertions. Every original semigroup
generator, full inverse monomial map, physical weight, source order,
canonical order and period-face set is retained in its JSON.

The \(\eta\) in every displayed point-module arrow is exactly
\(a_T \Theta'_Q(1) \Theta_Q(P^{2})/\Theta_Q(P)^{2}\), with \(\eta(0)=-1\)
and \(a_T(0)=1\). In \(c_{12}\) its full sixth power includes
\(\Theta'_Q(1)^{6}\). The source SI2 certificate retains this entire
identity and the reciprocal frame morphism.


% End expanded verified cp36 input


% Expanded verified cp36 input: notes/part3_e6_selected_i3_m2_return_support_07.tex
\par\bigskip\noindent\textbf{P3K.111 supporting proof 7: Complete generic common fan}\par\nobreak
\par

\medskip\noindent\textbf{Selected-I3 \(M_{2}\) N12: the complete generic common fan}

\par\nobreak

We retain SI2.1-\/-6 and the full period dictionary I2M.1-\/-3.
The original point line is \(M_{2}=\mathcal O(2G-O)\), and its source
translation is minus three G. This construction gives the
common generic boundary used by the three endpoint families.
The point-section graph extensions and their compact normal
ledger remain subsequent constructions.

\par

\medskip\noindent\textbf{I2F.1. The actual torus and the three parent morphisms}

\par\nobreak

Use the characters \((X_c,Z_c,P)=(V_{\rm reg}/Z,Z,P)\) and retain
the entire functions p and \(f_{0}\) on the punctured exceptional
parameter curve. The complete period maps are

\[
 L_c:(X_c,Z_c,P)\mapsto
 (P^{23}p^{-1}f_0^2X_c,P^3p^3Z_c,P),
\]
\[
 T_c:(X_c,Z_c,P)\mapsto
 (P^{11}p^{-1}f_0X_c,P^3p^2Z_c,P).
 \tag{I2F.1}
\]

They fix p and every coefficient in \(f_{0}\), so the maps commute
with their full analytic functions. The valuation lattice
has ordered basis \((\operatorname{ord} X_c,\operatorname{ord} Z_c,\operatorname{ord} P)\), and a period word
\(L_c^i T_c^j\) adds \(((23i+11j)h,(3i+3j)h,0)\) to a ray \((a,b,h)\).
Its full unit columns are

\[
 u_p(i,j)=(-i-j,3i+2j,0)^t,\qquad
 u_f(i,j)=(2i+j,0,0)^t.
 \tag{I2F.2}
\]

The original zero-parent character map is
\((X_c,Z_c)=(1/(X_{0}Y_{0}),Y_{0})\), with inverse
\((X_{0},Y_{0})=(1/(X_cZ_c),Z_c)\). The one-parent map is
\((X_c,Z_c)=(Y_{1},1/X_{1})\), with inverse
\((X_{1},Y_{1})=(1/Z_c,X_c)\). Their complete period changes
are the integral matrices I2M.2 and I2M.2, with the
written inverse matrices. At infinity the characters
already agree, while the transverse parameter is
the full function \(R=pP\). There \(f_{\rm infinity}=p^{-12} f_{0}\).
These identities give the full parent morphisms in
both directions on the common torus.

The generic period matrix is
\[
                   M=\begin{pmatrix}23&11\\3&3\end{pmatrix}.
 \tag{I2F.3}
\]
The polygon H with ordered vertices
(0,0),(-12,0),(-23,-3),(-11,-3) has area thirty-six.
I2I.1 proves that its translates by \(M \mathbb Z^{2}\) have
disjoint interiors and cover the plane. the zero-parent construction
and I2O.1-\/-3 give the other two fundamental domains
with the literal same period lattice through the
maps just specified. Each parent has seventy-two regular
upper triangles in its own lattice. Every one of
those triangles and its ordered rays remains in
the common-fan certificate.

\par

\medskip\noindent\textbf{I2F.2. Exact finite overlay and boundary closure}

\par\nobreak

For each zero-parent triangle, apply its character
matrix and every period translate that meets H.
The possible translate indices are bounded exactly
by the two coordinates \(M^{-1}(x,y)\): on each convex
polygon, their extrema occur at vertices. Taking
integer ceilings and floors of the differences of
the extremal intervals gives the complete finite
list. The code uses rational arithmetic for all
these bounds and intersections.

Intersect those translated triangles with the
seventy-two infinity-parent triangles in H. Then
intersect each nonempty polygon with every
translated one-parent triangle meeting it. Each
intersection is constructed by clipping against
the three oriented rational half-planes of the
triangle. The interpolation parameter is exactly
the ratio of the signed values of the defining
linear form at the two endpoints. This gives
6387 polygons with positive area; their areas sum
to thirty-six. Their interiors are disjoint because
each of the three original periodic triangulations
has that property.

Every polygon vertex is inserted in all incident
polygon edges. Its period images lying in H are
also inserted. The resulting vertex set has 6043
points. The two boundary identifications are
\[
 (a,b,h)\mapsto(a+11h,b+3h,h)
 \quad\hbox{from }b+3h=0\hbox{ to }b=0,
\]
\[
 (a,b,h)\mapsto(a+12h,b,h)
 \quad\hbox{from }3a-11b+36h=0
                  \hbox{ to }3a-11b=0.
 \tag{I2F.4}
\]
These are \(T_c\) and \(L_c T_c^{-1}\), with their full
coefficients I2F.1-\/-2. The certificate verifies
equality of the entire subdivided edge sets in
both identifications. These two words generate
the original period lattice.

Each convex intersection polygon is triangulated
by its ordered ears. A chosen ear has positive
oriented area and contains no other remaining
polygon vertex. Removing it preserves a polygon
with the same ordered boundary. The finite
recorded removals therefore partition the polygon
without overlapping interiors or an omitted edge.
All triangles are cones over primitive integer
rays: clear the two rational denominators and
divide by the gcd of the resulting triple.

\par

\medskip\noindent\textbf{I2F.3. The executed periodic regularization}

\par\nobreak

For a simplicial integer cone B of determinant d,
the coefficients of each ambient coordinate
vector in B give generators of the quotient
lattice \(B^{-1} \mathbb Z^{3}/\mathbb Z^{3}\). The code enumerates that
finite subgroup by exact addition modulo d,
and checks that it has precisely d elements.
Each nonzero residue gives an integral vector
in the half-open fundamental parallelepiped.
Its primitive ray and all its period images
in H are inserted in every containing cone.

The record retains every selected ray, all its
period images, the number of changed cones and
the selected determinant. On every changed cone
the new ray has nonnegative barycentric entries;
every positive entry is strictly less than one.
Replacing each positive-coordinate ray by the
inserted ray gives the ordinary stellar
subdivision. Its children cover the parent,
have disjoint interiors and agree on every
unchanged face. Their determinants are positive
integers strictly smaller than the parent\textquotesingle s.
The same operation on every incident cone and
period image retains the entire common boundary.

The exact regularization has completed: 62,281 recorded stellar
steps give 136,442 regular cones. Every final determinant is one.
There are 204,429 interior edges, each incident to exactly two cones,
and 468 boundary edges, each incident to exactly one cone before
period gluing. The exact height-one area is thirty-six. Both entire
period boundary edge sets agree under I2F.4. Every insertion history,
primitive height and original parent is retained in the certificate.
These computations establish the stated common generic atlas; no
numerical characteristic pair is assigned at this stage.

\par

\medskip\noindent\textbf{I2F.4. Every affine parent map with its complete units}

\par\nobreak

Let B be an ordered final ray matrix. Its
chart coordinates theta give
\[
                      (X_c,Z_c,P)=\theta^B.
 \tag{I2F.5}
\]
For each zero or one parent, the certificate
retains its translated ray matrix Bk, original
triangle index and exact word(i,j). The actual
parent coordinates are
\[
 \theta_{k,\alpha}=
 p^{-(B_k^{-1}u_p)_\alpha}
 f_0^{-(B_k^{-1}u_f)_\alpha}
 \prod_\beta\theta_\beta^{(B_k^{-1}B)_{\alpha\beta}}.
 \tag{I2F.6}
\]
Every entry of \(B_{k}^{-1} B\) is an integer at least
zero. All unit exponents are integers. The
affine formula therefore extends over the
complete child chart, and its inverse on the
original torus is the full character formula
I2F.1-\/-2 with the opposite word. The p and \(f_{0}\)
factors have their entire source-parameter
dependence in I2F.6.

For the infinity parent, let Bi denote its
original ordered ray matrix in (\(X_c\),\(Z_c\),R).
The literal parameter equality \(R=pP\) gives
\[
 \theta_{\infty,\alpha}=
 p^{(B_i^{-1}(0,0,1)^t)_\alpha}
 \prod_\beta\theta_\beta^{(B_i^{-1}B)_{\alpha\beta}}.
 \tag{I2F.7}
\]
Again all character exponents are nonnegative
integers, and the complete coefficient is a
unit on the generic exceptional domain. The
saved orientation permutation relates Bi to
the exact parent triangle order. These three
families of maps preserve every original
character. Character localizations give their
overlap morphisms, with the same full functions
and inverse period words. The subdivisions
cover each parent support; hence these affine
maps assemble to the proper toric refinements
of the three specified generic parent atlases.

\par

\medskip\noindent\textbf{I2F.5. Original source, line, physical circle and volume}

\par\nobreak

The original source node coordinates are
\[
 Z_1=pP^3Z_c,\qquad W_1=Z_c^{-1},\qquad
 Z_1W_1=Q=pP^3.
 \tag{I2F.8}
\]
On a ray \((a,b,h)\), their exponents are b+3h
and minus b. Both are nonnegative everywhere
on H. The coefficient p in the first coordinate
is retained. E81\textquotesingle s full inverse Tate map carries
these node coordinates to the original
Weierstrass source. The point-line morphism is
\[
 v_2=\frac{\eta P^2X_cZ_c}{f_{P^2}(pP^3Z_c)}.
 \tag{I2F.9}
\]
This follows from the full SI2 point identity and the full theta
identity \(f_{P^{2}}(QZ)=P^{2} f_{P^{2}}(Z)\). Both the
point and zero Cartier modules, including the
reciprocal line arrow, remain specified in
I2F.9 for their subsequent common graphs.

The point coefficient here is the actual Miller unit
\(\eta=a_T \Theta'_Q(1) \Theta_Q(P^{2})/\Theta_Q(P)^{2}\). Its value at the
source origin is −1. Both forward and reciprocal maps retain it.

The physical circle cocharacter is (1,0,0),
so its exact three chart weights are \(B^{-1}\)
times that vector. The program checks every weight is integral and retains every
primitive ray height. No finite-index
change of the physical circle has been made.

The full original source volume becomes
\[
 \Omega=-\frac{2J_{ts}}{a_T}\,
 P\,dP\wedge dp\wedge d\log X_c\wedge d\log Z_c.
 \tag{I2F.10}
\]
The minus sign is the determinant of
(\(X_{0}\),\(Y_{0}\)) to (\(X_c\),\(Z_c\)), namely minus one.
The complete coefficient \(2J_{\rm ts}/a_T\) is the
nonvanishing analytic source coefficient retained in SI2 and I2I.3. Inserting I2F.5 gives original
canonical order 2h-1 along a ray of height h;
each such integer is stored in the certificate.

The executing program is
\path{code/part3_root/e6_selected_i3_m2_common_fan.py}. It retains each
ordered matrix, the three parent morphisms, every full unit vector,
the periodic stellar history and the locked source files. Its completed
certificate is \path{code/part3_root/e6_selected_i3_m2_common_fan.json},
SHA256 \path{2999175af858e837e202908b15f4743fd7b5caf40af321673a05925ebbfe98cf}.
Every support, incidence, integral inverse and original source-map
gate has passed. All five original source locks were independently
recomputed after the certificate closed and matched. The full
corner lifts, global point graphs and compact normal attachments
remain the next calculations.


% End expanded verified cp36 input


% Expanded verified cp36 input: notes/part3_e6_selected_i3_m2_return_support_08.tex
\par\bigskip\noindent\textbf{P3K.111 supporting proof 8: Complete zero lift}\par\nobreak
\par

\medskip\noindent\textbf{Selected-I3 \(M_{2}\) N12: full zero-end lift of the common atlas}

\par\nobreak

This working proof retains the original source, selected height,
\(M_2=\mathcal O([2G]-[O])\), translation by \(-3G\), and
the cubic good-cap and four-root-turn maps of SI2. The source-compatible
zero parent is I2Z\textquotesingle s executed 420-chart fan. The common input is
\path{code/part3_root/e6_selected_i3_m2_common_fan.json}, with SHA256
\path{2999175af858e837e202908b15f4743fd7b5caf40af321673a05925ebbfe98cf}:
it contains 136,442 regular generic cones and all full parent labels.
The companion of this note computes the complete four-dimensional
lift and verifies its source, point, period and determinant maps.
No characteristic pair or return is assigned by this construction.

\par

\medskip\noindent\textbf{I2ZC.1. Actual characters and their full coefficients}

\par\nobreak

Use the original zero-end torus
\[
 (X_0,Y_0,P,p)=(V_{\rm reg}^{-1},Z,P,R/P),\qquad
 V_{\rm pt}=\eta V_{\rm reg},\qquad Q=P^3p.
\]
Here
\[
 f_0=(p-1)^{12}c_{12},\qquad
 c_{12}=\lambda_0c_q\eta^6/(u_s^{12}u_t^3),\qquad
 \eta=a_T\Theta'_Q(1)\Theta_Q(P^2)/\Theta_Q(P)^2.
\]
All these analytic units and their inverses are retained on the
zero-end domain. The original source functions are
\(s=P(p-1)/u_s\), \(t=Pu_t\), and \(q_{\rm old}=pPc_q\).

The zero period basis is \(L_0=B^{-1}\), \(T_0=hB^{-1}\):
\[
 L_0:(X_0,Y_0)\mapsto(P^{10}p f_0X_0,P^3Y_0),\qquad
 T_0:(X_0,Y_0)\mapsto(P^8p f_0X_0,P^6pY_0).
 \tag{I2ZC.1}
\]
Both keep \(P,p\). Their inverse functions use the reciprocal
coefficients at this same base point. In the common torus
\((X_c,Z_c)=(V_{\rm reg}/Z,Z)\), the complete periods are
\[
 L_c=(P^{23}p^{-1}f_0^2X_c,P^3p^3Z_c),\qquad
 T_c=(P^{11}p^{-1}f_0X_c,P^3p^2Z_c).
\]
For a saved common-parent word \(L_c^iT_c^j\), its full map is
\[
 \begin{aligned}
 X_c&=P^{23i+11j}p^{-i-j}f_0^{2i+j}X_0^{-1}Y_0^{-1},\\
 Z_c&=P^{3i+3j}p^{3i+2j}Y_0.
 \end{aligned} \tag{I2ZC.2}
\]
Thus the integral four-character matrix is
\[
 F_{i,j}=\begin{pmatrix}
 -1&-1&23i+11j&-i-j\\
 0&1&3i+3j&3i+2j\\
 0&0&1&0\\0&0&0&1
 \end{pmatrix},\qquad\det F_{i,j}=-1,
 \tag{I2ZC.3}
\]
with full \(f_0\)-column \((2i+j,0,0,0)^t\). The inverse is
the literal inverse matrix and reciprocal transported coefficient;
both are exported. On a generic ray \((x,y,h)\), the native ray is
\[
 (-x-y+(26i+14j)h,\ y-(3i+3j)h,\ h).
 \tag{I2ZC.4}
\]
The original physical vector \((-1,0,0,0)\) maps to
\((1,0,0,0)\). The genuine point coefficient acquires exactly
\(P^{26i+14j}p^{2i+j}f_0^{2i+j}\); its inverse retains every
power. In particular the p-character in this formula is carried
through the zero-end collar as a divisor character.

\par

\medskip\noindent\textbf{I2ZC.2. The relative four-cone construction}

\par\nobreak

For every saved common cone, I2ZC.4 and its recorded parent word
give a primitive regular native three-cone. The inverse of its
native upper-parent matrix sends all three columns to nonnegative
integers. The certificate checks uniqueness and the full rational
area of every one of the 72 native upper triangles.

The existing 420-chart zero parent already includes both actual
source bands and all periodic source-cut subdivisions. Its rays
are in \((\operatorname{ord}X_0,\operatorname{ord}Y_0,
\operatorname{ord}P,\operatorname{ord}p)\). The new relative
fan is constructed inside each such parent as follows. An upper
three-face is replaced by its pulled-back common triangles and
joined to the same actual lower ray. An upper two-face is replaced
by every common edge it contains and joined to its same two actual
lower rays. Parents with upper dimension at most one retain their
original rays. All lower rays are those of the executed source-cut
parent, including their nonzero original P-heights.

For each resulting four-cone the certificate computes its full
matrix \(N\), its integral inverse, and its determinant. It orders
the columns with determinant one. With parent matrix \(N_0\),
the coefficient map is literally
\[
 z_{0,a}=\prod_b z_b^{(N_0^{-1}N)_{ab}}.
 \tag{I2ZC.5}
\]
Every exponent is a nonnegative integer. On a common-face
localization the inverse is the matrix \(N^{-1}N_0\); all
coordinates that occur with negative exponents there are
invertible on that localization. The records retain both matrices.
These maps prove the asserted refinement to this actual regular
420-chart parent.

For each parent, the sums of the child slice volumes agree exactly.
For a regular cone on rays \(v_a\), its volume on
\(\operatorname{ord}P+\operatorname{ord}p=1\) is
\[
 \frac1{6\prod_a(v_{a,P}+v_{a,p})}.
 \tag{I2ZC.6}
\]
The common-face incidence calculation checks that every internal
three-face occurs twice and every boundary three-face once. The
generic boundary is exactly the entire pulled-back common fan.
The two original lower facets remain. These support and incidence
equalities give the proper chart gluing, with no additional face
selected by its characteristic value.

\par

\medskip\noindent\textbf{I2ZC.3. Full periodic source and line maps}

\par\nobreak

The four-dimensional valuation maps of I2ZC.1 are
\[
 H_L=\begin{pmatrix}1&0&10&1\\0&1&3&0\\0&0&1&0\\0&0&0&1\end{pmatrix},
 \quad
 H_T=\begin{pmatrix}1&0&8&1\\0&1&6&1\\0&0&1&0\\0&0&0&1\end{pmatrix}.
\]
The paired parent boundaries are those for \(L_0,T_0,L_0T_0^{-1}\),
with their actual full source and target equations. For a paired
source cone \(N_s\), target cone \(N_t\), and saved word
\((a,b)\), the complete map is
\[
 z'_\alpha=
 f_0^{[N_t^{-1}(a+b,0,0,0)^t]_\alpha}
 \prod_\beta z_\beta^{[N_t^{-1}H_L^aH_T^bN_s]_{\alpha\beta}}.
 \tag{I2ZC.7}
\]
Its inverse uses \(N_s^{-1}H_T^{-b}H_L^{-a}N_t\) and the
coefficient \(N_s^{-1}(-a-b,0,0,0)^t\). Since the periods fix
\(P,p\), these inverse unit columns cancel pointwise. The
certificate checks both localized exponent maps and the complete
unit inverse for every paired face, and accounts for the entire
remaining boundary by these periods and the two lower facets.

Every child retains its parent\textquotesingle s actual source band \(j=1\) or
\(j=2\). The original source functions on that band are
\[
 x_j=Q^jY_0,\qquad y_j=Q^{1-j}Y_0^{-1},\qquad x_jy_j=P^3p.
 \tag{I2ZC.8}
\]
Their exponent rows on a native ray \((x,y,r,z)\) are
\(y+3jr+jz\) and \(-y+3(1-j)r+(1-j)z\). Both rows are
nonnegative on every child. Their sum is exactly \(3r+z\).
The original analytic functions \(s,t,q_{\rm old}\) above are
then carried in these same charts and period squares.

The marked point graph remains the literal graph of
\((x_j,P^2)\), with its original two point-module generators
\((a-1)^{-1}\), \((c-1)^{-1}\), and the zero generator
\(x_j-1\). Its coefficient formulas are the actual I2P maps:
\[
 V_{\rm reg}=-\eta^{-1}P^{-2j}EA
 =\eta^{-1}P^{-2j}c^{-1}EC,\quad C=-cA,
\]
\[
 V_{\rm reg}=\eta^{-1}P^{-2j-2}(x_j-P^2)E_OC_O.
 \tag{I2ZC.9}
\]
Their pullbacks by I2ZC.5 are retained in every child record.
The finite marked comparison graphs use these actual generators
and the corresponding inverses; the full coefficient \(\eta\)
includes its theta derivative. The graph and coefficient arrows
give the original point-line boundary map used in the later
compact comparison.

\par

\medskip\noindent\textbf{I2ZC.4. Original determinant and exact scope}

\par\nobreak

The original four-form in the zero marking is
\[
 \Omega_{\rm orig}=
 \frac{2J_{ts}}{a_T}P\,dP\wedge dp\wedge
 d\log X_0\wedge d\log Y_0.
 \tag{I2ZC.10}
\]
The displayed sign and full nonzero unit are retained. On a ray
\((x,y,r,z)\), its canonical order is \(2r+z-1\). This is
recorded on every child together with the physical weights
\(N^{-1}(-1,0,0,0)^t\). The full derivatives of the base units
in I2ZC.2 and I2ZC.7 wedge to zero against the two base
differentials; their nonzero multiplicative units remain in the
actual chart and normal maps. No Spin extension is deduced by
discarding a canonical divisor.

The companion exports the complete native generic pullbacks,
four-cones, parent/inverse matrices, source rows, point/O modules,
period words and full coefficient columns. Input hashes are
verified before and after execution. Its exact counts and pass
status are recorded in its output; the earlier 72-triangle
preflight remains separately labelled as a preflight. The fixed
normal and moving-cap calculations use the actual new output and
are subsequent constructions.


% End expanded verified cp36 input


% Expanded verified cp36 input: notes/part3_e6_selected_i3_m2_return_support_09.tex
\par\bigskip\noindent\textbf{P3K.111 supporting proof 9: Complete infinity and one lifts}\par\nobreak
\par

\medskip\noindent\textbf{Selected-I3 \(M_{2}\) N12: full infinity and one-end lifts}

\par\nobreak

\emph{Historical construction record.} All pending-execution statements in this supporting proof describe its preparation stage. The completed 145160-chart infinity and 149774-chart one-end execution, including the full-object comparison, is recorded in \path{records/e6_selected_i3_m2_corner_complete_equivalence.json}. The completed common input, compact coverage and return are I2F, CFL and ASI2 below.

Execution status: the exact common input is still running. The complete
coordinate construction below is prepared for that input; its final counts,
incidence gate and executed certificate are pending. No characteristic
pair or completed global return is assigned.

We retain every ray and full function of I2F.1-\/-10.
The generic boundary is the entire I2F common regular fan.
This construction extends that exact boundary across
the infinity and one endpoints. The zero endpoint
is the independently constructed I2ZC.1-\/-4 atlas.
All three use the original \(M_{2}\) point line and
translation by minus three G. Point-section graphs
and the compact return ledger remain further work.

\par

\medskip\noindent\textbf{I2L.1. The literal generic boundary at each endpoint}

\par\nobreak

At infinity the native characters are already
(X,Z,R) with \(X=X_c,Z=Z_c\). Attach the ray
\(A_{0}=(0,0,0,1)\) to every common three-cone, in the
ordered lattice (ord X,ord Z,ord R,ord u).
The full parameter arrow is
\[
 (X_c,Z_c,P,p)=(X,Z,Ru,u^{-1}),\qquad
 f_0=u^{-12}f_\infty.
 \tag{I2L.1}
\]
Its inverse is \(R=pP,u=p^{-1}\), with X,Z fixed.
All these equalities are the entire holomorphic
functions on their stated parameter overlaps.

At one, take a common cone with recorded one-parent
word(i,j). Its common ray(a,b,h) has native ray
\[
 \bigl(-b+(3i+3j)h,\ a-(23i+11j)h,\ h\bigr)
 \tag{I2L.2}
\]
in (ord \(X_{1}\),ord \(Y_{1}\),ord P). The common-input gate checks that these cones are
distinct, have total area thirty-six, and have
precisely the edge incidence of the native
parallelogram I2O.1 at z zero. The common period
edge identifications I2F.4 give the same
subdivisions from the two incident parent
triangles. Thus I2L.2 unwraps the actual common
quotient into the complete native fundamental
parallelogram. Join each cone to \(E_{0}=(0,0,0,1)\).

The full one-end character arrow, including the
additional \(v=p-1\) powers, is
\[
 \begin{pmatrix}\log X_c\\\log Z_c\\\log P\\\log v\end{pmatrix}
 =
 \begin{pmatrix}
 0&1&23i+11j&12(2i+j)\\
 -1&0&3i+3j&0\\
 0&0&1&0\\0&0&0&1
 \end{pmatrix}
 \begin{pmatrix}\log X_1\\\log Y_1\\\log P\\\log v\end{pmatrix}
 +u_p\log p+u_1\log c_{12},
 \tag{I2L.3}
\]
where
\(u_p=(-i-j,3i+2j,0,0)^t\) and
\(u_1=(2i+j,0,0,0)^t\).
This logarithmic notation denotes the literal
integer monomial and unit product; it introduces
no logarithm choice. The inverse character
matrix is integral. Every shared edge retains
all its incident common-owner words. Differences
of those words give exactly the full I2F.1
period arrows on the common open part.

\par

\medskip\noindent\textbf{I2L.2. The finite lower joins and their subdivisions}

\par\nobreak

At infinity, extract the common boundary intervals
on \(y+3h=0\) and on \(3x-11y=0\). Join each interval to
\(A_{0},A_{1}=(0,-1,0,1)\). On the first edge the joins
have determinant one; on the second they have
determinant eleven. Their support is precisely the
lower part of I2I.1, adjoining the upper stars
of I2L.1. Every join lies in one of the original
thirteen lower parents of I2I.2, and its parent is
determined by testing its four rays against
that parent\textquotesingle s inverse ray matrix.

At one, extract every upper boundary interval on
\(3y-x-36h=0\). For each \(j=0,...,11\) join it to
\(E_j=(0,j,0,1)\) and \(E(j+1)\). Every such join has
determinant three and lies in the actual parent
I2O.2 indexed by j. They cover the extrusion of
the native top edge through the full twelve-unit
v height. Together with the upper stars they
have exactly the support I2O.1.

For each singular lower cone, enumerate the
integral points of its half-open fundamental
parallelepiped, divide each nonzero vector by
its gcd, and insert the chosen primitive ray
in every containing lower cone. Each positive
barycentric coefficient is strictly less than
one on every changed cone; every child has
positive determinant strictly smaller than its
parent. All coefficients, old cones and children
are retained in chronological order.

Every new infinity ray satisfies either
\(y<-3h\) or \(3x-11y>0\). Every new one-end ray satisfies
\(3y-x>36h\). These strict inequalities put them
outside every upper star, whose corresponding
inequalities are weak in the opposite direction.
The generic boundary is therefore unchanged by
all these lower subdivisions.

The executed counts of lower stellar insertions and regular four-cones
will be inserted after both complete endpoint certificates pass. All ordered ray
matrices have determinant one and all cones
are distinct. On each original raw parent,
the sum of its child simplex volumes in
\(h+t=1\) is exactly its volume, using
\[
 \operatorname{vol}(B)
 =\frac{\det B}{6\prod_{v\text{ a ray of }B}(v_h+v_t)}.
 \tag{I2L.4}
\]
The explicit stellar operations establish
non-overlap and support; I2L.4 independently
checks the complete coverage of every parent.

\par

\medskip\noindent\textbf{I2L.3. Full affine morphisms to the original raw parents}

\par\nobreak

Write B for a final ordered four-ray matrix.
The child chart is the actual monomial chart
\((X,Z,R,u)=\theta^B\) at infinity, or
\((X_{1},Y_{1},P,v)=\theta^B\) at one. Its inverse ray
matrix has integral entries. For the raw
parent retain every primitive dual ray and
every integral half-open-parallelepiped
representative listed in I2I.2 or I2O.2.
They generate its entire dual semigroup by
subtracting the four coefficient floors from
any lattice point in the dual cone.

For each retained dual-semigroup generator m,
its exact pullback is \(\theta^{mB}\). The
certificate verifies every exponent is an
integer at least zero. Equal character sums
have equal products after this substitution,
so every binomial relation of the parent
semigroup is preserved. This is the full
affine morphism to the original parent,
including each determinant-eleven or
determinant-three singular parent. The
identity of characters gives the overlap
localizations and their inverse maps.

These finite subdivisions cover each raw
parent support; the resulting toric morphisms
are proper. This identifies their precise
targets. The earlier 149-cone and 168-cone
resolutions retain their own recorded maps
to the same raw parents. The construction
here does not use an unverified affine
morphism between the two regular refinements.

\par

\medskip\noindent\textbf{I2L.4. Every period face and its complete inverse}

\par\nobreak

At infinity use the actual lattice matrices
L and T of I2I.3, with full coefficient
\(f_{\rm infinity}^{2i+j}\) in the X coordinate for
the word \(L^iT^j\). The entire opposite-face
pairings are
\[
 L:\ x+23h=0\longrightarrow x=0,\qquad
 T:\ y+3h+t=0\longrightarrow y=0,
\]
\[
 L/T:\ 3x-11y+36h=0\longrightarrow3x-11y-11t=0.
 \tag{I2L.5}
\]
At one use I2O.3. The two pairings are
\[
 L:\ x=0\longrightarrow x-3h=0,\qquad
 T:\ 3y-x=0\longrightarrow3y-x-36h-36t=0.
 \tag{I2L.6}
\]
Their complete character-unit columns for
\(L^iT^j\) are \(p^{i-j,i,0,0}\) and
\(c_{12}^{0,j,0,0}\).

For each paired face, the record specifies
the unique adjacent source and target charts
Bs,Bt, the actual lattice map N, and the
integer matrix \(E=B_{t}^{-1} N B_{s}\). Every column
belonging to a shared face ray is exactly
the corresponding target coordinate vector.
Thus all negative entries occur only in
coordinates inverted on the stated overlap.
The inverse \(E^{-1}\) is integral and has the
same property in the reverse direction.
For every full character-unit column u the
actual unit exponent vector is \(B_{t}^{-1} u\);
the inverse is \(-B_{s}^{-1} N^{-1} u\). Multiplication
by \(E^{-1}\) verifies their complete cancellation,
including all parameter-dependent units.

The complete internal, generic and period-facet counts are pending
the executed arrays. Every internal
facet belongs to two cones. Every other
facet is either one of the prescribed
generic cones or one of I2L.5-\/-6; there is
no unpaired artificial boundary. The final incidence and opposite-face checks certify these
literal localized morphisms on the two complete atlases.

\par

\medskip\noindent\textbf{I2L.5. Source, point-line, physical weights and orientation}

\par\nobreak

At infinity the original source coordinates
are \(Z_{0}=R^{3}u^{2}Z,W_{0}=Z^{-1}\). Their exponent rows
on a ray (x,y,h,t) are y+3h+2t and minus y.
At one they are \(Z_{0}=pP^{3}/X_{1},W_{0}=X_{1}\), with
rows \(3h-x\) and x and the retained p unit
in Z0. Every displayed row is nonnegative
on every complete cone. Their products
are the actual \(Q=R^{3}u^{2}\) or \(Q=pP^{3}\),
respectively. The full inverse Tate map
of E81 gives the original Weierstrass
source in these charts.

The point-line maps are exactly
\[
 v_2=\frac{\eta P^2XZ}{f_{P^2}(R^3u^2Z)}
 \quad\hbox{and}\quad
 v_2=\frac{\eta P^2Y_1}{X_1 f_{P^2}(pP^3/X_1)}.
 \tag{I2L.7}
\]
They retain \(\eta=a_T*\Theta'_Q(1)*\Theta_Q(P^{2})/\Theta_Q(P)^{2}\),
its reciprocal and central value minus one, both point/O
Cartier generators and the reciprocal
arrow. Their normalized section graphs
will supply the subsequent fixed-normal
construction.

The physical cocharacters are (1,0,0,0)
at infinity and (0,1,0,0) at one. All
chart weights are the exact integer
vectors \(B^{-1} \delta\). The original volume is
\[
 \Omega=\frac{2J_{ts}}{a_T}\,
 dP\wedge dR\wedge d\log Z\wedge d\log V_{2T}.
 \tag{I2L.8}
\]
At infinity X=\(V_{2,T}\) and \(P=Ru\). The base
identity \(\mathrm dP\wedge \mathrm dR=-R \mathrm dR\wedge \mathrm du\)
and the character swap each contribute
a minus sign, giving
\[
 \Omega=+\frac{2J_{ts}}{a_T}R\,
 dR\wedge du\wedge d\log X\wedge d\log Z.
 \tag{I2L.9}
\]
At one \(R=(1+v)P\), \(Z=X_{1}^{-1},V_{2,T}=Y_{1}\);
the character map has determinant minus
one and the base Jacobian is P. Thus
\[
 \Omega=-\frac{2J_{ts}}{a_T}P\,
 dP\wedge dv\wedge d\log X_1\wedge d\log Y_1.
 \tag{I2L.10}
\]
The coefficient \(2J_{\rm ts}/a_T\) is the entire
nonvanishing source unit. Both canonical
orders are exactly \(2h+t-1\), retained on
every ray. The signs are computed from the displayed base and character
Jacobians of this original selected-I3 source.

The prepared executable is
\path{code/part3_root/e6_selected_i3_m2_corner_lifts.py}.
It records both full atlases, every raw-parent semigroup map,
every period face map and inverse, all full coefficient columns,
physical weights and original canonical orders. The numerical
certificate, source locks and counts will be recorded after the
common fan closes and the two complete lifts execute. The point
and O graph attachment and compact normal ledger remain subsequent
steps of the selected-I3 \(M_{2}\) calculation.


% End expanded verified cp36 input


% Expanded verified cp36 input: notes/part3_e6_selected_i3_m2_return_support_10.tex
\par\bigskip\noindent\textbf{P3K.111 supporting proof 10: Generic curve normal maps}\par\nobreak
\par

\medskip\noindent\textbf{Selected-I3 \(M_{2}\): actual generic fixed curves and all three normal clutches}

\par\nobreak

This proof concerns the generic one-dimensional fixed components in the
executed N12 three-end completion. It uses the full 136,442-cone common
fan, the actual zero common lift and the root\textquotesingle s one/infinity lifts. The
companion records every curve separately, with its three original
character morphisms, their inverses, all base units, the point-line
arrow and the original canonical character. It does not assign a
global return or include any moving, lower-reference or compact
surface contribution.

\par

\medskip\noindent\textbf{I2CV.1. Components and actual end charts}

\par\nobreak

Let N be the matrix whose columns are an ordered regular common
three-cone. The physical vector in the original common lattice
\((X_c,Z_c,P)\) is \((1,0,0)\). The normal weights are therefore exactly
\(w=N^{-1}(1,0,0)\). A cone with all three weights nonzero gives a fixed
curve. The executed minimal-face/period inventory supplies its
component identifier, a single common chart occurrence and trivial
period stabilizer. These assertions are checked anew for every
component by the companion, against the complete fixed-face output.

The base is the actual exceptional parameter sphere with coordinate p,
and its three specified discs have coordinates p, \(v=p-1\) and \(u=p^{-1}\).
For every common cone the three completed endpoint atlases each give
one four-cone owner. Three of its columns lie over height zero for the
disc coordinate. Its remaining column is the actual transverse ray \(l_e\).
The companion checks that its base multiplicity is one. In particular,
the zero source-band subdivision can give \(l_e\) positive original
P-height. This column is retained in the calculation below.

Write \(B_e\) for the full endpoint ray matrix, \(F_e\) for its exact three-row
character map to the common chart, and

\[
                         E_e=N^{-1}F_eB_e.
\tag{I2CV.1}
\]

The three generic columns of \(E_e\) are the actual permutation matrix
identifying the common normals. The fourth column is denoted \(c_e\).
Thus, in the full chart arrows, the common normal coefficients are
the endpoint coefficients multiplied by the base monomials with
orders \(c_e\) and the units listed in I2CV.2. The physical derivative
square \(E_e w_e=w\) is checked including the zero weight of the base
column. No normal line is identified merely by its weight.

\par

\medskip\noindent\textbf{I2CV.2. Every period word, unit and inverse}

\par\nobreak

For the zero end let (i,j) be its recorded common L,T translation.
In native coordinates \((X_0,Y_0,P,p)\) its complete character matrix is

\[
 F_0=\begin{pmatrix}
 -1&-1&23i+11j&-i-j\\
 0&1&3i+3j&3i+2j\\
 0&0&1&0
 \end{pmatrix},
\tag{I2CV.2}
\]

and the full additional coefficient is \(f_0^{2i+j}\) in its first
component. Here
\[
 f_0=(p-1)^{12}\lambda_0c_q\eta^6/(u_s^{12}u_t^3),\qquad
 \eta=a_T\Theta'_Q(1)\Theta_Q(P^2)/\Theta_Q(P)^2.
\tag{I2CV.3}
\]
All of these functions are those of the actual SI2 source. The
inverse full four-character matrix and its inverse \(f_0\) powers are
stored by the zero lift and carried verbatim into each curve record.
Both actual source bands of that lift remain recorded.

For the one end, in \((X_1,Y_1,P,v)\), the full character matrix is
\[
 F_1=\begin{pmatrix}
 0&1&23i+11j&12(2i+j)\\
 -1&0&3i+3j&0\\
 0&0&1&0
 \end{pmatrix}.
\tag{I2CV.4}
\]
The remaining units have exponent vectors
\((-i-j,3i+2j,0)\) for p and \((2i+j,0,0)\) for \(f_1\),
with \(f_0=v^{12}f_1\). The full four-character inverse has rows
\[
 \begin{pmatrix}
 0&-1&3i+3j&0\\
 1&0&-23i-11j&-12(2i+j)\\
 0&0&1&0\\0&0&0&1
 \end{pmatrix},
\tag{I2CV.5}
\]
and inverse unit vectors \((3i+2j,i+j,0,0)\) for p and
\((0,-2i-j,0,0)\) for \(f_1\). Multiplication checks both the
character identity and the cancellation of each entire unit vector.

At infinity the native height is R and \(P=Ru\). The full map is
\[
 (X_c,Z_c,P)=(X_\infty,Z_\infty,Ru),\qquad
 F_\infty=\begin{pmatrix}1&0&0&0\\0&1&0&0\\0&0&1&1\end{pmatrix}.
\tag{I2CV.6}
\]
Its inverse is \(R=pP\), \(u=p^{-1}\). Thus the extra u in the original
P-height is retained in every normal order. All three maps are
actual holomorphic monomials in their recorded regular chart,
with the displayed extending units; the complete parent semigroup
checks were made in the endpoint construction.

The point-line coefficient is also carried through every map. In
the actual source band b,
\[
 v_2=\eta P^{2b}V_{\rm reg}/f_{P^2}(Q^bZ).
\tag{I2CV.7}
\]
The I2P point and origin module frames extend this equation at its
zeros and poles. In particular the extra origin factor \(P^{-2}\)
is retained. The source-period coefficient of (i,j) is
\(P^{3i+3j}p^{3i+2j}\); the point coefficient in comparison to common
band one is
\[
 P^{26i+14j+2-2b}p^{2i+j}f_0^{2i+j}
 \frac{f_{P^2}(Q^bZ_{\rm native})}{f_{P^2}(QZ_{\rm common})}.
\tag{I2CV.8}
\]
Its inverse is the reciprocal at the literal inverse source
coordinate. Both point/O module generators and the selected-height
unit remain in the chart-to-boundary arrow. The companion checks
the whole period vectors before using their normal derivatives.

\par

\medskip\noindent\textbf{I2CV.3. Signed line degrees and original determinant}

\par\nobreak

The direction of I2CV.1 is endpoint coefficient to common
coefficient. Consequently the three actual normal line degrees on
the parameter sphere are
\[
                              d_i=-\sum_{e=0,1,\infty}c_{e,i}.
\tag{I2CV.9}
\]
This is the divisor calculation in the three explicit parameter
discs: a common coefficient equal to \(t_e^{c_{e,i}}\) times the
regular endpoint coefficient has frame transition with the
reciprocal power. The full unit factors contribute through the
same boundary maps. They extend on their respective discs with
their fixed nonzero constants, so their exact based logarithms
give degree zero there. They are retained in each exported map,
including the p and \(f_1\) factors at the one end.

Let \(h_i\) be the common primitive heights and \(r_e\) the original
native height of the actual lower ray. The full character arrows
give the identity
\[
                     \sum_i h_i d_i=-1-\sum_e r_e.
\tag{I2CV.10}
\]
The original four-form has its original base-blowup factor. In
these charts its generic order is \(2h_i-1\), and the lower-ray
contribution is \(2r_e\). Therefore its restriction to the completed
curve is exactly
\[
 \sum_i(2h_i-1)d_i+2\sum_e r_e=-2-\sum_i d_i.
\tag{I2CV.11}
\]
The right side is the actual adjunction degree on this parameter
sphere. The companion checks this equality independently from
the full normal monomial maps. This preserves the original
determinant map and every positive lower source height.

\par

\medskip\noindent\textbf{I2CV.4. Scoped characteristic residues}

\par\nobreak

The fixed curve is the actual parameter sphere and its three
complex normal lines have the weights \(w_i\) and degrees \(d_i\) above.
Apply the orbifold Stokes expression SZ.57--59 / IE.16 with the
physical line \(M=U^{-1}\). Define
\[
 W=\prod_i w_i,\quad S=\sum_i d_i/w_i,\quad
 Q_w=\sum_iw_i^2,\quad T=\sum_iw_id_i.
\]
Expanding the stated universal expression to first order on the
curve gives its physical M contribution
\[
               k_M=S/W,\qquad \ell_M=(Q_wS-2T)/W,
 \qquad A_M=(4k_M-\ell_M)/24.
\tag{I2CV.12}
\]
The companion independently expands the universal-U numerator
and verifies the Hermitian-dual sign in I2CV.12 for every curve.
All rational arithmetic is exact. It records each triple of
weights and degrees, all three endpoint maps and inverses, the
full point/source coefficient vectors and the original volume
identity before summing this scoped curve contribution.

The final machine result is source and input locked at
\path{code/part3/e6_selected_i3_m2_fixed_curve_normals.json}. It includes
all generic curve components certified by the complete face
inventory. Compact surfaces, endpoint interiors, lower reference
doubles, moving caps and the outer physical transfer remain
separately recorded constructions.


% End expanded verified cp36 input


% Expanded verified cp36 input: notes/part3_e6_selected_i3_m2_return_support_11.tex
\par\bigskip\noindent\textbf{P3K.111 supporting proof 11: Whole compact surface/curve/corner ledger}\par\nobreak
\par

\medskip\noindent\textbf{Selected-I3 \(M_{2}\): complete compact fixed-component ledger}

\par\nobreak

This proof concerns the actual N12 periodic completion of SI2, with
\(M_2=\mathcal O([2G]-[O])\), translation by \(-3G\), the original
selected height \(q=t-t_c(s)\), and the cubic good cap and four-root-turn
reference retained. Its three endpoint atlases are I2ZC and I2L, glued
by the entire I2F common fan. The full characteristic contribution
of their compact generic and corner-interior fixed components is

\[
                  (k_M,\ell_M)_{\mathrm{compact}}=(2723,19124),
 \qquad A_{M,\mathrm{compact}}=-343.
 \tag{CFL.1}
\]

Here \(M=U^{-1}\) is the original physical nonzero-vector line.
The twelve lower/reference surfaces, the selected moving cap with
its whole-reference fan comparison, and the original outer
parameter transfer are the remaining pieces of the full boundary
ledger. CFL.1 is the contribution of the compact strata specified
below.

\par

\medskip\noindent\textbf{CFL.1. The actual three-ended family and fixed components}

\par\nobreak

Retain the complete source functions

\[
 Q=P^2R,\quad R-P=s u_s,\quad t=P u_t,\quad q=R c_q,
 \quad p=R/P,\quad
 \eta=a_T\Theta'_Q(1)\frac{\Theta_Q(P^2)}{\Theta_Q(P)^2},
\]
\[
 c_{12}=\frac{\lambda_0c_q\eta^6}{u_s^{12}u_t^3},
 \qquad f_0=(p-1)^{12}c_{12}.
 \tag{CFL.2}
\]

The full \(\eta\) and its reciprocal remain in the coefficients;
its central value is \(-1\), while \(a_T(0)=1\). On the common
torus set
\((X_c,Z_c,P)=(V_{\mathrm{reg}}/Z,Z,P)\), with
\(V_{\mathrm{pt}}=\eta V_{\mathrm{reg}}\). Its actual periods are

\[
 L_c:(X_c,Z_c,P)\longmapsto
 (P^{23}p^{-1}f_0^2X_c,P^3p^3Z_c,P),
\]
\[
 T_c:(X_c,Z_c,P)\longmapsto
 (P^{11}p^{-1}f_0X_c,P^3p^2Z_c,P).
 \tag{CFL.3}
\]

Both fix all the functions in CFL.2. Thus the word \(L_c^iT_c^j\)
has full unit vectors
\((-i-j,3i+2j,0)\) for \(p\) and \((2i+j,0,0)\) for \(f_0\).
Its original lattice shift is
\(((23i+11j)h,(3i+3j)h,0)\).

I2F constructs the complete regular common fan from the three
specified parent decompositions. Its 136,442 cones cover the
original area -36 period domain; all 204,429 interior edges and
468 boundary edges have their recorded incidence and full period
arrows. The two opposite-face words generate the original period
lattice. Every common affine coordinate is the primitive monomial
of its ordered determinant-one ray matrix. The parent maps
I2F.6--7 include the entire \(p,f_0\) coefficients and their inverse
characters.

The zero, one and infinity families are the proper subdivisions
I2ZC and I2L of their stated raw endpoint parents. Their recorded
semigroup maps pull every parent generator back to a monomial
with nonnegative integer exponents. Equal sums of parent
characters have identical products after pullback, so all parent
relations hold. The support and facet checks give the exact
gluing to the common fan, with every other boundary accounted
for by an original period face or a labelled lower face.
The full zero construction retains both source bands. Infinity
retains \(P=Ru\), and one retains \(f_0=v^{12}c_{12}\), \(v=p-1\).

For an ordered local ray matrix \(B\), the actual circle weights
are \(B^{-1}\delta\). The vectors are
\(\delta_c=(1,0,0)\), \(\delta_0=(-1,0,0,0)\),
\(\delta_1=(0,1,0,0)\), and \(\delta_\infty=(1,0,0,0)\)
in their respective named lattices. A fixed coordinate stratum
sets precisely the nonzero-weight coordinates to zero. Its
normal face is their span. Localization along a common toric
face gives the exact overlap of these fixed strata. Across a
period facet the full chart map and its inverse identify the
face, its tangent coordinates and its normal derivative. The
complete fixed-face inventory performs these identifications,
preserving every period word and stabilizer.

The components used here have trivial period stabilizer, checked
from their full orbit cycles. There are 106,166 generic fixed
curves and 15,138 generic fixed surfaces. The corner-interior
components are listed in CFL.8. Their complete stars close in the
three executed atlases. The twelve one-end components with a
lower source boundary remain separately labelled in that inventory.

\par

\medskip\noindent\textbf{CFL.2. Every full endpoint normal morphism}

\par\nobreak

Let \(N\) be the ordered matrix of a common three-cone, and let
\(B_e\) be its actual four-cone owner at \(e=0,1,\infty\).
Its three generic columns have disc-coordinate height zero;
its fourth column is the original lower ray \(\lambda_e\),
of disc multiplicity one. The common-character map \(F_e\) gives

\[
                         E_e=N^{-1}F_e B_e.
 \tag{CFL.4}
\]

The three generic columns are the recorded coordinate permutation.
Write \(c_e\) for its lower column. The full local morphism is

\[
 z_i=\varepsilon_{e,i}t_e^{c_{e,i}}z_{e,i},\qquad
 z_{e,i}=\varepsilon_{e,i}^{-1}t_e^{-c_{e,i}}z_i,
 \tag{CFL.5}
\]

on its specified coefficient localization. The matrices and
units extend by the endpoint semigroup maps at the regular
coordinate boundary. The base-disc coordinates are \(p,v=p-1\),
and \(u=p^{-1}\).

For a zero owner with common word \((i,j)\),

\[
 F_0=
 \begin{pmatrix}
 -1&-1&23i+11j&-i-j\\
 0&1&3i+3j&3i+2j\\
 0&0&1&0
 \end{pmatrix},\qquad u_{f_0}=(2i+j,0,0)^t.
 \tag{CFL.6}
\]

Here \(p\) is a divisor coordinate; its powers remain in the
matrix. The full four-character inverse and reciprocal
\(f_0\)-column are recorded by I2ZC. Every owner retains the
actual lower ray of the source-band subdivision, including
all positive original \(P\) heights.

For a one owner,

\[
 F_1=
 \begin{pmatrix}
 0&1&23i+11j&12(2i+j)\\
 -1&0&3i+3j&0\\
 0&0&1&0
 \end{pmatrix}.
 \tag{CFL.7}
\]

Its unit vectors are \((-i-j,3i+2j,0)\) for \(p\) and
\((2i+j,0,0)\) for \(c_{12}\). Its full inverse character
matrix is

\[
 \begin{pmatrix}
 0&-1&3i+3j&0\\
 1&0&-23i-11j&-12(2i+j)\\
 0&0&1&0\\0&0&0&1
 \end{pmatrix},
 \tag{CFL.8}
\]

with inverse coefficient columns
\((3i+2j,i+j,0,0)\) and \((0,-2i-j,0,0)\).
Multiplication gives the identity and cancels both complete
unit columns. At infinity,

\[
 (X_c,Z_c,P)=(X_\infty,Z_\infty,Ru),\qquad
 F_\infty=
 \begin{pmatrix}
 1&0&0&0\\0&1&0&0\\0&0&1&1
 \end{pmatrix}.
 \tag{CFL.9}
\]

The inverse is \(R=pP,u=p^{-1}\). The \(u\) factor in the
original \(P\) is retained in the normal and determinant maps.

For a period facet with source and target matrices \(B_s,B_t\),
original lattice map \(H\), and full coefficient column \(u\),
the affine overlap is

\[
 z'_\alpha=
 \varepsilon^{(B_t^{-1}u)_\alpha}
 \prod_\beta z_\beta^{(B_t^{-1}HB_s)_{\alpha\beta}}.
 \tag{CFL.10}
\]

Its inverse uses \(B_s^{-1}H^{-1}B_t\) and
\(-B_s^{-1}H^{-1}u\). The recorded localizations make every
negative-power coordinate invertible. Composing both maps
gives the identity as the full analytic function. Thus CFL.5
is the normal derivative of the actual source/line chart map.

All coefficient functions in CFL.10 are invariant under the
physical circle. Along a fixed component their derivative on
a nonzero-weight normal direction vanishes by equivariance.
Differentiating the complete coordinate arrow therefore gives
its written coefficient restricted to that component, with
the recorded monomial normal factors. The zero-weight tangent
block and the signed normal eigenlines retain their ordering.

\par

\medskip\noindent\textbf{CFL.3. Original point modules and determinant}

\par\nobreak

The source band \(b\) has

\[
 x_b=Q^bZ,\qquad y_b=Q^{1-b}/Z,\qquad x_by_b=P^2R.
 \tag{CFL.11}
\]

The proper point graph is the blowup of the literal ideal
\((x_b,P^2)\), with charts

\[
 x_b=P^2a,\quad ay_b=R;\qquad
 P^2=x_bc,\quad y_b=Rc.
 \tag{CFL.12}
\]

Their overlap is \(c=a^{-1}\), and the point section is
\(a=1\) or \(c=1\). The second chart has the explicit
two-chart resolution of \(x_bc=P^2\) in I2P.2. Pulling
these equations back by the endpoint monomials gives the
finite marked source graphs and their original source maps.

Let \(E,E_O\) be the entire convergent theta-product units of
I2P.3 and I2P.7. With the point-divisor generators
\(e_a=(a-1)^{-1},e_c=(c-1)^{-1}\) and zero-divisor generator
\(e_O=x_b-1\), the full coefficient arrows are

\[
 V_{\rm reg}=-\eta^{-1}P^{-2b}EA
       =\eta^{-1}P^{-2b}c^{-1}EC,\qquad C=-cA,
\]
\[
 V_{\rm reg}=
       \eta^{-1}P^{-2b-2}(x_b-P^2)E_OC_O.
 \tag{CFL.13}
\]

Each inverse uses the reciprocal at the actual inverse source
point. The \(c^{-1}\) has its specified Cartier type, and
the extra \(P^{-2}\) at \(O\) is retained. The source period
has its simultaneous band change \(b\mapsto b-1\); substitution
leaves \(A,C,C_O\) unchanged. The full suspension is

\[
 (Z,V_{\rm reg})\longmapsto(P^{-3}Z,d_BV_{\rm reg}),
 \quad d_B=\lambda_0s^{12}q\eta^6t^{-3},
\]
\[
 v_2'=d_BP^{2(b'-b)}
       \frac{f_{P^2}(x_b)}{f_{P^2}(x'_{b'})}v_2,
 \quad x'_{b'}=Q^{b'-b}P^{-3}x_b.
 \tag{CFL.14}
\]

These are the literal SI2 Miller Hom and I2P source/Cartier
squares. For the endpoint-to-common word \((i,j)\), the
genuine point coefficient relative to common band one is

\[
 P^{26i+14j+2-2b}p^{2i+j}f_0^{2i+j}
 \frac{f_{P^2}(Q^bZ_{\rm native})}
      {f_{P^2}(QZ_{\rm common})}.
 \tag{CFL.15}
\]

Equations CFL.12--14 cancel its point and origin zeros and
poles in their specified modules. They give the actual
boundary line morphism with exactly the \(p,f_0,\eta\) and
selected-height functions of CFL.5--10. The normal calculation
here concerns the stated completed periodic atlas. Its
original point-line boundary identification is this finite
graph pullback, with its reciprocal.

The original volume forms are

\[
 \Omega_0=+\frac{2J_{ts}}{a_T}
       P\,dP\wedge dp\wedge d\log X_0\wedge d\log Y_0,
\]
\[
 \Omega_1=-\frac{2J_{ts}}{a_T}
       P\,dP\wedge dv\wedge d\log X_1\wedge d\log Y_1,
\]
\[
 \Omega_\infty=+\frac{2J_{ts}}{a_T}
       R\,dR\wedge du\wedge d\log X_\infty\wedge d\log Z_\infty.
 \tag{CFL.16}
\]

Here \(J_{ts}=\det\partial(t,s)/\partial(P,R)\); the signs
are the exact base and character Jacobians. A native ray
\((a,b,r,z)\) has original canonical order \(2r+z-1\);
a generic ray \((a,b,h)\) has order \(2h-1\).
The base-unit derivatives wedge to zero against the two
base differentials in CFL.16, while their nonzero
multiplicative units remain in every coordinate and normal map.

For every common coordinate line, the completed degree is

\[
                         d_i=-\sum_e c_{e,i}.
 \tag{CFL.17}
\]

This follows from the actual direction of CFL.5: the frame
transition is the reciprocal of the coefficient transition.
Each remaining unit has its actual logarithm on the endpoint
disc after fixing its nonzero constant. Its degree is zero,
and its full boundary value stays in the bundle arrow.
The independently checked original adjunction identity is

\[
 \sum_i(2h_i-1)d_i+2\sum_e r_e=-2-\sum_i d_i,
 \tag{CFL.18}
\]

where \(r_e\) is the original native height of the actual
lower ray. This includes every positive lower height and
the full infinity factor \(P=Ru\).

\par

\medskip\noindent\textbf{CFL.4. Generic fixed surfaces and saturated quotient fans}

\par\nobreak

Each of the 15,138 generic surfaces has precisely two common
chart occurrences, two ordered normal rays \(n_-,n_+\), and
weights \(w_-<0<w_+\). Every period cycle is trivial.
Align each occurrence by the inverse of its full recorded
period word. The remaining rays of the two aligned charts
are \(v_0,v_\infty\). Their quotient by the normal span is
the saturated rank-one lattice, because either complete
three-ray matrix is unimodular. The exact wall equation is

\[
 v_0+v_\infty+d_-^Fn_-+d_+^Fn_+=0.
 \tag{CFL.19}
\]

The two projected rays are opposite primitive generators.
The two actual affine fibre coordinates are reciprocal,
with the full base unit of CFL.10, and their union is the
\(\mathbb P^1\) fibre. The normal fibre degrees are
\(d_\pm^F\). Its two coordinate sections \(S_0,S_\infty\)
are the tangent-zero curves in these charts. Their
self-intersections in the surface and their two normal
degrees are the three entries of CFL.17, with the actual
recorded tangent and normal indices.

At an endpoint, choose a complete aligned chart
\(B=(n_-,n_+,t_0,t_1)\). The last two rows of \(B^{-1}\)
give the saturated quotient lattice \(\mathbb Z^2\).
Every projected tangent ray is integral and primitive,
and consecutive rays have determinant one. The whole
projected fan is a chain between two opposite boundary
section rays of base height zero. Every other ray has
positive base height. The interior valence is two and
the boundary valence one; the exact traversal visits
all charts and all rays.

For an internal tangent ray \(r\) with neighbours \(r_l,r_r\),
the complete four-dimensional relation is

\[
 r_l+r_r+d_-^r n_-+d_+^r n_+ + s_r r=0.
 \tag{CFL.20}
\]

The inverse of \((n_-,n_+,r,r_l)\) computes it; its fourth
independent coefficient is zero. In the two adjacent
charts it gives \(z'=z^{-1}\) and the normal coefficient
changes \(a_\pm'=z^{-d_\pm^r}a_\pm\), including the entire
units of CFL.10. Reciprocal frames prove the degrees
\(d_\pm^r\); the tangent transverse coefficient proves
the self-intersection \(s_r\). The full fibre is
\(\sum_r m_r C_r\), with \(m_r\) the actual base height.
Every relation \(m_rs_r+m_{r_l}+m_{r_r}=0\) is checked,
with boundary multiplicity zero. Thus the original
base-function divisor has zero intersection with each
vertical curve.

\par

\medskip\noindent\textbf{CFL.5. Proper blowdowns and integral intersection basis}

\par\nobreak

Each executed chain contains a multiplicity-one component.
The saved anchor is one of these. Each other ray is
contracted only after its saturated quotient vector is
the sum of its two neighbours. The actual proper
ordinary blowup has affine maps

\[
       (u,v)\mapsto(u,uv),\qquad
       (u',v')\mapsto(u'v',v'),
 \tag{CFL.21}
\]

with overlap \(u'=v^{-1},v'=uv\). The full parent
monomials and normal units compose with these maps.
The anchor remains, and the final three-ray half-plane
is the ordinary \(\mathbb P^1\)-fibre chart over the disc.
There are exactly 52 recorded contractions: 40 at zero,
none at one and 12 at infinity. Every removed ray and
its two then-current neighbours are saved.

After contraction, the fibre charts are reciprocal and
their base transition is \(z_\beta=g_{\alpha\beta}(p)z_\alpha\),
with the full nonzero function from CFL.10. Thus the
contracted surface is the actual
\(\mathbb P(\mathcal O\oplus L)\) over the parameter sphere,
with \(L\) defined by that cocycle and with the two
specified sections. CFL.21 reconstructs the executed
surface by the recorded blowups.

Use the integral divisor basis

\[
 \mathcal B=\bigl(S_0,F,\{C_{e,r}:r\ne r_{e,\rm anchor}\}\bigr).
 \tag{CFL.22}
\]

The multiplicity-one anchors express the omitted components
integrally by the full fibre relation. Conversely the
chronological maps CFL.21 express this basis in the section,
fibre and integral exceptional classes. Hence it is an
integral basis. Its matrix \(I\) is given by
\(S_0^2=s_0,\ S_0F=1,\ F^2=0\), the full self-intersections
CFL.20, the chain adjacency and the actual section incidences.
Different endpoints are disjoint. Its inverse is integral,
which is also checked directly.

With normal degrees \(d_\pm^0\) on \(S_0\), the actual
normal Chern vectors are

\[
 \mathbf n_\pm=I^{-1}
       (d_\pm^0,d_\pm^F,(d_\pm^{e,r})_{e,r})^t.
 \tag{CFL.23}
\]

The second section is reconstructed from its pairings
\((0,1,\text{its component incidences})\). Its computed
self-intersection and both normal pairings agree with
the independently obtained values from CFL.17.
This verifies the full three-end normal-line gluing.

Adjunction on the listed rational curves gives

\[
 \mathbf K=I^{-1}
       (-2-s_0,-2,(-2-s_{e,r})_{e,r})^t.
 \tag{CFL.24}
\]

For \(b\) recorded blowups, the explicit maps give
\(K^2=8-b,\ c_2=4+b,\ p_1(TF)=-3b\). The matrix
values in CFL.24 satisfy these identities for every
surface. The normal intersections are the actual
products \(\mathbf n_-^tI\mathbf n_-\),
\(\mathbf n_-^tI\mathbf n_+\), and
\(\mathbf n_+^tI\mathbf n_+\), with all section and
wall data retained.

\par

\medskip\noindent\textbf{CFL.6. Compact corner points, curves and surfaces}

\par\nobreak

Each corner-interior component is aligned by the inverse of
its entire recorded period word. The full affine maps
CFL.10 identify its tangent and normal coordinates.
All these components have trivial period stabilizer.

An isolated point has one complete regular four-chart
star. Its four physical weights are the exact entries
of \(B^{-1}\delta_e\). A corner curve has two complete
charts. Its normal rays \(n_1,n_2,n_3\) and tangent rays
\(v_0,v_\infty\) obey

\[
                  v_0+v_\infty+\sum_{i=1}^3d_i n_i=0.
 \tag{CFL.25}
\]

The tangent rays project to opposite primitive generators
of the saturated quotient. Their two affine coordinates
give the actual compact \(\mathbb P^1\), and the coefficient/
frame calculation of CFL.20 gives its normal degrees
\(d_i\), with all aligned period and unit factors.

There are 17 compact corner surfaces, all at zero.
Project their full aligned stars by the saturated
rank-two normal quotient. The projected rays in cyclic
order have determinant one for every consecutive pair,
including the last-to-first pair. The list of those
pairs is exactly the list of chart cones. Thus the
projected fan covers the plane and its ordinary
affine gluing gives the complete compact toric surface.

On these corner-interior components both original endpoint
base functions vanish. The full unit coefficients restrict
to their fixed nonzero constants. Aligning a chart by its
recorded period word carries those constants as part of the
coordinate map. The trivial period-cycle check says that
their product around every alignment cycle is the identity.
Thus the aligned coordinate and normal maps give the stated
toric surface and its actual equivariant normal bundles.
This proves the torus lift used in the following calculation.

For an original integral auxiliary cocharacter \(a\),
each chart has the exact weights \(B^{-1}a\).
Its two normal and two tangent entries give

\[
 \int_F c_1(N_i)c_1(N_j)=
 \sum_\sigma
 \frac{n_{\sigma,i}(a)n_{\sigma,j}(a)}
      {t_{\sigma,1}(a)t_{\sigma,2}(a)},\qquad
 \int_Fp_1(TF)=
 \sum_\sigma
 \frac{t_{\sigma,1}(a)^2+t_{\sigma,2}(a)^2}
      {t_{\sigma,1}(a)t_{\sigma,2}(a)}.
 \tag{CFL.26}
\]

The auxiliary vectors lie in the original four-lattice.
Two distinct choices among
\((1,11,103,1009),(2,17,89,701),(5,13,29,53),(7,31,113,809)\)
have nonzero tangent entries throughout each complete
star. Both computations give identical integral values.
The actual torus action on the aligned star gives these
normal lifts, with coordinate and period maps CFL.10.
The tangent number also satisfies \(p_1=12-3r\), where
\(r\) is the number of rays of the complete smooth fan.

The independent root replay constructs the whole cyclic
divisor intersection matrix of each of the 17 surfaces.
It retains every divisor and every principal-character
relation, represents each normal line by its actual
support-function divisor \(D_i=-\sum_r a_i(r)D_r\), and
evaluates the products in that full matrix. This
independently reproduces all three normal intersections
and the tangent \(p_1\) of CFL.26.

\par

\medskip\noindent\textbf{CFL.7. Physical orientation and weighted residues}

\par\nobreak

The chart morphisms are holomorphic with the original
complex Jacobians CFL.16. The physical circle is the
displayed primitive cocharacter throughout. Use the
actual circle-first quotient and line \(M=U^{-1}\)
of SZ.57--59 and IE.16. A finite-stabilizer quotient
chart has its actual integration factor \(1/d\).
These strata lie in the quotient interior. Removing
the invariant normal tube of each fixed component gives
the fixed-boundary residue of

\[
                         \frac{4x^3-p_1(TZ)x}{24}.
 \tag{CFL.27}
\]

The formulas below expand the original equivariant
Euler denominator for every recorded nonzero integer
normal weight.

A point with four weights \(w_i\) contributes

\[
 k_M=-\frac1{\prod_iw_i},\qquad
 \ell_M=-\frac{\sum_iw_i^2}{\prod_iw_i}.
 \tag{CFL.28}
\]

For a fixed curve set \(W=\prod_iw_i\),
\(S=\sum_id_i/w_i\), \(Q_w=\sum_iw_i^2\) and
\(T=\sum_iw_id_i\). Its pair is

\[
                  k_M=\frac SW,\qquad
                  \ell_M=\frac{Q_wS-2T}{W}.
 \tag{CFL.29}
\]

For a surface use weights \(a,b\), intersections
\(x=\int c_1(N_-)^2,\ y=\int c_1(N_-)c_1(N_+),\
z=\int c_1(N_+)^2\), and \(p=\int p_1(TF)\).
Set

\[
                   B_2=\frac{x}{a^2}+\frac{y}{ab}+\frac{z}{b^2}.
 \tag{CFL.30}
\]

Then

\[
 k_M=-\frac{B_2}{ab},\qquad
 \ell_M=-\frac{p-x-z-2(a/b+b/a)y+(a^2+b^2)B_2}{ab}.
 \tag{CFL.31}
\]

Indeed, the \(U\) Euler denominator is
\(\prod_i(c_1(N_i)+w_i u)\). Its reciprocal has the
factor \(1-S/u\) in complex dimension one, and

\[
 1-\frac{x_-/a+x_+/b}{u}
 +\frac{x_-^2/a^2+x_-x_+/(ab)+x_+^2/b^2}{u^2}
\]

in complex dimension two. The \(u^{-1}\) coefficients
of \(u^3\) and
\(u[p_1(TF)+\sum_i(x_i+w_i u)^2]\) give the \(U\)
pair. The literal physical dual changes both signs,
giving CFL.28--31. This uses the actual physical line
and all finite normal weights.

\par

\medskip\noindent\textbf{CFL.8. Complete compact sums and ownership}

\par\nobreak

The 106,166 generic curves have the full normal clutches
of I2CV and CFL.4--18. Their exact pair is

\[
 k_{\rm curves}=
 -\frac{5884517677661433720832516179200147039685023720717}
 {692517323612787928519526807587126340256000000},
\]
\[
 \ell_{\rm curves}=
 -\frac{83730187907160166050180995681772221650262376569}
 {53270563354829840655348215968240487712000000}.
 \tag{CFL.32}
\]

The 15,138 completed generic surfaces have pair

\[
 k_{\rm surfaces}=
 \frac{5943727908830327088720935721248846341776911720717}
 {692517323612787928519526807587126340256000000},
\]
\[
 \ell_{\rm surfaces}=
 \frac{551192698986630725261025427857074561444990376569}
 {53270563354829840655348215968240487712000000}.
 \tag{CFL.33}
\]

Thus

\[
                  (k_M,\ell_M)_{\rm generic}
                           =(171/2,35101/4).
 \tag{CFL.34}
\]

The complete corner-interior table is

\begin{longtable}[]{@{}lrrrrr@{}}
\toprule\noalign{}
End & Points & Curves & Surfaces & \(k_M\) & \(\ell_M\) \\
\midrule\noalign{}
\endhead
\bottomrule\noalign{}
\endlastfoot
zero & 6163 & 72 & 17 & \(889/18\) & \(-212705/36\) \\
one & 4764 & 2340 & 0 & \(2736\) & \(17092\) \\
infinity & 7304 & 701 & 0 & \(-1331/9\) & \(-7513/9\) \\
\end{longtable}

Their sum is

\[
                    (k_M,\ell_M)_{\rm corners}
                            =(5275/2,41395/4).
 \tag{CFL.35}
\]

Equations CFL.34--35 prove CFL.1. Both generic types
have already been completed across the three parameter
discs. The corner-interior tags refer to their disjoint
complete fixed stars, with ownership proved by the
full period-orbit inventory. Each component in the
table is counted once.

\par

\medskip\noindent\textbf{CFL.9. Exact verification and remaining boundary}

\par\nobreak

The complete curve certificate is
\path{code/part3/e6_selected_i3_m2_fixed_curve_normals.json},
SHA256
\path{219b4314608a7d0c94f8586624a6e9643eeb6324e40a197502bae6c095f6e62b},
with 7,294,790 checks. It records every full endpoint
and inverse map, actual lower ray, point-period
coefficient and original determinant identity.

The complete surface certificate is
\path{code/part3_root/e6_selected_i3_m2_surface_stars.json},
SHA256
\path{000e22cf37479c760c6d9a5708cfe0f36e6d8887d5ad37732649825e76f177c4},
with 2,378,991 checks. It retains every common section
clutch, aligned surface star, saturated quotient ray,
chronological blowdown, integral intersection matrix,
normal Chern vector and actual tangent number.

The complete corner certificate is
\path{code/part3_root/e6_selected_i3_m2_corner_normals.json},
SHA256
\path{c8de0d7e5ef18f5a65bf648629cc0c03260449fb19a597aa1fa61dd656437cbc},
with 50,430 checks. It retains every aligned star,
exact wall relation, complete quotient fan and both
auxiliary-cocharacter computations. All three
executions have empty error streams.

The independent root arithmetic replay is
\path{code/part3_root/e6_selected_i3_m2_compact_arithmetic_peer.json},
SHA256
\path{a3804a52bad8b979ec072b3e2722a1b1fee6ae22a398543fb58764cb36eb7cf6}.
Its 856,099 checks multiply the original equivariant
Euler factors independently for all components and
reconstruct all 17 corner surface divisor matrices.
Every component value and the exact totals agree.

Surface and corner loading uses the new read-only
filtered JSON reader. Every mathematical function AST
and the complete geometric-computation suffix of both
programs remain unchanged. Its independent certificate
passes 151 checks, including standard-JSON decoding
of complete actual records and exact retained global
identifiers. The predecessor files, large-load failures
and equivalence records remain preserved under
\path{records/attempts/e6_selected_i3_m2_normal_filtered_loaders/.}

The compact calculation retains the source and point-line
arrows CFL.11--15, original physical action and complete
determinant CFL.16--18. The full original return also
requires the twelve lower doubles, selected moving
component and same-fan reference, and the outer
source/root/annular/parameter boundary. The whole
scalar-reference correction RFB acts on that entire
reference word. CFL.1 supplies the compact fixed term
for the stated final attachment.

The source-locked companion
\path{code/part3/e6_selected_i3_m2_compact_ledger_checks.py}
independently verifies the full character/unit inverses,
original physical vectors, both point and origin module
directions, formal Euler-denominator expansions, exact
scoped sums and all cited large-result hashes. Its fresh
JSON passes 109 checks. The complete large-row replay
remains the separately identified root certificate above.


% End expanded verified cp36 input


% Expanded verified cp36 input: notes/part3_e6_selected_i3_m2_return_support_12.tex
\par\bigskip\noindent\textbf{P3K.111 supporting proof 12: Actual whole scalar reference body}\par\nobreak
\par

\medskip\noindent\textbf{Selected-I3 \(M_{2}\): the complete fixed-source reference body and its boundary line}

\par\nobreak

Working proof, 18 November 2026. This construction retains the original
source and point line of SI2.1--4, its actual selected height
\(q=t-t_c(s)\), the cubic good cap, and the four-root-turn square
SI2.31--37. The source parameter is \(s\). The independent scalar-disc
coordinate below is \(\sigma\); its root is \(r_0\sigma^4\), so its
suspension scalar is \(\lambda_0\sigma^{12}\). We construct the
reference over one fixed nonzero source value \(s_*\), with the
original small \(t\)-disc containing both the \(I_2\) fibre at zero
and the selected \(I_1\) fibre at \(t_c(s_*)\).

The calculation identifies the entire source pieces on the scalar
central face, their original line, all three retained vanishing
discs, and their prescribed outer boundary. It supplies the relative
reference normal square used when the twelve one-end pieces are
doubled. The collision normal data, the moving selected fixed surface,
its Hopf cap, and the total outer characteristic transfer are retained
as their own geometric pieces. No return is assigned here.

\par

\medskip\noindent\textbf{RB.1. The original source cover and an exact frame at the unselected I2 fibre}

\par\nobreak

Fix the actual functions of SI2 and E81:
\[
 Q=P^2R,\qquad R-P=s u_s,\qquad t=P u_t,
 \qquad q=R c_q,\qquad Q=\kappa q,\quad\kappa=P^2/c_q.
 \tag{RB.1}
\]
The Miller coefficient is the full
\(\eta=a_T\Theta'_Q(1)\Theta_Q(P^2)/\Theta_Q(P)^2\), with
\(\eta(0)=-1\). Put
\[
 V_{\rm pt}=f_{P^2}(Z)v_2,\quad
 \alpha=P^2,\quad g=P^{-3},\quad
 \Gamma=\lambda_0\eta^6t^{-3}.
\]
On the selected neighbourhood of the fixed source, \(P,t,\kappa,
\alpha,g,\Gamma\) are nonzero analytic functions. Its complete
periods are
\[
 h:(Z,V_{\rm pt})\longmapsto(QZ,\alpha V_{\rm pt}),\qquad
 B_\sigma:(Z,V_{\rm pt})\longmapsto
                 (gZ,\sigma^{12}q\Gamma V_{\rm pt}).
 \tag{RB.2}
\]
Every coefficient is evaluated at the same fixed-source base point;
the source translation is the actual chronological \(-3G\).

Near \(t=0\), \(R\) is a unit because \(s_*\ne0\). The full
theta identity \(\Theta_Q(Q^{-1}z)=-Q^{-1}z\Theta_Q(z)\) gives
\[
 f_{P^2}(Z)=-\frac{Z}{P^2}f_{R^{-1}}(Z),\qquad
 V_I=-P^2Z^{-1}V_{\rm pt}=f_{R^{-1}}(Z)v_2.
 \tag{RB.3}
\]
This is an equality of the original Cartier modules, with inverse
\(V_{\rm pt}=-ZP^{-2}V_I\). On their common smooth source annulus
the factors are nonzero. At the marked sections the two sides use
their actual point or zero generators; hence the identity extends
as a module morphism there as well. Its two period squares are
\[
 h:(Z,V_I)\longmapsto(QZ,R^{-1}V_I),\qquad
 B_\sigma:(Z,V_I)\longmapsto
       (P^{-3}Z,\sigma^{12}\Gamma_I V_I),
 \quad \Gamma_I=\lambda_0q\eta^6/u_t^3.
 \tag{RB.4}
\]
In particular \(\Gamma_I\) is an actual unit on the fixed-source
\(I_2\) neighbourhood. The cancellation of the apparent \(P^{-3}\)
in the line coefficient follows from the displayed \(Z^{-1}\)
factor in RB.3; it is not a replacement of the Hom by its leading term.

The two-node source is covered by
\[
 x_i=Z/P^i,\qquad y_i=P^{i+1}/Z,\qquad x_i y_i=P.
 \tag{RB.5}
\]
On consecutive charts \(x_{i+1}=y_i^{-1}\),
\(y_{i+1}=x_i y_i^2\). The full periods are
\[
 h:(i,x_i,y_i)\longmapsto(i+2,Rx_i,R^{-1}y_i),\qquad
 B_\sigma:(i,x_i,y_i)\longmapsto(i-3,x_i,y_i).
 \tag{RB.6}
\]
These maps retain the unit \(R\), and RB.4 gives their line lifts.
The original base coordinate is still \(t=Pu_t\); all maps keep
the original source parameter. The source graph I2P.1--2 restricts
here to these two ordered \(P\)-factor charts because its third
factor \(R\) is a unit. The strict point and zero sections are those
of I2P.5--7. Thus RB.3--6 identify the two-node source and its point
line with the original resolved family, including its marked sections.

\par

\medskip\noindent\textbf{RB.2. Fourteen ordered factors at the selected node}

\par\nobreak

On a selected Tate source band put
\[
 x_i=Z/Q^i,\qquad \bar y_i=Q^{i+1}/(\kappa Z),
                  \qquad x_i\bar y_i=q.
 \tag{RB.7}
\]
This retains the original height rather than substituting \(Q\)
for it. Its adjacent-band arrow is
\[
 x_{i+1}=(\kappa\bar y_i)^{-1},\qquad
 \bar y_{i+1}=Q\bar y_i.
 \tag{RB.8}
\]
The source period \(h\), accompanied by \(i\mapsto i+1\), fixes
\(x_i,\bar y_i\); its line coefficient is \(\alpha\). The source
part of \(B_\sigma\) sends \((x_i,\bar y_i)\) to
\((gx_i,g^{-1}\bar y_i)\). This is the actual translation on the
fixed selected node chart, including the full base-dependent units.

Use the associated line-node charts
\[
 U_j=V_{\rm pt}/(\sigma^{12}q)^j,\qquad
 W_j=(\sigma^{12}q)^{j+1}/V_{\rm pt},\qquad
               U_jW_j=\sigma^{12}x_i\bar y_i.
 \tag{RB.9}
\]
The chronological factor order is
\[
 \underbrace{\sigma,\ldots,\sigma}_{12\ {
m factors}},
                       x_i,\bar y_i.
 \tag{RB.10}
\]
The fourteen charts of its ordered graph have the following literal
maps. For \(1\le r\le12\), use coordinates
\((u_r,v_r,x_i,\bar y_i)\) and
\[
 \sigma=u_rv_r,\quad
 U_j=\sigma^{r-1}u_r,\quad
 W_j=\sigma^{12-r}x_i\bar y_i v_r.
 \tag{RB.11}
\]
For factor thirteen use
\[
 x_i=u_xv_x,\quad U_j=\sigma^{12}u_x,\quad
 W_j=\bar y_i v_x,
 \qquad (u_x,v_x,\sigma,\bar y_i).
 \tag{RB.12}
\]
For factor fourteen use
\[
 \bar y_i=u_yv_y,\quad U_j=\sigma^{12}x_i u_y,\quad
 W_j=v_y,
 \qquad (u_y,v_y,\sigma,x_i).
 \tag{RB.13}
\]
These are the standard affine charts of the successive graph closures
of the written ordered factor ideals; their displayed products give
the original hypersurface exactly. Each chart is an affine smooth
fourfold. Consecutive charts overlap by the actual reciprocal arrow
\[
 u_{r+1}=v_r^{-1},\qquad v_{r+1}=\sigma v_r
                         \quad(1\le r<12),
 \tag{RB.14}
\]
and the last two transitions are
\[
 u_x=v_{12}^{-1},\quad v_x=x_i v_{12},\quad
 \sigma=u_{12}v_{12};\qquad
 u_y=v_x^{-1},\quad v_y=\bar y_i v_x,\quad x_i=u_xv_x.
 \tag{RB.15}
\]
Their inverses are obtained from these same equations, on the stated
nonzero overlap coordinate. This proves both the graph maps and their
directions, without identifying two different factor orders.

The physical action has weights \((+1,-1,0,0)\) on each displayed
chart. In RB.11 its fixed surface is exactly
\[
 u_r=v_r=0,
\]
with the entire source chart \((x_i,\bar y_i)\) left over. It contains
the selected source node \(x_i=\bar y_i=0\). In RB.12--13 the fixed
surfaces are instead \(u_x=v_x=0\) and \(u_y=v_y=0\), with
\((\sigma,\bar y_i)\) or \((\sigma,x_i)\) as coordinates.
The first RB.15 overlap requires \(v_{12}\ne0\); neither fixed
surface lies in that overlap. The same statement for every earlier
overlap follows from RB.14. Thus the scalar-source surfaces and the
moving selected fixed surface are disjoint in the resolved graph.
This conclusion follows from their fixed equations and exact overlap,
not from the names assigned to the pieces.

\par

\medskip\noindent\textbf{RB.3. All source, line and period squares on the reference pieces}

\par\nobreak

The full \(B_\sigma\) arrow on RB.9 has target label \(j+1\):
\[
 (x_i,\bar y_i,U_j,W_j)\longmapsto
        (gx_i,g^{-1}\bar y_i,\Gamma U_j,\Gamma^{-1}W_j).
 \tag{RB.16}
\]
On the first twelve factor charts it is
\((u_r,v_r)\mapsto(\Gamma u_r,\Gamma^{-1}v_r)\).
On the last two charts it is respectively
\[
 (u_x,v_x)\longmapsto(\Gamma u_x,g\Gamma^{-1}v_x),
 \qquad
 (u_y,v_y)\longmapsto(\Gamma g^{-1}u_y,\Gamma^{-1}v_y).
 \tag{RB.17}
\]
The inverse uses the reciprocal coefficients at the inverse source
point. Substitution in RB.11--13 verifies the source, line, height and
every product. The \(h\) arrow has \(i\mapsto i+1\), fixes the
source coordinates, and sends each normal pair to
\((\alpha u,\alpha^{-1}v)\). These maps carry RB.14--15 to themselves.
They do not change the factor index \(r\).

The moving fixed charts also have an explicit full source-band
overlap. On \(\bar y_i\ne0\), the old x-factor chart and the
next band\textquotesingle s y-factor chart are related by
\[
 x_{i+1}=(\kappa\bar y_i)^{-1},\qquad
 u_y'=\kappa\bar y_i u_x,\qquad
 v_y'=\bar y_i v_x,\qquad \sigma'=\sigma.
 \tag{RB.17a}
\]
Indeed \(u_y'v_y'=\kappa x_i\bar y_i^2=Q\bar y_i\),
and the two expressions for \(U_j,W_j\) in RB.12--13 agree
literally. The inverse has
\(\bar y_i=(\kappa x_{i+1})^{-1}\),
\(u_x=x_{i+1}u_y'\), and
\(v_x=\kappa x_{i+1}v_y'\), with the same original base \(q\).
The scalar-factor normal coordinates remain unchanged under
this source-band transition. Substitution of RB.16--17 shows
that RB.17a commutes with the full B and h arrows, including
\(g,\Gamma,\alpha,\kappa\).

On the overlap \(q\ne0\) with an ordinary scalar chain
\(V_{\rm pt}=\sigma^k A_k\), \(A_kB_k=\sigma\), the exact map is
\[
 k=12j+r-1,\qquad A_k=q^j u_r,\qquad B_k=q^{-j}v_r.
 \tag{RB.18}
\]
For example RB.16 becomes \(k\mapsto k+12\) with multiplier
\(q\Gamma\) on \(A_k\), exactly RB.2. Both integer powers of
\(q\) remain in the arrow. At \(q=0\) the full factor charts
RB.11--13 supply its continuation; RB.18 is used only on its
specified nonzero overlap.

On the unselected neighbourhood use the ordinary chain in \(V_I\)
from RB.3. The period \(h\) has source label \(i+2\), normal
factor \(R^{-1}\); \(B_\sigma\) has source label \(i-3\), scalar
label \(k+12\), and normal factor \(\Gamma_I\). On the common
smooth source region the normal arrow is precisely
\[
 (A_k^I,B_k^I)=(-P^2Z^{-1}A_k,-ZP^{-2}B_k).
 \tag{RB.19}
\]
Its product is unchanged and the source coordinate is the identity.
Equations RB.3--4 show its two full period squares. Thus both
unselected source nodes and the selected source node belong to the
same retained reference source pieces.

The genuine normal maps are obtained before differentiating by
the original point/O modules. On a source band \(\ell\), I2P gives
\[
 V_{\rm pt}=-P^{-2\ell}E_{P^2,R}A
   =P^{-2\ell}c^{-1}E_{P^2,R}C,\quad C=-cA,
 \tag{RB.20}
\]
and at the zero section
\[
 V_{\rm pt}=P^{-2\ell-2}(x-P^2)E_OC_O.
 \tag{RB.21}
\]
The full normally convergent units and their inverses occur here.
The two module generators are \((a-1)^{-1},(c-1)^{-1}\), and the
zero generator is \(x-1\). These equations carry RB.18--19 into
the actual coefficient of \(M_2\), not merely into a line with
matching degree. In a genuine local vector frame \(e_i\) write
the coefficient as \(v_i\). For \(v_b=g_{ba}v_a\), the scalar
normal coordinates satisfy
\[
 A_b=g_{ba}A_a,\qquad B_b=g_{ba}^{-1}B_a.
 \tag{RB.22}
\]
The source graph and all negative normal maps are therefore
reciprocal to the positive ones. They identify the two eigenlines,
in physical order \((-1,+1)\), with \((M_2^{-1},M_2)\).

For each \(r=1,\ldots,12\), the \(h\)-glued scalar-source piece
is the entire compact resolved fixed-source body \(S_{s_*}\).
The \(B_\sigma\) action changes \(j\) by one, so a representative
\(j=0\) has no further source identification by a nonzero B word.
Equivalently, at the unselected fibre the joint labels change by
\((i,k)\mapsto(i+2,k)\) and \((i-3,k+12)\), leaving one full
two-node source for each \(k\bmod12\). These statements concern
the pieces inside the retained \(t\)-disc. The cubic finite cap
has its own label shift four, written next; we do not identify these
twelve pieces with twelve closed components of the fully capped space.

\par

\medskip\noindent\textbf{RB.4. The actual finite cap, parameter closing and selected-tube boundary}

\par\nobreak

On the original infinity annulus retain \(q_\infty=t^{-1}\), the
root \(q_\infty=z^3\), and
\[
 h_\infty=1-t_c(s)z^3=q_\infty q.
\]
The root source, the correction \(E_P-E_O\), the normalized chord,
finite character, and sector maps are exactly SI2.22--28. In particular
\(A_0^3=B_0\). Write its full genuine line map as
\(A_0:(p,v)\mapsto(a_0p,a_{\rm line}(p)v)\). On the scalar chain
\(v=\sigma^kA_k\), \(A_kB_k=\sigma\), the actual primitive cap map is
\[
 (p,k,A_k,B_k)\longmapsto
 (a_0p,k+4,r_0a_{\rm line}(p)A_k,
                    r_0^{-1}a_{\rm line}(p)^{-1}B_k).
 \tag{RB.23}
\]
Its chronological cube is the full twelve-step B map; the coefficient
is the product of \(a_{\rm line}\) at its three successive source
points. The inverse uses \(a_0^{-1}p\). Every sector overlap keeps
the original word \(A^{j_0}B^{j_{\alpha\beta}}\) of SI2.28.

On the outer scalar circle the prescribed forward square is
\[
 \Phi_{\rho+4}q^4=h_\infty^4\Phi_\rho.
 \tag{RB.24}
\]
It has the reciprocal endpoint-to-initial square. At the selected
collar \(D_{\rm ord}^4=B^4R_4\), with
\(R_4=(g^{-4},\eta_*^{-4})\),
\(\eta_*=\lambda_{\rm base}\eta^6t^{-3}\). The exact native
comparison is
\[
 C_\rho=(e^{\rho G},e^{\rho J+3\pi i\rho(\rho-1)}),
 \quad G=\operatorname{Log}g,\quad J=\operatorname{Log}\eta_*,
 \quad C_4=R_4^{-1}.
 \tag{RB.25}
\]
The full source-cut increments of \(G,J\) are \(-6\pi i\).
Their Cox map, including both odd seam signs, is the actual
\(q_H\mapsto e^{3\pi i\rho}q_He^{9\pi i\rho}\) of PPH.2.
Its left adjoint extends over the specified three-balls and its
written quaternion map carries the principal Hopf line. The complete
DZ factor is the based exponential \(e^{48\pi it(t-1)}\) followed
by the retained physical loop \(e^{60\pi it}\). Thus the boundary
map contains the winding thirty, its two cap lifts, and the original
sector integers. The q-dependent maps and their derivatives use the
protected-tube velocity and inverse flow of PPH.2.

These are boundary and collar maps at nonzero scalar radius. We use
RB.11--13 at the scalar-zero selected corner and RB.23 at the actual
scalar-zero finite cap. In particular we do not apply \(q^4\) as
a unit at the selected zero or claim that it fixes every central
factor chart. Its entire height shift is retained by RB.24--25 and
the source-period factorization. The full module map, not an omitted
height factor, joins these boundary products.

The corresponding actual coupled-source endpoint has the complete
coordinate form I2C.1--3. With
\(f_0=(p-1)^{12}\lambda_0c_q\eta^6/(u_s^{12}u_t^3)\) and
\(A=(c_q/f_0)^4\), it is
\[
 B^{-4}q^4:(Z,V_{\rm pt})\longmapsto
                   (P^{12}Z,P^{-36}A V_{\rm pt}).
 \tag{RB.25a}
\]
Indeed \(\lambda^{-4}(\eta^6t^{-3})^{-4}
=P^{-36}u_s^{48}u_t^{12}/
[\lambda_0^4\eta^{24}(p-1)^{48}]\). In the actual selected
native chart \(q=Za_Nc_N\), \(c_N=V_{\rm pt}\), this sends
\((Z,a_N,c_N)\) to
\((P^{12}Z,P^{24}A^{-1}a_N,P^{-36}Ac_N)\).
The reciprocal chart carries the same map with the two normal
powers \(36,-24\). The positive spinor \((\bar a_N,c_N)\)
therefore has determinant phase \((P/|P|)^{-60}(A/|A|)^2\).
These are exact boundary arrows; they do not change RB.33\textquotesingle s
intrinsic source cube. The complete cut field, its inverse, the
Hermitian-ball map and the held-Cox normal-sphere extension are
I2C.5--11 in
\path{notes/part3_e6_selected_i3_m2_moving_cap_working.md}.
Their cap square is the very same RB.24, with the original
reference integer on each sector. The resolved moving normal
matrix and its Chern classes remain to be obtained from the
new atlas.

The moving fixed surface meets the outer scalar circle. The scalar
source pieces lie on \(\sigma=0\). Their only possible common
source-node overlap was RB.15, where the actual nonzero coordinate
excludes both fixed sets. On the rest of the reference source cover
the fixed equations and reciprocal chain overlaps give the same
separation. In the compact retained source body these are disjoint
closed fixed subsets. Choose an invariant metric in the written
charts and an invariant small tubular neighbourhood of the moving
fixed subset disjoint from all twelve source pieces; a finite cover
and the positive separation of disjoint compact subsets give a
positive admissible radius. Its outer-radius part is transported
through RB.25 and the actual Hopf cap. Hence excising this moving
selected tube leaves the whole source body, including its selected
vanishing disc, in each scalar-source piece.

The finite cap lies outside the retained small \(t\)-disc. Its
common ordinary-annulus product is the prescribed SI2.28 map;
RB.23--24 give its central and boundary arrows. Connections and
products on this cap are transported from the same ordinary annulus.
The source-period and scalar-label faces are paired by RB.2, RB.4,
RB.16 and RB.23 with their inverses. Thus no artificial boundary
around the selected source node is introduced into \(S_{s_*}\).
The moving surface and its cap remain separate terms in the later
fixed-set and outer-boundary ledger.

\par

\medskip\noindent\textbf{RB.5. Relative source and original-line transport, with all three discs}

\par\nobreak

Choose the original source boundary \(|t|=\epsilon_t\), with
\(|t_c(s)|<\epsilon_t/3\) on the source path from zero to \(s_*\).
Retain the proper I2P graph over the resulting compact source domain.
Its central source is the minimal \(I_3\) source. Its three local
charts have equations
\[
 R=ay,\quad s=(ay-P)/u_s(P,ay);
 \qquad P=uv,\quad s=(R-uv)/u_s(uv,R)
 \tag{RB.26}
\]
in the first graph chart and the two \(A_{1}\) charts, respectively.
The partial derivatives \(\partial_Ps\), respectively
\(\partial_Rs\), are nonzero units with central values
\(-1,+1,+1\). The local complex source lifts are explicitly
\[
 X_0=(\partial_Ps)^{-1}\partial_P,
 \qquad X_1=(\partial_Rs)^{-1}\partial_R,
 \qquad X_2=(\partial_Rs)^{-1}\partial_R.
 \tag{RB.27}
\]
On the retained boundary use the actual E81 marked source charts
and the prescribed ordinary boundary product, holding \(t\)
fixed. Their maps have inverses and are smooth over the source
path, so differentiating this product gives a specified boundary
lift \(X_\partial\) with \(ds(X_\partial)=1\).

For clarity the source-period part of that boundary map is explicit.
On its marked Tate cover write \(Z=\exp(2\pi i(u+\tau b))\),
\(\tau=\operatorname{Log}Q/(2\pi i)\). The quotient
\(Q(s,t)/Q(0,t)\) has a logarithm on this boundary source path:
\[
 \frac{Q(s,t)}{Q(0,t)}=
 \frac{u_Q(s,t)}{u_Q(0,t)}\left(1-\frac{t_c(s)}t\right).
 \tag{RB.28}
\]
Use the value zero at \(s=0\). Keeping the marked real coordinates
\((u,b)\) sends \(Z_0\) to
\(\exp(b\operatorname{Log}(Q_s/Q_0))Z_0\). Under the source
deck \(b\mapsto b+1,Z_0\mapsto Q_0Z_0\), its target is multiplied
by precisely \(Q_s\). Before the genuine point/O module
compositions, the line coefficient on this same cover is
\[
 V_{{\rm pt},s}=\exp\!\left(
      b\operatorname{Log}\frac{P_s^2}{P_0^2}\right)V_{{\rm pt},0}.
 \tag{RB.28a}
\]
The logarithm is the full one with value zero at \(s=0\);
\(P_s/P_0=u_P(s,t)/u_P(0,t)\) is a unit on the entire boundary
source path. Under \(b\mapsto b+1\), the old coefficient acquires
\(P_0^2\), and RB.28a acquires exactly the target factor
\(P_s^2\). The inverse uses both reciprocal exponentials at the
recovered marked coordinates \((u,b)\). The full point/O coefficient on each local
chart is then carried by RB.20--21 and its inverse. At a prescribed
boundary-product chart this construction is composed with that
written chart map; its derivative is the boundary lift just fixed.
Thus the integral source-period lift is retained rather than reset.

Choose a finite partition \(\chi_i\) subordinate to the actual
source charts, equal to the prescribed boundary partition on a
collar. Since each local lift projects to one,
\(X=\sum\chi_iX_i\) also projects to one. Along a source path
\(s(a)\), use the real vector field
\(\dot s(a)X+\overline{\dot s(a)X}\). It equals the prescribed
boundary velocity and is tangent to \(|t|=\epsilon_t\).
Properness and compactness make its flow defined on the whole
interval. The inverse is its backward flow. This is an explicit
relative diffeomorphism from the complete central source to
\(S_{s_*}\); no node chart or zero-section disc has been removed.

Use actual vector frames with \(e_j=g_{ij}e_i\). Define connection
coefficients for the lift by
\[
 a_i=-\sum_k\chi_k\,Xg_{ik}/g_{ik},\qquad
 a_j=a_i+Xg_{ij}/g_{ij}.
 \tag{RB.29}
\]
The horizontal vector is \(X-a_iv_i\partial_{v_i}\), with the
corresponding real conjugate. All \(g_{ij}\) are the full original
Cartier transitions RB.20--21; derivatives of the units are retained.
A relative cutoff changes the coefficients to the prescribed full
boundary derivative, which obeys the same transition identity.
Its linear ODE supplies the line transport and its exact inverse.
The reciprocal ODE transports the dual normal. This construction
takes place in the genuine point/O frames, so no meromorphic
coefficient is treated as a nonsingular vector frame.

On a period translate define the transported chart by the full
RB.16 or RB.4 map and its inverse; differentiating the same square
includes the source-dependent period derivative. On the finite
cap use the SI2.28 arrow on its ordinary annulus, extend the
relative lift through its actual proper cap charts, and retain
the central map RB.23. This is the same prescribed-annulus
construction as OG.4, now with \(h_\infty=1-t_c(s)z^3\) and
the complete selected cochain. It preserves RB.24 and its inverse.
All changes used to define these transports are outside the
already separated selected fixed tube or are its specified
quaternion/Hopf transport. Consequently they do not delete any
part of the scalar-source surfaces.

The three central node charts therefore give three actual
vanishing discs in every transported reference source body.
The two-node and one-node descriptions RB.5--8 locate their
fixed-source continuations. In the central plumbing each disc
has the same primitive boundary cycle \(a\); the complementary
elliptic cycle \(b\) survives. The retained zero-section disc
kills the base meridian \(\alpha_t\). Equivalently the three
literal plumbing arrows give
\[
 [a,b]=1,\quad\alpha_t a\alpha_t^{-1}=a,\quad
 \alpha_t b\alpha_t^{-1}=a^3b.
\]
Attaching the three discs and the zero-section disc gives the
actual inclusion
\[
 H_1(\partial S)=\mathbb Z\langle b,\alpha_t\rangle
       \oplus\mathbb Z/3\langle a\rangle
 \longrightarrow H_1(S)=\mathbb Z\langle b\rangle,
 \qquad (a,\alpha_t,b)\longmapsto(0,0,b).
 \tag{RB.30}
\]
Its use here is justified by RB.7--29: the selected source disc
has been explicitly retained and the outer boundary is the
prescribed one throughout the source and line transport.

\par

\medskip\noindent\textbf{RB.6. The twelve one-end normal squares and the intrinsic cube frame}

\par\nobreak

In the actual one-end chart use \(v=p-1\),
\(X=Z^{-1}\), \(Y=V_{\rm reg}/Z\), and
\[
 \alpha_l=Yv^{-l},\qquad\beta_l=v^{l+1}/Y,
                \quad0\le l\le11.
\]
The full directed lower-to-reference map on the retained source
boundary is
\[
 A_l=H_{+,l}\alpha_l,\quad B_l=H_{-,l}\beta_l,
\quad H_{+,l}=\eta u_s^lP^{-l}X^{-1},\quad
 H_{-,l}=\eta^{-1}XP^{l+1}u_s^{-l-1}.
 \tag{RB.31}
\]
It satisfies \(H_{+,l}H_{-,l}=P/u_s\), taking
\(\alpha_l\beta_l=v\) to \(A_lB_l=s\). The inverse keeps
the reciprocal functions at the same source point. The source
and reference parameters in this boundary comparison are the
ones in \(s=Pv/u_s\); \(\sigma\) was the independent parameter
used to construct its fixed-source reference filling.

Here the final fixed-source target includes the actual transport
of RB.5. Denote its source map from \(S_0\) to \(S_s\) by
\(T_s\), and put \(K_s=T_{s_*}T_s^{-1}\). With its transported
line and reciprocal dual maps, the complete boundary morphism is
\[
 (p,\alpha_l,\beta_l)\longmapsto
 \bigl(K_s p,\ K_s^{\rm line}(H_{+,l}(p)\alpha_l),
                 K_s^{\rm dual}(H_{-,l}(p)\beta_l)\bigr).
 \tag{RB.31a}
\]
Its inverse first applies \(K_s^{-1}\) and then the reciprocal
coefficients at that recovered source point. The normal product
remains \(s\), since the two transported line maps are dual.
Thus RB.31 is the full coefficient arrow before this specified
source transport, and RB.31a names its fixed-source target.

Under \(L=(P^3pX,PpY)\), the complete positive and negative
reference factors are \(P^{-2},P^2\). Under
\(T=(p^{-1}X,P^{12}v^{12}c_{12}Y)\) use \(l\mapsto l+12\).
The positive factor is
\[
 p u_s^{12}c_{12}
   =\lambda_0p c_q\eta^6/u_t^3
   =P^2\bigl(\lambda_0q\eta^6/t^3\bigr),
 \tag{RB.32}
\]
and the negative one is its reciprocal. This is exactly the
original \(hB\) point-line coefficient after dividing by the
additional \(s^{12}\). At \(p=1\), the actual central Miller
identity gives \(B_{2,3}=\eta^6P^2=4t^2\), hence
\(\eta^6=4u_t^2\), and RB.32 is \(4\lambda_0\).
All units and derivatives remain in the equality before this
central restriction.

The actual minimal-source point section is on \(C_1\), and
the zero section is on \(C_0\), by the complete I2P/MR source
and module arrows. Thus the intrinsic line degrees are
\((-1,1,0)\). Its source-period cube coefficient is
\(X^2V_{\rm pt}^3\). With
\(\tau_L=\log|X|/\log|Q|\), the full boundary frame is
\[
 W_p=e^{-2\tau_L\operatorname{Log}p}X^2V_{\rm pt}^3.
 \tag{RB.33}
\]
Here \(\operatorname{Log}p\) is the actual small logarithm on
the retained outer source boundary. Its two exact multipliers
under \(L\) cancel. The inverse vector frame at \(p=1\) has
the full divisor \(6F-2C_1-C_2\); multiplication by \(P^{-6}\)
gives coefficient zero on \(C_0\), retaining the entire fibre
clutch. This is the original cube frame of \(M_2\). The
four-turn endpoint factor \(P^{-12}\) in \(R_4\) is a
different, explicitly retained boundary map in RB.24--25; it
does not replace this intrinsic cube coefficient.

The integral matrix on \((C_1,C_2)\) is
\(I=\left(\begin{smallmatrix}-2&1\\1&-2\end{smallmatrix}\right)\).
The vector \((1,0)^t\) has order three in its cokernel and
\(I^{-1}(1,0)^t=(-2/3,-1/3)^t\). Consequently the relative
lift is \(D_+=-(2C_1+C_2)/3\) and \(D_+^2=-2/3\).
For the actual frame RB.33 it is \(2F+D_+\), with the same
square because \(F^2=FC_i=0\).

Finally RB.30 gives
\(H^1(\partial S;\mathbb Z)/\operatorname{im}H^1(S;\mathbb Z)
=\mathbb Z\langle\alpha_t^*\rangle\). Its connecting class
is the primitive fibre, realized by the literal vector-frame
clutch \(e_R=(P/\epsilon)e_L\). The full \(u_t\) carries
the same statement for the original \(t\)-meridian. The
\(b\)-character extends; the zero-period remainder has a
real logarithm extended by a collar cutoff. Thus every actual
boundary ratio from RB.31, including the full Miller unit,
point/O coefficients and the prescribed phase maps, is retained.
It can change the relative lift by a fibre multiple only.

Taking actual third roots of RB.33 gives the order-three flat
boundary connection; its dual is reciprocal. The interpolation
to the transported boundary connection on the other half of
the double is the identical full interpolation with reversed
collar coordinate. Its transgressions cancel pointwise.
The reference normal pair on each retained source piece is
therefore, in physical order,
\[
 (N_-,N_+)=(M_2^{-1},M_2),\qquad
 (c_-^2,c_-c_+,c_+^2)=(-2/3,+2/3,-2/3).
 \tag{RB.34}
\]
This conclusion uses the actual fourteen-factor source inclusion,
the two-node frame RB.3, the full original cap/parameter square,
and the relative source and line transport above. It leaves the
moving selected fixed surface and its cap in the subsequent
ledger, and assigns no total characteristic pair.


% End expanded verified cp36 input


% Expanded verified cp36 input: notes/part3_e6_selected_i3_m2_return_support_13.tex
\par\bigskip\noindent\textbf{P3K.111 supporting proof 13: Collision selected native star and true fan}\par\nobreak
\par

\medskip\noindent\textbf{Selected-I3 \(M_{2}\) N12: the protected native star and its full phase square}

\par\nobreak

This proof uses the executed zero endpoint and its actual common lift,
the original selected height, and the complete point-line coefficient.
It identifies the native tube on the outer source annulus by explicit
chart maps, their inverses and their source/period squares. It retains
the old I2C and RB sources as records of the preceding point-coordinate
calculation. The correction to their moving-tube interpretation is
specified in NS.6 below. No moving Chern number or return is assigned.

\par

\medskip\noindent\textbf{NS.1. The two protected facets and the original base}

\par\nobreak

Keep the SI2 source and line data
\[
 X=V_{\rm reg}^{-1},\quad Y=Z_{\rm old},\quad
 V_{\rm pt}=\eta/X,\quad R=Pp,\quad Q=P^3p,
\]
\[
 q=q_{\rm old}=Pp c_q,\qquad
 s=P(p-1)/u_s,\qquad t=Pu_t,
\]
\[
 f_0=(p-1)^{12}\frac{\lambda _0c_q\eta^6}{u_s^{12}u_t^3},
 \qquad
 \eta=a_T\Theta'_Q(1)\frac{\Theta_Q(P^2)}{\Theta_Q(P)^2}.
 \tag{NS.1}
\]
All functions are evaluated at the original source point. In particular
the theta derivative, the value \(\eta(0)=-1\), and the complete sixth
power remain present. On the selected outer collar \(P\ne0\), the
functions \(c_q,u_s,u_t,\eta,f_0\) are nonzero. The actual source map
\((P,R)\mapsto(t,s)\) has determinant \(J_{ts}\), with
\(J_{ts}(0)=1/2\). Since \(q=t-t_c(s)\),
\[
 \det\frac{\partial(s,q)}{\partial(P,p)}=-P J_{ts}.
 \tag{NS.2}
\]
Thus the original inverse source map gives the holomorphic inverse
\((s,q)\mapsto(P,p)\) on this collar. The formulas below always use
that inverse, including all its source-dependent units.

The original zero lattice is ordered as
\(\operatorname{ord}(X,Y,P,p)\). Write
\[
 A=(0,0,0,1),\quad B_1=(-1,0,0,1),\quad
 B_2=(-2,-1,0,1),\quad B_3=(-1,-1,0,1).
 \tag{NS.3}
\]
The complete common lift has exactly the two lower facets
\[
 \langle A,B_1,B_2\rangle,\qquad
 \langle A,B_2,B_3\rangle.
 \tag{NS.4}
\]
They are the unchanged lower facets of the 420-chart source-compatible
parent. This follows both from the relative construction I2ZC.2 and
from the persisted 143,270-cone output. The companion extracts every
cone containing either physical face \(\langle A,B_1\rangle\) or
\(\langle B_2,B_3\rangle\); there are seven and four, respectively.
It checks every parent map and exports the original source band,
normal characters, original canonical orders and full period units.

The actual physical vector is \(\delta=(-1,0,0,0)\). On the first
face \(\delta=B_1-A\); on the second \(\delta=B_2-B_3\).
The actual zero period
\[
 T_0=hB^{-1}:(X,Y,P,p)\longmapsto
        (P^8p f_0 X,P^6pY,P,p)
 \tag{NS.5}
\]
sends \(B_2\) to \(B_1\) and \(B_3\) to \(A\), with the same
physical vector. This is the specific period identification used below.
No source or fibre character is selected from a fan label.

\par

\medskip\noindent\textbf{NS.2. Both regular native charts and their literal overlap}

\par\nobreak

On the first lower facet use its actual primitive dual characters
\[
 u_-=pX/Y,\qquad u_+=Y^2/X,\qquad z=Y^{-1}.
 \tag{NS.6}
\]
Their inverse is
\[
 X=u_+^{-1}z^{-2},\quad Y=z^{-1},\quad p=u_-u_+z.
 \tag{NS.7}
\]
These inverse torus formulas specify the affine chart by its forward
semigroup generators; negative powers in NS.7 are used only on the
corresponding torus localization. With \(P\) invertible, NS.6 are
the three regular coordinates of this lower chart. The physical
weights of \((u_-,u_+,z)\) are \((-1,+1,0)\). Define
\[
 a_\infty=\frac{Pc_q}{\eta}u_-,\qquad
 c_\infty=\eta u_+,\qquad z=z.
 \tag{NS.8}
\]
Then \(q=z a_\infty c_\infty\). Given \((s,z,a_\infty,c_\infty)\),
first compute that literal \(q\), recover \((P,p)\) by NS.2, and set
\[
 u_-=\frac{\eta}{Pc_q}a_\infty,\qquad
 u_+=\eta^{-1}c_\infty.
 \tag{NS.9}
\]
Both compositions are identities, including their dependence on
\(q\). In particular the coefficients in NS.9 are evaluated at the
recovered source point, not at a substituted new normal product.

On the second lower facet use its primitive dual coordinates
\[
 u_A=pY,\qquad u_2=Y/X,\qquad u_3=X/Y^2,
 \tag{NS.10}
\]
with inverse \(X=u_2^{-2}u_3^{-1}\),
\(Y=u_2^{-1}u_3^{-1}\), \(p=u_Au_2u_3\).
Their physical weights are \((0,+1,-1)\).
Apply the full period NS.5 before comparing with the first chart.
The resulting native coordinates are
\[
 Z=P^6u_A,\qquad
 a_0=P^{-3}c_q f_0\eta^{-1}u_3,\qquad
 c_0=\eta P^{-2}f_0^{-1}u_2.
 \tag{NS.11}
\]
Again \(q=Za_0c_0\). Its full inverse, after the same actual base
recovery from \((s,q)\), is
\[
 u_A=P^{-6}Z,\quad
 u_2=P^2 f_0\eta^{-1}c_0,\quad
 u_3=P^3\eta(c_qf_0)^{-1}a_0.
 \tag{NS.12}
\]
These are holomorphic inverse maps on the protected outer collar.

On their common localization the complete period NS.5 gives
\[
 u_-^{\infty}=P^2f_0u_Au_3,\qquad
 u_+^{\infty}=P^4f_0^{-1}u_Au_2,\qquad
 z=P^{-6}u_A^{-1}.
 \tag{NS.13}
\]
Its inverse is
\[
 u_A=P^{-6}z^{-1},\quad
 u_2=P^2f_0u_+^{\infty}z,\quad
 u_3=P^4f_0^{-1}u_-^{\infty}z.
 \tag{NS.14}
\]
Substitution into NS.8 and NS.11 yields the exact two-chart map
\[
 z=Z^{-1},\qquad a_\infty=Za_0,\qquad c_\infty=Zc_0,
 \qquad Za_0c_0=z a_\infty c_\infty=q.
 \tag{NS.15}
\]
This proves the regular native tube and its fixed sphere directly
from the two protected lower facets and their specified period.
The original selected height, the physical circle and every analytic
coefficient are carried by the displayed maps.

On the aligned dense torus the same map reads
\[
                    c_0=ZV_{\rm pt},\qquad
                    a_0=\frac{q}{Z^2V_{\rm pt}}.
 \tag{NS.16}
\]
NS.8--15 prove its extension to both native vertex charts. They also
specify its inverse there, rather than using NS.16 alone at a zero
or a pole of \(Z\).

\par

\medskip\noindent\textbf{NS.3. Full point modules, source periods and original determinant}

\par\nobreak

For source band \(j\), retain \(x=Q^jZ\),
\(y=Q^{1-j}/Z\) and the actual point graph
\(\operatorname{Bl}_{(x,P^2)}(xy=P^3p)\). In the I2P point
generators and the zero generator the coefficient is literally
\[
 V_{\rm pt}=-P^{-2j}E A_{\rm pt}
 =P^{-2j}c^{-1}E C_{\rm pt}
 =P^{-2j-2}(x-P^2)E_OC_O.
 \tag{NS.17}
\]
Here \(C_{\rm pt}=-cA_{\rm pt}\), and \(E,E_O\) are the entire
I2P theta products. Multiplication of each expression by \(Z\)
gives \(c_0\); its reciprocal module map, with \(q/Z\), gives
\(a_0\). In the other native chart the coefficients are those of
NS.15. Thus every point/O graph uses the actual I2P Cartier module
and its inverse. The factor \(P^{-2}\) at the zero section is
retained. The use of \(\eta\) in NS.8--12 comes from
\(V_{\rm pt}=\eta V_{\rm reg}\); it cancels in NS.17 only by that
exact stated module equality.

The original point periods are
\[
 h:(Z,V_{\rm pt})\mapsto(P^3pZ,P^2V_{\rm pt}),\qquad
 B:(Z,V_{\rm pt})\mapsto(P^{-3}Z,P^{10}pf_0V_{\rm pt}).
 \tag{NS.18}
\]
Under the map NS.16 their complete native forms are
\[
 \begin{aligned}
 h:&(Z,a_0,c_0)\mapsto
 (P^3pZ,P^{-8}p^{-2}a_0,P^5p c_0),\\
 B:&(Z,a_0,c_0)\mapsto
 (P^{-3}Z,P^{-4}p^{-1}f_0^{-1}a_0,P^7pf_0c_0).
 \end{aligned}\tag{NS.19}
\]
Their products are \(q\), and their inverses divide by the same
full factors at the unchanged \((P,p)\). In particular
\[
 \begin{aligned}
 B^{-1}:&(P^3Z,P^4pf_0a_0,P^{-7}p^{-1}f_0^{-1}c_0),\\
 hB^{-1}:&(P^6pZ,P^{-4}p^{-1}f_0a_0,P^{-2}f_0^{-1}c_0).
 \end{aligned}\tag{NS.20}
\]
These squares include both native charts through NS.15. The
source-band change and point-module factors are exactly I2P.8--9;
the formulas do not erase the selected-height powers at \(p=0\).
Their continuation there is the named period map between the
appropriate affine charts, including NS.13--14.

The original four-form is
\[
 \Omega_{\rm orig}=\frac{2J_{ts}}{a_T}P\,
 dP\wedge dp\wedge d\log X\wedge d\log Y.
 \tag{NS.21}
\]
Using \(X=\eta Z/c_0\), \(Y=Z\), the derivatives of \(\eta\)
wedge to zero against the two original base differentials. NS.2 and
\(q=Za_0c_0\) then give, with the displayed orders of coordinates,
\[
 \Omega_{\rm orig}=\frac2{a_T}\,ds\wedge dZ\wedge da_0\wedge dc_0
 =-\frac2{a_T}\,ds\wedge dz\wedge da_\infty\wedge dc_\infty.
 \tag{NS.22}
\]
The last sign is the actual determinant of NS.15, namely \(-1\).
The full unit \(a_T\) and these signs are retained. At the interior
\(P=0\) divisors the original orders remain the exported
\(2r+z-1\); NS.22 describes the protected outer collar and does
not discard those divisors.

For every one of the eleven resolved moving-star charts, the
companion exports the actual parent matrix and its inverse. Align
the second face by NS.5. In the aligned lattice the normal characters
for the zero native chart are
\[
 m_-=(1,-2,1,1),\qquad m_+=(-1,1,0,0),
 \tag{NS.23}
\]
with coefficient units \(c_q\eta^{-1}f_0^n\) and
\(\eta f_0^{-n}\), respectively, where \(n=0\) on the first
face and \(n=1\) after NS.5. The matrix NS.5 carries all P powers;
these are not included a second time in that unit column. For the
infinity native chart add \((0,1,0,0)\) to both rows. Evaluating
these rows on each executed cone gives the full tangent monomial
coefficient of its two normal derivatives. On the two normal columns
the rows are exactly \((1,0)\) and \((0,1)\), in physical order
\((-1,+1)\). Both P and the native fixed-curve coordinate have
zero normal order; p contains the product of the two normals.
Consequently derivatives of the full base units have zero first
normal term on the fixed surface. This proves that the exported
coefficient matrix is its actual normal derivative.

More explicitly, let \((v_-,v_+)\) be the two actual local normal
coordinates of one executed chart, in physical order \((-1,+1)\),
and let \(e_-,e_+\) be the full tangent monomials and units obtained
from NS.23. In the positive complex coefficient convention the
resolved-to-held-Cox derivative is
\[
 N_\theta^+=G_0(Z)\operatorname{diag}(\overline{e_-},e_+),
 \qquad G_0(Z)=\begin{pmatrix}1&-Z\\\bar Z&1\end{pmatrix}
 \tag{NS.23a}
\]
on the zero native chart. On the other chart replace \(G_0\) by
\(G'_0(z)=\left(\begin{smallmatrix}\bar z&-1\\1&z\end{smallmatrix}\right)\)
and use the two infinity rows. The identity
\(G'_0(1/Z)\operatorname{diag}(\bar Z,Z)=G_0(Z)\)
is the exact normal overlap. The full transported Hermitian norm is
used on its sphere. These maps define \(N_\theta\) on the actual
outer boundary; a tangent monomial with nonzero central divisor
remains that divisor in the interior atlas.

The quotient tangent rays of the eleven charts, in the aligned
\((\operatorname{ord}Z,\operatorname{ord}P)\) lattice, are exactly
\[
 (1,0),(3,1),(2,1),(1,1),(0,1),(-1,1),(-2,1),(-3,1),
                 (-4,1),(-5,1),(-6,1),(-1,0).
 \tag{NS.23b}
\]
Every adjacent pair is a primitive cone and occurs once. Remove
successively \((3,1),(2,1),(1,1)\), then
\((-6,1),(-5,1),\ldots,(-1,1)\). Each removed ray is the sum of
its two then-adjacent rays. The inverse operation is the ordinary
two-coordinate blowup, with actual maps \((u,v)\mapsto(u,uv)\)
and \((u',v')\mapsto(u'v',v')\) in that retained lattice basis.
The result is precisely the three-ray fan
\((1,0),(0,1),(-1,0)\) of the native fixed sphere times its
parameter disc. Thus these are nine explicitly recorded interior
blowups. Their full normal monomial rows and original canonical
orders travel through the same displayed chart maps. This supplies
the resolved surface and normal input without substituting a
characteristic number for a chart or a clutch.

\par

\medskip\noindent\textbf{NS.4. The coupled endpoint in the actual native tube}

\par\nobreak

Put \(\mathcal A=(c_q/f_0)^4\). The exact point-coordinate endpoint
from SI2 and I2C is
\[
 E:(Z,V_{\rm pt})\mapsto(P^{12}Z,P^{-36}\mathcal A V_{\rm pt}).
 \tag{NS.24}
\]
Conjugation by the two actual native maps proved above gives
\[
 E_0:(Z,a_0,c_0)\mapsto
 (P^{12}Z,P^{12}\mathcal A^{-1}a_0,P^{-24}\mathcal A c_0),
 \tag{NS.25}
\]
\[
 E_\infty:(z,a_\infty,c_\infty)\mapsto
 (P^{-12}z,P^{24}\mathcal A^{-1}a_\infty,
                         P^{-12}\mathcal A c_\infty).
 \tag{NS.26}
\]
Both preserve the literal height, and NS.15 conjugates one to the
other. All inverses use the reciprocal functions at the same source
point. For the positive spinor \((\overline{a_0},c_0)\) the full
normal matrix is
\(\operatorname{diag}(\overline{P}^{12}\overline{\mathcal A}^{-1},
P^{-24}\mathcal A)\). Its determinant phase is therefore
\[
             (P/|P|)^{-36}(\mathcal A/|\mathcal A|)^2.
 \tag{NS.27}
\]
The full constant and based-unit phases remain in this equality.
On the selected fixed locus \(s=-P/u_s(P,0)\), the original positive
source meridian has positive degree one in P. The entire unit
\(\mathcal A(P,0)\) has its chosen logarithm relative to its nonzero
constant. Thus NS.27 has winding \(-36\) in these actual native
normal frames. Its use in a cap requires the exact Cox product below.

\par

\medskip\noindent\textbf{NS.5. Every phase-square edge and the full DZ factor}

\par\nobreak

Retain the relative root coordinate \(\rho\), the actual
\(r_{\rm base}=r_0s_*^4\), and the full SI2 logarithms
\[
 G=\operatorname{Log}(P^{-3}),\qquad
 J=\operatorname{Log}(\lambda_{\rm base}\eta^6t^{-3}).
\]
The point-coordinate product is
\[
 a_\rho=e^{\rho G},\quad
 d_\rho=e^{\rho J+3\pi i\rho(\rho-1)}.
 \tag{NS.28}
\]
Its actual native conjugate, using NS.16, is
\[
 C^0_\rho:(Z,a_0,c_0)\mapsto
 (a_\rho Z,(a_\rho^2d_\rho)^{-1}a_0,a_\rho d_\rho c_0).
 \tag{NS.29}
\]
In the standard native notation its second factor is
\(\beta_\rho=a_\rho d_\rho\), whose full logarithm is
\(\rho(G+J)+3\pi i\rho(\rho-1)\). These formulas preserve q
and have their reciprocal inverses at the same q and source parameter.

Write \(g=e^G\), \(\eta_*=e^J\),
\(\eta_\rho=e^{6\pi i\rho}\eta_*\). Subtracting the full
logarithms proves
\[
 C^0_{\rho+1}\,S(g^{-1},\eta_\rho^{-1})=C^0_\rho,
 \qquad C^0_0=1,\qquad C^0_4=S(g^4,\eta_*^4),
 \tag{NS.30}
\]
where the exact conjugation of a point map \((a,d)\) is
\[
 S(a,d):(Z,a_0,c_0)\mapsto(aZ,(a^2d)^{-1}a_0,ad c_0).
 \tag{NS.31}
\]
The map at four is the inverse of the same endpoint NS.25 at the
identical root endpoint. The map at the initial source cut uses the
full values of G and J there. Both source-cut logarithm increments
are \(-6\pi i\); hence its native mismatch is
\[
 (a,\beta)=(e^{-6\pi i\rho},e^{-12\pi i\rho}).
 \tag{NS.32}
\]
This identifies all four edges and every intermediate integer seam.

Use the right-complex held Cox coefficient
\[
 T_Z=\begin{pmatrix}1&-\overline Z\\Z&1\end{pmatrix},\qquad
 q_H=T_Z(a_0,\overline{c_0})^t.
 \tag{NS.33}
\]
The second chart uses
\(T'_z=\left(\begin{smallmatrix}z&-1\\1&\bar z\end{smallmatrix}\right)\);
the literal identity
\(T'_{1/Z}\operatorname{diag}(Z,\bar Z)=T_Z\)
proves that this is the same held quaternion. Under NS.32 put
\(w=e^{\pi i\rho}\). Direct matrix multiplication gives
\[
 T_{w^{-6}Z}\operatorname{diag}(w^{18},w^{12})
       =\operatorname{diag}(w^{18},w^{12})T_Z.
\]
Left and right multiplication on the right-complex coefficient pair
therefore give exactly
\[
                  q_H\longmapsto w^3q_Hw^{15}.
 \tag{NS.34}
\]
At each integer seam both factors have sign \((-1)^\rho\).
Their full Spin signs are retained. The left adjoint is the actual
base and normal-orbit rotation and extends linearly over the specified
three-ball; the same quaternion arrow supplies its principal Hopf
line map. The right factor is part of that map.

For arbitrary full functions \(a_\rho,d_\rho\), the normal sphere
map is given, with the transported Hermitian norm, by the invertible
matrix
\[
 T_{a_\rho Z}
 \operatorname{diag}((a_\rho^2d_\rho)^{-1},
                    \overline{a_\rho d_\rho})T_Z^{-1}.
 \tag{NS.35}
\]
Its inverse is the inverse product at the recovered point. The full
q-dependent native velocity preserves q. Keeping the prescribed
q-preserving source lift, multiply only that vertical correction by
the protected invariant cutoff. Its flow and backward flow provide
the collar comparison on both vertex charts. This carries every
source derivative and all endpoint jets. The finite cap is outside
the cutoff support, so its unchanged full square is
\[
              \Phi_{\rho+4}q^4=h_\infty^4\Phi_\rho.
 \tag{NS.36}
\]
It retains the root chart \(k\mapsto k+4\), its three chronological
steps, every point correction, sector integer and finite character.

On the diagonal \(\rho=4t\), the native second-factor logarithm
from NS.29 is
\[
       4t(G+J)+3\pi i(16t^2-4t).
\]
The original DZ product has the same source logarithm \(4tG\)
and native second-factor logarithm
\[
                         4t(G+J)+96\pi it^2.
\]
This follows by adding the unchanged source logarithm to the actual
point-line DZ logarithm \(4tJ+96\pi it^2\). Thus the entire G and J,
including all analytic units and their constant values, cancel
pointwise in the comparison. The remaining directed map is precisely
the physical scalar
\[
 F(t)=e^{3\pi i(16t^2+4t)}
     =e^{48\pi it(t-1)}e^{60\pi it},\qquad
 (a_0,c_0)\mapsto(F(t)^{-1}a_0,F(t)c_0).
 \tag{NS.37}
\]
The first logarithm has zero values at both endpoints. The second is
the retained positive physical winding thirty. In the held Cox
quaternion it is the inverse right phase. The explicit intertwining
\(S\delta_F=\delta_F S\) also follows directly from NS.16, so this
loop is carried on the same physical line and both finite-cap
interfaces. No contraction of its nonzero winding is used.

\par

\medskip\noindent\textbf{NS.6. Exact scope bridge for I2C, PPH and RB}

\par\nobreak

Keep I2C at SHA256
\path{e9b6b65b1ef64c7c086fca00bde3c937bb34d6b07d0b2e7ce8b422b2c80cd072}
and RB at SHA256
\path{a8f69a012bc2ec528d2da2181ae6049803b41d74710e3d7e262b4974a90fdb79}.
Their point-coordinate endpoint NS.24 is exact. On its dense common
torus the old coefficients satisfy
\[
 c_N=V_{\rm pt},\quad a_N=q/(ZV_{\rm pt}),\qquad
 a_0=Z^{-1}a_N,\quad c_0=Zc_N.
 \tag{NS.38}
\]
This is the exact morphism relating the two coefficient systems on
that localization. At the regular native point \(Z=0,a_0=0,c_0=1\),
the old expression \(c_N=c_0/Z\) has a pole. Therefore that displayed
old point-coordinate map does not extend to this native vertex chart.
NS.8--15 supply the actual regular extension selected by the protected
lower atlas.

In positive normal coefficients the common-locus comparison is
\(S_+(Z)=\operatorname{diag}(\bar Z^{-1},Z)\). Its directed endpoint
square is
\[
 S_+(P^{12}Z)
 \operatorname{diag}(\bar P^{24}\bar{\mathcal A}^{-1},
                         P^{-36}\mathcal A)
 S_+(Z)^{-1}
 =\operatorname{diag}(\bar P^{12}\bar{\mathcal A}^{-1},
                         P^{-24}\mathcal A).
 \tag{NS.39}
\]
This proves the change from the old phase exponent \(-60\) to the
actual native exponent \(-36\) without declaring the intervening
meromorphic character an extending gauge. Likewise NS.29--35 are
the full native replacement for the old point-coordinate interpretation
of the SI2.36/PPH.2 Cox product; its right phase becomes fifteen.
NS.37 proves the complete DZ comparison through that exact map.

For RB.25a the point endpoint stays NS.24, while its native-tube
assertion uses NS.25--27 and its cut/Cox maps use NS.29--36. The disjoint
fourteen-factor source pieces, their source/line/cap transports, and
the intrinsic cube and relative pairings in RB.31--34 are retained.
Their comparison with the moving boundary is now made through this
actual native tube. The compact resolved normal bundle, its cap
clutch and the global outer characteristic ledger remain subsequent
calculations. In particular this proof supplies their coordinate and
orientation input and assigns none of their characteristic numbers.


% End expanded verified cp36 input


% Expanded verified cp36 input: notes/part3_e6_selected_i3_m2_return_support_14.tex
\par\bigskip\noindent\textbf{P3K.111 supporting proof 14: Directed native Hopf cap and five-dimensional bundle bordism}\par\nobreak
\par

\medskip\noindent\textbf{Selected-I3 \(M_{2}\) N12: the actual native cap and its directed bundle bordism}

\par\nobreak

This construction uses the eleven actual native-star charts NS.1-\/-6,
the original SI2 height and point/O modules, the cubic cap with norm
one, its fixed finite character, the complete annular reference, and
the coupled scalar \(\lambda=\lambda_0 s^{12}\) with its chosen root \(r_{0} s^{4}\).
The source and period maps in NS are retained literally. The following
cap attaches to that regular native tube. The older I2C source remains
preserved at its stated hash; NS.38-\/-39 give the exact localization map
and its pole at the native vertex. No Chern number is inferred from
the older coefficient system.

\par

\medskip\noindent\textbf{NC.1. The full native endpoint, its inverse, and the cut collar}

\par\nobreak

Write \(q=q_{\rm old}\) and keep the complete functions

\[
 Q=P^3p,\quad q=Pp c_q,\quad
 s=P(p-1)/u_s,\quad t=Pu_t,\quad
 f_0=(p-1)^{12}\frac{\lambda_0c_q\eta^6}{u_s^{12}u_t^3},
 \qquad {\cal A}=(c_q/f_0)^4.                       \tag{NC.1}
\]

Here \(\eta\) includes the actual theta derivative and \(\eta(0)=-1\). All
functions and their reciprocals are evaluated at the original source
point recovered by the inverse of NS.2. On the protected outer source
collar P is nonzero and each displayed unit is nonzero. The regular
native charts proved in NS satisfy

\[
 q=Za_0c_0=z a_\infty c_\infty,\qquad
 z=Z^{-1},\quad a_\infty=Za_0,\quad c_\infty=Zc_0.
                                                        \tag{NC.2}
\]

The actual endpoint and its inverse in these charts are

\[
 E_0=(P^{12}Z,P^{12}{\cal A}^{-1}a_0,
                         P^{-24}{\cal A}c_0),\qquad
 E_0^{-1}=(P^{-12}Z,P^{-12}{\cal A}a_0,
                         P^{24}{\cal A}^{-1}c_0),       \tag{NC.3}
\]
\[
 E_\infty=(P^{-12}z,P^{24}{\cal A}^{-1}a_\infty,
                            P^{-12}{\cal A}c_\infty).
                                                        \tag{NC.4}
\]

These formulas fix s and q. Multiplying their three coordinate
factors gives one; composing either direction with NC.2 proves the
whole overlap square, including the full analytic functions. Their
inverses use the reciprocal factors at the same recovered source
point. NS.8-\/-14 already prove their regular extension at both native
vertices, so the dense expressions \(c_{0}=Z V_{\rm pt}\) and \(a_{0}=q/(Z^{2} V_{\rm pt})\) are
used here with those extending chart maps.

Cut the positive original s circle. Choose its continuous logarithm
of P and the based logarithm of the unit A on this cut. For a smooth
\(\chi(\theta)\) constant zero on the initial collar and constant one on
the terminal collar, set

\[
 \alpha_\theta=\exp[-\chi(\theta)\operatorname{Log}(P^{12})],
 \qquad
 \beta_\theta=\exp[-\chi(\theta)
                          \operatorname{Log}(P^{-24}{\cal A})].
                                                        \tag{NC.5}
\]

The logarithms include their constant values and their entire source
dependence. Define

\[
 C_\theta^0(Z,a_0,c_0)=
 (\alpha_\theta Z,(\alpha_\theta\beta_\theta)^{-1}a_0,
                                             \beta_\theta c_0),
                                                        \tag{NC.6}
\]
\[
 C_\theta^\infty(z,a_\infty,c_\infty)=
 (\alpha_\theta^{-1}z,\beta_\theta^{-1}a_\infty,
                              \alpha_\theta\beta_\theta c_\infty).
                                                        \tag{NC.7}
\]

NC.2 intertwines these maps exactly. Both fix q; hence the inverse
recovers that same q, recovers P,p from s,q, and divides by precisely
the written factors. This proves both compositions as maps of the
original source and the actual line modules. The endpoint values are
\(C_{0}=1\) and \(C_{1}=E^{-1}\). On the following cut define \(C_{\theta+1} E=C_\theta\).
Because \(\chi\) is constant on collars, this is an equality of full
collar maps and all their endpoint derivatives.

The real vector field of NC.6 has holomorphic-coordinate part

\[
 (\partial_\theta\log\alpha)Z\partial_Z
 -(\partial_\theta\log(\alpha\beta))a_0\partial_{a_0}
 +(\partial_\theta\log\beta)c_0\partial_{c_0}.            \tag{NC.8}
\]

It annihilates q by direct differentiation. NC.7 gives the same field
on the other vertex chart. Add the prescribed parameter-horizontal
lift fixing q, and multiply only this vertical correction by a smooth
physical-circle-invariant cutoff equal to one on the protected tube
and zero off a slightly larger tube. Its support is disjoint from the
good cap, its annular interface and the outer source boundary. The
real field is smooth on the compact cut fibres; its ordinary flow and
backward flow exist for the full compact parameter interval. They
give the exact NC.6 map on the smaller tube and the prescribed
unchanged comparison on each protected interface. The field commutes
with the physical action because its coefficients depend on the
invariant s,q and its coordinate weights are diagonal.

Thus this is an actual interval identification of the family and its
line bundles. The finite cap retains the original full square

\[
                     \Phi_{\rho+4}q^4=h_\infty^4\Phi_\rho.
                                                        \tag{NC.9}
\]

The root chart shift by four, its three chronological comparisons,
the point correction at each step, every sector integer and the
finite character are all the original SI2/NS arrows, since the flow
is the identity on that interface. The physical scalar of winding
thirty remains the full NS.37 comparison. In particular this collar
construction preserves that scalar and does not remove its phase.

\par

\medskip\noindent\textbf{NC.2. The whole Cox map and its extension over the actual ball}

\par\nobreak

Use the right-complex coefficient and the other sphere chart

\[
 T_Z=\begin{pmatrix}1&-\overline Z\\ Z&1\end{pmatrix},\quad
 q_H=T_Z(a_0,\overline{c_0})^t,\qquad
 T'_z=\begin{pmatrix}z&-1\\1&\overline z\end{pmatrix}.
                                                        \tag{NC.10}
\]

The literal identity
\(T'_{1/Z} \operatorname{diag}(Z,\overline{Z})=T_Z\) proves that both expressions give the same
held quaternion. The physical circle acts by its inverse right
phase. In this coefficient the entire first-normal map of NC.6 is

\[
 G_\theta(Z)=T_{\alpha_\theta Z}
 \operatorname{diag}((\alpha_\theta\beta_\theta)^{-1},
                                      \overline{\beta_\theta})
 T_Z^{-1}.                                             \tag{NC.11}
\]

In NC.11 \(\alpha\) and \(\beta\) are their actual restrictions to \(q=0\). Their
normal derivatives vanish at the fixed surface: \(q=Za0c_{0}\) has zero
first normal differential there and P,p are functions of s,q.
Thus NC.11 is the full first-normal derivative, including the
original source units. Its determinant is nonzero and its inverse
is the reversed matrix product at the recovered point. On the other
sphere chart the NC.10 identity and NC.7 give exactly this same
matrix. It is smooth at infinity; the other chart gives there the
diagonal value \(\operatorname{diag}(\overline{\alpha \beta},\beta^{-1})\). Sphere normalization
uses the transported Hermitian norm. The inverse applies \(G^{-1}\) and
then normalizes, so both sphere compositions are identities.

Represent the closed three-ball as the trace-one positive
semidefinite Hermitian matrices Y. Its boundary point is

\[
 Y_Z=\frac1{1+|Z|^2}
       \begin{pmatrix}|Z|^2&Z\\\overline Z&1\end{pmatrix}.
                                                        \tag{NC.12}
\]

Put \(D_\theta=\operatorname{diag}(\alpha_\theta|_{q=0},1)\). The actual ball map is

\[
 \Psi_{D_\theta}(Y)=
 \frac{D_\theta YD_\theta^*}
      {\operatorname{tr}(D_\theta YD_\theta^*)}.           \tag{NC.13}
\]

The denominator is positive on the whole ball: D is invertible and
Y is nonzero positive semidefinite. Its inverse is \(\Psi_{D^{-1}}\), as
direct substitution shows. It preserves rank, and on rank-one \(Y_Z\)
it is \(Y_{\alpha Z}\). Composition is \(\Psi_D \Psi_F=\Psi_{DF}\), including
the normalization factors. Hence NC.13 has exactly the boundary
source map and its directed endpoint seam.

Write \(Y=I/2+r H\) with \(\operatorname{tr} H=0\) and H on the unit direction sphere. Let
\(\tau(r)\) be smooth, zero near \(r=0\) and one near the boundary. In NC.11
replace \(\alpha\),\(\beta\) by

\[
 \exp(\tau(r)\log\alpha),\qquad
 \exp(\tau(r)\log\beta),                                \tag{NC.14}
\]

using the fixed cut logarithms and the sphere direction H for Z.
This defines \(\widehat G_\theta(Y)\). It is independent of direction near
the centre because it is the identity there. Away from the centre,
the exact other-chart identity in NC.10 proves smoothness across
the infinity direction. Every factor is invertible; the inverse
product is smooth by the same statements. The actual full cap map is

\[
 \widehat C_\theta(Y,q_H)=
 \left(\Psi_{D_\theta}(Y),
 \frac{\widehat G_\theta(Y)q_H}
      {\|\widehat G_\theta(Y)q_H\|}\right).              \tag{NC.15}
\]

Its inverse first recovers Y by \(\Psi_{D^{-1}}\) and then uses the inverse
matrix at that recovered Y. This order proves the two compositions
even though \(\widehat G\) varies with Y. It restricts to NC.6-\/-11 on the
common sphere boundary. The forward cap seam is the actual map

\[
                    \widehat E=\widehat C_1^{-1}\widehat C_0.
                                                        \tag{NC.16}
\]

It restricts to NC.3-\/-4 with the full source and line functions.
The equality \(\widehat C_{1} \widehat E=\widehat C_{0}\) proves that \(\widehat C\) identifies this
specific ball/sphere mapping torus with the product cap. The inverse
is the inverse just constructed on the cut and agrees at the seam.
Thus the construction fixes the cap product, its source ball and
its normal sphere, rather than merely selecting a determinant degree.

\par

\medskip\noindent\textbf{NC.3. The executed normal derivative and the exact cap product}

\par\nobreak

On each of the eleven actual resolved moving-star charts let
\((v_{\rm minus},v_{\rm plus})\) be the local normal coordinates in physical order
\((-1,+1)\). NS.23 evaluates the original native coefficient characters
on the actual executed cone and obtains their tangent monomials and
full units \(e_{\rm minus},e_{\rm plus}\). Their normal exponents are exactly the
two standard basis vectors, so the map is the actual derivative.
In particular the point/O formulas NS.17 and the full \(T_{0}\) period are
carried in those coefficients. No source band or zero-section factor
is removed. In positive complex coefficients the exact map is

\[
 N_\theta^+=G_0(Z)
          \operatorname{diag}(\overline{e_-},e_+),\qquad
 G_0(Z)=\begin{pmatrix}1&-Z\\\overline Z&1\end{pmatrix},
                                                        \tag{NC.17}
\]

and on the other native chart it uses

\[
 G_1(z)=\begin{pmatrix}\overline z&-1\\1&z\end{pmatrix},
 \qquad G_1(1/Z)\operatorname{diag}(\overline Z,Z)=G_0(Z).
                                                        \tag{NC.18}
\]

The two charts and the NS normal-coordinate inverses prove agreement
on every common localization. The meromorphic tangent monomials at
interior divisors keep their actual divisor orders; NC.17 is used as
an invertible boundary map where P is nonzero.

Let \(N_\theta\) be the corresponding right-complex map before conjugating
its coefficient pair. The map from the resolved normal sphere into
the actual cut family is \(B_\theta=C_\theta^{-1} N_\theta\). Composing the
specified cap map, in its actual order and at its actual source point,
gives

\[
                     \widehat C_\theta B_\theta=N_\theta.
                                                        \tag{NC.19}
\]

Its inverse equality is \(B_\theta^{-1} \widehat C_\theta^{-1}=N_\theta^{-1}\). The
base equality follows from NC.13, and the full normal equality
follows from NC.11 and NC.15. Conjugating the two right-complex
coefficient columns gives the same identity for the positive complex
rank-two bundle NC.17. Thus the completed bundle has precisely the
periodic full \(N_\theta\) as its clutch in the product cap of NC.16.
Every P power and unit of C, B and \(\widehat C\) has occurred in a displayed
map and its inverse before this exact cancellation.

\par

\medskip\noindent\textbf{NC.4. The actual oriented replacement bordism}

\par\nobreak

The fixed surface in the cap product is \(C_{0}=-D^3_Z\times S^1_P\). Take
\(C=\mathbb P^1_Z\times D^2_P\) with its product orientation. The oriented product
boundary is

\[
 \partial(D^3_Z\times D^2_P)=C+C_0.                    \tag{NC.20}
\]

Round the corner in its two radial coordinates. Attach its rounded
facet to the boundary collar of the resolved fixed surface times an
interval. Carry the exact map N of NC.19 constantly across that
interval, with its full inverse. Over \(D^3\times D^2\) the positive
rank-two bundle is trivial in the conjugated held Cox frame. This
constructs a compact smooth five-dimensional bundle bordism from
the actual cap-completed fixed surface and normal bundle to the
closed replacement \(F_{\rm star}-C\) and its attached bundle. Equality of
their \(c_2\) and \(c_1^2\) follows by extending a Hermitian connection
over this explicit bundle bordism and applying Stokes to its closed
characteristic forms.

On the C face the two ordered line subspaces in that frame are
\(C(1,\overline{Z})\) and \(C(-Z,1)\). NC.18 gives their coefficient transitions
\(\overline{Z}\) and \(Z\), so their first Chern numbers on the complex Z sphere
are respectively +1 and -1. Retain these exact subspaces on the
replacement face. The diagonal derivative coefficients NC.17
preserve them and therefore glue the two positive normal lines.

For a P-disc radius epsilon, the opposite face -C has the coordinate
\(\eta=\overline{\xi}/\epsilon^{2}\) on its abstract disc. This real map has negative
determinant and equals 1/P on the common boundary \(\xi=P\). Thus -C
has its \(\eta\)-complex orientation. The closed replacement is the
actual nine-fold blowup of \(\mathbb P^1_Z\times \mathbb P^1_{\rm parameter}\) given by the
NS.23b fan and its recorded chronological coordinate blowup maps.
All nine centres lie in the original P disc. The meromorphic normal
dual sections and their chart inverses determine the divisor lines
on these blowups; they are retained for the separate intersection
calculation.

The sign of each parameter degree is fixed by this actual attachment.
For the boundary coefficient arrow \(z_{\rm cap}=P^n z_F\), the meromorphic
section with coefficient one on the original disc has coefficient
\(\eta^-n\) on the new cap chart. Its divisor has degree -n. For a
conjugated P coefficient, \(\overline{P}=P^{-1} |P|^{2}\) on the boundary; the
positive modulus has its real logarithm and is retained by the
explicit nonzero homotopy. This proves the sign convention directly
for the arrows in NC.17.

For completeness every unit \(u(P)\) used in these normal coefficients
has its actual nonzero value \(u(0)\) and a unique based logarithm \(L(P)\)
on the source disc: \(u(P)=u(0)exp L(P)\). The coefficient path
\(u(0)exp(t L(P))\) and its reciprocal are invertible throughout. On
the cap it extends by the radial cutoff of NC.14 in the global
projector onto the corresponding NC.18 line subspace. This is smooth
at the ball centre and has the reciprocal exponential as inverse.
The constant value remains in the attached frame. Therefore these
particular full units have zero parameter degree by an explicit
extending bundle map. No meromorphic P or Z character is included
among those units; their divisor orders stay in the eleven-chart
calculation.

\par

\medskip\noindent\textbf{NC.5. Relation to the preserved point-coordinate calculation}

\par\nobreak

The dense localization map from the preserved I2C coefficients is
\(a_{0}=Z^{-1} a_N\) and \(c_{0}=Z c_N\). Its normal positive matrix is
\(\operatorname{diag}(\overline{Z}^{-1},Z)\). At the regular point \(Z=0,a_{0}=0,c_{0}=1\) the old inverse
\(c_N=c_{0}/Z\) has a pole. The exact endpoint conjugation is NS.39, giving
the actual normal exponent pair \((-12,-24)\) in positive phases and
their sum \(-36\). NC.3-\/-19 use this actual extending native pair at
both vertices. The original source endpoint, full finite-cap square,
the NS phase-square maps and the physical winding thirty retain
their original arrows. This specifies precisely which local maps are
used for the new cap and which old formulas remain statements on
the common torus localization.


% End expanded verified cp36 input


% Expanded verified cp36 input: notes/part3_e6_selected_i3_m2_return_support_15.tex
\par\bigskip\noindent\textbf{P3K.111 supporting proof 15: Two actual own-cap moving normal bundles}\par\nobreak
\par

\medskip\noindent\textbf{Selected-I3 \(M_{2}\): the two actual moving normal bundles and their own caps}

\par\nobreak

This calculation uses the executed eleven-chart collision star NS and
the original fourteen-factor reference RB. Each normal bundle is first
attached to its own held-Cox cap. The collision cap is the complete
NC.1--4 construction. The reference cap uses the same displayed ball
construction with the reference coefficients derived in MN.3 below.
The final section gives the exact dense comparison and its vertex
pole. The finite fan reset that compares the two whole completed
families is retained as a separate construction. No total return is
assigned here.

\par

\medskip\noindent\textbf{MN.1. Every collision normal coefficient and its divisor}

\par\nobreak

Keep NS.1--23, including the actual point modules, full
\(\eta,c_q,f_0\), the literal \(T_0\) alignment and the original
selected height. The moving surface has quotient tangent rays
\[
 (1,0),(3,1),(2,1),(1,1),(0,1),(-1,1),(-2,1),(-3,1),
                  (-4,1),(-5,1),(-6,1),(-1,0).
 \tag{MN.1}
\]
The two entries are \((\operatorname{ord}Z,\operatorname{ord}P)\).
The actual NS matrices give all adjacent cones, their inverse maps
and the normal columns. Let \(C_j\), \(-6\le j\le3\), denote
the compact curve belonging to \((j,1)\). Write \(F=\sum_jC_j\),
the actual divisor of P. The two remaining rays are the horizontal
zero and infinity sections of the fixed sphere.

In physical order \((-1,+1)\), let \(e_-,e_+\) be the coefficient
map from the executed normal coordinates to the regular native
coordinates. On the infinity charts the full units are
\((c_q/\eta,\eta)\); on the zero charts they are
\((c_qf_0/\eta,\eta/f_0)\). Their entire tangent monomials are
the rows NS.23 evaluated on the actual aligned cone. The derivative
has no suppressed source term: P and the native fixed-curve
coordinate have zero normal order, and p contains the product of
the two normals. Hence differentiating any of these full base units
has zero first normal term on the fixed surface.

The divisors \(D_- =\operatorname{div}(e_-)\) and
\(D_+ =\operatorname{div}(e_+)\), interpreted as the actual maps
from the executed lines to the pulled-back native lines, have the
following full coefficient list. The two horizontal orders are zero.

\begin{longtable}[]{@{}
  >{\raggedright\arraybackslash}p{(\columnwidth - 20\tabcolsep) * \real{0.1000}}
  >{\raggedleft\arraybackslash}p{(\columnwidth - 20\tabcolsep) * \real{0.0900}}
  >{\raggedleft\arraybackslash}p{(\columnwidth - 20\tabcolsep) * \real{0.0900}}
  >{\raggedleft\arraybackslash}p{(\columnwidth - 20\tabcolsep) * \real{0.0900}}
  >{\raggedleft\arraybackslash}p{(\columnwidth - 20\tabcolsep) * \real{0.0900}}
  >{\raggedleft\arraybackslash}p{(\columnwidth - 20\tabcolsep) * \real{0.0900}}
  >{\raggedleft\arraybackslash}p{(\columnwidth - 20\tabcolsep) * \real{0.0900}}
  >{\raggedleft\arraybackslash}p{(\columnwidth - 20\tabcolsep) * \real{0.0900}}
  >{\raggedleft\arraybackslash}p{(\columnwidth - 20\tabcolsep) * \real{0.0900}}
  >{\raggedleft\arraybackslash}p{(\columnwidth - 20\tabcolsep) * \real{0.0900}}
  >{\raggedleft\arraybackslash}p{(\columnwidth - 20\tabcolsep) * \real{0.0900}}@{}}
\toprule\noalign{}
\begin{minipage}[b]{\linewidth}\raggedright
Curve
\end{minipage} & \begin{minipage}[b]{\linewidth}\raggedleft
\(C_3\)
\end{minipage} & \begin{minipage}[b]{\linewidth}\raggedleft
\(C_2\)
\end{minipage} & \begin{minipage}[b]{\linewidth}\raggedleft
\(C_1\)
\end{minipage} & \begin{minipage}[b]{\linewidth}\raggedleft
\(C_0\)
\end{minipage} & \begin{minipage}[b]{\linewidth}\raggedleft
\(C_{-1}\)
\end{minipage} & \begin{minipage}[b]{\linewidth}\raggedleft
\(C_{-2}\)
\end{minipage} & \begin{minipage}[b]{\linewidth}\raggedleft
\(C_{-3}\)
\end{minipage} & \begin{minipage}[b]{\linewidth}\raggedleft
\(C_{-4}\)
\end{minipage} & \begin{minipage}[b]{\linewidth}\raggedleft
\(C_{-5}\)
\end{minipage} & \begin{minipage}[b]{\linewidth}\raggedleft
\(C_{-6}\)
\end{minipage} \\
\midrule\noalign{}
\endhead
\bottomrule\noalign{}
\endlastfoot
\(D_-\) & −4 & −4 & −4 & −4 & −4 & −5 & −6 & −5 & −4 & −3 \\
\(D_+\) & 5 & 5 & 5 & 5 & 5 & 5 & 4 & 3 & 2 & 1 \\
\end{longtable}

The companion recomputes each entry from the original four-cone
matrix and the full NS character, and verifies agreement from both
adjacent charts. This uses all eleven charts. On the common native
overlap the two line coefficients are carried by
\(a_\infty=Za_0\), \(c_\infty=Zc_0\); these are the exact native
line frames used to interpret the displayed divisors.

The original point/O identities are
\[
 V_{\rm pt}=-P^{-2j}E A_{\rm pt}
 =P^{-2j}c^{-1}E C_{\rm pt}
 =P^{-2j-2}(x-P^2)E_OC_O.
 \tag{MN.2}
\]
Thus their horizontal point or zero divisor is cancelled by its
actual Cartier generator before the normal derivative is taken.
The full products E and \(E_O\) and the extra \(P^{-2}\) at O
remain in the coefficient map. The map to the native coefficient
is exactly NS.16, not a degree-zero replacement of MN.2. On the
fixed locus the remaining units are the entire functions of
\((P,p=0)\) above. Each has its nonzero constant value and its
based logarithm on the P-disc. NC.4 gives the reciprocal extending
bundle homotopy in the held-Cox line projectors. The meromorphic
P and Z factors are excluded from that homotopy and have exactly
the orders in the table.

\par

\medskip\noindent\textbf{MN.2. Actual intersection matrix and the collision NC bundle}

\par\nobreak

The wall relation for the two endpoints of the compact chain is
\(v_{\rm left}+v_{\rm right}=v_j\), and at each of its other
vertices it is \(v_{\rm left}+v_{\rm right}=2v_j\).
Consequently the full ten-by-ten intersection matrix in the table\textquotesingle s
order is the tridiagonal matrix
\[
 M=\begin{pmatrix}
 -1&1&0&\cdots&0\\
 1&-2&1&\cdots&0\\
 0&1&-2&\ddots&\vdots\\
 \vdots&\vdots&\ddots&\ddots&1\\
 0&0&\cdots&1&-1
 \end{pmatrix}.
 \tag{MN.3}
\]
All adjacent determinants are one. This matrix has kernel generated
by the full fibre F. A generic fixed-sphere point divisor h meets
\(C_0\) once and every other \(C_j\) zero times, by the actual
projection given by the native Z character. Therefore
\[
 D_-^2=-5,\qquad D_-D_+=2,\qquad D_+^2=-4,
 \qquad hD_-=-4,\qquad hD_+=5.
 \tag{MN.4}
\]
These equations are direct products with MN.3 and its specified
marked section; no cap or degree is selected to obtain them.

The actual nine blowdowns in NS.23b remove, in order,
\[
 (3,1),(2,1),(1,1),(-6,1),(-5,1),(-4,1),(-3,1),(-2,1),(-1,1).
 \tag{MN.5}
\]
Each removed ray is the sum of its current neighbours. The two
coordinate blowup maps are the written NS maps \((u,v)\mapsto(u,uv)\)
and \((u',v')\mapsto(u'v',v')\). Let \(E_1,\ldots,E_9\) be
the pullbacks to the full surface of the exceptional curves contracted
at steps 1 through 9 of MN.5, respectively. Their matrix is \(-I_9\), and they are
orthogonal to the two pullback classes h and u from the fixed sphere
and the closed parameter sphere. Here h and u are classes on the explicit closed replacement
\(\overline F=F_{\rm star}-C\) constructed in NC.4, with
\(h^2=u^2=0\) and \(hu=1\); the original-disc divisor F
represents u there. NC.19--20 carries its rank-two bundle and tangent
stable bundle across the specified five-dimensional bordism to the
actual cap-completed surface \(F_{\rm star}+C_0\). All characteristic
numbers below are transferred along that bordism. Removing each actual ray and
subtracting the coefficients of its neighbours gives
\[
 \begin{aligned}
 D_-&=-4u+E_4+E_5+E_6-E_7-E_8,\\
 D_+&= 5u-E_4-E_5-E_6-E_7.
 \end{aligned}\tag{MN.6}
\]
This independently reproduces every product in MN.4 through the
actual nine coordinate blowups.

NC.17--20 constructs the cap product with the full resolved normal
map \(G_0\operatorname{diag}(\bar e_-,e_+)\), its other native
chart and the literal \(\widehat C B=N\) equality. Its explicit
five-dimensional bundle bordism replaces that cap by the opposite
fixed-sphere/parameter-disc face. The two positive line subspaces
have native classes \(+h,-h\). The coefficient map from the
executed negative line to its native line is \(e_-\), so its
unconjugated class is \(-h-D_-\); its complex conjugate has
class \(h+D_-\). The positive line has class \(-h-D_+\).
This fixes the signs from the coefficient direction. Equivalently,
NC.4\textquotesingle s actual opposite parameter-disc map assigns degree \(-n\)
to the target coefficient \(P^n\) on the boundary.

The two ordered line classes of the positive normal bundle on the
NC.4 closed replacement are thus
\[
 \begin{aligned}
 x&=h-4u+E_4+E_5+E_6-E_7-E_8,\\
 y&=-h-5u+E_4+E_5+E_6+E_7.
 \end{aligned}\tag{MN.7}
\]
Their exact products are
\[
 x^2=-13,\qquad xy=-3,\qquad y^2=6,
 \qquad c_1=x+y=-9u+2E_4+2E_5+2E_6-E_8,
\]
\[
                 c_1^2=-13,\qquad c_2=-3.
 \tag{MN.8}
\]
On the closed replacement the tangent Chern numbers are
\(c_1(T\overline F)^2=8-9=-1\) and
\(c_2(T\overline F)=4+9=13\), from the same nine coordinate
blowups of \(\mathbb P^1\times\mathbb P^1\). Hence
\(\int_{\overline F}p_1=-1-2(13)=-27\). The explicit oriented
bundle bordism NC.19--20 transfers \(c_1(E)^2,c_2(E)\) and
\(p_1(TF)\) to the specified actual NC-capped moving component;
its numbers are therefore exactly \((-13,-3,-27)\). The global orbifold
Stokes and return orientations are not assigned by MN.8.

\par

\medskip\noindent\textbf{MN.3. The actual RB moving reference and its own cap}

\par\nobreak

Keep the full fixed-source selected functions of RB:
\[
 Q=\kappa q,\quad \alpha=P^2,\quad g=P^{-3},\quad
 \Gamma=\lambda_0\eta^6t^{-3}.
\]
Here the original source parameter equals \(s_*\ne0\), and the
independent scalar-disc coordinate is \(\sigma\), with suspension
\(\lambda_0\sigma^{12}\). The complete source and line maps
are exactly RB.2 and RB.16--17, with their genuine Cartier frames.

Take the y-factor chart of RB.13 with \(i=j=0\). Then
\[
 Z=x_0,\quad q=Zu_yv_y,\quad
 V_{\rm pt}=\sigma^{12}Z u_y.
\]
Its regular native coefficients are
\[
 c_R=Z^{-1}V_{\rm pt}=\sigma^{12}u_y,\qquad
 a_R=q/V_{\rm pt}=\sigma^{-12}v_y,
 \qquad q=Za_Rc_R.
 \tag{MN.9}
\]
The powers in MN.9 are boundary coefficient maps at nonzero scalar
radius. The resolved normal variables \(u_y,v_y\) themselves
remain regular at \(\sigma=0\).

On the x-factor chart of band \(i=-1,j=0\), put
\[
 z=\kappa\bar y_{-1}=Z^{-1},\qquad
 c_{R,\infty}=\sigma^{12}u_x,\qquad
 a_{R,\infty}=\kappa^{-1}\sigma^{-12}v_x.
 \tag{MN.10}
\]
The exact RB.17a overlap, including its inverse, gives
\[
 u_y'=Z^{-1}u_x,\qquad v_y'=(\kappa Z)^{-1}v_x,
\qquad
 (a_{R,\infty},c_{R,\infty})=Z(a_R,c_R).
 \tag{MN.11}
\]
Both native products are q. The inverse recovers q first, then
the same actual \(P(q),\kappa(q),\Gamma(q),\eta(q)\) and the
displayed reciprocal coefficients. Thus this is the actual regular
reference native sphere, derived before any numerical comparison.

Changing the source or chain representative uses the full RB maps.
In the genuine point coefficient the actual period is
\[
 B_\sigma:(Z,V_{\rm pt})\longmapsto
       (gZ,\sigma^{12}q\Gamma V_{\rm pt}).
\]
For the target label \(j+1\), division by
\((\sigma^{12}q)^{j+1}\) gives \(U_{j+1}'=\Gamma U_j\),
and the reciprocal coordinate gives \(W_{j+1}'=\Gamma^{-1}W_j\).
Thus the original q factor cancels only through this literal chain
arrow. At \(q=0\) its two resolved continuations are exactly RB.17.
The h map has \(i\mapsto i+1\) and line coefficient \(\alpha\).
The B map has \(j\mapsto j+1\), source coefficients
\((g,g^{-1})\), and the normals in RB.17. Their inverses carry
the same full \(\alpha,g,\Gamma,\kappa\) at the inverse source
point. On the moving fixed surface \(q=0\) these are nonzero constants
on the scalar disc. Thus their actual logarithms extend over that
disc, while the powers \(\sigma^{-12},\sigma^{12}\) in
MN.9--10 remain unchanged. The twelve other scalar-source pieces
are disjoint by the explicit \(v_{12}\ne0\) overlap of RB.15;
none of them is included in this moving surface.

For completeness its own cap is the NC.6--19 construction applied
to MN.9--11. The actual point endpoint is
\[
 (Z,V_{\rm pt})\longmapsto
        (g^{-4}Z,\sigma^{-48}\Gamma^{-4}V_{\rm pt}).
\]
Its native reference factors are
\[
 Z\mapsto g^{-4}Z,\quad
 a_R\mapsto\sigma^{48}\Gamma^4a_R,\quad
 c_R\mapsto g^4\sigma^{-48}\Gamma^{-4}c_R.
 \tag{MN.12}
\]
Their product is one. Use the full logarithms of
\(g^{-4}\) and \(g^4\sigma^{-48}\Gamma^{-4}\) in the two
NC cut factors. The two native chart equations, the q-preserving
velocity and its inverse, the positive Hermitian-ball map and the
whole held-Cox normal matrix are exactly those NC formulas with
these explicitly substituted functions. Their cap seam is
\(\widehat C_1^{-1}\widehat C_0\), and their resolved boundary
map obeys the same full identity \(\widehat C B=N_R\).
The actual original finite-cap square remains
\(\Phi_{\rho+4}q^4=h_\infty^4\Phi_\rho\), with all root
characters and sector integers, because the cut field is supported
in the selected tube. This constructs the reference\textquotesingle s own cap.

Its positive normal coefficient map is
\(G_0\operatorname{diag}(\overline{\sigma^{-12}},\sigma^{12})\)
in the zero native chart. The infinity chart has the additional
full nonzero coefficient \(\kappa(0)^{-1}\) on the negative line.
That constant and the full point/O units extend by the explicit NC
unit map. Applying the same opposite parameter-disc bundle bordism,
the closed reference replacement is
\(\mathbb P^1\times\mathbb P^1\), and its two ordered line
classes are exactly
\[
                    x_R=h-12u,\qquad y_R=-h-12u.
 \tag{MN.13}
\]
In particular
\[
 x_R^2=-24,\quad x_Ry_R=0,\quad y_R^2=24,\qquad
 c_1(E_R)^2=0,\quad c_2(E_R)=0,\quad p_1(TF_R)=0.
 \tag{MN.14}
\]
The same explicit bundle bordism transfers these three characteristic
numbers to the actual reference chart with its own specified cap. It has not been substituted for the collision
cap by a dense character map.

\par

\medskip\noindent\textbf{MN.4. The exact ordered bases and the remaining whole-pair map}

\par\nobreak

In the original \((\operatorname{ord}Z,\operatorname{ord}V_{\rm pt})\)
plane at selected height one, the physical edge is
\(0\to(0,1)\). After the literal collision \(T_0\) alignment,
its opposite triangle vertices are \((1,-1)\) and \((-1,2)\).
The reference charts MN.9--10 give vertices \((1,1)\) and
\((-1,0)\). Thus the full ordered matrices are
\[
 P_N=\begin{pmatrix}1&-1\\-1&2\end{pmatrix},\qquad
 P_R=\begin{pmatrix}1&-1\\1&0\end{pmatrix},\qquad
 P_N(1,1)^t=P_R(1,1)^t=(0,1)^t.
 \tag{MN.15}
\]
Both determinants are one. Their dense coefficient relation is
\[
 (Z,a_0,c_0)=(Z,Z^{-2}a_R,Z^2c_R),\qquad
 \begin{pmatrix}1&0\\2&1\end{pmatrix}P_N=P_R.
 \tag{MN.16}
\]
At the regular reference vertex \(Z=0,a_R=1,c_R=0\), the displayed
\(a_0\) has a pole. Therefore MN.16 does not itself supply the
regular vertex or cap identification. The original point/O arrows
cannot remove that pole on this selected node collar: their
numerator/denominator divisors already cancel in the actual frames
MN.2, and their remaining local factors are nonzero.

The full point product \((a_\rho,d_\rho)\) acts in the reference
native chart by
\[
 (Z,a_R,c_R)\mapsto
       (a_\rho Z,d_\rho^{-1}a_R,(d_\rho/a_\rho)c_R).
 \tag{MN.17}
\]
Its source-cut mismatch has \(a=e^{-6\pi i\rho}\) and second
native factor one. Its held quaternion is consequently
\(q_H\mapsto e^{3\pi i\rho}q_He^{3\pi i\rho}\), with both
odd signs retained. The collision right factor in NS.34 is fifteen;
the exact dense conjugation MN.16 accounts for their difference.
Both complete diagonal products retain the physical DZ factor
NS.37, by the same pointwise subtraction of the full logarithms.

The finite whole-pair continuation can use the exact q-preserving
Cox/ruling maps IF3.16--19 and the signed source/period resets
VI3.14--15, with the full current multipliers substituted. Those
maps, their native Hopf products, finite exterior and cap endpoints
must be composed to compare MN.8 with MN.14 in one prescribed
selected fan. This final composition is separate from the two
normal-bundle computations proved here.


% End expanded verified cp36 input


% Expanded verified cp36 input: notes/part3_e6_selected_i3_m2_return_support_16.tex
\par\bigskip\noindent\textbf{P3K.111 supporting proof 16: Twelve full lower/reference doubles with original cube}\par\nobreak
\par

\medskip\noindent\textbf{Selected-I3 \(M_{2}\): all twelve lower source surfaces and their full reference doubles}

\par\nobreak

Working proof, 18 November 2026. The source is
\(y^2=x^3+r_{\rm src}x^2+tx+t^2\), with \(G=(0,t)\),
\(M_2=\mathcal O([2G]-[O])\), chronological translation \(-3G\),
the good cubic cap \((m,n)=(3,1)\), and its prescribed finite character.
The source parameter is \(s=1/4-r_{\rm src}\). We retain SI2\textquotesingle s coupled
scalar \(\lambda=\lambda_0s^{12}\), scalar root \(r_0s^4\), original
selected height \(q=t-t_c(s)\), and complete annular cochain.

The twelve surfaces here are the actual \texttt{lower\_source\_boundary} face
orbits of the executed one-end atlas. They are closed against the
twelve scalar-source reference pieces in RB by the actual source and
line maps below. This gives the scoped physical characteristic pair
\[
 (k_M,\ell_M)=(-226,-1591)\quad\hbox{per lower double},\qquad
             (-2712,-19092)\quad\hbox{for all twelve}.
 \tag{IL.1}
\]
The moving selected surface, its Hopf cap, the whole reference fan
change RFB, and the final exterior/Jin transfer retain their own
specified arrows in the subsequent assembly.

\par

\medskip\noindent\textbf{IL.1. Original source, point line, normal coordinates and every period}

\par\nobreak

Use the full functions from SI2 and I2O:
\[
 Q=P^2R=P^3p,\quad p=1+v,\quad v=su_s/P,\quad
 t=Pu_t,\quad q=R c_q,\quad
 c_{12}=\frac{\lambda_0c_q\eta_2^6}{u_s^{12}u_t^3}.
 \tag{IL.2}
\]
Here
\(\eta_2=a_T\Theta'_Q(1)\Theta_Q(P^2)/\Theta_Q(P)^2\).
All products and inverse coefficients retain these entire functions
on their specified source domains. With \(X=Z^{-1}\) and
\(Y=V_{\rm reg}/Z\), the full two periods are
\[
 L:(X,Y)\longmapsto(P^3pX,PpY),\qquad
 T:(X,Y)\longmapsto(p^{-1}X,P^{12}v^{12}c_{12}Y).
 \tag{IL.3}
\]
They fix \(P,v\); their inverses use the reciprocal complete factors.
The two actual original line coefficients are
\[
 V_{\rm pt}=\eta_2Y/X,\qquad
 v_2=\frac{\eta_2P^2Y}{Xf_{P^2}(pP^3/X)}.
 \tag{IL.4}
\]
The latter is a morphism in the original point/O Cartier modules.
Its reciprocal is \(Y=\eta_2^{-1}P^{-2}X
f_{P^2}(pP^3/X)v_2\), at the same recovered source point.

In the ordered lattice
\((\operatorname{ord}X,\operatorname{ord}Y,
\operatorname{ord}P,\operatorname{ord}v)\), the twelve minimal
physical fixed faces are
\[
 E_l=(0,l,0,1),\quad E_{l+1}=(0,l+1,0,1),
                  \qquad 0\leq l\leq11.
 \tag{IL.5}
\]
Their stabilizer is generated by the full \(L\) arrow. The independent
\(T\) arrow replaces \(l\) by \(l+12\), with its entire line multiplier;
the twelve residue labels therefore remain twelve separate surfaces.
The physical cocharacter is \((0,1,0,0)\). In a chart ordered by
\((E_l,E_{l+1},r_1,r_2)\), its weights are exactly \((-1,+1,0,0)\).
The source projection is the integral map
\((x,y,r,z)\mapsto(x,r)\). Both source functions \(X,P\) have zero
normal order.

Write the meromorphic parent normal coefficients as
\[
 \alpha_l=Yv^{-l},\qquad \beta_l=v^{l+1}/Y,\qquad
                     \alpha_l\beta_l=v.
 \tag{IL.6}
\]
The physical weights of \(\alpha_l,\beta_l\) are \(+1,-1\).
For a full ray \((x,y,r,z)\), their two source exponents are
\(y-lz,(l+1)z-y\). On every executed chart these are the two rows
of the full monomial map to the parent coefficients, with normal entries respectively
\((0,1)\) and \((1,0)\). These rows are retained for all 3,888 charts.
In particular the normals being calculated are the original total-line
normal bundles restricted to the fixed surfaces.

The full normal period squares are
\[
 L:(\alpha_l,\beta_l)\longmapsto
       (Pp\,\alpha_l,(Pp)^{-1}\beta_l),
 \tag{IL.7}
\]
\[
 T:(\alpha_l,\beta_l,l)\longmapsto
   (P^{12}c_{12}\alpha_l,P^{-12}c_{12}^{-1}\beta_l,l+12).
 \tag{IL.8}
\]
The saved period-unit key \(f_1\) denotes the full \(c_{12}\)
of IL.2 in this one-end execution. Its character columns are
\((0,0,0,0)\) for \(L\) and \((0,1,0,0)\) for \(T\);
the \(p\) columns are \((1,1,0,0)\) and \((-1,0,0,0)\).
Every executed lower endpoint arrow carries those columns through
its full source and target inverse matrices, with its reciprocal
unit vector also retained.
The \(p\) and \(c_{12}\) factors are retained before restriction to
the fixed surface. A derivative of a source unit multiplying a
vanishing normal coefficient contributes zero additional normal
term at that surface. Also \(v=\alpha_l\beta_l\) has zero first
normal derivative there. Thus the actual normal derivatives are
precisely the displayed diagonal coefficients, together with the
full recorded source derivative.

\par

\medskip\noindent\textbf{IL.2. All source charts, normal walls and ordinary blowdowns}

\par\nobreak

Each of the twelve executed stars has 324 source charts and 325
ordered rays in \(0\leq x\leq3r\), with endpoints \((0,1),(3,1)\).
Every consecutive pair has determinant one and is an actual
executed source chart. The full endpoint lifts are
\[
       (0,12,1,0),\qquad(3,13,1,0)
                    =L(0,12,1,0).
 \tag{IL.9}
\]
The endpoint identification includes the coefficient \(p\) in
IL.3. It is evaluated at the actual target source coordinates.
The quotient identifies these ends and gives a central cycle of
324 compact curves. Its three surviving minimal source rays are
\((0,1),(1,1),(2,1)\); the last endpoint \((3,1)\) is their
periodic copy of \((0,1)\).

For a central ray \(r\), let \(r_-,r_+\) be its actual neighbours,
using the full inverse \(L\) at the first endpoint. The unimodular
four-dimensional matrix with columns
\((E_l,E_{l+1},r,r_-)\) gives the exact wall equation
\[
 -r_- -r_+=d_-E_l+d_+E_{l+1}+c\,r.
 \tag{IL.10}
\]
The last inverse-matrix entry is zero. The normal degrees on that
compact curve are \((d_-,d_+)\) and its source self-intersection is
\(c\). To see the direction, express the common meromorphic
coefficients IL.6 in the two smooth adjacent monomial charts.
Their coefficient transition is the corresponding integer power
of the inverse base coordinate. The divisor of that transition is
exactly \(d_-\), respectively \(d_+\), with the physical order in
IL.5. This also gives the reciprocal transition on the opposite
overlap. The complete nonzero unit factors are the restrictions of
IL.3; their two defining discs and reciprocal maps remain in that
calculation.

The executed blowdown sequence removes 321 rays from each source
star. At each step the removed primitive ray is the sum of its
two present neighbours, and its self-intersection is \(-1\).
The actual maps on its two affine charts are
\[
       (u,w)\longmapsto(u,uw),\qquad
       (u,w)\longmapsto(uw,w).
 \tag{IL.11}
\]
Their inverse maps on the complementary principal opens divide
the second, respectively first, coefficient by the retained
nonzero coordinate. The exceptional projective line is contracted
by these proper maps. The ordered list of the two neighbours,
the ray, both normal degrees, and both coefficient maps is saved
for every step, independently for every label \(l\).
Their composition is the actual source blowdown to the minimal
\(I_3\) source, with identity map on the retained outer boundary.
Every original source function is its pullback through the full
I2P/E81 map.

If \(d_\pm=N_\pm\cdot E\) at one step, tensoring \(N_\pm\) by
\(\mathcal O(d_\pm E)\) makes its exceptional degree zero. Its
two exact Cartier coefficient maps descend to the parent charts:
in the chart where \(E\) is \(u=0\), the map to the parent
coefficient is multiplication by \(u^{d_\pm}\), with reciprocal
module arrow on the common principal open. Consequently
\[
 N_\pm=\pi^*N_{\pm,0}\otimes\mathcal O(-d_\pm E),\qquad
 \Delta(c_-^2,c_-c_+,c_+^2)=(-d_-^2,-d_-d_+,-d_+^2).
 \tag{IL.12}
\]
Degrees on each adjacent strict curve add \(d_\pm\) in the
parent. The executable performs this integral operation together
with each source blowdown.

The final minimal-source degree vectors are
\[
 \deg N_-=(-1,0,1),\qquad \deg N_+=(-1,1,0)
                  \quad\hbox{on }(C_0,C_1,C_2).
 \tag{IL.13}
\]
The 321 chronological exceptional degree pairs consist of 49
copies of \((-1,0)\), 45 of \((0,-1)\), 133 of \((-1,-1)\),
and 94 of \((0,0)\). Summing the actual IL.12 maps gives
\[
       \Delta(c_-^2,c_-c_+,c_+^2)=(-182,-133,-178).
 \tag{IL.14}
\]
The source fan agrees for the twelve labels, while their full
normal lifts and period representatives remain separately stored.

\par

\medskip\noindent\textbf{IL.3. Relative lifts and the complete original tangent comparison}

\par\nobreak

With coefficient zero on \(C_0\), the integral source matrix is
\[
 I=\begin{pmatrix}-2&1\\1&-2\end{pmatrix}.
 \tag{IL.15}
\]
The degree vectors on \(C_1,C_2\) in IL.13 give the rational
relative lifts \(I^{-1}(0,1)^t=(-1/3,-2/3)^t\) and
\(I^{-1}(1,0)^t=(-2/3,-1/3)^t\). Their three pairings are
\[
        (-2/3,-1/3,-2/3).
 \tag{IL.16}
\]
These lifts retain the integral line bundles and their order-three
boundary characters. The full fibre \(F\), whose multiplicities
are the \(P\)-orders of the rays, is the intersection radical.
Both normal bundles have degree zero on \(F\). Adding its
coefficient to a relative lift therefore changes none of the
three pairings.

Combining IL.14 and IL.16 gives the collision-side pairings
\[
        (c_-^2,c_-c_+,c_+^2)=(-548/3,-400/3,-536/3).
 \tag{IL.17}
\]
An independent computation uses the complete 324-component cycle.
After the marked \(C_0\) is omitted its integral intersection
matrix is tridiagonal and negative definite. Exact successive
elimination of its diagonal pivots, followed by backward
substitution, solves the two full degree equations. Every row
is checked in the original matrix. The three resulting products
are exactly IL.17 for each of the twelve labelled stars.

The same maps IL.11 determine the tangent framing. Their
Jacobians have the single exceptional factor \(u\), respectively
\(w\), so \(K_{\widetilde S}=\pi^*K_S+E\), with
\(E^2=-1\) and zero pairing against the pullback class. The
exceptional projective line increases the relative Euler number
by one. Hence each step changes \(p_1=c_1^2-2c_2\) by \(-3\),
with the identity outer boundary map, and
\[
            \Delta\!\int p_1(TS)=-963.
 \tag{IL.18}
\]
The original volume in I2O is
\[
 -\frac{2J_{ts}}{a_T}P\,dP\wedge dv\wedge d\log X\wedge d\log Y.
 \tag{IL.19}
\]
Its entire unit, order \(2r+z-1\) on each ray, and all
coefficient derivatives are retained in the chart records.
IL.18 follows from the actual Jacobians relative to that fixed
outer volume and does not change its boundary marking.

\par

\medskip\noindent\textbf{IL.4. Exact point/O and source maps to the fixed reference}

\par\nobreak

The source bands \(x=Q^jZ,y=Q^{1-j}/Z\), with \(xy=P^2R\),
carry the proper graph of \((x,P^2)\):
\[
 x=P^2a,\quad ay=R;\qquad xc=P^2,\quad y=Rc.
 \tag{IL.20}
\]
The second chart is resolved by
\[
 P=uw,\ x=u,\ c=uw^2,\ y=Ruw^2;\qquad
 P=uw,\ x=u^2w,\ c=w,\ y=Rw.
 \tag{IL.21}
\]
These arrows preserve \(P,R\), and hence the original \(s,t\).
Their reciprocal overlaps are the actual blowup overlaps. At
\(s=0\), \(R=P\), the source coordinates
\[
       x_i=X/P^i,\qquad y_i=P^{i+1}/X,\qquad x_iy_i=P
 \tag{IL.22}
\]
identify their three node charts with the marked minimal
\(I_3\) source. The point section lies on \(C_1\): in IL.20
it is \(a=1\), so \(X=P\) after the actual \(L\) reindex.
The zero section lies on \(C_0\), at \(X=1\); the alternate
representative at \(X=P^3\) is carried back by the same complete
source and line period.

For \(b=P^2,e=R\), the full theta expression has the factorization
\[
 b f_b(x)=(b-x)E_{b,e}(x,y),
 \tag{IL.23}
\]
where \(E_{b,e}\) is I2P.3\textquotesingle s complete normally convergent product
on this node chart. In the original point modules
\[
 v_2=(a-1)^{-1}A=(c-1)^{-1}C,\quad C=-cA,
\]
\[
 V_{\rm reg}=-\eta_2^{-1}P^{-2j}E_{b,e}A
     =\eta_2^{-1}P^{-2j}c^{-1}E_{b,e}C .
 \tag{IL.24}
\]
At the zero section the actual generator is \(x-1\), and
\[
 V_{\rm reg}=\eta_2^{-1}P^{-2j-2}(x-P^2)E_OC_O.
 \tag{IL.25}
\]
The full \(E_O\), the additional \(P^{-2}\), every reciprocal,
and the original Miller arrow \(r_2=f_2^{-1}\) remain in these
formulas. Their point/O horizontal factors are interpreted in
their indicated line modules through coincident section points.

On the retained boundary of a fixed small original \(t\)-disc,
the directed normal map to the scalar-source chain is
\[
 A_l=H_{+,l}\alpha_l,\quad B_l=H_{-,l}\beta_l,\qquad
 H_{+,l}=\eta_2u_s^lP^{-l}X^{-1},\quad
 H_{-,l}=\eta_2^{-1}XP^{l+1}u_s^{-l-1}.
 \tag{IL.26}
\]
Their product is \(P/u_s\), so the map takes
\(\alpha_l\beta_l=v\) to \(A_lB_l=s\). Its inverse uses the
two reciprocal coefficients at the recovered source point.
Under \(L\) its two target factors are \(P^{-2},P^2\).
Under \(T,l\mapsto l+12\), they are
\[
 p u_s^{12}c_{12}
       =P^2\lambda_0q\eta_2^6/t^3
 \quad\hbox{and its reciprocal}.
 \tag{IL.27}
\]
Thus the whole selected height, source phase, Miller coefficient
and line period are retained. At \(p=1\) the exact central
identity is \(\eta_2^6=4u_t^2\), making IL.27 equal to
\(4\lambda_0\); the equality before this restriction remains
the actual boundary arrow. For an arbitrary source-cell representative
write the original chain label as \(12j+l\). The actual inverse period is
\[
 X_*=p^jX,\quad
 Y_*=P^{-12j}v^{-12j}c_{12}^{-j}Y,\quad
 V_{{\rm pt},*}=(P^{12}v^{12}pc_{12})^{-j}V_{\rm pt}.
 \tag{IL.27a}
\]
It takes the positive parent coefficient to
\(\alpha_*=(P^{12}c_{12})^{-j}\alpha_{12j+l}\).
Thus the two original reference-chain coefficients acquire,
in addition to IL.26 evaluated at \(X_*\), the respective
factors \(\gamma^j,\gamma^{-j}\), where
\[
                   \gamma=pu_s^{12}c_{12}
                         =P^2\lambda_0q\eta_2^6/t^3.
 \tag{IL.27b}
\]
This follows by substituting
\(P^{12}v^{12}pc_{12}=s^{12}\gamma\) in
IL.27a, and retains every chronological source and line power.

The final target of IL.26 is the actual fixed-source reference
at \(s_*\ne0\). We spell out its source transport. In the
first graph chart of IL.20,
\(s=(ay-P)/u_s(P,ay)\) has \(\partial_Ps\) a unit of
central value \(-1\). In either IL.21 chart,
\(s=(R-uw)/u_s(uw,R)\) has \(\partial_Rs\) a unit of
central value \(+1\). Take respectively
\((\partial_Ps)^{-1}\partial_P\) and
\((\partial_Rs)^{-1}\partial_R\) as local lifts.
On the original \(t\)-boundary the prescribed source-period map is
\[
 Z_s=\exp\!\left(b\,\operatorname{Log}(Q_s/Q_0)\right)Z_0,
 \quad
 V_{{\rm pt},s}=\exp\!\left(
 b\,\operatorname{Log}(P_s^2/P_0^2)\right)V_{{\rm pt},0}.
 \tag{IL.28}
\]
The marked real coordinate \(b\) increases by one under the
source deck. The two target factors are then exactly \(Q_s\)
and \(P_s^2\). The full ratios have their source-path logarithms
fixed to zero at \(s=0\), and the complete IL.24--25 module
arrows are composed with IL.28.

A relative partition with sum one joins these local source
lifts, agreeing with the derivative of IL.28 on a boundary
collar. The real field along a source path \(s(a)\) is
\(\dot s X+\overline{\dot s X}\). It is tangent to the
retained \(t\)-boundary and exists for the whole interval
because the source graph over the closed domain is proper.
Its backward flow is the inverse. Write the resulting source
map as \(T_s:S_0\to S_s\), and \(K_s=T_{s_*}T_s^{-1}\).

For original vector frames \(e_j=g_{ij}e_i\), set
\[
 a_i=-\sum_k\chi_k\,Xg_{ik}/g_{ik},\qquad
 a_j=a_i+Xg_{ij}/g_{ij}.
 \tag{IL.29}
\]
The horizontal vector in the coefficient \(w_i\) is
\(X-a_iw_i\partial_{w_i}\). A relative cutoff makes it agree
with the complete boundary derivative of IL.24--28. Its
linear ODE, and its inverse backward ODE, give the actual
line transport; the reciprocal ODE gives the dual.
Consequently the full boundary map is
\[
 (z,\alpha_l,\beta_l)\longmapsto
 \bigl(K_sz,\ K_s^{\rm line}(H_{+,l}\alpha_l),
                   K_s^{\rm dual}(H_{-,l}\beta_l)\bigr).
 \tag{IL.30}
\]
The inverse first applies \(K_s^{-1}\), then the reciprocal
IL.26 factors. On every period translate, the construction is
defined by conjugating with the full IL.3 map and its inverse;
this includes its source-dependent derivative.

The actual good cubic cap retains SI2.28\textquotesingle s ordinary-to-cap
module map, its deck character and all sector integers.
For the original height \(q\) its exact four-turn square is
\[
           \Phi_{\rho+4}\,q^4=h_\infty^4\Phi_\rho,
                 \qquad h_\infty=1-t_c(s)z^3.
 \tag{IL.31}
\]
The inverse closing uses \(q^{-4},h_\infty^{-4}\).
IL.30 is composed with that point-line cap map and its
reciprocal on the prescribed annulus. Its relative source
lift extends through the actual proper cap charts with that
fixed boundary derivative. Hence the full parameter square
is preserved.

All three vanishing discs remain in these source bodies.
The two-node unselected source and the selected node are the
three actual charts IL.22. The fourteen-factor RB resolution
has their scalar-source fixed components at
\(\alpha_l=\beta_l=0\); its overlap with the moving selected
fixed curve requires the nonzero coordinate \(v_{12}\).
It consequently retains each selected vanishing disc in the
scalar-source component. These explicit overlaps give the
source boundary inclusion
\[
 H_1(\partial S)=\mathbb Z\langle b,\alpha_t\rangle
                  \oplus\mathbb Z/3\langle a\rangle
 \longrightarrow H_1(S)=\mathbb Z\langle b\rangle,\qquad
                  (a,\alpha_t,b)\mapsto(0,0,b).
 \tag{IL.32}
\]
Indeed the three plumbing maps give
\(\alpha_t b\alpha_t^{-1}=a^3b\), and each vanishing
disc kills the same primitive \(a\); the retained O-disc
kills \(\alpha_t\). This also specifies the boundary on
which IL.30 is used.

\par

\medskip\noindent\textbf{IL.5. The actual cube frame and the twelve closed normal pairings}

\par\nobreak

The reference positive normal in the scalar chain is the
genuine point line \(M_2\), and its negative normal is
its reciprocal. IL.24--25 show the central point/O degree
vector \((-1,1,0)\). To retain its full boundary class,
write
\[
 \tau_L=\frac{\log|X|}{\log|Q|},\qquad
 W_p=e^{-2\tau_L\operatorname{Log}p}X^2V_{\rm pt}^3.
 \tag{IL.33}
\]
The logarithm of \(p\) is the actual small boundary logarithm.
Under \(L\), \(\tau_L\) increases by one and
\(X^2V_{\rm pt}^3\) acquires \(p^2\), so IL.33 is exactly
invariant. Its inverse is a vector frame of \(M_2^{\otimes3}\).
At \(p=1\), the complete theta identity in the actual
point/O modules gives this frame divisor
\[
                    6F-2C_1-C_2.
 \tag{IL.34}
\]
More explicitly it is
\[
 \frac{P^6X}{(X-P)^3}
       E_{P^2,P}(P^3/X,X)^{-3}\sigma_2^{\otimes3},
 \tag{IL.35}
\]
where \(\sigma_2\) is the original canonical meromorphic
section. The arguments of the entire product are exactly
the band-one source coordinates \((x,y)=(P^3/X,X)\).
Indeed IL.23 and the literal \(Q=P^3\) reindex give
\[
 f_{P^2}(1/X)
      =P^{-2}(1-P/X)E_{P^2,P}(P^3/X,X).
 \tag{IL.35a}
\]
Its cube, multiplied by \(X^2\), has IL.35 as its inverse
coefficient in the canonical section. The point/O horizontal factors cancel through
IL.24--25. Multiplication by \(P^{-6}\) gives the literal
coefficient-zero-on-\(C_0\) frame with divisor
\(-2C_1-C_2\), retaining that entire fibre clutch.

Thus the positive reference relative lift is
\[
 D_+=-(2C_1+C_2)/3,\qquad D_+^2=-2/3.
 \tag{IL.36}
\]
Its boundary character has exact order three, since
\(3(1,0)^t=I(-2,-1)^t\) while \(I^{-1}(1,0)^t\)
has denominators three. For the actual frame IL.33 the
lift is \(2F+D_+\), with the same square. The negative
reference line uses the reciprocal frame. Therefore
\[
 (c_-^2,c_-c_+,c_+^2)_{\rm ref}=(-2/3,+2/3,-2/3).
 \tag{IL.37}
\]

This computation retains the remaining full boundary gauges.
By IL.32,
\[
 H^1(\partial S;\mathbb Z)/\operatorname{im}H^1(S;\mathbb Z)
                    =\mathbb Z\langle\alpha_t^*\rangle.
 \tag{IL.38}
\]
Its connecting class is the primitive fibre, realized by
the vector-frame clutch \(e_R=(P/\epsilon)e_L\).
The original \(t\)-meridian has the same class through
\(t=Pu_t\), with its full \(u_t\) factor retained.
The \(b\)-character extends over \(S\), and a zero-period
boundary ratio has a logarithm extended by a collar cutoff.
Every actual gauge in IL.26--31 therefore changes a relative
lift only by the retained fibre multiple. It changes no
square or mixed pairing because \(F^2=F\cdot N_\pm=0\).

Choose the norm making the inverse IL.33 cube frame unit,
and take its actual third roots on a simply connected
boundary cover. Their transition functions are constant
elements of \(\mu_3\); declaring these frames parallel
gives the actual flat order-three boundary connection.
The reciprocal connection belongs to the negative line.
Transport both through IL.30. The interpolation to any
previous prescribed boundary connection is placed on the
opposite collar with its coordinate reversed. This is the
identical full interpolation pulled back by the reversal;
its complete transgression cancels pointwise in the glued
double.

The collision negative degree vector on \(C_1,C_2\)
differs from the reference negative vector by
\((1,1)^t=I(-1,-1)^t\). Thus their boundary torsion
characters agree by the actual coefficient arrow IL.26.
The positive characters agree as well. Gluing through
IL.30, using the cube connections and the retained fibre
gauges, subtracts IL.37 from IL.17 and gives
\[
       (c_-^2,c_-c_+,c_+^2)_{\rm double}
                           =(-182,-134,-178).
 \tag{IL.39}
\]
For the tangent bundle, the minimal collision source is
carried to the entire reference body by the explicit
relative source flow. Its derivative and inverse determine
matching boundary tangent connections. Their identical
opposite-collar interpolation cancels. The remaining
contribution is exactly the 321 actual blowups in IL.18:
\[
                 \int_{\rm double}p_1(TS)=-963.
 \tag{IL.40}
\]
On the fixed surface \(s=A_lB_l=0\), the first normal
derivative of \(s\) is zero. A circle-averaged source
partition has no linear term of nonzero physical weight.
The normal derivative of the transport is consequently
the complex diagonal map IL.26, in the original Cartier
modules, and adds no unrecorded tangent-normal mixing.

\par

\medskip\noindent\textbf{IL.6. Physical characteristic evaluation and exact provenance}

\par\nobreak

Let \(a=c_1(N_-)\), \(b=c_1(N_+)\), with physical weights
\((-1,+1)\). Use \(u=c_1(U)\), \(U=M^*\), so the
physical class is \(c_1(M)=-u\). The actual equivariant
normal Euler factor is \((-u+a)(u+b)\). Expanding its
inverse and the normal tangent Pontryagin factors gives
the surface contributions
\[
 k_M=\int(a^2-ab+b^2),\qquad
 \ell_M=\int\bigl(p_1(TS)+(a+b)^2\bigr).
 \tag{IL.41}
\]
In detail, the degree-four coefficient of
\((-u)^3/[(-u+a)(u+b)]\) is \(a^2-ab+b^2\).
The corresponding coefficient of
\[
 \frac{(-u)\bigl(p_1(TS)+(-u+a)^2+(u+b)^2\bigr)}
                {(-u+a)(u+b)}
\]
is \(p_1(TS)+(a+b)^2\). These are evaluated with
the complex fixed-surface orientation already used in
IL.10--40.

Substitution of the three complete normal pairings and
the tangent value gives
\[
 k_M=-182+134-178=-226,\qquad
 \ell_M=-963-182-268-178=-1591.
 \tag{IL.42}
\]
All twelve fixed-face orbits have the same pair by their
separately executed full maps, hence IL.1. Their distinct
labels and period representatives are retained.

The primary \path{code/part3/e6_selected_i3_m2_lower_normals.py}
records all 3,888 inverse chart matrices, original
source/normal/volume rows, all central wall maps, and
3,852 chronological ordinary blowdowns. It also solves
both full intersection systems independently. The initial
collision-only execution passed 66,128 exact assertions;
its complete source/output is preserved under
\path{records/attempts/e6_selected_i3_m2_lower_collision_only/.}
The companion relative certificate checks the full
IL.26--38 period and cube-frame identities, the exact
reference subtraction, the physical Euler expansion,
and each of the twelve separate resulting rows.

The whole scalar-reference fan comparison RFB changes
the complete reference pair from \(T_1\) to \(T_{-1}\)
by \((-72,0)\). IL.1 uses the actual RB \(T_1\)
reference pieces and its own cap maps. The whole RFB
correction is applied once in the complete outer
assembly, after these lower pieces and the separately
computed moving/reference cap have been included.


% End expanded verified cp36 input


% Expanded verified cp36 input: notes/part3_e6_selected_i3_m2_return_support_17.tex
\par\bigskip\noindent\textbf{P3K.111 supporting proof 17: Whole reference fan bridge with exact Jin type}\par\nobreak
\par

\medskip\noindent\textbf{Selected-I3 \(M_{2}\): the complete scalar-reference fan and cap comparison}

\par\nobreak

Working proof, 18 November 2026. This note retains the fixed original
source \(s=s_*\ne0\), \(M_2=\mathcal O(2G-O)\), the chronological
translation \(-3G\), the cubic good cap of norm one with its prescribed
finite character, and every SI2 annular integer. The independent
scalar-disc coordinate is \(\sigma\), its root is \(r_0\sigma^4\),
and the suspension scalar is \(\lambda_0\sigma^{12}\).

The reference constructed in RB has fan \(T_1\). The executed
collision native star NS has fan \(T_{-1}\). We construct the two
actual reverse \(C_1\) changes on the nonzero scalar boundary circle,
including their Cox correspondences, and compare their completed
circle quotients with the complete source, line, cap and outer maps.
The resulting \textbf{whole physical characteristic-pair correction} is

\[
 (k_{M,-1}-k_{M,1},\ell_{M,-1}-\ell_{M,1})=(-72,0).
 \tag{RFB.1}
\]

The pair in RFB.1 belongs to this independent scalar word \((0,4)\).
All reference fixed components and the reference moving cap are
included in its whole pair. No individual moving-cap value is obtained
by subtracting RFB.1 from a partial fixed-component ledger. The
comparison is a construction on the actual nonzero-parameter family
and its completed quotient. The analytic scalar-disc reference RB,
including its fourteen-factor charts, is kept unchanged.

\par

\medskip\noindent\textbf{RFB.1. The named periods, fan labels and the full endpoint}

\par\nobreak

Write \(q=t-t_c(s_*)\), and retain the complete RB functions

\[
 \begin{gathered}
 Q=\kappa q,\quad \kappa=P^2/c_q,\quad \alpha=P^2,\quad g=P^{-3},\\
 \eta_2=a_T\Theta'_Q(1)\frac{\Theta_Q(P^2)}{\Theta_Q(P)^2},\qquad
 \Gamma=\lambda_0\eta_2^6t^{-3},\qquad \eta=\sigma^{12}\Gamma .
 \end{gathered}
 \tag{RFB.2}
\]

On the selected neighbourhood of this fixed source, all functions
except the displayed \(q\) and \(\sigma\) factors are nonzero
holomorphic functions of the original \(q\). Their logarithms are
chosen on one common smaller \(q\)-disc. In particular \(\eta_2\)
is the full Miller/theta unit; none of its factors is replaced by
its central value. The original periods are

\[
 h:(Z,V)\mapsto(\kappa qZ,\alpha V),\qquad
 B:(Z,V)\mapsto(gZ,\eta qV),\qquad V=V_{\rm pt}.
 \tag{RFB.3}
\]

The positive scalar-circle coordinate is \(\sigma=\rho e^{2\pi i\theta}\),
with fixed \(\rho>0\). Hence the two line windings of RFB.3 are
\((0,12)\), while \(\kappa,g\) have winding zero.

For each of the three integers \(j=1,0,-1\), use the exact characters
and period generators

\[
 Z_j=Z,\qquad V_j=Z^{-j}V,\qquad h_j=hB^j,\quad B_j=B.
 \tag{RFB.4}
\]

Their inverses are \(V=Z_j^jV_j\) and \(h=h_jB_j^{-j}\). Substitution
in both original periods gives the full, untruncated functions

\[
 \kappa_j=\kappa g^j,\qquad
 \alpha_j=\alpha\eta^j\kappa^{-j}g^{-j^2},\qquad
 \eta_j=\eta g^{-j},
 \tag{RFB.5}
\]
\[
 h_j:(Z_j,V_j)\mapsto(\kappa_jqZ_j,\alpha_jV_j),\qquad
 B_j:(Z_j,V_j)\mapsto(gZ_j,\eta_jqV_j).
 \tag{RFB.6}
\]

The line windings are \((12j,12)\). All negative powers in RFB.5
are retained on the nonzero scalar circle. The actual affine
extensions are the primitive character charts on the indicated
fan, with full six-chart rows

\[
\begin{array}{c|ccc}
 &u&v&w\\\hline
 C_0&\kappa_j\alpha_jbc&\alpha_j^{-1}ab&\kappa_j^{-1}b^{-1}\\
 C_1&g\eta_jbc&\eta_j^{-1}c^{-1}&g^{-1}ac\\
 C_2&a^{-1}&ab&ac .
\end{array}
\qquad Z_j=b,\quad V_j=c,\quad q=abc=uvw.
 \tag{RFB.7}
\]

For each row recover the original \(q\) from the target triple
product and evaluate every inverse coefficient at that recovered
value. These maps and inverses retain the full normal jets.

Put \(u_j=(1,j)\), \(v_j=(-1,1-j)\), so \(u_j+v_j=\delta=(0,1)\).
The two positive-diagonal triangles are
\(\{0,u_j,\delta\}\), \(\{0,v_j,\delta\}\).
Translating the second by the actual \(h_j\) sends it to
\(\{u_j,\delta,u_j+\delta\}\). These are precisely the two
triangles of \(T_j=T(e_1+je_2,e_2)\); the translation carries the
complete RFB.6 line coefficient. Thus the RB matrix
\(\left(\begin{smallmatrix}1&-1\\1&0\end{smallmatrix}\right)\)
is \(T_1\), and the NS matrix
\(\left(\begin{smallmatrix}1&-1\\-1&2\end{smallmatrix}\right)\)
is \(T_{-1}\), with their exact point marking.

The actual ordinary forward endpoint is \(q^4\). After its specified
\(B^{-4}\) reindex the bounded endpoint is

\[
 E:(Z,V)\mapsto(AZ,DV),\qquad A=g^{-4},\quad D=\eta^{-4}.
 \tag{RFB.8}
\]

In the \(j\)-th native \(C_0\) chart set

\[
 a_j=q/(Z_jV_j),\quad c_j=V_j,\quad q=Z_ja_jc_j,\qquad
 D_j=D A^{-j}=\eta^{-4}g^{4j}.
 \tag{RFB.9}
\]

The full endpoint and its reciprocal inverse are

\[
 E_j:(Z_j,a_j,c_j)\mapsto
       (AZ_j,(AD_j)^{-1}a_j,D_jc_j).
 \tag{RFB.10}
\]

At \(j=1\), this is exactly MN.12. At \(j=-1\) the native
characters are exactly NS.16. The two normal factors in RFB.10
are, respectively,

\[
\begin{array}{c|cc}
j&a_j&c_j\\\hline
1&\eta^4&g^4\eta^{-4}\\
0&g^4\eta^4&\eta^{-4}\\
-1&g^8\eta^4&g^{-4}\eta^{-4}.
\end{array}
 \tag{RFB.11}
\]

In particular the full \(\sigma^{-48}\) factor remains in the
positive normal. This table is a direct substitution in the named
forward endpoint; it does not use the source-word return.

\par

\medskip\noindent\textbf{RFB.2. The two extending Cox changes and the unchanged exterior}

\par\nobreak

In the \(C_1\) row of RFB.7 the two native triples are

\[
 (z,n_0,n_1)=(c,a,b),\qquad
 (z',n'_0,n'_1)=(\eta_jv,gw,(g\eta_j)^{-1}u).
 \tag{RFB.12}
\]

They satisfy \(z'=z^{-1},n'_i=zn_i\) literally. Take Cox
coordinates \((a_0,a_1,c_0,c_1)\) with auxiliary weights
\((1,1,-1,-1)\), \(z=a_1/a_0\), \(n_i=a_0c_i\), and
\(q=a_0a_1c_0c_1\). On the opposite ruling\textquotesingle s \(c_1\ne0\)
chart the regular coordinates are

\[
 (a',b',c')=(c_0/c_1,a_0c_1,a_1c_1).
 \tag{RFB.13}
\]

On the common torus they are \((a/b,b,bc)\); their torus inverse
is \((a,b,c)=(a'b',b',c'/b')\). Thus the actual target characters
are \(Z_{j-1}=Z_j,V_{j-1}=Z_jV_j\), not a scalar-only substitution.
The opposite \(c_0\ne0\) chart and its reciprocal overlap give the
other target vertex. Its complete generators are \(h_{j-1}=h_jB^{-1}\),
and the full coefficients obey

\[
 \kappa_{j-1}=\kappa_jg^{-1},\quad
 \alpha_{j-1}=\alpha_j\eta_j^{-1}\kappa_jg^{-1},\quad
 \eta_{j-1}=\eta_jg .
 \tag{RFB.14}
\]

Here is also the common proper graph, with both actual morphisms.
On the total space of \(\mathcal O(-1,-1)\) use homogeneous vectors
\(A=(A_0,A_1),C=(C_0,C_1)\) and a coefficient \(\tau\), with

\[
 (A,C,\tau)\sim(lA,mC,(lm)^{-1}\tau).
\]

Its two projections are

\[
 p_+(A,C,\tau)=[A_0,A_1,\tau C_0,\tau C_1],\qquad
 p_-(A,C,\tau)=[\tau A_0,\tau A_1,C_0,C_1].
 \tag{RFB.15}
\]

Both pull the original height back to
\(\tau^2A_0A_1C_0C_1\). On the complement of the respective
zero section, the inverse to \(p_+\) takes
\([C]=[c_0:c_1]\) and the unique line coefficient \(\tau\)
with \(c=\tau C\); the inverse to \(p_-\) uses
\([A]=[a_0:a_1]\) and \(a=\tau A\). Their expressions in any
nonzero homogeneous coordinate are holomorphic reciprocal maps.
At \(\tau=0\) the displayed projections retain the two separate
projective rulings. No inverse across that exceptional divisor is
claimed.

Every bounded torus arrow \((Z_j,V_j)\mapsto(AZ_j,DV_j)\)
has the full \(C_1\) Cox lift

\[
 (a_0,a_1,c_0,c_1)\mapsto
 (a_0,Da_1,(AD)^{-1}c_0,Ac_1).
 \tag{RFB.16}
\]

It preserves \(q\), so its inverse divides by the same full
functions at that same \(q\). It commutes with the auxiliary and
physical circles. In the target chart RFB.13 its factors are
\(((A^2D)^{-1},A,AD)\), exactly the \(D\mapsto AD\) reset
of RFB.14. The same diagonal map on \(A,C\), with \(\tau\)
unchanged, intertwines both projections RFB.15. Thus every bounded
seam and every unit arrow has an actual lift on the common graph.
The period translates are defined by RFB.7 and RFB.14 on their
labelled affine charts; the written cancellations of their \(q\)
powers take place before the unit Cox maps are used.

Transport the actual Hermitian norms by RFB.16. The two common-link
Hopf inverses are \(v=u^{-1}q_H\), \(u=q_Hv^{-1}\); their auxiliary
circle ambiguity is exactly the written Cox action. The exterior
arrow and its inverse are

\[
 T(y,q_H)=(q_H^{-1}yq_H,q_H),\qquad
 T^{-1}(y,q_H)=(q_Hyq_H^{-1},q_H).
 \tag{RFB.17}
\]

These are the full IF3.17--19 common-link maps, with all four Cox
weights and paired half-phase endpoints retained. The explicit
normal-rescaling flow of IF3.2, using its \(q\)-height inverse,
carries the full nonlinear germ on a uniform smaller tube. Apply
RFB.12--17 first for \(j=1\) and then \(j=0\). On the original
outside source and finite-cap charts use the identity, with the
inverse character arrow RFB.4 on their torus overlap. This gives
the actual two fan families and their same original endpoint.

The analytic reference\textquotesingle s first reverse change can also be seen in
RB\textquotesingle s factor graph: interchange the final \(x_i,\bar y_i\) in
\((\sigma^{12},x_i,\bar y_i)\). On the nonzero scalar circle this
is exactly RFB.15. The subsequent change is the next periodic
\(C_1\) cell, with the complete \(B\) label and RFB.14; its scalar
factors are retained. The whole-pair calculation below is on this
completed nonzero-circle family and does not require extending a
dense \(Z^2\) multiplier through a scalar-zero vertex.

\par

\medskip\noindent\textbf{RFB.3. The exact cut, native sphere and finite-cap squares}

\par\nobreak

Choose an endpoint-flat cut function \(\chi\), and actual cut
logarithms of \(g\) and \(\eta\). Set

\[
 \mathcal A_\theta=\exp(4\chi\operatorname{Log}g),\quad
 \mathcal D_\theta=\exp(4\chi\operatorname{Log}\eta),\quad
 \mathcal D_{j,\theta}=\mathcal D_\theta\mathcal A_\theta^{-j}.
 \tag{RFB.18}
\]

The two native chart maps are

\[
 C^0_{j,\theta}(Z,a,c)=
 (\mathcal AZ,(\mathcal A\mathcal D_j)^{-1}a,\mathcal D_jc),
\]
\[
 C^\infty_{j,\theta}(z,a_\infty,c_\infty)=
 (\mathcal A^{-1}z,\mathcal D_j^{-1}a_\infty,
                         \mathcal A\mathcal D_jc_\infty).
 \tag{RFB.19}
\]

Both preserve the original \(q\); their inverse uses its unchanged
value before evaluating the reciprocal coefficients. They obey
\(C_{j,1}E_j=C_{j,0}\) on complete endpoint collars. Extend to the
next cut by \(C_{j,\theta+1}E_j=C_{j,\theta}\), retaining every
derivative. The angular part of \(\operatorname{Log}\eta\) remains
\(24\pi i\theta\), including its \(48\)-fold endpoint normal phase.

Write \(S_j(Z,a,c)=(Z,Z^ja,Z^{-j}c)\) on its stated character
localization, with reciprocal \(S_j^{-1}(Z,a,c)=(Z,Z^{-j}a,Z^jc)\).
There is the exact full cut square

\[
 C_{j,\theta}S_j=S_jC_{0,\theta},\qquad
 S_j^{-1}C_{j,\theta}=C_{0,\theta}S_j^{-1}.
 \tag{RFB.19a}
\]

For example the left side\textquotesingle s negative normal coefficient is
\(\mathcal A^{j-1}\mathcal D^{-1}Z^ja\), and its positive
coefficient is \(\mathcal D\mathcal A^{-j}Z^{-j}c\);
the right side has exactly these two coefficients. The identity
\(\mathcal D_j=\mathcal D\mathcal A^{-j}\) retains the complete
\(\eta\) phase. RFB.15--16 and the directed Cartier modules below
are its extending chart arrows. Its first-normal derivative and
the two cap charts therefore carry this same square, with its
reciprocal, through RFB.21--23. This is the shared inner gauge used
in RFB.42; no endpoint scalar phase is suppressed.

The vertical velocity is
\((\partial_\theta\log\mathcal A)Z\partial_Z
 -\partial_\theta\log(\mathcal A\mathcal D_j)a\partial_a
 +(\partial_\theta\log\mathcal D_j)c\partial_c\).
It annihilates \(q\). Add the prescribed horizontal parameter lift,
and cut off only this vertical difference on a larger selected tube.
Its flow and backward flow give the full relative collar map;
the actual cubic cap and its annular reference are fixed there.
In particular their square remains

\[
                 \Phi_{\rho+4}q^4=h_\infty^4\Phi_\rho,
 \qquad h_\infty=1-t_c(s_*)z^3,
 \tag{RFB.20}
\]

with the reciprocal square in the opposite direction. The actual
root action \(A_0^3=B_0\), the three chronological line factors,
the finite character and every sector integer are exactly RB.23--24.

For clarity, the selected Hopf cap is also fixed by a literal map.
Use NC\textquotesingle s matrices

\[
 T_Z=\begin{pmatrix}1&-\bar Z\\Z&1\end{pmatrix},\quad
 T'_z=\begin{pmatrix}z&-1\\1&\bar z\end{pmatrix},\quad
 T'_{1/Z}\operatorname{diag}(Z,\bar Z)=T_Z.
\]

The complete first-normal derivative of RFB.19 is

\[
 G_{j,\theta}(Z)=T_{\mathcal AZ}
 \operatorname{diag}((\mathcal A\mathcal D_j)^{-1},
                                  \overline{\mathcal D_j})T_Z^{-1}.
 \tag{RFB.21}
\]

Here the functions are their actual \(q=0\) restrictions; their
first normal derivatives vanish since \(dq\) vanishes on the
fixed sphere. The other chart equality above proves smoothness
at infinity, with the full inverse matrix. On the trace-one
positive Hermitian three-ball use

\[
 \Psi_{\mathcal A}(Y)=
 \frac{\operatorname{diag}(\mathcal A,1)Y
       \operatorname{diag}(\mathcal A,1)^*}
      {\operatorname{tr}(\operatorname{diag}(\mathcal A,1)Y
       \operatorname{diag}(\mathcal A,1)^*)}.
 \tag{RFB.22}
\]

Its inverse is \(\Psi_{\mathcal A^{-1}}\), and on the boundary
it is exactly \(Z\mapsto\mathcal AZ\). In RFB.21 replace the two
cut logarithms by their product with a radial function zero near
the ball centre and one near its boundary. The resulting
\(\widehat G_{j,\theta}\) is independent of the sphere direction
near the centre and has the exact other-chart expression. Hence

\[
 \widehat C_{j,\theta}(Y,q_H)=
 \left(\Psi_{\mathcal A}Y,
 \frac{\widehat G_{j,\theta}(Y)q_H}
      {\|\widehat G_{j,\theta}(Y)q_H\|}\right)
 \tag{RFB.23}
\]

is the actual invertible cap map. Its inverse first applies
\(\Psi_{\mathcal A^{-1}}\), then the inverse matrix at that recovered
point and the positive normalization. The cap forward seam is
\(\widehat E_j=\widehat C_{j,1}^{-1}\widehat C_{j,0}\).
The resolved boundary map is \(B_{j,\theta}=C_{j,\theta}^{-1}N_{j,\theta}\);
the exact square is \(\widehat C_{j,\theta}B_{j,\theta}=N_{j,\theta}\),
with reciprocal inverse. At \(j=1\) these are precisely MN.12--14
and NC.10--20. The \(j=0,-1\) caps are specified by this same
construction with their actual RFB.11 coefficients.

\par

\medskip\noindent\textbf{RFB.4. An actual completed orbit-piece and relative-line morphism}

\par\nobreak

We give the full named application of K72 SI.6--27, rather than
identifying the pieces from their winding matrix. For \(j\in\mathbb Z\)
use \(Q_j=\kappa qg^j\) and the actual contracting coefficient

\[
 \gamma_j=
 \frac{\log|Q_j|\log|V_j|-\log|\alpha_j|\log|Z_j|}
 {\log|Q_j|\log|q\eta_j|-\log|\alpha_j|\log|g|}.
 \tag{RFB.24}
\]

The denominator has leading term \((\log|q|)^2\) and is nonzero
on a uniform small disc. The displayed full positive norms are
independent of \(\theta\). Direct substitution gives
\(\gamma_j(h_jp)=\gamma_j(p)\) and
\(\gamma_j(Bp)=\gamma_j(p)+1\). The F9 quotient coordinate has

\[
 2b=r-|y|^2/r,\quad2b'=|x|^2r-r^{-1},\quad
 db'/db=(|x|^2+r^{-2})/(1+|y|^2r^{-2})>0.
 \tag{RFB.25}
\]

It is matched to the contracting coordinate by the proper decreasing
map \(1/2-\epsilon(r-1)/(r+1)\), with \(0<\epsilon<1/4\).
Both branch derivatives are \(-4\epsilon\). The second fixed-sphere
endpoint is the actual \(h_j^{-1}\) translate of the second
conifold endpoint and therefore has the same lifted \(\gamma_j\).
The endpoint heights stay distinct in that short interval; no
integer is inserted. Convex interpolation of the proper decreasing
maps, fixed on the outside collar, gives the full matching and its
inverse. Untwisting by \(\exp(\gamma\log g)\) and the complete
positive line norm gives the actual orbit piece
\(W_j=E_j\times S^1_\gamma\), with the marked outer coefficients
unchanged. The line phases remain in its principal circle bundle.

The explicit source comparison to \(j=0\) uses

\[
 Q_{a}=\kappa q\exp(aj\log g),\qquad
 X=\partial_a+j\log g(q)\,y\partial_y
          \quad\hbox{on }xy=Q_a,\quad0\le a\le1.
 \tag{RFB.26}
\]

It fixes the original \(q\). The inverse local base map recovers
\(q\) by the inverse of \(q\mapsto Q_a(q)\); its derivative at zero
is the nonzero \(\kappa(0)g(0)^{aj}\), so one smaller disc works
uniformly. On the outside collar prescribe instead the real marked
period map \(u+\tau_av\mapsto u+\tau_0v\). Both vector fields
project to \(\partial_a\); their difference is vertical. A finite
partition, fixed on the protected collars, glues them. The compact
boundary-tangent field has a complete flow, with the backward flow
as inverse. It is independent of the scalar phase.

The fixed normalization spheres have their lifted graphs in the
same short \(\gamma\) interval. Interpolate them with the distinct
endpoint heights fixed. The source normalization is injective away
from its two node preimages, so this interpolation is an embedding.
Extend its velocity over a tube and cut it off. This gives the actual
relative ambient isotopy, not a replacement of the nodal source by
real logarithmic coordinates.

The native normal arrow is

\[
 (Z,a_0,c_0)\mapsto(Z,Z^j a_0,Z^{-j}c_0),\qquad q=Za_0c_0.
 \tag{RFB.27}
\]

It is interpreted on the two \textbf{specified} fan stars. On their
normalization sphere its full positive-spinor map is the Cartier
module arrow

\[
 E_j^{\rm spin}=E_0^{\rm spin}\otimes
       \mathcal O(j(D_0-D_\infty))\longrightarrow E_0^{\rm spin},
                 \qquad w\longmapsto Z^j w.
 \tag{RFB.28}
\]

The defining Cartier generators cancel the written zeros and poles
at \(D_0,D_\infty\). RFB.27 is the character arrow from the \(0\)
star to the \(j\) star; its positive spinor coefficient matrix is
\(\operatorname{diag}(\overline Z^{\,j},Z^{-j})\). RFB.28 is its
directed inverse from the \(j\) spinor bundle to the \(0\) bundle,
with matrix
\(\operatorname{diag}(\overline Z^{-j},Z^j)\). After the transported
positive norms this inverse has the common \(Z^j\) phase. The
remaining positive factors are the explicit diagonal norm
rescalings, with their reciprocal inverse. This proves the
directed spinor map, including the physical circle. Its adjoint
carries the normal three-plane. The
source derivative and the graph velocity give their actual positive
triangular maps. Extend the sphere velocity over the ball by a
collar cutoff and use RFB.23 on its Hopf normal sphere. These maps
give the completed orbit-piece diffeomorphism \(I_j:C_j\to C_0\),
where

\[
 C_j=(W_j\setminus\operatorname{int}(F_j\times D^3))
                      \cup(D^3\times S^2).
 \tag{RFB.29}
\]

All inverses are the displayed reciprocal bundle maps and backward
flows. In the original cut coordinates the cap map is
\(\widehat C_{0,\theta}^{-1}I_j\widehat C_{j,\theta}\).
Its endpoint identity follows directly from
\(\widehat C_{j,1}\widehat E_j=\widehat C_{j,0}\) and the
periodic closed-coordinate \(I_j\). On the principal circle line
the additional tensor is retained next; it is trivialized on this
cap by making the graph map \(b\) equal to zero there.

Use the actual coefficient decomposition

\[
 (\log Z,\log V)/(2\pi i)=
 u(1,0)+v(\tau,\mu)+d(0,1)+\gamma(G,\zeta),
 \tag{RFB.30}
\]

where \(\tau=\log(\kappa q)/(2\pi i)\),
\(\mu=\log\alpha/(2\pi i)\), \(G=\log g/(2\pi i)\),
and \(\zeta=\log(q\eta)/(2\pi i)\). These are the full
functions in RFB.2. The exact inverse period matrix determines the
real coefficients and retains their complete derivatives. The
integer deck maps are

\[
 q\hbox{-turn}:(u,d)\mapsto(u-v,d-\gamma),\qquad
 \theta\hbox{-turn}:(u,d)\mapsto(u,d-12\gamma).
 \tag{RFB.31}
\]

RFB.4 gives exactly

\[
 u_j=u,\quad v_j=v,\quad\gamma_j=\gamma-jv,\quad d_j=d-ju.
 \tag{RFB.32}
\]

Write \(\alpha_q=\arg q/(2\pi)\). For the meridian-one line
\(U=M^*\), the actual coefficient connections are

\[
 A_0=dd+(\alpha_q+12\theta)d\gamma,\qquad
 A_j=dd_j+(\alpha_q+12\theta)d\gamma_j+12j\theta\,dv.
 \tag{RFB.33}
\]

They are invariant under both full deck maps, and their prescribed
outer difference is exactly

\[
                  \beta_j=-j(du+\alpha_qdv).
 \tag{RFB.34}
\]

Let \(b:C_0\to\mathbb R/\mathbb Z\) extend the source graph
coordinate \(v\). It extends across the node\textquotesingle s vanishing disc
because its attaching circle is the \(u\)-circle. Its restriction
to \(F_0\) has a real lift; a supported real homotopy makes it zero
on a tube containing the surgery and hence on the cap. Its exact
theta-cut gauge is \(\exp(2\pi i12j\theta h)\), where \(h\)
is this real change of graph map. Define the actual line
\(\mathcal T\) by the \textbf{endpoint-to-initial} coefficient clutch
\(w\mapsto e^{2\pi ib}w\). On its theta cut the coefficient
connection is \(A=dd+\theta\,db\); at the endpoint this is the
pullback of \(dd\) by that exact clutch. The reciprocal
initial-to-endpoint arrow multiplies by \(e^{-2\pi ib}\).
Then \(c_1(\mathcal T)=\Xi=d\theta\wedge db\), \(\Xi^2=0\).

The Thom and Mayer-\/-Vietoris maps of this displayed surgery give
\(H_1(C_0)=\mathbb Z\langle b,\gamma\rangle\) and
\(H_2(C_0)=\mathbb Z\langle\text{Hopf meridian},b\times\gamma\rangle\).
The primitive meridian has \(U\)-value one; the real-period torus
has value zero. The theta clutch fixes these generators and adds
the \(b\times\theta\) and \(\gamma\times\theta\) evaluations.
RFB.32 gives difference \(12j\) on the first and zero on the other
three. Thus the actual line morphism is

\[
          J_j:U_j\longrightarrow I_j^*U_0\otimes\mathcal T^{12j}.
 \tag{RFB.35}
\]

Its prescribed boundary arrow is also constructed, not inferred
from those evaluations. The boundary first cohomology has basis
\((\alpha_q,\theta,v,\gamma)\); the extending image is
\((\theta,b,\gamma)\). The only remaining obstruction is therefore
the \(q\)-angle. On one parameter arc choose a generic \(Z_*\)
with \(|Z_*|=1\), \(V=1\), avoiding the finite marked sections.
The literal free section disc is

\[
 (q/Z_*,Z_*,1)\mapsto(qZ_*^{j-1},Z_*,Z_*^{-j}).
 \tag{RFB.36}
\]

It misses the surgery and has \(v=\gamma=0\); both line factors
and their reciprocals extend over \(q=0\). The source flow can fix
this disc on that parameter arc. Thus the sole boundary obstruction
is zero, as an integer. On a good cover take logarithms of the
actual ratio-line transitions, including
\(j\log\Gamma-j\log\kappa-j^2\log g\) and \(-j\log g\),
and subtract the integral one-cochain supplied by the four preceding
evaluations. In frame convention \(e_j=e^{2\pi ir_{ij}}e_i\),
set \(a_i=\sum_k\rho_kr_{ki}\); then \(a_j-a_i=r_{ij}\).
The reciprocal exponentials give RFB.35 and its inverse. Remove
the remaining extending boundary windings by the actual functions
\(\theta,b,\gamma\), then extend its real logarithm with a collar
cutoff. This makes the boundary arrow literally RFB.32--34. The
cap gauge extends because \(b=0\) there. Every original annular
integer, finite character, and full nonzero point/O module is
carried throughout this construction.

\par

\medskip\noindent\textbf{RFB.5. The whole Pontryagin and Euler integrals}

\par\nobreak

The full outer stable coframe is
\((du+(\alpha_q+\phi_\kappa)dv,dv,d\gamma,d\alpha_q,d\theta,dr_q)\).
The actual source map just constructed has matrix

\[
 G_j=I+j\phi_gE_{12}-jE_{32},\qquad
 G_j^{-1}dG_j=j\,d\phi_gE_{12},\qquad
 \phi_g=\operatorname{Im}\log g/(2\pi).
 \tag{RFB.37}
\]

Here \(\phi_g\) is independent of scalar phase. The Maurer form
squares to zero. Interpolating the two actual boundary connections
by a cutoff multiple of this form gives \(\operatorname{tr}F^2=0\)
on the whole collar. No line-phase term has been added to a tangent
frame; its phase remains in RFB.33.

Choose nested tubes containing all surgery and connection changes.
Away from them take the tangent connection pulled from the actual
nodal source, with the trivial \(\gamma\) direction and the marked
flat outside collar. A regular level of the original source graph
is \(N^3\times S^1_\gamma\), on which that source \(p_1\) form
vanishes. The supported homotopy making \(b\) constant near the
larger tube gives a relative regular-level cobordism; choose a
regular value outside that constant. Its new level misses the
surgery. Stokes and the zero full outer transgression preserve the
pairing with the closed \(p_1\) form. This proves

\[
                    \langle p_1,\Xi\rangle=0.
 \tag{RFB.38}
\]

The exact boundary primitive RFB.34 contributes zero to this
pairing because RFB.37 vanishes throughout its collar. Therefore
\(\ell_{M,j}=\ell_{M,0}\) for these actual completed pieces,
with the original exterior unchanged.

For the Euler cube choose the affine interior curvature supplied
by RFB.35: \(F_j=F_0+12j\Xi\). Since \(\Xi^2=0\), its entire
interior difference is linear in \(j\), with one fixed coefficient.
Use the piece coefficient \(\gamma=\gamma_j\), so the retained
exterior coefficients are \(\gamma_{\rm ext}=\gamma+jv\) and
\(d_{\rm ext}=d_j+ju\). Its exact curvature is

\[
 F_{{\rm ext},j}=(d\alpha_q+12d\theta)\wedge(d\gamma+j\,dv),
 \qquad d\beta_j=-j\,d\alpha_q\wedge dv.
 \tag{RFB.39}
\]

Orient the outward boundary by
\(\Omega=d\alpha_q\wedge du\wedge dv\wedge d\gamma\wedge d\theta\).
The marked period cell has integral one and the exact deck maps
preserve its volume. Direct exterior algebra gives

\[
 3\int_0^1\beta_j\wedge(F_{{\rm ext},j}+a\,d\beta_j)^2\,da
                         =36j^2\Omega.
 \tag{RFB.40}
\]

The actual outward connection runs from piece to exterior, so its
cutoff runs from one to zero. Its contribution is \(-36j^2\).
The physical line is \(M=U^*\), and the K68 original-Jin parameter
arrow gives \(k_M=\int F^3\) in this positive coefficient
orientation. Both signs are the same actual Hopf/physical arrow
used in NC and MN, not a choice of positive-hyperplane generator.
Consequently, for one fixed interior coefficient \(L\),

\[
       k_{M,j}=k_{M,0}+Lj-36j^2,
       \qquad\ell_{M,j}=\ell_{M,0}.
 \tag{RFB.41}
\]

The coefficient is fixed by an actual integer comparison. In the
\(j=0\) family\textquotesingle s \(C_2\) Cox bordism the first outer fixed-branch
map is the identity, and the second is

\[
        \text{base }\kappa^{-1},\qquad
        \text{normals }\operatorname{diag}(\alpha^{-1},\kappa\alpha).
 \tag{RFB.42}
\]

These are exactly the full RFB.2 functions and are independent of
the scalar phase. The two inner charged boundaries have the same
physical-circle gauge from RFB.18--23; it extends through these
constant outer maps. Thus the line and normal-three-plane
transitions on the additional sphere boundary pull back from the
three-dimensional fixed factor. Their \(b^2,p_1(\nu)\), and tangent
\(p_1\) vanish there. The actual sphere-bundle formulas
\(4\int x^3=3b^2+p_1(\nu)\) and
\(\int p_1x=p_1(T)+p_1(\nu)\) give zero separately. This is the
full K62 SCN.9 integer map with the present \(e=4\) seam; the finite
cap, scalar phase and integral cochain have been fixed by RFB.20.
The target fan is \(T_1\): its triangle identification is the
literal \(h^{-1},B\) translation with all RFB.3 units. Its native
cap is exactly MN.12 by RFB.10 and RFB.23. Hence
\(k_{M,1}=k_{M,0}\), so \(L=36\).

Substitution gives the complete named relative formula

\[
 (k_{M,j}-k_{M,0},\ell_{M,j}-\ell_{M,0})
                         =(-36j(j-1),0).
 \tag{RFB.43}
\]

In the two requested directions the differences are

\[
 T_1\longrightarrow T_0:\ (0,0),\qquad
 T_0\longrightarrow T_{-1}:\ (-72,0).
 \tag{RFB.44}
\]

This is RFB.1. The induced whole scalar-reference index change is
\((4(-72)-0)/24=-12\). Its integer numerator was obtained from
the actual interior line, outer primitive and both cap maps. No
individual central class or free coordinate was divided to obtain it.

\par

\medskip\noindent\textbf{RFB.6. Exact scope and reproducibility}

\par\nobreak

The completed scalar boundary family with fan \(T_{-1}\) is the
genuine same-fan reference for the collision NS boundary. RFB.1
converts the \textbf{whole} RB \(T_1\) characteristic pair to that
reference. The boundary source/point/root square is still SI2,
with original root \(r_0\sigma^4\), finite character and annular
integers. The supplied correspondence, quotient diffeomorphism,
relative line morphism and cap maps have explicit inverses on
their stated domains. The reference\textquotesingle s original analytic filling
and its individual fixed-component data remain RB and MN.

The companion \path{code/part3/e6_selected_i3_m2_reference_fan_bridge.py}
checks the full character and generator reset, both Cox charts,
common-graph projections, three endpoint/cap overlaps, exact
real-period inverse and integer deck square, nilpotent tangent
collar, whole exterior-algebra cube, and the two signed pair
differences. Its check count and final source hash are recorded in
the generated JSON. These identities supplement the actual map
construction above; they do not replace its compact flow,
Cartier, cap or relative line arguments.


% End expanded verified cp36 input


% Expanded verified cp36 input: notes/part3_e6_selected_i3_m2_return_support_18.tex
\par\bigskip\noindent\textbf{P3K.111 supporting proof 18: Independent actual doubled scalar native arrows}\par\nobreak
\par

\medskip\noindent\textbf{Selected-I3 \(M_{2}\): the actual scalar-reference native products and signed flops}

\par\nobreak

Working proof, 18 November 2026. Fix the original source \(s=s_*\ne0\),
its point line \(M_2\), chronological translation \(-3G\), cubic
good cap with norm one and prescribed finite character. The scalar
root is \(r_0\sigma^4\), and \(\lambda=\lambda_0\sigma^{12}\).
All functions, inverse source charts, point/O modules, root actions,
sector integers and forward/inverse closing arrows are exactly
RFB.1--3. This is the independent scalar word \((0,4)\).

For the actual \(T_j\) family, the two single native free pairs are

\[
 (m_{C_1},n_{C_1})=(12(1-j),12(j-1)),\qquad
 (m_{C_2},n_{C_2})=(12j,-12j).
 \tag{RFS.1}
\]

These coordinates belong to K23\textquotesingle s actual relative exterior group.
The signed Cox maps give the central changes

\[
 R_{j-1}-R_j=12(j-1),\qquad R_{j+1}-R_j=-12j
                       \quad\hbox{in }\Theta_7.
 \tag{RFS.2}
\]

The first and second reverse changes \(T_1\to T_0\to T_{-1}\)
therefore contribute \(0,-12\). The whole physical integer-pair
comparison \((-72,0)\) is proved independently in RFB.4--5; no
integer Chern number is inferred from RFS.2.

\par

\medskip\noindent\textbf{RFS.1. The full single products and their literal doubles}

\par\nobreak

Keep the complete functions

\[
 \kappa=P^2/c_q,\quad\alpha=P^2,\quad g=P^{-3},\quad
 \eta=\sigma^{12}\lambda_0\eta_2^6t^{-3},\qquad
 \eta_2=a_T\Theta'_Q(1)\Theta_Q(P^2)/\Theta_Q(P)^2,
\]
\[
 \kappa_j=\kappa g^j,\quad
 \alpha_j=\alpha\eta^j\kappa^{-j}g^{-j^2},\quad
 \eta_j=\eta g^{-j}.
 \tag{RFS.3}
\]

The original selected height is \(q=t-t_c(s_*)\). The two vertex
triples satisfy \(q=abc=uvw\), with the three complete rows RFB.7.
The \(C_1\) native charts are
\((c,a,b)\) and \((\eta_jv,gw,(g\eta_j)^{-1}u)\);
the \(C_2\) native charts are \((a,b,c)\) and \((u,v,w)\).
Direct substitution gives \(z'=z^{-1},n'_i=zn_i\) in each case.

The forward endpoint is the actual \(q^4\) arrow reindexed by
\(B^{-4}\), so its bounded point map is
\((Z_j,V_j)\mapsto(g^{-4}Z_j,\eta_j^{-4}V_j)\).
Use exactly the RFB.18--23 cut maps in these two native charts.
They preserve \(q\), have reciprocal inverses evaluated at that
same \(q\), and close on full collars with all jets. The original
cubic cap retains \(\Phi_{\rho+4}q^4=h_\infty^4\Phi_\rho\), its
three chronological source/line factors, character and sector integers.

This specifies a single native product before taking a cover. Its
double is its literal pullback. On the second lap use
\(C_\theta E^{-1}\), with the original \(E\), and at the final cut
\(E^2\). Thus the capped endpoint is the actual square \(f^2\).
The root/ordinary square is pulled back to
\(\Phi_{\rho+8}q^8=h_\infty^8\Phi_\rho\), with the same two
chronological laps of the original cochain. No independent double
seam is selected.

Write the doubled scalar circle as \(\sigma=\rho e^{2i\vartheta}\),
\(0\le\vartheta\le2\pi\). The winding matrix of the full periods
in their displayed \(h_j,B\) order is

\[
 \mathsf B_j^{(2)}=2\begin{pmatrix}0&0\\12j&12\end{pmatrix}.
 \tag{RFS.4}
\]

All four remaining factors are full \(q\)-dependent units with
zero parameter winding. In particular the native closing is still
the one just specified; RFS.4 does not define a new closing.

\par

\medskip\noindent\textbf{RFS.2. The full ordinary product and its inverse}

\par\nobreak

At the second vertex use the full arrow

\[
          (u,v,w)\longmapsto(u,\eta_j(q,\vartheta)v,
                                      \kappa_j(q)w).
 \tag{RFS.5}
\]

Its target height is \(H_\vartheta(q)=q\eta_j(q,\vartheta)\kappa_j(q)\).
Recover the original \(q\) by \(H_\vartheta^{-1}\) and then divide
by these complete functions. The derivative at zero is nonzero;
compactness of the scalar circle gives one uniform injective disc.
For the normal orders \((a,c),(a,b),(b,c)\), the three whole
matrices are

\[
 D_0=\operatorname{diag}(\eta_j/\alpha_j,\kappa_j\alpha_j),\quad
 D_1=\operatorname{diag}(\kappa_j/g,g\eta_j),\quad
 D_2=\operatorname{diag}(\eta_j,\kappa_j).
 \tag{RFS.6}
\]

On the literal double their determinant winding is \(2t=24\),
and their first-entry windings are

\[
              t=12,\qquad l=(24(1-j),0,24).
 \tag{RFS.7}
\]

Divide each entry by its full reference value at the same \(q\)
and by the displayed angular factor. The quotient has a unique
based periodic logarithm on this protected disc. Retain those
entire logarithms and all their \(q\) derivatives.

For any integer \(d\), use the endpoint-flat based path

\[
 H_d(x,\vartheta)=
 (\sin\varphi(x)e^{id\vartheta}+j_{\rm quat}\cos\varphi(x))
 (\sin\varphi(x)+j_{\rm quat}\cos\varphi(x))^{-1},
 \quad\varphi(0)=0,\quad\varphi(1)=\pi/2.
 \tag{RFS.8}
\]

It equals one at \(x=0\) and both parameter ends, and equals
\(e^{id\vartheta}\) at \(x=1\). On
\(Q=n_0+j_{\rm quat}\bar n_1\) use
\(Q\mapsto H_tQH_{l_i-t}\), then compose the two full based-log
exponentials with the same radial cutoff. The radial endpoints
are exactly identity and \(D_i(q,0)^{-1}D_i(q,\vartheta)\).

For completeness, its interior inverse retains the original height.
On a bounded curve tube this map is \(F(n)=A n+G(n)\), where
\(A\) is its invertible value at \(q=0\). Since \(q=zn_0n_1\),
the complete unit dependence gives \(G=O(|n|^3)\) uniformly with
its first derivative. Choose the tube with
\(\|A^{-1}DG\|<1/2\). Iterating
\(n\mapsto A^{-1}n'-A^{-1}G(n)\) gives the unique smooth inverse,
and the source height is the actual \(zn_0n_1\) at that recovered
point. At the vertex use the \(H_\vartheta\) inverse from RFS.5.
The radial endpoint equalities glue the two constructions.

The normal rescaling \((z,n)\mapsto(z,\tau n)\), followed by the
cutoff flow of its conjugated family velocity, restores the full
nonlinear germs, with the compact inverse bound just proved. It
fixes the protected vertex collars, and the three open curve tubes
are disjoint. Carry the rounded boundary by this same isotopy.
The flow is in the closed doubled coordinates; it therefore
preserves the literal pulled-back seam and all source/cap arrows.
This constructs the actual ordinary double product, including its
inverse, before its free integers are evaluated.

\par

\medskip\noindent\textbf{RFS.3. The actual action and stable-derivative prisms}

\par\nobreak

The physical tangent weights are \(w=(0,1,-1)\) and the ordered
normal differences are \(d=(-2,-1,-1)\). Under RFS.8 the normal
action is

\[
 Q\longmapsto
 (H_t e^{-iw_i\phi/2}H_t^{-1})Q
 (H_{l_i-t}^{-1}e^{id_i\phi/2}H_{l_i-t}).
 \tag{RFS.9}
\]

The full unit maps and normal flow give its based homotopy to the
original action. The actual trace at \(q_H=1\), with
\(U=\cos x+j_{\rm quat}\sin x\), \(E_i=e^{iw_i\phi/2}\), is

\[
 E_i^{-1}UE_iH_tE_i^{-1}H_t^{-1}U^{-1}
                 H_{l_i-t}^{-1}e^{id_i\phi/2}H_{l_i-t}.
 \tag{RFS.10}
\]

Differentiating the displayed \(H_d\) gives the oriented degrees
\(-d,d\) of \(H_diH_d^{-1}\) and \(H_d^{-1}iH_d\).
The positive quaternion volume pairs angle \(c\phi/2\) and axis
degree \(e\) by \(ce\). The Maurer product boundary terms cancel
on the periodic parameter charts and have central constant radial
ends. Thus the normal-meridian and stable-derivative degrees are

\[
 I_i=tw_i+d_i(l_i-t),\qquad W_i=tw_i-d_i(l_i-t),
\]
\[
 I=(24(2j-1),24,-24),\qquad W=(24(1-2j),0,0).
 \tag{RFS.11}
\]

The stable lift is the lifted polar map to Spin: left quaternion
multiplication contributes its third degree and right multiplication
its negative. The base rotation depends only on \(\phi\), and
its endpoint Spin sign cancels the normal sign. Close the traces
by the actual vertex normal sections and link paths, capping their
parameter-independent end orbits by \(D^2_\vartheta\times S^1_\phi\).
The vertex maps factor through dimension two. The full
coordinate-to-ambient frame arrow extends over the annulus times
\(D^4\), so contributes zero on these capped three-cycles. This
gives the actual marked prism with the full RFS.6 frames.

For \(\gamma_i=\partial x_i\), the oriented positive-circle
identity \(\partial\Delta=-\Delta\partial\) yields

\[
 V(w_ie_i-w_0e_0)=I_ie_i-I_0e_0,\qquad
 w_i\delta_i-w_0\delta_0=W_i-W_0.
 \tag{RFS.12}
\]

Hence the actual two selected entries are

\[
 V_{11}=24,\quad\delta_1=24(2j-1),\qquad
 V_{22}=24,\quad\delta_2=24(1-2j).
 \tag{RFS.13}
\]

The ambient-to-local frame is literally \(c=u^{-1}q_H\).
Its derivative contains \(-u^{-1}(du)u^{-1}q_H\), removed by the
invertible triangular homotopy. Its normal-meridian third class
is \(+1\); in the K23 oriented relative disc it contributes
\(-V_{ii}\). Thus, in these same products,

\[
 2n_i=\delta_i-V_{ii},\qquad m_i=-V_{ii}-n_i.
 \tag{RFS.14}
\]

The ordinary double pairs are therefore

\[
 C_1:(-24j,24(j-1)),\qquad
 C_2:(24(j-1),-24j).
 \tag{RFS.15}
\]

No fixed \(C_0\) value has entered these equations.

\par

\medskip\noindent\textbf{RFS.4. The native correction, actual square and full flop}

\par\nobreak

The native quaternion is
\(q_H=U(z)(n_0+j_{\rm quat}\bar n_1)\), with
\(U(z)=(1+j_{\rm quat}z)/(1+|z|^2)^{1/2}\).
The actual native-to-ordinary boundary arrow of RFS.6--8 is

\[
 q_H\longmapsto(UH_tU^{-1})q_HH_{l_i-t}.
 \tag{RFS.16}
\]

For \(C_1\), its native second normals are
\((gw,(g\eta_j)^{-1}u)\), and the ordinary normals are
\((\kappa_jw,u)\). Their exact ratio is \(D_1\), so RFS.16
has retained every \(g,\kappa_j,\eta_j\) factor. The \(C_2\)
ratio is directly \(D_2\). The full based unit logarithms and
the inverse height maps are those in RFS.2.

The left factor in RFS.16 has degree \(2t=24\) on the
Hopf-oriented sphere times parameter; direct substitution of
RFS.8 in its Maurer boundary integral gives twice the angular
winding. The right factor has only radial and parameter variables
and third degree zero. Extend the inverse factors over the
exterior boundary cylinder by its explicit radial retraction;
the endpoint degrees are \((-24,0)\). They extend over
\(D^3\times D^4\), so no disc-supported central class is added.
Thus the ordinary return is the native return followed by
\(L^{-24}\), and the native doubles are

\[
 C_1:(24(1-j),24(j-1)),\qquad C_2:(24j,-24j).
 \tag{RFS.17}
\]

Use the connection defining the single native product of RFS.1.
Its literal pullback has endpoint \(f^2\), with both laps of the
same original and cap closing. The corrected ordinary connection
has the same native boundary connection. Affine interpolation
relative to that boundary gives an isotopy of relative endpoints.
The free-coordinate homomorphism to \(\mathbb Z^2\) takes
\(f^2\) to twice the coordinates of \(f\); multiplication by two
in this free group is injective. This proves RFS.1. The central
coordinate is not divided.

The full C1 opposite-ruling chart is RFB.13, with Cox lift

\[
 (a_0,a_1,c_0,c_1)\mapsto
 (a_0,Da_1,(AD)^{-1}c_0,Ac_1).
\]

It has the full reciprocal inverse at the same \(q\), retains the
positive moduli by the transported Hermitian metric, and induces
the actual target \(T_{j-1}\) through RFB.14. The C2 lift is
\((a_0,(AD)^{-1}a_1,Ac_0,Dc_1)\), and its actual target is
\(T_{j+1}\). For arbitrary four Cox angular weights the exact
link is IF3.17: it retains the full factor
\(e^{hS\vartheta/2}\), or equivalently \(e^{kS\vartheta/2}\),
and both paired central endpoints. Its exterior is precisely
RFB.17, with the written inverse and Hopf inverses.

In these actual relative products the full K23 identity is

\[
 T(L^mR^nC^a)T^{-1}=L^{-n}R^{m+2n}C^{a+n}.
 \tag{RFS.18}
\]

The fixed anti-quaternion reset restores the first-ruling complex
normal convention, swaps the two free entries and preserves the
central entry. Hence the true central change is \(n\), giving
RFS.2 in the two stated directions. For example the C1 target
free pair after that full reset is exactly the C2 pair at \(j-1\);
the reciprocal C2/C1 check is identical by the displayed equations.
All exterior source/root/cap/sector arrows are those in RFB.3,
extended by identity away from the selected tube. This proves the
signed scalar fan comparison on the actual named reference.

The companion \path{code/part3/e6_selected_i3_m2_scalar_native.py} checks
the entire trace, variation/derivative relations, ordinary/native
corrections, both target resets, and signed recurrence, including
the full coefficient squares on the common Cox chart. It locks the
complete geometric inputs and the RFB source separately.


% End expanded verified cp36 input


% Expanded verified cp36 input: notes/part3_e6_selected_i3_m2_return_support_19.tex
\par\bigskip\noindent\textbf{P3K.111 supporting proof 19: Full original attachment, parameter square and outer transfer}\par\nobreak
\par

\medskip\noindent\textbf{Selected-I3 \(M_{2}\): full compact attachment, same-fan reference and physical source return}

\par\nobreak

Working proof, 18 November 2026. We retain the original source
\[
 y^2=x^3+r_{\rm src}x^2+tx+t^2,\quad G=(0,t),\quad
 s=1/4-r_{\rm src},\quad M_2=\mathcal O([2G]-[O]),
 \tag{ASI2.1}
\]
translation by \(-3G\), the selected height \(q=t-t_c(s)\),
the good cubic cap \((m,n)=(3,1)\), and its prescribed finite
line character and full reference cochain. The positive original
\(s\)-meridian is the one in SI2.1. The actual selected native
star is the one derived from the two protected zero-end facets
in NS.1--3. Its point-coordinate matrix is
\[
 P_N=\begin{pmatrix}1&-1\\-1&2\end{pmatrix};
 \tag{ASI2.2}
\]
in the retained convention \(T_j=T(e_1+je_2,e_2)\), this is
\(T_{-1}\), by the literal triangle and period map RFB.4--7.

The auxiliary filling uses \(\lambda=\lambda_0s^{12}\) and
scalar root \(r_0s^4\), \(r_0^3=\lambda_0\), with \(\lambda_0\)
in the specified common strict contraction range. Choose a
small positive basepoint \(s_*\), and retain
\(r_{\rm base}=r_0s_*^4\), \(\lambda_{\rm base}=r_{\rm base}^3\).
The return computed below keeps this original scalar fixed.
The completed Jin physical pair and index are
\[
 (k_M,\ell_M)_{\rm source}=(71,-4),\qquad
 \frac{4k_M-\ell_M}{24}=12,\qquad
                 [\Phi_{\rm source}]=12\ \hbox{in }\mathbb Z/28.
 \tag{ASI2.3}
\]
We prove the attachment and directed reference comparison before
using the pair in ASI2.3. All auxiliary and physical lines retain
their stated Hopf orientations.

\par

\medskip\noindent\textbf{ASI2.1. The actual completed periodic piece and compactness}

\par\nobreak

The complete functions and point coefficient are
\[
 Q=P^2R,\quad R-P=su_s,\quad t=Pu_t,\quad q=Rc_q,\quad
 \eta_2=a_T\Theta'_Q(1)\frac{\Theta_Q(P^2)}{\Theta_Q(P)^2},
 \quad V_{\rm reg}=\eta_2^{-1}V_{\rm pt}.
 \tag{ASI2.4}
\]
The full Miller telescope in SI2.6--9 is
\[
 B_{2,3}=f_2(p)f_2(p-3G)f_2(p-G)^2f_2(p-2G)^2
 =\eta_2^6\frac{f_{P^2}(Z)}{f_{P^2}(P^{-3}Z)},\quad
 f_2=\frac{y-t-x/2}{x-s}.
 \tag{ASI2.5}
\]
Thus the genuine point periods are
\[
 h:(Z,V_{\rm pt})\mapsto(P^2RZ,P^2V_{\rm pt}),\quad
 B:(Z,V_{\rm pt})\mapsto
       (P^{-3}Z,\lambda q\eta_2^6t^{-3}V_{\rm pt}).
 \tag{ASI2.6}
\]
Their inverses retain the same full functions at the recovered
source point. The entire \(\Theta'_Q(1)\), with central value
\(-1\), remains in \(\eta_2\) and its reciprocal.

Take precisely the executed common I2F fan and its proper zero,
one and infinity lifts I2ZC/I2L. The complete common boundary has
136,442 primitive regular three-cones. The actual parent semigroup
maps, full character inverses, paired period facets and protected
source cuts are the maps in CFL.2--3. Every refined cone is the
specified proper toric subdivision of its raw parent; the original
monomial relations pull back identically. The actual base maps are
\[
       (P,R)=(P,pP),\quad(P,(1+v)P),\quad(Ru,R).
 \tag{ASI2.7}
\]
Their inverses on the overlap are \(p=R/P,v=R/P-1,u=P/R\).
The original \(s,t,q\) are the entire functions in ASI2.4.
We use this executed total-line completion with its physical
circle and derivative maps. The source/Cartier graphs below
specify its actual finite boundary attachment.

Here is an analytic compactness check with those full coefficients.
In the balanced tensor coefficient \(V_{2,T}=V_{\rm reg}/Z\),
the real logarithmic period columns in the order
\((\log|Z|,\log|V_{2,T}|)\) are
\[
 \begin{gathered}
 \binom{-2A-C}{C},\qquad\binom{3A}{-D_B},\\
 A=-\log|P|,\quad C=-\log|R|,\quad
 D_B=-\log|\lambda qK|,\quad K=(P/t)^3\eta_2^6 .
 \end{gathered}
 \tag{ASI2.8}
\]
Their determinant is
\[
                    (2A+C)D_B-3AC .
 \tag{ASI2.9}
\]
The full units in \(K,c_q,u_s\), and their reciprocals,
are bounded on one smaller closed original bidisc. The exact
equation \(R-P=su_s\) gives
\[
 D_B\geq12\min(A,C)+C+C_*
 \tag{ASI2.10}
\]
after choosing the stated strict contraction range for
\(\lambda_0\) and then shrinking the bidisc; \(C_*>0\) is
fixed. Specifically a sufficient choice of
\(-\log|\lambda_0|\) is larger than the supremum of
\(\log|c_qK|-12\log|u_s|+12\log2\), plus \(C_*\).
Consequently ASI2.9 is bounded below by
\[
 \begin{cases}
 23AC+13C^2+C_*(2A+C),&A\geq C,\\
 24A^2+11AC+C^2+C_*(2A+C),&A\leq C.
 \end{cases}
 \tag{ASI2.11}
\]
The exact inverse matrix therefore gives the marked real period
coefficients. Their values in \([0,1]^2\), together with the
compact angular variables, give a fundamental cut.

For completeness the finite-support comparison is explicit.
In the three endpoint lattices, the valuation period matrices are
\[
 M_0(r,z)=
 \begin{pmatrix}10r+z&8r+z\\3r&6r+z\end{pmatrix},\quad
 M_1(r,z)=
 \begin{pmatrix}3r&0\\r&12r+12z\end{pmatrix},\quad
 M_\infty(r,z)=
 \begin{pmatrix}23r&11r\\3r&3r+z\end{pmatrix}.
 \tag{ASI2.12}
\]
The zero support is the six-vertex polygon
\[
 (0,0),(-5r-z,0),(-10r-z,-3r),
 (-13r-2z,-6r-z),(-8r-z,-6r-z),(-3r-z,-3r-z),
 \tag{ASI2.13}
\]
the one support is the parallelogram generated by the two columns
of \(M_1\), and the infinity support is
\[
 (0,0),(-12r,0),(-23r,-3r),(-23r,-3r-z),
                         (-11r,-3r-z),(0,-z).
 \tag{ASI2.14}
\]
The matrices and polygons retain their original lattices; the
index-six and index-thirty-three inspection maps used to derive
them are inverted only on their stated image lattices.
Direct multiplication by \(M_0^{-1}\), respectively
\(M_\infty^{-1}\), puts every polygon vertex in \([-2,2]^2\).
For example the nonintegral zero coordinates have numerators
\(-30r^2-11rz-z^2,\ 15r^2+3rz,\ -21r^2-10rz-z^2,\ 6r^2+2rz\)
over \(36r^2+13rz+z^2\). At infinity they have numerators
\(-36r-12z,\ 36r,\ -23z,\ 11z\) over \(36r+23z\).
The asserted inequalities follow coefficient by coefficient
after multiplication by these positive denominators. The one
parallelogram has period coordinates exactly \([0,1]^2\).
Their period translates tile the stated native support, with the
full unit arrows recorded in the parent construction. Thus a
fixed finite set of translates, for example indices in
\([-2,3]^2\), covers the cut in valuation coordinates.

The full logarithms perturb those cuts through the exact matrix
ASI2.8. On each endpoint domain, the inverse of ASI2.12 is
uniformly bounded by a constant times the reciprocal of the
smaller retained nonzero original log-radius. Shrinking that
radius bounds the product of this inverse with the complete
unit-log perturbation by \(1/2\). The inverse
\((I+M^{-1}\Delta)^{-1}=\sum_{k\geq0}(-M^{-1}\Delta)^k\)
then gives the actual norm-coordinate comparison. Hence only
a fixed finite enlargement of those period translates meets
the exact cut. All their coordinate moduli lie in closed
polydiscs of the finitely many proper parent charts.
On the one lower rank-one face the remaining source period is
the full \(L:(X,Y)\mapsto(P^3pX,PpY)\); its source multiplier
is a strict contraction. On the infinity lower rank-one face
the remaining period is the full \(L:(X,Z)\mapsto
(R^{23}f_\infty^2X,R^3Z)\), again with a strict source
contraction. The chain labels and their source/line factors
are retained. Those contractions prove local finiteness on
these faces and give their actual completed source cycles.
Away from them the inverse real period matrix proves local
finiteness directly. The separated parent charts, their
proper subdivisions and these locally finite quotient maps
therefore give a Hausdorff compact cut.

Choose a closed original \(t\)-disc \(U\) containing \(0\)
and \(t_c(s)\), and excluding the other simple singular value
and the actual infinity cap. Its boundary misses all source
critical values for the closed small source-parameter disc.
It retains every horizontal point/O section chart. Round
its two base corners in the actual monomial coordinates
with an invariant radial function. The compact resulting
collision piece is denoted \(W_c\).

\par

\medskip\noindent\textbf{ASI2.2. Whole point/O boundary, all source discs and the reference body}

\par\nobreak

For a source band \(j\), put \(x=Q^jZ,y=Q^{1-j}/Z\).
The full graph and point/O coefficients are
\[
 x=P^2a,\ ay=R;\qquad xc=P^2,\ y=Rc,\qquad
 P=uw,\ (x,c)=(u,uw^2)\ \hbox{or }(u^2w,w),
 \tag{ASI2.15}
\]
\[
 V_{\rm reg}=-\eta_2^{-1}P^{-2j}E A
   =\eta_2^{-1}P^{-2j}c^{-1}EC,\quad C=-cA,
 \qquad
 V_{\rm reg}=\eta_2^{-1}P^{-2j-2}(x-P^2)E_OC_O.
 \tag{ASI2.16}
\]
Here \(E,E_O\) are the complete convergent products SI2.17--19.
The first graph charts have \(ds/dP\) a unit with value \(-1\);
the two second charts have \(ds/dR\) a unit with value \(+1\).
Every inverse recovers the original \(P,R\) first, and then
uses the reciprocal coefficient in the same Cartier module.
The additional \(P^{-2}\) at \(O\) remains. On the finite
original \(t\)-boundary \(P\ne0\), these are precisely the
original source and line maps through the section charts.
At a section coincidence the module numerator and denominator
are carried as a line-module family.

Let \(W_{r,1}\) be the actual RB scalar-disc reference on the
fixed source \(s_*\), with scalar \(\lambda_0\sigma^{12}\),
root \(r_0\sigma^4\), and the same finite character and
sector integers. The ordinary chain is \(U_kW_k=\sigma\).
At the selected node its full equation is
\[
             U_jW_j=\sigma^{12}q,\qquad
                       x_i\bar y_i=q,
 \tag{ASI2.17}
\]
with the ordered factors
\((\sigma,\ldots,\sigma,\bar y_i,x_i)\), twelve copies
of \(\sigma\). For factors \(f_1,\ldots,f_{14}\) its
fourteen charts are
\[
 f_r=s_rt_r,\quad
 U=s_r\prod_{\nu<r}f_\nu,\quad
 W=t_r\prod_{\nu>r}f_\nu .
 \tag{ASI2.18}
\]
On each adjacent overlap the inverse divides by the
retained nonzero factor. RB.14--17a are their literal
source-band, point-line and period maps. This includes the
full \(q\), \(\kappa,\alpha,g,\Gamma\) factors.
On a source-unit patch it reduces by those displayed
coordinate divisions to the ordinary chain. The twelve
fixed source components have \(\alpha_l=\beta_l=0\);
the overlap with the selected moving curve requires the
specified nonzero \(v_{12}\). Thus every source component
retains all three source vanishing discs. The reference\textquotesingle s
moving component is the distinct fixed sphere times its
parameter disc supplied by the last two factor charts.
Its selected fan is \(T_1\), with matrix
\(\left(\begin{smallmatrix}1&-1\\1&0\end{smallmatrix}\right)\).

The side boundary from \(W_c\) to this actual reference is
the whole IL.26--31 map. In the one characters
\(X=Z^{-1},Y=V_{\rm reg}/Z,v=p-1\), it is
\[
 \alpha_l=Yv^{-l},\quad\beta_l=v^{l+1}/Y,\qquad
 A_l=H_{+,l}\alpha_l,\quad B_l=H_{-,l}\beta_l,
\]
\[
 H_{+,l}=\eta_2u_s^lP^{-l}X^{-1},\quad
 H_{-,l}=\eta_2^{-1}XP^{l+1}u_s^{-l-1},\quad
                    H_{+,l}H_{-,l}=P/u_s.
 \tag{ASI2.19}
\]
For an arbitrary label \(12j+l\), the actual inverse
source period is
\[
 X_*=p^jX,\quad Y_*=P^{-12j}v^{-12j}c_{12}^{-j}Y.
 \tag{ASI2.20}
\]
The two target coefficients in ASI2.19, evaluated at
that recovered source point, also carry
\(\gamma^j,\gamma^{-j}\), with
\[
                 \gamma=pu_s^{12}c_{12}
                        =P^2\lambda_0q\eta_2^6/t^3.
 \tag{ASI2.21}
\]
These follow from the full \(T=hB\) arrow, including
its additional source \(h\) factor. They preserve
\(A_lB_l=s\) and have their stated reciprocal inverse.

The source transport is the actual compact relative
flow \(T_s:S_0\to S_s\) of IL.28--30, with final target
\(K_s=T_{s_*}T_s^{-1}\). Its boundary formula holds
the marked real periods fixed:
\[
 Z_s=e^{b\operatorname{Log}(Q_s/Q_0)}Z_0,\qquad
 V_{{\rm pt},s}
      =e^{b\operatorname{Log}(P_s^2/P_0^2)}V_{{\rm pt},0}.
 \tag{ASI2.22}
\]
Both logarithms have value zero at \(s=0\); under the
source deck \(b\mapsto b+1\), their targets acquire
exactly \(Q_s,P_s^2\). The local vector fields
\((\partial_Ps)^{-1}\partial_P\) and
\((\partial_Rs)^{-1}\partial_R\) on ASI2.15
project to one. A finite relative partition, agreeing
with the derivative of ASI2.22 on the boundary,
gives the complete source flow; its backward flow is
the inverse.

For genuine vector frames \(e_j=g_{ij}e_i\), its line
connection has
\[
 a_i=-\sum_k\chi_k\,Xg_{ik}/g_{ik},\quad
 a_j=a_i+Xg_{ij}/g_{ij},\quad
              X^{\rm hor}=X-a_iw_i\partial_{w_i}.
 \tag{ASI2.23}
\]
A relative affine correction makes this equal to the
full ASI2.16 boundary derivative. The line ODE and
its reciprocal dual ODE supply the entire target
\((K_sz,K_s^{\rm line}H_+\alpha,
K_s^{\rm dual}H_-\beta)\), with inverse first
\(K_s^{-1}\) and then the reciprocal coefficient.
On every period translate it is conjugated by the
full original period and its inverse. This retains
the derivative of every unit and integral source
deck. Since \(v=\alpha\beta\), its first normal
derivative on the fixed locus is zero. Circle
averaging keeps the source lift and its complex
diagonal normal derivative. This matches the entire
side-boundary tangent and line frames.

The actual three source discs give
\[
 H_1(\partial S)=\mathbb Z\langle b,\alpha_t\rangle
                 \oplus\mathbb Z/3\langle a\rangle
       \longrightarrow H_1(S)=\mathbb Z\langle b\rangle,
 \quad (a,\alpha_t,b)\mapsto(0,0,b).
 \tag{ASI2.24}
\]
The full reference cube frame is
\(e^{-2\tau_L\operatorname{Log}p}X^2V_{\rm pt}^3\),
\(\tau_L=\log|X|/\log|Q|\), with inverse divisor
\(6F-2C_1-C_2\). Its order-three flat boundary connection,
the primitive \(P/\epsilon\) fibre clutch, the complete
point/O modules and all \(\gamma^j\) are those of IL.33--38.
Thus ASI2.19 is the actual boundary map used in the
twelve closed normal doubles, not a choice of their
boundary degree.

\par

\medskip\noindent\textbf{ASI2.3. The original finite cap and full exterior cylinder}

\par\nobreak

On the infinity root keep the literal source map
\[
 q_\infty=z^3,\quad t=z^{-3},\quad
 x=z^{-2}X',\quad y=z^{-3}Y',\quad
 Y'^2=X'^3+r_{\rm src}z^2X'^2+zX'+1.
 \tag{ASI2.25}
\]
The deck is
\((z,X',Y')\mapsto(\zeta_3z,\zeta_3^{-4}X',Y')\).
The marked sections are
\(G'=(0,1)\), \(2G'=(sz^2,-1-sz^3/2)\).
The actual point correction on the two quotient graph
rays \((2/3,1/3),(1/3,2/3)\) is
\[
       \beta^*\mathcal O([2G']-[O])
                    =\pi^*M_2+E_P-E_O.
 \tag{ASI2.26}
\]
Its local monomial coefficient is
\(\min(v_w,2v_z)\) at the two marked sections,
with the specified Cartier generators and reciprocal.
The small division \(3R'_0=3G'\) has
\[
 \tau(3G')=
 z^2\frac{-2s(2+sz^3)}
 {8+6sz^3+(s^2-s^3)z^6},\qquad
 \tau(R'_0)=sz^2\ell_R,\quad\ell_R(0,0)=-1/6.
 \tag{ASI2.27}
\]
It preserves the actual graph ideals
\((w,z),(w,z^2),(z,w^2)\), with its inverse.
The normalized chord from SI2.26 has value one at
\(\tau=0\), has its stated horizontal divisor, and
commutes with the deck because the chord and
\(\tau\) weights cancel. Write this full Hom as \(F\).
The complete root coefficient is
\[
 B'_{\rm sel}=h_\infty B'_{2,3},\quad
 h_\infty=1-t_c(s)z^3,\quad
 \beta=B'_{\rm sel}/F^3,\quad
 \widetilde F=\beta^{1/3}F,\quad
 A_\lambda=r_\lambda\zeta_3^d\widetilde F D_3 .
 \tag{ASI2.28}
\]
The cube is chronological. The analytic invariant
cube root is fixed by its central value, and
\(A_\lambda^3=B_\lambda=\lambda B'_{\rm sel}\)
as a full source and line map.

On an original infinity sector set
\(\ell_\alpha=\log_\alpha q_\infty/(2\pi i)\),
\(z_\alpha=e^{2\pi i\ell_\alpha/3}\). The actual
annular arrow is
\[
 \Phi_\alpha=
 \operatorname{Exp}\{\ell_\alpha\operatorname{Log}_\alpha
       \widetilde F+
       \ell_\alpha(\log r_\lambda/(2\pi i)+d/3)\delta
       +h_\alpha\}\,I_{z_\alpha}.
 \tag{ASI2.29}
\]
It satisfies
\(\Phi_\beta=A_\lambda^{j_0}
B_\lambda^{j_{\alpha\beta}}\Phi_\alpha\) and
\(\Phi_\alpha B_\lambda=B_\lambda\Phi_\alpha\).
The map \(I_z\) carries ASI2.26; the original reference
cochain \(h_\alpha\) is the fixed-contour, reference-zero
solution SI2.28/OG.12--13 with the entire source periods
and line lattice. No integer sector term is reset.
Its inverse reverses the actual Lie, root and
Cartier arrows at the recovered source point.

On the outer scalar circle the finite cap is a compact
proper family over the source interval. Extend the
prescribed side-boundary source field through its
actual root charts, agreeing on its annulus by the
derivative of ASI2.29. The difference of two such
lifts is vertical; a relative partition and finite
deck average preserve the prescribed interval
component. The full line connection is ASI2.23 in
the actual corrected modules. Compactness gives
its interval flow and backward inverse. On the
ordinary exterior use the original smooth source
charts and the same construction. On a period
translate use the actual conjugating map and its
inverse, so the complete functions and derivatives
remain. This produces the actual exterior cylinder
between the coupled and fixed-source families.
It is constructed on the nonzero outer scalar
circle, where all original cap maps are defined.

The scalar root coordinate \(\rho\) has the exact
forward square
\[
             \Phi_{\rho+4}q^4=h_\infty^4\Phi_\rho,
 \tag{ASI2.30}
\]
with endpoint-to-initial arrows \(q^{-4},h_\infty^{-4}\).
It follows from the entire identity
\(q_\infty q=h_\infty\), including ASI2.29\textquotesingle s
integral cochain. The field just constructed
transports this square literally. At \(q=0\) the
ordinary arrow is its actual \(B^4\) reindex followed
by the full bounded source/line factor NS.24. Thus
the selected zero is retained rather than divided out.
This closes every original ordinary/root boundary
of the relative piece \(W_c-W_{r,1}\).

\par

\medskip\noindent\textbf{ASI2.4. Both native Hopf caps and the entire inner ledger}

\par\nobreak

The two actual collision native charts are
\[
 q=Za_0c_0=z a_\infty c_\infty,\quad
 z=Z^{-1},\quad(a_\infty,c_\infty)=Z(a_0,c_0),
 \quad c_0=ZV_{\rm pt},\quad a_0=q/(Z^2V_{\rm pt}).
 \tag{ASI2.31}
\]
NS.8--15 give their holomorphic extension and inverse
at both vertices, recovering the original \(s,q\)
and then every full unit. Their original volume is
\[
 \Omega_{\rm orig}=\frac2{a_T}\,
 ds\wedge dZ\wedge da_0\wedge dc_0
 =-\frac2{a_T}\,ds\wedge dz\wedge da_\infty\wedge dc_\infty .
 \tag{ASI2.32}
\]
The endpoint is
\[
 E_0=(P^{12}Z,P^{12}\mathcal A^{-1}a_0,
                         P^{-24}\mathcal A c_0),
       \qquad\mathcal A=(c_q/f_0)^4 .
 \tag{ASI2.33}
\]
It preserves the whole \(q\), and its reciprocal
inverse evaluates the coefficients at that same
recovered original \(s,q\).

The NC cut map has
\[
 C_\theta(Z,a,c)=
      (\alpha_\theta Z,(\alpha_\theta\beta_\theta)^{-1}a,
                                       \beta_\theta c),
 \quad C_1E=C_0,
 \tag{ASI2.34}
\]
where its two full logarithms are the NC.5 logarithms
of the corresponding two factors in ASI2.33, with
the specified endpoint-flat cutoff. Its velocity
annihilates \(q\). Cut off only that vertical
velocity, keeping the original parameter-horizontal
lift, and integrate on both native charts. The
backward flow is the inverse. Its support misses
the source boundary and the actual infinity cap.

Use the held Cox quaternion
\[
 T_Z=\begin{pmatrix}1&-\bar Z\\Z&1\end{pmatrix},\quad
 q_H=T_Z(a,\bar c)^t,\quad
 T'_{1/Z}\operatorname{diag}(Z,\bar Z)=T_Z.
 \tag{ASI2.35}
\]
The entire normal matrix is
\[
 G_\theta=T_{\alpha_\theta Z}
       \operatorname{diag}((\alpha_\theta\beta_\theta)^{-1},
                              \overline{\beta_\theta})T_Z^{-1}.
 \tag{ASI2.36}
\]
Its inverse is the reversed matrix product at the
recovered base point. On the Hermitian three-ball
the exact fixed-sphere map is
\[
 \Psi_\alpha(Y)=
 \frac{\operatorname{diag}(\alpha,1)Y
                    \operatorname{diag}(\alpha,1)^*}
 {\operatorname{tr}(\operatorname{diag}(\alpha,1)Y
                    \operatorname{diag}(\alpha,1)^*)}.
 \tag{ASI2.37}
\]
Its denominator is positive and its inverse is
\(\Psi_{\alpha^{-1}}\). Its rank-one boundary is
the actual \(Z\mapsto\alpha Z\) map. In ASI2.36
multiply the two full cut logarithms by a radial
function zero near the ball centre and one near
the boundary. The resulting \(\widehat G_\theta\)
is direction independent near the centre. The
map
\[
 \widehat C_\theta(Y,q_H)=
 \left(\Psi_\alpha Y,
       \widehat G_\theta(Y)q_H/
                         \|\widehat G_\theta(Y)q_H\|\right)
 \tag{ASI2.38}
\]
has inverse first \(\Psi_{\alpha^{-1}}\), then
\(\widehat G_\theta^{-1}\) at that recovered
point. Its seam is
\(\widehat E=\widehat C_1^{-1}\widehat C_0\).
For the actual resolved normal derivative \(N_\theta\)
of NS.23, the directed equality is
\[
          B_\theta=C_\theta^{-1}N_\theta,\qquad
                    \widehat C_\theta B_\theta=N_\theta.
 \tag{ASI2.39}
\]
The reciprocal equality uses the inverse maps in
the reversed order. This attaches the actual
product Hopf cap with every normal coefficient
and positive spinor frame retained.

The reference has its own regular native
coefficients \(c_R=V_{\rm pt}/Z,a_R=q/V_{\rm pt}\),
its exact second chart from RB.17a, and its own
endpoint MN.12. Applying ASI2.34--39 with those
full reference factors is exactly its specified
own cap, as proved in MN.3. We do not identify
the two native vertex charts by the dense
\(Z^2\) expression; their complete reference
continuation is constructed in ASI2.5 below.

The fixed-face inventory examines every regular
cone and every actual paired period facet. Its
minimal nonzero-weight faces, exact stabilizer
words and complete stars give the following
disjoint list:

\begin{itemize}

\item
  106,166 compact generic curves, with all three
  complete endpoint normal and determinant maps;
\item
  15,138 compact generic surfaces with their
  saturated quotient stars and 52 actual blowups;
\item
  zero corner interiors: 6,163 points, 72 curves
  and 17 complete compact surfaces;
\item
  one corner interiors: 4,764 points and 2,340 curves;
\item
  infinity corner interiors: 7,304 points and 701 curves;
\item
  the twelve lower/reference source doubles of IL;
\item
  the single NS moving face and the RB moving
  reference face, with their respective caps.
\end{itemize}

The first five rows are exactly CFL.1--9. All have
their complete period alignments and full analytic
normal units. Every generic component has already
been completed across all three ends; the listed
corner-interior tags are disjoint from those
generic boundary pieces. The twelve lower
stabilizers are their actual source \(L\) groups.
The moving face is the single orbit identified
by the full \(T_0=hB^{-1}\) map NS.5. No additional
fixed orbit remains in the exhaustive face
inventory or in the exterior cylinder.

Their first five rows give \((2723,19124)\).
IL gives the next row \((-2712,-19092)\).
For the moving row, retain the actual positive
normal bundle \(E=\overline{N_-}\oplus N_+\).
NC.19--20 is an explicit five-dimensional bundle
bordism from its own cap-completed surface to
the opposite-disc closed replacement. On that
replacement the nine actual source blowups
give the MN.7 roots
\[
 x=h-4u+E_4+E_5+E_6-E_7-E_8,\quad
 y=-h-5u+E_4+E_5+E_6+E_7,
 \tag{ASI2.40}
\]
with \(h^2=u^2=0,hu=1,E_iE_j=-\delta_{ij}\).
Thus \(c_1(E)^2=-13,c_2(E)=-3,p_1(TF)=-27\),
transferred through that actual bordism.
The reference roots are \(h-12u,-h-12u\);
its own cap gives \(c_1^2=c_2=p_1(TF)=0\).
From the actual physical Euler denominator,
writing \(x=-c_1(N_-),y=c_1(N_+)\), the
physical surface contributions are
\[
 k_M=c_1(E)^2-c_2(E),\qquad
 \ell_M=p_1(TF)+c_1(E)^2-4c_2(E).
 \tag{ASI2.41}
\]
These are IL.41 after the literal negative-line
conjugation. Hence the moving collision-minus-
own-reference cap pair is \((-10,-28)\).
The full inner ledger with the RB \(T_1\) reference is
\[
 (2723,19124)+(-2712,-19092)+(-10,-28)=(1,4).
 \tag{ASI2.42}
\]
This ledger includes its entire reference exactly once.

\par

\medskip\noindent\textbf{ASI2.5. Positive outer Stokes, Jin reversal and the actual reference change}

\par\nobreak

On the free quotient use the physical nonzero-vector
line \(M\). On a finite stabilizer chart use its
actual covering slice and reciprocal covering degree.
Average the tangent and line connections over the
physical circle and finite deck group. Across each
identified boundary transport those connections
through the full derivative and line arrow above;
place the identical interpolation on the opposite
collar with its interval reversed. Its complete
transgression cancels pointwise. The stable tangent
splitting into the quotient tangent and the
physical orbit direction is the actual one in
these slice charts. Stokes therefore applies to
\(c_1(M)^3\) and \(p_1c_1(M)\), including every
finite covering factor in CFL and IL.

The orientation is fixed by the original complex
source and the physical circle:
\[
 \operatorname{or}(W)=d\phi_\delta\wedge
           \operatorname{or}(Y)\wedge dr_s\wedge d\theta.
 \tag{ASI2.43}
\]
The outward parameter boundary in the circle-first
quotient is \(-\operatorname{or}(Y)\wedge d\theta\),
as the radial normal crosses five directions.
The removed fixed tubes have the two opposite
reversals of the actual normal-Hopf boundary map
ASI2.35. These give exactly the complex
fixed-surface orientations used in CFL, IL and
MN. Equations ASI2.32 and NC.20 fix the native
and cap signs, including the original nonzero
volume unit. Consequently, with \(O\) denoting
the completed positive-parameter physical pair,
\[
                 O_c-O_{r,1}=(1,4).
 \tag{ASI2.44}
\]

For the same actual lifted relative return \(f\),
the map
\[
 \begin{gathered}
 Q_f^+=(Y\times[0,1])/((y,1)\sim(f(y),0))\\
 \longrightarrow M_f=(Y\times[0,1])/((y,0)\sim(f(y),1)),\\
 [y,t]\longmapsto[y,1-t].
 \end{gathered}
 \tag{ASI2.45}
\]
has line map \([v,t]\mapsto[v,1-t]\), with the same
full endpoint lift. It sends the two displayed
endpoint relations to one another. Its inverse
is the identical formula. It preserves the
physical line and reverses the parameter
orientation. Write \(J=-O\) for the resulting
Jin-oriented pair.

The whole RFB construction has the precise type
\[
                    J_{r,-1}-J_{r,1}=(-72,0).
 \tag{ASI2.46}
\]
Indeed RFB.33--40 computes with \(U=M^*\);
its physical dual and original-to-Jin map give
\(k_M=+\int F_U^3\) in that coefficient
orientation, explicitly as in RFB.40--43.
Thus in the positive outer orientation
\[
                    O_{r,-1}-O_{r,1}=(72,0).
 \tag{ASI2.47}
\]
These are directed equalities on the entire
fixed-source scalar reference word \((0,4)\).

We recall its actual morphisms to specify its
use here. For \(j=1,0,-1\),
\[
 Z_j=Z,\quad V_j=Z^{-j}V_{\rm pt},\quad h_j=hB^j,\quad
 \kappa_j=\kappa g^j,\quad
 \alpha_j=\alpha\eta^j\kappa^{-j}g^{-j^2},\quad
 \eta_j=\eta g^{-j}.
 \tag{ASI2.48}
\]
Here \(\eta=\sigma^{12}\lambda_0\eta_2^6t^{-3}\),
\(\kappa=P^2/c_q,\alpha=P^2,g=P^{-3}\), all at
the fixed original source. Every inverse is
the displayed reciprocal character and period
substitution. Each reverse \(C_1\) change has
the full Cox map
\[
 (a',b',c')=(a/b,b,bc),\quad
 (a,b,c)=(a'b',b',c'/b'),
 \tag{ASI2.49}
\]
on its dense overlap, and the actual two
projections of \(\operatorname{Tot}\mathcal O(-1,-1)\)
\[
 p_+(A,C,\tau)=[A_0,A_1,\tau C_0,\tau C_1],\quad
 p_-(A,C,\tau)=[\tau A_0,\tau A_1,C_0,C_1].
 \tag{ASI2.50}
\]
They pull \(q\) back to
\(\tau^2A_0A_1C_0C_1\); their inverses off the
zero sections are the specified projective
direction and nonzero fibre coefficient.
The entire bounded Cox lift and its inverse
are RFB.17. The six chart rows RFB.7 retain
every \(\kappa_j,\alpha_j,g,\eta_j\).
The common link/exterior map is
\((y,q_H)\mapsto(q_H^{-1}yq_H,q_H)\), with
its conjugation inverse and unrestricted
four-weight Hopf lift.

Their exact cut square is
\(C_jS_j=S_jC_0\), where
\[
 S_j=(Z,Z^ja,Z^{-j}c),\quad
 \mathcal D_j=\mathcal D\mathcal A^{-j}.
 \tag{ASI2.51}
\]
The full inverse square and the same
Hermitian-ball construction ASI2.37--39
carry it through both Hopf caps. The
ordinary cubic cap and ASI2.30 are fixed.
RFB.24--38 construct the completed quotient
map and inverse by the actual log-norm
period coefficients, source/graph flows,
Cartier spinor arrow and relative Cech
line map. The latter is
\(U_j\to I_j^*U_0\otimes\mathcal T^{12j}\)
with the exact endpoint-to-initial clutch
\(w\mapsto e^{2\pi ib}w\), all outer
primitive \(-j(du+\alpha_qdv)\) retained.
The full nilpotent tangent collar has
zero \(p_1\) transgression. Its relative
regular-level cobordism and the entire
Euler collar give RFB.43, and the actual
integer \(0\to1\) Cox comparison fixes
its linear coefficient. These maps prove
ASI2.46 for this reference and its own
three native caps. They are the required
finite continuation through \(T_1,T_0,T_{-1}\).

At \(j=-1\), \(a_j=q/(Z^2V_{\rm pt})\),
\(c_j=ZV_{\rm pt}\). Hence this final
reference has exactly ASI2.31\textquotesingle s native
chart, point marking, full endpoint
and cap product. Combining the directed
whole-reference equality with ASI2.44 gives
\[
 O_c-O_{r,-1}=(-71,4),\qquad
                    J_c-J_{r,-1}=(71,-4).
 \tag{ASI2.52}
\]
The correction is made once on the whole
reference after all its components and
its own cap have been included.

\par

\medskip\noindent\textbf{ASI2.6. All four parameter edges, retained winding and the source word}

\par\nobreak

Use the independent source cut \(\tau\) and
relative scalar-root coordinate \(\rho\),
based at the actual \(r_{\rm base}\).
The source, point and root family has
\(A_\rho=e^{2\pi i\rho}A_{\rm base}\),
\(B_\rho=e^{6\pi i\rho}B_{\rm base}\).
On its original sector, the complete
cochain is the same ASI2.29 with the
corresponding \(\log r_\lambda\), and
its retained integer shifts. Equivalently
its varying part is \(\rho\ell_\infty\delta\).
For every integer root turn the full
ordinary/cap forward arrows are
\(q^j,h_\infty^j\), with reciprocal
endpoint-to-initial arrows and the
actual \(B^j\) reindex at the selected zero.
At four this is exactly ASI2.30.

Put
\[
 G=\operatorname{Log}(P^{-3}),\quad
 J=\operatorname{Log}(\lambda_{\rm base}\eta_2^6t^{-3}),
 \quad
 A_\rho=e^{\rho G},\quad
 D_\rho=e^{\rho J+3\pi i\rho(\rho-1)}.
 \tag{ASI2.53}
\]
Here \(A_\rho,D_\rho\) denote the two point-coordinate
factors in this display; the root deck action
above remains the specified three-dimensional
cap map. Their types distinguish the two
uses. The native current-to-reference arrow is
\[
 C_\rho^N:(Z,a,c)\mapsto
      (A_\rho Z,(A_\rho^2D_\rho)^{-1}a,
                                  A_\rho D_\rho c).
 \tag{ASI2.54}
\]
It preserves the entire \(q\); its inverse
recovers that \(q\) and divides by the same
full functions. With
\(\eta_\rho=e^{6\pi i\rho}e^J\), subtraction
of the complete logarithms gives
\[
 C_{\rho+1}^N\,S(g^{-1},\eta_\rho^{-1})=C_\rho^N,
 \quad C_0^N=1,\quad C_4^N=S(g^4,e^{4J}),
 \tag{ASI2.55}
\]
where \(S(a,d)=(aZ,(a^2d)^{-1}a_N,ad\,c_N)\)
with \(a_N,c_N\) its independent normal
coordinates. This is the inverse of the
actual forward endpoint at the same
scalar-root endpoint. At the source-cut
edges, both \(G,J\) increase by \(-6\pi i\).
The native mismatch factors are therefore
\[
          (A,\beta)=(e^{-6\pi i\rho},e^{-12\pi i\rho}).
 \tag{ASI2.56}
\]
They are retained for every intermediate
integer corner, including the full
constant and based-unit values at the
initial source edge.

Writing \(w=e^{\pi i\rho}\), the full
held-Cox matrix identity is
\[
 T_{w^{-6}Z}\operatorname{diag}(w^{18},w^{12})
       =\operatorname{diag}(w^{18},w^{12})T_Z.
 \tag{ASI2.57}
\]
In its right-complex quaternion coefficient
this is exactly
\[
                         q_H\mapsto w^3q_Hw^{15}.
 \tag{ASI2.58}
\]
Both half-phase endpoints have sign
\((-1)^\rho\) at every integer corner.
The left adjoint is the actual rotation
on the fixed sphere and normal orbit
sphere, and extends linearly over their
specified three-balls. The full right
factor is its principal Hopf-line lift.
Thus this is an actual cap boundary map,
with closed paired Spin endpoints.
For arbitrary complex moduli the complete
matrix ASI2.36 and its positive Hermitian
ball map provide the same full extension,
without replacing the moduli by phases.

On the coupled diagonal \(\rho=4\tau\),
the original DZ product and ASI2.54
differ by the directed physical scalar
from NS.37:
\[
 F(\tau)=e^{3\pi i(16\tau^2+4\tau)}
       =e^{48\pi i\tau(\tau-1)}e^{60\pi i\tau},
 \quad(a,c)\mapsto(F^{-1}a,Fc).
 \tag{ASI2.59}
\]
The first displayed logarithm vanishes at
both endpoints and contracts by scaling
that logarithm. The second is the actual
positive physical loop of winding thirty.
Its held-Hopf factor is the inverse right
phase; it acts by the same global physical
circle on every original point, root,
Cartier and cap chart. Its degree-three
pullback is zero because it depends only
on \(\tau\). It commutes with the full
source and finite-deck maps, extends on
the product Hopf cap, and has paired
Spin endpoint \(+1\). The nonzero winding
is retained in this map. Replacing the
unflattened \(\tau\) by the NC endpoint-flat
cutoff adds only the displayed full
zero-endpoint logarithm differences;
the cutoff vertical flow transports all
normal jets and leaves ASI2.30 fixed.

The whole scalar reference at \(j=-1\)
uses the same native conjugation ASI2.54:
in RFB.19 its \(D_j=D A^{-j}\) is \(DA\).
Its source-cut quaternion is therefore
exactly ASI2.58. This is the actual
same-fan product obtained after RFB,
including its cap and original point-line
coefficient. All four edges of the
source/root parameter rectangle are now
specified by ASI2.29--30 and ASI2.53--59.

The coupled word is \((1,4)\), the
fixed-source scalar word is \((0,4)\),
and the fixed-scalar source word is
\((1,0)\). Pull back the displayed
family by the endpoint-fixed polygonal
homotopy from the diagonal to the source
edge followed by the scalar edge.
On each corner the full functional
equation ASI2.55 and the finite-cap
square ASI2.30 agree. The Hermitian
ball and full normal matrices extend
those maps over the selected caps,
with the retained winding ASI2.59.
Thickening these intervals gives the
actual pair-of-pants bordism of the
completed six-manifolds and their
physical line. Connections extend by
the same boundary derivatives.
Stokes on this bordism is additivity
of the two integer characteristic
numbers. Consequently ASI2.52 is
exactly the Jin physical pair of the
fixed-scalar source word, proving
the first equality in ASI2.3.

\par

\medskip\noindent\textbf{ASI2.7. The oriented sphere, native pair and marked generator}

\par\nobreak

The outer fibres are the original
oriented spheres of F41/F46 with the
proved \(M_{2}\) good cubic cap. An explicit
path retains the selected colliding
branch and \(r_{\rm src}\) real from
\(0\) to \(1/4-s_*\), with
\(0<s_*<1/12\). K96 gives the full
\(M_{2}\) source/Cartier/good-root family at
the initial unselected III value.
The remaining interval avoids the
source collisions and selected cusp.
The exact SI2 Miller, point/O, source
and cubic-root maps continue along
that interval with the prescribed
finite character and integral
annular reference. A finite compact
parameter cover gives a common strict
contraction bound. Its invariant
horizontal fields give the actual
equivariant fibre diffeomorphism,
with backward inverse. The original
Hopf tube and primitive attaching
maps of F41/F35 travel through this
same source/cap family; their
Gysin and van Kampen maps retain
the sphere topology for the stated
selected fan.

The physical line and held quaternion
are the ones in ASI2.35--39. Thus
K68\textquotesingle s explicit physical-line inverse
clutch identifies the class of the
return with
\[
        (4k_M-\ell_M)/24=(284+4)/24=12.
 \tag{ASI2.60}
\]
The actual fixed-\(C_0\) free-coordinate
map in that same product is
\(k_M=\mu+\nu,\ell_M=4\mu\).
It is obtained by the same directed
Hopf transfer and negative-normal
conjugation, with both current cap
products fixed in ASI2.39. Hence
the source native pair is
\[
                         (\mu,\nu)=(-1,72).
 \tag{ASI2.61}
\]
This division by four is the stated
integral Hopf coordinate formula;
no division of a homotopy-sphere
class is used.

Finally the original marked-node
cover has \(v=w^2-1\),
\(D=v-2\), \(s=v(v+4)/(4D^2)\).
At \(w=-1\), \(Z_T=Z_w\), and
\(ds/dw\ne0\); the positive small
\(w=-1\) circle is exactly one
positive original \(s\)-meridian.
The full point/source map at the
other mark first sends
\((Z,V_{\rm pt})\mapsto(Z^{-1},P^2V_{\rm pt})\),
then uses the actual target period
\((Q,P^{-2})\), giving
\((Q/Z,V_{\rm pt})\).
Its reciprocal and all \(r_2\)
Cartier factors are SI2.21.
Thus the result transports to the
positive \(w=+1\) loop on the
literally reflected fan \(JT_{-1}\).
The original source and good cubic
cap depend on \(w^2\), and their
finite character and annular
cochain are unchanged by that
specified source identity. The
fan marking in ASI2.2 is retained
through the entire statement.


% End expanded verified cp36 input

\end{proof}


\subsection{The complete selected-I3 M2 return on every good cap and fan}
\begin{theorem}[Source result P3K.112 under continuation review]\label{night:P3K.112}
\StatementRecord{night:P3K.112}{open}{Historical events: GPT-6 Astra (OpenAI), ChatGPT Pro mode, effort Pro, S6\_06\_BLIND, 16--18 November 2026: targeted proof check. Also checked in this writing. Continuation mathematical and source checks: GPT-6 Astra (OpenAI), ChatGPT Pro mode, effort Pro, S6\_10\_PAPERS continuation, 19 November 2026: computation or formal certificate; used and rechecked in this run. Continuation metadata events: GPT-6 Astra (OpenAI), ChatGPT Pro mode, effort Pro, S6\_10\_PAPERS continuation, 19 November 2026: used and rechecked in this run; receiving source placement; metadata reconciliation only; no new mathematical check. GPT-6 Astra (OpenAI), ChatGPT Pro mode, effort Pro, S6\_10\_PAPERS continuation, 19 November 2026: used and rechecked in this run; source-row intake; metadata reconciliation only; no new mathematical check. GPT-6 Astra (OpenAI), ChatGPT Pro mode, effort Pro, S6\_10\_PAPERS continuation, 19 November 2026: used and rechecked in this run; explicitly scoped metadata reconciliation; metadata reconciliation only; no new mathematical check. \emph{Pre-continuation recorded proof scope.} Historical event: GPT-6 Astra (OpenAI), ChatGPT Pro mode, effort Pro, S6\_06\_BLIND, 16--18 November 2026: targeted proof check, recorded producing-run re-derivation. Also checked in this writing by GPT-6 Astra (OpenAI), ChatGPT Pro mode, effort Pro, S6\_10\_PAPERS continuation, 19 November 2026: computation or formal certificate for the explicitly listed algebraic substeps in the event ledger; full continuation rederivation and renewed credit comparison remain pending.}{\cref{night:P3K.111}; the original source, point line, roots, characters, reference, selected fan and scalar range stated here}
Retain the original marked source
\[
 y^2=x^3+r_{\rm src}x^2+tx+t^2,\qquad G=(0,t),\qquad
 M_2=\mathcal O([2G]-[O]),\qquad s=\tfrac14-r_{\rm src},
\]
the chronological translation by \(-3G\), and the selected node
\(q=t-t_c(s)\) tending to \(t=0\) at \(s=0\).  The parameter
is one positive original \(s\)-circle, and the original scalar is
fixed in the common strict contraction range of the compared cap
and the cubic reference.  Retain every good cap
\[
 m=3l,\quad l\ge1,\quad n\in\{1,-1\},\quad
 d_0\in\mathbb Z/m,
\]
with its chosen scalar root, finite character and complete annular
reference, and every prescribed selected lattice-basis fan.

On the actual initial native fan \(T_{-1}\) of P3K.111, the
completed physical Jin pair, native free coordinates and return are
\[
 (k_M,\ell_M)=(71,-4),\qquad(\mu,\nu)=(-1,72),\qquad
 [\Sigma]=12\pmod{28}.
\]
For every signed vertical fan \(T_j=T(e_1+je_2,e_2)\), the
whole characteristic and native pairs are
\[
 (k_j,\ell_j)=(11-6j^3+24j^2-30j,-4),
\]
\[
 (\mu_j,\nu_j)=(-1,12-6j^3+24j^2-30j),\qquad
 S_j=2-j^3+4j^2-5j.
\]
The return is \(S_j\bmod28\).  Its vertical image is exactly
\[
 \{0,2,8,10,12,14,16,22,24,26\}.
\]
The full coefficient, native-product and prism algorithm in
AFM2.1--5 below computes the actual return on every selected
lattice-basis fan. Its complete image is
\(\Theta_7=\mathbb Z/28\), witnessed by the 28 full paths and
integer sums reproduced below.  Every path starts with the actual
\(T_{-1}\to T_0\) native arrow with pair \((15,-10)\), hence
with the proved \(T_0\) value \(2\).  The selected source/point
reflection at \(w=+1\) is the full \((Q/Z,V)\) map; its fan is
\(JP\) and its vertical value is \(S_{1-j}\).

The proof retains the original source, both point/O modules, full
theta and Miller coefficients, finite-root and character maps,
both norm signs, all annular integers, and the original closing.
The integer comparison precedes reduction modulo 28.
\end{theorem}
\begin{proof}[Supplied source argument; continuation review in progress]

% Expanded verified cp36 input: notes/part3_e6_selected_i3_m2_continuation_proof.tex
\par\noindent
The entire signed native proof VM2.1--7 is retained first.
Its source-stage statements that the absolute seed and scalar-reference
attachment were then open record that stage faithfully. P3K.111
subsequently supplies the complete actual initial attachment and pair;
IVM2.1--7 below evaluates the whole integer comparison from that pair.
IGC.1--5 gives every signed good-cap map, and AFM2.1--5 gives the
complete all-fan construction. No proof paragraph or mathematical
display from these four sources is omitted.


% Expanded verified cp36 input: notes/part3_e6_selected_i3_m2_continuation_vm2_proof.tex
Retain the original SI2 source \(y^2=x^3+rx^2+tx+t^2\), \(G=(0,t)\),
\(M_2=\mathcal O(2G-O)\), chronological translation \(-3G\), the selected branch
\(t_c(s)\) tending to zero, \(s=1/4-r\), cubic good cap with norm one,
its prescribed finite character, and its complete annular reference.
The scalar here is the fixed original \(\lambda_0=r_0^3\). The parameter
is one positive original \(s\) circle. All arrows below are on this
actual fixed-scalar family; the independently filled scalar-reference
disc and the coupled \((1,4)\) phase square retain their own endpoints.

For the periodic fan \(T_j=T(e_1+je_2,e_2)\) in the original point-coordinate
marking, the two free native pairs are

\[
\begin{array}{c|cc}
 &m&n\\\hline
 C_1&-(3j-5)(j-1)&(3j-5)(j-2)\\
 C_2&(3j-2)(j-2)&-(3j-2)(j-1).
\end{array}                                                    \tag{VM2.1}
\]

Their actual signed Cox arrows give

\[
 S_{j+1}-S_j=-(3j-2)(j-1),\qquad
 S_{j-1}-S_j=(3j-5)(j-2),                                      \tag{VM2.2}
\]
\[
                 S_j-S_0=-j^3+4j^2-5j\quad\hbox{in }\Theta_7.
                                                                  \tag{VM2.3}
\]

Here \(S_j\) is the actual return of the indicated already constructed
family, in the presentation \(\Theta_7=\mathbb Z/28\text{ with }\mu(\Sigma_a)=a/28\).
These equalities evaluate differences of its actual returns. This
note does not assign the still unevaluated constant \(S_0\). Its complete
source, native products and signs are proved below.

\par

\medskip\noindent\textbf{VM2.1. The original full point periods and identity source closing}

\par\nobreak

Write \(V=V_{\rm pt}\) and retain the full SI2 coefficient

\[
 \eta_2=a_T\Theta'_Q(1)\frac{\Theta_Q(P^2)}{\Theta_Q(P)^2},\qquad
 Q=P^2R,\quad q=Rc_q,\quad t=Pu_t.
\]

The actual periods on a uniform selected \(q\)-disc over the source
circle are

\[
 h:(Z,V,q)\longmapsto(\kappa qZ,\alpha V,q),\quad
 B:(Z,V,q)\longmapsto(gZ,\eta_*qV,q),                            \tag{VM2.4}
\]
\[
 \kappa=P^2/c_q,\quad \alpha=P^2,\quad g=P^{-3},\quad
 \eta_*=\lambda_0\eta_2^6t^{-3}.
                                                                  \tag{VM2.5}
\]

Every function is the original analytic function of \(q\) and the source
parameter. In particular the sixth power includes the theta derivative
and its actual constant value; the point/O modules are SI2.18-\/-20.
The original positive source meridian has \(P\)-degree one. The complete
units \(c_q\),\(u_t\),\(\eta_2\) extend over the small source disc with nonzero values.
Their based logarithms have winding zero. The winding matrix of
VM2.4, in the written \(h\),\(B\) order, is therefore

\[
                         \begin{pmatrix}2&-3\\2&-3\end{pmatrix}.
                                                                  \tag{VM2.6}
\]

The full finite cap is SI2.22-\/-28. Its literal root source is
\(Y^2=X^3+rz^2X^2+zX+1\) with \(t=z^{-3}\). Its marked sections, two graph
corrections \(E_P-E_O\), the small division point \(R'_0\), normalized chord \(F\),
invariant cube root of \(B'_{\rm sel}/F^3\), and chosen finite factor \(\zeta_3^d\)
all extend jointly in \(s\) and z. The displayed inverse modules and
the preserved refinement ideals give their extending inverses.
The selected multiplier \(h_\infty=1-t_c(s)z^3\) is the same full
function throughout. At the fixed scalar, \(r_\lambda=r_0\) is constant.
The actual sector maps SI2.28 use fixed contours for their reference
Cauchy operators and the same integer cochain. Consequently those
source, line, cap, norm-root and sector maps are single-valued in \(s\)
on the selected circle, and the original source/cap closing is the
identity. The Tate coordinates use the same marked Weierstrass
scale and point coefficient, so they are periodic as well. Thus
the following products have this literal original identity closing;
no scalar-period reindex is introduced into that source loop.

\par

\medskip\noindent\textbf{VM2.2. Every signed character and period reset}

\par\nobreak

In fan \(T_j\) use

\[
 Z_j=Z,\quad V_j=Z^{-j}V,\quad q_j=q,\qquad h_j=hB^j,
 \quad B_j=B.                                                 \tag{VM2.7}
\]

The inverse character arrow is \(V=Z_j^jV_j\) and the inverse generator
arrow is \(h=h_jB_j^{-j}\). These are identities on the common torus.
For each fan the affine extension is defined by its actual primitive
character semigroups, so VM2.7 specifies the family on that fan; it
is not an assertion that the character changes extend across the
old fan\textquotesingle s vertices. The signed adjacent extending maps are given
in VM2.6 below.

Substitution in both original periods gives the complete functions

\[
 \kappa_j=\kappa g^j,\qquad
 \alpha_j=\alpha\eta_*^j\kappa^{-j}g^{-j^2},\qquad
 \eta_j=\eta_*g^{-j},                                       \tag{VM2.8}
\]
\[
 h_j:(Z_j,V_j)\mapsto(\kappa_jqZ_j,\alpha_jV_j),\qquad
 B_j:(Z_j,V_j)\mapsto(gZ_j,\eta_jqV_j).                        \tag{VM2.9}
\]

For example, the line factor of \(hB^j\) is \(\alpha(\eta_*q)^j\), and
the \(Z\) factor is \(\kappa qg^j\); its minus-\(j\) power gives exactly
\(\alpha_j\). This proves the cancellation of \(q^j\) before any winding
is taken. All functions and their reciprocal powers remain in
VM2.8. The exact winding matrix is

\[
 \mathsf B_j=
 \begin{pmatrix}2-3j&-3\\3j^2-5j+2&3j-3\end{pmatrix},
 \qquad\operatorname{tr}\mathsf B_j=-1.                      \tag{VM2.10}
\]

The two inverse integral changes preserve the original real period
lattice and its finite period-cell quotient. For each fixed finite
word of these fan changes the original strict contraction domain
supplies one uniform smaller selected tube. The exterior uses the
inverse VM2.7 map back to the original source, point line and cap.
It therefore has the unchanged original cap, finite character,
integer cochain and endpoint map. This constructs the actual family
and its original closing for every indicated signed integer \(j\).

\par

\medskip\noindent\textbf{VM2.3. The full vertex charts and ordinary double product}

\par\nobreak

The two regular toric vertex triples \((a,b,c)\) and \((u,v,w)\) have height
\(q=abc=uvw\). Their complete overlap rows are

\[
\begin{array}{c|ccc}
 &u&v&w\\\hline
 C_0&\kappa_j\alpha_jbc&\alpha_j^{-1}ab&\kappa_j^{-1}b^{-1}\\
 C_1&g\eta_jbc&\eta_j^{-1}c^{-1}&g^{-1}ac\\
 C_2&a^{-1}&ab&ac.
\end{array}                                                   \tag{VM2.11}
\]

The local torus dictionary on the first vertex is
\(Z_j=b,\ V_j=c,\ q=abc\), with inverse \(a=q/(Z_jV_j)\). Its physical
weights are \((-1,0,+1)\). The rows follow by the literal \(h_j\),\(B_j\)
period translations in VM2.9. The coordinate products equal \(q\).
Each row has the inverse obtained by recovering that same \(q\) and
dividing its complete functions there. Thus all \(q\) derivatives and
both native vertex collars are retained.

The \(C_2\) native charts are \((z,n_0,n_1)=(a,b,c)\) and \((u,v,w)\).
The \(C_1\) native charts are \((c,a,b)\) and
\((\eta_jv,gw,(g\eta_j)^{-1}u)\). Direct substitution gives \(z'=1/z\)
and \(n'_i=zn_i\) in both cases. They define periodic single native
products. Their doubles are their literal pullbacks under the
degree-two source-circle map, with the original identity source
closing and all cap-sector arrows pulled back at the same time.

Use \(s=s_*e^{2i\theta},\ 0\le\theta\le2\pi\), on that double. At its
second vertex take \((u,\eta_j(q,s)v,\kappa_j(q,s)w)\), compared with
the same arrow at \(s_*\). The actual new height is
\(H_s(q)=q\eta_j(q,s)\kappa_j(q,s)\). Its derivative at zero is a
unit. Recover \(q\) by \(H_s^{-1}\) and then divide by these exact functions
at the recovered \(q\). Compactness gives a uniform small disc for
this inverse. In the ordered normal pairs \((a,c),(a,b),(b,c)\),
the three full diagonal matrices are

\[
 D_0=\operatorname{diag}(\eta_j/\alpha_j,\kappa_j\alpha_j),
 \quad D_1=\operatorname{diag}(\kappa_j/g,g\eta_j),\quad
 D_2=\operatorname{diag}(\eta_j,\kappa_j).                      \tag{VM2.12}
\]

Their determinant winding is \(2t=-2\). Their first-entry windings
on this literal double are

\[
 l_0=-6j^2+16j-10,\quad l_1=10-6j,\quad l_2=6j-6,
 \qquad t=-1.                                                \tag{VM2.13}
\]

Divide each entry by its reference value at the same \(q\) and by its
written angular factor. The resulting winding-zero loop has its
unique based periodic logarithm on that entire \(q\)-disc. Retain both
full logarithms, including all their \(q\) derivatives. The actual
normal interpolation in \(Q=n_0+j_{\rm quat}\overline{n_1}\) uses

\[
 H_d(x,\theta)=
 (\sin\vartheta(x)e^{id\theta}+j_{\rm quat}\cos\vartheta(x))
 (\sin\vartheta(x)+j_{\rm quat}\cos\vartheta(x))^{-1},
 \qquad Q\longmapsto H_tQH_{l_i-t},                          \tag{VM2.14}
\]

where \(\vartheta\) is endpoint-flat with values zero and \(\pi/2\).
Compose with the two full based-log exponentials using a radial
cutoff. The two radial endpoints are identity and exactly
\(D_i(q,s_*)^{-1}D_i(q,s)\). The \(q\)-dependent map on a bounded curve
chart is \(F(n)=An+G(n)\), where A is its invertible value at \(q=0\).
Since \(q=zn_0n_1\), \(G=O(|n|^3)\) with its uniform first derivative.
Choose the compact tube so that \(\|A^{-1}DG\|<1/2\). The iteration
\(n\longmapsto A^{-1}n'-A^{-1}G(n)\) gives the unique smooth inverse. Its
source height is \(q=zn_0n_1\) at that recovered \(n\). At the second
vertex this is the \(H_s\) inverse already written. This proves the
whole inverse through the quaternion mixing.

The endpoint equalities glue to the full vertex maps. The normal
rescaling \(R_\tau(z,n)=(z,\tau n)\) and the cutoff flow of
\(R_\tau^{-1}FR_\tau\) restore the full nonlinear germs. The preceding
inverse bound is uniform on the compact collars, and the flow
fixes the protected endpoint collars. The three open curve tubes
are disjoint. Thus this is the actual based ordinary plumbing
product with all \(q\)-normal jets, original source/cap closing and
complete line coefficients, for this literal double.

\par

\medskip\noindent\textbf{VM2.4. The action and stable-derivative prisms}

\par\nobreak

The physical tangent weights and ordered normal differences are
\(w=(0,1,-1)\text{ and }d=(-2,-1,-1)\), by VM2.11. Under VM2.14 the action
on \(Q\) is

\[
 Q\longmapsto
 (H_t e^{-iw_i\phi/2}H_t^{-1})Q
 (H_{l_i-t}^{-1}e^{id_i\phi/2}H_{l_i-t}).                    \tag{VM2.15}
\]

The retained logarithms and full normal-rescaling flow give its
based homotopy to the original action. The corresponding trace,
with \(U=\cos x+j_{\rm quat}\sin x\) and \(E_i=e^{iw_i\phi/2}\), is exactly

\[
 E_i^{-1}UE_iH_tE_i^{-1}H_t^{-1}U^{-1}
 H_{l_i-t}^{-1}e^{id_i\phi/2}H_{l_i-t}.                       \tag{VM2.16}
\]

For the explicit \(H_d\), differentiation on its oriented path-parameter
sphere gives \(\text{degrees }-d\text{ and }d\) for \(H_diH_d^{-1}\) and \(H_d^{-1}iH_d\).
The positive quaternion volume evaluates an axis of degree \(e\) and
angle \(c\phi/2\) as \(c\) \(e\). In the Maurer product identity the radial
endpoint terms are constant central signs and the \(\theta\) endpoint
terms cancel in these periodic charts. The resulting normal and
stable-derivative degrees are therefore, with their actual ordered
Spin lifts,

\[
 I_i=t w_i+d_i(l_i-t),\qquad W_i=t w_i-d_i(l_i-t).             \tag{VM2.17}
\]

At \(\phi=2\pi\) the base Spin sign \((-1)^{w_i}\) cancels the normal sign.
Close the traces by generic normal sections at the two vertices,
paths in their actual links, and the parameter-independent orbit
caps \(D_\theta^2\times S_\phi^1\). The vertex maps have two-dimensional
image. The coordinate-to-ambient stable-frame map extends over
each annulus times \(D^4\), whose \(H^3\) vanishes. Hence these specified
closures add zero to the three-degrees. These are precisely the
same oriented prism domains and maps as IF3.7-\/-11, with VM2.12\textquotesingle s
complete coefficients and VM2.13\textquotesingle s newly derived windings.

Substitution gives

\[
 I=(12j^2-32j+18,\ 6j-12,\ 6-6j),
\]
\[
 W=(-12j^2+32j-18,\ 10-6j,\ 6j-4).                          \tag{VM2.18}
\]

Let \(\mathrm V\) be the actual variation matrix and \(\delta\) the actual
ambient stable derivative in this ordinary double product. The
oriented orbit-prism identity \(\partial\Delta=-\Delta\partial\),
applied to the retained relative discs, gives

\[
 \mathrm V(w_i e_i-w_0e_0)=I_i e_i-I_0e_0,\qquad
 w_i\delta_i-w_0\delta_0=W_i-W_0.                           \tag{VM2.19}
\]

As \(w_0=0\), its selected free entries are determined exactly:

\[
\begin{array}{c|cc}
 &\mathrm V_{ii}&\delta_i\\\hline
 C_1&6j-12&12j^2-38j+28\\
 C_2&6j-6&-12j^2+26j-14.
\end{array}                                                   \tag{VM2.20}
\]

The map to the held native quaternion has its literal derivative
\(dc=-u^{-1}(du)u^{-1}q_H+u^{-1}dq_H\). Remove the first term by its
invertible triangular shear homotopy. Its oriented normal-meridian
third class is +1 in the retained K23 product. The relative disc
orientation is \(D_i=-[D_y^3]\). Therefore its ordinary free coordinates
obey \(2n_i=\delta_i-\mathrm V_{ii}\) and \(m_i=-\mathrm V_{ii}-n_i\). This evaluates

\[
\begin{array}{c|cc}
 &m^{ord}&n^{ord}\\\hline
 C_1&-2(3j-2)(j-2)&2(3j-5)(j-2)\\
 C_2&2(3j-5)(j-1)&-2(3j-2)(j-1).
\end{array}                                                   \tag{VM2.21}
\]

These selected coordinates do not require a value for the fixed
\(C_0\) characteristic pair or for the absolute return.

\par

\medskip\noindent\textbf{VM2.5. The actual native product and the literal square}

\par\nobreak

On either native star the complete native-to-ordinary boundary is

\[
 q_H\longmapsto(UH_tU^{-1})q_HH_{l_i-t},                     \tag{VM2.22}
\]

with the full logarithms and nonlinear maps of VM2.3. Its left
third degree is \(2t=-2\): the boundary Maurer integral of the explicit
VM2.14 path over the Hopf-oriented sphere times \(\theta\) is twice
its angular winding. The right factor depends only on radial and
parameter variables and has degree zero. The inverse-cylinder
extension over \(D^3\times I\) consequently has endpoint \((2,0)\). Its
same factors extend over \(D^3\times D^4\), preserving the cap boundary
and adding no disc-supported class. Thus the ordinary return is
the native return followed by \(L^2\). Subtract two from each first
entry in VM2.21 to obtain the native doubled pairs

\[
\begin{array}{c|cc}
 C_1&-2(3j-5)(j-1)&2(3j-5)(j-2)\\
 C_2&2(3j-2)(j-2)&-2(3j-2)(j-1).
\end{array}                                                   \tag{VM2.23}
\]

Choose a connection compatible with the actual single native
product VM2.11. Its literal pullback to the double circle has
endpoint \(f^2\), with the same identity source closing and full cap
maps on both sheets. The ordinary connection, changed back by
VM2.22, has that identical native boundary connection. Their affine
interpolation with fixed boundary gives an isotopy of relative
endpoints. The integral free-coordinate homomorphism sends \(f^2\)
to twice the coordinates of \(f\). Injectivity of multiplication by
two in \(\mathbb Z\) now gives exactly VM2.1. The central finite-group
coordinate is retained on the single actual family throughout.

\par

\medskip\noindent\textbf{VM2.6. Both signed full Cox arrows and the original exterior}

\par\nobreak

For a vertex arrow \(((AD)^{-1}a,Ab,Dc)\), its full native \(C_2\) Cox
lift is

\[
 (a_0,a_1,c_0,c_1)\longmapsto
 (a_0,(AD)^{-1}a_1,A c_0,D c_1),\qquad
 q=a_0a_1c_0c_1.                                             \tag{VM2.24}
\]

In \(C_1\) order \((c,a,b)\) its normal factors are \(((AD)^{-1},A)\), so
the same full lift is
\[
(a_0,Da_1,(AD)^{-1}c_0,Ac_1).
\] The fixed-\(C_0\)
order \((b,a,c)\) has normal factors \(((AD)^{-1},D)\). These are literal
coordinate permutations of VM2.11. All maps preserve \(q\) and have
their reciprocal inverses at the same source point. Their four
weights sum to zero. They commute with the auxiliary Cox action
and the physical circle. Transporting their Hermitian norms gives
the same positive radial normalization on the common conifold.

For arbitrary angular weights \(A_0,A_1,C_0,C_1\) retain the full identity
of IF3.17: with \(u=a_0+j_{\rm quat}a_1\), \(v=c_0+j_{\rm quat}\overline{c_1}\),
\(q_H=uv\), \(h=uiu^{-1}\) and \(k=v^{-1}iv\), the four coefficient phases
give \(h'=aha^{-1}\), \(k'=b^{-1}kb\) and
\(q_H'=a e^{hS\theta/2}q_Hb=aq_He^{kS\theta/2}b\),
where \(S\) is their sum. Hence \((q_H')^{-1}h'q_H'=k'\). This is a
literal product identity; it retains each half-phase and its paired
central endpoint. The positive moduli and all based unit logarithms
in VM2.24 preserve that same identity by their stated invertible
paths. The common normal-rescaling flow carries all higher germs
on both rulings simultaneously.

The actual link and exterior maps are consequently

\[
 \tau(h,q_H)=(q_H^{-1}h q_H,q_H),\quad
 T(y,q_H)=(q_H^{-1}yq_H,q_H),\quad
 T^{-1}(y,q_H)=(q_Hyq_H^{-1},q_H).                            \tag{VM2.25}
\]

They extend smoothly at \(y=0\), preserve the product orientation and
have the displayed inverse. The two Hopf inverses choose a unit
\(u\) or \(v\) and recover \(v=u^{-1}q_H\) or \(u=q_Hv^{-1}\); changing the unit
representative is exactly the auxiliary circle. Thus these maps
carry the same held quaternion and complete native Hopf boundary,
not only its circle weight. Period translates use VM2.9 with every
unit. Outside the small selected neighbourhood the original source,
Cartier line, cubic cap and full annular closing are the identity
comparison of VM2.1-\/-2.

The \(C_2\) change is \(T_j\longmapsto T_{j+1}\), with full reset

\[
 Z_{j+1}=Z_j,\quad V_{j+1}=Z_j^{-1}V_j,\quad
 h_{j+1}=h_jB_j,\quad B_{j+1}=B_j,                           \tag{VM2.26}
\]
\[
 \kappa_{j+1}=\kappa_jg,\quad
 \alpha_{j+1}=\alpha_j\eta_j\kappa_j^{-1}g^{-1},\quad
 \eta_{j+1}=\eta_jg^{-1}.                                   \tag{VM2.27}
\]

The reverse \(C_1\) change is \(T_j\longmapsto T_{j-1}\), with

\[
 V_{j-1}=Z_jV_j,\quad h_{j-1}=h_jB_j^{-1},\quad
 \kappa_{j-1}=\kappa_jg^{-1},\quad
 \alpha_{j-1}=\alpha_j\eta_j^{-1}\kappa_jg^{-1},\quad
 \eta_{j-1}=\eta_jg.                                       \tag{VM2.28}
\]

The actual two-ruling coordinate maps are
\((n_0/n_1,zn_1,n_1)\) and \((n_1/n_0,zn_0,n_0)\), whose products are
\(q=zn_0n_1\). Their shared rank-one matrix has rows
\((n_0,n_1)\) and \((zn_0,zn_1)\). The \(C_1\) normal interchange is the constant
unitary path with eigenspaces \((1,1),(1,-1)\) and eigenvalues
\(1,e^{\pi it}\); it preserves the real normal orientation and the
scalar transition. Thus these are full local resolution maps and
their common-link continuation, including the actual inverses on
their stated localizations.

For the fan labels put \(u=e_1+je_2,\ v=e_2\). The \(C_2\) operation replaces
the diagonal between \(n+u\) and \(n+v\) by that between \(n\) and \(n+u+v\);
the two resulting triangles are exactly \(T(u+v,v)\), one directly
and one after \(n\longmapsto n+v\). The inverse adjacent-cell operation gives
\(T(u-v,v)\). These literal translates produce VM2.26 and VM2.28,
including their period-generator changes. This proves both signed
directions in the original lattice.

The actual K23 generators \(L\),\(R\) are the left and right degree-one
quaternion maps on \((D^3,S^2)\) and extend over \(D^3\times D^4\). The
disc-supported generator \(C\) has the positive Eells-\/-Kuiper convention
fixed by the same oriented commutator plumbing and physical line
map in K68. The literal full-link conjugations are
\(TLT^{-1}=R\) and \(TRT^{-1}=L^{-1}R^2C\). Hence the common exterior
arrow sends \((m,n,a)\text{ to }(-n,m+2n,a+n)\). The fixed anti-quaternion
reset of the first complex-normal convention interchanges the two
free entries and retains the central entry. Its full result is

\[
                        (m,n,a)\longmapsto(m+2n,-n,a+n).
                                                                  \tag{VM2.29}
\]

Substitution of VM2.1 gives the \(C_1\) pair at \(j+1\) from the \(C_2\) pair
at \(j\), and the \(C_2\) pair at \(j-1\) from the \(C_1\) pair at \(j\). Thus the
complete free pairs are checked on both ends of both directed
arrows before their central changes are read. Since \(L\) and \(R\) extend
over the seven-disc, the central entry is the actual absolute return.
Its changes are precisely VM2.2. Summing the polynomial differences,
with the same reverse arrows for negative \(j\), gives VM2.3.

\par

\medskip\noindent\textbf{VM2.7. The two actual markings occurring in the current atlas}

\par\nobreak

The protected collision star NS has \(c_0=ZV_{\rm pt}\), while the RB moving
reference has \(c_R=Z^{-1}V_{\rm pt}\). From the literal first-vertex dictionary
in VM2.11 and the period-aligned other vertex, these are respectively
\(T_{-1}\) and \(T_1\). Their positive-diagonal ordered column matrices are

\[
 P_N=\begin{pmatrix}1&-1\\-1&2\end{pmatrix},\qquad
 P_R=\begin{pmatrix}1&-1\\1&0\end{pmatrix}.
\]

Put \(H=\left(\begin{smallmatrix}0&1\\-1&2\end{smallmatrix}\right)\) as a two-by-two matrix. Then
\(P_RH=\left(\begin{smallmatrix}1&-1\\0&1\end{smallmatrix}\right)\) and \(P_RH^2=P_N\), by direct multiplication.
The complete two reverse-\(C_1\) continuation is VM2.28 followed by
VM2.28, with its original exterior, cap and full Cox maps at both
steps. Its common torus coefficient is exactly \(c_0=Z^2c_R\) and
\(a_0=Z^{-2}a_R\). At the vertices it uses the two actual ruling changes
VM2.24-\/-29; the pole of the torus inverse is retained as the reason
to use those charts. The fixed-source cap over its independent
scalar disc still requires its own full continuation; these fixed-
scalar source-loop arrows do not substitute its endpoint word.

For the present fixed-scalar source loop, the two reverse changes
have complete free inputs

\[
 (m_{C1},n_{C1})\big|_{j=1}=(0,2),\qquad
 (m_{C1},n_{C1})\big|_{j=0}=(-5,10).
\]

Their central increments are +2 and +10. Therefore the actual
fixed-scalar returns in these two marked fans satisfy

\[
                              S_{-1}-S_1=12\quad(\bmod 28).
                                                                  \tag{VM2.30}
\]

This is a proved return difference with all source and cap arrows
specified. It leaves the absolute seed and the independent scalar-
reference attachment to the ongoing compact characteristic ledger.

% End expanded verified cp36 input

\par

% Expanded verified cp36 input: notes/part3_e6_selected_i3_m2_continuation_ivm2_proof.tex
Complete proof. Retain SI2\textquotesingle s original source
\(y^2=x^3+r_{\rm src}x^2+tx+t^2\), \(G=(0,t)\),
\(M_2=\mathcal O([2G]-[O])\), chronological translation
\(-3G\), \(s=1/4-r_{\rm src}\), the selected branch
\(q=t-t_c(s)\), good cap \((m,n)=(3,1)\), its finite
character and complete annular reference. The original scalar
is fixed. The parameter is one positive original \(s\)-circle.
For every signed \(j\), the fan is
\(T_j=T(e_1+je_2,e_2)\) in the original point marking.

The actual ASI2 completed Jin pair on \(T_{-1}\) is
\((71,-4)\), with native coordinates \((-1,72)\).
The comparison below evaluates the integers on every \(T_j\):
\[
 k_j=11-6j^3+24j^2-30j,\qquad \ell_j=-4,
\]
\[
 \mu_j=-1,\quad \nu_j=12-6j^3+24j^2-30j,\qquad
 S_j=\nu_j/6=2-j^3+4j^2-5j .
 \tag{IVM2.1}
\]
This evaluation uses the physical line and the parameter
reversal of K68 on the complete fixed-\(C_0\) quotient.
The integer changes are calculated before reduction modulo
28. The local source/node comparison is given in full here;
the VZ, BW and IVP constructions are used only through
their explicit maps repeated below with the present functions.

\par

\medskip\noindent\textbf{IVM2.1. Full source map, inverse and both period squares}

\par\nobreak

On a smaller uniform selected \(q\)-disc over the original
source circle, put
\[
 Q=q\kappa,\quad \kappa=P^2/c_q,\quad
 \alpha=P^2,\quad g=P^{-3},\quad
 \eta=\lambda_{\rm base}\eta_2^6t^{-3}.
 \tag{IVM2.2}
\]
Here \(\eta_2=a_T\Theta'_Q(1)\Theta_Q(P^2)/\Theta_Q(P)^2\),
including its entire theta derivative and the original
Weierstrass differential. The source periods are
\(h=(q\kappa Z,\alpha V)\), \(B=(gZ,q\eta V)\).
Their source-circle windings are \((2,2,-3,-3)\)
in the displayed order. All four full functions and
their derivatives are retained.

For a fixed signed \(j\), solve
\[
                  Q(q',s)=Q(q,s)g(q,s)^j.
 \tag{IVM2.3}
\]
Since \(\partial_qQ(0,s)=\kappa(0,s)\ne0\), the holomorphic
inverse of \(Q\) gives \(q'=a_j(q,s)q\) with
\(a_j(0,s)=g(0,s)^j\). Compactness of the source circle
gives one radius for this \(j\), and the inverse on the
actual image collar. Uniqueness makes both maps periodic
in \(s\). This uses that image collar, with no change
of its marked radial boundary.

Set \(c(q,s)=g(q',s)/g(q,s)\). It equals one at \(q=0\).
Its holomorphic logarithm zero there is single-valued
in \(s\). On the suspension cover the source map is
\[
 (q,Z,\gamma_j)\longmapsto
       (q',c(q,s)^{\gamma_j}Z,\gamma_j).
 \tag{IVM2.4}
\]
The \(h_j=hB^j\) square is exactly IVM2.3, and the \(B\)
square is
\[
 c^{\gamma_j+1}g(q,s)=g(q',s)c^{\gamma_j}.
 \tag{IVM2.5}
\]
The inverse first uses the inverse \(q\)-map and then
the reciprocal \(c\) coefficient. At the node the two
branch multipliers are reciprocal units, with \(c-1=O(q)\);
they extend over the actual \(xy=Q\) source and preserve
both marked branches.

The character and period maps, with their inverses, are
\[
 Z_j=Z,\quad V_j=Z^{-j}V,\quad h_j=hB^j,\quad B_j=B,
\]
\[
 \kappa_j=\kappa g^j,\qquad
 \alpha_j=\alpha\eta^j\kappa^{-j}g^{-j^2},\qquad
 \eta_j=\eta g^{-j}.
 \tag{IVM2.6}
\]
Recover \(V=Z_j^jV_j\), \(h=h_jB_j^{-j}\). Substitution
cancels the entire \(q^j\) in \(hB^j\), and preserves
the exact point/O and Miller factors of SI2.
The exterior uses the recovered original source and
line. The finite source cap, its root, finite character
and full annular cochain are unchanged.

\par

\medskip\noindent\textbf{IVM2.2. Fixed graph and actual normal map}

\par\nobreak

Let \(\gamma_j\) be the real coefficient of the contracting
\(B\)-period. For the actual positive logarithms of
\((Q_j,\alpha_j,g,q\eta_j)\), its exact torus formula is
\[
 \gamma_j=
 \frac{L_Q\log|V_j|-L_\alpha\log|Z_j|}
      {L_QL_\eta-L_\alpha L_g}.
 \tag{IVM2.7}
\]
The denominator has leading term \((\log|q|)^2\) and
is nonzero on the chosen collar. The actual \(h_j\)
adds \((L_Q,L_\alpha)\), hence has lift zero in
\(\gamma_j\), while \(B\) has lift one.
The outward F9 direction is \(-\gamma_j\).
In its two conifold charts the moment equations are
\[
 2b=r-|y|^2/r,\qquad 2b'=|x|^2r-1/r,\qquad
 \frac{db'}{db}=
 \frac{|x|^2+r^{-2}}{1+|y|^2r^{-2}}>0.
 \tag{IVM2.8}
\]
On the exceptional interval use
\[
 \gamma_j=\tfrac12-\epsilon(r-1)/(r+1),\quad
 0<\epsilon<\tfrac14.
 \tag{IVM2.9}
\]
Its two endpoint lifts are \(1/2+\epsilon,1/2-\epsilon\),
with derivatives \(-4\epsilon\) in both endpoint moment
coordinates. Extend on the two half-lines by the actual
exterior coefficient, using the relative positive-vector
field and its inverse; repeat by \(B\)-translation.
The full positive unit norms remain in this construction.

IVM2.4 preserves the source normalization and its two
marked points. Interpolate the transported graph height
and the reference height in this common short arc.
The endpoint order remains strict. Away from the two
node preimages the normalization is injective, so this
is an embedded-sphere isotopy. Extend its velocity to
a tube by a cutoff, periodically in \(s\), and integrate.
Backward integration is the inverse. This constructs
the graph comparison relative to the original exterior.

In the fixed-star charts \(q=ZAB\), the normal map is
\[
                  (Z,A,B)\longmapsto(Z,a_j(q,s)A,B).
 \tag{IVM2.10}
\]
It intertwines the second chart \(A_1=ZA,B_1=ZB\);
its inverse first recovers \(q\) by IVM2.3.
The \(c^{\gamma_j}=1+O(q)\) correction has no extra
normal linear term. The new-to-old spinor-line
Cartier map multiplies by \(Z^j\): at \(x/y=Z\)
it sends
\[
        (x^{-j-1}y^j,x^{-j}y^{j+1})\longmapsto(1/x,y).
 \tag{IVM2.11}
\]
The determinant frames and the Hom/adjoint normal
arrows therefore agree literally. The remaining
normal coefficient is the entire \(g(0,s)^j\),
with winding \(-3j\). The graph-normal derivative
is the positive triangular derivative of the
constructed isotopy. The base sphere map extends
over its three-ball by that velocity and a collar
cutoff. Its full adjoint Hopf map and physical
principal-circle lift remain part of the comparison.

\par

\medskip\noindent\textbf{IVM2.3. Full theta-cut line morphism and outer primitive}

\par\nobreak

Write \((u,v,d,\gamma)\) for the real coefficients of the
source vanishing period, \(h\), physical \(\delta\), and
\(B\). The character map gives
\[
 u_j=u,\quad v_j=v,\quad
 \gamma_j=\gamma-jv,\quad d_j=d-ju.
 \tag{IVM2.12}
\]
The positive \(q\)-deck is
\((u,d)\mapsto(u-v,d-\gamma)\). The original positive
\(s\)-deck is
\[
       u\mapsto u-2v+3\gamma,\qquad
       d\mapsto d-2v+3\gamma.
 \tag{IVM2.13}
\]
Conjugating by IVM2.12 gives the full new integral decks
\[
 q:\ (u,d_j)\mapsto(u-v,d_j-\gamma_j),
\]
\[
 s:\quad
 u\mapsto u+(3j-2)v+3\gamma_j,\qquad
 d_j\mapsto d_j-(3j^2-5j+2)v-(3j-3)\gamma_j.
 \tag{IVM2.14}
\]
All zero-winding unit displacements are the derivatives
of the complete based logarithms in IVM2.2 and IVM2.6.
They are retained with these integral terms.

For \(U=M^*\), let \(\alpha_q,\theta\) have period one.
The integral parts of the actual coefficient connections are
\[
 A_0=dd+(\alpha_q-3\theta)d\gamma+2\theta\,dv,
\]
\[
 A_j=dd_j+(\alpha_q+(3j-3)\theta)d\gamma_j
                         +(3j^2-5j+2)\theta\,dv.
 \tag{IVM2.15}
\]
Their difference in original coefficients is
\[
 A_j-A_0=-j[du+(\alpha_q+2\theta)dv-3\theta\,d\gamma]
                      +\hbox{the full unit derivative}.
 \tag{IVM2.16}
\]
The target \(q'\) angle is
\(\alpha_q-3j\theta\) plus its full unit term.
The reference at \(q'\), expressed in the new coefficients,
is
\[
 A_{0,q'}=dd_j+(\alpha_q-(3j+3)\theta)d\gamma_j
                                      +2\theta\,dv.
 \tag{IVM2.17}
\]
Subtracting leaves the exact integral residual
\[
           (3j^2-5j)\theta\,dv+6j\theta\,d\gamma_j.
 \tag{IVM2.18}
\]

Choose nested fixed-sphere tubes \(U_0\Subset U_1\Subset U_2\)
with all surgery and connection changes in \(U_1\).
Use their real lifts to make the graph maps \(b,\gamma\),
equal to \(v,\gamma_j\) outside, constant on \(U_2\).
They extend constantly over the same Hopf cap.
The endpoint-to-initial theta-cut coefficient arrows
\(w\mapsto e^{2\pi ib}w\) and
\(w\mapsto e^{2\pi i\gamma}w\) define actual lines
\(\mathcal T_b,\mathcal T_\gamma\), with curvatures
\(\Xi=d\theta\wedge db,\Psi=d\theta\wedge d\gamma\).
Their squares and product vanish. The full relative line map is
\[
 U_j\simeq U_{0,q'}\otimes
        \mathcal T_b^{\,3j^2-5j}\otimes\mathcal T_\gamma^{\,6j}.
 \tag{IVM2.19}
\]

Here is its prescribed-boundary extension. The \(q\)-monodromy
kills \(u\) primitively, while \(\theta,b,\gamma\) extend over
the piece. The sole remaining boundary character is the \(q\)
angle. At one theta choose a smooth free source section
\(|Z_*|=1,v=\gamma_j=0\), away from the marked sections
and fixed graph. Its actual coefficient \(Z_*^{-j}\),
the \(q'\) map and both tensor frames extend over that
disc, proving zero \(q\)-degree. Choose ratio-line
transitions \(e^{2\pi ir_{il}}\). Remove their zero
integral two-cocycle, solve \(r_{il}=a_l-a_i\) by a
partition, and use \(e^{-2\pi ia_i}e_i\) as frames.
Remove the extending theta, \(b,\gamma\) characters.
The remaining real boundary logarithm extends by
a collar cutoff. This gives exactly the prescribed
outer arrow, including the full unit factors and
IVM2.16; its inner cap arrow is computed below.

\par

\medskip\noindent\textbf{IVM2.4. Tangent collar and the four relative pairings}

\par\nobreak

Set \(\tau=\log Q/(2\pi i)\), \(G=\log g/(2\pi i)\),
\(\Delta G=\log c/(2\pi i)\). The first two are cut
logarithms, and the last is the actual periodic
logarithm zero at \(q=0\). For the covariant coframe
\(\omega=du+\operatorname{Re}\tau\,dv+
\operatorname{Re}G\,d\gamma\), the full source block is
\[
 \omega'=\omega+\operatorname{Re}(\Delta G)(d\gamma-j\,dv),
 \quad dv'=dv,\quad d\gamma'=d\gamma-j\,dv.
 \tag{IVM2.20}
\]
All base derivatives of the full periods are retained
in the complementary off-diagonal block. In order
\((\omega,\gamma,v)\), the source block is unipotent
upper triangular. Its strictly upper triangular
connection and curvature have zero trace-square.
The base map is \(r'=r+\log|a_j|\),
\(\alpha_q'=\alpha_q+\arg(a_j)/(2\pi)\), \(\theta'=\theta\).
It retains winding \(-3j\). Its collar connection
depends on two boundary angles and the interval,
so its curvature-square vanishes. The full derivative
is block triangular; the trace-square is the sum
of these two diagonal blocks. Thus the full tangent
collar has zero Pontryagin transgression. No logarithm
of \(g\) is contracted on the source circle.

With the carried \(U\)-curvature \(F\) and orientation
\(\Omega_+=dr\wedge d\alpha_q\wedge du\wedge dv
\wedge d\gamma\wedge d\theta\), the actual pairings are
\[
 \int F^2\Xi=1,\quad \int F^2\Psi=0,\quad
 \int p_1\Xi=0,\quad \int p_1\Psi=2.
 \tag{IVM2.21}
\]
For the first, freeze theta. The full units have
holomorphic logarithms on the \(q\)-disc. The BW
logarithmic path with the same closing exponent,
its source-node field, real-period boundary field,
F9 height isotopy, Cartier/Hopf arrow and invariant
line lift carry this actual source piece to the
SCN.9 local comparison. Its two outer normal arrows
and shared inner physical gauge extend over the
complete common link and cap. Its separate integer
corrections vanish; its full Euler collar is
\(3\int F^2\Xi-3\). This fixes the pairing to one.
The carried connections retain the theta decks
and make the same calibration on every slice.

For \(F^2\Psi\), take a gamma level outside the
fixed graph\textquotesingle s short arc. It is the actual selected
\(I_1\) source \(E\), disjoint from the surgery and cap.
Its full boundary phase is
\[
 \phi_\alpha(q',s)=2\theta+
  (2\pi)^{-1}\operatorname{Im}
       \operatorname{Log}[\alpha(q',s)e^{-4\pi i\theta}].
 \tag{IVM2.22}
\]
At frozen theta the connection is \(dd+\phi_\alpha\,dv\).
The \(q\)-deck has \(d\)-shift \(-\gamma\); at gamma zero
its \(q\) holonomy is zero. The line has degree zero
on the normalization sphere. Extend the source graph
\(b\) over the vanishing handle by its constant attaching
lift, and extend the connection as
\(dd+\phi_\alpha\,db\). Its full curvature is
\(d\phi_\alpha\wedge db\), with square zero.
The free section disc and the Cech/logarithm map
above retain its exact boundary frame. Terms with
\(d\theta\) vanish when wedged with \(\Psi\).

For \(p_1\Xi\), the regular graph level, chosen outside
\(U_2\), has its relative level cobordism to the original
three-dimensional source level times gamma. Its tangent
connection is source-pulled, so its four-form integral
is zero. The complete collar transgression just proved
is zero. The modified regular level is used through
this cobordism, rather than assigned a product chart.

For \(p_1\Psi\), on the node \(xy=q\) keep
\(\Omega=dx\wedge dy\), \(V=(x,-y)\),
\(H=(\bar y,\bar x)/R^2\), \(R^2=|x|^2+|y|^2\).
Then \(dq(H)=\Omega(V,H)=1\). The Tate horizontals
\[
 H_x=H-\bar y(xR^2)^{-1}V=(0,1/x),\qquad
 H_y=H+\bar x(yR^2)^{-1}V=(1/y,0)
 \tag{IVM2.23}
\]
are the literal determinant-one shear homotopies.
The normalized frame
\(R^{-1}\left(\begin{smallmatrix}x&\bar y\\-y&\bar x
\end{smallmatrix}\right)\) has degree one on the node
sphere, hence relative \(c_2=1\). The polar frame
\((V,qH)\) has determinant \(q\), the primitive fibre
with square zero, so relative \(p_1=-2\).
The marked real-period frame differs by the displayed
shears, a positive matrix and a polar angle-circle
rotation, whose third transgressions vanish by these
explicit paths and factorization. Finally
\(\Psi=d\theta\wedge d\gamma\) reverses the last two
directions of \(\Omega_+\), giving \(+2\).

\par

\medskip\noindent\textbf{IVM2.5. Full Hopf cap with degree \(-3j\)}

\par\nobreak

On the fixed tube put
\[
 u=(1+\mathbf j Z)/\sqrt{1+|Z|^2},\quad h=u i u^{-1},
 \quad Q_N=A+\mathbf j\bar B,\quad q_H=uQ_N.
 \tag{IVM2.24}
\]
Carry the entire initial factor \(g(0,s_*)^j\) into
both normal reference products. For the actual remaining
coefficient write
\[
 \frac{g(0,s)^j}{g(0,s_*)^j}
       =r_j(s)e^{i\alpha(s)},\qquad
 r_j(s)=\left|\frac{g(0,s)^j}{g(0,s_*)^j}\right|>0.
 \tag{IVM2.24a}
\]
The full positive factor is carried by the held-Cox path
\[
 \mathcal R_u=T_Z\operatorname{diag}(r_j(s)^u,1)T_Z^{-1},
 \quad 0\le u\le1,\qquad
 \mathcal R_u^{-1}
     =T_Z\operatorname{diag}(r_j(s)^{-u},1)T_Z^{-1}.
 \tag{IVM2.24b}
\]
Both source-cut endpoints have \(r_j=1\). On the
normal Hopf sphere this is
\(q_H\mapsto\mathcal R_uq_H/\|\mathcal R_uq_H\|\),
with its normalized reciprocal. Its three-ball
extension is explicit: the boundary Hermitian
projector \(P_Z=T_Z\operatorname{diag}(1,0)T_Z^{-1}\)
extends as the trace-one Hermitian ball coordinate
\(Y\); use \(\exp(u\chi\,\log r_j\,Y)\), with a
radial cutoff \(\chi=1\) near the boundary and
\(\chi=0\) near the centre. The reciprocal has
the negative exponent. The centre is the identity
and every matrix is positive invertible. Thus
this is an actual extending normal-bundle map.
The phase equality now concerns the ratio in
IVM2.24a divided by \(r_j\). On its cut choose
\(\alpha(0)=0\), with increment \(-6\pi j\).
The initial constant phase \(\alpha_0\) is retained
by its cap map
\(v\mapsto \exp(v\alpha_0/2)\exp(-i\alpha_0/2)\)
on the pure-imaginary unit three-ball.

Let \(H=e^{i\alpha/2}\). The actual normal arrow is
\[
 Q_N'=HQ_NH,\qquad
 q_H'=\exp(h\alpha/2)q_H\exp(i\alpha/2).
 \tag{IVM2.25}
\]
The paired endpoints are both \((-1)^{3j}\).
The constant-in-\(h\) map \(C(q_H)=Hq_HH\) extends
on the ball. Compose the actual map on the right
by \(C^{-1}\), leaving the left coordinate clutch
\[
 A_j(h,\theta)=\exp(h\alpha/2)\exp(-i\alpha/2).
 \tag{IVM2.26}
\]
Its zero-winding part contracts by its endpoint-zero
real logarithm. The remaining map is
\(\exp(-3\pi j\theta h)\exp(3\pi j\theta i)\).
For \(N=-3j\), its \(-1\) preimages occur at
\(h=-i,\theta=(2r+1)/(2N)\), with the directed
choices in the unit interval. The differential
in outward basis \((\mathbf j,-\mathbf k,\partial_\theta)\)
has target components
\((\mathbf k,\mathbf j,2\pi N i)\).
The positive target basis at \(-1\) is
\((i,\mathbf j,-\mathbf k)\). The local determinant
is \(2\pi N\), and the full degree is \(N=-3j\).

The complex fixed sphere is \(-S^2_h\). The fixed-tube
quotient has orientation \(-S^2_h\wedge dr\wedge S^2_\eta\),
and the inserted cap is \(-D^3_h\times S^2_\eta\).
The coordinate direction \(v_{\rm ref}=A_jv_j\)
gives the difference base
\(-(D^3_j\cup-D^3_{\rm ref})\times S^1_\theta\).
The radial connections \(fA_j^{-1}dA_j\), or
\(-f\,dA_j A_j^{-1}\) in the other frame, with
\(\int\operatorname{tr}(A^{-1}dA)^3=-24\pi^2\deg A\),
give \(c_1(E)=0,c_2(E)=3j,p_1(\operatorname{ad}E)=-12j\).
The physical tautological line has fibre value \(-1\).
Its projective relation therefore gives positive pair
\((3j,12j)\); the actual K68 Jin reversal gives
\[
                (\Delta k,\Delta\ell)_{\rm cap}=(-3j,-12j).
 \tag{IVM2.27}
\]
All initial phases, positive normal factors and the base
sphere flow remain in this full cap map.

\par

\medskip\noindent\textbf{IVM2.6. Original exterior Euler collar and integer assembly}

\par\nobreak

In the piece coefficients \(\gamma=\gamma_j\), retain
the actual exterior connection and its boundary primitive:
\[
 A_{\rm ext}=dd_j+(\alpha_q-3\theta)d\gamma+j\,du
                          +[j\alpha_q+(2-3j)\theta]dv,
\]
\[
 \beta=-j[du+(\alpha_q+(2-3j)\theta)dv-3\theta\,d\gamma].
 \tag{IVM2.28}
\]
The full zero-winding corrections are derivatives of
the original coefficient maps and deform with this
prescribed primitive. For
\(\Omega_B=d\alpha_q\wedge du\wedge dv\wedge d\gamma
\wedge d\theta\), exact exterior multiplication gives
\[
 3\int_0^1\beta(F_{\rm ext}+a\,d\beta)^2\,da
                 =(12j-15j^2+6j^3)\Omega_B.
 \tag{IVM2.29}
\]
The actual outward interpolation is piece one to
exterior zero, so its Euler contribution is
\(-12j+15j^2-6j^3\). Its Pontryagin term vanishes
by the full tangent collar of IVM2.4.

IVM2.19 and IVM2.21 give the interior changes
\((9j^2-15j,12j)\). Adding the directed cap and collar:
\[
 \Delta\ell=12j-12j=0,
\]
\[
 \Delta k=(9j^2-15j)-3j+(-12j+15j^2-6j^3)
                       =-6j^3+24j^2-30j.
 \tag{IVM2.30}
\]
The physical dual and parameter reversal are exactly
the K68 arrow used in ASI2. With \(j=-1\) this
difference is \((60,0)\); the actual ASI2 pair
\((71,-4)\) therefore determines \((k_0,\ell_0)=(11,-4)\).
This proves IVM2.1 without choosing an integer from
its residue.

The signed increments of \(S_j\) are
\[
 S_{j+1}-S_j=-(3j-2)(j-1),\qquad
 S_{j-1}-S_j=(3j-5)(j-2).
 \tag{IVM2.31}
\]
They are the actual VM2 native products with their
full Cox resets and original closing. The independent
integer comparison has retained that entire marking.

\par

\medskip\noindent\textbf{IVM2.7. Complete ordinary-double seed for further fan maps}

\par\nobreak

For the literal double of the single native product,
VM2 has trace half-winding \(t=-1\). The inverse
native-to-ordinary fixed-\(C_0\) endpoint carries
the pure-left degree \(+2\). Consequently
\[
 V_{00}=-2(\mu_j+\nu_j)-2=-12S_j,\quad
 \delta_0=V_{00}+4\nu_j=12S_j.
 \tag{IVM2.32}
\]
The complete nonfixed traces of VM2.17 yield
\[
 I_0=12j^2-32j+18,\qquad
 V=\begin{pmatrix}
 -12S_j&-I_0&I_0\\
 -I_0&6j-12&0\\
 I_0&0&6j-6
 \end{pmatrix},
\]
\[
 \delta=(12S_j,\ 12j^2-38j+28,\ -12j^2+26j-14).
 \tag{IVM2.33}
\]
For a selected ordinary diagonal \(a=V_{ii},b=\delta_i\),
the actual single native coordinates are
\[
               m_i=-1-(a+b)/4,\qquad n_i=(b-a)/4.
 \tag{IVM2.34}
\]
The division uses the free integral relative group
of the literal double and its original source/cap
closing. Both signed VM2 pairs are recovered from
IVM2.33 before any central reduction. The source
and full coefficient maps for subsequent fan changes
are their unchanged VM2 Cox morphisms.

% End expanded verified cp36 input

\par

% Expanded verified cp36 input: notes/part3_e6_selected_i3_m2_continuation_igc_proof.tex
Complete proof. Fix the original source
\(y^2=x^3+r_{\rm src}x^2+tx+t^2\), \(G=(0,t)\),
\(M_2=\mathcal O([2G]-[O])\), translation \(-3G\),
and the selected node tending to \(t=0\) at
\(s=1/4-r_{\rm src}=0\). The parameter is its positive
original \(s\)-circle and the scalar \(\lambda\) is fixed.
We retain every admitted good cap
\[
             m=3l,\quad l\ge1,\quad n\in\{1,-1\},
             \quad d_0\in\mathbb Z/m,
 \tag{IGC.1}
\]
its chosen scalar root \(R_\lambda^m=\lambda\), full
finite character, source lattice and annular reference.
The scalar is in a common strict contraction range for
the compared cap and the cubic reference. Every selected
fan uses its actual finite collection of charts.

The following maps identify the two characteristic
integers and the selected native products with those
of the cubic norm-one family, while preserving its
entire ordinary source/point line. This proves, on
the actual \(T_{-1}\) native fan, the physical Jin
pair \((71,-4)\) and class \(12\) for all caps in
IGC.1. On every \(T_j\), the pair is the IVM2 pair
\[
       (11-6j^3+24j^2-30j,-4).
 \tag{IGC.2}
\]
The full nonvertical algorithm is retained on the
identical selected exterior in the companion fan proof.

\par

\medskip\noindent\textbf{IGC.1. Exact selected heights and the full root source}

\par\nobreak

Choose the analytic square root equal to one at zero
and put
\[
 t_c(s)=\frac{1-36s-(1+12s)^{3/2}}{108},\qquad
 t_o(s)=\frac{1-36s+(1+12s)^{3/2}}{108}.
 \tag{IGC.3}
\]
They solve
\(27t^2+(18s-\tfrac12)t-4s(\tfrac14-s)^2=0\);
\(t_c(0)=0,t_o(0)=1/54\). With
\(v=8s/(1+4s+\sqrt{1+12s})\), \(D=v-2\),
the literal SI2 coordinate map is
\(s=v(v+4)/(4D^2)\), \(t_c=v/D^3\).
Thus \(q_c=t-t_c(s)\) is the same original height,
including its full source derivatives.

At infinity set
\[
 q_\infty=z^{3l},\quad t=z^{-3l},\quad
 x=z^{-2l}X,\quad y=z^{-3l}Y,
\]
\[
 Y^2=X^3+r_{\rm src}z^{2l}X^2+z^lX+1,\qquad G'=(0,1).
 \tag{IGC.4}
\]
The inverse on the punctured domain is displayed in
this equation, with the original projective \(O\)
chart. The deck is
\[
 D_m:(z,X,Y)\mapsto(\zeta_m z,\zeta_m^{-4l}X,Y).
 \tag{IGC.5}
\]
Put \(u=z^l\). The exact \(\mu_l\) quotient is
the smooth cubic-root equation
\(Y^2=X^3+r_{\rm src}u^2X^2+uX+1\), \(q_\infty=u^3\).
At the two actual marked sections use
\[
 2G'=(su^2,-1-\tfrac{s}{2}u^3),\qquad
 w_P=X-su^2,\qquad w_O=-X/Y.
 \tag{IGC.6}
\]
Their residual deck weights are \((w,u)=(2,1)\)
modulo three. The invariants \(A=w^3,B=u^3,C=wu\)
obey \(AB=C^3\). Their three resolution charts are
\[
 (A,B,C)=(a^{j+1}b^j,a^{2-j}b^{3-j},ab),
 \quad j=0,1,2;\qquad a'=b^{-1},\ b'=ab^2.
 \tag{IGC.7}
\]
The original base is exactly \(B\) on every chart.
The primitive quotient rays in \((w,u)\) are
\((2/3,1/3),(1/3,2/3)\); their actual \((w,z)\)
pullbacks are
\((2/3,1/(3l)),(1/3,2/(3l))\) in
\(\mathbb Z^2+\mathbb Z(2l,1)/(3l)\), whose
primitive boundary is \((0,1/l)\).

The strict-section support is
\(\max(v_w-2v_u,0)\), and the pulled point
section has order \(v_w\). Their exact difference
is \(\min(v_w,2v_u)=\min(v_w,2l v_z)\).
Retain its monomial Cartier generators separately
at \(P=2G'\) and \(O\), denoting the corrections
\(E_P,E_O\). This gives the actual module map
\[
 \beta^*\mathcal O([2G']-[O])=\pi^*M_2+E_P-E_O,
 \tag{IGC.8}
\]
with reciprocal coefficient maps in every chart
of IGC.7. The section equations IGC.6 and all
generators are holomorphic jointly in \(s\).
These are the full MG.3--6 graphs with the present
original \(r_{\rm src}=1/4-s\) substituted.

\par

\medskip\noindent\textbf{IGC.2. Chronological translation, full Miller norm and both signs}

\par\nobreak

The exact formal coordinate of \(3G'\) is
\[
 \tau(3G')=
 u^2\frac{-2s(2+su^3)}
 {8+6su^3+(s^2-s^3)u^6}.
 \tag{IGC.9}
\]
Let \(R=[m]^{-1}(3G')\), using the inverse formal
multiplication with derivative \(1/m\).
Its entire series is retained:
\(\tau(R)=-su^2/(2m)+O(u^5)\).
It is jointly analytic at \(s=u=0\), equals \(O\)
at \(s=0\), and commutes with the deck. Near each
marked point its actual translation is
\[
       w\mapsto w\,\varepsilon(w,u,s)+u^2a(w,u,s),
       \qquad \varepsilon(0)\ne0.
 \tag{IGC.10}
\]
The inverse has the same form. Both preserve
\((w,u),(w,u^2),(u,w),(u,w^2)\), their invariant
product and normalized blowup. They consequently
carry the full correction divisors and canonical
sections of IGC.8, not only their valuations.

For the actual \(P=2G'\), the full line morphism
covering \(p\mapsto p-R\) is
\[
 F(p)=\tau(R)\frac{\ell_{P,R}(p)}{X(p)-X(P+R)}.
 \tag{IGC.11}
\]
Its chord divisor is
\((P)+(R)-(P+R)-(O)=D_P-t_{-R}^*D_P\).
Write \(X_R=A/\tau^2,Y_R=-A/\tau^3\), \(\tau=\tau(R)\).
Its complete numerator after multiplication by
\(\tau\) is
\[
 \tau(Y-Y_P)+
       \frac{A+\tau^3Y_P}{A-\tau^2X_P}(X-X_P).
 \tag{IGC.12}
\]
The denominator \(A-\tau^2X_P\) is a unit.
At \(\tau=0\), this and the actual \(X-X_{P+R}\)
denominator are the same typed generator.
Thus the Hom and its reciprocal have coefficient
one there. Their horizontal divisors are the
displayed Cartier-module divisors; their generic
vertical orders are zero. On the smooth cubic
root there is one central component. Normality,
together with the actual graph ideals and maps
above, extends the morphism and inverse over
every joint chart. The chord weight
\(\zeta_m^{-2l}\) cancels the \(\tau(R)\) weight
\(\zeta_m^{2l}\), so it commutes with \(D_m\).

Every original Miller factor is exactly
\(f_{2,\rm old}=z^{-l}f_2'\). Its six chronological
factors give
\[
 B_{2,3,\rm old}=z^{-6l}B_{2,3}',\qquad
 (t-t_c)/t^3=z^{6l}h_c,\qquad
 B_c'=h_c B_{2,3}',\quad h_c=1-t_c(s)z^m.
 \tag{IGC.13}
\]
The central product of the three actual factors
\((Y-1)/X,-(Y+1)/X,2/X\) is \(-2\);
the full six-factor product is four.
At \(s=0\), \(3G'=O\) and \(h_c=1\), so
the actual Hom coefficient is also four.
The two parameter axes and every horizontal
Cartier divisor are retained in these equations.

For the distant height use \(h_o=1-t_o(s)z^m\).
The exact ratio
\[
                  H=h_c/h_o,\qquad
       H^{1/m}=\exp(m^{-1}\operatorname{Log}H)
 \tag{IGC.14}
\]
is deck invariant, analytic with analytic inverse
on a common smaller joint disc; its log has value
zero at \(z=0\). The original Hom changes by
this entire \(H\), including its derivatives.
Both \(B_c'\) and the chronological norm \(F^m\)
cover \(-3G'\). Their quotient \(\beta_c\)
is a nonzero base function: on each compact
connected source fibre it is a holomorphic
invertible scalar. It has central value four.
Fix its analytic root by the prescribed reference:
\[
 \widetilde F_c=\beta_c^{1/m}F,\qquad
 \widetilde F_c^{\,m}=B_c',\qquad
 \widetilde F_c=H^{1/m}\widetilde F_o .
 \tag{IGC.15}
\]
The remaining constant root choice is the
specified finite character. Define
\[
 B_\lambda=\lambda B_c',\qquad
 A=R_\lambda^n\zeta_m^{d_0}\widetilde F_c^{\,n}D_m .
 \tag{IGC.16}
\]
Chronological source and coefficient composition gives
\[
                     A^m=B_\lambda^n,\qquad[A,B_\lambda]=1.
 \tag{IGC.17}
\]
For \(n=-1\) every inverse coefficient is evaluated
at the inverse source point. The primitive
contracting generator is \(J=A^n\), with
\[
 J=R_\lambda\zeta_m^{nd_0}\widetilde F_c D_m^n,\qquad
                         J^m=B_\lambda.
 \tag{IGC.18}
\]
This retains the actual finite character for both signs.

\par

\medskip\noindent\textbf{IGC.3. Full annular reference and its inverse}

\par\nobreak

On each original infinity sector put
\[
 s_\alpha=\log_\alpha q_\infty/(2\pi i),\quad
 z_\alpha=e^{2\pi is_\alpha/m},\quad
 f_\alpha=\operatorname{Log}_\alpha\widetilde F_c,\quad
 c_\lambda=n\log R_\lambda/(2\pi i)+d_0/m .
\]
The source variable remains \(s\); \(s_\alpha\) is
the separate annular coordinate. The original
ordinary-to-cap map is
\[
 \Phi_\alpha=\operatorname{Exp}
       (ns_\alpha f_\alpha+s_\alpha c_\lambda\delta+h_\alpha)
                         I_{z_\alpha}.
 \tag{IGC.19}
\]
The map \(I_z\) is IGC.4--8, including its two
separate correction divisors. The source
displacement is the signed \(-ns_\alpha\)
Abel integral of the entire \(R\); its scalar
is \(R_\lambda^{ns_\alpha}e^{2\pi id_0s_\alpha/m}\).

For the actual overlap integers \(k,j_{\alpha\beta}\),
and full source/line lattices \(\Pi,\Lambda\), use
\[
 s_\beta=s_\alpha+k,\qquad
 f_\beta=D_{m*}^k f_\alpha+\Pi_{\alpha\beta},
\]
\[
 h_\beta-D_{m*}^k h_\alpha
  =-ns_\alpha\Pi_{\alpha\beta}
       +j_{\alpha\beta}\operatorname{Log}_\beta B_\lambda
       +\Lambda_{\alpha\beta}.
 \tag{IGC.20}
\]
Extend the specified reference cochain on fixed
annular contours, form its complete overlap residual
\(E_{\alpha\beta}\), and set
\(k_\alpha=-\sum_\beta\chi_\beta E_{\alpha\beta}\).
Subtract the common Cauchy transform of
\(\bar\partial k_\alpha\), using
\(\pi^{-1}\int\omega(w,s)/(q-w)\,dA(w)\).
The residual and correction vanish at the
reference. They depend holomorphically on the
full source disc. The added \(m^{-1}\log H\)
is retained in this same solve.
Direct substitution yields the exact squares
\[
 \Phi_\beta=A^kB_\lambda^{j_{\alpha\beta}}\Phi_\alpha,
       \qquad \Phi_\alpha B_\lambda=B_\lambda\Phi_\alpha.
 \tag{IGC.21}
\]
Their inverses reverse the Lie arrow and then
apply \(I_z^{-1}\) with the recovered source
coordinates. This retains every integer deck,
finite character, line coefficient and derivative.
All functions are single-valued on the closed
original \(s\)-disc. The source-circle closing
therefore is the actual identity on the original
source/root/reference family, for either norm.

\par

\medskip\noindent\textbf{IGC.4. Proper cap products with the identical prescribed boundary}

\par\nobreak

On the compact source-line cut choose a transported
Hermitian norm and \(\rho_0<1\) such that
\(\|B_\lambda v\|\le\rho_0^m\|v\|\), uniformly
over the closed source disc. Use the entire
iterates in
\[
       \|v\|'^2=\sum_{a=0}^{m-1}\rho_0^{-2a}\|J^av\|^2 .
 \tag{IGC.22}
\]
This is an equivalent norm on that compact cut, and
its actual shift is
\[
 \|Jv\|'^2=\rho_0^2[
       \|v\|'^2-\|v\|^2+\rho_0^{-2m}\|B_\lambda v\|^2]
                       \le\rho_0^2\|v\|'^2 .
 \tag{IGC.23}
\]
The norm shell, compact source charts and finite
group supply a proper quotient with boundary.
The physical circle is free there: a nonzero
power of the contracting generator changes
this norm, whereas a physical unit scalar
preserves it. The original source derivative
in \(s\) is one in IGC.4, and all corrections
and normalized translations are joint maps.
Hence these are actual smooth proper cap
families over the full source disc.

Prescribe one product on their common original
ordinary infinity annulus, with the whole
source and point/O line frames. Pull its
horizontal field into each cap by the full
\(\Phi_\alpha\) derivative. Extend by a relative
partition and physical-circle average,
retaining its original \(s\)-component.
Integrate along the radial segments of the
source disc. Compactness makes these flows
complete; backward integration is the inverse.
This yields a product of each actual cap
family agreeing exactly with that prescribed
ordinary boundary product. Its full tangent
and physical-line arrows are the integrated
derivatives and the corresponding line
connection. The central coincidence
of the source sections is treated by their
Cartier-module family throughout.

The original source outside these caps, its
point line, selected height, finite fan
charts and chosen Hopf completion are identical.
The signed meridian columns in the common
ordinary basis are, literally,
\[
 C_{m,n}=\begin{pmatrix}m&1\\n&0\end{pmatrix},\quad
 C_{m,n}^{-1}=\begin{pmatrix}0&n\\1&-mn\end{pmatrix},
\]
\[
 T_{m,n}=\begin{pmatrix}n&0\\3-mn&1\end{pmatrix},\quad
 C_{m,n}T_{m,n}=C_{3,1}.
 \tag{IGC.24}
\]
These are the full signed original cap meridian
maps of HC/SN/MG, in the same original generators
used in IGC.19--21. They preserve the physical
line; their determinants retain the norm sign.
The actual M2 source transport and good-root
attachment of K96/K105 carry the F41/F46
primitive sphere construction to this source
disc. On the identical selected Hopf boundary,
the physical Euler class in
(fixed sphere, complex meridian) is \((0,-1)\).
The actual orbit-cap and conifold difference
matrices are
\[
 \begin{pmatrix}0&-1\\-1&0\end{pmatrix},\qquad
 \begin{pmatrix}0&0\\-1&-1\end{pmatrix},
 \tag{IGC.25}
\]
with primitive conifold kernel \((1,-1)\).
Their complete Hopf/normal arrows are unchanged.
These integral Gysin and van Kampen maps give
the same oriented homotopy six-sphere property
for both signed cap attachments. The selected
fan words carry them by their actual primitive
common-link Cox maps.

Choose the line and tangent connections on
each cap quotient pulled back from its
five-dimensional body, agreeing with the
prescribed collar connection. On the product
with the parameter circle the six-forms
\(c_1(M)^3\) and \(p_1c_1(M)\) vanish separately,
because their full coefficients pull back
from that body. On the common exterior use
identical connections. The whole derivatives
and line arrows IGC.21 match their collars;
the opposite interpolation transgressions
cancel pointwise. Thus the two completed
characteristic integers agree individually.
Their physical line, selected quaternion,
own Hopf cap and Jin reversal are identical.

\par

\medskip\noindent\textbf{IGC.5. Native and fan consequences in the retained marking}

\par\nobreak

IGC.4 carries ASI2\textquotesingle s actual \(T_{-1}\) pair
\((71,-4)\) to every cap in IGC.1. Its
held-Cox coordinate map is unchanged, so
the native pair is \((-1,72)\).
For any signed \(T_j\), all IVM2 source,
graph, Cartier, line, tangent and Hopf maps
occur in the same selected piece, with this
identical exterior and infinity cap product.
The same comparison gives IGC.2 and the
complete IVM2 ordinary-double seed before
reduction modulo 28.

For a prescribed nonvertical finite fan
word, first choose the common smaller
selected tube for all its full coefficient
and inverse-height maps. It is disjoint
from the actual infinity cap. Every
common-link and native reset is performed
there with the unchanged ordinary source,
line and closing. Induction through those
explicit maps carries the full companion
all-fan algorithm to every cap and norm
in IGC.1.

At \(w=-1\), the actual point Tate marking
is \(Z_T=Z_w\); its positive loop has degree
one over \(s\). At \(w=+1\), the source
and point arrow is first
\((Z,V)\mapsto(Z^{-1},P^2V)\), then the
actual target period, giving \((Q/Z,V)\).
It sends the chosen fan to its literal
reflection \(J\), with all original
Cartier factors. Both the original
source and IGC.4 cap depend on \(w^2\).
The finite character and full reference
stay fixed through this identity.

% End expanded verified cp36 input

\par

% Expanded verified cp36 input: notes/part3_e6_selected_i3_m2_continuation_afm2_proof.tex
Complete proof. Keep SI2\textquotesingle s original source, section \(G\),
point line \(M_2=\mathcal O([2G]-[O])\), translation \(-3G\),
selected node approaching \(t=0\) at \(s=1/4-r_{\rm src}=0\),
positive original source circle, fixed scalar, and the complete
finite-character/annular reference. The proved cubic norm-one
seed is ASI2 on the actual \(T_{-1}\) native star:
\((k_M,\ell_M)=(71,-4)\), \((\mu,\nu)=(-1,72)\), class \(12\).
The exact IGC cap maps carry the same construction to every
\((m,n)=(3l,\pm1)\), \(l\ge1\), with its prescribed character.
All fan maps below use their unchanged selected exterior.

For each unimodular integer basis, the algorithm below gives
the actual return in \(\Theta_7=\mathbb Z/28\). Its image over
all these fans is the full set of 28 classes. The finite
witness table records actual paths from the same source family,
with the complete coefficient maps and unreduced increments.

\par

\medskip\noindent\textbf{AFM2.1. Entire characters and periods}

\par\nobreak

Let \(q=t-t_c(s)\) and retain
\[
 h:(Z,V)\mapsto(q\kappa Z,\alpha V),\qquad
 B:(Z,V)\mapsto(gZ,q\eta V),
\]
\[
 \kappa=P_G^2/c_q,\quad\alpha=P_G^2,\quad
 g=P_G^{-3},\quad
 \eta=\lambda_{\rm base}\eta_2^6t^{-3},\qquad
 \eta_2=a_T\Theta'_Q(1)\frac{\Theta_Q(P_G^2)}{\Theta_Q(P_G)^2}.
 \tag{AFM2.1}
\]
The genuine point coordinate \(V=V_{\rm pt}\), its two
Cartier-module maps, the whole theta derivative and every
original source scale are those of SI2. On this source
circle the full winding matrix is
\[
                         B_0=\begin{pmatrix}2&-3\\2&-3\end{pmatrix}.
 \tag{AFM2.2}
\]

For a basis \(P\in GL_2(\mathbb Z)\), write \(R=P^{-1}\).
The actual new characters and periods are
\[
 Z'=Z^{R_{11}}V^{R_{12}},\quad
 V'=Z^{R_{21}}V^{R_{22}},\quad q'=q,\qquad
 h'_j=h^{P_{1j}}B^{P_{2j}} .
 \tag{AFM2.3}
\]
The inverses use \(P\) for characters and \(R\) for
period generators. Their exact coefficients are
\[
 A_{ij}=(\kappa^{P_{1j}}g^{P_{2j}})^{R_{i1}}
           (\alpha^{P_{1j}}\eta^{P_{2j}})^{R_{i2}},
\]
\[
 h'_1=(qA_{11}Z',A_{21}V'),\qquad
 h'_2=(A_{12}Z',qA_{22}V').
 \tag{AFM2.4}
\]
The \(q\) exponents are \(RP=I\) exactly. These full
functions and all reciprocal powers, including their
normal derivatives, are retained in every chart.
Each prescribed finite word admits one uniform smaller
selected tube over the compact original source circle.
The exterior arrow is identity on the original
absolute torus, with the displayed inverse characters.
It carries the original source, point line, cap and
complete annular reference. The original fixed-scalar
source closing is identity, by the joint disc maps
of SI2 and IGC.1--4.

\par

\medskip\noindent\textbf{AFM2.2. Full ordinary double, action prisms and native reset}

\par\nobreak

In the two vertex triples \((a,b,c)\), \((u,v,w)\),
the height is \(q=abc=uvw\). The complete three
overlap rows are
\[
 \begin{array}{c|ccc}
 &u&v&w\\ \hline
 C_0&A_{11}A_{21}bc&A_{21}^{-1}ab&A_{11}^{-1}b^{-1}\\
 C_1&A_{12}A_{22}bc&A_{22}^{-1}c^{-1}&A_{12}^{-1}ac\\
 C_2&a^{-1}&ab&ac
 \end{array}.
 \tag{AFM2.5}
\]
Recover the same original \(q\), and use the reciprocal
coefficient at that recovered point for each inverse.
The actual second-vertex arrow
\((u,A_{22}(q,s)v,A_{11}(q,s)w)\) has target height
\(qA_{11}A_{22}\). Its inverse first solves this
one-variable height equation, then divides by the full
units. The derivative at \(q=0\) is a unit, uniformly
on the smaller compact tube.

Take the literal double pullback of every single
native star product, including its original closing
and full cap maps. On that double put
\[
 \mathsf B=2P^{-1}B_0P,\qquad t=-1,\qquad
 l=(\mathsf B_{22}-\mathsf B_{21},
             \mathsf B_{11}-\mathsf B_{12},\mathsf B_{22}).
 \tag{AFM2.6}
\]
The actual three normal matrices are
\[
 D_0=\operatorname{diag}(A_{22}/A_{21},A_{11}A_{21}),\
 D_1=\operatorname{diag}(A_{11}/A_{12},A_{12}A_{22}),\
 D_2=\operatorname{diag}(A_{22},A_{11}).
 \tag{AFM2.7}
\]
Their determinant winding is \(2t=-2\). Remove only
the written angular powers and reference values:
the resulting loops have their unique based
periodic logarithms on the full \(q\)-disc.
The explicit quaternion path is
\[
 H_d(x,\theta)=
  (\sin\vartheta(x)e^{id\theta}+\mathbf j\cos\vartheta(x))
  (\sin\vartheta(x)+\mathbf j\cos\vartheta(x))^{-1},
 \qquad Q\mapsto H_tQH_{l_i-t}.
 \tag{AFM2.8}
\]
The endpoint-flat \(\vartheta\) has endpoints
zero and \(\pi/2\). Compose this path with the
two complete logarithmic unit exponentials.
Its endpoints are identity and exactly
\(D_i(q,s_*)^{-1}D_i(q,s)\).

In the normal interior the full map is
\(n'=An+G(n)\), with \(G=O(|n|^3)\), since
\(q=zn_0n_1\). Choose the compact tube with
\(\|A^{-1}DG\|<1/2\). Iteration
\(n=A^{-1}n'-A^{-1}G(n)\) gives its unique
smooth inverse; the recovered original height
is \(zn_0n_1\) at that point. At the two
vertices this is exactly the height inverse
above. The normal-rescaling flow restores
the complete nonlinear germs, fixing the
protected vertex collars. Thus AFM2.8 gives
the actual ordinary double product with
every unit, physical action and source/cap
closing retained.

For \((x,y)^t=P^{-1}e_2\), define
\[
 w=(x,y,-x-y),\qquad d=(-x-2y,-2x-y,x-y),
\]
\[
 I_i=-w_i+d_i(l_i+1),\qquad
 W_i=-w_i-d_i(l_i+1).
 \tag{AFM2.9}
\]
These are the actual action and stable-derivative
traces of AFM2.8. The VM2 quaternion-axis calculation,
with its oriented sphere and paired Spin endpoints,
gives the whole prism equations
\[
 V(w_je_j-w_0e_0)=I_je_j-I_0e_0,\qquad
 w_j\delta_j-w_0\delta_0=W_j-W_0,\quad j=1,2.
 \tag{AFM2.10}
\]
Here \(V\) is symmetric and integral.
The ambient-to-star derivative gives
\(2n_i^{\rm ord}=\delta_i-V_{ii}\).
The exact native-to-ordinary Hopf map is
\((UH_tU^{-1})q_HH_{l_i-t}\), with the full
unit and positive-norm paths. Its left
degree is \(2t=-2\), and its right degree
is zero. Its inverse-cylinder endpoint is
the ordinary pure-left degree \(+2\).
Consequently, for \(a=V_{ii},b=\delta_i\),
the single native coordinates are
\[
                m_i=-1-\frac{a+b}{4},\qquad
                n_i=\frac{b-a}{4}.
 \tag{AFM2.11}
\]
These divisions take place in the free integral
relative group of the literal double, before
any homotopy-sphere class is used.

For every directed flop retain the complete
two-ruling Cox morphism and inverse on its
common link. On the reverse \(C_1\) chart
they are
\[
 (a',b',c')=(a/b,b,bc),\qquad
 (a,b,c)=(a'b',b',c'/b').
 \tag{AFM2.12}
\]
The two contractions of
\(\operatorname{Tot}\mathcal O(-1,-1)\)
are
\([A_0,A_1,\tau C_0,\tau C_1]\) and
\([\tau A_0,\tau A_1,C_0,C_1]\), both
pulling \(q\) back to
\(\tau^2A_0A_1C_0C_1\).
Their full coefficient lifts are AFM2.4--5
with the corresponding cyclic vertex order.
The unrestricted four-weight common link
\((y,q_H)\mapsto(q_H^{-1}yq_H,q_H)\)
and its conjugation inverse carry the
physical line and original Hopf cap.
The native reset is exactly
\((m,n)\mapsto(m+2n,-n)\), and the signed
central change is \(+n\), by the actual
K23 directed Hopf transfer. This is the
same full reset used in VM2, with every
coefficient replaced by AFM2.4.

\par

\medskip\noindent\textbf{AFM2.3. Complete integer vertical seed}

\par\nobreak

IVM2 gives the actual source/graph/Cartier
and Hopf comparison for every signed \(j\).
Its relative line is
\[
 U_j\simeq U_{0,q'}\otimes
       \mathcal T_b^{3j^2-5j}\otimes\mathcal T_\gamma^{6j},
 \quad Q(q',s)=Q(q,s)g(q,s)^j.
 \tag{AFM2.13}
\]
The full four pairings are \((1,0,0,2)\);
the normal source phase is \(-3j\);
the cap is \((-3j,-12j)\), and the entire
outward Euler collar is
\(-12j+15j^2-6j^3\). Its tangent collar
has zero Pontryagin transgression.
Thus IVM2 proves, integrally,
\[
 S_j=2-j^3+4j^2-5j,\qquad
 (k_j,\ell_j)=(6S_j-1,-4),\qquad
 (\mu_j,\nu_j)=(-1,6S_j).
 \tag{AFM2.14}
\]
It uses the actual ASI2 \(j=-1\) seed,
not a residue lift. The exact ordinary
double is
\[
 I_0=12j^2-32j+18,\qquad
 V=\begin{pmatrix}-12S_j&-I_0&I_0\\
 -I_0&6j-12&0\\ I_0&0&6j-6\end{pmatrix},
\]
\[
 \delta=(12S_j,12j^2-38j+28,-12j^2+26j-14).
 \tag{AFM2.15}
\]
Its nonfixed entries are independently
the full VM2 action/derivative prisms.
Both signed vertical resets reproduce
\(S_{j+1}-S_j=-(3j-2)(j-1)\) and
\(S_{j-1}-S_j=(3j-5)(j-2)\).
The exterior, finite cap, source
orientation and full original closing
remain identical in this comparison.

\par

\medskip\noindent\textbf{AFM2.4. Exact finite word and full target reconstruction}

\par\nobreak

Put
\[
 R_0=\begin{pmatrix}-1&-1\\1&0\end{pmatrix},\
 Q_0=\begin{pmatrix}1&-1\\0&1\end{pmatrix},\
 Q_1=\begin{pmatrix}1&0\\-1&1\end{pmatrix}.
 \tag{AFM2.16}
\]
The bases \(P,PR_0,PR_0^2\), their negatives,
and the column swap describe the identical
infinite two-triangle fan. Substitution
in the vertices \(\{0,u,v\}\) and
\(\{u,v,u+v\}\) gives the actual integral
translations for each equality. Identity
on the original absolute characters and
the inverse period basis therefore gives
the literal affine semigroup maps, with
all full coefficients, source/cap and
closing unchanged.

There is a unique representative either
\(T_j\), or
\[
 C=\begin{pmatrix}p&-r\\q&s\end{pmatrix},
             \qquad p,r>0,\quad\det C=1.
 \tag{AFM2.17}
\]
Indeed the primitive physical vector has
a unique oriented cone in the fan;
on its boundary the primitive edge
gives the written \(T_j\), and in its
interior it gives exactly the positive
\(p,r\) representative. These are the
six explicit vertex translations above,
after a column swap for determinant
negative one. To reduce the interior
case, for \(r\ge p\) multiply on the
right by \(Q_0^{-1}\); for \(p>r\)
multiply by \(Q_1^{-1}\). The new
positive pair has smaller sum \(p+r\).
Since the determinant is one, it
terminates at \((p,r)=(1,0)\), hence
at the exact \(T_j\). Reversing this
finite word uses only the actual
\(C_0,C_1\) flops \(Q_0,Q_1\).
Every new interior target has
physical weights \((r,p,-r-p)\),
all nonzero. Thus no target fixed
coordinate is guessed from the prism.

For one such step keep the selected
old \(a=V_{ii},b=\delta_i\), and compute
the target \((w',I',W')\) from AFM2.4--9.
The same ruled curve is target \(C_2\);
its full complex reset is
\(V'_{22}=a,\delta'_2=2a-b\).
Set
\[
 c=(w'_2)^2a-w'_2 I'_2,\qquad
 e=w'_2(2a-b)-W'_2.
\]
The whole target is uniquely
\[
 V'_{kl}=\delta_{kl}\frac{I'_k}{w'_k}
                  +\frac{c}{w'_kw'_l},\qquad
 \delta'_k=\frac{W'_k+e}{w'_k}.
 \tag{AFM2.18}
\]
To derive uniqueness, multiply each
column of AFM2.10 by its weight:
subtracting the two equations makes
\(w'_kw'_l(V'_{kl}-\delta_{kl}I'_k/w'_k)\)
the same scalar for every \(k,l\);
symmetry identifies it with \(c\).
The derivative equations give the
single scalar \(e\) in the same way.
The selected two entries determine
both scalars as written. These are
the actual integral characteristic
entries of the target full prism;
the finite execution checks their
integrality before any reduction.

Add the native right coordinate
\((b-a)/4\) to the integer path sum,
carry the complete \(V',\delta'\),
and repeat. This constructs the
actual return on the canonical
fan by the finite original family
maps. The identical-fan arrows
then give it on the input basis.

\par

\medskip\noindent\textbf{AFM2.5. Explicit complete image and retained actual seed}

\par\nobreak

The actual forward \(C_2\) arrow
\(T_{-1}\to T_0\) has native pair
\((15,-10)\), by VM2. Its target
reset is \((-5,10)\). It sends
the proved class \(12\) to \(2\),
with the full \(P_G,\eta_2,t,c_q\)
coefficient maps, source and cap
closing retained. Every witness
below starts with this same
explicit prefix; it then follows
the displayed \(C_0,C_1\) word
from \(T_0\).

At \(T_0\) the complete seed is
\[
 V=\begin{pmatrix}-24&-18&18\\-18&-12&0\\18&0&-6\end{pmatrix},
       \qquad \delta=(24,28,-14),\qquad S=2 .
 \tag{AFM2.19}
\]
The first \(C_0\) increment is \(12\).
The next \(C_1\) increment is \(102\),
giving the actual two-flop basis
\(\left(\begin{smallmatrix}2&-1\\-1&1\end{smallmatrix}\right)\)
and unreduced value \(116\), class \(4\).
These are full coefficient/prism
computations by AFM2.18.

The companion witness table lists
all 28 residues, their complete
bases, finite words and integer
sums. Each JSON step retains the
old single native pair, selected
reset, target \(V,\delta\), and
all four full multiplier-exponent
vectors. Consequently every
residue is realized by a constructed
source family on its displayed fan.
Every return already lies in
\(\mathbb Z/28\) by the original
oriented sphere and K68 Hopf
identification. These two facts
prove that the image is exactly
all 28 classes.

For any cap in IGC.1, its explicit
fixed-boundary product leaves
every selected map above unchanged.
The same 28 families and their
exact characteristic/prism maps
therefore realize that image for
each such cap and norm. The
positive \(w=+1\) marking uses
the actual source/point reflection
\((Q/Z,V)\), so its algorithm is
evaluated on the reflected fan
\(JP\). The original parameter
degrees, finite character and
complete annular reference remain
those of SI2 and IGC.

% End expanded verified cp36 input


\par\medskip\noindent\textbf{All 28 actual witness paths.}\par
In the table, \(0\) means the directed \(C_0\) arrow and \(1\)
the directed \(C_1\) arrow.  Each word is traversed in its displayed
order after the common \(T_{-1}\to T_0\) prefix of AFM2.5. The
integer sum includes the actual \(T_0\) seed \(2\). The full old
and target matrices, native reset and coefficient-exponent maps at
every step are retained in the \texttt{all28\_witnesses} field of
\path{code/part3/e6_selected_i3_m2_all_fans.json}.

% Expanded verified cp36 input: notes/part3_e6_selected_i3_m2_all_fan_witness_table.tex
\begin{center}
\small
\begin{longtable}{rlll}
Class & Basis $P$ & Native word & Integer sum\\\hline
0 & $\left(\begin{smallmatrix}6&-11\\-1&2\end{smallmatrix}\right)$ & 0,0,0,0,0,1,0 & 34944\\
1 & $\left(\begin{smallmatrix}7&-3\\-2&1\end{smallmatrix}\right)$ & 0,0,0,1,1 & 6329\\
2 & $\left(\begin{smallmatrix}1&0\\0&1\end{smallmatrix}\right)$ & empty & 2\\
3 & $\left(\begin{smallmatrix}21&-8\\-13&5\end{smallmatrix}\right)$ & 0,1,0,1,0,1,1 & 1000107\\
4 & $\left(\begin{smallmatrix}2&-1\\-1&1\end{smallmatrix}\right)$ & 0,1 & 116\\
5 & $\left(\begin{smallmatrix}1&-3\\0&1\end{smallmatrix}\right)$ & 0,0,0 & 145\\
6 & $\left(\begin{smallmatrix}5&-2\\-2&1\end{smallmatrix}\right)$ & 0,0,1,1 & 2162\\
7 & $\left(\begin{smallmatrix}5&-8\\-3&5\end{smallmatrix}\right)$ & 0,1,0,1,0 & 36379\\
8 & $\left(\begin{smallmatrix}3&-1\\-2&1\end{smallmatrix}\right)$ & 0,1,1 & 484\\
9 & $\left(\begin{smallmatrix}4&-3\\-1&1\end{smallmatrix}\right)$ & 0,0,0,1 & 1325\\
10 & $\left(\begin{smallmatrix}9&-5\\-7&4\end{smallmatrix}\right)$ & 0,1,1,1,0,1 & 70150\\
11 & $\left(\begin{smallmatrix}2&-3\\-1&2\end{smallmatrix}\right)$ & 0,1,0 & 767\\
12 & $\left(\begin{smallmatrix}6&-1\\-5&1\end{smallmatrix}\right)$ & 0,1,1,1,1,1 & 6564\\
13 & $\left(\begin{smallmatrix}9&-4\\-2&1\end{smallmatrix}\right)$ & 0,0,0,0,1,1 & 14685\\
14 & $\left(\begin{smallmatrix}1&-1\\0&1\end{smallmatrix}\right)$ & 0 & 14\\
15 & $\left(\begin{smallmatrix}4&-7\\-1&2\end{smallmatrix}\right)$ & 0,0,0,1,0 & 8135\\
16 & $\left(\begin{smallmatrix}7&-10\\-2&3\end{smallmatrix}\right)$ & 0,0,0,1,1,0 & 47924\\
17 & $\left(\begin{smallmatrix}1&-4\\0&1\end{smallmatrix}\right)$ & 0,0,0,0 & 325\\
18 & $\left(\begin{smallmatrix}3&-4\\-2&3\end{smallmatrix}\right)$ & 0,1,1,0 & 3686\\
19 & $\left(\begin{smallmatrix}7&-19\\-4&11\end{smallmatrix}\right)$ & 0,1,0,0,1,0,0 & 592275\\
20 & $\left(\begin{smallmatrix}7&-2\\-3&1\end{smallmatrix}\right)$ & 0,0,1,1,1 & 6656\\
21 & $\left(\begin{smallmatrix}3&-5\\-1&2\end{smallmatrix}\right)$ & 0,0,1,0 & 2989\\
22 & $\left(\begin{smallmatrix}3&-7\\-2&5\end{smallmatrix}\right)$ & 0,1,1,0,0 & 16262\\
23 & $\left(\begin{smallmatrix}8&-13\\-3&5\end{smallmatrix}\right)$ & 0,0,1,0,1,0 & 143215\\
24 & $\left(\begin{smallmatrix}23&-10\\-16&7\end{smallmatrix}\right)$ & 0,1,1,0,0,0,1,1 & 1907524\\
25 & $\left(\begin{smallmatrix}1&-2\\0&1\end{smallmatrix}\right)$ & 0,0 & 53\\
26 & $\left(\begin{smallmatrix}10&-7\\-7&5\end{smallmatrix}\right)$ & 0,1,1,0,0,1 & 126194\\
27 & $\left(\begin{smallmatrix}2&-5\\-1&3\end{smallmatrix}\right)$ & 0,1,0,0 & 3023\\
\end{longtable}
\end{center}

% End expanded verified cp36 input


\par\medskip\noindent\textbf{Exact execution and independent full audits.}\par
The signed-native primary certificate passes 727 checks, with an
independent full-prism and map certificate of 531 checks. The whole
integer comparison passes 685 checks; its independent full geometric,
source/line/cap/exterior certificate passes 487. The good-cap
certificate passes 2435 checks and its independent symbolic root,
Cartier, norm, character and meridian replay passes 981.

The complete fan certificate passes 2964452 checks on 6504 input
unimodular bases and 814 canonical fans, with independent all-three-edge
traversals from five signed vertical seeds. The independent root
replay passes 325236 checks by Gaussian elimination of the entire
nine-unknown target prism system on every bounded fan and every step
of all 28 witness paths. It does not use the primary closed-form
target reconstruction. These finite replays certify the displayed
exact identities and witnesses; AFM2.1--4 supplies the terminating
construction and proof for every integral basis.

P3K.111 contains the complete original source, corner, normal,
reference, cap and outer-transfer proof of the actual initial pair.
The present four complete sources, their exact certificates, full
independent geometric audits, and every exported support are locked
in \path{records/k112_export_manifest.json}. The mathematical
sources remain byte-for-byte unchanged by this export.

% End expanded verified cp36 input

\end{proof}


\subsection{The intrinsic elliptic quotient and all fixed-source endpoints on the noncollision E6 locus}
\begin{theorem}[Source result P3K.113 under continuation review]\label{night:P3K.113}
\StatementRecord{night:P3K.113}{open}{Historical events: GPT-6 Astra (OpenAI), ChatGPT Pro mode, effort Pro, S6\_06\_BLIND, 16--18 November 2026: targeted proof check. Also checked in this writing. Continuation mathematical and source checks: GPT-6 Astra (OpenAI), ChatGPT Pro mode, effort Pro, S6\_10\_PAPERS continuation, 19 November 2026: computation or formal certificate; targeted proof check; used and rechecked in this run. Continuation metadata events: GPT-6 Astra (OpenAI), ChatGPT Pro mode, effort Pro, S6\_10\_PAPERS continuation, 19 November 2026: used and rechecked in this run; explicitly scoped metadata reconciliation; metadata reconciliation only; no new mathematical check. GPT-6 Astra (OpenAI), ChatGPT Pro mode, effort Pro, S6\_10\_PAPERS continuation, 19 November 2026: used and rechecked in this run; receiving source placement; metadata reconciliation only; no new mathematical check. GPT-6 Astra (OpenAI), ChatGPT Pro mode, effort Pro, S6\_10\_PAPERS continuation, 19 November 2026: used and rechecked in this run; source-row intake; metadata reconciliation only; no new mathematical check. \emph{Pre-continuation recorded proof scope.} Historical event: GPT-6 Astra (OpenAI), ChatGPT Pro mode, effort Pro, S6\_06\_BLIND, 16--18 November 2026: targeted proof check, recorded producing-run re-derivation. Also checked in this writing by GPT-6 Astra (OpenAI), ChatGPT Pro mode, effort Pro, S6\_10\_PAPERS continuation, 19 November 2026: computation or formal certificate for the explicitly listed algebraic substeps in the event ledger; full continuation rederivation and renewed credit comparison remain pending.}{\cref{night:P3K.96,night:P3K.109,night:P3K.110}; the original source, point line, roots, characters, reference, selected fan and scalar range stated here}
Retain the original marked source
\[
 E_r:\quad y^2=x^3+rx^2+tx+t^2,\qquad G=(0,t),\qquad
 r\in\mathcal V=\mathbb C\setminus\{0,\tfrac14,\tfrac13\}.
\]
Choose either simple root \(t_c\) of
\(27t^2+(4-18r)t+r^2(4r-1)=0\), with its original selected
height \(q=t-t_c\), branch marking and prescribed lattice-basis
fan. Use the genuine point line \(M_a=\mathcal O([aG]-[O])\),
translation by \(-bG\), and either full cap progression
\[
 (a,b;m)=(2,3;3l),\quad l\ge1,
 \qquad (a,b;m)=(6,1;3l+2),\quad l\ge0.
\]
Retain either primitive norm \(n=\pm1\), every prescribed finite
character \(\zeta_m^d\), the full annular reference of P3K.96,
the ordinary attachments at the other finite fibres, and a fixed
nonzero scalar in the strict contraction range. Let \(X\) be the
resulting completed complex threefold.

Every holomorphic isomorphism between two completed scalar endpoints
of this fixed marked source preserves the original base pointwise.
On each regular fibre its lift carries the actual cyclic cover
\(M_a^\times\) and its physical fibres to the corresponding
target cover and physical fibres. The induced source map is a
degree-one affine holomorphic isomorphism of the original elliptic
fibres. The complete proof includes its affine inverse and the full
integral kernel lattice.

Multiplication on the genuine point line, carried through all root,
reference and fan charts, gives the exact group isomorphism
\[
 \mathbb C^*_{\rm physical}
       \xrightarrow{\ \cong\ }\operatorname{Aut}_{\rm hol}(X).
\]
Every such automorphism is isotopic to the identity through this
physical action. The ratio of two endpoint identifications at the
same scalar winding is exactly a physical scalar, with its completed
isotopy and inverse retained.

On the scalar cover keep
\[
 \lambda_\theta=e^{2\pi i\theta}\lambda_*,\qquad
 R_\theta=e^{2\pi i\theta/m}R_*,\qquad R_*^m=\lambda_*.
\]
With the full SEF reference restoration, every holomorphic endpoint
has equal original scalars at its two ends and integral winding
\(J\). All endpoint linear parts on the ordinary torus are \(+I\).
The complete groups of possible windings are
\[
 \begin{array}{c|c}
 (a,b;m)&\text{scalar winding group}\\\hline
 (2,3;3l)&l\mathbb Z\\
 (6,1;3l+2)&m\mathbb Z.
 \end{array}
\]
For M2 put \(s=1/4-r\) and retain the full source lift
\[
 f_2(p)=\frac{y-t-x/2}{x-s},\qquad
 B_{2,1}(p)=f_2(p)f_2(p-G),\qquad U_G=B_{2,1}^{-1},\qquad S=tU_G.
\]
The inverse map evaluates its coefficient at the actual translated
source point. On the good root retain
\(U_{G,\rm old}=z^{2l}U'_G\). The primitive forward endpoint
at \(J=l\) is
\[
 E_R^{\rm ord}=(R^{-l}S)^n,\qquad
 E_R^{\rm cap}=(R^{-l}U'_G)^n.
\]
Their three chronological composites, at
\(R,\zeta_3R,\zeta_3^2R\), are exactly
\[
 E_2^{\rm ord}E_1^{\rm ord}E_0^{\rm ord}=q^nB_\lambda^{-n},
 \qquad
 E_2^{\rm cap}E_1^{\rm cap}E_0^{\rm cap}=h^nB_\lambda^{-n},
 \quad h=q_\infty q=1-t_cz^m.
\]
For M6 the primitive winding \(J=m\) is the canonical forward
\(q^n/h^n\) quotient closing. Negative windings use the exact
chronological inverses. Every source, Cartier, norm, character and
annular multiplier remains in these endpoint maps. The physical
return uses the inverse endpoint. This theorem assigns no new return
class. Its source domain is exactly \(\mathcal V\).
\end{theorem}
\begin{proof}[Supplied source argument; continuation review in progress]

% Expanded verified cp36 input: notes/part3_e6_intrinsic_base_proof.tex
\par\bigskip\noindent\textbf{Complete supporting proof 1: EB.}\par


% Expanded verified cp36 input: notes/part3_e6_intrinsic_base_eb_proof.tex
\emph{Audit provenance.} The submission sentence retained in the opening source paragraph below records its preparation stage. The complete EB, EBR and EBP geometric audits have passed, with 1098 primary, 5103 root-independent and 2396 further independent exact checks. All original source text is retained.

\par

\medskip\noindent\textbf{The intrinsic elliptic quotient on the full noncollision E6 source locus}

\par\nobreak

This separate continuation leaves the audited K109/EF source unchanged.
It constructs the elliptic quotient carried by an arbitrary endpoint
map before using its Weierstrass invariant to fix the source base.
It thereby treats the finite branch-permutation candidates excluded
by EF\textquotesingle s auxiliary centroid polynomials. No source or scalar return is
computed here. The proof is submitted for independent audit.

\par

\medskip\noindent\textbf{EB.1. The exact source and completed objects}

\par\nobreak

Retain
\[
 E_r:\quad y^2=x^3+rx^2+tx+t^2,\qquad G=(0,t),\qquad
 \mathcal V=\{r:r(4r-1)(3r-1)\ne0\}.
 \tag{EB.1}
\]
Choose either simple I1 as the selected one, with its original height
\(q=t-t_c\), complete branch and fan marking. Use either
\((a,b;m)=(6,1;3l+2)\), \(l\ge0\), or
\((a,b;m)=(2,3;3l)\), \(l\ge1\), either norm \(n=\pm1\),
and the prescribed finite character and full K96 reference.

The original source has precisely IV*, I2 and two simple I1 fibres
throughout EB.1, by the exact discriminant in EF.8. The root,
point-Cartier and annular maps EX.1-\/-19 and EF.28-\/-38 use
\(s=1/4-r\ne0\); the finite source charts and their full line
modules remain through every point of EB.1. A compact path in the
marked simple-node cover from the concrete EF.5 source to a given
object has a common smaller strict scalar contraction. All actual
source/root/line maps, the prescribed selected fan and the reference
continue on this path. The F46 and F35 integral source/flop topology
maps therefore give the same sphere, with \(H^2=H^4=0\).
The scalar can then be moved radially inside its contraction domain
to the specified value using the same complete cap square. This
gives the topology on EB.1 without any condition on \(P_r,Q_r\).

The horizontal-divisor and transverse \(t=u^e\) Laurent argument
EF.2 now applies literally. It proves
\(\mathbb C(X)=\mathbb C(t)\), with its specified pullback.
Every holomorphic isomorphism between two completed endpoints of the fixed
source thus descends to an actual Möbius base map \(h\).

\par

\medskip\noindent\textbf{EB.2. The full period module forces the original elliptic quotient}

\par\nobreak

On the regular base keep the full actual lattice of EF.11,
\[
 \gamma=(-bg,\ell),\qquad e_1=(1,0),\qquad
 e_2=(\tau,ag),\qquad\delta=(0,1),
 \tag{EB.2}
\]
and the complete monodromy EF.18-\/-19, including its central cochain.
The nonzero Abel-class proof EF.4 uses the nontrivial I1 monodromy
and the actual infinity logarithm of \(3G\), whose original
additive value is \(-s/2\). Both remain exactly as written on EB.1.
In particular the full invariant lattice line and invariant dual line
are intrinsically
\[
       \mathbb Z\delta,\qquad\mathbb Z\psi,
       \quad\psi(\gamma)=1,\quad\psi(e_1)=\psi(e_2)=\psi(\delta)=0.
 \tag{EB.3}
\]
The calculation establishing these lines does not assume that an
endpoint fixes the base: it is a computation of the invariant
subgroups of the local system itself. An isomorphism covering
\(h\) carries them to the corresponding subgroups of the pulled
back local system. Primitivity therefore gives
\(F_*\delta=\pm\delta\) and \(F^*\psi=\pm\psi\).

The subgroup \(\ker\psi\) gives the actual cyclic cover of a
regular torus fibre,
\[
 \widetilde T_t=
 \mathbb C^2/\langle e_1,e_2,\delta\rangle
           \simeq M_a^\times|_{E_{r,t}}.
 \tag{EB.4}
\]
The displayed identification is precisely the point-theta arrow
\(V_{\rm pt}=f_{P^a}(Z)v_M\), including its full constant factor
and the inverse Cartier-module maps EF.10-\/-17. A local lift of
\(F\) between these covers exists because it carries this exact
subgroup, and it is affine in the displayed universal-cover
coordinates: its periodic derivative is constant on the compact
torus fibre. Since the derivative takes the real vector
\(\delta=(0,1)\) to \(\pm(0,1)\) and is complex linear, it
takes its full complex span to itself. Its first coordinate is
therefore exactly
\[
                       z'=A(t)z+z_0(t),\qquad A(t)\ne0.
 \tag{EB.5}
\]
The two elliptic periods and the primitive \(\delta\) are carried
bijectively inside the actual kernel lattice. Quotienting that
lattice by \(\mathbb Z\delta\) gives
\[
          A(t)\langle1,\tau(t)\rangle
                =\langle1,\tau(h(t))\rangle.
 \tag{EB.6}
\]
Thus EB.5 descends to a degree-one affine holomorphic isomorphism
of the original elliptic source fibres. In multiplicative terms the
map first lifts to the genuine \(M_a^\times\), then preserves its
physical \(\mathbb C^*\) fibres with exponent \(\pm1\), and
then descends by the actual source projection. Its inverse follows
the same three arrows. The full line constants and the central
integer cochain remain in the lifted second coordinate; this
quotient assertion does not delete them from the completed map.

\par

\medskip\noindent\textbf{EB.3. The exact source invariant fixes three actual base points}

\par\nobreak

The Weierstrass coefficients give
\[
 c_4=16(r^2-3t),\quad
 \Delta=-16t^2[27t^2+(4-18r)t+r^2(4r-1)],
\]
\[
 j_r(t)=\frac{-256(r^2-3t)^3}
 {t^2[27t^2+(4-18r)t+r^2(4r-1)]}.
 \tag{EB.7}
\]
Because EB.5-\/-6 are actual elliptic isomorphisms, their source
invariant satisfies the literal equality
\(j_r(h(t))=j_r(t)\). Every relevant valuation is exact on EB.1.
The pole at zero has order two; the other two finite poles have
order one. The numerator has its triple zero at
\(t_z=r^2/3\), and at that point its remaining denominator factor is
\[
 27t_z^2+(4-18r)t_z+r^2(4r-1)
                  =\frac{r^2(3r-1)^2}{3}\ne0.
 \tag{EB.8}
\]
The unique other zero is infinity, of order one; in its exact
coordinate \(q_\infty=1/t\),
\(j_r(1/q_\infty)=256q_\infty+O(q_\infty^2)\).
Hence every base map preserving EB.7 fixes the unique double pole
zero, the unique simple zero infinity, and the unique triple zero
\(r^2/3\). These three points are distinct. The first two make the
Möbius map \(h(t)=ct\), and the third makes \(c=1\).

This proves pointwise preservation of the original base on all of
EB.1, including every candidate permutation of the two marked
simple-node branches. It is derived from the actual cyclic-cover
and elliptic-source maps, before any invariant is compared.

\par

\medskip\noindent\textbf{EB.4. Completion of the fixed-source group and endpoint scope}

\par\nobreak

With the base now fixed, the full centralizer and real-rank proof
EF.4 apply with all their cochains. The completed fixed-curve
obstruction EF.5 excludes its negative automorphisms. The entire
finite-source, selected-fan, integral source-generator, Miller and
root divisor proof EF.6 applies verbatim: its hypotheses are the
source fibre types of EB.1 and \(s\ne0\), including the separately
checked \(s=1\) point/O specialization. Both exact cap and annular
arrows EF.7 retain every norm sign and finite character.
Consequently
\[
              \operatorname{Aut}_{\rm hol}(X)
                    =\mathbb C^*_{\rm physical}
 \tag{EB.9}
\]
on EB.1, with the same actual injection and inverse-endpoint isotopy.

The scalar endpoint construction SEF.1-\/-4 uses those same source,
Cartier, root and reference arrows. Its integer winding restriction
uses precisely EB.3 and the full period intertwiner. The NE
selected vertex maps, genuine root cover, and retained full
annular integer \(n+m\nu\) are unchanged. Thus the complete
endpoint maps and winding groups SEF/NE.11 extend to EB.1 by these
literal identifications. The points \(r=0,1/4,1/3\) have different
actual source degenerations; no endpoint at those sources is
assigned by this note.

% End expanded verified cp36 input


\par\bigskip\noindent\textbf{Complete supporting proof 2: EBR.}\par


% Expanded verified cp36 input: notes/part3_e6_intrinsic_base_ebr_proof.tex
\par

\medskip\noindent\textbf{Root proof of the intrinsic base and the full noncollision endpoint locus}

\par\nobreak

We use the original source
\[
 E_r:\ y^2=x^3+rx^2+tx+t^2,\qquad G=(0,t),\qquad
 s=\tfrac14-r,\qquad
 \mathcal V=\mathbb C\setminus\{0,\tfrac14,\tfrac13\}.
 \tag{EBR.1}
\]
The completed threefold, its point line, its selected simple node,
its integral basis fan, and its entire annular reference are the
objects of EF.1, with either
\((a,b;m)=(2,3;3l)\), \(l\geq1\), or
\((a,b;m)=(6,1;3l+2)\), \(l\geq0\), and \(n=\pm1\).
Every endpoint in this proof is a holomorphic isomorphism between
the specified completed endpoints. The correction from ``map'' to
``isomorphism'' in EB.1 is retained with its original source and
certificate under the recorded preservation directory.

\par

\medskip\noindent\textbf{EBR.1. The marked parameter cover and the actual topology comparison}

\par\nobreak

Write
\[
 F_r(t)=27t^2+(4-18r)t+r^2(4r-1).
\]
The cover which chooses the selected simple node is exactly
\[
 w\longmapsto
 \left(r=\frac{1-w^2}{3},
 \quad t_c=\frac{(w-1)^2(2w+1)}{27}\right),
 \qquad
 w\in\mathbb C\setminus\{0,1,-1,\tfrac12,-\tfrac12\}.
 \tag{EBR.2}
\]
Its inverse on \(F_r(t_c)=0\), \(r\in\mathcal V\), is
\[
 w=\frac{54t_c-18r+4}{4(1-3r)}.
 \tag{EBR.3}
\]
Substitution proves both composites are the identity. In particular
the other selected-node choice is \(w\mapsto-w\), with the same
source \(r\). The displayed punctured plane is path connected:
a polygonal path between two points can be perturbed in finitely
many disjoint small discs to avoid its five excluded points.
This proves the precise connectedness assertion used in EB.1.

Along such a compact path the two simple nodes remain distinct
and nonzero; their exact discriminant is \(16(1-3r)^3\).
At either node the parameter-dependent analytic Morse coordinates
give the same original height \(q=t-t_c\). The I2 coordinate
charts retain \(s\ne0\). The infinity root has equation
\[
 Y^2=X^3+r z^2X^2+zX+1,
 \quad q_\infty=z^3,\quad x=z^{-2}X,\quad y=z^{-3}Y,
 \tag{EBR.4}
\]
and its central curve and deck weights are unchanged. The higher
good and reduced roots, their point/O graph maps, and their complete
chord lifts are precisely EF.32-\/-38, with the same varying \(r\).
All finite characters and the chosen reference cochain are continued
through those actual charts.

The strict contraction inequalities are open on each compact norm
shell. A finite cover of the compact path gives one smaller nonzero
scalar for every chart in that path. The source, root, selected-fan
and annular squares then glue a proper smooth family of the completed
spaces over the path. An explicit smooth comparison is obtained by
lifting the path tangent field in these charts and patching the lifts
with a partition of unity on the common compact shell. On each glued
boundary the lifts are transported by its actual gluing differential;
thus the patched field descends to the completed family. Its flow
exists throughout the compact parameter interval, with inverse given
by the reverse field, and identifies its endpoint fibres. The
integral F35/F46 maps give the sphere at the chosen EF.5 starting
source and every prescribed integral fan. Transport the prescribed
fan marking backwards along the path to that starting source; the
same flow then applies to that marking. Radial transport within the
strict contraction domain reaches the required scalar. Consequently
the actual completed objects throughout EBR.1 retain
\(H^2(X;\mathbb Z)=H^4(X;\mathbb Z)=0\).

The horizontal-divisor argument EF.7 uses exactly this vanishing
and a flat Kähler form on a regular two-torus fibre. Its equality
between the divisor class pairing and positive divisor area excludes
horizontal polar divisors. The transverse map at every special
component is \(u\mapsto u^e=t-t_0\); invariance under its \(e\)
branches makes the descended Laurent exponents integral. Thus the
specified pullback identifies \(\mathbb C(X)\) with
\(\mathbb C(t)\). An endpoint isomorphism \(F\) therefore satisfies
\(\pi_{\rm target}F=h\pi_{\rm source}\) for an actual Möbius map
\(h\), including at the special fibres by holomorphic continuation.
The derivative of this equation and its inverse show that \(h\)
permutes the critical values and carries the regular locus onto
the regular locus.

\par

\medskip\noindent\textbf{EBR.2. The complete lattice gives the original elliptic morphism}

\par\nobreak

Retain the entire period module, including the source-dependent
suspension logarithm and the integer central cochain:
\[
 \gamma=(-bg,\ell),\quad e_1=(1,0),\quad
 e_2=(\tau,ag),\quad \delta=(0,1),\qquad
 M_\alpha=
 \begin{pmatrix}
 1&0&0\\ bv_\alpha&A_\alpha&0\\
 c_\alpha&a v_\alpha^tJA_\alpha&1
 \end{pmatrix}.
 \tag{EBR.5}
\]
The periods are mapped to the point-Cartier line by the full
\(V_{\rm pt}=f_{P^a}(Z)v_M\) arrow EF.10, with the constant
in EF.12 and the reciprocal module maps at both divisors.
No coefficient is replaced by a central value.

The nonzero rational Abel-class proof EF.4 applies at every
parameter in EBR.1. At the finite multiplicative fibres an invariant
period is the extending annular period. At infinity the actual
logarithm of \(3G'\) is \(-sz^2/2+O(z^5)\); division by \(z^2\)
gives the original extending Lie frame. A rational coboundary would
therefore give a global section of the Lie line \(\mathcal O(-1)\).
The explicit coefficient transition
\(u_0(q)=q^{-1}u_t(q^{-1})\) makes this section zero. The original
central additive coordinate of \(3G\) is \(-s/2\ne0\), so its
nonzero integer multiples cannot be the zero section. This proves
the required nonzero Abel class on the whole stated domain.
The infinity elliptic matrix has no invariant vector. In EBR.5 an
invariant vector with \(\gamma\)-coefficient \(z\) would satisfy
\((A_\alpha-I)u+bzv_\alpha=0\); hence \(z=0,u=0\).
The dual calculation, using the retained
\(a v_\alpha^tJA_\alpha\), gives the dual invariant line.
Primitivity of the displayed lattice generators yields exactly
\[
 \Lambda^{\pi_1}=\mathbb Z\delta,\qquad
 (\Lambda^*)^{\pi_1}=\mathbb Z\psi,\qquad
 \psi(\gamma)=1,\quad\psi(e_1)=\psi(e_2)=\psi(\delta)=0.
 \tag{EBR.6}
\]

The endpoint intertwines this local system with its pullback by
\(h\). The invariant subgroups in EBR.6 are preserved under this
intertwining, without assuming \(h=1\). Thus
\(F_*\delta=\epsilon\delta\) and \(F^*\psi=\epsilon'\psi\),
with \(\epsilon,\epsilon'\in\{1,-1\}\). In particular it carries
\[
 H_t=\ker\psi_t=\langle e_1,e_2,\delta\rangle_{\mathbb Z}
 \quad\hbox{bijectively to}\quad H_{h(t)}.
 \tag{EBR.7}
\]
The actual corresponding cyclic cover is
\(\mathbb C^2/H_t=M_a^\times|_{E_{r,t}}\), by the point-theta
arrow just specified.

Lift the endpoint on a regular fibre to its universal cover.
The derivative of that lift is periodic for the full rank-four
lattice; each matrix entry therefore descends to a holomorphic
function on a compact complex torus and is constant. The lift is
affine. Complex linearity and \(F_*\delta=\epsilon\delta\)
give its exact form
\[
 (z,w)\longmapsto
 (A z+z_0,\ C z+\epsilon w+w_0),\qquad A\ne0.
 \tag{EBR.8}
\]
Its inverse is
\[
 (z',w')\longmapsto
 \left(\frac{z'-z_0}{A},
 \epsilon\left(w'-w_0-C\frac{z'-z_0}{A}\right)\right).
 \tag{EBR.9}
\]
The full second-coordinate constants remain in these arrows.
On the physical nonzero-line fibre the multiplicative map is
\(V'\!=e^{2\pi i(Cz+w_0)}V^\epsilon\), interpreted in its
actual source and target point-theta frames.

The lattice map on EBR.7 has the exact integral form
\[
 \begin{pmatrix}p&q&0\\ r_0&s_0&0\\ u&v&\epsilon\end{pmatrix},
 \qquad
 A=p+r_0\tau(h(t)),\quad
 A\tau(t)=q+s_0\tau(h(t)).
 \tag{EBR.10}
\]
It is unimodular. The first-coordinate complex map preserves real
orientation, so \(ps_0-qr_0=1\). These equalities prove
\[
 A\langle1,\tau(t)\rangle_{\mathbb Z}
 =\langle1,\tau(h(t))\rangle_{\mathbb Z}.
 \tag{EBR.11}
\]
Consequently \([z]\mapsto[Az+z_0]\) and its displayed inverse
are degree-one holomorphic maps between the original elliptic
source fibres. They are induced by EBR.8-\/-9 through the literal
cyclic cover and the physical \(\mathbb C^*\) projection.
The central integers \(u,v,c_\alpha\) remain in the lifted line
map. This supplies both directions of the geometric comparison
needed to use the elliptic invariant.

\par

\medskip\noindent\textbf{EBR.3. The exact invariant fixes the entire original base}

\par\nobreak

Translate the cubic variable by \(x=X-r/3\). Its exact short
Weierstrass coefficients are
\[
 p=t-r^2/3,\qquad q=t^2-rt/3+2r^3/27.
 \tag{EBR.12}
\]
The formula \(1728\,4p^3/(4p^3+27q^2)\) therefore gives
\[
 j_r(t)=
 -\frac{256(r^2-3t)^3}{t^2F_r(t)}.
 \tag{EBR.13}
\]
This calculation is independently replayed from the translated
cubic in the companion program. EBR.8-\/-11 give the actual
elliptic isomorphisms, hence \(j_r\circ h=j_r\).

At zero the Laurent coefficient is
\[
 \lim_{t\to0}t^2j_r(t)=-\frac{256r^4}{4r-1}\ne0.
 \tag{EBR.14}
\]
The other poles, the two distinct roots of \(F_r\), are simple.
Indeed the only finite numerator zero is \(t_z=r^2/3\), and
\[
 F_r(t_z)=r^2(3r-1)^2/3,\qquad
 t_z^2F_r(t_z)=r^6(3r-1)^2/27,
 \tag{EBR.15}
\]
both nonzero on EBR.1. Thus \(t_z\) is the unique triple zero.
At infinity the original coordinate gives
\[
 j_r(1/q_\infty)=256q_\infty+O(q_\infty^2),
 \tag{EBR.16}
\]
so infinity is the unique simple zero. Pullback by a Möbius map
preserves each local order, by substitution of its invertible
local coordinate. Therefore \(h\) fixes zero, infinity and
\(r^2/3\). The first two equalities give \(h(t)=ct\); the
third, with \(r\ne0\), gives \(c=1\).

This argument includes the actual auxiliary-polynomial zeros
excluded in EF.3. As a concrete branch-permutation candidate,
\(r=2/9\) has \(F_r(t)=27t^2-4/729\); the map \(t\mapsto-t\)
permutes its two simple poles but moves the unique triple zero
\(4/243\). Substitution into EBR.13 also gives an explicit
nonzero rational difference. The geometric quotient above,
followed by these exact orders, excludes that candidate on the
completed objects.

\par

\medskip\noindent\textbf{EBR.4. Every remaining endpoint map retains its full domain}

\par\nobreak

After \(h=1\), the entire EF.18 centralizer computation applies
to EBR.5, including \(c_\alpha\). It gives
\(eI+c\delta\psi\); complex linearity removes its real rank-one
term and gives the two signs. For an automorphism of one completed
object, the negative sign has sixteen transverse fixed points in
every regular fibre. Their proper analytic curve closure has
positive degree over the original base, whereas pairing with
\(\pi^*[\omega_{\mathbb P^1}]=0\) gives zero. The full completed
fixed-set argument EF.5 thus excludes this sign throughout EBR.1.

For the positive sign, the same actual cyclic-cover lifts and
selected-vertex unit arrows EF.24-\/-25 give a global original
source section. Its component labels remain those carried by
the finite I2 and infinity root charts. The full ten-class Gram
matrix EF.26a has determinant \(-1\); the two disjoint section
charts and the E6/A1 correction columns give height \(1/6\).
Thus the source section is \(jG\), \(j\in\mathbb Z\), on all
of EBR.1. Its actual line lift is \(H(t)U_G^j\).

No discarded auxiliary polynomial enters its three divisor
orders. At the I2 the M2 full coefficient is
\(t(\beta-1-\alpha/2)/(-s+t\alpha)\). The M6 full coefficient,
after the exact original \(t^3\) factor, is
\(8((1-s)w-s)/s^3\). Its separately computed \(s=1\)
point/O specialization is \(-8\), with the original two
Cartier equations retained. The root maps EF.32-\/-35 have
generic inverse orders \(2m/b\) and all their point/O correction
monomials. At the selected node the actual \(B^k\) reindex
gives its integer order \(-k\).
The full orders of \(H\) are therefore \(3j/b,-2j/b,-k\)
at zero, infinity and the selected node. Its divisor degree
gives \(j=bk\) and
\[
 H=c(t^3/q)^k,\qquad H U_G^{bk}=cB_{\rm sel}^{-k}.
 \tag{EBR.17}
\]
On the quotient this is precisely multiplication by
\(c\lambda^k\). It commutes with every actual cap and annular
arrow EF.34-\/-38. The additive value \(-s/2\) also proves the
physical scalar injection has trivial kernel. Thus
\(\operatorname{Aut}_{\rm hol}(X)=\mathbb C^*_{\rm physical}\)
on the entire stated locus.

For scalar endpoints, the full ordinary intertwiner and its
two complex-independent vectors \(e_1,\delta\) give SEF.21:
the two original scalars are equal and the path winding \(J\)
is integral. The root-side comparison is the full
\(z^{nJ}\) in SEF.5. The third divisor equation becomes
\[
 m\,\operatorname{ord}_\infty H+2mj/b+nJ=0.
 \tag{EBR.18}
\]
Together with the other two actual orders it gives
\(j=bk+bnJ/m\), hence \(m\mid bJ\).
For M2 the full primitive forward maps are
\[
 (R^{-l}tU_G)^n\quad\hbox{on the ordinary piece},\qquad
 (R^{-l}U'_G)^n\quad\hbox{on the root piece},
 \quad R^{3l}=\lambda .
 \tag{EBR.19}
\]
The complete SEF.14 Cartier comparison contains \(z^{2l}\);
its product with \(z^{nl}\) and \(t^n=z^{-3ln}\) is one.
This is the actual annular endpoint square, for either signed
inverse and every finite character. Its three chronological
composites are \(q^n\) and \(h^n\), with
\(h=1-t_cz^m\), as in SEF.20. For M6 the primitive maps
are these \(q^n/h^n\) arrows themselves.

The negative scalar-endpoint obstruction uses the actual
ordinary and root covers of NE.2-\/-3. The complete reference
monodromy on one root circle is \(B_\lambda^{n+m\nu}\),
with its retained integer \(\nu\). The two single-valued
negative lifts would give
\(B_\lambda^{2(n+m\nu)}=1\); strict contraction and
\(n=\pm1,m\ge2\) exclude this equality. The proof applies
to the quadratic M6 cap as well. All root, selected-vertex
and Cartier charts used there have just been retained over
EBR.1. Hence the full scalar winding groups are
\(l\mathbb Z\) for M2 and \(m\mathbb Z\) for M6, with the
explicit primitive arrows above and physical scalar ambiguity.
No scalar return class is computed or divided by this endpoint
classification. The three removed source parameters are outside
the scope of every assertion in this proof.

% End expanded verified cp36 input


\par\bigskip\noindent\textbf{Complete supporting proof 3: EBP.}\par


% Expanded verified cp36 input: notes/part3_e6_intrinsic_base_ebp_proof.tex
\par

\medskip\noindent\textbf{Independent intrinsic-base and endpoint-domain audit}

\par\nobreak

This peer proof concerns the complete EB and EBR constructions and their
use of EF, SEF and NE. The source is
\[
 y^2=x^3+rx^2+tx+t^2,\qquad G=(0,t),\qquad
 s=\tfrac14-r,\qquad r(4r-1)(3r-1)\ne0.
 \tag{EBP.1}
\]
The point line, selected simple node, ordered fan, scalar, norm sign,
finite character and complete annular reference are the original ones.
The cap data are \((a,b;m)=(6,1;3l+2)\), \(l\ge0\), and
\((2,3;3l)\), \(l\ge1\), with \(n=\pm1\). Every endpoint in
this argument is a holomorphic isomorphism of the specified completed
endpoints. Its inverse is retained in each quotient construction below.

\par

\medskip\noindent\textbf{EBP.1. The actual marked cover and topology transport}

\par\nobreak

Put \(F_r(t)=27t^2+(4-18r)t+r^2(4r-1)\). The selected-node cover is
\[
 r=(1-w^2)/3,\qquad t_c=(w-1)^2(2w+1)/27,
 \qquad
 w=\frac{54t_c-18r+4}{4(1-3r)}.
 \tag{EBP.2}
\]
Direct substitution gives both inverse composites, using \(F_r(t_c)=0\)
for the second. The original excluded product pulls back exactly to
\[
 -\frac{w^2(w^2-1)(4w^2-1)}9.
 \tag{EBP.3}
\]
Thus the parameter cover is the plane with precisely
\(0,\pm1,\pm\tfrac12\) removed. Its polygonal paths, detoured in
disjoint discs around these five points, join the concrete EF.5 source
to either marked branch of every object in EBP.1. The involution
\(w\mapsto-w\) gives the other root of the same quadratic; their sum
and product are the exact coefficients of \(F_r\).

Along a compact such path, the two selected candidates are distinct and
nonzero, since the discriminant is \(16(1-3r)^3\) and the constant
coefficient is \(r^2(4r-1)\). Analytic Morse charts keep the actual
height \(q=t-t_c\). The I2 charts use the unit \(s\). At infinity the
good cubic source root is literally
\[
 q_\infty=z^3,\qquad x=z^{-2}X,
 \quad y=z^{-3}Y,\qquad
 Y^2=X^3+rz^2X^2+zX+1.
 \tag{EBP.4}
\]
The higher roots are the full EF.32-\/-38 maps reviewed in EBP.6 below.
Their characters and reference integers are continued as part of those
maps. Transport the prescribed terminal fan backwards to the starting
object and then use these actual cone and period identifications.

A common smaller scalar gives strict contraction on a finite collection
of compact norm shells covering the path. In each source/root/fan chart,
lift the path vector field; on common boundary charts use the displayed
gluing derivative to identify the lifts. A partition of unity on the
compact glued family preserves the projection of the field to the path
tangent. Its flow and reverse flow identify the completed endpoint
fibres. This is the precise smooth comparison used with the original
F35/F46 integral topology maps. Radial scalar transport uses the same
full cap and reference square. Hence the resulting completed spaces
retain the sphere cohomology, including \(H^2=H^4=0\), throughout
EBP.1. The argument carries each prescribed discrete marking from the
corresponding initial object.

The horizontal-divisor pairing EF.7 now excludes every horizontal polar
divisor on a regular compact two-torus fibre. Removal across its remaining
codimension-two indeterminacy makes a meromorphic function constant on
that fibre. At each special component the actual transverse map
\(u\mapsto u^e=t-t_0\), after its extending unit root, forces the
Laurent exponents to be multiples of \(e\). The descended function is
therefore meromorphic at the special base value. This proves the precise
pullback equality \(\mathbb C(X)=\mathbb C(t)\). An endpoint and its
inverse induce an automorphism of this rational function field and give
\(\pi_{\rm target}F=h\pi_{\rm source}\) for a Möbius map. Taking
derivatives in that equality and its inverse carries the regular locus
onto the regular locus. No assertion that \(h\) fixes the base has yet
been used.

\par

\medskip\noindent\textbf{EBP.2. The two primitive invariants before base rigidity}

\par\nobreak

In the actual point-theta coordinates, retain
\[
 \gamma=(-bg,\ell),\quad e_1=(1,0),\quad e_2=(\tau,ag),
 \quad\delta=(0,1),\qquad V_{\rm pt}=f_{P^a}(Z)v_M.
 \tag{EBP.5}
\]
The full constant relating this coefficient to the tensor coefficient,
the two Cartier generators and their inverse maps are EF.10-\/-17.
In particular \(\ell\) is the logarithm of the entire suspension
coefficient. It retains the scalar, original selected height, Miller
function and point/O factors.

The monodromy, with \(J=\left(\begin{smallmatrix}0&-1\\1&0\end{smallmatrix}\right)\), is
\[
 M_\alpha=\begin{pmatrix}
 1&0&0\\ bv_\alpha&A_\alpha&0\\
 c_\alpha&a v_\alpha^tJA_\alpha&1
 \end{pmatrix}.
 \tag{EBP.6}
\]
Multiplication gives the full central term
\(c_{\alpha\beta}=c_\alpha+c_\beta+
abv_\alpha^tJA_\alpha v_\beta\), as well as
\(v_{\alpha\beta}=v_\alpha+A_\alpha v_\beta\). The peer
calculator verifies this identity with symbolic central coefficients.

Here is the domain of the nonzero Abel-class argument. A rational
coboundary, after clearing denominators and multiplying by six, gives
a single-valued Lie logarithm of an integer multiple of \(G\) lying in
every identity component. At the finite multiplicative fibres the local
Tate coordinate is an extending nonzero unit. The difference between
its logarithm and the global logarithm is an invariant annular period,
which extends in the same Lie frame. At infinity the actual chords give
\[
 x(3G)=4t^2/s^2+2t/s,\qquad
 y(3G)=8t^3/s^3+6t^2/s^2+t/s-t.
 \tag{EBP.7}
\]
On EBP.4 the formal parameter has leading term \(-sz^2/2\).
The original minimal differential is \(z^{-2}\omega_{\rm good}\),
whereas the affine differential is \(z\omega_{\rm good}\).
Dividing that root logarithm by \(z^2\) gives the extending original
Lie coefficient. The infinity elliptic monodromy has polynomial
\(x^2+x+1\), so it supplies no invariant-period ambiguity there.

The full Lie coefficient transition is
\(u_0(q_\infty)=q_\infty^{-1}u_t(q_\infty^{-1})\). Its global
holomorphic section is zero, because the entire function \(u_t\)
decays at infinity. But the actual original central point has
\((x_0,y_0)=(4/s^2,8/s^3)\), so its additive coordinate
\(-x_0/y_0=-s/2\) and all its nonzero integer multiples are nonzero.
The differential identity on \((x_0,y_0)=(u^2,u^3)\) is
\(dx_0/(2y_0)=d(-1/u)\). This proves the nonzero rational Abel class
for every parameter in EBP.1, using only \(s\ne0\) and the retained
fibre types.

An invariant vector with first coefficient \(z\) satisfies
\((A_\alpha-I)u=-bzv_\alpha\). The nonzero Abel class gives
\(z=0\); the infinity matrix gives \(u=0\). Its last coefficient
is free. For the dual equation write the middle row as \(w^tJ\).
The exact identity
\[
 w^tJ(A-I)=(w-Aw)^tJA
 \tag{EBP.8}
\]
turns the dual invariant equation into the same Abel coboundary
equation. Its final coefficient is zero and its middle row is zero.
Consequently the full integral invariant groups are precisely
\[
 \Lambda^{\pi_1}=\mathbb Z\delta,
 \qquad (\Lambda^*)^{\pi_1}=\mathbb Z\psi,
 \qquad\psi(\gamma)=1,
 \quad\psi(e_1)=\psi(e_2)=\psi(\delta)=0.
 \tag{EBP.9}
\]
The displayed basis proves primitivity. The equations include the entire
central cochain; it has not been assigned a different value.
An endpoint covering \(h\) intertwines this local system with its
pullback. It carries these intrinsic invariant groups bijectively,
so it sends \(\delta\) and \(\psi\) to their signed primitive
generators. This statement precedes the determination of \(h\).

\par

\medskip\noindent\textbf{EBP.3. The actual cyclic cover, affine arrow and elliptic inverse}

\par\nobreak

The subgroup
\(H=\ker\psi=\langle e_1,e_2,\delta\rangle\) defines the
literal cyclic cover \(\mathbb C^2/H=M_a^\times|_{E_t}\), through
the full point-theta and inverse Cartier-module arrows in EBP.5.
The endpoint carries \(H\) bijectively to the corresponding target
subgroup. Lift it to the universal covers. Its derivative is periodic
for the full rank-four lattice; every entry descends to a holomorphic
function on the compact torus and is constant. The lift is affine.
Complex linearity and the primitive physical period give its exact form
\[
 (z,w)\longmapsto(Az+z_0,Cz+\epsilon w+w_0),\qquad
 \epsilon\in\{1,-1\},\quad A\ne0,
 \tag{EBP.10}
\]
with inverse
\[
 (z',w')\longmapsto
 \left((z'-z_0)/A,
 \epsilon\{w'-w_0-C(z'-z_0)/A\}\right).
 \tag{EBP.11}
\]
In multiplicative point coefficients the full line arrow is
\(V'=\exp(2\pi i(Cz+w_0))V^\epsilon\). Its target and source
Cartier frames are the ones already specified; their constant factors
and integral transitions remain in \(C,w_0\) and the period rows.

The map on the kernel lattice is an integral matrix
\[
 \begin{pmatrix}p&q&0\\r_0&s_0&0\\u&v&\epsilon\end{pmatrix},
 \quad A=p+r_0\tau',\quad A\tau=q+s_0\tau'.
 \tag{EBP.12}
\]
It is unimodular. The area identity
\[
 \operatorname{Im}\bigl(\overline{p+r_0\tau'}(q+s_0\tau')\bigr)
 =(ps_0-qr_0)\operatorname{Im}\tau'
 =|A|^2\operatorname{Im}\tau
 \tag{EBP.13}
\]
gives \(ps_0-qr_0=1\). Thus the quotient period lattices agree
exactly. The maps \([z]\mapsto[Az+z_0]\) and EBP.11\textquotesingle s first
coordinate are inverse holomorphic isomorphisms of the original
elliptic source fibres. They are obtained through the specified
cyclic-cover and physical-fibre projections. The integers \(u,v\)
and the full central cochain remain in the line map, where they belong.

\par

\medskip\noindent\textbf{EBP.4. Exact base rigidity on the whole stated locus}

\par\nobreak

Translating \(x=X-r/3\) gives the short cubic coefficients
\(p=t-r^2/3\), \(q=t^2-rt/3+2r^3/27\). Direct computation gives
\[
 c_4=16(r^2-3t),\qquad\Delta=-16t^2F_r(t),\qquad
 j_r(t)=-\frac{256(r^2-3t)^3}{t^2F_r(t)}.
 \tag{EBP.14}
\]
The actual source isomorphisms EBP.10-\/-13 prove \(j_r\circ h=j_r\).
Its double-pole coefficient at zero is \(-256r^4/(4r-1)\).
The other two poles are simple. The unique finite zero is the triple
zero \(t_z=r^2/3\), since
\[
 F_r(t_z)=r^2(3r-1)^2/3,\qquad
 t_z^2F_r(t_z)=r^6(3r-1)^2/27.
 \tag{EBP.15}
\]
Infinity is the unique simple zero, with coefficient
\(j_r(1/q_\infty)=256q_\infty+O(q_\infty^2)\).
All these nonzero coefficients are valid on EBP.1. An invertible local
coordinate preserves the order of a zero or pole, so \(h\) fixes
zero, infinity and \(r^2/3\). The first two give \(h(t)=ct\), and
the third gives \(c=1\).

This includes the auxiliary centroid-polynomial zeros from EF.3.
For example, at \(r=2/9\), \(F_r(t)=27t^2-4/729\); its pole
permutation \(t\mapsto-t\) moves \(t_z=4/243\). The exact
rational functions \(j_r(-t)\) and \(j_r(t)\) differ. The source
isomorphism was established before making this invariant comparison.

\par

\medskip\noindent\textbf{EBP.5. The remaining automorphism hypotheses and source generator}

\par\nobreak

After the base is fixed, the full centralizer calculation is
\(eI+c\delta\psi\). Its commutator with EBP.6 vanishes for every
retained central coefficient. The real rank-one remainder is eliminated
by complex linearity, and unimodularity gives \(e=\pm1\). For an
automorphism of a completed object, the negative sign gives sixteen
transverse fixed points per regular two-torus. Their closed analytic
equalizer in the compact completed space has a horizontal compact curve
component. Its positive base degree contradicts the already transported
\(H^2=H^4=0\). This is the full completed fixed-curve argument EF.5.

For the positive sign, the ordinary cyclic-cover extension, genuine
root pullback, and selected vertex unit maps EF.24-\/-25 give the actual
global source section. All vanishing cycles lie in \(\ker\psi\);
the finite cap uses \(A^m=B^n\) and strict contraction to remove
the commutator integer. The lift torsor over the one-puncture orbifold
base has no homomorphism from its finite cyclic fundamental group to
the suspension deck group. These are the explicit extension steps,
with both signs and all reference integers retained.

The rational source chart is
\(U=y-t-x/2\), \(V=y+t+x/2\), \(UV=x^2(x-s)\), with rational
inverse \(t=(V-U-x)/2\), \(y=(V+U)/2\). The I2 tangent cone
\(UV=-sx^2\) and the two distinct infinity sections retain the
component labels used in EF.26a. The peer calculator reconstructs the
full ten-class Gram matrix from those labels and obtains determinant
\(-1\). The E6 correction is \(4/3\), the A1 correction is
\(1/2\), and the height of \(G\) is \(1/6\). Hence the actual
global section is an integral multiple \(jG\), with the same
primitive generator and no omitted section torsion.

The full endpoint lift is \(H(t)U_G^j\). For M2 the exact finite chart
gives \(B_{2,1}/t=2v/(s-vh)\). For M6 its coefficient after the
actual \(t^3\) factor is \(8((1-s)w-s)/s^3\), which is not
identically zero for any \(s\ne0\). The special value \(s=1\)
is checked from the full chord before taking a limit. The independent
calculation retains the actual tangent at \(3G\), the doubled point,
and both translated quotients, and obtains
\[
 \frac{f_3(p-G)}{h f_3(p)}\longrightarrow
 \frac{w^3}{(w-1)(w^2-w+1)},\qquad
 \frac{f_{2,3G}(p-G)}{h f_{2,3G}(p)}\longrightarrow
 -\frac{(w-1)^2}{w^2-w+1}.
 \tag{EBP.16}
\]
Their complete product gives \(B_{6,1}/t^3\to-8\), and
\(\tau(6G)/t\to-6\). The two actual point/O Cartier equations
are retained at that coincidence. The Hom divisor and its inverse
therefore extend through the resolved chart by the same normality
argument as EF.30g. This explicitly covers the only extra specialization
in that finite-source factor calculation.

The three original divisor orders of \(H\) are
\(3j/b,-2j/b,-k\) at zero, infinity and the selected node.
Their degree gives \(j=bk\) and the literal equality
\(H U_G^{bk}=cB_{\rm sel}^{-k}\). On the actual quotient this
is multiplication by \(c\lambda^k\). The central additive value
from EBP.7 makes the physical scalar injection faithful. Its full
source/root/fan extension and inverse are the EF.34-\/-38 maps.

\par

\medskip\noindent\textbf{EBP.6. Both root families and every downstream endpoint square}

\par\nobreak

The higher-root substitutions retain the following complete exponents:
\[
\begin{array}{c|c|c}
(a,b;m)&(q_\infty,x_0,y_0)&\text{root equation}\\ \hline
(6,1;3l+2)&(z^{3l+2},z^{4l+2}X,z^{6l+3}Y)
 &Y^2=X^3+rz^{2l+2}X^2+z^{l+2}X+z^2\\
(2,3;3l)&(z^{3l},z^{4l}X,z^{6l}Y)
 &Y^2=X^3+rz^{2l}X^2+z^lX+1.
\end{array}
 \tag{EBP.17}
\]
The reduced source division and all graph rays are EF.32; the good
division and both point/O correction sums are EF.33. The parameter
\(r\) is an ordinary holomorphic coefficient in these equations.
The formal sections and normalized chord lifts use \(s\ne0\) and
the same original central points. Their full norm unit has central
value four. The prescribed root, character \(\zeta_m^d\), signed
norm and entire reference give the same \(A_\lambda\), \(B_\lambda\)
and sector maps EF.35-\/-38. The nonzero selected factor is exactly
\(h=1-t_cz^m=q_\infty q\), with value one at the root centre.
No centroid polynomial occurs in any of these maps.

For SEF the actual finite map is \(S=tU_G\). On the good root,
\(f_{2,\rm old}=z^{-l}\phi\),
\(U_{G,\rm old}=z^{2l}U'_G\), with
\(\phi=(Y-1-z^lX/2)/(X-sz^{2l})\).
Its central product is \(-X\) in the retained meromorphic point
frames, with exactly their Cartier divisor. The actual graph correction
\(E_P-E_O\) remains in both directions of that equality.
The full deck calculation gives
\(U'_G D=\zeta_m^{-2l}D U'_G\). Consequently the primitive
forward maps are
\[
 (R^{-l}tU_G)^n\quad\hbox{and}\quad(R^{-l}U'_G)^n,
 \tag{EBP.18}
\]
with the directed finite square SEF.18. For \(n=-1\) the
coefficients are evaluated at the actual inverse source point.

The full reference restoration SEF.3 retains every sector integer.
At integral winding \(J\) its endpoint ratio is exactly
\(z^{nJ}\), since its additional scalar factor on an integral
suspension label is one. For the primitive good endpoint the three
annular powers are \(nl,-3nl,2nl\); their sum is zero. The three
chronological endpoint maps retain the successive root scalars and
compose to \(q^nB_\lambda^{-n}\) and \(h^nB_\lambda^{-n}\).
All finite characters commute through these squares and retain their
specified norm. This proves the primitive and iterated maps on the
entire domain EBP.1.

For an arbitrary endpoint the complete ordinary periods first force
the two original scalars to agree and their path winding to be an
integer \(J\). The three full divisor equations then give
\(j=bk+bnJ/m\), hence \(m\mid bJ\). With the explicitly
constructed primitive maps this yields \(l\mathbb Z\) for M2
and \(m\mathbb Z\) for M6.

The negative-sign exclusion uses the full NE construction. Its selected
inversion carries the two actual vertex charts by NE.4, preserving
their product \(q\), their full unit coefficients and their inverse.
Composition with that inversion supplies the single-valued ordinary
cyclic-cover lift across the selected puncture. The original source/O
section and point/O module frames extend the other ordinary lifts.
The genuine root fibre product similarly gives a single-valued root
lift fixing \(z\). Its full annular deck is
\(A_\lambda B_\lambda^\nu\), with the actual carried integer
\(\nu\). Its chronological \(m\)-fold product is
\(B_\lambda^{n+m\nu}\); the same integer occurs at the two
scalar endpoints by the explicit SEF restoration.
Transporting the negative lift around the root circle gives the
conjugate factor \(B_\lambda^{2(n+m\nu)}\). Single-valuedness
and strict contraction imply \(n+m\nu=0\), whereas this integer
is congruent to \(\pm1\) modulo \(m\ge2\). Thus the full NE
obstruction, including the quadratic reduced cap, retains its proof
on EBP.1. The actual winding maps and physical scalar ambiguity are
therefore exactly the stated ones.

\par

\medskip\noindent\textbf{EBP.7. Independent verification and scope}

\par\nobreak

\path{code/part3_fourfold/e6_intrinsic_base_peer.py} imports neither the
author calculator nor the root calculator. It derives the original
cubic invariant, both marked-cover composites, exact zero/pole orders,
full symbolic monodromy product and dual equations, affine inverses,
integral kernel inverses, original source/chord/Lie identities,
ten-class determinant, the full specialized M6 chord ratios, both
root substitutions and the signed endpoint/annular words. Its finite
matrix cases supplement the symbolic identities; the analytic cover,
Lie extension, compact flow and proper fixed-curve arguments are the
proof above. The corresponding JSON records every source hash and
all complete tested kernel matrices.

The completed independent run passes 2,396 exact checks, including 264
complete integral kernel matrices with both inverse products. The
symbolic identities cover the arbitrary central cochain and arbitrary
signed root exponent; the finite executions separately retain all
displayed coefficients, source inverses and endpoint words.

The full geometric read of EB, EBR, EF, SEF and NE passes on EBP.1.
The earlier centroid-polynomial restriction is replaced at its sole
base-rigidity use by EBP.2-\/-4\textquotesingle s actual elliptic quotient and invariant.
Every remaining construction has its original source, Cartier module,
root, period and endpoint hypotheses explicitly retained above.
This audit assigns no new scalar or source return.

% End expanded verified cp36 input



\par\medskip\noindent\textbf{Completed dependencies and exact provenance.}\par
The complete EF.1--8 source, source/root extension and physical-group
proof is P3K.109; the complete SEF.1--6 and NE.1--4 scalar endpoint,
reference and negative-lift proof is P3K.110. Their original sources are
\path{notes/part3_e6_fixed_source_endpoints_working.md},
\path{notes/part3_e6_fixed_source_scalar_endpoints_working.md}, and
\path{notes/part3_e6_fixed_source_negative_endpoint_working.md}.
The three complete proofs printed above carry the extension of every
one of those hypotheses to the stated noncollision locus; EBR.4 and
EBP.5--6 write its source, Cartier, root, character, period and endpoint
maps explicitly. Their completed proof inputs and all mathematical
source hashes are locked in \path{records/k113_export_manifest.json}.

The primary calculation is
\path{code/part3/e6_fixed_source_intrinsic_base.py}; the full root peer
and its independent executable are
\path{notes/part3_e6_intrinsic_base_root_peer.md} and
\path{code/part3_root/e6_intrinsic_base_peer.py}.
The further independent proof and executable are
\path{notes/part3_e6_intrinsic_base_fourfold_peer.md} and
\path{code/part3_fourfold/e6_intrinsic_base_peer.py}.
Their 1098, 5103 and 2396 exact checks retain the full affine inverses,
primitive lattice and dual equations, full source and root coefficients,
signed annular integers and chronological endpoint words. The analytic
cover, Lie extension and compact comparison arguments are the complete
proofs above. The source correction from an arbitrary map to an
endpoint isomorphism and its original predecessor remain in
\path{records/attempts/e6_intrinsic_base_before_isomorphism_scope/}.


% End expanded verified cp36 input

\end{proof}


\subsection{Full holomorphic endpoints on the original type-III E6 source}
\begin{theorem}[Source result P3K.114 under continuation review]\label{night:P3K.114}
\StatementRecord{night:P3K.114}{open}{Historical events: GPT-6 Astra (OpenAI), ChatGPT Pro mode, effort Pro, S6\_06\_BLIND, 16--18 November 2026: targeted proof check. Also checked in this writing. Continuation mathematical and source checks: GPT-6 Astra (OpenAI), ChatGPT Pro mode, effort Pro, S6\_10\_PAPERS continuation, 19 November 2026: computation or formal certificate; targeted proof check; used and rechecked in this run. Continuation metadata events: GPT-6 Astra (OpenAI), ChatGPT Pro mode, effort Pro, S6\_10\_PAPERS continuation, 19 November 2026: used and rechecked in this run; explicitly scoped metadata reconciliation; metadata reconciliation only; no new mathematical check. GPT-6 Astra (OpenAI), ChatGPT Pro mode, effort Pro, S6\_10\_PAPERS continuation, 19 November 2026: used and rechecked in this run; source-row intake; metadata reconciliation only; no new mathematical check. GPT-6 Astra (OpenAI), ChatGPT Pro mode, effort Pro, S6\_10\_PAPERS continuation, 19 November 2026: used and rechecked in this run; receiving source placement; metadata reconciliation only; no new mathematical check. \emph{Pre-continuation recorded proof scope.} Historical event: GPT-6 Astra (OpenAI), ChatGPT Pro mode, effort Pro, S6\_06\_BLIND, 16--18 November 2026: targeted proof check, recorded producing-run re-derivation. Also checked in this writing by GPT-6 Astra (OpenAI), ChatGPT Pro mode, effort Pro, S6\_10\_PAPERS continuation, 19 November 2026: computation or formal certificate for the explicitly listed algebraic substeps in the event ledger; full continuation rederivation and renewed credit comparison remain pending.}{\cref{night:P3K.96,night:P3K.109,night:P3K.110,night:P3K.113}; the original source, point line, roots, characters, reference, selected fan and scalar range stated here}
Retain the original marked type-III source
\[
 E_0:\ y^2=x^3+tx+t^2,\qquad G=(0,t),\qquad
 t_c=-4/27,\quad q=t+4/27,
\]
the indicated simple selected node with its actual branch marking,
and any prescribed lattice-basis fan. Use the genuine point line
\(M_a=\mathcal O([aG]-[O])\), chronological source translation by
\(-bG\), and either complete infinity-cap progression
\[
 (a,b;m)=(6,1;3l+2),\ l\ge0,\qquad
 (a,b;m)=(2,3;3l),\ l\ge1.
\]
Keep both primitive norms \(n=\pm1\), every prescribed finite
character \(\zeta_m^d\), the whole P3K.96 annular reference, and a
nonzero original scalar in the common strict contraction domain.
At the finite III fibre retain the inherited order-one, norm-zero
filling of the P3K.96 family.
Let \(X\) be the resulting completed threefold.

Every holomorphic automorphism fixes the original base pointwise,
and multiplication on the actual point line induces the exact
group isomorphism
\[
       \mathbb C^*_{\rm physical}\xrightarrow{\ \cong\ }
                       \operatorname{Aut}_{\rm hol}(X).
\]
For the full scalar family with \(\lambda=R^m\), every holomorphic
endpoint identification has integral scalar winding, and its
complete winding group and primitive forward maps are
\[
\begin{array}{c|c|c|c}
 (a,b;m)&\text{winding group}&\text{ordinary map}&\text{cap map}\\ \hline
 (6,1;3l+2)&m\mathbb Z&q^n&h^n\\
 (2,3;3l)&l\mathbb Z&(R^{-l}tU_G)^n&(R^{-l}U'_G)^n ,
\end{array}
\]
where \(h=1+4z^m/27=q_\infty q\),
\(U_G(p,v)=(p+G,B_{a,1}(p+G)^{-1}v)\), and
\(U'_G\) is its full good-root counterpart. The point coefficient
is \(B_{a,1}(p)=f_2(p)^a f_a(p-G)/f_a(p)\), with the full
Miller and Cartier arrows specified below. Their actual inverses
are the inverse-endpoint return maps.

The three chronological primitive M2 endpoints have composite
\(q^nB_\lambda^{-n}\) on the ordinary cover and
\(h^nB_\lambda^{-n}\) on the root cover; thus they give precisely
the canonical \(q^n/h^n\) pair on the quotient. Every negative
linear endpoint is excluded by the complete annular deck
\(B_\lambda^{n+m\nu}\), retaining its arbitrary reference integer
\(\nu\). The ratio of any two endpoints with the same winding
is exactly a physical scalar. All source, root, selected-fan,
Cartier and finite-character maps in these assertions are the
displayed original maps. No numerical return is assigned.
\end{theorem}
\begin{proof}[Supplied source argument; continuation review in progress]

% Expanded verified cp36 input: notes/part3_e6_iii_fixed_source_endpoint_proof.tex
This separate continuation keeps K109, K110 and EB unchanged. The source
is the actual \(r=0\) member of the K96 disc family,

\[
 E_0:\quad y^2=x^3+tx+t^2,\qquad G=(0,t),\qquad
 t_c=-4/27,\quad q=t+4/27.
 \tag{R0E.1}
\]

\par

\medskip\noindent\textbf{Retained finite-stage history.}

\par

The following scope paragraph and R0E.1-\/-3 record the completed initial finite-chart stage. Their former global-proof obligation is resolved by the full R0E.4-\/-10 proof reproduced below. The historical scope paragraph is retained verbatim in content.

The following calculation concerns the finite III source and its point
line \(M_2=\mathcal O(2G-O)\). It records the full coefficient before
any central restriction. The completed global period and endpoint
classification will be written separately after these charts are
checked; this note presently assigns no automorphism group or return.

\par

\medskip\noindent\textbf{R0E.1. The genuine quartic source cover and its line coefficient}

\par\nobreak

Use the actual map

\[
 t=z^4,\quad x=z^2X,\quad y=z^3Y,\qquad
 Y^2=X^3+X+z^2,\quad G'=(0,z).
 \tag{R0E.2}
\]

Its deck is \(D(z,X,Y)=(iz,-X,iY)\). The differential is

\[
 \omega_t=\frac{dx}{2y}=z^{-1}\frac{dX}{2Y}
 \tag{R0E.3}
\]

on each fibre. Doubling the actual section gives

\[
 2G=(1/4,-t-1/8),\qquad
 \tau(2G')=-\frac{X(2G')}{Y(2G')}
            =\frac{2z}{1+8z^4}.
 \tag{R0E.4}
\]

The finite point-to-tensor coefficient and its entire pullback are

\[
 f_2=\frac{y-t-x/2}{x-1/4}=z^2\phi,
 \qquad
 \phi=\frac{zY-z^2-X/2}{z^2X-1/4}.
 \tag{R0E.5}
\]

For the actual translation \(T:p\mapsto p-G'\), set

\[
 L=\frac{Y+z}{X},\qquad
 X'=L^2-X=\frac1X+\frac{2z(Y+z)}{X^2},\qquad
 Y'=L(X-X')-Y.
 \tag{R0E.6}
\]

These are rational expressions for the source map, with their actual
Cartier-module extensions. They preserve the whole equation R0E.2.
The entire coefficient of the point-line lift is

\[
 \frac{B_{2,1}}{t}
     =\phi(X,Y,z)\,\phi(X',Y',z).
 \tag{R0E.7}
\]

At \(z=0\), the source arrow is exactly

\[
 (X,Y)\longmapsto(X^{-1},-Y/X^2),\qquad
 \phi|_{z=0}=2X,
 \qquad (B_{2,1}/t)|_{z=0}=4.
 \tag{R0E.8}
\]

The equality is in the Hom line; the zeros and poles of \(2X\)
are retained by the two point/O Cartier modules. The original-source
components require the graph charts below, rather than a replacement
of those modules by the number \(4\).

\par

\medskip\noindent\textbf{R0E.2. Both original III components and the actual point frames}

\par\nobreak

On the strict cuspidal component use \(p=-x/y\), so

\[
 x=p^{-2},\quad y=-p^{-3}\qquad(t=0).
 \tag{R0E.9}
\]

On the other component use the literal resolution chart

\[
 x=t\alpha,\quad y=t\beta,\qquad
 \beta^2=1+\alpha+t\alpha^3.
 \tag{R0E.10}
\]

The two source arrows induced by \(p\mapsto p-G\), before the
point-line factor is read, are

\[
 C_0\longrightarrow C_1:\quad \beta'=p-1,
 \qquad
 C_1\longrightarrow C_0:\quad p'=\beta-1.
 \tag{R0E.11}
\]

For the first arrow keep original \(x=p^{-2}\) fixed and solve the
source equation analytically in \(t\). Its first derivative is

\[
 y_t(0)=-p/2,\quad x(p-G)=t\,p(p-2)+O(t^2),
 \quad y(p-G)=t(p-1)+O(t^2).
 \tag{R0E.12}
\]

The full second arrow has central

\[
 x(p-G)=(\beta-1)^{-2},\qquad
 y(p-G)=-(\beta-1)^{-3}.
 \tag{R0E.13}
\]

Substituting the full rational coefficient R0E.5 at both the original
and the translated source point gives

\[
 \left.\frac{B_{2,1}}t\right|_{C_0}
       =4\frac{p-2}{p},\qquad
 \left.\frac{B_{2,1}}t\right|_{C_1}
       =4\frac{\beta-1}{\beta-3}.
 \tag{R0E.14}
\]

The actual point \(2G\) has \(p=2\); \(O\) is \(p=0\).
The rational point section therefore has the following coefficients
in genuine nonzero line frames:

\[
 r_0=(p-2)/p,\qquad r_1=1,
 \qquad e_i=\sigma_2/r_i.
 \tag{R0E.15}
\]

For a raw coefficient \(B\) covering \(T\), the corresponding
nonzero-frame coefficient is exactly \(B\,r(Tp)/r(p)\).
Equations R0E.11,14,15 give \(4\) on both components. In particular
all four section points \(O,G,2G,3G\) remain in the module maps;
their apparent raw zeros or poles are their displayed generators.

\par

\medskip\noindent\textbf{R0E.3. The source graph that retains those modules}

\par\nobreak

At the root point \(O\), put \(\eta=-X/Y\) and

\[
 U=\eta^2X,\qquad
 U^2-U^3-\eta^4U-z^2\eta^6=0,\qquad U(0,0)=1.
 \tag{R0E.16}
\]

The two charts of the actual blowup \((\eta,z)\) are

\[
 \eta=zp:\quad\tau_{\rm old}=p;
 \qquad z=\eta v:\quad
 (x,y,t)=(v^2U,-v^3U,\eta^4v^4).
 \tag{R0E.17}
\]

The first point generator is exactly

\[
 \frac{p-2/(1+8z^4)}p.
 \tag{R0E.18}
\]

At \(T=(0,0)\) in the central good root, the equation gives

\[
 X(1+X^2)=Y^2-z^2.
 \tag{R0E.19}
\]

Thus the source graph is the blowup \((Y,z)\). On \(Y=z\beta\),

\[
 X=z^2\alpha,\qquad
 \beta^2=1+\alpha+z^4\alpha^3,
 \qquad (x,y,t)=(z^4\alpha,z^4\beta,z^4),
 \tag{R0E.20}
\]

which is exactly R0E.10. The other chart is \(z=Yv\),

\[
 X=Y^2C,\qquad C(1+Y^4C^2)=1-v^2,
 \qquad(x,y,t)=(Y^4v^2C,Y^4v^3,Y^4v^4).
 \tag{R0E.21}
\]

Every formula preserves the original \(t\) and the actual point
divisor. The source translation R0E.6 exchanges these two graph
centres and carries their strict point generators. Its inverse is
obtained by using \(+G'\). These are the graph and Cartier arrows
used for the next full finite-piece extension; the cap at original
infinity remains the complete K96 cap on its own actual domain.

\par

\medskip\noindent\textbf{R0E.4. The full M6 coefficient and the tangential source chart}

\par\nobreak

We now retain both actual data
\[
 (a,b;m)=(6,1;3l+2),\ l\geq0,
 \qquad (a,b;m)=(2,3;3l),\ l\geq1,
 \qquad n=\pm1,
 \tag{R0E.22}
\]
with every finite character \(\zeta_m^d\), the prescribed selected
fan and the whole K96 annular reference. In the second row the
suspension is the chronological third power of the one-step line
map. All root and source formulas in this section use the original
\(r=0\) source R0E.1.

The full Miller functions on that source are
\[
 \begin{gathered}
 s=1/4,\quad U=y-t-x/2,\quad V=y+t+x/2,\quad
 f_2=U/(x-s),\\
 f_3=\frac{s xU}{sV-(s+2t)x},\qquad
 f_6=f_3^2 f_{2,3G},\\
 x_3=8t+64t^2,\quad y_3=3t+96t^2+512t^3,\quad
 \mu_3=\frac{3x_3^2+t}{2y_3},\\
 x_6=\mu_3^2-2x_3,\quad
 y_6=\mu_3(x_3-x_6)-y_3,\qquad
 f_{2,3G}(p)=\frac{y-y_3-\mu_3(x-x_3)}{x-x_6}.
 \end{gathered}                                                   \tag{R0E.23}
\]
The quotients in R0E.23 are coefficients between their actual
Cartier modules. The numerator tangent and denominator vertical
line retain their point divisors also where individual displayed
rational functions have a zero or a pole. Direct cancellation in
\(\mu_3\) gives an analytic function at \(t=0\), with
\[
 \mu_3(0)=1/6,\qquad (x_6(0),y_6(0))=(1/36,-1/216),
 \qquad -x_6(0)/y_6(0)=6.
 \tag{R0E.24}
\]

Here is its entire quartic pullback. Put
\[
 N=zY-z^2-X/2,\quad u_3=8+64z^4,\quad
 v_3=3+96z^4+512z^8,
\]
\[
 \begin{split}
 \phi_2&=N/(z^2X-1/4),\\
 \phi_3&=\frac{XN}{zY+z^2-X/2-8z^4X},\\
 \phi_{23}&=
 \frac{zY-z^2v_3-\mu_3(z^4)(X-z^2u_3)}{z^2X-x_6(z^4)},
 \qquad \phi_6=\phi_3^2\phi_{23}.
 \end{split}                                                     \tag{R0E.25}
\]
Substitution in the original functions gives the exact identities
\[
 f_2=z^2\phi_2,\quad f_3=z^2\phi_3,\quad
 f_{2,3G}=z^2\phi_{23},\quad f_6=z^6\phi_6,
 \qquad
 \frac{B_{6,1}}{t^3}=\phi_2(p)^6\frac{\phi_6(Tp)}{\phi_6(p)}.
 \tag{R0E.26}
\]
In particular \(\phi_3|_0=X\), \(\phi_{23}|_0=6X\),
\(\phi_6|_0=6X^3\). The full source translation in R0E.6 therefore
gives central coefficient \(64\) in the actual Hom line. Each
\(\phi_2,\phi_3,\phi_{23},\phi_6\) changes sign under the quartic
deck. The coefficient in R0E.26 is invariant, including its source
argument. The powers \(z^2,z^6\) remain in the point-to-tensor
arrows throughout.

For the original two components the full functions give the following
limits. The first column uses R0E.9; the second uses R0E.10, with
\(\alpha\) the analytic solution of
\(\alpha+t\alpha^3=\beta^2-1\):
\[
\begin{array}{c|c|c}
 &C_0& C_1\\ \hline
 f_2&2/[p(p-2)]&t\,2(\beta-1)^2+O(t^2)\\
 f_3&p^{-2}&t(\beta-1)^3/(\beta-3)+O(t^2)\\
 f_{2,3G}&6/[p(p-6)]&6t(\beta-3)^2+O(t^2)\\
 f_6&6/[p^5(p-6)]&6t^3(\beta-1)^6+O(t^4).
\end{array}                                                       \tag{R0E.27}
\]
These are obtained from R0E.23 before multiplying the factors.
Combining them with the two directed translations R0E.11 gives
\[
 \left.\frac{B_{6,1}}{t^3}\right|_{C_0}
       =64\frac{p-6}{p},\qquad
 \left.\frac{B_{6,1}}{t^3}\right|_{C_1}
       =64\frac{\beta-1}{\beta-7}.
 \tag{R0E.28}
\]
For either even value \(a=2,6\), keep the actual point frames
\[
 r_{a,0}=(p-a)/p,\quad r_{a,1}=1,\quad
 e_{a,i}=\sigma_a/r_{a,i},\qquad
 \mathcal B_{a,1}=t^{-a/2}B_{a,1}\,r_a(Tp)/r_a(p).
 \tag{R0E.29}
\]
Equations R0E.14,28 give \(\mathcal B_{a,1}=2^a\) on both central
components. The chronological \(b=6/a\) product has central
coefficient \(64\). Its actual source map is \(p\mapsto p-bG\).
All four horizontal divisors \(O,aG,G,(a+1)G\), and all their
successive translates in the third power, remain in R0E.29.

The tangency point is also covered by an explicit source morphism.
On \(X\ne0\), put
\[
 T_s=t/x=z^2/X,\qquad Y_s=y/x=zY/X,\qquad
 x=Y_s^2-T_s-T_s^2=z^2X.
 \tag{R0E.30}
\]
The original resolution map is
\((x,Y_s,T_s)\mapsto(x,xY_s,xT_s)\), with the displayed smooth
source equation. Its two central components are \(T_s=0\) and
\(x=0\); their intersection has length two, since on \(T_s=0\)
the equation is \(x=Y_s^2\). R0E.30 includes the root points
\(X=\pm i,Y=z=0\). Its coordinates are invariant under the full
quartic deck. Together with R0E.17,20,21 these charts cover the
entire compact central root. At the tangency the coefficients in
R0E.14,28 tend respectively to \(4,64\), as \(p\) or \(\beta\)
tends to infinity. In the normalization used in EF.30b,
\(w=1/2-1/p\); its literal coefficient
\(128(3w-1)\) is exactly R0E.28, including this limit.

The root translation carries the graph itself. At \(O\), substitution
\(X=U/\eta^2,Y=-U/\eta^3\) in R0E.6 gives
\[
 X'=\eta(\eta-2z)/U+2z^2\eta^4/U^2,\qquad
 Y'=z+(1/\eta-z\eta^2/U)X'.
 \tag{R0E.31}
\]
The last expression is holomorphic and has first jet \(\eta-z\),
so \((Y',z)=(\eta,z)\) as analytic ideals. At \(T\) write
\[
 A=1-\frac{X(Y-z)^2}{(1+X^2)^2},\qquad
 \eta'=-X'/Y'
  =\frac{(Y-z)A(1+X^2)^2}{A(1+X^2)^3-z(Y-z)^3}.
 \tag{R0E.32}
\]
Its denominator is a unit and \((\eta',z)=(Y,z)\). The inverse
\(+G'\) has the inverse ideal squares. These equations lift both
source maps to the two projective graph charts. On the remaining
\(X\)-unit chart they agree with the actual rational translation
and extend by their inverse morphisms on the smooth minimal source.

For \(a=2,6\), the full marked-section equation in the \(O\)-graph is
\[
 \tau(aG')=z\tau_{a,\rm old}(t),\quad
 \frac{\eta-\tau(aG')}{\eta}
 =\frac{p-\tau_{a,\rm old}(z^4)}p
 =1-v\tau_{a,\rm old}(z^4),
 \tag{R0E.33}
\]
where the last expression is on the other projective chart,
\(\tau_{2,\rm old}=2/(1+8t)\) and
\(\tau_{6,\rm old}=-x_6/y_6\). Both are analytic, with limits
\(2,6\), and both marked sections lie in this \(O\)-chart.
The numerator and denominator in R0E.33 have equal exceptional
orders; their two strict directions remain distinct. The root point
line and the pullback of the original point line are thus identified
by these actual ratios on the common graph. On the \(T\)-graph and
R0E.30 neither marked section is present, so the coefficient is one.
All these maps are quartic-deck invariant.

The original III component matrix is
\(\left(\begin{smallmatrix}-2&2\\2&-2\end{smallmatrix}\right)\)
and its primitive kernel is \((1,1)\). The even marked point \(aG\)
and \(O\) lie on the identity component, so the line has degree zero
on each component. The Hom multidegrees are zero. R0E.14,28 give
order \(a/2\) on each component. After removing precisely that
whole-fibre factor, the rational Hom section has no vertical divisor;
its horizontal divisor is exactly canceled by the two modules in
R0E.29,33. On the smooth resolved source, a rational section with
zero divisor and its reciprocal extend in every local line frame.
This proves the full map through the tangent point and every marked
section, with no omitted codimension-two extension. In particular
\[
 U_G(p,v_M)=(p+G,B_{a,1}(p+G)^{-1}v_M),\qquad
                 t^{a/2}U_G
 \tag{R0E.34}
\]
is an actual invertible finite-source line map. The inverse always
uses the reciprocal at its actual inverse source point.

\par

\medskip\noindent\textbf{R0E.5. The entire period module and the intrinsic original base}

\par\nobreak

The completed objects in R0E.22 are the central fibres of the actual
K96 disc family. F46 gives their integral sphere topology on its
reference fan; the finite common blowups, both ruling maps and the
primitive integral Gysin and van Kampen maps of F35 carry that topology
to each prescribed lattice-basis fan. The root and annular maps used
in this passage are stated again in R0E.8 below, specialized to
\(r=0\). Thus the constructed compact threefold has
\(H^2(X;\mathbb Z)=H^4(X;\mathbb Z)=0\). Its scalar belongs to the
actual common contraction domain, and its full character and reference
remain fixed. Its ordinary fibre is the compact complex two-torus
obtained from the given point-line suspension.

A horizontal effective divisor would restrict to a nonzero effective
divisor on a generic two-torus. Its class is the restriction of a
class in \(H^2(X)=0\). For a flat Kähler form \(\omega\) on that
torus the divisor formula would give
\(0=\int c_1(D)\wedge\omega=\int_D\omega>0\), with all positive
multiplicities included. Consequently every polar divisor of a
meromorphic function is vertical. On a regular compact torus the
function is holomorphic after removal of codimension-two indeterminacy
and hence constant. It descends to the regular base. At a generic
smooth point of a reduced special component choose the exact
transverse coordinate \(t-t_0=u^e\), after taking the extending root
of its base unit. The finite Laurent principal part is invariant
under \(u\mapsto\zeta_eu\), since the function is constant on its
generic fibres. All exponents are divisible by \(e\), proving the
meromorphic extension on the base. This proves the actual pullback
\(\mathbb C(t)\xrightarrow{\pi^*}\mathbb C(X)\) is an isomorphism.
An endpoint isomorphism therefore has
\(\pi_{\rm target}F=h\pi_{\rm source}\) for a Möbius map \(h\),
with the same equality across all special fibres.

We identify the full source quotient before restricting this base map.
On a regular source disc choose its actual Abel representative
\(g\), period \(\tau\), and write
\[
 Z=e^{2\pi iz},\quad Q=e^{2\pi i\tau},\quad P=e^{2\pi ig},
 \quad f_P(Z;Q)=\frac{\Theta(Z/P;Q)}{\Theta(Z;Q)},\quad
 V_{\rm pt}=f_{P^a}(Z;Q)v_M.
 \tag{R0E.35}
\]
Here \(v_M\) is the coefficient in the canonical meromorphic point
section; its actual two Cartier modules, and the reciprocal arrow in
each, define R0E.35 through the section points. The theta product
satisfies \(f_P(QZ)=P f_P(Z)\). In particular the four periods of
the actual completed torus, in the order
\((\gamma,e_1,e_2,\delta)\), are
\[
 \gamma=(-bg,\ell),\quad e_1=(1,0),\quad
 e_2=(\tau,ag),\quad\delta=(0,1).
 \tag{R0E.36}
\]
The entire suspension multiplier defining \(\ell\) is
\[
 \lambda\frac q{t^3}B_{a,b}(p)
       \frac{f_{P^a}(P^{-b}Z;Q)}{f_{P^a}(Z;Q)}.
 \tag{R0E.37}
\]
Its ratio with the translation lift having the same constant factors
is an invertible function on the compact elliptic fibre, so it is a
base constant. R0E.37 specifies that whole constant, including every
Miller and selected-height factor. The tensor-frame dictionary is
\[
 V_{\rm tensor}=f_P(Z)^a f_a(p)^{-1}v_M
        =C_a(t)V_{\rm pt},\qquad
 C_a(t)=\frac{f_P(Z)^a}{f_a(p)f_{P^a}(Z)}.
 \tag{R0E.38}
\]
The full divisor identity proves independence of the source point;
its value and its logarithm continuation remain in \(\ell\).

For completeness the exact period-coordinate changes are recorded.
For \(g'=g+\mu+\nu\tau\), \(\mu,\nu\in\mathbb Z\), the entire
point coefficient changes by
\[
 V'_{\rm pt}=(-Z/P^a)^{a\nu}
              Q^{-a\nu(a\nu+1)/2}V_{\rm pt}.
 \tag{R0E.39}
\]
Its inverse pulls the new vectors to
\[
 e'_1=e_1-a\nu\delta,\quad e'_2=e_2+a\mu\delta,\quad
 \gamma'=\gamma-b\mu e_1-b\nu e_2+c\delta,\quad
 c=\ell'-\ell+ab\nu(2g+\mu+\nu\tau)\in\mathbb Z.
 \tag{R0E.40}
\]
For \(A=\left(\begin{smallmatrix}p&q_0\\r_0&s_0\end{smallmatrix}\right)
\in SL_2(\mathbb Z)\), set \(L=p+r_0\tau\). The full source/line map is
\[
 z'=z/L,\quad \tau'=(q_0+s_0\tau)/L,\quad g'=g/L,
 \quad w'=w-a r_0gz/L+c_A,
 \tag{R0E.41}
\]
where its entire constant is fixed by
\[
 e^{2\pi i c_A}=
 \frac{f_{e^{2\pi iag'}}(e^{2\pi iz'};e^{2\pi i\tau'})}
 {f_{P^a}(Z;Q)}\exp(2\pi i a r_0gz/L).
 \tag{R0E.42}
\]
The divisor and both period factors cancel, making this the stated
constant; evaluating it at an ordinary point fixes its full value.
The inverse solves R0E.41 with its actual target parameters. The
vectors \(p e_1+r_0e_2,q_0e_1+s_0e_2\) go exactly to the two new
elliptic periods. The resulting entire lattice monodromy is
\[
 M_\alpha=\begin{pmatrix}
 1&0&0\\ b v_\alpha&A_\alpha&0\\
 c_\alpha&a v_\alpha^t J A_\alpha&1
 \end{pmatrix},\quad J=\begin{pmatrix}0&-1\\1&0\end{pmatrix},
 \tag{R0E.43}
\]
\[
 v_{\alpha\beta}=v_\alpha+A_\alpha v_\beta,\qquad
 c_{\alpha\beta}=c_\alpha+c_\beta+
                  ab v_\alpha^tJA_\alpha v_\beta.
 \tag{R0E.44}
\]
The integral central cochain is the actual one from the full functions;
none of its entries is prescribed to vanish.

The rational Abel class in R0E.43 is nonzero. Indeed a rational
coboundary, multiplied by an integer and its full period cochain,
would give a single-valued holomorphic Lie logarithm of \(NG\).
Multiply \(N\) further by six to put it in every identity component.
At the III fibre R0E.4 puts \(2G'\) in the formal neighbourhood of
\(O\), with \(\tau(2G')=2z/(1+8z^4)\). The root formal logarithm
has deck weight \(i\); its quotient by \(z\), by R0E.3, is an
extending logarithm in the original finite Lie frame. The difference
from the proposed single-valued logarithm is an invariant elliptic
period. The central quartic action on \(dX/(2Y)\) has eigenvalue
\(i\), so it has no nonzero invariant period. At the selected I1,
the invariant annular period extends, by integration of the actual
invariant differential in \(xy=q\).

At original infinity the exact good cubic is
\(Y^2=X^3+zX+1\). The section \(3G'\) has root logarithm
\(-z^2/8+O(z^5)\). The original minimal differential is
\(\omega_0=z^{-2}\omega_{\rm good}\); thus division by \(z^2\)
gives its extending original logarithm. There is no invariant
elliptic period under its order-three central action. Consequently
the proposed logarithm is a global section of the actual Lie line.
Since \(x_0=q_\infty^2x,y_0=q_\infty^3y\),
\(\omega_0=q_\infty^{-1}\omega_t\). Its coefficient transition
is \(u_0(q)=q^{-1}u_t(q^{-1})\). An entire \(u_t\) satisfying this
is \(O(t^{-1})\) at infinity, hence zero. The Lie line is
\(\mathcal O(-1)\) by this explicit transition. On the other hand
\(3G\) has original central additive coordinate
\(-x_0/y_0=-1/8\); all its nonzero integral multiples are nonzero.
Thus \(G\) is not torsion and the rational Abel class cannot vanish.

The invariant integral subgroups of R0E.43 are now determined:
\[
 \Lambda^{\pi_1}=\mathbb Z\delta,\qquad
 (\Lambda^*)^{\pi_1}=\mathbb Z\psi,\quad
 \psi(\gamma)=1,\quad\psi(e_1)=\psi(e_2)=\psi(\delta)=0.
 \tag{R0E.45}
\]
An invariant vector with \(\gamma\)-coefficient \(z\) would solve
\((A_\alpha-I)u+bzv_\alpha=0\). The nonzero Abel class forces
\(z=0\), and infinity has no elliptic invariant vector, so \(u=0\).
The dual equation gives the second assertion. Primitivity follows
from the literal nonzero-line cover with lattice
\(\langle e_1,e_2,\delta\rangle\), whose primitive extra deck is
\(\gamma\).

An endpoint intertwines these invariant groups even before the base
map is fixed. On each regular fibre its lifted derivative is periodic
for the complete lattice, hence constant. Complex linearity and
\(F_*\delta=\epsilon\delta\) give the actual affine map and inverse
\[
 (z,w)\mapsto(Az+z_0,Cz+\epsilon w+w_0),\qquad
 (z',w')\mapsto
 ((z'-z_0)/A,\epsilon(w'-w_0-C(z'-z_0)/A)).
 \tag{R0E.46}
\]
Here \(A\ne0\) and \(\epsilon=\pm1\). The preserved subgroup
\(\ker\psi\) gives a degree-one elliptic source map, with complete
integer matrix
\[
 \begin{pmatrix}p&q_0&0\\ r_0&s_0&0\\u&v&\epsilon\end{pmatrix},
 \quad A=p+r_0\tau(h(t)),\quad
 A\tau(t)=q_0+s_0\tau(h(t)),\quad ps_0-q_0r_0=1.
 \tag{R0E.47}
\]
The determinant sign is positive because the first coordinate is a
nonzero complex-linear map. Thus the two period lattices are mapped
bijectively, with the displayed inverse. The original elliptic
quotient is identified through the actual point-line cyclic cover,
including all integers \(u,v,c_\alpha\).

Directly for R0E.1 the elliptic invariant is
\[
                         j_0(t)=\frac{6912}{27t+4}.
 \tag{R0E.48}
\]
R0E.46-\/-47 imply \(j_0(h(t))=j_0(t)\); substitution in this rational
function gives \(h(t)=t\). This treats every permitted base map,
including a proposed interchange of marked branches, using the actual
elliptic isomorphism just obtained.

\par

\medskip\noindent\textbf{R0E.6. Full linear centralizer and the primitive source section}

\par\nobreak

The rational elliptic monodromy centralizer is \(\mathbb QI\).
Its infinity matrix generates the quadratic field with polynomial
\(x^2+x+1\). The selected I1 has a nonzero unipotent nilpotent part,
so it cannot belong to that field. A non-scalar element of that field
generates it, proving the assertion for both monodromies together.
A matrix commuting with the full R0E.43 preserves R0E.45 and has form
\[
 X=\begin{pmatrix}s_0&0&0\\u&A&0\\c&w&t_0\end{pmatrix}.
\]
The middle block is \(A=eI\). Its middle-left equation is
\((I-A_\alpha)u=b(s_0-e)v_\alpha\), so the nonzero Abel class gives
\(s_0=e,u=0\). The lower-middle equation is the dual Abel coboundary
with coefficient \(a(e-t_0)\), giving \(t_0=e,w=0\). All remaining
central-cochain equations are satisfied exactly by
\[
                             X=eI+c\delta\psi.
 \tag{R0E.49}
\]
Integral unimodularity gives \(e=\pm1,c\in\mathbb Z\). Its complex
linear realization excludes \(c\ne0\): subtracting \(eI\) would
have real rank one, whereas a nonzero complex-linear map has even
real rank. Thus an endpoint has linear part precisely \(I\) or
\(-I\), with its complete affine displacement still retained.

For an automorphism of a single completed object the negative part
is excluded geometrically. It has sixteen transverse fixed points
on each regular two-torus, since its vertical fixed-point equation
is multiplication by two. The actual global equalizer with identity
is a closed analytic set in the compact threefold. An irreducible
component through a regular fixed point has dimension one and is the
proper analytic closure of a horizontal fixed curve. Its normalization
maps to the original base with positive degree. Pairing it with
\(\pi^*[\omega_{\mathbb P^1}]=0\) in \(H^2(X)\) gives zero, while
integration of this positive base form gives positive degree. This
contradiction uses the full cap-completed map. Negative isomorphisms
between two scalar cut endpoints will be treated by their actual
annular continuation in R0E.10.

An identity-linear endpoint lifts on the regular locus to the cyclic
nonzero-line cover, commuting with its actual generator \(B_\lambda\).
At the ordinary III disc, R0E.26-\/-34 give that cover as the invertible
contraction on the actual resolved point line. A nonzero line vector
over the \(O\)-section on the base disc kills the base meridian.
Generic-fibre loops therefore generate after the actual vanishing-cycle
relations, all in \(\ker\psi\). The preserved fibre homomorphism
extends and the covering-space lift is single valued through that
whole disc. Its source projection evaluated at this \(O\)-vector
is an extending section. The same argument applies to any ordinary
smooth or nodal source disc, retaining both Cartier modules.

At the selected node write its entire periods
\[
 h:(Z,V,q)\mapsto(QZ,\alpha V,q),\qquad
 B:(Z,V,q)\mapsto(gZ,q\eta V,q),\quad
 Q=\kappa q,\quad\alpha=P^a,\quad g=P^{-b}.
 \tag{R0E.50}
\]
All four displayed coefficient functions are the actual nonzero
units of R0E.35-\/-38. The source normalization of the selected fibre
is \(w=y/(x-2/9)\), with branches \(\pm\sqrt{2/3}\). Taking
\(Z=(w-\sqrt{2/3})/(w+\sqrt{2/3})\) gives
\[
             P(0)=\frac{2-\sqrt6}{2+\sqrt6}=-5+2\sqrt6.
 \tag{R0E.51}
\]
The full functions \(P(q),Q(q)\) remain in R0E.50. In a cone chart
\(q=z_0z_1z_2\), the image of a section has nonnegative coordinate
orders summing to one. Its two torus characters thus have integral
orders. The actual height vectors of \(h\) and \(B\) are the two
primitive coordinate generators, so a literal \(h^rB^k\) reindex
removes them. Keeping all its unit factors, the endpoint and inverse
then give extending units \(K(q),C(q)\). On every character their
map and inverse are
\[
 Z^uV^vq^d\mapsto K(q)^uC(q)^v Z^uV^vq^d,
 \qquad Z^uV^vq^d\mapsto K(q)^{-u}C(q)^{-v}Z^uV^vq^d.
 \tag{R0E.52}
\]
They are actual maps on every prescribed cone, with \(q\) unchanged.
The integer \(B^k\) is retained for the divisor equation below.

At infinity use the literal normalized fibre product with
\(z^m=q_\infty\). The maps in R0E.8 identify it with the root
quotient by \(B_\lambda\), since \(A_\lambda^m=B_\lambda^n\)
and \(n=\pm1\). Its canonical endpoint lift fixes \(z\).
The nonzero-line lift has commutator \(B_\lambda^k\) with the full
cap generator. Iterating \(m\) times gives \(B_\lambda^{mk}=1\)
because the lift already commutes with \(B_\lambda^n\). Strict
contraction gives infinite order, hence \(k=0\). This retains the
finite character. Evaluation extends the source section through
the actual root and its Cartier graph. The local lifts away from
the selected point differ by integer suspension powers. Their
orbifold base is \(\mathbb P^1\setminus\{t_c\}\) with its one
order-\(m\) cap. Its fundamental group is finite cyclic, so its
homomorphism to the integer deck group is zero. Removing this
cocycle by its actual deck powers gives a global source section.

We identify its section group integrally on this source. The rational
chart \(UV=x^2(x-1/4)\), with
\(t=(V-U-x)/2\), is the original surface. Resolving its A1 point
gives the rational elliptic surface. The actual fibres are IV*, III
and I1, with Euler sum \(8+3+1=12\); its intersection lattice is
unimodular of rank ten. In the A1 tangent cone
\(UV=-x^2/4\), the section \(G\) has direction
\([x:U:V]=[0:0:1]\), off the strict fibre plane \(V-U-x=0\).
It therefore meets the nonidentity III component once, away from
the tangency; \(O\) meets the identity component. The two sections
are disjoint. At infinity the actual cubic root has the distinct
fixed sections \((0,1)\) and \(O\), with quotient weights \((2,1)\)
and rays \((2/3,1/3),(1/3,2/3)\). Their two resolved arms are
disjoint. Label the E6 components as in the root charts, with edges
\((0,1),(1,2),(2,3),(3,4),(2,5)\) and \(G\) meeting component 0.
The positive Cartan matrix is
\[
 A=\begin{pmatrix}
 2&-1&0&0&0&0\\-1&2&-1&0&0&0\\
 0&-1&2&-1&0&-1\\0&0&-1&2&-1&0\\
 0&0&0&-1&2&0\\0&0&-1&0&0&2
 \end{pmatrix},\quad
 A^{-1}e_0=(4/3,5/3,2,4/3,2/3,1)^t.
 \tag{R0E.53}
\]
The nonidentity III block is \((2)\), with intersection vector 1.
The two corrections are \(4/3,1/2\); \(G\cdot O=0\), so
\(h(G)=2-4/3-1/2=1/6\). Both section self-intersections are \(-1\)
from the already identified Lie line. In the order
\((O,F,C_0,\ldots,C_5,C_{\rm III},G)\), take
\(O^2=G^2=-1,OF=GF=1\), block \(-A\oplus(-2)\),
\(GC_0=GC_{\rm III}=1\), with other unspecified pairings zero.
The determinant of this integral Gram matrix is \(-1\). The ten
actual divisor classes give the full lattice; restriction to the
generic fibre identifies its quotient by the first nine classes
with the free section group generated by \(G\). Equivalently the
height equations give \(\iota^2|\mathrm{MW}_{\rm tors}|^2=1\).
Thus the extending section above is exactly \(jG\), \(j\in\mathbb Z\).

\par

\medskip\noindent\textbf{R0E.7. The actual full automorphism factors through the physical scalar}

\par\nobreak

For \(a=2,6\), retain the entire point-line multiplier
\[
 B_{a,1}(p)=f_2(p)^a\frac{f_a(p-G)}{f_a(p)},\qquad
 B_{a,b}=U_G^{-b},\qquad B_{\rm sel}=(q/t^3)B_{a,b}.
 \tag{R0E.54}
\]
An ordinary endpoint lift covering \(+jG\) differs from the actual
\(U_G^j\) by a base function \(H(t)\), because their ratio on a
compact elliptic fibre is a holomorphic nonzero function. Its
extension is meromorphic in the actual finite and root Cartier
modules. R0E.52 gives its integral order at the selected node, so
\(H\) is a rational base function.

The exact coefficient orders, with the root order measured in
\(z^m=q_\infty\), are
\[
\begin{array}{c|ccc}
 &t=0&z=0&q=0\\ \hline
 U_G&-3/b&2m/b&0\\
 H&3j/b&-2mj/b&-k.
\end{array}                                                       \tag{R0E.55}
\]
The first column is R0E.26-\/-34, including the tangency and both
Cartier modules. The root column follows from the entire pullbacks
\[
 \begin{array}{c|c|c}
 (a,b;m)& B_{a,1,\rm old}&\text{central root Hom coefficient}\\ \hline
 (2,3;3l)&z^{-2l}B'_{2,1}&B'_{2,1}=-X\text{ in the canonical point frame}\\
 (6,1;3l+2)&z^{-6(l+1)}B'_6&B'_6/z^2\longrightarrow4 .
 \end{array}                                                       \tag{R0E.56}
\]
The root graph and strict-section arrows giving these expressions
are retained in R0E.8. Their generic orders do not replace those
extension conditions. The third column is the literal \(B^k\)
reindex in R0E.52. At every other finite irreducible fibre the Hom
numerator and denominator have equal generic order; none of the
remaining Miller factors contains that fibre. Thus there are no
other base divisors. Their degree equation is
\[
 3j/b-2j/b-k=0,\qquad j=bk,\qquad
 H=c(t^3/q)^k,\qquad H U_G^{bk}=c B_{\rm sel}^{-k}.
 \tag{R0E.57}
\]
The full quotient deck is \(B_\lambda=\lambda B_{\rm sel}\), so
this last map is exactly multiplication by \(c\lambda^k\) on
the actual point line. It extends through the original source,
root, fan and cap by their complete coefficient maps and commutes
with their annular square. Conversely this physical action gives
those actual holomorphic automorphisms. Its kernel is trivial:
an identity physical scalar would equal one fixed integral deck
power on the connected nonzero-line cover; that power translates
the source by \(-kbG\). The proven non-torsion of \(G\) forces
\(k=0\), and its scalar is one. Therefore on every object R0E.22,
\[
                  \operatorname{Aut}_{\rm hol}(X)
                        =\mathbb C^*_{\rm physical}.
 \tag{R0E.58}
\]
The group on the right is connected by its actual scalar paths.
All selected fan markings and finite character data in this equality
are the ones specified at the outset.

\par

\medskip\noindent\textbf{R0E.8. Both full infinity roots and the retained annular reference}

\par\nobreak

We give the explicit specialization of K96 rather than extending the
finite quartic root to infinity. In the reduced row put
\(m=3l+2,k=2l+1\). The full original-coordinate map is
\[
 \begin{gathered}
 q_\infty=z^m,\quad t=z^{-m},\quad
 x=z^{-2(l+1)}X,\quad y=z^{-3(l+1)}Y,\\
 Y^2=X^3+z^{l+2}X+z^2,\quad G'=(0,z).
 \end{gathered}\tag{R0E.59}
\]
Its deck is \(D(z,X,Y)=(\zeta_mz,\zeta_m^{-2k}X,
\zeta_m^{-3k}Y)\). With \(c=(1+z^lX)^{1/2}\), chosen with
constant term one, the exact coordinates
\(U=Y-zc,V=Y+zc\) have \(UV=X^3\). The equation
\(2zc=V-U\) has its analytic inverse for \(z\). The three regular
A2 charts are
\[
 X=\xi\eta,\quad U=\xi^{3-i}\eta^{2-i},\quad
 V=\xi^i\eta^{i+1},\qquad i=0,1,2.
 \tag{R0E.60}
\]
In the middle chart \(X=ab,U=a^2b,V=ab^2\), the entire \(+G'\)
map is
\[
 \begin{gathered}
 K=1+\frac{b(b-a)z^l}{2c(c+1)},\quad
 F=1-\frac{abz^l(K+1)}{2c(c+1)},\quad X'=a(a-b)F,\\
 c'=(1+z^lX')^{1/2},\quad
 H=K+\frac{b(a-b)z^l}{2c(1+c')},\\
 a'=(b-a)F/H,\quad b'=-aH.
 \end{gathered}                                                   \tag{R0E.61}
\]
It is the chord with slope \((Y-z)/X=aK\),
\(Y'=-z-aKX'\); hence \(V'=-aX'H\) and the inverse of
\(2zc=V-U\) retains exactly the same \(z\). The inverse chord
is obtained from \(Y\mapsto-Y\). At the other smooth group points
it is the ordinary translation, giving the inverse morphisms in
all three charts. The vertical field
\[
 W=a(2b-a-z^{l+1}/c)\partial_a+
   b(2a-b+z^{l+1}/c)\partial_b
 \tag{R0E.62}
\]
is holomorphic and vanishes at the triple point. Its flow is the
identity-component group flow on the dense smooth locus. Thus
\[
 R_{\rm div}=2G'-[m]^{-1}((2m-1)G'),\qquad mR_{\rm div}=G'
 \tag{R0E.63}
\]
is an actual extending translation: \(2m-1=3(2l+1)\) puts the
argument in the small identity formal group. The formal inverse
has derivative \(1/m\). The deck commutes with these arrows:
\(l-2k=-m\) makes \(c\) invariant; the two middle coordinates
have weight \(\zeta_m^{l+1}\), and every \(K,F,H\) is invariant.
Uniqueness of the formal inverse gives the same commutation for
R0E.63.

The point-line graph retains
\[
 \tau'=-X/Y=z^k\tau_0,\qquad
 \tau'(6G')=z^kL(q_\infty),\quad L(0)=-1/4.
 \tag{R0E.64}
\]
It is the normalized graph of \((\tau',z^k)\), resolved by the
ordered rays \((j,1)\), \(1\le j\le k\), together with its axes.
Adjacent determinants are one; both section equations have order
\(j\) on the ray \((j,1)\). On the equality ray \(j=k\) their
strict directions are separately \(0,L\). Their exact ratio is
\[
 \frac{\tau'}{\tau'-z^kL}=\frac u{u-L},\qquad
 \frac1{1-vL}\quad\text{on }\tau'v=z^k,\ uv=1.
 \tag{R0E.65}
\]
The remaining cyclic charts carry these same monomials and units.
The product of the actual commuting source/root translations
preserves the closure of the graph and its inverse. Consequently
\(\beta^*N_6=\pi^*M_6\) as actual Cartier maps. The full selected
coefficient is
\[
 B'_{\rm sel}=h z^{-2}B'_6,\quad
 h=1-t_cz^m=1+4z^m/27,\quad B'_6/z^2|_0=4.
 \tag{R0E.66}
\]
The last limit is the full K96 EX.13 chord calculation at
\(s=1/4\); all its factors are retained by R0E.23 and the graph.
Its zero vertical divisor after the displayed factor and its
literal horizontal Hom divisor give the extending map and inverse.

For the good row \(m=3l\), the original map is
\[
 q_\infty=z^{3l},\quad t=z^{-3l},\quad
 x=z^{-2l}X,\quad y=z^{-3l}Y,\qquad
 Y^2=X^3+z^lX+1,\quad G'=(0,1).
 \tag{R0E.67}
\]
Its deck is \((z,X,Y)\mapsto(\zeta_mz,\zeta_m^{-4l}X,Y)\).
Put \(u=z^l\). At \(2G'=(u^2/4,-1-u^3/8)\) use
\(w=X-u^2/4\); at \(O\) use \(w=-X/Y\). After the literal
\(\mu_l\) quotient these have weights \((2,1)\). The invariant
coordinates \((w^3,u^3,wu)\) have equation \(AB=C^3\), resolved
by the three charts of R0E.60 with these corresponding coordinates.
In the quotient lattice the two rays are
\((2/3,1/3),(1/3,2/3)\). The strict-section support is
\(\max(v_w-2v_u,0)\); its difference from the pulled-back root
section is \(\min(v_w,2v_u)\), hence
\[
 \beta^*N_2=\pi^*M_2+E_{2G'}-E_O,\qquad
 E_P(v_w,v_z)=\min(v_w,2l v_z).
 \tag{R0E.68}
\]
This equality uses the ratio of the actual monomial section
generators and their units on each cone, with inverse ratio in the
other direction. The small section
\(R_{\rm div}=[m]^{-1}(3G')\) has
\(\tau=-u^2/(8m)+O(u^5)\). Its translation has
\(w\mapsto w\varepsilon(w,u)+u^2a(w,u)\), with \(\varepsilon\)
a unit, and preserves \((w,u^j),(u,w^j)\), \(j=1,2\), with their
inverse graph arrows. Thus it carries both full corrections.
The original Miller factor is
\[
 f_{2,\rm old}=z^{-l}\phi,\quad
 \phi=\frac{Y-1-z^lX/2}{X-z^{2l}/4},\quad
 B'_{\rm sel}=hB'_{2,3},\quad B'_{2,3}|_0=4.
 \tag{R0E.69}
\]
Indeed on \(Y^2=X^3+1\) the three translated factors are
\((Y-1)/X,-(Y+1)/X,2/X\), whose product is \(-2\); their squared
norm is four. Their actual chord translations have order three
there. This verifies the entire central norm while preserving the
full functions away from the central curve.

Here are the exact norm lift and reference in both rows. For
\(N=\mathcal O(P_0-O)\), with \(P_0=aG'\), use the literal
chord-over-vertical coefficient
\[
 g_{P_0,R}(p)=\frac{\ell_{P_0,R}(p)}{x(p)-x(P_0+R)}.
 \tag{R0E.70}
\]
Its divisor is \((P_0)+(R)-(P_0+R)-(O)\), precisely the Hom
coefficient for the map covering \(p\mapsto p-R\). Set
\(F=\tau(R)g_{P_0,R}\) in the good row and
\(F=\tau(P_0)g_{P_0,R}\) in the reduced row. Their orders are
\(2l,k\), with leading coefficients \(-1/(8m),-1/4\).
Near \(O\), \(x=\tau^{-2}+O(1),y=-\tau^{-3}+O(\tau^{-1})\)
and the actual chord slope is \(-\tau^{-1}+O(1)\). Thus multiplying
that chord by the displayed \(\tau\) gives central coefficient
one. The horizontal Hom divisor and the primitive component kernel
extend \(F\) and its inverse through the entire root. Its chord
weight is \(\zeta_m^{-k}\), its \(\tau\)-factor weight
\(\zeta_m^k\), where \(k=2l\) in the good row. It therefore
commutes with the full deck on every Cartier chart.

The ratio \(\beta=B'_{\rm sel}/F^m\) covers identity and is an
invertible base function, central value four. It is deck invariant.
Its analytic \(m\)-th root is fixed by the full chosen reference;
its transform-to-itself ratio is a constant root of unity equal to
one at the central fibre. Define
\[
 \widetilde F=\beta^{1/m}F,\quad
 A_0=\zeta_m^d\widetilde F^{\,n}D,\quad
 R_\lambda^m=\lambda,\quad A_\lambda=R_\lambda^n A_0,\quad
 B_\lambda=\lambda B'_{\rm sel}.
 \tag{R0E.71}
\]
Then \(\widetilde F^m=B'_{\rm sel}\),
\(A_\lambda^m=B_\lambda^n\), and
\([A_\lambda,B_\lambda]=1\), as full source/line maps with
character \(\zeta_m^d\).

On the prescribed sectors set
\(s_\alpha=\log_\alpha q_\infty/(2\pi i)\),
\(z_\alpha=e^{2\pi i s_\alpha/m}\),
\(c_\lambda=n\log R_\lambda/(2\pi i)+d/m\), and
\(f_\alpha=\operatorname{Log}_\alpha\widetilde F\). The full map is
\[
 \Phi_\alpha=\operatorname{Exp}
 (n s_\alpha f_\alpha+s_\alpha c_\lambda\delta+h_\alpha)I_{z_\alpha}.
 \tag{R0E.72}
\]
It retains signed source displacement \(-n s_\alpha R_{\rm div}\),
scalar \(R_\lambda^{n s_\alpha}e^{2\pi i d s_\alpha/m}\), and the
literal root change. For \(s_\beta=s_\alpha+k\),
\(f_\beta=D_*^k f_\alpha+\Pi_{\alpha\beta}\), its complete
cochain equation is
\[
 h_\beta-D_*^kh_\alpha=-n s_\alpha\Pi_{\alpha\beta}
               +j_{\alpha\beta}\operatorname{Log}_\beta B_\lambda
               +\Lambda_{\alpha\beta}.
 \tag{R0E.73}
\]
All transported periods and all integers are those of the prescribed
K96 reference. In the fixed ordinary Lie frame take its extension
\(h^0\), residual \(E\), and fixed partition \(\chi_\beta\).
The explicit solve is
\(k_\alpha=-\sum\chi_\beta E_{\alpha\beta}\),
\(v(q)=\pi^{-1}\int\bar\partial k(w)/(q-w)\,dA(w)\),
\(h_\alpha=h_\alpha^0+k_\alpha-v\). The residual has the compact
support specified in K96 and vanishes at its reference; this solve
therefore preserves the full reference value and every integral
column. Subtracting the target-sheet logarithm
\(kn f_\beta+kc_\lambda\delta\) gives exactly
\[
 \Phi_\beta=A_\lambda^k B_\lambda^{j_{\alpha\beta}}\Phi_\alpha,
 \qquad \Phi_\alpha B_\lambda=B_\lambda\Phi_\alpha.
 \tag{R0E.74}
\]
The inverse is the inverse root map followed by the negative
translation at its actual inverse source point, with all reciprocal
line factors. These are the entire ordinary/cap gluing arrows used
above. Their finite roots extend by R0E.59-\/-71 on their own infinity
disc. The finite III quartic charts have original \(t\) near zero;
their domain comparison is precisely R0E.1 and R0E.59 or 67 on the
common ordinary annulus. Every base function and reference period
is pulled back by these stated original-coordinate maps.

\par

\medskip\noindent\textbf{R0E.9. The complete scalar family and its primitive positive endpoints}

\par\nobreak

Keep the full scalar root \(R^m=\lambda\) distinct from the source
coordinate \(t\). Starting at the prescribed reference, set
\[
 \lambda_\theta=e^{2\pi i\theta}\lambda_*,
 \quad R_\theta=e^{2\pi i\theta/m}R_*,
 \quad B_\theta=e^{2\pi i\theta}B_*,
 \quad A_\theta=e^{2\pi i n\theta/m}A_*.
 \tag{R0E.75}
\]
On its labelled reference sectors let \(J_\alpha\) be the actual
integer representative with \(J_\beta-J_\alpha=j_{\alpha\beta}\).
The entire scalar family is
\[
 \Phi_{\theta,\alpha}
 =e^{2\pi i n\theta s_\alpha/m}
    B_\theta^{J_\alpha}
      (B_*^{-J_\alpha}\Phi_{*,\alpha}).
 \tag{R0E.76}
\]
This formula retains the source-dependent periods \(\Pi,\Lambda\)
and their reference integers. On an overlap its two extra factors
are \(e^{2\pi i n\theta k/m}\) and
\(e^{2\pi i\theta j_{\alpha\beta}}\), giving exactly the
R0E.74 squares for \(A_\theta,B_\theta\). All reciprocal factors
give its inverse. At an integral winding \(J\),
\[
 B_{\theta+J}=B_\theta,\qquad
 \Phi_{\theta+J,\alpha}=z^{nJ}\Phi_{\theta,\alpha}.
 \tag{R0E.77}
\]
The retained suspension factors have ratio
\(e^{2\pi iJ J_\alpha}=1\), after their integers have been carried.

Every scalar endpoint fixes the original base by R0E.45-\/-48.
It intertwines the identical full monodromy matrices R0E.43,
since the scalar changes only the entire suspension logarithm.
Its lattice matrix has form R0E.49. The two complex-independent
vectors \(e_1=(1,0)\), \(\delta=(0,1)\) force its complex-linear
map to be \(eI\). Its suspension equation therefore says
\[
             e(\ell_{\rm target}-\ell_{\rm source})\in\mathbb Z.
 \tag{R0E.78}
\]
All source, height and Cartier factors in R0E.37 agree. Their
remaining difference is the logarithm of the original scalar
ratio, so the two scalar values are equal. Its scalar path has
integral winding \(J\), and its ordinary linear part is one of
the two signs. The integer in R0E.78 is retained as its suspension
lift change.

A positive endpoint has ordinary lift \(H(t)U_G^j\) by the complete
source/cover argument. R0E.77 gives its full extra root factor
\(z^{nJ}\). Consequently the three divisor equations become
\[
 \operatorname{ord}_0H=3j/b,\qquad
 \operatorname{ord}_{t_c}H=-k,\qquad
 m\operatorname{ord}_\infty H+2mj/b+nJ=0.
 \tag{R0E.79}
\]
There are no other base divisors, by R0E.26-\/-34 and 56.
Their degree yields
\[
                   j=bk+\frac{bnJ}{m},\qquad m\mid bJ.
 \tag{R0E.80}
\]
Thus the possible positive windings are multiples of \(m\) in the
M6 row and of \(l\) in the M2 row. We now construct their actual
primitive maps, including all cap characters.

For M2 the finite map \(S=tU_G\) extends by R0E.34. At the selected
node its full point coefficient is
\[
 C_S(t)=t\,B_{2,1}(p+G)^{-1}
              \frac{f_{P^2}(PZ;Q)}{f_{P^2}(Z;Q)}.
 \tag{R0E.81}
\]
It is constant on the compact source fibre by the full Hom argument,
and both it and its reciprocal extend at \(q=0\). Its actual
character map is \((Z,V)\mapsto(PZ,C_SV)\); on
\(Z^uV^vq^d\) the multiplier is \(P^uC_S^v\). It carries the same
fan with the unchanged height and complete units. Chronologically,
\[
                    S^3=t^3U_G^3=qB_{\rm sel}^{-1}.
 \tag{R0E.82}
\]
On the good root R0E.67, define \(U'_G=(B'_{2,1})^{-1}\) with
\(B'_{2,1}(p)=\phi(p)\phi(p-G')\) and the full \(\phi\) of
R0E.69. The coefficient of \(B'_{2,1}(p)\) at the central curve is
\(-X(p)\) in the canonical point frame. The actual inverse is
\[
 U'_G(p,v_M)=\bigl(p+G',B'_{2,1}(p+G')^{-1}v_M\bigr),
 \qquad B'_{2,1}(p+G')^{-1}\big|_0
   =-\frac1{X(p+G')}=\frac{Y+1}{2X}.
\]
Here \(X(p+G')=2(1-Y)/X^2=-2X/(Y+1)\); thus the coefficient
is evaluated at the literal inverse source point. The actual
Hom Cartier modules supply both maps\textquotesingle{} zeros and poles. Their
generic central order is zero, so each Hom map and its inverse
extend everywhere on the root. The complete
ordinary/root coefficient square is
\[
 f_{2,\rm old}=z^{-l}\phi,\quad
 B_{2,1,\rm old}=z^{-2l}B'_{2,1},\quad
                     U_{G,\rm old}=z^{2l}U'_G .
 \tag{R0E.83}
\]
It is read on the common graph R0E.68 with both actual correction
monomials, source and translated point modules retained.

The full character computation is
\[
 \phi(Dp)=\zeta_m^l\phi(p),\quad
 B'_{2,1}(Dp)=\zeta_m^{2l}B'_{2,1}(p),\quad
 U'_G D=\zeta_m^{-2l}D U'_G.
 \tag{R0E.84}
\]
The source translations of \(U'_G\) and the full norm lift
\(\widetilde F\) commute. In a common constant-factor point frame
both line coefficients are base constants, specified by their
entire rational maps; their lifted affine translations thus
commute, including their complete values. The root \(\beta^{1/m}\)
is also a base function. With
\(A_R=R^n\zeta_m^d\widetilde F^{\,n}D\), this gives the exact square
\[
       (U'_G)^n A_R=\zeta_m^{-2ln}A_R(U'_G)^n
                    =\zeta_m^{ln}A_R(U'_G)^n .
 \tag{R0E.85}
\]

For the positive scalar winding \(J=l\), so \(R\mapsto\zeta_3R\),
the forward endpoint is
\[
 E_R^{\rm ord}=(R^{-l}tU_G)^n,\qquad
 E_R^{\rm cap}=(R^{-l}U'_G)^n.
 \tag{R0E.86}
\]
For \(n=-1\) these are the actual inverse maps with the reciprocal
at the inverse source point. They satisfy, with the entire
character \(d\),
\[
 E_R^{\rm cap}A_R=A_{\zeta_3R}E_R^{\rm cap},\qquad
 \Phi_{\theta+l}E_R^{\rm ord}\Phi_\theta^{-1}
 =z^{nl}R^{-ln}t^n z^{2ln}(U'_G)^n
 =E_R^{\rm cap}.
 \tag{R0E.87}
\]
The last cancellation uses precisely the original \(t=z^{-3l}\).
Its inverse square retains the same units and all reference
periods. Therefore R0E.86 is a holomorphic endpoint on each
prescribed fan and on the full completed object.

The three chronological endpoints use
\(R,\zeta_3R,\zeta_3^2R\). Their scalar factors multiply to
\(\lambda^{-n}\), including the successive phases. Their full
composites are
\[
 \begin{split}
 E_{\zeta_3^2R}^{\rm ord}E_{\zeta_3R}^{\rm ord}E_R^{\rm ord}
     &=\lambda^{-n}S^{3n}=q^n B_\lambda^{-n},\\
 E_{\zeta_3^2R}^{\rm cap}E_{\zeta_3R}^{\rm cap}E_R^{\rm cap}
     &=\lambda^{-n}(U'_G)^{3n}=h^n B_\lambda^{-n}.
 \end{split}                                                     \tag{R0E.88}
\]
Thus on the actual quotient their endpoint is the canonical
\(q^n/h^n\) pair. In the M6 row this pair itself is the primitive
endpoint at \(J=m\): R0E.77 gives
\(z^{nm}q^n=(q_\infty q)^n=h^n\) on the overlap, with inverse
at the return. Its selected zero is carried by the actual
\(B_\lambda^n\) reindex and reciprocal full source/line unit.
Specifically in R0E.50,
\(q^n=B_\lambda^n(P^{bn},\eta^{-n})\); this latter ordered pair
means its actual source and line torus-unit map. Its character
arrows and inverses extend on every selected cone. Thus the zero
is retained by that literal quotient square, for either sign.
Chronological powers and inverses of these primitive maps give
every positive-part winding in R0E.80.

\par

\medskip\noindent\textbf{R0E.10. The full negative endpoint obstruction and exact winding group}

\par\nobreak

The negative argument uses the actual selected inversion and the
whole annular reference. On R0E.50 let
\[
                    I(Z,V,q)=(Z^{-1},V^{-1},q).
 \tag{R0E.89}
\]
It satisfies \(Ih=h^{-1}I\) and \(IB=B^{-1}I\), with all coefficient
functions retained at the unchanged \(q\). For a periodic fan of
basis \((u,v)\), negation sends its triangle
\(\operatorname{conv}(p,p+u,p+u+v)\) to
\(\operatorname{conv}(p',p'+v,p'+u+v)\), where
\(p'=-p-u-v\). This is the other triangle of the same periodic fan.
The full chart arrows can be read directly. With
\(\mathcal A=\kappa g,\mathcal C=\alpha\eta\), the representative
\(hBI\) sends
\[
 (a,b,c)\mapsto(\mathcal A\mathcal C a,\mathcal C^{-1}c,
                         \mathcal A^{-1}b),\qquad
 (u_0,v_0,w_0)\mapsto((\mathcal A\mathcal C)^{-1}u_0,
                          \mathcal A w_0,\mathcal C v_0).
 \tag{R0E.90}
\]
The first source chart has \(q=abc,Z=b,V=c\), the second
\(q=u_0v_0w_0,Z=u_0v_0,V=u_0w_0\). The displayed maps interchange
these charts and are literal inverses. Both preserve the full
product \(q\) and give \(Z'=\mathcal A q/Z,V'=\mathcal C q/V\).
For a general basis the actual integral character change carries
these equations with its complete period products. This proves
the inversion extends to the full selected quotient.

A negative endpoint lifts locally to the nonzero-line cover and
satisfies
\[
                        \widetilde F B_\lambda
                              =B_\lambda^{-1}\widetilde F.
 \tag{R0E.91}
\]
Any two local lifts differ by an integral suspension deck. At the
ordinary III disc the actual cover and Cartier maps are R0E.26-\/-34.
A nonzero line vector over the \(O\)-section of the entire disc
kills the base meridian. Generic fibre loops then generate modulo
the actual vanishing cycles in \(\ker\psi\). The preserved
negative cover class gives a single-valued holomorphic lift,
including its local biholomorphic inverse through the two
Cartier modules. At the selected point compose with R0E.89.
The result has positive linear part. R0E.50-\/-52 remove its
integer valuations using the entire \(h^rB^k\), and give a
single-valued bounded unit map and inverse on the punctured disc.
Composing back with R0E.89 gives a single-valued negative lift
there. The finite base with the selected point removed has
fundamental group generated by this one meridian. The ordinary
discs add no puncture. Continuing the local lifts therefore
gives one single-valued map on the whole ordinary nonzero-line
cover, retaining their constant joining deck integers.

On the actual normalized root fibre product, the kernel of the
finite generator is \(\langle B_\lambda\rangle\) because
\(A_\lambda^m=B_\lambda^n\), \(n=\pm1\). Its root nonzero-line
cover has the analogous \(O\)-section vector over the whole
\(z\)-disc. On a point/O intersection the frame is exactly
\(\sigma_a/[(\tau-\tau(aG'))/\tau]\), and elsewhere it is the
single-divisor frame; the transition is the actual Cartier unit
of R0E.65 or 68. The same meridian and vanishing-cycle argument
therefore supplies a single-valued root lift
\(\widetilde F_{\rm cap}\), fixing \(z\) and satisfying R0E.91.

Follow one positive original \(q_\infty\)-circle through every
reference sector in R0E.74. Its actual arrow is
\[
 A_\lambda B_\lambda^\nu,\qquad \nu\in\mathbb Z,\qquad
 (A_\lambda B_\lambda^\nu)^m=B_\lambda^N,\quad N=n+m\nu.
 \tag{R0E.92}
\]
The integer \(\nu\) is the retained sum of reference suspension
integers; all analytic source/line factors and period columns
are already in these full operators. No value is imposed on
\(\nu\). The scalar family R0E.76 retains it at both endpoints.
At winding \(J\), its extra \(z^{nJ}\) is single valued on the
root circle, so both endpoint annular maps have this same
monodromy.

Choose one constant suspension deck so that on a cut annulus
the actual lifts satisfy
\[
 \widetilde F_{\rm cap}
   =\Phi_{\rm target}\widetilde F_{\rm ord}\Phi_{\rm source}^{-1}.
 \tag{R0E.93}
\]
After one positive root circle its right side continues to
\[
 B_\lambda^N\widetilde F_{\rm cap}B_\lambda^{-N}
                   =B_\lambda^{2N}\widetilde F_{\rm cap}.
 \tag{R0E.94}
\]
Both actual lifts were proved single valued, so
\(B_\lambda^{2N}=1\). Strict contraction implies \(N=0\), whereas
\(n+m\nu=0\) has no integral solution for \(n=\pm1,m\ge2\).
This excludes every negative endpoint, including the quadratic
reduced cap, without dropping the reference integer.

The completed classification on the fixed original \(r=0\) source
is consequently
\[
 \begin{array}{c|c|c|c}
 (a,b;m)&\text{winding group}&\text{primitive ordinary forward map}
                                &\text{primitive cap map}\\ \hline
 (6,1;3l+2),\ l\ge0&m\mathbb Z&q^n&h^n\\
 (2,3;3l),\ l\ge1&l\mathbb Z&(R^{-l}tU_G)^n&(R^{-l}U'_G)^n .
 \end{array}                                                       \tag{R0E.95}
\]
It holds for every prescribed lattice-basis selected fan, both
norm signs and every retained finite character and reference.
The ratio of two endpoints with the same winding is exactly
the physical scalar group R0E.58. Each inverse return uses the
inverse of the displayed forward endpoint. R0E.88 is their
actual chronological threefold comparison. This proof assigns
no new numerical return.

\par

\medskip\noindent\textbf{R0E.11. Exact verification and preserved stages}

\par\nobreak

The original 47-check finite-chart stage, before this extension,
is preserved byte for byte under
\path{records/attempts/e6_iii_endpoint_before_global_extension/}.
Its recorded scope was R0E.1-\/-3. The added proof retains that
source map, proves the complete second marked-line and tangential
chart calculation, and supplies the entire global endpoint
classification in R0E.4-\/-10. The companion exact program checks
the full rational equations, ideal/chart and point-frame
identities, period-coordinate matrices, integral source lattice,
both root powers and characters, and the signed endpoint and
reference products. Its finite algebra is accompanied by the
explicit compactness, module-extension, covering and
single-valuedness arguments above. The final source and executed
certificate are submitted for independent full audit.

\par\medskip\noindent\textbf{Exact source and independent verification.}\par
The full mathematical source is
\path{notes/part3_e6_iii_fixed_source_endpoint_working.md}.
The primary exact certificate passes 9532 checks, including 116 full
integral elliptic basis changes and 4838 signed whole-annulus
instances. The independent finite-source proof and certificate pass
49 checks on the exact R0E.1--4 prefix:
\path{records/checks_audit_e6_iii_finite_source.json}.
All full-scope review gates and their final source locks are recorded
in \path{records/k114_export_manifest.json}.
The historical finite-stage files and the preserved algebraic-number
checker evaluation are retained in the two attempt directories named
in R0E.11. The original source bytes are unchanged by this export.

% End expanded verified cp36 input

\par\bigskip\noindent\textbf{Full independent root proof of the original III endpoint classification}\par\nobreak

% Expanded verified cp36 input: notes/part3_e6_iii_endpoint_root_peer.tex
We retain the original surface, line, selected height, finite character,
reference and fan of R0E.1-\/-95. In particular
\[
 y^2=x^3+tx+t^2,\quad G=(0,t),\quad q=t+4/27,
 \qquad (a,b;m)=(6,1;3l+2),\ l\ge0,
 \quad (2,3;3l),\ l\ge1,
 \tag{R0R.1}
\]
and the norm is \(n=\pm1\). The finite III disc retains the actual
order-one cap and zero norm of K96. Each scalar is in the actual contraction
domain of the complete K96 construction. The assertions below concern
these completed objects and their stated scalar paths. We prove their
full automorphism group and their scalar endpoint winding groups, with
the original-coordinate maps throughout. We assign no return integer.

\par

\medskip\noindent\textbf{R0R.1. Reconstructing the finite source and every point-line factor}

\par\nobreak

Put \(u=y-t-x/2,v=y+t+x/2\). Direct multiplication gives
\(uv=x^2(x-1/4)\), with inverse
\(y=(u+v)/2,t=(v-u-x)/2\). On \(x\ne0\), write
\(u=xu_1,v=xv_1\); then
\[
 x=u_1v_1+1/4,\qquad y=x(u_1+v_1)/2,
 \qquad t=x(v_1-u_1-1)/2.                       \tag{R0R.2}
\]
The reciprocals are the displayed \(u/x,v/x\). These are actual
birational source morphisms on their indicated open sets.

The A1 point of the original surface is resolved in its \(x\)-chart by
\[
 (x,Y_s,T_s)\longmapsto(x,xY_s,xT_s),\qquad
 x=Y_s^2-T_s-T_s^2.                              \tag{R0R.3}
\]
On \(x\ne0\) the inverse is \((x,y/x,t/x)\). The original fibre
function is \(t=xT_s\). Its two components have equations \(x=0\)
and \(T_s=0\); their intersection ideal has quotient
\(\mathbb C\{Y_s\}/(Y_s^2)\). Their intersection multiplicity is
two. Pairing the whole fibre with either component therefore gives
their intersection matrix \(\left(\begin{smallmatrix}-2&2\\2&-2
\end{smallmatrix}\right)\), with primitive kernel \((1,1)\).

The actual quartic source cover is
\[
 t=z^4,\quad x=z^2X,\quad y=z^3Y,\qquad
 Y^2=X^3+X+z^2,
 \quad D(z,X,Y)=(iz,-X,iY).                       \tag{R0R.4}
\]
The equation pulls back with the exact factor \(z^6\). The three
original coordinates are invariant under \(D\). In R0R.3 its map is
\(T_s=z^2/X,Y_s=zY/X\), defined on \(X\ne0\), including the two
central root points \((X,Y)=(i,0),(-i,0)\). Substitution gives
the full R0R.3 equation. Thus this chart includes the tangency in the
original fibre, with its full nonreduced intersection scheme.

We reconstruct the Miller functions from addition. The tangent slope
at \(G\) is \(1/2\). Its double and its chord with \(G\) give
\[
 2G=(1/4,-t-1/8),\quad L_{2G,G}=-8t-1/2,
 \quad 3G=(8t+64t^2,3t+96t^2+512t^3).             \tag{R0R.5}
\]
Write this last point as \((x_3,y_3)\), put
\(L_3=(3x_3^2+t)/(2y_3)\), and set
\(x_6=L_3^2-2x_3,y_6=L_3(x_3-x_6)-y_3\). The entire factors are
\[
 f_2=\frac{y-t-x/2}{x-1/4},\quad
 g_{2G,G}=\frac{y+t+1/8+(8t+1/2)(x-1/4)}{x-x_3},
 \quad f_3=f_2g_{2G,G},
 \tag{R0R.6}
\]
\[
 g_{3G,3G}=\frac{y-y_3-L_3(x-x_3)}{x-x_6},
 \qquad f_6=f_3^2g_{3G,3G}.                       \tag{R0R.7}
\]
Reduction by the original cubic identifies R0R.6 with the entire
\(f_3\) of R0E.23. Its divisor is \(3(G)-(3G)-2(O)\); the divisor
of R0R.7 is \(6(G)-(6G)-5(O)\). Thus these coefficients define
the required point-to-tensor Hom maps, with their actual point and
zero Cartier modules. The computation retains the rational functions
before taking any special-fibre value.

In R0R.4 the translation by \(-G'\), with \(G'=(0,z)\), is
\[
 L=(Y+z)/X,\quad X_T=L^2-X,
 \quad Y_T=z-LX_T.                               \tag{R0R.8}
\]
At the image point the slope for \(+G'\) is exactly \(-L\);
its two coordinates return \((X,Y)\). Both maps preserve the full
root equation. At \(O\), put \(X=U/\eta^2,Y=-U/\eta^3\), where
\(U^2-U^3-\eta^4U-z^2\eta^6=0,U(0,0)=1\). Then
\[
 X_T=\eta(\eta-2z)/U+2z^2\eta^4/U^2,
 \quad Y_T=z+(1/\eta-z\eta^2/U)X_T.              \tag{R0R.9}
\]
This last expression is holomorphic with first jet \(\eta-z\).
It identifies the ideals \((Y_T,z)=(\eta,z)\). At \(T=(0,0)\),
the expression R0E.32 has unit denominator and first jet \(Y-z\),
and identifies \((\eta_T,z)=(Y,z)\). Their inverse arrows have
the inverse ideal equalities. Hence these maps lift to both actual
blowups, in both directions, through every point of their exceptional
curves. R0E.17,20,21 and R0R.3 give the original-coordinate maps of
these charts. They cover the entire compact central root.

For \(a=2,6\), the marked point has
\(\eta(aG')=z\tau_{a,\mathrm{old}}(z^4)\), with finite limits
\(\tau_{2,\mathrm{old}}(0)=2\) and
\(\tau_{6,\mathrm{old}}(0)=6\). The two graph charts therefore
carry the exact point ratios
\[
 \frac{p-\tau_{a,\mathrm{old}}(z^4)}p,
 \qquad 1-v\tau_{a,\mathrm{old}}(z^4),
 \quad \eta=zp,\quad z=\eta v.                  \tag{R0R.10}
\]
Their inverses are the reciprocal Cartier-module arrows. Both marked
divisors have the same exceptional order; the two strict directions
remain separate. This proves the comparison of the two point lines
on the graph. The other graph centre and the tangency chart carry no
marked divisor, so their coefficient is one.

On the identity component use \((x,y)=(p^{-2},-p^{-3})\); on the
other use \((x,y)=(t\alpha,t\beta)\), with
\(\alpha+t\alpha^3=\beta^2-1\). R0R.8 gives the directed component
maps \(\beta'=p-1\) and \(p'=\beta-1\). Substituting the full
functions R0R.6-\/-7 first gives
\[
 \begin{array}{c|cc}
 &C_0&C_1\\\hline
 f_2&2/[p(p-2)]&2t(\beta-1)^2+O(t^2)\\
 f_3&p^{-2}&t(\beta-1)^3/(\beta-3)+O(t^2)\\
 g_{3G,3G}&6/[p(p-6)]&6t(\beta-3)^2+O(t^2)\\
 f_6&6/[p^5(p-6)]&6t^3(\beta-1)^6+O(t^4).
 \end{array}                                                       \tag{R0R.11}
\]
Consequently \(B_{a,1}=f_2^a f_a(Tp)/f_a(p)\), after the exact
whole-fibre factor \(t^{a/2}\), has central coefficients
\[
 2^a\frac{p-a}p\quad\hbox{and}\quad
 2^a\frac{\beta-1}{\beta-a-1}.                   \tag{R0R.12}
\]
Using the actual point-frame factors \(r_{a,0}=(p-a)/p,r_{a,1}=1\),
the Hom coefficient is the raw factor times \(r_a(Tp)/r_a(p)\).
It is \(2^a\) on both components. The exact graph generators in
R0R.10 retain this cancellation at the section points. The primitive
component kernel and the zero horizontal Hom divisor show that the
remaining rational section and its reciprocal have zero divisor on
the smooth resolved source. In each local regular ring they are
units, including at the tangency. This proves the complete finite
extension. Its chronological \(b=6/a\) product has coefficient 64.
The coefficient \(128(3w-1)\) in the earlier source chart equals
\(64(p-6)/p\) under the actual map \(w=1/2-1/p\).

\par

\medskip\noindent\textbf{R0R.2. The full invariant quotient, original base and primitive section}

\par\nobreak

The actual K96 disc family supplies the completed topology at this
source, with the full finite charts just proved and its own infinity
root, not the finite quartic root. The integral F35/F46 maps carry
\(H^2=H^4=0\) through each prescribed fan: their primitive normal
circle is unchanged by the common blowups and the two ruling maps.
On a compact scalar interval the original finite, root and annular
arrows have a common strict contraction; their compact gluing gives
the corresponding smooth transport. Thus these are the same sphere
objects used throughout R0R.1.

A horizontal effective divisor has positive pairing with a flat
Kähler form on a generic compact two-torus, but its class is the
restriction of a class in \(H^2(X)=0\). Hence all polar divisors of
meromorphic functions are vertical. Restriction to each regular
compact torus and removal of indeterminacy give constancy there.
At a special component the actual transverse map is
\(u\mapsto t_0+u^e\). Invariance of its Laurent expansion makes all
exponents multiples of \(e\), proving meromorphic descent on the
whole base. Thus \(\pi^*\mathbb C(t)=\mathbb C(X)\), and a
holomorphic endpoint isomorphism and its inverse induce a Möbius
base isomorphism \(h\).

We retain the complete point-theta module R0E.35-\/-44, whose periods
are \(\gamma=(-bg,\ell),e_1=(1,0),e_2=(\tau,ag),\delta=(0,1)\).
Its full monodromy is
\[
 M_\alpha=\begin{pmatrix}1&0&0\\bv_\alpha&A_\alpha&0\\
 c_\alpha&av_\alpha^tJA_\alpha&1\end{pmatrix},
 \qquad J=\begin{pmatrix}0&-1\\1&0\end{pmatrix}.   \tag{R0R.13}
\]
Both point Cartier modules, the full Miller-to-theta constant, and
every central cochain entry remain in this formula. Multiplication
gives \(v_{\alpha\beta}=v_\alpha+A_\alpha v_\beta\) and
\(c_{\alpha\beta}=c_\alpha+c_\beta+
abv_\alpha^tJA_\alpha v_\beta\). The companion checks these as
symbolic identities for arbitrary matrices of determinant one.

The Abel class is nonzero: a rational coboundary, after multiplication
by an integer and the finite component orders, gives a global
holomorphic Lie logarithm of a multiple of \(G\). At III the
logarithm of \(2G'\) divided by \(z\) extends by R0R.4-\/-5; the
quartic central action has no invariant elliptic period. At the
selected node its invariant annular period extends in \(xy=q\).
At infinity the actual good cubic is \(Y^2=X^3+zX+1\), the logarithm
of \(3G'\) has leading term \(-z^2/8\), and division by \(z^2\)
is the original Lie coefficient. Its order-three elliptic action
also has no invariant period. The full Lie transition is
\(u_0(q_\infty)=q_\infty^{-1}u_t(q_\infty^{-1})\). An entire
\(u_t\) satisfying it tends to zero at infinity and is zero.
But the original infinity value of \(3G\) is \((64,512)\), with
nonzero additive parameter \(-1/8\); every nonzero multiple remains
nonzero. This contradiction proves both non-torsion and the claimed
nonzero Abel class, with its actual finite extensions.

The vector and dual equations in R0R.13 now give exactly
\(\Lambda^{\pi_1}=\mathbb Z\delta\) and
\((\Lambda^*)^{\pi_1}=\mathbb Z\psi\), where
\(\psi(\gamma)=1\) and \(\psi(e_1)=\psi(e_2)=\psi(\delta)=0\).
Indeed any nonzero first vector coefficient makes the Abel cocycle a
coboundary; the remaining elliptic vector vanishes by the infinity
matrix. The dual equation has the same Abel obstruction. The extra
deck of the actual point-line cover is the primitive \(\gamma\).

An endpoint preserves these two integral subgroups with their signs
before fixing the base. The lifted derivative is periodic for the
whole compact torus lattice, so it is constant. Complex linearity
and \(\delta\mapsto\epsilon\delta\) give the complete affine map
\[
 (z,w)\mapsto(Az+z_0,Cz+\epsilon w+w_0),
 \quad (z',w')\mapsto
 ((z'-z_0)/A,\epsilon(w'-w_0-C(z'-z_0)/A)),         \tag{R0R.14}
\]
with \(A\ne0,\epsilon=\pm1\). The kernel cover is precisely the
actual nonzero point line, using both Cartier modules. On its lattice
the complete matrix is
\[
 K=\begin{pmatrix}p&q_0&0\\r&s&0\\u&v&\epsilon\end{pmatrix},
 \quad K^{-1}=\begin{pmatrix}s&-q_0&0\\-r&p&0\\
 \epsilon(vr-us)&\epsilon(uq_0-vp)&\epsilon\end{pmatrix},
 \quad ps-q_0r=1.                               \tag{R0R.15}
\]
The last determinant sign follows from the positive oriented real
area of a nonzero complex-linear map. The first coordinates of
R0R.14 and its inverse are actual inverse elliptic isomorphisms;
the integers \(u,v\) remain in the line arrow. For this exact
source \(j_0(t)=6912/(27t+4)\). Therefore
\(j_0(h(t))=j_0(t)\) gives \(h(t)=t\) on the regular locus and
hence on the whole base.

The full linear centralizer is \(eI+c\delta\psi\). For its elliptic
block, the order-three infinity matrix generates \(\mathbb Q(\zeta_3)\),
and the nontrivial selected unipotent matrix has a nonzero nilpotent
part. No nonzero nilpotent lies in that field. Every nonscalar field
element generates it, so the simultaneous elliptic centralizer is
\(\mathbb QI\). The two off-diagonal equations then use the nonzero
Abel class to set both adjoining scalars equal to \(e\) and both
elliptic vectors to zero. All central cochains remain arbitrary in
the surviving formula. Unimodularity gives \(e=\pm1\). Complex
linearity excludes the rank-one remainder when the fibre periods
coincide. A negative automorphism of one object has sixteen transverse
fixed points on each regular fibre. Its closed analytic equalizer in
the compact total space has a horizontal curve of positive base
degree. Pairing that curve with \(\pi^*H^2(\mathbb P^1)=0\) contradicts
that degree. Thus the automorphism has positive linear part.

The actual cyclic-cover, vanishing-cycle and selected-vertex unit maps
R0E.50-\/-52 extend its source displacement to a global source section.
At the finite III disc the nonzero \(O\)-section vector kills the
meridian, and all remaining vanishing cycles lie in \(\ker\psi\).
At the cap the normalized fibre product \(z^m=q_\infty\) has its
literal root lift. Its commutator with the cap generator is \(B^k\);
the \(m\)-fold relation makes \(B^{mk}=1\), so strict contraction
sets \(k=0\). The integer joining torsor on the one-punctured base
with one finite root has zero homomorphism to \(\mathbb Z\), since
its orbifold fundamental group is finite. These are the actual
single-valued lifts used to extend the section.

Its primitive generator is \(G\). The rational source R0R.2 has
fibres IV*, III and I1 and Euler number 12. The two sections are
disjoint, including their different infinity quotient arms. The
III A1 direction of \(G\) is \([x:u:v]=[0:0:1]\), off the strict
fibre plane \(v-u-x=0\), so its correction is \(1/2\). Its E6
arm correction is \(4/3\), giving height \(1/6\). More fully, the
ten actual classes \((O,F,C_0,\ldots,C_5,C_{\rm III},G)\), with the
pairings stated in R0E.53, have determinant \(-1\). They therefore
span the whole rank-ten unimodular divisor lattice. Quotienting by
the first nine classes gives the full free section group \(\mathbb ZG\),
with no torsion or index omitted. The extending section is \(jG\).

\par

\medskip\noindent\textbf{R0R.3. The actual line map and every positive scalar endpoint}

\par\nobreak

The lift over \(+jG\) is \(H(t)U_G^j\), where
\(U_G(p,v)=(p+G,B_{a,1}(p+G)^{-1}v)\). The full finite source
calculation gives its order \(-3/b\) at each III component. The
original infinity root maps are exactly
\[
 \begin{array}{c|c|c}
 m&t,x,y&\text{root equation}\\\hline
 3l+2&z^{-m},z^{-2(l+1)}X,z^{-3(l+1)}Y
     &Y^2=X^3+z^{l+2}X+z^2\\
 3l&z^{-3l},z^{-2l}X,z^{-3l}Y
     &Y^2=X^3+z^lX+1.
 \end{array}                                                       \tag{R0R.16}
\]
The reduced source has the actual A2 coordinates and inverse chord
maps R0E.60-\/-63; the point graph retains
\(\tau'/(\tau'-z^{2l+1}L)\), with \(L(0)=-1/4\). The good root
has the two full corrections
\(E_{2G'}-E_O\), with \(E_P(v_w,v_z)=\min(v_w,2lv_z)\), and the
actual ratio of their monomial generators. These are the complete
maps R0E.64-\/-70, including all strict-section units and their inverses.
Their full coefficient orders give \(2m/b\) for \(U_G\) at the root.
The normalized chord lift and its central-one factor yield the full
norm \(B'_{\rm sel}\), central value four; taking its prescribed root
retains the full finite character and gives
\(A_R=R^n\zeta_m^d\widetilde F^{\,n}D\),
\(A_R^m=B_R^n\), with \(R^m=\lambda\). The original selected
height factor is \(h=1+4z^m/27=z^mq\).

The full finite and root divisor equations for an automorphism are
\(\operatorname{ord}_0H=3j/b\),
\(\operatorname{ord}_\infty H=-2j/b\), and
\(\operatorname{ord}_{t_c}H=-k\). The last integer is its literal
selected \(B^k\) reindex. There are no other divisors. Their degree
gives \(j=bk,H=c(t^3/q)^k\). Consequently the full quotient map is
the physical scalar \(c\lambda^k\). These scalars extend through
every root and fan by the stated unit maps. Their kernel is trivial,
since an identity scalar deck would translate the source by
\(-kbG\) and \(G\) is not torsion. Thus
\(\operatorname{Aut}_{\rm hol}(X)=\mathbb C^*_{\rm physical}\).

For an endpoint the two independent period vectors \(e_1,\delta\)
force its complex-linear map to be \(eI\). Its suspension equation
retains an integral central difference and forces equality of the
two original scalar values. Write the resulting path winding as
\(J\in\mathbb Z\). The full annular reference R0E.72-\/-77 gives
the extra root factor \(z^{nJ}\); its integral suspension labels
have ratio \(e^{2\pi iJ J_\alpha}=1\). The divisor degree is now
\[
 3j/b-k+(-2mj/b-nJ)/m=0,
 \qquad j=bk+bnJ/m.                              \tag{R0R.17}
\]
Thus \(m\mid bJ\). All constants and reference cochains are retained
in the same ordinary/root comparison.

In the good row the finite map is \(S=tU_G\). On its good root,
\(f_{2,\rm old}=z^{-l}\phi\), where
\(\phi=(Y-1-z^lX/2)/(X-z^{2l}/4)\), and
\(B'_{2,1}(p)=\phi(p)\phi(p-G')\). The correct inverse is
\[
 U'_G(p,v)=(p+G',B'_{2,1}(p+G')^{-1}v).
 \tag{R0R.18}
\]
On the central cubic \(Y^2=X^3+1\), the two source maps have
\(X_+=2(1-Y)/X^2\) and \(X_-=2(1+Y)/X^2\).
Their actual line coefficients are respectively
\((Y+1)/(2X)=-1/X_+\) and \(-X\). Both compositions have coefficient
one at the inverse source point. Their three chronological products
have coefficients \(1/4\) and 4. This checks the inverse typing
without replacing its complete meromorphic Hom maps by those values.

The full root square is \(U_{G,\rm old}=z^{2l}U'_G\), read with
both corrections in R0R.16. Its deck relation is
\(U'_GD=\zeta_m^{-2l}DU'_G\). For \(J=l\) the actual forward
maps are
\[
 (R^{-l}tU_G)^n\quad\hbox{and}\quad(R^{-l}U'_G)^n.
 \tag{R0R.19}
\]
They intertwine the full finite generator with that at
\(R\mapsto\zeta_3R\), including \(\zeta_m^d\), since
\(-2ln\equiv ln\pmod {3l}\). The full annular square has product
\(z^{nl}t^nz^{2ln}=1\), using the original \(t=z^{-3l}\).
For the negative norm every factor is evaluated at its actual inverse
source point. The three chronological root phases give
\(\lambda^{-n}\); using \(S^3=qB_{\rm sel}^{-1}\), the composite
is the full quotient pair \(q^n,h^n\). In the reduced row that pair
itself is primitive at \(J=m\), since \(z^{nm}q^n=h^n\). Its
selected zero is carried by the literal \(B^n\) reindex and the
nonzero torus-unit map, with inverse, on every fan. Powers and inverses
give every multiple allowed by R0R.17.

\par

\medskip\noindent\textbf{R0R.4. Single-valued negative lifts and the completed endpoint group}

\par\nobreak

The selected inversion \((Z,V,q)\mapsto(Z^{-1},V^{-1},q)\) reverses
both full periods. Negating a triangle in a lattice-basis periodic fan
gives its opposite triangle with translated vertex \(-p-u-v\).
The actual chart maps R0E.90 exchange the two vertices, preserve
\(q\) and retain \(\kappa g,\alpha\eta\) and their reciprocals.
Thus this is an extending selected map for every stated fan.

A negative endpoint has a single-valued ordinary nonzero-line lift:
the entire finite III cover is proved in R0R.1, its \(O\)-section
vector kills the base meridian, and the other vanishing cycles lie
in \(\ker\psi\). Composing with the selected inversion gives the
actual positive selected-unit map and its inverse, so its continuation
has no additional selected monodromy. Constant integral deck joins
give the ordinary lift globally. On the normalized infinity root
fibre product, the same \(O\)-section vector in the actual point/O
frame R0E.65 or 68 gives a single-valued root lift. Both lifts reverse
the full suspension: \(\widetilde F B=B^{-1}\widetilde F\).

The complete annular reference around an original infinity circle
is \(AB^\nu\), with its actual retained integer \(\nu\). On a
root circle it is \(B^N\), where \(N=n+m\nu\). At both scalar
endpoints these are the same operators by the full reference
restoration. Choose one constant deck to equate the two lifts on
a cut annulus. Continuing that equality around the root circle gives
\[
 B^N\widetilde F B^{-N}=B^{2N}\widetilde F.
 \tag{R0R.20}
\]
Single-valuedness forces \(B^{2N}=1\), and strict contraction gives
\(N=0\). But \(N\equiv n\equiv\pm1\pmod m\), with \(m\ge2\),
so this is impossible for every retained reference integer. This
also treats the quadratic reduced cap. All negative endpoints are
excluded by their actual ordinary/root comparison.

The full winding group is consequently \(m\mathbb Z\) in the
reduced M6 row and \(l\mathbb Z\) in the good M2 row, with primitive
maps R0R.19 and the \(q^n,h^n\) pair just constructed. At a fixed
winding, the ratio of any two endpoints is exactly a physical scalar
automorphism. The proof retains both norm signs, every finite
character, every prescribed complete reference and every selected
lattice-basis fan. The executable peer independently derives the
Miller factors and both inverse-source compositions; its full kernel
and central-cochain equations are symbolic identities for all their
integer entries. These exact checks support the full geometric proof
above and the complete R0E source; they do not replace its extension,
compactness or single-valuedness arguments.

% End expanded verified cp36 input

\par\bigskip\noindent\textbf{Complete independent source, point-line, root and endpoint proof}\par\nobreak

% Expanded verified cp36 input: notes/part3_e6_iii_endpoint_fourfold_peer.tex
The source and its height are fixed literally as
\[
 E_0:\ y^2=x^3+tx+t^2,\qquad G=(0,t),\qquad
 t_c=-4/27,\qquad q=t+4/27.
 \tag{R0P.1}
\]
The completed objects retain the ordinary order-one, zero-norm III
attachment of K96 and either infinity cap
\((a,b;m)=(6,1;3l+2)\), \(l\ge0\), or
\((2,3;3l)\), \(l\ge1\), with \(n=\pm1\).
The genuine point line is \(M_a=\mathcal O([aG]-[O])\), the
suspension translates by \(-bG\), and all finite characters,
annular reference data, prescribed lattice-basis fans and admitted
strictly contracting scalars are retained. This independent proof
checks the complete R0E.1-\/-11 and R0R.1-\/-4 maps on that scope.

\par

\medskip\noindent\textbf{R0P.1. The whole finite source and its inverse graph maps}

\par\nobreak

The change
\(U=y-t-x/2\), \(V=y+t+x/2\) gives
\(UV=x^2(x-1/4)\), with inverse
\(t=(V-U-x)/2\), \(y=(U+V)/2\). At its A1 point the actual
\(x\)-chart of the resolution is
\[
 x=Y_s^2-T_s-T_s^2,\qquad
 (x,y,t)=(x,xY_s,xT_s).
 \tag{R0P.2}
\]
The fibre is \(xT_s=0\). The intersection algebra of its two
components is \(\mathbb C\{Y_s\}/(Y_s^2)\), so their intersection
number is two. Their total fibre has zero intersection with each
component, giving the matrix
\(\left(\begin{smallmatrix}-2&2\\2&-2\end{smallmatrix}\right)\)
and primitive kernel \((1,1)\).

The complete finite quartic map is
\[
 t=z^4,\qquad x=z^2X,\qquad y=z^3Y,
 \qquad Y^2=X^3+X+z^2,
 \quad D(z,X,Y)=(iz,-X,iY).
 \tag{R0P.3}
\]
The original equation pulls back with factor \(z^6\), and all three
original coordinate functions are invariant under this deck. On
\(X\ne0\), its map to R0P.2 is
\(T_s=z^2/X\), \(Y_s=zY/X\), including the two root points
\((X,Y)=(i,0),(-i,0)\). Substitution gives exactly R0P.2.

Translation by \(-G'\), \(G'=(0,z)\), has the full coefficients
\[
 L=(Y+z)/X,\qquad X_T=L^2-X,
 \qquad Y_T=L(X-X_T)-Y=z-LX_T.
 \tag{R0P.4}
\]
They preserve the whole equation. At the target point the slope for
\(+G'\) is \((Y_T-z)/X_T=-L\); its two coordinates return \(X,Y\).
This gives both actual source inverses before taking a central limit.

At \(O\), write \(X=U/\eta^2\), \(Y=-U/\eta^3\), where
\(U^2-U^3-\eta^4U-z^2\eta^6=0\) and \(U(0,0)=1\).
The exact pullback of R0P.4 is
\[
 X_T=\eta(\eta-2z)/U+2z^2\eta^4/U^2,
 \qquad
 Y_T=z+(1/\eta-z\eta^2/U)X_T.
 \tag{R0P.5}
\]
The first factor in \(X_T\) cancels the displayed \(1/\eta\).
The first jet of \(Y_T\) is \(\eta-z\), and its coefficient of
\(\eta\) is a unit. Thus \((Y_T,z)=(\eta,z)\) as analytic
ideals. At the other root centre set \(W=Y-z\) and
\(A=1-XW^2/(1+X^2)^2\). Reducing R0P.4 by the actual root equation
gives
\[
 -X_T/Y_T=
 \frac{W A(1+X^2)^2}{A(1+X^2)^3-zW^3}.
 \tag{R0P.6}
\]
Its denominator is a unit and its coefficient of \(Y-z\) is a
unit. Hence \((\eta_T,z)=(Y,z)\). The inverse translation supplies
the inverse ideal equalities. The universal property of the blowup
therefore gives morphisms and inverses on both entire projective
graphs, including all their exceptional directions.

The original-coordinate charts of these graphs are precisely
\(\eta=zp\), \(z=\eta v\), and \(Y=z\beta\), \(z=Yv\)
in R0E.17,20,21. In the last pair the source equations are
\(X=z^2\alpha\), \(\beta^2=1+\alpha+z^4\alpha^3\), and
\(X=Y^2C\), \(C(1+Y^4C^2)=1-v^2\), respectively. Together
with the \(X\)-unit chart R0P.2 these maps cover the whole compact
central root and its tangential source intersection.

\par

\medskip\noindent\textbf{R0P.2. Full Miller functions and genuine point-module cancellation}

\par\nobreak

Addition on R0P.1 gives
\[
 2G=(1/4,-t-1/8),\qquad
 3G=(8t+64t^2,3t+96t^2+512t^3).
 \tag{R0P.7}
\]
Write the second point as \((x_3,y_3)\), put
\(\mu_3=(3x_3^2+t)/(2y_3)\), and retain
\(x_6=\mu_3^2-2x_3\),
\(y_6=\mu_3(x_3-x_6)-y_3\). These are the actual doubled
point, not only its divisor. The entire functions are
\[
 \begin{gathered}
 f_2=\frac{y-t-x/2}{x-1/4},\qquad
 f_3=\frac{x(y-t-x/2)/4}{(y+t+x/2)/4-(1/4+2t)x},\\
 f_{2,3G}=\frac{y-y_3-\mu_3(x-x_3)}{x-x_6},\qquad
 f_6=f_3^2f_{2,3G}.
 \end{gathered}
 \tag{R0P.8}
\]
The actual chord formula for \(f_2g_{2G,G}\) reduces to the
displayed \(f_3\) by the original cubic equation. Its divisor and
that of \(f_6\) are the point-to-tensor divisors used in R0E.23.

Under R0P.3 these satisfy exactly
\(f_2=z^2\phi_2\), \(f_3=z^2\phi_3\),
\(f_{2,3G}=z^2\phi_{23}\), \(f_6=z^6\phi_6\), with the
entire rational functions R0E.25. All four \(\phi\)\textquotesingle s change sign
under the stated quartic deck. At its central curve they are
\(2X,X,6X,6X^3\). The full translated coefficient is therefore
\[
 t^{-3}B_{6,1}(p)=
 \phi_2(p)^6\frac{\phi_6(Tp)}{\phi_6(p)}.
 \tag{R0P.9}
\]
Its central Hom coefficient is 64, with its horizontal point and
zero modules retained. This identity follows before any such
restriction and retains the source argument of every factor.

On the identity component use \(x=p^{-2},y=-p^{-3}\); on the
other component use \(x=t\alpha,y=t\beta\), with
\(\alpha+t\alpha^3=\beta^2-1\). The actual source maps have
\(\beta'=p-1\) and \(p'=\beta-1\). Directly in the full
functions R0P.8 their leading coefficients are
\[
 \begin{array}{c|cc}
 & C_0&C_1\\\hline
 f_2&2/[p(p-2)]&2t(\beta-1)^2+O(t^2)\\
 f_3&p^{-2}&t(\beta-1)^3/(\beta-3)+O(t^2)\\
 f_{2,3G}&6/[p(p-6)]&6t(\beta-3)^2+O(t^2)\\
 f_6&6/[p^5(p-6)]&6t^3(\beta-1)^6+O(t^4).
 \end{array}
 \tag{R0P.10}
\]
Multiplication with those directed source maps gives
\(t^{-a/2}B_{a,1}=2^a(p-a)/p\) on \(C_0\), and
\(2^a(\beta-1)/(\beta-a-1)\) on \(C_1\), for \(a=2,6\).
The genuine frame is \(e_{a,i}=\sigma_a/r_{a,i}\), where
\(r_{a,0}=(p-a)/p\) and \(r_{a,1}=1\). Thus the directed
coefficient is the raw coefficient times
\(r_a(Tp)/r_a(p)\), giving \(2^a\) on each component. Its
chronological \(b=6/a\) product gives 64.

The two point divisors are retained globally by
\[
 \frac{\eta-z\tau_{a,\mathrm{old}}(z^4)}\eta
 =\frac{p-\tau_{a,\mathrm{old}}(z^4)}p
 =1-v\tau_{a,\mathrm{old}}(z^4),
 \tag{R0P.11}
\]
on the respective graph charts. Here
\(\tau_{2,\mathrm{old}}=2/(1+8t)\) and
\(\tau_{6,\mathrm{old}}=-x_6/y_6\), with respective limits
2 and 6. The equal exceptional orders cancel in this explicit
ratio, and their distinct strict directions remain. The reciprocal
ratio supplies the inverse module morphism. The other graph centre
and the tangency chart carry no point divisor. The two raw Hom
coefficients tend to \(2^a\) at the tangency; their vertical
orders on both components are exactly \(a/2\). Subtracting the
whole-fibre divisor removes those orders. The remaining horizontal
Hom divisor is canceled by R0P.11. A rational section and its
reciprocal with zero divisor in these smooth local rings are units.
This proves the entire finite extension, including every point
where a displayed meromorphic coefficient alone has a zero or pole.

\par

\medskip\noindent\textbf{R0P.3. The full period module and the primitive elliptic quotient}

\par\nobreak

The completed sphere topology is carried by the actual K96 disc
family and the F35/F46 integral source, common blowup and ruling
maps. Their primitive normal circles and the full root/reference
maps are the specified ones in R0E.5,8. They give
\(H^2(X;\mathbb Z)=H^4(X;\mathbb Z)=0\) for each prescribed
fan. A horizontal polar divisor would have positive area against
a flat Kähler form on a regular compact two-torus, while its
restricted divisor class is zero. Thus every polar divisor is
vertical. Removal of codimension-two indeterminacy and compactness
make the meromorphic function constant on each regular fibre.
At a special component, invariance under the branches of
\(t-t_0=u^e\) forces its Laurent exponents to be divisible by
\(e\). This proves the precise equality
\(\pi^*\mathbb C(t)=\mathbb C(X)\). An endpoint and its inverse
therefore give an actual Möbius map on the original base.

Keep the genuine point-theta coefficient and all four periods
\[
 V_{\rm pt}=f_{P^a}(Z)v_M,\qquad
 \gamma=(-bg,\ell),\quad e_1=(1,0),\quad
 e_2=(\tau,ag),\quad\delta=(0,1).
 \tag{R0P.12}
\]
The full \(\ell\) is the logarithm of R0E.37, including
\(\lambda q/t^3\), every Miller factor and the translated theta
ratio. The tensor comparison is the complete constant R0E.38,
with both reciprocal Cartier module arrows.

The point-representative change \(g'=g+\mu+\nu\tau\) multiplies
the line coefficient by
\((-Z/P^a)^{a\nu}Q^{-a\nu(a\nu+1)/2}\). The identity
\(\Theta(Q^{-N}x)=(-x)^NQ^{-N(N+1)/2}\Theta(x)\), with
\(N=a\nu\), proves this exact formula. Its inverse linear
matrix is \(\left(\begin{smallmatrix}1&0\\-a\nu&1\end{smallmatrix}\right)\).
Applying it to all three target periods gives R0E.40, including
\(c=\ell'-\ell+ab\nu(2g+\mu+\nu\tau)\). The modular
change R0E.41 sends \(p e_1+r e_2\) and \(q_0e_1+s e_2\)
to exactly the new elliptic periods, because \(ps-q_0r=1\).
Its theta quotient fixes the entire additive constant. These
maps give the full monodromy
\[
 M_\alpha=
 \begin{pmatrix}1&0&0\\ bv_\alpha&A_\alpha&0\\
 c_\alpha&av_\alpha^tJA_\alpha&1\end{pmatrix},
 \qquad J=\begin{pmatrix}0&-1\\1&0\end{pmatrix},
 \tag{R0P.13}
\]
with its complete central product R0E.44.

The Abel class is nonzero on this source. Clearing a rational
coboundary and multiplying by six puts a multiple of \(G\) in
the identity components and gives a single-valued Lie logarithm.
At the finite III point, \(\tau(2G')=2z/(1+8z^4)\), and
\(\omega_t=z^{-1}\omega_{\rm root}\). Dividing the root
logarithm by \(z\) gives the extending original Lie coefficient.
The quartic elliptic action has no invariant period. At the
selected I1 the invariant annular period extends in \(xy=q\).
At infinity the good cubic logarithm of \(3G'\) has leading
term \(-z^2/8\), and division by \(z^2\) gives the original
Lie frame; its order-three action again has no invariant period.
The full transition is \(u_0(q_\infty)=q_\infty^{-1}u_t(q_\infty^{-1})\).
A global holomorphic coefficient with that transition is zero.
The original infinity value of \(3G\) is \((64,512)\), whose
additive parameter is \(-1/8\); all nonzero integer multiples
remain nonzero. This proves the required contradiction and
non-torsion through the actual finite and infinity Lie charts.

The vector equation \((A_\alpha-I)u+bzv_\alpha=0\) and its
dual, using
\(w^tJ(A-I)=(w-Aw)^tJA\), now give exactly
\(\mathbb Z\delta\) and \(\mathbb Z\psi\) as the invariant
integral subgroups, with \(\psi(\gamma)=1\) and zero on the
other three generators. These are primitive in the displayed
lattice. An endpoint carries them to the corresponding pulled-back
groups before any base rigidity is used.

Its periodic derivative on the universal cover is constant, and
complex linearity gives the exact affine map
\((z,w)\mapsto(Az+z_0,Cz+\epsilon w+w_0)\), with
\(A\ne0\), \(\epsilon=\pm1\). Its inverse is R0E.46.
The full lattice map on \(\ker\psi\) is
\(\left(\begin{smallmatrix}p&q_0&0\\r&s&0\\u&v&\epsilon\end{smallmatrix}\right)\).
The complex area identity makes its elliptic determinant one.
Consequently its first coordinate and inverse descend to actual
inverse elliptic source isomorphisms, with the line constants and
integers retained in the second coordinate. Directly from the
original cubic,
\[
 j_0(t)=6912/(27t+4).
 \tag{R0P.14}
\]
Equality of this invariant under the proved elliptic map forces
the Möbius map to be \(h(t)=t\).

\par

\medskip\noindent\textbf{R0P.4. Global source sections and the entire automorphism group}

\par\nobreak

The elliptic monodromy centralizer is \(\mathbb QI\): its
order-three infinity matrix generates \(\mathbb Q(\zeta_3)\),
and the selected unipotent has a nonzero nilpotent part. A
nonscalar element of that quadratic field generates it, so the
simultaneous commutant is scalar. Applying the two Abel
coboundary equations to R0P.13 then gives exactly
\(eI+c\delta\psi\), with all central cochain entries retained.
Integral unimodularity gives \(e=\pm1\); complex linearity
eliminates the real rank-one remainder. On a single completed
object the negative sign gives sixteen transverse fixed points
on every regular two-torus. Their closed analytic equalizer has
a horizontal compact curve component. Its positive degree over
the base contradicts its pairing with the zero class
\(\pi^*H^2(\mathbb P^1)\). The automorphism has positive
linear part.

The actual ordinary nonzero-line cover has a nonzero line vector
over the \(O\)-section of each whole ordinary source disc, using
R0P.11 at point coincidences. This kills the base meridian.
The remaining vanishing cycles lie in \(\ker\psi\), so the
endpoint lift and its inverse extend there. At the selected
source, the integral orders in a chart \(q=z_0z_1z_2\) sum to
one. A literal \(h^rB^k\) reindex removes them, retaining all
four full units of R0E.50. The remaining map on each character
is multiplication by \(K(q)^uC(q)^v\), with reciprocal inverse,
as in R0E.52.

The normalized infinity fibre product has its actual root lift
because \(A_\lambda^m=B_\lambda^n\), \(n=\pm1\). The
commutator with the cap generator is \(B_\lambda^k\).
Its \(m\)-fold iteration is \(B_\lambda^{mk}=1\), which
strict contraction makes \(k=0\). The integer joins on the
one-punctured base with its single order-\(m\) orbifold point
have zero homomorphism to \(\mathbb Z\). Their actual deck
powers therefore join these source projections to a global
source section.

The rational surface R0P.2 has fibres IV*, III and I1, Euler
number 12 and rank-ten unimodular divisor lattice. The actual
III direction of \(G\) is \([x:U:V]=[0:0:1]\), off the
strict fibre plane \(V-U-x=0\); it meets the nonidentity
component once. At infinity the two section arms are the
distinct root points \((0,1)\) and \(O\), with the two
quotient rays specified in R0E.53. Thus the E6 and III
corrections are \(4/3\) and \(1/2\), and \(G\cdot O=0\).
The full ten-class matrix R0E.53 has determinant \(-1\).
These actual classes span the whole lattice; quotienting by
the fibre, zero section and seven components gives the free
section group \(\mathbb ZG\). In particular the source
section is \(jG\), with the actual primitive generator.

The full endpoint lift is \(H(t)U_G^j\), where
\(U_G(p,v)=(p+G,B_{a,1}(p+G)^{-1}v)\). Its finite orders
are \(-3/b\) on both III components by R0P.9-\/-11. The
infinity order is \(2m/b\) in its own root coordinate, and
the selected integer is the retained \(B^k\) reindex. The
three divisor orders of \(H\) are therefore
\(3j/b,-2j/b,-k\). Their sum gives \(j=bk\), and the
literal map is
\[
 H=c(t^3/q)^k,\qquad H U_G^{bk}=c B_{\rm sel}^{-k}.
 \tag{R0P.15}
\]
On the complete quotient this equals the physical scalar
\(c\lambda^k\). The reciprocal maps extend through the same
source, root, fan and reference charts. An identity scalar
deck would translate the source by \(-kbG\); non-torsion
forces \(k=0\). This proves precisely
\(\operatorname{Aut}_{\rm hol}(X)=\mathbb C^*_{\rm physical}\),
including its faithful action and connected scalar paths.

\par

\medskip\noindent\textbf{R0P.5. Both complete infinity-root and line morphisms}

\par\nobreak

The infinity coordinate is \(q_\infty=z^m\), so its two
original-coordinate maps are
\[
 \begin{array}{c|c|c}
 m&(t,x,y)&\text{root equation}\\\hline
 3l+2&(z^{-m},z^{-2(l+1)}X,z^{-3(l+1)}Y)
   &Y^2=X^3+z^{l+2}X+z^2\\
 3l&(z^{-3l},z^{-2l}X,z^{-3l}Y)
   &Y^2=X^3+z^lX+1.
 \end{array}
 \tag{R0P.16}
\]
These share the stated original \(t\) on the ordinary overlap
with R0P.3. Each root is used on its own original infinity
disc. The selected factor remains
\(h=1+4z^m/27=z^mq\).

For the reduced root put \(k=2l+1\) and
\(c=(1+z^lX)^{1/2}\), with its given unit branch. The
coordinates \(U=Y-zc,V=Y+zc\) have \(UV=X^3\), and
\(2zc=V-U\) has its actual analytic inverse for \(z\).
The three charts R0E.60 are regular and retain this source
map. In the middle chart \(X=ab,U=a^2b,V=ab^2\), write
\(v=z^l\). The whole R0E.61 chord has slope \(aK\),
\(X'=a(a-b)F\), \(Y'=-z-aKX'\), and
\(V'=-aX'H\). Our independent calculation uses the
relations \(c^2=1+vab\), \(2zc=ab(b-a)\) and
\((c')^2=1+vX'\) to prove these identities before taking
any central value. The coordinate map
\(a'=(b-a)F/H\), \(b'=-aH\) and the inverse chord
are therefore the actual two source morphisms, with
\(c,c+1,c'+1,K,F,H\) unit on these charts.

The vertical field in R0E.62 has
\(W(X)=2Y\) and \(W(Y)=3X^2+z^{l+2}\) on the fixed
\(z\) slice; these identities also hold in the independent
universal calculation. It is the extending invariant group
field. Since \(2m-1=3(2l+1)\), the small formal inverse
in R0E.63 has derivative \(1/m\) and supplies the stated
extending division translation with \(mR_{\rm div}=G'\).
The complete deck weights make every displayed unit
invariant. Its line graph is \((\tau',z^k)\), with all
rays \((j,1)\), \(1\le j\le k\). Adjacent determinants
are one, and both section equations have order \(j\) on
the corresponding ray. The two actual module ratios are
\(u/(u-L)\) and \(1/(1-vL)\), with \(L(0)=-1/4\).
Their inverses give the point-line identification carried by
the K96 common graph and commuting source/root maps.
The exact multiplier is \(h z^{-2}B'_6\); the full chord
gives \(B'_6/z^2=4\) at the centre, with its complete
horizontal modules retained.

For the good root set \(u=z^l\). The two marked charts
have quotient weights \((2,1)\), invariant equation
\(AB=C^3\), and rays \((2/3,1/3),(1/3,2/3)\).
The literal generator ratio gives
\(\beta^*N_2=\pi^*M_2+E_{2G'}-E_O\), with
\(E_P(v_w,v_z)=\min(v_w,2lv_z)\). The small translation
has \(w'=w\varepsilon(w,u)+u^2a(w,u)\), preserving
\((w,u^j)\) and \((u,w^j)\), \(j=1,2\), with their
inverse arrows. Thus both complete corrections are carried.
The whole pullback is
\(f_{2,\mathrm{old}}=z^{-l}\phi\), with the actual
\(\phi\) of R0E.69. The three central translated factors
are \((Y-1)/X,-(Y+1)/X,2/X\); their product is \(-2\),
and their squared norm is four. The selected multiplier
is the full \(hB'_{2,3}\).

For both rows the chord coefficient
\(g_{P_0,R}=\ell_{P_0,R}/(x-x(P_0+R))\), with
\(P_0=aG'\), has exactly the required horizontal Hom
divisor. Multiplication by the stated \(\tau(R)\) or
\(\tau(P_0)\) gives central coefficient one in its
actual module frames. Its two deck weights cancel. The
ratio \(\beta=B'_{\rm sel}/F^m\) is an invariant
extending base unit with central value four; its prescribed
analytic root retains the reference. Hence
\[
 \widetilde F=\beta^{1/m}F,\quad
 A_R=R^n\zeta_m^d\widetilde F^{\,n}D,\quad R^m=\lambda,
 \qquad A_R^m=B_R^n,\quad [A_R,B_R]=1
 \tag{R0P.17}
\]
are full source and line identities for both signs and every
finite character.

\par

\medskip\noindent\textbf{R0P.6. Full annular reference and the correct primitive inverses}

\par\nobreak

Use the actual root map \(I_{z_\alpha}\), transported
periods and reference integers from R0E.72-\/-74. With
\(s_\beta=s_\alpha+k\) and
\(f_\beta=D_*^kf_\alpha+\Pi_{\alpha\beta}\), the
entire residual is
\[
 h_\beta-D_*^kh_\alpha
 =-n s_\alpha\Pi_{\alpha\beta}
   +j_{\alpha\beta}\operatorname{Log}_\beta B_R
   +\Lambda_{\alpha\beta}.
 \tag{R0P.18}
\]
The written partition and Cauchy solution retain their
prescribed reference values and all integral columns.
Subtracting the target-sheet exponent
\(knf_\beta+kc_R\delta\) from the full exponent gives
exactly \(j_{\alpha\beta}\operatorname{Log}_\beta B_R+
\Lambda_{\alpha\beta}\). Thus
\(\Phi_\beta=A_R^kB_R^{j_{\alpha\beta}}\Phi_\alpha\),
with the displayed reciprocal inverse at its actual source
point. This verifies the directed annular square, including
the source-dependent period corrections.

The complete scalar family is R0E.76. At integral winding
\(J\) it has \(B_{\theta+J}=B_\theta\) and
\(\Phi_{\theta+J}=z^{nJ}\Phi_\theta\). The additional
sector scalar is \(e^{2\pi iJ J_\alpha}=1\) for its
retained integer, so no sector column is erased. The full
ordinary intertwiner and the independent period vectors
\(e_1,\delta\) force equality of the original scalar
values. The winding is therefore integral.

The three actual divisor orders for a positive endpoint are
R0E.79. Their degree gives \(j=bk+bnJ/m\), and hence
\(m\mid bJ\). For the good root the actual central
coefficient of \(B'_{2,1}\) is \(-X(p)\). Its inverse
is evaluated at \(p+G'\), giving precisely
\[
 U'_G(p,v)=\bigl(p+G',B'_{2,1}(p+G')^{-1}v\bigr),
 \qquad B'_{2,1}(p+G')^{-1}\big|_0
 =-1/X(p+G')=(Y+1)/(2X).
 \tag{R0P.19}
\]
The equality follows from the actual chord
\(X(p+G')=2(1-Y)/X^2=-2X/(Y+1)\) on
\(Y^2=X^3+1\). Our independent execution checks both
source inverses, and the three chronological plus-source
composites give the identity source and coefficient \(1/4\).
The complete root functions and their Cartier modules,
rather than these central values alone, supply the extension.

The literal ordinary/root square is
\(U_{G,\mathrm{old}}=z^{2l}U'_G\), with both graph
corrections from R0P.16. The entire deck relation is
\(U'_G D=\zeta_m^{-2l}D U'_G\). The normalized norm
lift commutes with it in the full constant-factor point
frame, retaining the base constants fixed by their
specified functions. Thus the primitive good forward maps
are
\[
 (R^{-l}tU_G)^n,\qquad (R^{-l}U'_G)^n,
 \tag{R0P.20}
\]
with \(-2ln\equiv ln\pmod {3l}\) and the exact
annular factor \(z^{nl}t^nz^{2ln}=1\).
For \(n=-1\) every reciprocal is evaluated at the
actual inverse source point.

Three chronological roots \(R,\zeta_3R,\zeta_3^2R\)
give scalar product \(\lambda^{-n}\). With
\(S=tU_G\), \(S^3=qB_{\rm sel}^{-1}\), their
complete endpoint product is
\[
 q^nB_\lambda^{-n}\quad\hbox{and}\quad
 h^nB_\lambda^{-n}.
 \tag{R0P.21}
\]
For M6 the primitive winding is \(m\), with the
canonical \(q^n/h^n\) quotient pair. Its selected
zero is carried by the actual \(B^n\) reindex and
the unit character map \((P^{bn},\eta^{-n})\) in
R0E.50, with reciprocal inverse on every cone.
This constructs every positive-part winding allowed
by the divisor equation, including all finite characters.

\par

\medskip\noindent\textbf{R0P.7. Negative lifts and the complete winding groups}

\par\nobreak

The selected inversion \((Z,V,q)\mapsto(Z^{-1},V^{-1},q)\)
reverses the two full periods. Negation sends a periodic
basis triangle to its other triangle with translated
vertex \(-p-u-v\). The full representative \(hBI\)
has the two chart maps R0E.90, with
\(\mathcal A=\kappa g\), \(\mathcal C=\alpha\eta\).
Their products preserve \(q\), their dense maps are
\(Z'=\mathcal A q/Z\), \(V'=\mathcal C q/V\), and
the two coordinate substitutions compose to the literal
identity. All coefficient functions remain evaluated at
the unchanged source height. Integral character changes
carry these arrows to every prescribed basis fan.

Composing a negative endpoint with that extending
inversion gives the positive selected unit map and
its inverse. It therefore supplies a single-valued
negative lift on the selected punctured disc.
On every other ordinary disc the actual point/O
frame gives the nonzero \(O\)-section vector killing
the meridian, and the vanishing cycles are in
\(\ker\psi\). The constant suspension joins give
one global ordinary lift. The normalized infinity
fibre product and its point/O Cartier frame give
the corresponding single-valued root lift fixing
\(z\). Both retain
\(\widetilde F B=B^{-1}\widetilde F\).

One original infinity circle has the full reference
word \(AB^\nu\), with the actual integer \(\nu\).
A root circle has word \(B^N\),
\(N=n+m\nu\). Both endpoints retain the same
integer by R0P.18 and the scalar restoration; the
extra \(z^{nJ}\) is single-valued on that root
circle. Equating the two lifts on a cut annulus,
with its one constant deck, and continuing around
the root gives
\[
 B^N\widetilde F B^{-N}=B^{2N}\widetilde F.
 \tag{R0P.22}
\]
Single-valuedness and strict contraction give
\(N=0\), while \(N\equiv\pm1\pmod m\) and
\(m\ge2\). This excludes every negative endpoint
with its actual annular integer, including the
quadratic reduced cap.

The complete winding groups are therefore
\(m\mathbb Z\) for \((6,1;3l+2)\) and
\(l\mathbb Z\) for \((2,3;3l)\), with the
primitive maps R0P.20-\/-21 and their chronological
inverses. At a fixed winding the full endpoint
ambiguity is the physical scalar group proved in
R0P.15. The physical return uses the inverse
endpoint. This proof assigns no numerical return.

\par

\medskip\noindent\textbf{R0P.8. Independent execution and audit scope}

\par\nobreak

The independent local executable is
\path{code/part3_fourfold/e6_iii_local_point_peer.py}.
Its 67 checks derive the full quartic source and
inverse, all four Miller functions, their component
limits and genuine-frame ratios, the two graph
ideal equations, tangency map and whole-fibre
intersection kernel. The independent global
executable is
\path{code/part3_fourfold/e6_iii_global_endpoint_peer.py}.
Its 1663 further checks retain the complete good
root source inverses and their threefold product,
the reduced chord in its full \(c,c'\) equations,
the invariant vertical field, every root exponent,
point graph, signed character, period coordinate,
target-sheet reference equation and endpoint word.
No author or root calculator is imported.

The full geometric read of R0E.1-\/-11 and R0R.1-\/-4
passes. The analytic extension, compact topology,
primitive invariant quotient, actual source section
and single-valued negative lifts are established
above using the retained explicit maps. The finite
integer executions supplement the universal
identities and those geometric proofs. The final
record locks all source and certificate bytes;
the two initial comparison-code stages are
preserved with their exact defining-equation
corrections. They changed no mathematical source.

% End expanded verified cp36 input

\par\bigskip\noindent\textbf{Independent proof of the full root, signed endpoint and annular comparisons}\par\nobreak

% Expanded verified cp36 input: notes/part3_e6_iii_endpoint_checks_peer.tex
\par

\medskip\noindent\textbf{R0CP.1. Exact scope and source dependency}

\par\nobreak

This is the independent checks audit of R0E.8-\/-10 at source SHA256
\texttt{21a7e21c\allowbreak 39ceedef\allowbreak 3abf4c14\allowbreak 8362ecd7\allowbreak a480d063\allowbreak fcc0c1c4\allowbreak 25fe6c2a\allowbreak 95c0c9ce}.
The original source is
\[
 y^2=x^3+tx+t^2,\qquad G=(0,t),\qquad
 t_c=-4/27,\qquad q=t+4/27.
\]
The genuine point line is \(M_a=\mathcal O([aG]-[O])\), and the
suspension translates by \(-bG\). Both entire cap progressions
\((a,b;m)=(6,1;3l+2)\), \(l\geq0\), and
\((2,3;3l)\), \(l\geq1\), retain \(n=\pm1\), the prescribed finite
character, the original selected-height fan and the full K96 reference.

The finite R0E.1-\/-4 source and Cartier maps have their separate
49-check audit. For the present read I also read all R0E.5-\/-7:
the whole point-theta period module, its full suspension coefficient,
the nonzero Abel class in the explicitly identified Lie line, both
primitive invariant lattice groups, the affine source map and inverse,
the rational source invariant, the centralizer, the actual cover lifts,
and the integral source-section lattice. The present numerical checks
do not replace those geometric arguments.

The stated negative and positive endpoint conclusions use those actual
cover lifts. At an ordinary III disc, a nonzero vector over the full
\(O\)-section disc kills the base meridian; the remaining generic-fibre
loops are subject to the actual vanishing-cycle relations in
\(\ker\psi\). The same argument uses the original point/O Cartier
modules on the normalized root disc. Thus the later annular comparison
uses single-valued lifts on both actual covers.

\par

\medskip\noindent\textbf{R0CP.2. Both original-coordinate root and line maps}

\par\nobreak

In the reduced row, the literal substitutions are
\[
 t=z^{-(3l+2)},\quad x=z^{-2(l+1)}X,\quad y=z^{-3(l+1)}Y,
 \qquad Y^2=X^3+z^{l+2}X+z^2.
\]
With \(c=(1+z^lX)^{1/2}\), the coordinates
\(U=Y-zc,V=Y+zc\) satisfy \(UV=X^3\). The three displayed A2
charts, the middle-chart chord and its inverse preserve the same
\(z\). The identity formal inverse in
\(R_{\rm div}=2G'-[m]^{-1}((2m-1)G')\) has derivative \(1/m\),
and \(mR_{\rm div}=G'\). The deck weights retain the invariant
\(c\) and the complete chord units.

The point/O graph has \(\tau'=z^{2l+1}\tau_0\) and
\(\tau'(6G')=z^{2l+1}L\), \(L(0)=-1/4\). Its actual coefficient
ratios are \(u/(u-L)\) and \((1-vL)^{-1}\) on the two charts.
The two sections have equal exceptional orders on each preceding
ray and distinct strict directions at the equality ray. These
are the directed Cartier maps in
\(\beta^*N_6=\pi^*M_6\). The complete selected Hom coefficient is
\(h z^{-2}B'_6\), with \(h=1+4z^m/27\); its central coefficient
after that displayed factor is four.

In the good row the full map is
\[
 t=z^{-3l},\quad x=z^{-2l}X,\quad y=z^{-3l}Y,
 \qquad Y^2=X^3+z^lX+1.
\]
The quotient and both section directions retain the monomial
correction
\(E_P(v_w,v_z)=\min(v_w,2l v_z)\), so the comparison is
\(\beta^*N_2=\pi^*M_2+E_{2G'}-E_O\).
The small formal translation has the actual form
\(w\mapsto w\varepsilon+u^2a\), \(u=z^l\), with inverse;
it preserves all four ideals \((w,u^j),(u,w^j)\), \(j=1,2\).
The Miller factor remains
\[
 f_{2,\mathrm{old}}=z^{-l}\phi,\qquad
 \phi=\frac{Y-1-z^lX/2}{X-z^{2l}/4}.
\]
Its three central translates have product \(-2\); the full squared
norm is four. The displayed quotient monomials remain in the
ordinary/root comparison.

In each row the normalized chord-over-vertical map has exactly the
Hom divisor for translation by \(-R_{\rm div}\). Multiplication by
the stated formal \(\tau\)-factor cancels its pole and gives central
coefficient one. Its opposite deck weights cancel before taking the
reference-fixed root of \(B'_{\rm sel}/F^m\). Consequently the full
maps satisfy
\[
 \widetilde F^m=B'_{\rm sel},\qquad
 A_R=R^n\zeta_m^d\widetilde F^{\,n}D,\qquad
 A_R^m=B_\lambda^n,
\]
with every base unit and finite character present. The finite III
quartic charts and these infinity charts are compared only through
their displayed original \(t,x,y\) maps on the ordinary overlap.

\par

\medskip\noindent\textbf{R0CP.3. The actual inverse coefficient and its chronological norm}

\par\nobreak

On the central good cubic \(Y^2=X^3+1\), the chord for translation
by \(+G'=(0,1)\) gives
\[
 X_1=-\frac{2X}{Y+1},\qquad Y_1=\frac{Y-3}{Y+1}.
 \tag{R0CP.1}
\]
The full forward \(-G'\) Hom coefficient has central expression
\(B'_{2,1}(p)=-X(p)\) in the canonical meromorphic point frame.
Its inverse is therefore
\[
 U'_G(p,v)=
 \left(p+G',\frac{Y+1}{2X}v\right),
 \qquad
 \frac{Y+1}{2X}=-\frac1{X(p+G')}.
 \tag{R0CP.2}
\]
The original point/O modules carry the zeros and poles of this
meromorphic expression and its reciprocal. Their Hom maps have
zero vertical order and extend on the actual root line.

The next source point is
\(X_2=2X/(Y-1),Y_2=-(Y+3)/(Y-1)\), and the third is the original
point. The three coefficients of the inverse arrow, evaluated at
their chronological source points, are
\[
 \frac{Y+1}{2X},\qquad -\frac{Y-1}{2X},\qquad-\frac1X.
\]
Their product is \((Y^2-1)/(4X^3)=1/4\). This gives a pointwise
check of the inverse direction and its central norm. The earlier
wording that identified this inverse coefficient with \(-X(p)\)
was located in the audit and preserved by the author before repair.
The complete ordinary/root square is still
\(U_{G,\mathrm{old}}=z^{2l}U'_G\), with its actual correction
monomials; the corrected source-point evaluation is now literal.

\par

\medskip\noindent\textbf{R0CP.4. Full annular reference, scalar family and signed endpoints}

\par\nobreak

The sector construction retains the complete source displacement,
physical scalar and root change in
\[
 \Phi_\alpha=\operatorname{Exp}
 (ns_\alpha f_\alpha+s_\alpha c_\lambda\delta+h_\alpha)I_{z_\alpha}.
\]
Its cochain equation carries both the full period \(\Pi\) and every
integer \(j,\Lambda\). The reference-zero Cauchy solve is performed
on that residual. Subtraction of the actual target-sheet logarithm
then gives \(\Phi_\beta=A_\lambda^kB_\lambda^j\Phi_\alpha\).
The inverse uses the actual inverse source point and reciprocal
Cartier maps.

Keeping the integer representatives \(J_\alpha\), the scalar family
is exactly
\[
 \Phi_{\theta,\alpha}=e^{2\pi in\theta s_\alpha/m}
 B_\theta^{J_\alpha}B_*^{-J_\alpha}\Phi_{*,\alpha}.
\]
The overlap exponent is \(n\theta k/m+\theta j\). At integer
winding \(J\), the extra exponent is
\(nJs_\alpha/m+JJ_\alpha\); the second term is an integer after
it has been carried through the square. Thus the full residual is
\(z^{nJ}\).

For a positive endpoint \(H(t)U_G^j\), the actual finite, selected
and root divisors give
\[
 \operatorname{ord}_0H=3j/b,\quad
 \operatorname{ord}_{t_c}H=-k,\quad
 m\operatorname{ord}_\infty H+2mj/b+nJ=0.
\]
Their degree gives \(j=bk+bnJ/m\) and \(m\mid bJ\).
This is the integral restriction on the full rational base
coefficient, with the selected reindex \(B^k\) retained.

For M2, \(S=tU_G\) has the entire selected coefficient of R0E.81,
an extending unit and reciprocal at unchanged \(q\), on every
prescribed cone. On the good root
\(U'_GD=\zeta_m^{-2l}DU'_G\). The genuine point-extension
translations commute with the full norm lift; their base constants
remain in the constant-factor point frame. Hence, for either sign,
\[
 (U'_G)^nA_R=A_{\zeta_3R}(U'_G)^n.
\]
The primitive forward maps are
\((R^{-l}S)^n\) and \((R^{-l}U'_G)^n\). Their annular square
cancels precisely
\(z^{nl}t^nz^{2ln}=1\), using the original \(t=z^{-3l}\).
The three successive root basepoints contribute the complete scalar
\(\lambda^{-n}\), so the chronological products are
\(q^nB_\lambda^{-n}\) and \(h^nB_\lambda^{-n}\).

For M6 the primitive forward \(q^n/h^n\) closing uses
\(z^{nm}q^n=h^n\). At the selected zero its literal
\(B_\lambda^n\) reindex leaves the source/line unit pair
\((P^{bn},\eta^{-n})\); each character and reciprocal extends
at unchanged height. Negative winding and negative primitive norm
use the corresponding chronological inverse maps.

\par

\medskip\noindent\textbf{R0CP.5. Complete negative continuation and audit conclusion}

\par\nobreak

The selected inversion \((Z,V,q)\mapsto(Z^{-1},V^{-1},q)\)
conjugates both full periods to their inverses. The two explicit
\(hBI\) vertex maps preserve \(q\), interchange the charts and
are literal inverses, with \(\mathcal A=\kappa g\) and
\(\mathcal C=\alpha\eta\) retained. Composing a negative lift
with this map gives the positive integral-height unit construction;
composing back supplies its single-valued selected lift.

On the normalized root the preserved kernel is exactly the
\(B_\lambda\) cyclic cover, since \(n=\pm1\). The actual Cartier
frames over the \(O\)-section give its single-valued root lift.
The entire original-base turn is \(A_\lambda B_\lambda^\nu\),
and its full root turn is \(B_\lambda^N\),
\(N=n+m\nu\). The integer \(\nu\) remains unrestricted.

Continuing the directed comparison
\(\widetilde F_{\rm cap}=\Phi_t\widetilde F_{\rm ord}\Phi_s^{-1}\)
around that root circle gives
\(B_\lambda^N\widetilde F_{\rm cap}B_\lambda^{-N}
=B_\lambda^{2N}\widetilde F_{\rm cap}\).
Both lifts are single valued and the contraction has infinite
order, so this requires \(N=0\), whereas
\(n+m\nu\) cannot vanish for \(n=\pm1,m\geq2\).
Every analytic coefficient is already in the displayed operators;
this contradiction does not remove a reference integer.

The complete R0E.8-\/-10 geometric read passes after the preserved
inverse-coefficient repair. The independent certificate checks the
actual central source/inverse, both vertex inverses, full root and
finite-character exponents, arbitrary sector representatives,
the three-divisor equation and the entire negative annular word.
Its finite checks accompany the explicit extension, covering and
cochain arguments above. This audit assigns no numerical return.

% End expanded verified cp36 input

\end{proof}


\subsection{Complete scalar returns on the fixed E6 source locus}
\begin{theorem}[Source result P3K.115 under continuation review]\label{night:P3K.115}
\StatementRecord{night:P3K.115}{open}{Historical events: GPT-6 Astra (OpenAI), ChatGPT Pro mode, effort Pro, S6\_06\_BLIND, 16--18 November 2026: targeted proof check. Also checked in this writing. Continuation mathematical and source checks: GPT-6 Astra (OpenAI), ChatGPT Pro mode, effort Pro, S6\_10\_PAPERS continuation, 19 November 2026: computation or formal certificate; targeted proof check; used and rechecked in this run. Continuation metadata events: GPT-6 Astra (OpenAI), ChatGPT Pro mode, effort Pro, S6\_10\_PAPERS continuation, 19 November 2026: used and rechecked in this run; source-row intake; metadata reconciliation only; no new mathematical check. GPT-6 Astra (OpenAI), ChatGPT Pro mode, effort Pro, S6\_10\_PAPERS continuation, 19 November 2026: used and rechecked in this run; explicitly scoped metadata reconciliation; metadata reconciliation only; no new mathematical check. GPT-6 Astra (OpenAI), ChatGPT Pro mode, effort Pro, S6\_10\_PAPERS continuation, 19 November 2026: used and rechecked in this run; receiving source placement; metadata reconciliation only; no new mathematical check. \emph{Pre-continuation recorded proof scope.} Historical event: GPT-6 Astra (OpenAI), ChatGPT Pro mode, effort Pro, S6\_06\_BLIND, 16--18 November 2026: targeted proof check, recorded producing-run re-derivation. Also checked in this writing by GPT-6 Astra (OpenAI), ChatGPT Pro mode, effort Pro, S6\_10\_PAPERS continuation, 19 November 2026: computation or formal certificate for the explicitly listed algebraic substeps in the event ledger; full continuation rederivation and renewed credit comparison remain pending.}{\cref{night:P3K.67,night:P3K.73,night:P3K.96}; the original source, point line, roots, characters, reference, selected fan and scalar range stated here}
Retain the original marked source
\[
 E_r:\quad y^2=x^3+rx^2+tx+t^2,\qquad G=(0,t),
 \qquad q=t-t_c,
\]
on the complete marked-node domain
\[
 \omega\in\mathbb C\setminus\{0,1,1/2,-1/2\},\qquad
 r=(1-\omega^2)/3,\qquad
 t_c=(\omega-1)^2(2\omega+1)/27.
\]
Prescribe the ordering of the selected simple-node branches and
any lattice-basis fan \(P\). Use the genuine point line
\(M_a=\mathcal O([aG]-[O])\), its full chronological translation
by \(-bG\), and either infinity-cap progression
\[
 (a,b;m)=(6,1;3l+2),\quad l\ge0,
 \qquad (a,b;m)=(2,3;3l),\quad l\ge1.
\]
Keep either norm \(n=\pm1\), every prescribed finite character,
the complete P3K.96 annular reference, and a nonzero scalar in its
actual strict contraction disc. The finite fibre at \(t=0\)
has its inherited order-one, norm-zero filling; at \(\omega=-1\)
this is the original type-III source filling. Put \(\lambda=R^m\).

Every scalar path with a holomorphic endpoint has equal original
scalar values at its ends and winding \(J=kJ_0\), \(k\in\mathbb Z\).
Its primitive positive forward endpoint and its return class are
\[
 \begin{array}{c|c|c|c}
 &J_0&E_R^{\rm ord}&E_R^{\rm cap}\\\hline
 M6&m&q^n&h^n\\
 M2&l&(R^{-l}tU_G)^n&(R^{-l}U'_G)^n,
 \end{array}
 \qquad h=1-t_cz^m=z^mq,
\]
where \(U_G\) and \(U'_G\) are the complete point-line lifts
at their actual source arguments. Every endpoint is the chronological
\(k\)-th power of this endpoint, composed with one physical scalar.
Its oriented physical sphere return is
\[
                         [\Phi]=kS(P)\pmod {28}.
\]
Here \(S(P)\) is the exact P3K.73 algorithm reproduced in
ES.29--31 below, with the actual inputs
\[
 \begin{array}{c|c|c|c}
 &J&\widehat H&\widehat\chi\\\hline
 M6&3l+2&6(3l+2)&l+1\\
 M2&l&2l/3&l/3.
 \end{array}
\]
For \(T_j=T(e_1+je_2,e_2)\) the primitive characteristic
pair and return integer are
\[
 \begin{array}{c|c|c|c}
 &k_{M,j}&\ell_{M,j}&S(T_j)\\\hline
 M6&18m-3mj(j-1)&24(l+1)&8l+5-mj(j-1)/2\\
 M2&2l-3lj(j-1)&8l&-lj(j-1)/2.
 \end{array}
\]
The native left entries are respectively \(6(l+1)\) and \(2l\).
The exact M2 primitive-family comparison with P3K.67 is
\(F_R(v_M)=R^{nl}v_M\), with inverse \(R^{-nl}\) and
\(F_{\zeta_3R}E_R=\widetilde D_RF_R\). Its threefold pullback
is the actual canonical \(q^n/h^n\) family on every prescribed
fan and cap.

At a fixed fan the full image over all scalar paths and holomorphic
endpoints is the cyclic subgroup generated by \(S(P)\).
For fixed \(l\), over all fans and scalar loops, the M2 image is
exactly \(l\mathbb Z/28\), realized by covers of \(T_2\), whose
primitive integer is \(-l\). For every allowed M6 index the image
is all \(\mathbb Z/28\): use \(T_0\) for \(l\not\equiv2\pmod7\)
and \(T_4\), of integer \(-10l-7\), otherwise.

The forward seam \((z,1)\sim(\Phi(z),0)\) gives the oriented
mapping torus \((S^1\times S^6)\#\Sigma_{kS(P)}\), with invariant
\(kS(P)/28\) for both spin structures. The reverse Jin seam gives
the opposite sphere summand and invariant. Every nonzero class
therefore has a mapping torus not diffeomorphic to
\(S^1\times S^6\) in either orientation. All conclusions retain
the original source, point line, roots, characters, reference,
height, selected branches, and prescribed fan.
\end{theorem}
\begin{proof}[Supplied source argument; continuation review in progress]

% Expanded verified cp36 input: notes/part3_e6_complete_scalar_returns_proof.tex
\par

\medskip\noindent\textbf{Completed proof and retained submission history.}

\par

The original introductory submission notice is retained below. The full ES.1-\/-6 mathematical proof has now passed the complete independent comparison reproduced afterward, together with its source-locked exact certificate.

This proof joins the complete endpoint classifications of EB and R0E to
the actual scalar-family and fan maps of P3K.67-\/-73. Its new comparison
is a holomorphic family isomorphism with a displayed inverse. The
full proof and its executable checks are submitted for independent
audit before the reserved statement P3K.115 is exported.

\par

\medskip\noindent\textbf{ES.1. The original source and the complete marked-node domain}

\par\nobreak

We retain the original equation, section, and base coordinate
\[
 E_r:\quad y^2=x^3+rx^2+tx+t^2,\qquad G=(0,t),\qquad
 s=\frac14-r .
 \tag{ES.1}
\]
The selected fibre has original height \(q=t-t_c\). Its two local
branches and its lattice-basis fan are prescribed. The coordinate
\(\omega\) below marks the selected fibre; the prescribed ordering
of its local branches is additional retained data. The exact
marked-node domain used in this statement is
\[
 \mathcal W=\mathbb C\setminus
 \left\{0,1,\frac12,-\frac12\right\},\qquad
 r=\frac{1-\omega^2}{3},\qquad
 t_c=\frac{(\omega-1)^2(2\omega+1)}{27}.
 \tag{ES.2}
\]
Throughout this domain we use the genuine point line
\(M_a=\mathcal O([aG]-[O])\), its full chronological translation
by \(-bG\), and either of the two infinity-cap progressions
\[
 \begin{array}{c|c|c|c}
 \text{row}&(a,b)&m&l\\ \hline
 M6&(6,1)&3l+2&l\geq0\\
 M2&(2,3)&3l&l\geq1 .
 \end{array}
 \tag{ES.3}
\]
Both primitive norms \(n=1,-1\), every prescribed finite character,
and the complete K96 annular reference are retained. At the finite
source fibre \(t=0\) we use its inherited ordinary filling, with
effective order one and norm zero. For \(\omega=-1\) this is the
original type-III filling of R0E. For the other points of
\(\mathcal W\) it is the original type-\(I_2\) filling of EF.
The scalar lies in an actual punctured disc of strict contraction.
The scalar-root coordinate will be \(R\), with \(\lambda=R^m\);
it is separate from both the source parameter \(\omega\) and the
source-base coordinate \(t\).

We verify the domain directly. The entire discriminant is
\[
 \Delta_r(t)=-16t^2F_r(t),\qquad
 F_r(t)=27t^2+(4-18r)t+r^2(4r-1),
 \qquad \operatorname{disc}_tF_r=16(1-3r)^3 .
 \tag{ES.4}
\]
On \(F_r(t_c)=0\), \(r\ne1/3\), the inverse to ES.2 is
\[
 \omega=\frac{54t_c-18r+4}{4(1-3r)}.
 \tag{ES.5}
\]
Substitution gives both identity composites and
\[
 F_r(t_c)=0,\qquad
 \partial_tF_r(t_c)=4\omega^3,\qquad
 s=\frac{4\omega^2-1}{12}.
 \tag{ES.6}
\]
Thus the selected root is simple and nonzero throughout ES.2.
The fibre at \(t=0\) is \(I_2\), except at
\(\omega=-1\), where \(r=0,t_c=-4/27\) and R0E supplies the
actual type-III source and point-line charts.

The excluded values have explicit original-source meanings.
At \(\omega=0\), \(r=1/3,t_c=1/27\), and the cubic at \(t_c\)
is \((x+1/9)^3\); the selected fibre is the type-II collision.
At \(\omega=1\), \(r=t_c=0\), and the selected height is the
type-III fibre itself. These heights do not give a selected
simple \(I_1\) whole zero. At \(\omega=\pm1/2\), \(r=1/4\).
The elliptic addition law gives
\[
 2G=(s,-t-s/2)=(0,-t)=-G,\qquad 3G=O.
 \tag{ES.7}
\]
Here \(G\) has exact order three on every regular generic fibre.
For completeness the original rational-surface arrow is
\[
 u=\frac{y-t-x/2}{x},\qquad
 v=\frac{y+t+x/2}{x},
\]
\[
 x=uv+s,\qquad
 t=\frac{x(v-u-1)}2,\qquad y=\frac{x(u+v)}2.
 \tag{ES.8}
\]
These rational maps are mutual inverses on their indicated dense
open sets. The undivided equation is
\((y-t-x/2)(y+t+x/2)=x^2(x-s)\). At \(s=0\) its finite
singularity is \(UV=x^3\); its two exceptional curves give the
finite \(I_3\) fibre. ES.4 gives the other finite \(I_1\), at
\(t=1/54\), and the unchanged infinity resolution is \(IV^*\).
The rational minimal elliptic surface therefore has Euler number
twelve and unimodular rank-ten divisor lattice. Its fibre, zero
section, \(E_6\) components and \(A_2\) components give a full-rank
sublattice of determinant \(-9\). Its index is three. The section
quotient consequently has rank zero and order three, and ES.7
generates it. This is the exact loss of the required
Mordell-\/-Weil group \(\mathbb ZG\) at both of these parameters.

For \(r\ne0,1/4,1/3\), EB proves the required group and all
completed endpoint statements through its intrinsic original
elliptic quotient. At \(r=0\), R0E proves them with its original
type-III source. Hence those proofs apply at every point of ES.2.
The formula ES.5 is valid also at \(\omega=-1\); its numerator
and denominator are respectively \(-4\) and \(4\).
No limiting identification of the two fibres is used.

\par

\medskip\noindent\textbf{ES.2. Full scalar endpoints and the exact family comparison}

\par\nobreak

Write \(U_G\) for the rational lift of translation by \(+G\) on
the actual point line, with its coefficient evaluated at the
translated source point. The entire chosen suspension is
\[
 B_0=B_{\mathrm{sel}}=\frac{q}{t^3}U_G^{-b},
 \qquad B_\lambda=\lambda B_0,\qquad
 \rho=B_0^{-1}/U_G^b=\frac{t^3}{q}.
 \tag{ES.9}
\]
The full infinity-root equation and line give
\[
 h=1-t_cz^m=z^mq,\qquad q_\infty=z^m,\qquad
 A_R^m=B_\lambda^n .
 \tag{ES.10}
\]
All these equalities are equalities of maps of the original
Cartier modules, including their chronological source arguments.
EF, SEF and R0E give the extending ordinary maps, the complete
root maps and their reciprocals at every divisor.

The complete EB/R0E endpoint classification states that every
holomorphic scalar endpoint has the same scalar at its two
ends and integral scalar winding. Its winding group and actual
primitive positive forward arrows are
\[
 \begin{array}{c|c|c|c}
 \text{row}&J_0& E_R^{\mathrm{ord}}&E_R^{\mathrm{cap}}\\ \hline
 M6&m&q^n&h^n\\
 M2&l&(R^{-l}tU_G)^n&(R^{-l}U'_G)^n .
 \end{array}
 \tag{ES.11}
\]
Every endpoint of winding \(kJ_0\) is the chronological \(k\)-th
power of this arrow, composed with a physical scalar.
For negative \(k\) the power uses the actual inverse maps at
their inverse source points. The ratio of two endpoints at
the same winding is exactly an element of
\(\mathbb C^*_{\mathrm{physical}}\).

We recall the geometric contents of this classification because
they are needed for exhaustiveness here. EB.2/R0E.5 identify the
primitive invariant line \(\mathbb Z\delta\) and dual line
\(\mathbb Z\psi\) in the full period module. The actual cyclic
cover \(\ker\psi\) is \(M_a^\times\) with its full point-theta
map. An endpoint and its affine inverse descend from this cover
to degree-one isomorphisms of the original elliptic source
fibres. The displayed \(j_r(t)\) in EB.7 fixes the base
pointwise for \(r\ne0\); at \(r=0\) the exact degree-one
function \(6912/(27t+4)\) does so. The full period intertwiner
then forces equality of the original scalars. The original
Cartier divisor equations give \(m\mid bJ\), and the displayed
maps attain its smallest positive solution \(J_0\).
The complete negative endpoint comparison retains the annular
integer \(N=n+m\nu\) and forces \(B_\lambda^{2N}=1\).
Strict contraction forces \(N=0\), whereas \(n=\pm1,m\ge2\)
gives no such integral \(\nu\). This includes the quadratic
reduced cap. No further endpoint component is left.

We now identify ES.11 with the exact K67 family, including its
constant phase. In K67 SE.4 take \(M_0=m\), \(\varepsilon=n\)
and the specified forward ordinary closing \(h_0=q^n\).
ES.9 gives the actual rational root
\[
 g_0^b=h_0\rho^n=t^{3n}.
 \tag{ES.12}
\]
For M2, \(b=3\), so choose \(g_0=t^n\). The resulting K67
arrow is exactly
\[
 D_R=R^{-nm/3}g_0U_G^n=(R^{-l}tU_G)^n=E_R.
 \tag{ES.13}
\]
Its cap arrow is the same \((R^{-l}U'_G)^n\), by the
ordinary/root coefficient identity \(U_G=z^{2l}U'_G\),
the source identity \(t=z^{-3l}\), and the full reference
change \(z^{nl}\). Their product is one:
\[
 z^{nl}t^nz^{2ln}=1 .
 \tag{ES.14}
\]
The root character square is
\((U'_G)^nA_R=\zeta_m^{ln}A_R(U'_G)^n\), with the
prescribed character in \(A_R\) retained. This is exactly
the square for \(R\mapsto\zeta_3R\).

K67 SE.8 uses the corrected forward endpoint
\[
 \widetilde D_R=cD_R,\qquad
 c=\exp(2\pi i nl/3).
 \tag{ES.15}
\]
There is an actual holomorphic family isomorphism between these
two endpoint presentations:
\[
 F_R(p,v_M)=(p,R^{nl}v_M),\qquad
 F_R^{-1}(p,v_M)=(p,R^{-nl}v_M).
 \tag{ES.16}
\]
Both are defined for every nonzero \(R\) in the family. They
commute with every original source/line suspension, root deck
action, finite character and full annular transition, because
their scalar is a nonzero base function and every such
transition is linear in its line vector. They satisfy the
directed endpoint identity
\[
 F_{\zeta_3R}E_R=\widetilde D_RF_R ,
 \tag{ES.17}
\]
since \(F_{\zeta_3R}=cF_R\).

Here is the extension of ES.16 on every selected fan. In the
original cusp characters it is \((Z,V,q)\mapsto
(Z,R^{nl}V,q)\). In the characters of a basis \(P\), the two
multipliers are
\[
 A=(R^{nl})^{(P^{-1})_{12}},\qquad
 D=(R^{nl})^{(P^{-1})_{22}}.
 \tag{ES.18}
\]
Their two vertex maps are
\[
 (a,b_1,c_1)\longmapsto((AD)^{-1}a,Ab_1,Dc_1),
 \qquad
 (u,v_1,w_1)\longmapsto(ADu,D^{-1}v_1,A^{-1}w_1).
 \tag{ES.19}
\]
Each preserves the full product \(q\). Substitution in the
complete period transitions of K73 SA.2 preserves all their
units. On each affine cone a character is multiplied by its
nonzero constant character value; this defines an invertible
map of the entire character-semigroup algebra. The inverse
uses the reciprocal multipliers. On each native star its Cox
lift is K73.7 with these same \(A,D\), so it carries both normal
collars as well. Thus ES.17 descends to a biholomorphic
isomorphism of the complete families, including every chosen
fan and the finite caps. Its fibre maps preserve complex
orientation.

The primitive descended parameter is \(W=R^3\), with
\(\lambda=W^l\). The uncorrected threefold phase in K67 is
\(\exp(-2\pi i n m(3-1)/(2\cdot3))=\exp(-2\pi i nl)=1\).
ES.15 also has \(c^3=1\). More fully, SEF.20/R0E.88 give
\[
 E_{\zeta_3^2R}^{\mathrm{ord}}
 E_{\zeta_3R}^{\mathrm{ord}}E_R^{\mathrm{ord}}
 =q^nB_\lambda^{-n},
 \qquad
 E_{\zeta_3^2R}^{\mathrm{cap}}
 E_{\zeta_3R}^{\mathrm{cap}}E_R^{\mathrm{cap}}
 =h^nB_\lambda^{-n}.
 \tag{ES.20}
\]
The quotient maps are \(q^n,h^n\). These identities retain
the deck reindex and give the actual degree-three pullback of
the shorter family.

For M6, \(b=1\) and the K67 phase is one. ES.12 gives
\(g_0=t^{3n}\), and its full ordinary arrow is
\[
 D_R=R^{-nm}t^{3n}U_G^n=q^nB_\lambda^{-n}.
 \tag{ES.21}
\]
On the quotient this is precisely \(q^n\), with cap \(h^n\)
by ES.10 and the full annular square. The primitive parameter
is \(W=R\), with \(\lambda=W^m\). Hence both rows of ES.11
have now been identified with the exact endpoints of K73
through actual full-family morphisms and their inverses.

\par

\medskip\noindent\textbf{ES.3. The actual page, physical line and positive reference return}

\par\nobreak

Choose the regular binding at an actual base value \(\tau\)
away from the source singular values. The retained K32 page
map for the canonical positive \(R\)-loop is
\[
 \eta-\tau\xi^m=1,\qquad
 t=\tau+\xi^{-m}.
 \tag{ES.22}
\]
Eliminating \(\eta\) compactifies its base to
\(\mathbb P^1_\xi\), and the displayed meromorphic map has
degree \(m\). Its branch points are \(\xi=0,\infty\).
On a source chart the page arrow is the actual sheet-zero
map \(p\mapsto[p,\gamma_{\mathrm{graph}}(t(p))]\).
Two graph lifts differ by the specified power of the original
source translation and line suspension. At the infinity root
the equality \(A_R^m=B_\lambda^n\) removes the residual
quotient exactly. At \(t=0\) its order-one, norm-zero map is
the original finite source and its Cartier line. The entire
root and point/O graph arrows in EF.28-\/-38 and R0E.59-\/-74
carry these sheet maps across every resolved component.
K32\textquotesingle s relative translation-torsor map, fixed on the binding
collar, carries the page to this actual source pullback.
It preserves the full physical line and tangent boundary
trivializations.

The source height is \(h(G)=1/6\) throughout ES.2. Its component
calculation is \(2-4/3-1/2=1/6\): the \(IV^*\) correction
is \(4/3\), and the finite \(I_2\) or III correction is \(1/2\).
The full ten-class primitive-source computations in EF and
R0E prove that this is the marked generator, including at
\(\omega=-1\). Pullback by ES.22 and its actual minimal root
resolution gives
\[
 H=\frac{ma^2}{6},\qquad
 c_1(M_{\mathrm{page}})_{\mathrm{rel}}^2=-H .
 \tag{ES.23}
\]
These are the retained K32/F2 relative intersection equalities
for the sheet-zero point line, including its exceptional
divisor corrections.

The total source Euler number away from infinity is four.
At \(\omega=-1\) it is \(3+1\); at the other parameters it
is \(2+1+1\). Each finite fibre has \(m\) unramified lifts
under ES.22. The minimal infinity root is IV for M6 and is
smooth for M2. Its Euler number is respectively four or zero.
The compact page therefore has
\[
 12\chi=4m+e'_\infty,\qquad
 \begin{array}{c|c|c|c}
 \text{row}&H&\chi&p_{\mathrm{rel}}\\ \hline
 M6&6m&l+1&-24(l+1)\\
 M2&2l&l&-24l .
 \end{array}
 \tag{ES.24}
\]
The finite III replacement has thus been counted on its actual
unramified sheets. Its contribution to the formula remains
four together with the selected \(I_1\).

The orientation is the original complex orientation on \(S^6\),
the parameter circle is positive, and
\(\mu(\Sigma_j)=j/28\) in \(\Theta_7=\mathbb Z/28\).
K68\textquotesingle s actual Hermitian dual map carries the physical line to
the meridian-one Thom line and reverses its circle character.
Its cap and inverse-clutch maps give
\[
 (k_M,\ell_M)=(3H,24\chi),\qquad
 A_M=\frac{4k_M-\ell_M}{24}=\frac H2-\chi
 \tag{ES.25}
\]
on the canonical positive \(R\)-loop. The norm sign \(n\)
is already retained in its signed angular page map; K32.18
and K69 give ES.25 for both signs. The direction of the
scalar loop in ES.25 is unchanged.

For M6 the primitive loop is this canonical loop. For M2,
ES.16-\/-20 identify the primitive loop with the K67/K73
shorter family. K73 SA.1 gives an oriented degree-three
cover of its actual completed orbit products. The physical
line, the full tangent bundle and their boundary
trivializations pull back under that cover. Integration of
their integral classes gives
\[
 3k_M^{\mathrm{prim}}=3H,\qquad
 3\ell_M^{\mathrm{prim}}=24\chi .
 \tag{ES.26}
\]
These are equalities of free integral characteristic numbers.
Their unique integer solutions, and the corresponding native
pair from K73.9, are
\[
 \begin{array}{c|c|c|c}
 \text{row}&(k_M,\ell_M)_{\mathrm{prim}}
 &(\mu_{\mathrm{nat}},\nu_{\mathrm{nat}})_{\mathrm{prim}}
 &A_{\mathrm{prim}}\\ \hline
 M6&(18m,24(l+1))&(6(l+1),18m-6(l+1))&8l+5\\
 M2&(2l,8l)&(2l,0)&0 .
 \end{array}
 \tag{ES.27}
\]
The equality ES.26 arises from the full-family cover. In
particular the zero in the last row is fixed by the actual
characteristic pair of the primitive loop.

\par

\medskip\noindent\textbf{ES.4. Every prescribed lattice-basis fan}

\par\nobreak

For a prescribed \(P\in GL_2(\mathbb Z)\), apply the complete
K73 algorithm with the following actual family data:
\[
 \begin{array}{c|c|c|c}
 \text{row}&J&\widehat H&\widehat\chi\\ \hline
 M6&3l+2&6(3l+2)&l+1\\
 M2&l&2l/3&l/3 .
 \end{array}
 \tag{ES.28}
\]
Denote its single-loop integer path sum by \(S(P)\).
ES.16-\/-21 carry the full endpoints and all source, root,
normal and character maps to the K73 family; ES.22-\/-26
identify its page and cover data. Thus every hypothesis of
that algorithm is satisfied by these specific completed
families.

For clarity we give its exact execution rule. The identity
on the absolute torus, extended on matching character
semigroups, identifies the six bases of the same fan in
K71 AF.1. Swapping columns first handles determinant minus
one. The terminating integer subtraction of AF.2 gives a
vertical seed
\(T_j=\left(\begin{smallmatrix}1&0\\j&1\end{smallmatrix}\right)\)
and its finite word in
\[
 Q_0=\begin{pmatrix}1&-1\\0&1\end{pmatrix},
 \qquad Q_1=\begin{pmatrix}1&0\\-1&1\end{pmatrix}.
\]
At that seed put
\[
 V=\begin{pmatrix}
 2J-6\widehat H+6Jj(j-1)&2J(1-2j)&2J(2j-1)\\
 2J(1-2j)&2J&0\\
 2J(2j-1)&0&2J
 \end{pmatrix},
\]
\[
 \delta=\bigl(
 2J+6\widehat H-24\widehat\chi-6Jj(j-1),
 2J(2j-1),2J(1-2j)\bigr)^t,\quad
 S=\widehat H/2-\widehat\chi-Jj(j-1)/2.
 \tag{ES.29}
\]
At a word step \(P\mapsto P'=PQ_i\), \(i=0,1\), take
\(a_0=V_{ii}\), \(b_0=\delta_i\), and compute
\[
 B=(b_{jk})=2J(P')^{-1}e_2e_2^tP',\qquad
 \ell=(b_{22}-b_{21},b_{11}-b_{12},b_{22}),
\]
\[
 (x_0,y_0)^t=(P')^{-1}e_2,\quad
 w=(x_0,y_0,-x_0-y_0),\quad
 d=(-x_0-2y_0,-2x_0-y_0,x_0-y_0),
\]
\[
 I_j=Jw_j+d_j(\ell_j-J),\qquad
 W_j=Jw_j-d_j(\ell_j-J).
 \tag{ES.30}
\]
All target weights are nonzero. Set
\[
 C=w_2^2a_0-w_2I_2,\qquad
 E=w_2(2a_0-b_0)-W_2,
\]
\[
 V'_{jk}=\frac{I_j}{w_j}\delta_{jk}
             +\frac C{w_jw_k},\qquad
 \delta'_j=\frac{W_j+E}{w_j},\qquad
 S'=S+\frac{b_0-a_0}{4}.
 \tag{ES.31}
\]
K73 proves integrality through the actual doubled ordinary
product and the single native right integer
\((b_0-a_0)/4\). Its Cox maps preserve \(q\), carry the
entire bounded unit and both collars, and retain the original
finite caps. ES.18-\/-19 commute with these maps. The final
absolute-fan identity therefore gives the original prescribed
family on \(P\). Its primitive positive physical return is
\[
                         [\Phi_P]=S(P)\pmod {28}.
 \tag{ES.32}
\]

In particular the vertical fan formula, with all characteristic
numbers taken on the actual primitive loop, is
\[
 \begin{array}{c|c|c|c}
 \text{row}&k_{M,j}&\ell_{M,j}&S(T_j)\\ \hline
 M6&18m-3mj(j-1)&24(l+1)&8l+5-\dfrac m2j(j-1)\\
 M2&2l-3lj(j-1)&8l&-\dfrac l2j(j-1).
 \end{array}
 \tag{ES.33}
\]
The native pairs are
\((6(l+1),18m-6(l+1)-3mj(j-1))\) and
\((2l,-3lj(j-1))\), respectively. Their sums and four
times their first entries give the two characteristic
columns in ES.33. This also verifies their integral spin
congruence directly.

\par

\medskip\noindent\textbf{ES.5. Exhausting scalar paths and endpoint choices}

\par\nobreak

We keep the source, cap, character, reference and fan data
fixed. Consider any closed scalar path with a holomorphic
endpoint identification in the stated contraction disc.
The endpoint classification in ES.11 first forces its two
original scalar values to agree. Write its continuous
logarithmic lift as \(L(\theta)\), \(0\le\theta\le1\).
Then
\[
 L(1)-L(0)=2\pi iJ,\qquad J=kJ_0,\quad k\in\mathbb Z.
 \tag{ES.34}
\]
Use parametrizations constant on small collars of the
endpoints. Choose a smooth nondecreasing
\(\beta:[0,1]\to[0,1]\), constant at zero and one on
these collars. The explicit based homotopy is
\[
 \begin{gathered}
 L_s(\theta)=(1-s)L(\theta)
       +s\bigl(L(0)+2\pi iJ\beta(\theta)\bigr),\\
 \lambda_s(\theta)=e^{L_s(\theta)},\qquad
 R_s(\theta)=R(0)
       \exp\bigl((L_s(\theta)-L(0))/m\bigr).
 \end{gathered}\tag{ES.35}
\]
Its endpoints and the root endpoint ratio
\(\exp(2\pi iJ/m)\) are fixed. The real part in ES.35
is a convex combination of the original real part
and \(\operatorname{Re}L(0)\); it stays in the same
strict-contraction disc. Both source-parameter endpoints
and every retained finite character therefore stay fixed.

The full family over this rectangle is obtained by
substituting ES.35 in the original operators
\(B_\lambda=\lambda B_0\), \(A_R=R^nA_0\), and the
SEF.3/R0E.76 sector cochains
\[
 \Phi_{\theta,\alpha}
 =\exp(2\pi in\theta s_\alpha/m)\,
 B_\theta^{J_\alpha}
       (B_*^{-J_\alpha}\Phi_{*,\alpha}).
 \tag{ES.36}
\]
Here the complex logarithmic parameter \(\theta\) in ES.36
is replaced by \((L_s-L(0))/(2\pi i)\), with the initial
reference restored at \(L(0)\). Every integer \(J_\alpha\)
and every source period column remains unchanged. On
sector overlaps the same entire equalities SEF.4 hold
with these substituted functions. Their inverses use
the reciprocal factors. At the endpoints ES.34 makes
the reindex factors \(\exp(2\pi iJJ_\alpha)=1\).
This constructs the proper smooth family of completed
fibres over the full rectangle, including its selected
fan, caps and fixed endpoint collars.

To transfer the return, take smooth local lifts of the
\(s\)-tangent field in these total-space charts and
transport the lifts on overlaps by the actual gluing
differentials. The vertical correction of two such
lifts is a tangent field of the compact fibre. A
partition of unity on the glued compact shells patches
them to a global lift. On the fixed parameter collars
choose the product lift first and patch relative to
those collars. Its flow exists for \(0\le s\le1\)
by compactness. The reverse field gives its inverse.
The endpoint gluing is fixed, so this flow gives a
smooth based isotopy of the return maps, conjugated
by the initial-fibre identification. These maps
preserve the complex orientation of the fibres.
Consequently the return of the original path equals
the return of the \(k\)-fold primitive circle.

There is also no unrecorded contribution from choosing
a different holomorphic endpoint. Write that endpoint
as \(cE^{(k)}\), \(c\in\mathbb C^*\), using ES.11.
Choose one complex number \(\log c\). The actual global
maps \(v_M\mapsto e^{s\log c}v_M\), \(0\le s\le1\),
give an isotopy of the target-fibre automorphism.
They extend through every cap and selected chart by
the same ES.18-\/-19 character maps. Their inverses
have reciprocal factors. Applying these maps on a
parameter collar changes the endpoint continuously
from \(E^{(k)}\) to \(cE^{(k)}\); the product near each
end supplies the smooth collar matching. This
changes the return by an orientation-preserving
isotopy.

Concatenating the primitive interval \(k\) times
concatenates its actual endpoint arrows with their
successive parameter values. The inverse intervals
use their inverse arrows. These give respectively
\(\Phi_P^k\) for positive \(k\), the identity for
\(k=0\), and the corresponding inverse powers for
negative \(k\). The complete classification at the
fixed data and fixed fan is therefore
\[
 \boxed{\quad J=kJ_0,\qquad [\Phi]=kS(P)\pmod {28},
       \qquad k\in\mathbb Z .\quad}
 \tag{ES.37}
\]
Every integer \(k\) is realized. On the actual concentric
annulus \(|W_*|e^{-\epsilon}<|W|<|W_*|e^\epsilon\),
with \(W_*\ne0\) and sufficiently small \(\epsilon>0\), pull back
the constructed primitive holomorphic family by
\(W\mapsto W_*(W/W_*)^k\). Its image stays in the
given contraction domain. For \(k=0\) this is the
constant family; for negative \(k\) it is a holomorphic
annular map with the reversed circle winding.
These are actual family pullbacks with the endpoint
maps just described. At a fixed fan the precise image
is the cyclic subgroup generated by \(S(P)\), of order
\(28/\gcd(28,S(P))\).

\par

\medskip\noindent\textbf{ES.6. Exact images over all fans in each cap progression}

\par\nobreak

For M2, the seed ES.29 and every subsequent quantity in
ES.30-\/-31 is linear in the three inputs \(J,\widehat H,
\widehat\chi\). The fan word depends only on \(P\).
Their inputs in ES.28 are \(l\) times the actual
\(l=1\) inputs. Induction along that same finite word
gives the exact integer equality
\[
 S_{M2,l}(P)=lS_{M2,1}(P).
 \tag{ES.38}
\]
Every scalar return at this fixed \(l\), on any fan,
is consequently a multiple of \(l\) modulo 28.
On the actual fan \(T_2\), ES.33 gives
\(S_{M2,l}(T_2)=-l\). Its positive parameter covers
realize every multiple of \(l\). The exact image
over all fans and scalar loops is thus
\[
             l\mathbb Z/28\mathbb Z,
 \qquad \text{of order }28/\gcd(28,l).
 \tag{ES.39}
\]
In particular the \(l=1\) family at \(T_2\) has
primitive class \(-1=27\), and its covers realize
all twenty-eight classes.

For M6, the reference fan \(T_0\) has integer
\(8l+5\), which is always odd. Its residue modulo
seven is \(l+5\). For each \(l\not\equiv2\pmod7\)
it is therefore relatively prime to 28. For each
\(l\equiv2\pmod7\), use the actual fan \(T_4\).
Its integer is
\[
 (8l+5)-6(3l+2)=-10l-7.
 \tag{ES.40}
\]
It is odd and has residue one modulo seven.
It too is relatively prime to 28. Thus at every
allowed \(l\) there is a specified fan whose
primitive scalar return generates \(\Theta_7\).
Positive parameter covers of that one family
realize all twenty-eight classes. The same
statement holds at every source in ES.2, for
both norm signs and all prescribed finite
characters and full references, because every
preceding comparison retained those data.

For each nonzero class in ES.37, define the forward
oriented mapping torus by the seam
\(T^+_\Phi:(z,1)\sim(\Phi(z),0)\).
P3F.24\textquotesingle s actual disc-and-collar insertion identifies this
torus with
\((S^1\times S^6)\#\Sigma_{kS(P)}\).
The physical inverse-clutch map of K68 is the
one used in that identification. Both spin
structures have Eells-\/-Kuiper invariant
\(kS(P)/28\). The reverse Jin seam defines \(T^-_\Phi\);
its corresponding sphere summand and both invariants
are \(-\Sigma_{kS(P)}\) and \(-kS(P)/28\).
On \(S^1\times S^6\) both
extend across \(S^1\times D^7\) and have
invariant zero. Reversing orientation changes
the displayed nonzero invariant by its negative.
Hence these mapping tori are not diffeomorphic
to \(S^1\times S^6\), in either orientation.

This result computes all scalar loops with
holomorphic endpoints for the complete fixed-source
domain and two cap progressions ES.2-\/-3. The
finite III site throughout has its specified
order-one, norm-zero filling. Other finite
III root progressions and loops moving the
source parameter retain their separately
stated proof obligations.

% End expanded verified cp36 input

\par\bigskip\noindent\textbf{Complete independent comparison proof.}\par\nobreak

% Expanded verified cp36 input: notes/part3_e6_complete_scalar_returns_blind_peer.tex
This is the independent geometric audit of ES.1-\/-6 in
\path{part3_e6_complete_scalar_returns_working.md}. It retains the actual
forward endpoint convention, the full physical line, both signs of the
cap norm, every finite character, and the original annular reference.
The exact certificate is \path{code/part3/e6_complete_scalar_returns_blind_peer.py}.
It imports no target-state reconstruction from the author certificate.
The source and dependency hashes are recorded before and after execution.

\par

\medskip\noindent\textbf{ESP.1. Original source, marked domain, and integral height}

\par\nobreak

Put
\[
 E_r:\quad y^2=x^3+rx^2+tx+t^2,\qquad G=(0,t),
 \qquad s=\frac14-r.
 \tag{ESP.1}
\]
The complete discriminant and the selected-root coordinate are
\[
 \Delta=-16t^2F_r(t),\quad
 F_r(t)=27t^2+(4-18r)t+r^2(4r-1),
\]
\[
 r=\frac{1-\omega^2}{3},\qquad
 t_c=\frac{(\omega-1)^2(2\omega+1)}{27},\qquad
 \omega=\frac{54t_c-18r+4}{4(1-3r)}.
 \tag{ESP.2}
\]
Substitution of the first two functions in the last one is the
identity. Conversely the differences of the reconstructed \(r,t_c\)
from their original values have numerators divisible by \(F_r(t_c)\).
The denominators are nonzero for \(r\ne1/3\). Thus these are the
two directed maps of the marked-node cover. Their derivative identity is
\[
                 F'_r(t_c)=4\omega^3,
 \qquad s=\frac{4\omega^2-1}{12}.
 \tag{ESP.3}
\]
The allowed domain is exactly the finite affine \(\omega\)-line with
\(0,1,1/2,-1/2\) removed. It has a simple, nonzero selected height.
The original fibre at \(t=0\) is \(I_2\), except at \(\omega=-1\),
where \(r=0,t_c=-4/27\) and it is III. The other finite fibres at a
generic allowed point are the selected and unselected \(I_1\) fibres.
At \(\omega=-1\) the finite Euler sum is \(3+1\). The actual R0E
quartic source, point/O, tangency, root, and annular maps apply at that
point; the inverse in ESP.2 has numerator \(-4\) and denominator
\(4\) there.

The excluded points can be checked on ESP.1 itself. At \(\omega=0\)
the selected cubic is \((x+1/9)^3\). At \(\omega=1\) the selected
height is the III fibre at zero. At \(\omega=\pm1/2\), \(r=1/4\)
and the exact doubling formula gives
\[
                   2G=(s,-t-s/2)=-G,
                   \qquad 3G=O.
 \tag{ESP.4}
\]
Its order is precisely three on the generic smooth fibre. The maps
\[
 u=\frac{y-t-x/2}{x},\quad v=\frac{y+t+x/2}{x},
 \qquad x=uv+s,\quad t=\frac{x(v-u-1)}2,
 \quad y=\frac{x(u+v)}2
 \tag{ESP.5}
\]
are inverse rational maps, as follows by their direct composites and
the undivided equation
\((y-t-x/2)(y+t+x/2)=x^2(x-s)\). At \(s=0\) this is the
\(A_2\) equation \(UV=x^3\); its two exceptional components give
the finite \(I_3\). The other root of \(F_{1/4}\) is \(1/54\).
The actual infinity resolution remains \(IV^*\). The trivial
divisor lattice \(U\oplus E_6(-1)\oplus A_2(-1)\) has full
rank ten and determinant \(-9\), inside the unimodular rational
surface lattice. Its index is three. Restriction to the generic fibre
identifies the quotient with the section group, which is therefore
generated by the exact order-three section ESP.4. This proves the
stated loss of the free marked generator at both excluded points.

On the allowed locus the actual EF and R0E component labels use the
positive \(E_6\) matrix
\[
 A=\begin{pmatrix}
 2&-1&0&0&0&0\\-1&2&-1&0&0&0\\
 0&-1&2&-1&0&-1\\0&0&-1&2&-1&0\\
 0&0&0&-1&2&0\\0&0&-1&0&0&2
 \end{pmatrix},\qquad
 A^{-1}e_0=(4/3,5/3,2,4/3,2/3,1)^t.
 \tag{ESP.6}
\]
The root section meets precisely the labelled component \(C_0\).
The original finite section meets the nonidentity \(I_2\), or III,
component once; its positive block is \((2)\). The actual section
charts at finite values and the two distinct root arms at infinity
give \(G\cdot O=0\). Consequently
\[
                          h(G)=2-4/3-1/2=1/6.
 \tag{ESP.7}
\]
In the ordered classes \((O,F,C_0,\ldots,C_5,C_{
m fin},G)\),
use \(O^2=G^2=-1\), \(OF=GF=1\), block
\(-A\oplus(-2)\), and \(GC_0=GC_{
m fin}=1\), with the other
unspecified entries zero. The complete integral Gram determinant is
\(-1\). These ten actual classes fill the source lattice and their
quotient by the first nine gives \(\mathbb ZG\). The calculation
holds on the original III source as well as on the \(I_2\) sources.

\par

\medskip\noindent\textbf{ESP.2. Exhaustive endpoints and the full ordinary/root map}

\par\nobreak

Use the two completed rows
\[
 (a,b;m)=(6,1;3l+2),\ l\ge0,
 \qquad (a,b;m)=(2,3;3l),\ l\ge1.
 \tag{ESP.8}
\]
Their physical point line is \(M_a=\mathcal O([aG]-[O])\).
Let \(U_G\) be its full rational lift of translation by \(+G\).
Every inverse coefficient is evaluated at the inverse source point.
For example, for a coefficient \(c(p)\) of \(U_G\), its inverse
has coefficient \(c(p-G)^{-1}\). All chronological products below
retain these arguments. The suspension and root norm are
\[
 B_0=\frac{q}{t^3}U_G^{-b},\quad B_\lambda=\lambda B_0,
 \quad q=t-t_c,\quad \lambda=R^m,
\]
\[
 q_\infty=z^m,\qquad h=1-t_cz^m=z^mq,
                         \qquad A_R^m=B_\lambda^n.
 \tag{ESP.9}
\]
The finite character in \(A_R\) and the complete annular cochains
are fixed data. The source, point/O modules, root graph corrections,
and their inverses are precisely EF/SEF/K96 on the noncollision locus
and R0E on the original III source.

For completeness, the exhaustiveness proof used here has actual
maps at each step. EB and R0E recover the primitive invariant lines
\(\mathbb Z\delta\) and \(\mathbb Z\psi\) from the whole period
module. The cover \(\ker\psi\) is the actual nonzero point line
with its full theta arrow. An endpoint lifts to that cover; preservation
of the complex span of \(\delta\) makes its first coordinate an affine
degree-one map of the original elliptic source. Its inverse gives the
inverse elliptic map. For \(r\ne0\) the precise invariant
\[
 j_r(t)=\frac{-256(r^2-3t)^3}{t^2F_r(t)}
 \tag{ESP.10}
\]
has its unique double pole at zero, unique simple zero at infinity,
and unique triple zero at \(r^2/3\). Its denominator at the latter
has factor \(r^2(3r-1)^2/3\ne0\). Every permitted base map fixes
these three points and hence is the identity. At \(r=0\) the exact
degree-one function \(6912/(27t+4)\) gives the same conclusion.
This includes the proposed permutations of the two selected-root
choices, after descending through the actual source map.

The full period intertwiner then forces the original scalars at the
two endpoints to agree. The source and Cartier divisor equations
give \(m\mid bJ\), with \(J\) the integral scalar winding. The
complete negative endpoint lift preserves the actual annular integer
\(N=n+m\nu\). Its continued square is
\(B_\lambda^N\widetilde F B_\lambda^{-N}
=B_\lambda^{2N}\widetilde F\). The original ordinary and normalized
root covers make both lifts single valued. Hence \(B_\lambda^{2N}=1\).
Strict contraction makes the generator have infinite order, while
\(n=\pm1\), \(m\ge2\) makes \(n+m\nu=0\) impossible. The
integer \(\nu\) has remained in the equation throughout.

The positive maps attaining the least allowed winding are
\[
 \begin{array}{c|c|c|c}
 &J_0&E_R^{\rm ord}&E_R^{\rm cap}\\\hline
 M6&m&q^n&h^n\\
 M2&l&(R^{-l}tU_G)^n&(R^{-l}U'_G)^n.
 \end{array}
 \tag{ESP.11}
\]
For \(n=-1\) these are the actual inverse maps on their inverse
source points. For M2 the good root has
\(t=z^{-3l}\) and \(U_G=z^{2l}U'_G\), as an equality of
rational coefficient arrows on the complete two-projection Cartier
graph. The scalar-reference shift at winding \(l\) is \(z^{nl}\).
Thus the entire ordinary-to-cap endpoint multiplier is
\[
                         z^{nl}t^nz^{2ln}=1.
 \tag{ESP.12}
\]
The deck commutator of the full root point-line map is
\((U'_G)^n A_R=\zeta_m^{-2ln}A_R(U'_G)^n
=\zeta_m^{ln}A_R(U'_G)^n\). Since
\(A_{\zeta_3R}=\zeta_3^nA_R\), this is exactly the required
directed cap square. Adding any prescribed finite character multiplies
both sides by its same value. The whole source and line periods and
every sector suspension integer remain in the reference square before
ESP.12 is applied.

At a fixed winding the quotient of two endpoints is the actual
holomorphic automorphism group \(\mathbb C^*_{\rm physical}\),
proved by EB and R0E on these completed objects. Therefore the maps
in ESP.11, their chronological powers, and that physical scalar
give every endpoint in the asserted scope.

\par

\medskip\noindent\textbf{ESP.3. The K67 phase, its explicit family inverse, and all fan charts}

\par\nobreak

The ratio in the original line is
\(\rho=B_0^{-1}/U_G^b=t^3/q\). In K67 take
\(M_0=m\), \(\varepsilon=n\), and forward ordinary closing
\(h_0=q^n\). Its rational root satisfies
\[
                         g_0^b=h_0\rho^n=t^{3n}.
 \tag{ESP.13}
\]
For M2 choose \(g_0=t^n\). The uncorrected K67 arrow is precisely
\(D_R=(R^{-l}tU_G)^n\), including its cap ESP.12. Its corrected
forward arrow is \(\widetilde D_R=cD_R\), where
\[
                     c=\exp(2\pi i nl/3).
 \tag{ESP.14}
\]
The actual maps between these two parameter-family presentations are
\[
 F_R(p,v_M)=(p,R^{nl}v_M),\qquad
 F_R^{-1}(p,v_M)=(p,R^{-nl}v_M).
 \tag{ESP.15}
\]
The exponent is an integer for both signs of \(n\); the parameter
is nonzero. Every original line-module transition is linear in its
vector. Hence these maps commute with the full source suspension,
every root deck and finite character, and the entire annular
reference, with all source arguments unchanged. Their directed
endpoint square is
\[
                  F_{\zeta_3R}E_R=\widetilde D_RF_R,
 \qquad F_{\zeta_3R}=cF_R.
 \tag{ESP.16}
\]
This proves the descent of ESP.15 to the primitive parameter
\(W=R^3\). Its inverse satisfies the reciprocal square. The scalar
\(R^{nl}\) is used on the parameter cover with this equivariance;
the descent is exactly the square ESP.16.

On the selected torus ESP.15 is
\((Z,V,q)\mapsto(Z,R^{nl}V,q)\). For a prescribed integral basis
\(P\), its two character multipliers are
\[
 A=(R^{nl})^{(P^{-1})_{12}},\qquad
 D=(R^{nl})^{(P^{-1})_{22}}.
 \tag{ESP.17}
\]
The two full affine vertex arrows and their inverses are
\[
 (a,b,c)\mapsto((AD)^{-1}a,Ab,Dc),\qquad
 (u,v,w)\mapsto(ADu,D^{-1}v,A^{-1}w),
 \tag{ESP.18}
\]
and the reciprocal multipliers, respectively. Each product remains
exactly \(q\). On every character semigroup its action is multiplication
by the specified nonzero character value, which gives the complete
affine-chart isomorphism. The full \(q\)-dependent period coefficients
are evaluated at the unchanged height and commute with this map.
The native Cox lift is
\[
 (a_0,a_1,c_0,c_1)\mapsto
          (a_0,Aa_1,(AD)^{-1}c_0,Dc_1).
 \tag{ESP.19}
\]
Its first native chart has factors \((A,(AD)^{-1},D)\); its second
has \((A^{-1},D^{-1},AD)\). The actual overlap
\((z,n_0,n_1)\mapsto(z^{-1},zn_0,zn_1)\) intertwines them.
The two normal collars, both rulings, and every original finite cap
are carried by these same maps. The integral basis change acts on
the complete logarithmic coefficient matrix by \(L\mapsto Q^{-1}LQ\);
its inverse is \(L\mapsto QLQ^{-1}\). This proves compatibility
with every full period reindex used by the fan algorithm.

The three chronological primitive endpoints use
\(R,\zeta_3R,\zeta_3^2R\). Their extra phase has exponent
\(-ln(0+1+2)=-3ln\), hence is one. They give exactly
\[
 E_{\zeta_3^2R}^{\rm ord}E_{\zeta_3R}^{\rm ord}E_R^{\rm ord}
       =q^nB_\lambda^{-n},\qquad
 E_{\zeta_3^2R}^{\rm cap}E_{\zeta_3R}^{\rm cap}E_R^{\rm cap}
       =h^nB_\lambda^{-n}.
 \tag{ESP.20}
\]
The quotient endpoints are therefore \(q^n,h^n\). K67\textquotesingle s correction
also has \(c^3=1\). The complete three-seam identity can be checked
before exponentiating. With the full source/line functions
\(\Gamma,C,U_0\), put
\[
 \sigma(v)=\left(n\Gamma v,
 v(U_0+nC)-\frac{nm}{2}(v^2-v)\right),
\]
\[
 \gamma(v)=(\Gamma,C-s-mv),\quad H=ns+U_0,
 \qquad d(v)=\frac{H\delta+n\gamma(v)}3+\frac{nm}{9}\delta.
 \tag{ESP.21}
\]
Direct subtraction gives \(\sigma(v+1/3)-\sigma(v)=d(v)\), and
the sum of its three successive values is \(H\delta+n\gamma(v)\).
These are the full K67 cut equations, including all units and the
quadratic term. K73 carries them on ESP.19 to the native product on
the whole parameter circle. Thus ESP.20 is the literal degree-three
pullback of the shorter completed family and product.

For M6, \(b=1\), the K67 phase is one and \(g_0=t^{3n}\).
Its entire ordinary arrow is
\(R^{-nm}t^{3n}U_G^n=q^nB_\lambda^{-n}\), with the same
root equality. The primitive parameter is \(W=R\). All family maps
just used are holomorphic on each fibre and preserve its complex
orientation.

\par

\medskip\noindent\textbf{ESP.4. Page, physical orientation, and the primitive integer pairs}

\par\nobreak

The actual canonical positive \(R\)-loop has the K32 page
\[
                   \eta-\tau\xi^m=1,
                   \qquad t=\tau+\xi^{-m},
 \tag{ESP.22}
\]
where \(\tau\) is a regular binding value. Its compact base is
\(\mathbb P^1_\xi\), and this map has degree \(m\), branching at
\(\xi=0,\infty\), over original infinity and \(\tau\). In a
source chart its sheet is literally
\(p\mapsto[p,\gamma_{\rm graph}(t(p))]\). Changing a graph lift
uses its actual source/line suspension. At infinity the full norm
\(A_R^m=B_\lambda^n\) removes its residual quotient. At zero the
order-one, norm-zero filling is exactly the inherited source and
point line. The full root and Cartier graphs proved in K96, EF,
and R0E carry those maps through every resolved component.

K32\textquotesingle s comparison removes the source translation torsor by the
actual lattice and Lie cochains relative to the binding collar.
For the line it uses the rigidified relative Picard evaluation. The
full cap component degrees are zero, so its remaining vertical
evaluation divisor lies in the kernel of every connected fibre
intersection matrix and is an integral sum of whole fibres. After
its base equations are divided, the Hom section and its reciprocal
extend. The residual base line is retained. Its class is \(kF\),
and \(F^2=c_1(M_{\rm page})F=0\). The same relative comparison
carries the tangent framing. These actual maps give
\[
                 c_1(M_{\rm page})_{
m rel}^2
                       =-ma^2/6=-H.
 \tag{ESP.23}
\]
They also show why a finite character or a different retained base
holonomy does not change this pairing: that line remains the same
pulled-back base line, with the displayed vanishing intersections.

Every finite fibre has \(m\) unramified sheets in ESP.22. The finite
Euler sum is four in both source types, by ESP.1. The actual minimal
root at infinity is IV in the M6 progression and smooth in the M2
progression. Hence
\[
 \begin{array}{c|c|c|c}
 &H&12\chi&p_{
m rel}\\\hline
 M6&6m&4m+4&-24(l+1)\\
 M2&2l&4m&-24l.
 \end{array}
 \tag{ESP.24}
\]
For the III source the same count is the explicit \(m(3+1)\),
with its original tangential fibre retained on each sheet.

The signed norm does not reverse the positive scalar loop. K32.18
gives the literal origin difference
\[
 F_n(v+\theta)-F_n(v)
       =n\theta\gamma-\frac{nm}{2}\theta(\theta-1)\delta,
 \qquad \beta=n\gamma',
 \tag{ESP.25}
\]
with the same full source-dependent origin in the two products.
The binding equation is \(\alpha=\phi-m\beta\); the actual page
and signed normal use this \(\beta\) for each \(n=1,-1\).
The global source/parameter factorization K32.28 retains the integer
graph winding and has the already computed zero normal Pontryagin
transgression. K68\textquotesingle s Hermitian dual and inverse-clutch arrows take
the meridian-one line to the physical line, giving
\[
 (k_M,\ell_M)=(3H,24\chi),\qquad
                   A_M=H/2-\chi.
 \tag{ESP.26}
\]
In detail, its physical Hopf map is
\(\lambda_M(F)=-\Sigma_{F^{-1}}=\Sigma_F\); the dual map is
\(\lambda_U(F)=\Sigma_{F^{-1}}\). The odd line pairings change
sign under the actual dual map. K69 therefore supplies ESP.26 with
the same positive parameter and complex sphere orientation for both
signed page choices in ESP.25.

For M2 the full oriented degree-three covering ESP.20-\/-21 pulls back
the physical line, stable tangent bundle, and their collar frames.
Their free integral pairings satisfy
\(3k_M^{\rm prim}=3H\), \(3\ell_M^{\rm prim}=24\chi\).
Together with the M6 primitive loop this gives
\[
 \begin{array}{c|c|c|c}
 &(k_M,\ell_M)_{\rm prim}&(\mu,\nu)_{\rm nat}&A_{\rm prim}\\\hline
 M6&(18m,24(l+1))&(6(l+1),18m-6(l+1))&8l+5\\
 M2&(2l,8l)&(2l,0)&0.
 \end{array}
 \tag{ESP.27}
\]
The equalities used to obtain the second row are in free integral
characteristic groups on the displayed covering. Its sphere class
is subsequently evaluated by \((4k_M-\ell_M)/24\).

\par

\medskip\noindent\textbf{ESP.5. Full fan propagation and the fresh prism reconstruction}

\par\nobreak

The precise K73 inputs are
\[
 (J,\widehat H,\widehat\chi)
 =\bigl(3l+2,6(3l+2),l+1\bigr)\quad(M6),
 \qquad
 =\bigl(l,2l/3,l/3\bigr)\quad(M2).
 \tag{ESP.28}
\]
The rational entries in the latter row are the actual page data
divided by the degree-three family cover; the characteristic and
native pairs in ESP.27 are integral. The full absolute torus
identity of K71 identifies the six bases of each fan on their
identical character semigroups. A column swap handles determinant
minus one. The stated decreasing integer subtraction then gives
the finite canonical word from a vertical \(T_j\). Every subsequent
target has physical weights \((r,p,-r-p)\), \(p,r>0\).

The two ruling maps at a step are the actual Cox maps
\((n_0/n_1,zn_1,n_1)\) and \((n_1/n_0,zn_0,n_0)\), with
their other affine charts, common rank-one matrix, and prescribed
constant normal interchange. They preserve the entire height and
carry the full coefficient matrix by the integral monomial change
in ESP.19. The source/cap exterior is the original one. Thus the
isomorphism ESP.15 commutes with these complete maps at every step.

On the doubled ordinary product let the selected values be
\(a=V_{ii}\), \(b_0=\delta_i\). K73\textquotesingle s actual single native pair
satisfies
\[
 a=2J-2\mu-2\nu,\qquad
 b_0=2J-2\mu+2\nu,
 \qquad \nu=(b_0-a)/4.
 \tag{ESP.29}
\]
The unrestricted common-link and its normal reset give
\((\mu,\nu,S)\mapsto(\mu+2\nu,-\nu,S+\nu)\).
Consequently the new selected star, labelled 2, has
\(V'_{22}=a\), \(\delta'_2=2a-b_0\).

For the full target weights \(w\), trace vector \(I\), and
derivative vector \(W\), the equations are
\[
 V'(w_ie_i-w_0e_0)=I_ie_i-I_0e_0,\qquad
 w_i\delta'_i-w_0\delta'_0=W_i-W_0\quad(i=1,2).
 \tag{ESP.30}
\]
All three weights are nonzero. Symmetry of \(V'\) makes the
difference of two solutions a scalar multiple of \(rr^t\),
where \(r_i=1/w_i\); the diagonal particular solution is
\(\operatorname{diag}(I_i/w_i)\). The derivative difference
is a scalar multiple of \(r\). The two selected reset values
fix the scalars uniquely. This derives the ES.31 formula from
the complete nine-unknown system. Its existence and integrality
come from the actual native products and Cox maps just retained.

The independent certificate solves that full linear system anew
for both actual input rows, three signed vertical seeds, and each
word of the tested two further depths after the first nonvertical
step. It verifies all equations, both native resets, the closed
formula, full character/inverse maps, and integrality at actual
positive cap indices. It then proves homogeneity symbolically
for the M2 row. The unbounded fan statement follows from the
finite canonical word and the same unique-solution argument
at every step.

In particular the exact vertical values are
\[
 \begin{array}{c|c|c|c}
 &k_{M,j}&\ell_{M,j}&S(T_j)\\\hline
 M6&18m-3mj(j-1)&24(l+1)&8l+5-mj(j-1)/2\\
 M2&2l-3lj(j-1)&8l&-lj(j-1)/2.
 \end{array}
 \tag{ESP.31}
\]
Their first native entry remains \(6(l+1)\) or \(2l\),
and the sum of the two entries is the first characteristic
number. These are the actual primitive loops of ESP.11.

\par

\medskip\noindent\textbf{ESP.6. Every scalar path and every holomorphic closing}

\par\nobreak

Take any allowed closed scalar path and its holomorphic endpoint.
ESP.2 forces equal endpoint scalars and winding \(J=kJ_0\).
For its continuous logarithm \(L(\theta)\), flatten small endpoint
collars and choose a smooth \(\beta\) equal to zero and one there.
The full based homotopy is
\[
 L_s=(1-s)L+s\bigl(L(0)+2\pi iJ\beta\bigr),\quad
 \lambda_s=e^{L_s},\quad
 R_s=R(0)\exp((L_s-L(0))/m).
 \tag{ESP.32}
\]
Its real part lies between the two original real parts and remains
in the same strict contraction disc. The original endpoint values
and the exact root endpoint ratio \(e^{2\pi iJ/m}\) are fixed.

For the actual labelled annular reference write
\(s_\alpha=\log_\alpha q_\infty/(2\pi i)\), with every
suspension integer \(J_\alpha\) retained. The family comparison is
\[
 \Phi_{\vartheta,\alpha}
 =e^{2\pi in\vartheta s_\alpha/m}
 B_\vartheta^{J_\alpha}
           (B_*^{-J_\alpha}\Phi_{*,\alpha}),
 \qquad \vartheta=(L_s-L(0))/(2\pi i).
 \tag{ESP.33}
\]
On an overlap with changes \((s_\beta-s_\alpha,
J_\beta-J_\alpha)=(h,j)\), the extra factors on the two sides
are exactly \(e^{2\pi in\vartheta h/m}\) and
\(e^{2\pi i\vartheta j}\). Thus the full cap/suspension square
holds with the same source period columns, not just with their
central values. Its inverse uses the reciprocal factors and the
same integers. At a closing winding \(J\), the reindex factor
is \(e^{2\pi iJJ_\alpha}=1\), giving the original endpoint square.
Equations ESP.18-\/-19 retain the chosen fan over this whole
rectangle. The source and root pieces are the unchanged compact
pieces; the strict contraction and compact parameter rectangle
give the common compact shell construction.

Local smooth lifts of the homotopy-direction vector field transport
by the differentials of these actual gluing maps. Their differences
are vertical fields on the compact fibre. A partition of unity
patches them, with the product lift fixed first on the endpoint
collars. The resulting field is complete for this compact interval,
and the reverse flow is its inverse. This supplies the based return
isotopy with the endpoint gluing fixed. The fibre identification
preserves the complex orientation. Hence the original path has
the return of the chronological \(k\)-fold primitive loop.

Two holomorphic closings at that winding differ by one global
physical scalar \(c\ne0\). A chosen \(\log c\) gives the actual
isotopy \(v_M\mapsto e^{s\log c}v_M\), with inverse scalar,
on all source, root, and character charts. Applied on a parameter
collar it changes the endpoint without changing the return isotopy
class. Concatenation uses the successive endpoint values in ESP.20;
negative powers use their directed inverse maps. Thus the full
fixed-data, fixed-fan classification is
\[
               J=kJ_0,\qquad [\Phi]=kS(P)\pmod {28}.
 \tag{ESP.34}
\]
Every integer \(k\) occurs on the explicit annulus
\(|W_*|e^{-\epsilon}<|W|<|W_*|e^{\epsilon}\), by the
holomorphic pullback \(W\mapsto W_*(W/W_*)^k\), with
\(\epsilon\) small enough for its image to remain in the
contraction domain. Its core circle winds \(k\) times. The
case \(k=0\) is constant. This also gives the negative winding
maps on their genuine annuli.

\par

\medskip\noindent\textbf{ESP.7. Unbounded images and both mapping-torus orientations}

\par\nobreak

The ES/K73 seed and every target equation are linear in
\(J,\widehat H,\widehat\chi\); the canonical word depends
only on \(P\). Uniqueness of ESP.30 therefore gives, on every
finite canonical word,
\[
                       S_{M2,l}(P)=lS_{M2,1}(P).
 \tag{ESP.35}
\]
All scalar returns for fixed \(l\) lie in \(l\mathbb Z/28\).
The actual vertical fan \(T_2\) has integer \(-l\), so its
positive covers give that whole subgroup, of order
\(28/\gcd(28,l)\).

For M6 the actual \(T_0\) integer is \(8l+5\). It is odd,
and its residue modulo seven is \(l+5\). It is a unit modulo
28 for every \(l\not\equiv2\pmod7\). For
\(l=7h+2\), the actual \(T_4\) integer is
\[
                  -10l-7=7(-10h-4)+1,
 \tag{ESP.36}
\]
which is also odd and a unit modulo 28. One of these two specified
fans therefore has a primitive generator for every allowed \(l\).
Its actual positive parameter covers give all twenty-eight classes.
Every preceding map retained the fixed source, both norm signs,
finite character, and full annular reference, so these image
statements hold at every stated set of those data.

Finally take the forward seam convention
\(T^+_\Phi=(S^6\times[0,1])/((z,1)\sim(\Phi(z),0))\).
P3F.24 first isotopes \(\Phi\) to be supported in an oriented
six-disc. Regluing the rounded seven-ball crossing the circle seam
by \(\Phi\) on its upper face gives exactly this mapping torus;
the same collar insertion gives
\[
                 T^+_\Phi\cong
                    (S^1\times S^6)\#\Sigma_{kS(P)}.
 \tag{ESP.37}
\]
The inverse seam defining the reverse Jin convention replaces
\(\Phi\) by \(\Phi^{-1}\) and the summand by
\(-\Sigma_{kS(P)}\). K68\textquotesingle s physical inverse-clutch identification
fixes this same sphere orientation. Both spin structures extend
on the product summand across \(S^1\times D^7\). Boundary connected
sum with the sphere\textquotesingle s spin bound adds its intersection and
Pontryagin data, so their invariants are respectively
\(kS(P)/28\) on the forward torus and \(-kS(P)/28\) on the
reverse torus. Both values on the product itself are zero.
The nonzero cases consequently exclude a diffeomorphism with
the product for either orientation.

The audit covers all ES.1-\/-6 maps and conclusions in the stated
two cap progressions, with the inherited finite order-one,
norm-zero filling. The independent execution records its finite
map and prism coverage explicitly. The source-domain, endpoint
exhaustiveness, compact-family isotopy, and unbounded image
arguments are the full geometric arguments above.

% End expanded verified cp36 input

\par\medskip\noindent\textbf{Exact execution and source locks.}\par
The primary certificate passes 1367 exact checks. The independent certificate passes 6129 checks, including 104 full integral bases, 468 signed finite-character instances, and 42 fresh symbolic nine-unknown prism solves. Every mathematical section and tagged display of both proofs is included above. The complete proof, code, input, audit, and page-review locks are recorded in \path{records/k115_export_manifest.json}.
\end{proof}


\subsection{The marked OS45 source cover and the twelve unselected-II disc returns}
\begin{theorem}[Source result P3K.116 under continuation review]\label{night:P3K.116}
\StatementRecord{night:P3K.116}{open}{Historical events: GPT-6 Astra (OpenAI), ChatGPT Pro mode, effort Pro, S6\_06\_BLIND, 16--18 November 2026: targeted proof check. Also checked in this writing. Continuation mathematical and source checks: GPT-6 Astra (OpenAI), ChatGPT Pro mode, effort Pro, S6\_10\_PAPERS continuation, 19 November 2026: computation or formal certificate; targeted proof check; used and rechecked in this run. Continuation metadata events: GPT-6 Astra (OpenAI), ChatGPT Pro mode, effort Pro, S6\_10\_PAPERS continuation, 19 November 2026: used and rechecked in this run; source-row intake; metadata reconciliation only; no new mathematical check. GPT-6 Astra (OpenAI), ChatGPT Pro mode, effort Pro, S6\_10\_PAPERS continuation, 19 November 2026: used and rechecked in this run; explicitly scoped metadata reconciliation; metadata reconciliation only; no new mathematical check. GPT-6 Astra (OpenAI), ChatGPT Pro mode, effort Pro, S6\_10\_PAPERS continuation, 19 November 2026: used and rechecked in this run; explicitly scoped metadata reconciliation; metadata reconciliation only; no new mathematical check. GPT-6 Astra (OpenAI), ChatGPT Pro mode, effort Pro, S6\_10\_PAPERS continuation, 19 November 2026: used and rechecked in this run; receiving source placement; metadata reconciliation only; no new mathematical check. \emph{Continuation recorded proof scope.} Two written prereading derivations cover the marked source curve, its full puncture set, source submersion, literal rational map and inverse, unit-preserving infinity charts, directed Hom degree, and the disc/scalar-annulus argument. The complete supplied arguments were then compared. A separate root derivation, frozen before the residual targets were opened, checks O45.18--23, O45P.14--19 and O45P.23--27 at the exact source, Miller, two-Cartier-frame, cusp-ring, full-theta and primitive-period scopes recorded in \path{analysis/continuation_root/K116_residual_comparison.md}; its fresh program passes81 exact checks. The fixed initial coefficient now explicitly defines the continuation family; an independently prescribed varying linearization retains its actual evaluation ratio. Earlier fresh28-check and later42-check programs retain their respective before/after-reading chronology. The original K10 topology, all-fan input chain, and remaining whole printed-proof gates are still pending. This is a collection of scoped checks and does not close the whole receiving proof gate; the receiving status remains open.}{the original source, point line, roots, characters, reference, selected fan and scalar range stated here}
Retain the literal OS45 source, marked section and point line
\[
\begin{gathered}
 R=t^2+a,\qquad f=R^2+t,\qquad
 g=R^3+\frac32tR+\frac38,\\
 S_a:\ y^2=x^3-3fx+2g,\qquad
 G=(R,\sqrt3/2),\qquad M_8=\mathcal O([8G]-[O]).
\end{gathered}
\]
The suspension source map is translation by \(-G\); its simple zero is
the whole selected \(I_1\) fibre. Retain the original multiplicity-one
cap at the ordinary value \(t=0\), the complete prescribed line and
annular reference, and any prescribed selected lattice-basis fan.

The ordered selected-node source cover and its literal source map are
\[
\overline{\mathcal C}:\ v^2=4u^3+3,\qquad
a=-\frac u3-\frac{u^4}{9},\qquad
t_c=\frac1{2u}-\frac{u^2}{3},\qquad
\xi=\frac1{4u^2}+\frac u3.
\]
On the selected normalization the marked branches, section and inverse
torus coordinate are
\[
x=S^2-2\xi,\qquad y=S(S^2-3\xi),\qquad
c=\frac v{2u},\qquad
Z=\frac{S-c}{S+c},\qquad
P_G=\frac{\sqrt3+v}{\sqrt3-v}.
\]
The open family is obtained by removing exactly the 48 points specified
in O45.2. Its fundamental group has the full marked presentation
\[
[A,B]\mu_1\cdots\mu_{48}=1,
\]
with the two handle loops and the individually specified meridians.

There are exactly twelve points of \(\overline{\mathcal C}\) with
\[
16u^6+120u^3+243=0.
\]
At each of them two unselected nodal fibres combine into a type-II
fibre, while the selected fibre remains nodal, its branch order remains
marked, and the original cap value remains ordinary. On a sufficiently
small disc about each point, the complete original source, \(M_8\)
linearization, prescribed selected fan, infinity attachment and
ordinary cap extend to a proper holomorphic submersion for a fixed
common strict-contraction scalar. Its boundary objects are the original
OS45 objects through the full coefficient maps and their inverses.
The original six-sphere topology is carried by the explicit path family
and the resulting equivariant horizontal flow.

Each of the twelve marked unselected-II meridians has physical return
isotopic to the identity, and hence physical return class \(0\).
For every scalar admitted on the original boundary family, the
constant-argument scalar interval and the full cap logarithm give the
actual annular comparison with this disc family, so the same conclusion
holds at that scalar. Every source, point-line, finite character,
period integer and endpoint marking carried by the original
multiplicity-one construction is retained. The other marked generators
retain their explicit source labels in O45.2; their return classes are
outside this result.
\end{theorem}
\begin{proof}[Supplied source argument; continuation review in progress]

% Expanded verified cp36 input: notes/part3_os45_unselected_ii_proof.tex
\par\bigskip\noindent\textbf{Complete supporting proof 1: O45.}\par


% Expanded verified cp36 input: notes/part3_os45_unselected_ii_o45_proof.tex
\emph{Audit provenance.} The preparation-stage sentence in O45.7 is retained verbatim as part of the original proof. The completed full geometric audits and their exact input locks are recorded with the complete supporting proofs below and in \path{records/k116_export_manifest.json}. The 76 primary, 88 root-independent and 225 further independent exact checks and all full geometric audit gates have passed.

\par

\medskip\noindent\textbf{The OS45 marked-source cover and the unselected-II disc attachments}

\par\nobreak

This proof retains the literal source family used in K10 and its
original multiplicity-one cap at the ordinary value \(t=0\). The line
is \(M_8=\mathcal O([8G]-[O])\), the suspension source map is
translation by \(-G\), and the simple zero is the whole selected
\(I_1\) fibre. The scalar-loop calculation is already covered by
K73. Here the parameter moves the marked source. All source,
point-line, selected-height and cap arrows below are specified on
that family. Historical K10 return values are not used.

\par

\medskip\noindent\textbf{O45.1. The literal source and its marked-node cover}

\par\nobreak

Write
\[
\begin{gathered}
 R=t^2+a,\qquad f=R^2+bt,\qquad
 g=R^3+\frac32btR+\frac38b^2,\\
 S_{a,b}:\ y^2=x^3-3fx+2g,\qquad
 G=(R,\sqrt3 b/2).
 \end{gathered}
\tag{O45.1}
\]
The fixed square root is the one in K10. The following scaling is an
actual map of the source, section and relative invariant differential:
\[
 (t,x,y)=(sT,s^2X,s^3Y),\qquad
 a=s^2A,\quad b=s^3B,\qquad
 \frac{dx}{y}=s^{-1}\frac{dX}{Y}.
 \tag{O45.2}
\]
The reciprocal substitutions are its inverse. We work on the
explicit \(B=1\) slice after adjoining the specified cube root of
\(b\ne0\). Thus no scalar or source-coordinate identification is
being made without this map.

Set \(D=64(g^2-f^3)/b^2\). Direct expansion and elimination give
\[
\begin{split}
 D(t;a,b)&=48a(t^2+a)^2+8bt^3+72abt+9b^2,\\
 \operatorname{disc}_tD&=331776b^2(256a^3-3b^2)^3,\\
 D(0;a,b)&=3(16a^3+3b^2).
 \end{split}
\tag{O45.3}
\]
The Weierstrass discriminant is \(-27b^2D\). For \(a\ne0\),
the minimal infinity coefficients are units and its order is eight.
Away from \(256a^3=3b^2\), the other four fibres are simple \(I_1\)\textquotesingle s.
At a finite repeated discriminant value, the equations below give
the actual cusp rather than a double nodal total-space singularity.

For a selected double cubic root \(x=\xi\), put
\(u=\xi-(t_c^2+a)\). The equations \(f=\xi^2\), \(g=\xi^3\)
are equivalent, with \(b,u\ne0\), to
\[
 a=-\frac u3-\frac{u^4}{9b^2},\qquad
 t_c=\frac b{2u}-\frac{u^2}{3b},\qquad
 \xi=\frac{b^2}{4u^2}+\frac u3.
 \tag{O45.4}
\]
To see the inverse directly, translate \(x=R+u\) in the cubic.
The double-root equations become
\(bt=u(u+2R)\) and \(u^2(2u+3R)=3b^2/4\). They yield
(O45.4), while \(u=\xi-(t_c^2+a)\) recovers the original marked
point. This is the actual four-sheeted selected-node cover of
the \(a\)-line on \(b=1\), with no discriminant root renamed.
Its derivatives and selected simplicity are
\[
 \frac{da}{du}=-\frac{4u^3+3b^2}{9b^2},\qquad
 \frac{dt_c}{du}=-\frac{4u^3+3b^2}{6bu^2},\qquad
 D_t(t_c)=-\frac{2(4u^3+3b^2)^3}{27b^3u^2}.
 \tag{O45.5}
\]

The entire selected nodal normalization is
\[
 x=S^2-2\xi,\qquad y=S(S^2-3\xi),\qquad
 S_G=-\frac{\sqrt3 b}{2u}.
 \tag{O45.6}
\]
Adjoin \(v\) with \(v^2=4u^3+3b^2\), and set \(c=v/(2u)\).
Then \(c^2=3\xi\), the two marked branches are \(S=\pm c\),
and the actual zero-section-normalized torus coordinate is
\[
 Z=\frac{S-c}{S+c},\qquad
 P_G=\frac{\sqrt3 b+v}{\sqrt3 b-v}.
 \tag{O45.7}
\]
It sends \(O\) to one and \(G\) to the displayed point. Its
inverse is \(S=c(1+Z)/(1-Z)\). Branch reversal \(v\mapsto-v\)
therefore gives \(Z\mapsto Z^{-1}\), by these literal maps.
No same-fan endpoint is being assigned to that reversal here.

On \(b=1\) the compact ordered cover is the smooth elliptic curve
\(\overline{\mathcal C}:v^2=4u^3+3\). Its residual cube deck is
\[
 (u,v)\longmapsto(\zeta_3u,v),\quad
 (a,t,x,y)\longmapsto(\zeta_3a,\zeta_3^2t,\zeta_3x,y).
 \tag{O45.8}
\]
Substitution in (O45.1) proves this source isomorphism and its
reciprocal. The quotient coordinate is \(v\); its source equation
is still used through the full cubic cover and this endpoint map.

\par

\medskip\noindent\textbf{O45.2. Every boundary point on this ordered cover}

\par\nobreak

Pulling back (O45.3) gives the exact factorization
\[
 256a^3-3b^2
 =-\frac{(4u^3+3b^2)^2
 (16u^6+120b^2u^3+243b^4)}{729b^6}.
 \tag{O45.9}
\]
The discriminant of the three unselected roots, after dividing
the actual polynomial \(D\) by \(t-t_c\), is
\[
 -\frac{1024u^4(16u^6+120b^2u^3+243b^4)^3}{6561b^{10}}.
 \tag{O45.10}
\]
Thus on \(b=1\) the following list gives the complete removed
points for the open source family with four simple nodes,
ordered selected branches and ordinary cap \(t=0\).

\begin{longtable}[]{@{}
  >{\raggedright\arraybackslash}p{(\columnwidth - 4\tabcolsep) * \real{0.3000}}
  >{\raggedleft\arraybackslash}p{(\columnwidth - 4\tabcolsep) * \real{0.4000}}
  >{\raggedright\arraybackslash}p{(\columnwidth - 4\tabcolsep) * \real{0.3000}}@{}}
\toprule\noalign{}
\begin{minipage}[b]{\linewidth}\raggedright
Exact locus on \(\overline{\mathcal C}\)
\end{minipage} & \begin{minipage}[b]{\linewidth}\raggedleft
Number of points
\end{minipage} & \begin{minipage}[b]{\linewidth}\raggedright
Source/marking event
\end{minipage} \\
\midrule\noalign{}
\endhead
\bottomrule\noalign{}
\endlastfoot
\(u=0, v=\pm\sqrt3\) & 2 & The selected node approaches infinity at \(a=0\). \\
\(u^3=-3, v=\pm3i\) & 6 & The selected finite node is retained while the infinity fibre becomes \(I_9\). \\
\(4u^3+3=0, v=0\) & 3 & The selected cubic becomes cuspidal. \\
\(16u^6+120u^3+243=0\) & 12 & Two unselected nodes combine into a type-II fibre. \\
\(2187-16u^3(u^3+3)^3=0\) & 24 & The original ordinary cap value \(t=0\) becomes singular. \\
The point at infinity & 1 & The normalized source parameter leaves the finite \(b=1\) chart. \\
\end{longtable}

All rows are disjoint, and every polynomial in the finite rows
has simple roots after the indicated branch multiplicities.
The cap row factors as
\[
 2187-16w(w+3)^3
 =-(2w-3)(8w^3+84w^2+342w+729),\qquad w=u^3.
 \tag{O45.11}
\]
Its \(2u^3=3\) part is exactly \(t_c=0\). The other part retains
the selected node and places an unselected node at the cap.

In the \(v\)-quotient the removed set consists of
\[
 \begin{gathered}
 v=\pm\sqrt3,\quad\infty,\quad\pm3i,\quad0,\\
 v^4+24v^2+162=0,\\
 (v^2-9)(v^6+33v^4+459v^2+4131)=0.
 \end{gathered}
\tag{O45.12}
\]
There are 18 points there. The three points of ramification
of (O45.8) are already removed; the degree-three map between
these open curves is unramified. Its Euler identity is
\(-48=3(2-18)\). On the ordered cover a complete source-loop
system consists of two handle loops and meridians at the 48
listed points, with the single relation
\[
 [A,B]\,\mu_1\cdots\mu_{48}=1.
 \tag{O45.13}
\]
This is an explicit surface presentation, hence 49 free
generators after eliminating one meridian. A cut polygon for
the double cover of the \(u\)-sphere, branched at the three
roots of \(4u^3+3\) and infinity, supplies \(A,B\); paths to
the specified finite algebraic points supply the meridians.
Their markings are transported by (O45.4),(O45.6),(O45.7).
The twelve unselected-II meridians are filled below. The
selected-cusp, \(I_9\), ordinary-cap and remaining cover
generators stay individually labelled; this proof assigns
no value to them.

\par

\medskip\noindent\textbf{O45.3. The actual unselected-II source disc and translation}

\par\nobreak

Fix one point \((u_*,v_*)\) in the unselected-II row. Choose
a small closed analytic disc \(\Delta\) around it avoiding
all other rows. Since \(da/du\ne0\) and \(v_*\ne0\), the
full source and ordered selected branch vary holomorphically
on this disc. Write \(a_*=a(u_*)\), \(c_*=-4a_*\). Then
\[
 c_*^3=-3/4,\qquad a_*=-c_*/4,\qquad
 t_{\mathrm{II}}=-c_*^2=-16a_*^2,\qquad R_{\mathrm{II}}=-c_*.
 \tag{O45.14}
\]
At this point \(f=g=0\). The full coefficient Jacobian is
\[
 \det\frac{\partial(f,g)}{\partial(a,t)}
 =\frac32c_*^2\ne0,\qquad
 \partial_t(x^3-3fx+2g)|_{x=0}=\frac32c_*\ne0.
 \tag{O45.15}
\]
Consequently \((a,x,y)\) are actual local coordinates on the
source total space, with \(t=t(a,x,y)\) the unique holomorphic
solution of its full displayed equation. In particular the
source surface and its projection to \(\Delta\) are smooth
through the cusp. The value \(D(0)=153/16\) retains the
ordinary cap. The selected value is separate because
\[
 t_c+16a^2=\frac{(2u^3+9)(4u^3+3)^2}{162u}
 \tag{O45.16}
\]
is nonzero at every unselected-II point.

The actual translation \(T=-G\), on its affine chord chart,
is
\[
 L=\frac{y+\sqrt3/2}{x-R},\qquad
 x_T=L^2-x-R,\qquad
 y_T=L(x-x_T)-y.
 \tag{O45.17}
\]
At the target the plus-\(G\) slope is \(-L\); its full chord
formula returns \(x,y\). At the cusp \(x-R=c_*\ne0\), so
both (O45.17) and its inverse are holomorphic. They fix the
cusp. In the \((x,y)\) source coordinates their derivative is
\(\left(\begin{smallmatrix}1&2g_G\\0&1\end{smallmatrix}\right)\),
where \(g_G=x(G)/y(G)=-2c_*/\sqrt3\). Thus this is a smooth
source morphism through the complete cusp, with its full
source-parameter derivatives retained. Away from this chart
the usual point/O addition charts and their actual reciprocal
chords apply. At infinity the fixed \(I_8\) simultaneous
resolution and its section marking are retained, as specified
in O45.5 below.

\par

\medskip\noindent\textbf{O45.4. The complete \(M_8\) coefficient, including the cusp module}

\par\nobreak

Use the actual doubling points \(2G,4G,8G\) from (O45.1).
For \(j=1,2,4\), write \(\ell_{jG}\) for the tangent line
at \(jG\), evaluated at the source point, and put
\[
 f_2=\frac{\ell_G}{x-x(2G)},\qquad
 f_4=f_2^2\frac{\ell_{2G}}{x-x(4G)},\qquad
 f_8=f_4^2\frac{\ell_{4G}}{x-x(8G)}.
 \tag{O45.18}
\]
These are complete rational functions, with every addition
coefficient taken from (O45.1). Their horizontal divisors
are \(m[G]-[mG]-(m-1)[O]\). Define
\[
 B_{8,1}(p)=f_2(p)^8\frac{f_8(Tp)}{f_8(p)}.
 \tag{O45.19}
\]
For a meromorphic point section \(\sigma_8\), with divisor
\([8G]-[O]\), its directed line map is
\(\sigma_8(p)\mapsto B_{8,1}(p)\sigma_8(Tp)\).
Indeed its horizontal divisor is exactly
\([G]+[8G]-[9G]-[O]=D-T^*D\), \(D=[8G]-[O]\).
Thus it is the actual Hom map \(M_8\to T^*M_8\).
It is also the pullback by \(T\) of the original
\(\operatorname{Hom}(t_G^*M_8,M_8)\) linearization: its
evaluation point is \(p-G\). This specifies the direction.

No finite whole-fibre coefficient is suppressed in (O45.19).
At each finite irreducible fibre the orders of \(f_8\circ T\)
and \(f_8\) are equal. The denominator of the \(G\)-doubling
has constant \(2y(G)=\sqrt3\), so \(f_2\) has no finite
whole-fibre zero or pole. The actual point and O generators
cancel every remaining horizontal coefficient, including
section coincidences. At \(I_8\), \(8G\) and \(O\) meet the
identity component, so the Hom has multidegree zero. The
cyclic intersection matrix has kernel generated by the full
fibre. Its degree is one: in the present direction the height
formula gives \(G\cdot O=0\), \(8G\cdot O=3\),
\(9G\cdot O=5\), hence
\(\deg\operatorname{Hom}(M_8,T^*M_8)=5-(3+1)=1\).
Equivalently the original direction uses
\(8G\cdot O+1-7G\cdot O+G\cdot O=3+1-3=1\).
The kernel and degree identify its full vertical divisor as
the single infinity fibre. Therefore the literal map
\[
 B_\lambda(p,v_M)=
 \bigl(Tp,\lambda C_* (t-t_c)B_{8,1}(p)v_M\bigr)
 \tag{O45.20}
\]
extends across infinity and has precisely the selected
whole-fibre zero. Here is the definition of the continuation
family used in this proof. Let \(a_0\) denote its initial marked
source and \(L_{a_0}\) the prescribed directed whole-zero line map.
At a protected regular point \(s_0\), evaluate both maps in the
same Cartier frames and set
\[
 C_*=
 \frac{L_{a_0}(s_0)}
 {(t(s_0)-t_c(a_0))B_{8,1,a_0}(s_0)}\ne0.
\]
The ratio is an invertible function on the compact initial source,
so it is constant on that source. Thus (O45.20) agrees with
\(\lambda L_{a_0}\) on the entire initial source. For every
subsequent marked parameter, the continuation family is defined
by (O45.20) with this fixed \(C_*\). The full coefficient maps
and inverses establish the analytic extension of this defined
family. Its boundary family is its restriction with the prescribed
initial cap and reference continued by the displayed maps.

The initial evaluation proves equality at \(a_0\); it does not
identify an independently prescribed varying linearization
\(L_a\) with the defined continuation. Its exact fibrewise ratio,
evaluated in the same frames on a protected section \(s(a)\), is
\[
 c(a)=\frac{L_a(s(a))}
 {(t(s(a))-t_c(a))B_{8,1,a}(s(a))}.
\]
Such an additional comparison retains \(c(a)\) in the full
suspension coefficient. No constancy or disc extension of that
additional ratio is deduced from the one initial evaluation.

Here is the local cusp coefficient in full, beyond its divisor.
On the central cuspidal curve put \(p=x/y\), so
\(x=p^{-2},y=p^{-3}\), and \(G\) has additive parameter
\(g_G\ne0\). Its multiples have parameters \(mg_G\), and
the literal source map is \(p\mapsto p-g_G\). The full
addition formulas (O45.18) give
\[
 f_m(p)=\frac{(-1)^{m-1}m}{g_G^{m-1}}
 \frac{(p-g_G)^m}{(p-mg_G)p^{m-1}},\quad m=2,4,8,
 \qquad
 B_{8,1}(p)=\frac{256}{g_G^8}
 \frac{(p-g_G)(p-8g_G)}{p(p-9g_G)}.
 \tag{O45.21}
\]
On the finite normalization chart the point/O generator ratio
is \(r(p)=(p-8g_G)/p\). Hence its genuine coefficient is
\[
 B_{8,1}(p)\frac{r(p-g_G)}{r(p)}=\frac{256}{g_G^8}.
 \tag{O45.22}
\]
This chart includes the two strict section directions through
their actual Cartier generators. At the cusp itself the frame
is \(\sigma_8\), because neither horizontal divisor passes
through it. Its coefficient, in the actual cusp local ring,
is
\[
 B_{8,1}(x,y)=\frac{256}{g_G^8}
 \frac{1-73g_G^2x+72g_G^3y}{1-81g_G^2x}.
 \tag{O45.23}
\]
Both numerator and denominator are units at \((0,0)\).
The overlap between these two frames is the full \(r(p)\).
In particular the point/O frame is not being extended across
the cusp by deleting its linear normalization coefficient.
Equations (O45.22),(O45.23) give the precise gluing instead.
The inverse coefficient is evaluated at the actual inverse
source point. Thus the source, point line and entire Hom map
extend on this disc with no charged or deleted cusp point.

\par

\medskip\noindent\textbf{O45.5. Both retained collars and the original ordinary cap}

\par\nobreak

At the selected value put \(q=t-t_c(u)\), preserving this
original height. Let \(\xi(q,u)=\sqrt{f(t_c+q,u)}\) be the
branch with (O45.4) at \(q=0\). It is a unit throughout the
small selected collar. Write

\[
 E(q,u)=g(t_c+q,u)-\xi(q,u)^3,\quad
 h(q,u)=2E(q,u)/q.
 \tag{O45.24}
\]
Here \(h(0,u)=-2u^4\), hence \(h\) is an actual extending
unit. With the ordered square root of \(x+2\xi\), put
\[
 X=y-(x-\xi)\sqrt{x+2\xi},\qquad
 Y=\frac{y+(x-\xi)\sqrt{x+2\xi}}{h(q,u)}.
 \tag{O45.25}
\]
Then \(XY=q\) literally. Conversely, \(q=XY\),
\(y=(X+hY)/2\), and the implicit equation
\(2(x-\xi)\sqrt{x+2\xi}=hY-X\) has derivative
\(2\sqrt{3\xi}\ne0\), giving its unique holomorphic
inverse. Thus the actual nodal source chart and every unit
are specified before applying the cusp fan.

Let \(Q=\kappa(q,u)q\) be the actual Tate parameter from
the original \(j\)-function, with \(Qj(Q)=1+O(Q)\), and
let \(Z\) be the inverse marked uniformization agreeing
with (O45.7). Define \(P=Z(G)\) on the full smooth collar.
Both its differential scale and the inverse uniformization
are chosen to agree with the original reference; their
unique analytic continuations over \(\Delta\) retain all
higher coefficients. For the complete theta product
\(f_b(Z)=\Theta_Q(Z/b)/\Theta_Q(Z)\), the point coefficient
\(V_{\rm pt}=f_{P^8}(Z)v_M\) has exact periods
\[
 h:(Z,V_{\rm pt})\mapsto(QZ,P^8V_{\rm pt}),\qquad
 B:(Z,V_{\rm pt})\mapsto(P^{-1}Z,\lambda q\beta V_{\rm pt}),
 \tag{O45.26}
\]
where, as a point-module map,
\[
 \beta=C_* B_{8,1}(p)
 \frac{f_{P^8}(P^{-1}Z)}{f_{P^8}(Z)}.
 \tag{O45.27}
\]
The horizontal divisors cancel in this displayed direction,
so (O45.27) is independent of the point \(p\). Evaluation
at one protected regular \(Z_*\) gives its full analytic
coefficient and inverse. It is a unit by (O45.20). The
usual selected toric charts, their prescribed fan and every
period translate therefore vary holomorphically over the
entire disc. Their source height is the same \(q\); the
units \(\kappa,P,\beta\) have not been discarded.

At infinity retain \(z=1/t\) and the minimal coordinates
\((z^2x,z^3y)\). The exact source has a split node with
smoothing parameter \(z^8\kappa_\infty(z,u)\), where
\(\kappa_\infty(0,u)=-a(u)/2304\). Choose its branch
so \(P_\infty=z^3\pi(z,u)\), \(\pi(0,u)=-1/24\).
These functions are the full inverse Tate uniformization,
not their leading terms. The source coordinates have their
actual consecutive charts \(X_iY_i=z\), obtained from
\(X_i=Z/z^i,Y_i=z^{i+1}/Z\), with the corresponding full
source units. Their overlaps and translation are exactly
the monomial arrows obtained by the displayed \(Q,P\).
The section remains on component three; its inverse
translation remains the shift by minus three.

The actual identity-component point of \(8G\) is
\[
 b_0=P_\infty^8/Q_\infty^3,\qquad
 b_0(0,u)=-1/(9a(u)^3).
 \tag{O45.28}
\]
Three original group-law doublings verify this value, as
well as the full minimal coordinates
\[
 X_{8,0}=\frac{81a^6-90a^3+1}{(9a^3+1)^2},\quad
 (Y_{8,0}/\sqrt3)=
 \frac{108a^3(9a^3-1)}{(9a^3+1)^3}.
 \tag{O45.29}
\]
At the II centre \(b_0=-256/27\ne1\), so both point/O
modules are separated on this infinity collar. The full
theta identity is
\[
 f_{Q^3b_0}(Z)=(-Z/b_0)^3Q^{-6}f_{b_0}(Z).
 \tag{O45.30}
\]
It gives the actual reciprocal line-frame arrow, including
the base power \(Q^{-6}\). In the \(b_0\) point frame the
suspension coefficient is \(P_\infty^3\) times its original
point-theta coefficient, and the source/line valuation
vectors are exactly \((8,0),(-3,0)\). The primitive periods
are the full words
\[
 H=h^{-1}B^{-3},\qquad K=h^3B^8,\qquad
 h=H^8K^3,\quad B=H^{-3}K^{-1}.
 \tag{O45.31}
\]
Their integer matrix has determinant one. The line
coefficient of \(K\) is \(\lambda^8\) times its complete
analytic unit. Thus the original nonzero-line \(I_8\)
quotient and every source component persist on the entire
disc, with literal source and line inverses. This is the
same rank-one filling as the source line complement:
on each consecutive source chart its map is the above
identity on the actual point module and the indicated
integer period reindex. Their equations and overlaps
agree, so the maps glue across the whole cycle.

Finally choose the fixed original cap disc around \(t=0\),
disjoint from the selected and unselected source collars.
All its source fibres are smooth over \(\Delta\). Continue
the prescribed full periods and their initial logarithms
over this contractible parameter disc. In particular
\(\gamma(t,u,\lambda)\) is the actual suspension period of
(O45.20), with its full source, line and scalar components.
The original cap-to-ordinary map and inverse are
\[
 \Psi(\zeta,t,u)=
 \left(\zeta+\frac{\log t}{2\pi i}\gamma(t,u,\lambda),t,u\right),
 \qquad
 \Psi^{-1}(\zeta,t,u)=
 \left(\zeta-\frac{\log t}{2\pi i}\gamma(t,u,\lambda),t,u\right).
 \tag{O45.32}
\]
A logarithm branch change adds exactly the original
\(\gamma\)-period. Thus this is the unchanged order-one
cap with its marked meridian and physical line, now
holomorphic over \(\Delta\). All initial annular
integers are carried by analytic continuation from the
reference. There is no source-parameter deck around the
filled disc: every displayed source, module, theta and
cap map returns to itself.

\par

\medskip\noindent\textbf{O45.6. Compact disc family and the filled source meridians}

\par\nobreak

The punctured ordered curve of O45.2 is connected. Choose
a path in it from one of the original K10 \(a=b=1\)
markings to the chosen boundary point of \(\Delta\).
Cover that path by finitely many small discs avoiding
all rows of the table. On each disc the same original
equation, source inverses, simultaneous \(I_8\) charts,
selected fan and cap map (O45.32) apply. On overlaps
the rational point line and (O45.20) are literally
the same; source period changes are their actual
integral changes of basis, with the initial cap
period and reference continued along the path.
A smaller common scalar and its constant-argument
interval identify the original admitted boundary
family with these proper submersions. This carries
the known K10 sphere to the boundary of the present
disc by explicit source/line/cap families; no K10
return value enters this transport.

Choose a Hermitian metric on the actual point line over
the compact source family outside the selected tube.
Equations (O45.17)-\/-(O45.23) give a uniform finite norm
for the full unscaled lift, including the II centre.
Choose the fixed scalar with norm times this bound less
than one. The full suspension is then strictly
contracting there. The quotient is free: a physical
circle equality with \(B^k\) forces \(k=0\) by the norm
inequality. A compact annular fundamental domain gives
properness and Hausdorffness uniformly over the closed
source parameter disc.

Use the actual selected fan charts (O45.25)-\/-(O45.27)
inside the excluded tube, the unchanged infinity
charts (O45.28)-\/-(O45.31), and the original cap
(O45.32). Their overlaps are the full source, point
module and period arrows already written. There are
finitely many chart and period representatives on
each compact collar. They glue to a proper holomorphic
family over \(\Delta\). The source projection is a
submersion at the unselected cusp by (O45.15), at the
selected fan by its actual toric coordinates, and at
the other pieces by their unchanged source or cap
coordinates. Hence the completed family is a smooth
proper submersion over the whole disc.

On its punctured boundary these are exactly the
original OS45 objects with the selected whole \(I_1\),
\(M_8\), translation \(-G\), fixed fan and original
ordinary cap, by the explicit line-reference match
in (O45.20). The original sphere fibres transport
through this actual proper submersion. A connection
on this disc family, averaged over the physical
circle, retains its line action and provides the
disc trivialization. Its boundary return is isotopic
to the identity. Thus each of the twelve marked
unselected-II meridians has zero physical return;
this follows from its displayed filling, with all
source and cap markings retained.

The same conclusion holds for every scalar admitted
on the original boundary family. To compare it with
the common disc-contraction scalar, decrease its
modulus along the actual constant-argument interval.
On the boundary all intermediate lifts remain
contracting. Carry the full \(\log\lambda\) in
\(\gamma\), (O45.27), (O45.31), and (O45.32) along
that interval. This gives the exact relative
parameter family and inverse, with no extra scalar
turn. The smaller endpoint fills by the disc just
constructed, and the boundary annulus transports
that filling to the original admitted scalar.

This closes the unselected-II source generators in
the ordered cover. The selected-cusp, \(I_9\),
ordinary-cap-collision and remaining source-cover
generators of (O45.13) remain outside this completed
disc calculation. Their precise source equations
and parameter labels have been retained above.

\par

\medskip\noindent\textbf{O45.7. Independent exact execution}

\par\nobreak

The executable
\path{code/part3_fourfold/os45_marked_source_cover.py}
derives the discriminant, selected inverse cover,
ordered normalization, boundary factorization,
full source inverse, actual cusp Jacobian, complete
addition-derived \(M_8\) functions and their two
Cartier charts. It also computes the original
threefold doubling limit of \(8G\), both directed
Hom degrees, the integral \(I_8\) kernel and full
primitive-period inverse. The local algebra is
supplemented by the complete compact source,
quotient and cap construction above. Existing
source files and scalar-return certificates are
unchanged. This working proof is awaiting full
independent geometric audit before numbered export.

% End expanded verified cp36 input


\par\bigskip\noindent\textbf{Complete supporting proof 2: O45R.}\par


% Expanded verified cp36 input: notes/part3_os45_unselected_ii_o45r_proof.tex
\par

\medskip\noindent\textbf{Independent root proof of the OS45 unselected-II attachments}

\par\nobreak

We retain the whole original source, point line, selected height,
ordinary cap and fixed scalar of O45.1-\/-32. In particular
\[
 R=t^2+a,\quad f=R^2+t,\quad
 g=R^3+\frac32tR+\frac38,\quad
 y^2=x^3-3fx+2g,\quad G=(R,\sqrt3/2).
 \tag{O45R.1}
\]
The chosen square root is fixed. The point line is
\(M_8=\mathcal O([8G]-[O])\), the source translation is
\(T:p\mapsto p-G\), and the ordinary logarithmic cap
at \(t=0\) has its full inward suspension clutch.
The initial source is the literal K10 construction.
Its old scalar-return sign is not used in this source-disc proof.

\par

\medskip\noindent\textbf{O45R.1. Original parameter maps and the exact collision discs}

\par\nobreak

The weighted map
\((t,x,y)=(sT,s^2X,s^3Y)\), \(a=s^2A,b=s^3B\)
has inverse given by division by those powers. Substitution
in the equation multiplies it by \(s^6\); it carries the
actual section and scales \(dx/y\) by \(s^{-1}\).
Thus the \(b=1\) slice is taken on the explicitly marked
cube-root cover of \(b\ne0\). We keep that cover.

Direct differentiation of the original cubic at a selected
double root \(\xi\), followed by putting
\(u=\xi-(t_c^2+a)\), gives
\[
 t_c=u(u+2R_c),\qquad
 u^2(2u+3R_c)=\frac34,\qquad
 R_c=t_c^2+a .
 \tag{O45R.2}
\]
Solving these two equations gives
\[
 a=-u/3-u^4/9,\quad
 t_c=1/(2u)-u^2/3,\quad
 \xi=1/(4u^2)+u/3.
 \tag{O45R.3}
\]
Their inverse is the original difference \(u=\xi-R_c\).
The polynomial
\[
 D=48a(t^2+a)^2+8t^3+72at+9
 \tag{O45R.4}
\]
satisfies \(\Delta=-27D\). Eliminating \(t\) gives
\(\operatorname{disc}_tD=331776(256a^3-3)^3\).
The selected-root derivative is
\[
 D_t(t_c)=-\frac{2(4u^3+3)^3}{27u^2}.
 \tag{O45R.5}
\]
Consequently the branch cover is the literal elliptic curve
\(v^2=4u^3+3\), with the excluded loci specified in O45.2.
On its selected fibre the maps
\[
 x=S^2-2\xi,\quad y=S(S^2-3\xi),\qquad
 Z=\frac{S-v/(2u)}{S+v/(2u)}
 \tag{O45R.6}
\]
have inverses \(S=y/(x-\xi)\) and
\(S=(v/(2u))(1+Z)/(1-Z)\) on their indicated charts.
They send the original \(G\) to
\((\sqrt3+v)/(\sqrt3-v)\). Inversion of the branch marking
is literally \(Z\mapsto Z^{-1}\). We do not use it as a
closing of a fixed fan.

The unselected collision polynomial is
\[
 H(u)=16u^6+120u^3+243.
 \tag{O45R.7}
\]
It has six simple roots, disjoint from \(u=0\),
\(u^3=-3\), \(4u^3+3=0\) and
\(2187-16u^3(u^3+3)^3=0\). Each has two distinct
values of \(v\), so there are precisely twelve
marked disc centres. These are all the unselected
collisions: division of O45R.4 by its actual factor
\(t-t_c\) gives cubic discriminant
\(-1024u^4H(u)^3/6561\). The distinct selected
root and the ordinary cap therefore stay separate.

The other loci in O45.2 have, respectively, two, six,
three, twenty-four and one points. Direct polynomial
gcd calculations show that their rows are disjoint.
Their total with the twelve centres is forty-eight.
The curve \(v^2=4u^3+3\) has four simple branch
points over the \(u\)-sphere, hence genus one by
its actual double-cover Euler calculation. Deleting
the forty-eight points gives a connected surface.
The map to \(v\) is a degree-three cover, with its
three ramification points among the deleted points.
Its eighteen-puncture target has Euler number
\(-16\), and the source has Euler number \(-48\).
The surface presentation with two handle loops and
the forty-eight meridians has the single boundary
relation written in O45.13. The present disc maps
fill exactly the twelve meridians of O45R.7.

Fix one centre and a closed disc \(\Delta\) around it.
The maps in O45R.3 and the chosen branch \(v\) are
holomorphic on \(\Delta\); \(da/du\ne0\) there.
Write \(c=-4a_*\). Then the collision point is
\[
 c^3=-3/4,\quad a_*=-c/4,\quad
 t_*=-c^2,\quad R_*=-c,\quad f_*=g_*=0 .
 \tag{O45R.8}
\]
The independent original-equation derivatives are
\[
 \det\frac{\partial(f,g)}{\partial(a,t)}
    =\frac32c^2,\qquad
 \partial_t(x^3-3fx+2g)|_{x=y=0}
    =\frac32c .
 \tag{O45R.9}
\]
Both are nonzero. The full source near its cusp is
therefore the graph of a holomorphic function
\(t=t(a,x,y)\); projection to \(a\) is a submersion.
This is a chart on the original total source, with
inverse the original \((a,x,y)\) projection.
The cusp occurs at a finite value different from
the selected value, since
\[
 t_c+16a^2=
 \frac{(2u^3+9)(4u^3+3)^2}{162u}\ne0
 \quad\text{on }H(u)=0.
 \tag{O45R.10}
\]
The original cap remains ordinary because its
discriminant factor is \(D(0)=153/16\).

\par

\medskip\noindent\textbf{O45R.2. Full source translation and the two Cartier frames}

\par\nobreak

The actual translation by \(-G\) is
\[
 L=\frac{y+\sqrt3/2}{x-R},\qquad
 x_T=L^2-x-R,\qquad y_T=L(x-x_T)-y .
 \tag{O45R.11}
\]
Substitution reduces its source equation to the
original source equation. Its inverse uses the
plus-\(G\) slope
\((y_T-\sqrt3/2)/(x_T-R)=-L\); applying the
same chord formula gives \(x,y\) exactly.
At the cusp both denominators are \(c\ne0\),
both maps fix \(x=y=0\), and their derivatives
in the graph coordinates of O45R.9 are
\[
 DT=\begin{pmatrix}1&2g_G\\0&1\end{pmatrix},
 \qquad
 DT^{-1}=\begin{pmatrix}1&-2g_G\\0&1\end{pmatrix},
 \qquad g_G=-2c/\sqrt3 .
 \tag{O45R.12}
\]
This proves invertibility on the full cusp chart.
The original section and zero-section charts give
the same inverse on the rest of the source.

Let \(f_m\) be the actual Miller function obtained
by consecutive doublings of \(G\), for \(m=2,4,8\).
Its horizontal divisor is \(m[G]-[mG]-(m-1)[O]\).
The full rational coefficient
\[
 B_{8,1}=f_2^8\,\frac{f_8\circ T}{f_8}
 \tag{O45R.13}
\]
has horizontal divisor
\([G]+[8G]-[9G]-[O]\). The point-line frame
at the target has divisor \(T^*([8G]-[O])
=[9G]-[G]\). Thus multiplication by O45R.13
carries the actual two Cartier modules in the
stated direction. All horizontal zeros and
poles cancel in those modules, at their
intersections as well as away from them.

We reconstruct its central value as a function
before restricting it to a constant. Put
\(p=x/y\) on \(y^2=x^3\), so that
\(x=p^{-2},y=p^{-3}\) and \(G\) has additive
coordinate \(g_G\). The tangent at \(jG\)
has slope \(3/(2jg_G)\). Substituting this
tangent and \(x(2jG)=(2jg_G)^{-2}\) in its
chord quotient gives all three identities
\[
 f_m(p)=\frac{(-1)^{m-1}m}{g_G^{m-1}}\,
       \frac{(p-g_G)^m}{(p-mg_G)p^{m-1}},
 \qquad m=2,4,8 .
 \tag{O45R.14}
\]
Using \(Tp=p-g_G\) in O45R.13 now gives
\[
 B_{8,1}(p)=\frac{256}{g_G^8}
       \frac{(p-g_G)(p-8g_G)}{p(p-9g_G)}.
 \tag{O45R.15}
\]
The finite normalization chart uses the true
point/O ratio \(r(p)=(p-8g_G)/p\).
Its coefficient is
\[
 B_{8,1}(p)\frac{r(p-g_G)}{r(p)}
       =\frac{256}{g_G^8}.
 \tag{O45R.16}
\]
The cusp chart uses the original meromorphic
point frame, whose horizontal divisor avoids
the cusp. Multiplying numerator and denominator
of O45R.15 by the exact opposite factor
\(1+9g_G/p\) gives the literal original-ring expression
\[
 B_{8,1}(x,y)=\frac{256}{g_G^8}
       \frac{1-73g_G^2x+72g_G^3y}
            {1-81g_G^2x}.
 \tag{O45R.17}
\]
Both factors have value one at the cusp.
The overlap of these two frames is the full
function \(r\). In particular O45R.17, and
its reciprocal at the inverse source point,
are actual holomorphic module maps through
the cusp. The parameter graph O45R.9 and the
nonzero denominators in these formulas make
the maps holomorphic on a neighborhood in
the total family.

At every finite irreducible fibre the two
terms \(f_8\circ T\) and \(f_8\) have the same
whole-fibre order. The \(G\)-doubling denominator
is the fixed nonzero number \(\sqrt3\), and
the divisor of \(f_2\) has no finite
whole-fibre coefficient. Thus O45R.13 has
no finite vertical divisor after its
horizontal module cancellation. At infinity
\(8G\) and \(O\) lie on the identity
component, so the Hom has zero component
degrees. The \(I_8\) intersection matrix
has kernel \(\mathbb Z(1,\ldots,1)\).
The height \(h(G)=1/8\) and component
index \(3\) give
\[
 G\cdot O=0,\quad 7G\cdot O=3,\quad
 8G\cdot O=3,\quad 9G\cdot O=5 .
 \tag{O45R.18}
\]
Restriction of the directed Hom to \(O\)
has degree \(5-0-3-1=1\). Its entire
vertical divisor is consequently one
infinity fibre. Multiplication by \(t-t_c\)
replaces this by the selected whole zero,
giving exactly the original suspension
\[
 B_\lambda(p,v_M)=
  (Tp,\lambda C_*(t-t_c)B_{8,1}(p)v_M).
 \tag{O45R.19}
\]
The prescribed nonzero \(C_*\) is fixed by
one original regular-point evaluation.
The matching ratio is an invertible
function on the compact source and is
constant there. Continuing that constant
retains the initial scalar normalization.

\par

\medskip\noindent\textbf{O45R.3. The selected and infinity collars with every unit}

\par\nobreak

At the selected fibre use its original
\(q=t-t_c(u)\). Choose
\(\xi(q,u)=\sqrt{f(t_c+q,u)}\) with the
marked value O45R.3, and put
\[
 h(q,u)=\frac{2(g(t_c+q,u)-\xi(q,u)^3)}q.
 \tag{O45R.20}
\]
The numerator vanishes at \(q=0\), and its
quotient has value \(-2u^4\ne0\). Thus
\[
 X=y-(x-\xi)\sqrt{x+2\xi},\qquad
 Y=h^{-1}\bigl(y+(x-\xi)\sqrt{x+2\xi}\bigr)
 \tag{O45R.21}
\]
has exact equation \(XY=q\). The inverse
first puts \(q=XY\), \(y=(X+hY)/2\),
then solves
\(2(x-\xi)\sqrt{x+2\xi}=hY-X\).
Its derivative at the node is
\(2\sqrt{3\xi}\ne0\). Both directions are
therefore holomorphic on the complete
selected collar over \(\Delta\).

The full inverse Tate uniformization gives
\(Q=q\kappa(q,u)\), the marked \(Z\), and
\(P=Z(G)\), with their original reference
scales. The point line is carried by
\[
 V_{\mathrm{pt}}=
 \frac{\Theta_Q(Z/P^8)}{\Theta_Q(Z)}\,v_M.
 \tag{O45R.22}
\]
Its full period arrows are
\[
 h:(Z,V_{\mathrm{pt}})\mapsto(QZ,P^8V_{\mathrm{pt}}),
 \qquad
 B:(Z,V_{\mathrm{pt}})\mapsto
        (P^{-1}Z,\lambda q\beta V_{\mathrm{pt}}),
 \tag{O45R.23}
\]
where
\[
 \beta=C_* B_{8,1}(p)
       \frac{f_{P^8}(P^{-1}Z)}{f_{P^8}(Z)} .
 \tag{O45R.24}
\]
The displayed Cartier divisors cancel, and
the ratio is independent of \(p\) by
compactness of a regular source fibre.
Evaluation at a protected regular point
defines its complete analytic base function.
The same module divisor calculation and
its reciprocal prove that it is an
extending unit. Thus every coefficient of
the prescribed selected fan is analytic on
\(\Delta\). On an affine cone the transition
acts on each character by its actual
monomial times these full units; the
inverse monomial and reciprocal units
give its inverse. The ordered branch is
single-valued on \(\Delta\), so this
extends the original prescribed fan
without a fan permutation at its centre.

There is a separate check of the infinity
normalization. Put \(z=1/t\) and use the
original minimal coordinates
\(X=z^2x,Y=z^3y\). The leading inverse
\(j\)-coefficient gives
\[
 Q_\infty=z^8\kappa_\infty(z,u),\qquad
 \kappa_\infty(0,u)=-a/2304 .
 \tag{O45R.25}
\]
The central cubic is
\(Y^2=(X-1)^2(X+2)\). Its exact normalized
torus coordinate has
\[
 X=1+\frac{12Z}{(1-Z)^2},\qquad
 \frac{Y}{\sqrt3}
      =\frac{12Z(Z+1)}{(Z-1)^3}.
 \tag{O45R.26}
\]
The zero section has \(Z=1\); this fixes
the coefficient twelve and the branch
sign. For the original section \(G\),
\[
 X_G=1+az^2,\quad
 Y_G=(\sqrt3/2)z^3,\quad
 \sqrt{f_{\min}}=1+az^2+z^3/2+O(z^5).
 \tag{O45R.27}
\]
The marked component is three. The full
Tate expression for its branch has leading
term \(X_G-\sqrt{f_{\min}}=12P_\infty\);
the next Laurent terms have orders at least
five, from \(Q_\infty/P_\infty\), and six,
from \(P_\infty^2\). Its analytic scaling
has central value one by O45R.26.
It follows that
\[
 P_\infty=z^3\pi(z,u),\qquad
 \pi(0,u)=-1/24,\qquad
 \left.\frac{P_\infty^8}{Q_\infty^3}\right|_{z=0}
          =-\frac1{9a^3}.
 \tag{O45R.28}
\]
The full functions, rather than their
leading terms, define the subsequent
maps. Formula O45R.26 applied to the
last value also gives independently
\[
 X_{8,0}=\frac{81a^6-90a^3+1}{(9a^3+1)^2},
 \quad
 \frac{Y_{8,0}}{\sqrt3}
       =\frac{108a^3(9a^3-1)}{(9a^3+1)^3}.
 \tag{O45R.29}
\]
This agrees with the full original
addition-law calculation in O45.29.
At each unselected-II centre the point
value is \(-256/27\), and hence is
distinct from the zero section.

Put \(b_0=P_\infty^8/Q_\infty^3\).
The product definition of theta gives
the entire identity
\[
 f_{Q_\infty^3b_0}(Z)=
       (-Z/b_0)^3Q_\infty^{-6}f_{b_0}(Z).
 \tag{O45R.30}
\]
Its reciprocal is the actual change
from the original point frame to the
identity-component point frame. Under
translation by \(P_\infty^{-1}\), its
factor changes the full line coefficient
by \(P_\infty^3\). This retains the
base factor \(Q_\infty^{-6}\) in the
frame arrow. The valuation vectors
in that frame are \((8,0)\) and
\((-3,0)\).

The whole period reindex is
\[
 H=h^{-1}B^{-3},\quad K=h^3B^8,\qquad
 h=H^8K^3,\quad B=H^{-3}K^{-1}.
 \tag{O45R.31}
\]
Its determinant is one and its
source valuations are \(1,0\).
The second multiplier of \(K\)
is \(\lambda^8\) times its full
extending unit. The consecutive
source charts
\(X_i=Z/z^i,Y_i=z^{i+1}/Z\)
have equation \(X_iY_i=z\).
Their full \(Q_\infty/z^8\) and
\(P_\infty/z^3\) units give all
period and translation overlaps,
with the inverse overlaps given
by their reciprocals. Thus the
rank-one completed quotient and
the original source line
complement are identified chart
by chart, including the entire
eight-component cycle.

\par

\medskip\noindent\textbf{O45R.4. The original cap and the proper disc family}

\par\nobreak

The disc at \(t=0\) is smooth over
\(\Delta\), by \(D(0)=153/16\).
Its original period basis and
suspension logarithm continue
holomorphically on this
contractible parameter disc.
Retain their full vector
\(\gamma(t,u,\lambda)\).
The cap-to-ordinary map is
\[
 \Psi(\zeta,t,u)=
  \left(\zeta+\frac{\log t}{2\pi i}
           \gamma(t,u,\lambda),t,u\right).
 \tag{O45R.32}
\]
The inverse subtracts precisely
this full vector. A logarithm
branch change adds one actual
\(\gamma\)-period, so both maps
are single-valued on the
original torus quotient. This
preserves the original
multiplicity-one cap, its
nonzero meridian clutch and
every initial annular integer.
It also preserves the
physical line: the
coefficient vector contains
its original line coordinate
and the complete
\(\log\lambda\).

All the full source and line
maps above are holomorphic
over the closed parameter
disc, and the selected
coefficient has its sole
zero on the removed selected
fibre. On the complementary
compact source choose a
Hermitian metric on \(M_8\).
The entire unscaled lift has
a uniform finite norm.
A smaller common nonzero
scalar makes that norm
strictly less than one.
The inverse lift on its
image expands. A fundamental
annular shell and its two
identified faces therefore
give a compact proper
Hausdorff quotient, uniformly
over \(\Delta\).

The cusp chart O45R.9 is
smooth over the source
parameter, and O45R.17
extends the line map there.
The selected toric charts
O45R.21-\/-24 are smooth
over that parameter.
The infinity charts
O45R.25-\/-31 and the
ordinary cap O45R.32
are smooth over it as well.
Their displayed inverses
and common source/line
overlaps glue these pieces
to a proper holomorphic
submersion over
\(\Delta\).

The topology is transported
from the original K10
source by an actual path
in the connected punctured
ordered cover. On finitely
many discs along that path
the same source equations,
point modules, selected
fan, infinity charts and
ordinary cap apply.
Their original period
changes are integral
basis changes with their
full source/line
coefficients. A smaller
common scalar gives
proper submersions
over the whole compact
path. A horizontal lift
of its tangent field,
patched on the compact
quotient by a partition
of unity, integrates to
an oriented fibre
diffeomorphism. Its
inverse is the reverse
flow. It carries the
actual K10 six-sphere
and its prescribed
fan to the boundary
of the current disc.

The same construction on
the filled disc gives a
smooth trivialization:
choose a connection,
average over the compact
physical circle, and
transport radially
from the centre.
Its inverse uses the
reverse radial paths.
All endpoints on a
boundary circle are
therefore identified
by a product whose
closing is the
identity. The original
analytic family has
this same closing,
because every source,
point, fan and cap
coefficient above is
single-valued on
\(\Delta\).
Its physical return
is isotopic to the
identity, hence has
class zero.

Finally take any scalar
admitted on the original
boundary family. The
constant-argument radial
interval to the smaller
common scalar stays in
the strict contraction
domain on that boundary.
Substitution of this
interval in the complete
periods, O45R.19,
O45R.24, O45R.31 and
O45R.32 gives an
actual proper annular
comparison family.
It has no scalar
winding and retains
the endpoint maps.
Horizontal flow on
that compact annulus,
with its reverse as
inverse, gives an
isotopy of the two
boundary returns.
Thus every admitted
fixed scalar gives
class zero on each
of the twelve marked
unselected-II
meridians. The other
source-cover
generators remain
outside this
calculation.

% End expanded verified cp36 input


\par\bigskip\noindent\textbf{Complete supporting proof 3: O45P.}\par


% Expanded verified cp36 input: notes/part3_os45_unselected_ii_o45p_proof.tex
\par

\medskip\noindent\textbf{Complete independent comparison for the twelve OS45 unselected-II meridians}

\par\nobreak

This proof retains the original source, section, point line, suspension,
ordinary logarithmic cap, selected branch order, and prescribed lattice-basis
fan of O45.1-\/-7. It verifies the complete author construction and the
independent O45R.1-\/-4 comparison. Its conclusion is the zero return of the
twelve specified source meridians. The other source generators retain their
separate obligations. The earlier numerical scalar-return sign in K10 is
not an input to this proof: the actual source, point line, local-system
maps, integral topology and ordinary cap of that construction are used.

\par

\medskip\noindent\textbf{O45P.1. Original source maps and the exact marked discs}

\par\nobreak

Write
\[
 R=t^2+a,\quad f=R^2+bt,\quad
 g=R^3+\frac32btR+\frac38b^2,\quad
 y^2=x^3-3fx+2g,\qquad G=(R,\sqrt3b/2).
 \tag{O45P.1}
\]
The chosen square root is fixed. The weighted arrow
\((t,x,y)=(sT,s^2X,s^3Y)\), \(a=s^2A,\ b=s^3B\), and its inverse
divide the equation by \(s^6\) and carry \(dx/y\) to \(s^{-1}dX/Y\).
We use its explicitly marked cube-root cover to put \(b=1\).
This keeps the deck arrow
\((u,v)\mapsto(\zeta_3u,v)\) and the original source map
\((a,t,x,y)\mapsto(\zeta_3a,\zeta_3^2t,\zeta_3x,y)\).

For clarity, rationality and the direction of the translation can be
checked on the original source before choosing any singular fibre.
Put \(Y=y/\sqrt3\) and use the two independent coordinates \(A,B\) with
\[
 v_0=3(AB-a),\quad
 t=\frac{B-A}{2}+\frac{b}{2v_0},\quad
 x=t^2+a+v_0,\quad Y=\frac{v_0(A+B)}2.
 \tag{O45P.2}
\]
The inverse is
\(v_0=x-t^2-a,\ A=Y/v_0-t+b/(2v_0),\
B=Y/v_0+t-b/(2v_0)\).
Substitution gives the original equation \(3Y^2=x^3-3fx+2g\).
Thus these are mutual rational source maps on the stated open sets;
the minimal regular compact source is the rational elliptic surface
used in K10.

The full finite discriminant and its discriminant in the base coordinate
are
\[
 D=48a(t^2+a)^2+8bt^3+72abt+9b^2,\quad
 \Delta=-27b^2D,\quad
 \operatorname{disc}_tD=331776b^2(256a^3-3b^2)^3.
 \tag{O45P.3}
\]
At a selected double root \(\xi\) put \(u=\xi-(t_c^2+a)\).
Solving the two original double-root equations gives the actual inverse
marked cover
\[
 a=-\frac u3-\frac{u^4}{9b^2},\qquad
 t_c=\frac b{2u}-\frac{u^2}{3b},\qquad
 \xi=\frac{b^2}{4u^2}+\frac u3.
 \tag{O45P.4}
\]
The inverse is precisely \(u=\xi-(t_c^2+a)\). In particular
\(f(t_c)=\xi^2,\ g(t_c)=\xi^3\), and
\[
 \frac{da}{du}=-\frac{4u^3+3b^2}{9b^2},\qquad
 D_t(t_c)=-\frac{2(4u^3+3b^2)^3}{27b^3u^2}.
 \tag{O45P.5}
\]
The original selected normalization and its inverse are
\[
 x=S^2-2\xi,\quad y=S(S^2-3\xi),\quad S=\frac y{x-\xi},
 \qquad
 Z=\frac{S-v/(2u)}{S+v/(2u)},\quad
 S=\frac v{2u}\frac{1+Z}{1-Z}.
 \tag{O45P.6}
\]
Here \(v^2=4u^3+3b^2\). Substitution of
\(S_G=-\sqrt3b/(2u)\) sends the original section to
\[
 P=\frac{\sqrt3b+v}{\sqrt3b-v}.
 \tag{O45P.7}
\]
Changing the ordered branch is the displayed source arrow
\(Z\mapsto Z^{-1}\). This proof uses the fixed ordered branch on each
disc and its continued fan; it does not use that arrow as a new closing.

Now put \(b=1\) and \(H(u)=16u^6+120u^3+243\). The two exact
factorizations are
\[
 256a(u)^3-3=-\frac{(4u^3+3)^2H(u)}{729},\qquad
 \operatorname{disc}_t\!\left(\frac{D}{t-t_c}\right)
        =-\frac{1024u^4H(u)^3}{6561}.
 \tag{O45P.8}
\]
The six roots of \(H\) are simple. Polynomial gcds with
\(u,\ u^3+3,\ 4u^3+3,\ 2187-16u^3(u^3+3)^3\) are one.
Thus each root has two distinct ordered values of \(v\), and gives one
of the twelve specified unselected-II disc centres.
All forty-eight deleted points are retained: respectively two over
\(u=0\), six over \(u^3=-3\), three with \(v=0\), twelve over \(H=0\),
twenty-four over the ordinary-cap polynomial, and the point at infinity.
The curve \(v^2=4u^3+3\) has four simple branch points over the
\(u\)-sphere; the Euler calculation gives genus one. Removing the
forty-eight points leaves a connected surface with free fundamental
group of rank \(2+48-1=49\). The degree-three map to \(v\) is unramified
on this open surface, and its Euler identity is
\(-48=3(2-18)\). These actual covering maps specify which twelve
meridians are filled.

\par

\medskip\noindent\textbf{O45P.2. Smooth original cusp source and both translation directions}

\par\nobreak

Fix one of the twelve points, and choose a closed parameter disc
\(\Delta\) avoiding all other deleted values. The functions in
O45P.4 and the chosen \(v\) are single-valued on this disc.
At its centre write \(c=-4a_*\). Then
\[
 c^3=-\frac34,\quad a_*=-c/4,\quad
 t_*=-c^2,\quad R_*=-c,\quad f_*=g_*=0.
 \tag{O45P.9}
\]
Direct differentiation of the full coefficients gives
\[
 \det\frac{\partial(f,g)}{\partial(a,t)}=\frac32c^2,\qquad
 \left.\partial_t(x^3-3fx+2g)\right|_{x=y=0}=\frac32c.
 \tag{O45P.10}
\]
Consequently the original source is locally the holomorphic graph
\(t=t(a,x,y)\). Its inverse chart is the original \((a,x,y)\)
projection, so its projection to \(\Delta\) is a submersion.
The centre is disjoint from the selected fibre and the ordinary cap:
\[
 t_c+16a^2=\frac{(2u^3+9)(4u^3+3)^2}{162u}\ne0,\qquad
 D(0)=153/16.
 \tag{O45P.11}
\]
The nonvanishing assertion is checked by the gcd with \(H\), so it
holds at every one of the twelve marked centres.

The full translation \(T=-G\) on the original source is
\[
 L=\frac{y+\sqrt3/2}{x-R},\qquad
 x_T=L^2-x-R,\qquad y_T=L(x-x_T)-y.
 \tag{O45P.12}
\]
At the image point its plus-\(G\) slope is
\((y_T-\sqrt3/2)/(x_T-R)=-L\). Applying the same chord formula
with this slope returns the original \(x,y\).
On O45P.2 the scaled slope is simply \(t+A\); substituting these
rational source coordinates verifies both composites and the full
source equation without introducing another source family.
At the cusp the denominators are \(c\ne0\). The point is fixed and
the derivatives in the actual graph chart are
\[
 DT=\begin{pmatrix}1&2g_G\\0&1\end{pmatrix},\qquad
 DT^{-1}=\begin{pmatrix}1&-2g_G\\0&1\end{pmatrix},
 \qquad g_G=-\frac{2c}{\sqrt3}.
 \tag{O45P.13}
\]
The derivative of \(t(a,x,y)\) in \(x,y\) is zero at this point, so
these are the derivatives on the full source graph, not on a
different transverse surface. The ordinary addition charts at \(G\)
and \(O\), and the simultaneous infinity charts below, extend the
same mutually inverse translation over the rest of the source.

\par

\medskip\noindent\textbf{O45P.3. The complete point-line coefficient and its two cusp frames}

\par\nobreak

Let \(M_8=\mathcal O([8G]-[O])\), with its canonical rational point
section \(\sigma_8\). The actual tangent/chord construction gives
\[
 f_2=\frac{\ell_G}{x-x(2G)},\qquad
 f_4=f_2^2\frac{\ell_{2G}}{x-x(4G)},\qquad
 f_8=f_4^2\frac{\ell_{4G}}{x-x(8G)},\qquad
 B_{8,1}=f_2^8\,\frac{f_8\circ T}{f_8}.
 \tag{O45P.14}
\]
Every tangent is evaluated at the indicated original section.
The horizontal divisor of \(f_m\) is
\(m[G]-[mG]-(m-1)[O]\). Pullback by \(T\) adds \(G\) to each
section in that divisor. It follows directly that
\[
 \operatorname{div}_{\rm hor}B_{8,1}
      =[G]+[8G]-[9G]-[O]
      =([8G]-[O])-T^*([8G]-[O]).
 \tag{O45P.15}
\]
Thus \(\sigma_8(p)\mapsto B_{8,1}(p)\sigma_8(Tp)\) is the
specified point-module arrow. Its inverse at \(p\) uses
\(B_{8,1}(p+G)^{-1}\); the inverse-source argument is part of the map.

The coefficient through the cuspidal fibre can be computed before
taking a constant. Put \(p=x/y\), so \(x=p^{-2},y=p^{-3}\);
the original \(G\) has additive coordinate \(g_G\), and \(T\) is
\(p\mapsto p-g_G\). The actual tangent at \(jG\) has slope
\(3/(2jg_G)\), and its doubled point is
\(((2jg_G)^{-2},(2jg_G)^{-3})\).
The three chronological tangent quotients give
\[
 f_m(p)=-\frac{m}{g_G^{m-1}}\,
       \frac{(p-g_G)^m}{(p-mg_G)p^{m-1}},
       \qquad m=2,4,8,
 \tag{O45P.16}
\]
and hence
\[
 B_{8,1}(p)=\frac{256}{g_G^8}
       \frac{(p-g_G)(p-8g_G)}{p(p-9g_G)}.
 \tag{O45P.17}
\]
There are two explicitly related frames. On the finite
normalization point/O chart use
\(r(p)=(p-8g_G)/p\) and \(e=\sigma_8/r\).
The arrow on the coefficients of \(e\) is
\[
 B_{8,1}(p)\frac{r(p-g_G)}{r(p)}
          =\frac{256}{g_G^8}.
 \tag{O45P.18}
\]
This equation includes the zero-section and translated section
directions in their actual Cartier modules.

At the cusp \(\sigma_8\) itself is a nonzero frame, since its
horizontal divisor avoids that point. Multiplying the numerator
and denominator of the dimensionless quotient in O45P.17 by
the exact opposite factor \(1+9g_G/p\) gives
\[
 B_{8,1}(x,y)=\frac{256}{g_G^8}
       \frac{1-73g_G^2x+72g_G^3y}{1-81g_G^2x}.
 \tag{O45P.19}
\]
Both numerator and denominator are units at \((x,y)=(0,0)\).
The overlap between this cusp frame and O45P.18 is precisely the
full \(r(p)\). No extension of \(r\) as an invertible cusp-ring
function is used. The reciprocal in the inverse-source chart
extends as well.

The full meromorphic Hom coefficient has no other finite vertical
divisor. Each finite fibre in the disc family is irreducible;
\(T\) preserves it, so \(f_8\circ T\) and \(f_8\) have equal
whole-fibre orders. The original \(G\)-tangent denominator is
the nonzero constant \(\sqrt3\), and its chord quotient \(f_2\)
has whole-fibre order zero. After O45P.15 cancels the actual
horizontal module divisors, the coefficient and its reciprocal
are meromorphic functions without divisorial poles on the smooth
total source. The normal local ring therefore extends both
holomorphically. The explicit nonzero value in O45P.19 also
excludes a parameter-divisor zero at the new II fibre.
This proves the actual total-family point-line map through the
coincident-fibre parameter, with all its module arguments retained.

\par

\medskip\noindent\textbf{O45P.4. The directed infinity Hom degree and every period unit}

\par\nobreak

The original infinity minimal coordinates are
\(z=1/t,\ X=z^2x,\ Y=z^3y\).
The two original source branches at \(G\) have orders three and
five. Indeed \(X_G=1+az^2,\ Y_G=(\sqrt3/2)z^3\), while
\(\sqrt{f_{\min}}=1+az^2+z^3/2+O(z^5)\); the marked branch
cancels the cubic term on exactly one side.
The product has the original discriminant order eight.
Thus \(G\) meets the marked component three of \(I_8\).

Here is an integral verification of the height antecedent used
for the Hom degree. The positive \(A_7\) matrix \(C\) has diagonal
two and adjacent entries minus one. It has determinant eight and
\((C^{-1})_{33}=15/8\). In the ordered actual source classes
\((O,F,C_1,\ldots,C_7,G)\), use
\[
 O^2=G^2=-1,\quad OF=GF=1,\quad GO=0,\quad
 (C_iC_j)=-C,\quad GC_3=1,
 \tag{O45P.20}
\]
with all other unspecified entries zero. The resulting ten-class
Gram determinant is \(-1\). These classes therefore fill the
unimodular rational-surface lattice, and their quotient by the
first nine classes is generated freely by the original \(G\).
Its height is \(2-15/8=1/8\).
For \(jG\) the actual component is \(3j\bmod8\); the height
intersection formula then gives
\[
 G.O=0,\qquad 7G.O=3,\qquad 8G.O=3,\qquad 9G.O=5.
 \tag{O45P.21}
\]
The directed \(-G\) Hom restricted to \(O\) has degree
\(5-0-3-1=1\). Equivalently, the original \(t_G\)-pullback
linearization has degree \(3+1-3+0=1\); these are the two
directions with their own source arguments.

The line \(M_8\) has component degree zero at infinity, since
\(8G\) and \(O\) meet its identity component. The full cycle
intersection matrix has kernel
\(\mathbb Z(1,\ldots,1)\). Its entire remaining vertical
Hom divisor is consequently exactly \(F_\infty\), with
coefficient one fixed by the preceding degree. Multiplication
by the original \(q=t-t_c\) changes it to the selected whole
zero. Thus the original suspension is literally
\[
 B_\lambda(p,v_M)=
       (Tp,\lambda C_*q B_{8,1}(p)v_M).
 \tag{O45P.22}
\]
The continuation in (O45P.22) is defined using the single
nonzero \(C_*\) obtained from the prescribed initial line map
\(L_{a_0}\) at a protected regular point \(s_0\), in the same
Cartier frames:
\[
 C_*=
 \frac{L_{a_0}(s_0)}
 {(t(s_0)-t_c(a_0))B_{8,1,a_0}(s_0)}.
\]
Equality of the divisors makes this ratio an invertible function
on the compact initial source, hence constant on that source.
This proves equality of the entire initial map with
\(C_*(t-t_c)B_{8,1}\). At later marked parameters (O45P.22),
with this same specified \(C_*\), defines the continuation
family whose coefficient maps and inverses extend over the disc.
The boundary identification in this proof concerns this defined
continuation of the initial source, line, cap and reference.

For an independently prescribed varying line map \(L_a\), the
exact comparison retains the analytic evaluation ratio
\[
 c(a)=\frac{L_a(s(a))}
 {(t(s(a))-t_c(a))B_{8,1,a}(s(a))}
\]
wherever these maps and the protected section are defined. Its value
\(c(a_0)=C_*\) alone proves neither constancy nor extension
through a missing parameter. That further family comparison
is not asserted here.

The marked infinity Tate quantities are the full functions
\[
 Q_\infty=z^8\kappa_\infty,\quad
 \kappa_\infty(0)=-a/2304,\qquad
 P_\infty=z^3\pi,\quad \pi(0)=-1/24.
 \tag{O45P.23}
\]
Their constants follow from the inverse \(j\)-coefficient and
the original normalized nodal parametrization
\[
 X=1+\frac{12Z}{(1-Z)^2},\qquad
 \frac Y{\sqrt3}=\frac{12Z(Z+1)}{(Z-1)^3}.
 \tag{O45P.24}
\]
In particular \(b_0=P_\infty^8/Q_\infty^3\) extends and has value
\(-1/(9a^3)\). Substitution in O45P.24 gives the full
identity-component eighth-point values
\[
 X_{8,0}=\frac{81a^6-90a^3+1}{(9a^3+1)^2},\qquad
 \frac{Y_{8,0}}{\sqrt3}
      =\frac{108a^3(9a^3-1)}{(9a^3+1)^3}.
 \tag{O45P.25}
\]
These are also obtained by the three original tangent doublings.
At the II centres \(a^3=3/256\), so \(b_0=-256/27\ne1\).
The point and zero module charts therefore remain separated in
the infinity collar.

For \(f_b(Z)=\Theta_Q(Z/b)/\Theta_Q(Z)\), three actual theta
shifts give
\[
 f_{Q^3b_0}(Z)=(-Z/b_0)^3Q^{-6}f_{b_0}(Z).
 \tag{O45P.26}
\]
The inverse of this entire coefficient changes the original
point coordinate to its identity-component coordinate.
Under \(Z\mapsto P_\infty^{-1}Z\) it changes the line
multiplier by \(P_\infty^3\); under \(Z\mapsto Q_\infty Z\)
it changes \(P_\infty^8\) to \(b_0\).
Thus the full \(Q_\infty^{-6}\) divisor and every base unit
remain in the frame arrow. The valuation columns in this frame
are exactly \((8,0)\) and \((-3,0)\).
Their actual period words and inverses are
\[
 H=h^{-1}B^{-3},\quad K=h^3B^8,\qquad
 h=H^8K^3,\quad B=H^{-3}K^{-1}.
 \tag{O45P.27}
\]
The integer matrix
\(\left(\begin{smallmatrix}-1&3\\-3&8\end{smallmatrix}\right)\)
has determinant one and sends the original source valuations to
\((1,0)\). The second multiplier of \(K\) is \(\lambda^8\)
times its entire extending unit.
The source charts \(X_i=Z/z^i,\ Y_i=z^{i+1}/Z\) have
\(X_iY_i=z\); the complete \(Q_\infty/z^8\) and
\(P_\infty/z^3\) coefficients give their period and translation
overlaps. The identity on the actual point module, combined
with O45P.27, is an isomorphism to the rank-one completed
quotient. Its inverse uses the displayed inverse words and
reciprocal units. These maps include all eight source components.

\par

\medskip\noindent\textbf{O45P.5. The selected collar, prescribed fan and original cap}

\par\nobreak

The selected node stays simple on \(\Delta\). Define
\(\xi(q,u)=\sqrt{f(t_c+q,u)}\) with its marked value in O45P.4,
and put
\[
 h_s(q,u)=\frac{2(g(t_c+q,u)-\xi(q,u)^3)}q,\qquad
 h_s(0,u)=-2u^4.
 \tag{O45P.28}
\]
This is an actual extending unit. The full source chart is
\[
 X=y-(x-\xi)\sqrt{x+2\xi},\qquad
 Y=h_s^{-1}\bigl(y+(x-\xi)\sqrt{x+2\xi}\bigr),\qquad XY=q.
 \tag{O45P.29}
\]
Its inverse sets \(q=XY,\ y=(X+h_sY)/2\) and solves
\(2(x-\xi)\sqrt{x+2\xi}=h_sY-X\). The derivative
\(2\sqrt{3\xi}\) is nonzero, so this is the inverse on the
whole node collar over the source parameter disc.

Keep the full inverse Tate coordinate
\(Q=q\kappa(q,u)\) and \(P=Z(G)\), with their original
reference scales. The actual point-module arrow and periods are
\[
 V_{\rm pt}=f_{P^8}(Z)v_M,\quad
 h:(Z,V_{\rm pt})\mapsto(QZ,P^8V_{\rm pt}),\quad
 B:(Z,V_{\rm pt})\mapsto(P^{-1}Z,\lambda q\beta V_{\rm pt}),
 \tag{O45P.30}
\]
\[
 \beta=C_*B_{8,1}(p)
       \frac{f_{P^8}(P^{-1}Z)}{f_{P^8}(Z)}.
 \tag{O45P.31}
\]
The complete horizontal divisors cancel. On every regular
source fibre this is independent of \(p\); evaluation at a
protected regular point gives the full analytic base function.
The module maps and their reciprocals just proved show that
it is an extending unit. Its value and logarithm are never
replaced by their central limits.

For a prescribed integral basis, each source/line transition
therefore acts on the affine character algebra by its original
monomial times the complete coefficient in O45P.30-\/-31.
For example its two vertex maps, with actual coefficient
characters \(A,D\), are
\[
 (a,b,c)\mapsto((AD)^{-1}a,Ab,Dc),\qquad
 (u,v,w)\mapsto(ADu,D^{-1}v,A^{-1}w).
 \tag{O45P.32}
\]
They preserve the original height product \(q\).
Their reciprocals are the inverse maps, and the same
character evaluation carries every cone of the given fan.
Since \(v\) and all coefficient functions are single-valued
on \(\Delta\), this is the prescribed fan at the centre,
with its unchanged branch order and period arrows.

The original cap disc at \(t=0\) is smooth throughout
\(\Delta\). Continue its complete source and line period
vector \(\gamma(t,u,\lambda)\), with all initial logarithms
and annular integers. Its directed map and inverse are
\[
 \Psi(\zeta,t,u)=
   \left(\zeta+\frac{\log t}{2\pi i}\gamma(t,u,\lambda),t,u\right),
 \qquad
 \Psi^{-1}(\zeta,t,u)=
   \left(\zeta-\frac{\log t}{2\pi i}\gamma(t,u,\lambda),t,u\right).
 \tag{O45P.33}
\]
A logarithm branch change adds precisely the whole original
\(\gamma\)-period in every coordinate. Thus this is the
single-valued map on the original torus quotient.
Its cap meridian maps to the base meridian plus \(\gamma\),
retaining its multiplicity-one nonzero clutch. Its line
coordinate includes the original \(\log\lambda\).
There is no parameter-deck residue on this filled disc:
all the displayed source, point, fan and cap maps return
to themselves by their actual single-valued coefficients.

\par

\medskip\noindent\textbf{O45P.6. Reference sphere, compact family and identity closing}

\par\nobreak

We make the topology antecedent for every prescribed fan
explicit. At the original K10 \(a=b=1\) source, K10\textquotesingle s
integral local-system and cap maps give its standard
six-sphere. The full five source matrices in that proof
have product one; the homology images of \(T_i^{-t}-I\)
generate a primitive rank-three sublattice. Its quotient
is measured by \(\psi\), and the actual cap meridian in
O45P.33 has \(\psi\)-value one. The complete primitive
Gysin/Leray calculation in K10 leaves only \(H^0,H^6\).
This establishes the stated reference sphere from its
actual cap. Its earlier scalar-return value is unused.

For each prescribed lattice-basis fan \(P\), the finite
canonical word of K71/K73 starts at this actual source.
At each adjacent change the common blowup and its two
ruling collapses, with the primitive integral Gysin
and van Kampen maps of F35/K23, preserve simple
connectivity and the homology-sphere groups. The
complete characteristic/period arrows in K71 keep
the same source, line, cap and all full coefficients.
Applying these finitely many actual maps gives the
reference sphere on \(P\); the six-dimensional
homotopy sphere is the standard sphere. This supplies
the fan-specific starting topology before any source
path is taken.

The open ordered cover is connected. Choose a path
from that K10 marking to the chosen boundary point
of \(\Delta\), and cover its compact image by finitely
many small discs avoiding the listed deleted values.
The original source, line, selected height, fixed
fan, infinity charts and cap maps O45P.22-\/-33 apply
on every disc. Their overlaps are the actual source
identities, point modules and integral changes of
period basis. Continue the original cap vector and
all reference integers along the path.

On the compact source outside the selected tube,
choose a smooth Hermitian metric on the actual
point line. The full unscaled lift and its base
source translation have a uniform finite norm over
the compact parameter discs, including the II centre.
Choose a common nonzero scalar whose product with
that bound is less than one. The suspension is
strictly contracting and its inverse expands.
This gives a compact annular fundamental domain.
Only finitely many iterates can meet a fixed compact
subannulus, by the strict norm inequality. Thus the
quotient is proper and Hausdorff, uniformly in the
parameter. A physical-circle equality with an iterate
of \(B\) forces that iterate to be zero by the same
inequality; the quotient line action is the original
physical one.

Inside the selected tube use the actual finite
period representatives of O45P.29-\/-32; at infinity
use the entire cycle and O45P.27; at the ordinary
cap use O45P.33. Their overlaps are the displayed
source/module/period isomorphisms with reciprocal
inverses. The compact collar pieces cover the
fundamental domain and glue to a proper family.
The source projection is a submersion at the cusp
by O45P.10, at the selected fan by its full toric
coordinates, and at the other source/cap pieces
by their continued smooth charts. Therefore the
completed family is a smooth proper holomorphic
submersion over the whole \(\Delta\).

The same construction along the reference path
carries the reference sphere to the disc boundary.
To describe its inverse explicitly, choose local
lifts of the path tangent vector in these charts,
transport them by the displayed differential
overlaps, and patch by a partition of unity on
the compact total space. Their vertical differences
are tangent to the compact fibre. The resulting
horizontal flow exists along the whole path; its
inverse is the reverse flow. It preserves the
orientation induced by the complex fibres.

On the disc itself choose such a connection and
average it over the physical circle. Radial
transport from the centre, with inverse transport
along the reversed radial segment, trivializes
the completed family smoothly over the disc.
On its boundary the analytic closing is exactly
the identity because every source, line, fan,
cap and reference map is single-valued on
\(\Delta\). The boundary return is therefore
isotopic to the identity through these actual
fibre maps. Its physical return class is zero.

Finally take any fixed scalar admitted on the
original boundary family. Decrease its modulus
to the common disc-contraction scalar along the
constant-argument interval. The strict
contraction inequality remains valid on the
boundary throughout this interval. Substitute
the full \(\log\lambda\) in O45P.22,
O45P.27, O45P.30-\/-31 and O45P.33.
This is a proper annular family with the
same endpoint maps and no added scalar turn.
Horizontal flow relative to the endpoint
collars and its reverse give an isotopy of
the two boundary returns. Hence the zero
return holds for every admitted fixed scalar.
This constructs the claimed filling for
each of the twelve marked meridians, with
the original source, cap and prescribed fan.

\par

\medskip\noindent\textbf{O45P.7. Independent exact execution and scope}

\par\nobreak

The executable
\path{code/part3/os45_unselected_ii_blind_peer.py}
imports no author or root certificate.
It reconstructs the translation and its inverse
on the rational source chart O45P.2; checks
the whole discriminant, marked cover and
each deleted locus; derives the actual cusp
graph derivatives and chronological tangent
functions; verifies both complete Cartier
frames and their reciprocal-source arrows;
and computes the integral ten-class Gram,
directed Hom degrees, whole-cycle kernel,
theta reindex and primitive period inverse.
It also checks the full character and inverse
maps on all 104 integral bases with entries
between minus two and two, the complete
ordinary-cap branch vector, and the original
primitive van Kampen/cap relation lattice.
The full analytic proof above supplies the
unbounded fan and compact-disc conclusions.

Every read reusable source is scanned for
control bytes other than ordinary newlines
and tabs. The output locks the original
author source, corrected root proof, retained
K10 source, this full peer proof and the
actual fan antecedents. No previous source
or return derivation is overwritten.
The conclusion is confined to the twelve
unselected-II source meridians. The selected
cusp, infinity \(I_9\), ordinary-cap collision
and remaining handle/source generators
retain their separate exact obligations.

% End expanded verified cp36 input



\par\medskip\noindent\textbf{Completed construction and audit dependencies.}\par
The original K10 source construction is retained through its literal
equation, point line, selected whole zero, original ordinary cap and
topological family. Its source is
\path{notes/part3_os45.md}; the generic source catalogue is
\path{notes/part3_generic_source_models.md}. The present proof carries
their actual source and line morphisms and every retained cap and
period coefficient. The scalar-return calculation is the completed
P3K.73 result; no historical scalar value enters the disc transfer.

The primary executable is
\path{code/part3_fourfold/os45_marked_source_cover.py}.
The independent root executable is
\path{code/part3_root/os45_unselected_ii_peer.py}.
The further independent executable is
\path{code/part3/os45_unselected_ii_blind_peer.py}.
Their 76, 88 and 225 exact checks retain the original source and
all full coefficient maps. The complete independent supporting texts
and all executed checks are
locked in \path{records/k116_export_manifest.json}. The two original
Cartier frames at the unselected cusp are retained separately with
their exact overlap; the full infinity theta shift, primitive period
inverse and ordinary-cap logarithm appear in each geometric proof.
The compact disc and scalar-annular flows, with their reverse flows,
give the claimed physical return directly.


% End expanded verified cp36 input

\end{proof}


\subsection{The complete OS45 selected-cusp return on every fan}
\begin{theorem}[Source result P3K.117 under continuation review]\label{night:P3K.117}
\StatementRecord{night:P3K.117}{open}{Historical events: GPT-6 Astra (OpenAI), ChatGPT Pro mode, effort Pro, S6\_06\_BLIND, 16--18 November 2026: targeted proof check. Also checked in this writing. Continuation mathematical and source checks: GPT-6 Astra (OpenAI), ChatGPT Pro mode, effort Pro, S6\_10\_PAPERS continuation, 19 November 2026: computation or formal certificate; used and rechecked in this run. Continuation metadata events: GPT-6 Astra (OpenAI), ChatGPT Pro mode, effort Pro, S6\_10\_PAPERS continuation, 19 November 2026: used and rechecked in this run; explicitly scoped metadata reconciliation; metadata reconciliation only; no new mathematical check. GPT-6 Astra (OpenAI), ChatGPT Pro mode, effort Pro, S6\_10\_PAPERS continuation, 19 November 2026: used and rechecked in this run; source-row intake; metadata reconciliation only; no new mathematical check. GPT-6 Astra (OpenAI), ChatGPT Pro mode, effort Pro, S6\_10\_PAPERS continuation, 19 November 2026: used and rechecked in this run; receiving source placement; metadata reconciliation only; no new mathematical check. \emph{Pre-continuation recorded proof scope.} Historical event: GPT-6 Astra (OpenAI), ChatGPT Pro mode, effort Pro, S6\_06\_BLIND, 16--18 November 2026: targeted proof check, recorded producing-run re-derivation. Also checked in this writing by GPT-6 Astra (OpenAI), ChatGPT Pro mode, effort Pro, S6\_10\_PAPERS continuation, 19 November 2026: computation or formal certificate for the explicitly listed algebraic substeps in the event ledger; full continuation rederivation and renewed credit comparison remain pending.}{the original source, point line, roots, characters, reference, selected fan and scalar range stated here}
Retain the original marked OS45 source
\[
 R=t^2+a,\quad f=R^2+t,\quad
 g=R^3+\tfrac32tR+\tfrac38,\qquad
 y^2=x^3-3fx+2g,\qquad G=(R,\sqrt3/2),
\]
with the genuine point line \(M_8=\mathcal O([8G]-[O])\),
its actual translation lift by \(-G\), the original order-one,
positive-norm-one ordinary cap at \(t=0\), and its complete
prescribed annular reference. On the ordered marked cover write
\[
 a=-u/3-u^4/9,\qquad t_c=1/(2u)-u^2/3,\qquad
 v^2=4u^3+3,\qquad q=t-t_c .
\]
For each of the three points \(u_0^3=-3/4\), take one positive
small \(v\)-circle, with the scalar fixed at any value admitted by
the original boundary family. This ordered loop has four positive
turns in the original \(a\)-plane and closes by the actual
original source, line and cap identity.

For every integral unimodular selected fan \(T(P)\), this family
has the physical return determined by the finite algorithm below.
The reference manifold is \(S^6\), carried to the prescribed fan
by the actual common-blowup and two-ruling maps. At
\(T_0=T(e_1,e_2)\), the characteristic and native pairs are
\[
 (k_M,\ell_M)=(-6,0),\qquad(\mu,\nu)=(0,-6),\qquad
 A_M=-1\pmod {28}.
\]
For every signed integer \(j\), on \(T_j=T(e_1+je_2,e_2)\),
\[
 \begin{aligned}
 (k_M,\ell_M)_j&=(-6-18j(j-1),0),\\
 (\mu_j,\nu_j)&=(0,-6-18j(j-1)),\\
 A_M(T_j)&=-1-3j(j-1)\pmod {28}.
 \end{aligned}
\]
The native coordinates in this held Cox product satisfy
\(k_M=\mu+\nu,\ell_M=4\mu\).

For a basis \(P\), let \(R=P^{-1}\) and retain the full periods
\[
 h=(\kappa qZ,\alpha V),\qquad B=(gZ,\eta qV),\qquad
 g=P_G^{-1},\quad\alpha=P_G^8,\quad\eta=\lambda\beta .
\]
Their ordered-parameter windings are \((-6,0,0,0)\). Put
\[
 \begin{gathered}
 A_{ik}=(\kappa^{P_{1k}}g^{P_{2k}})^{R_{i1}}
        (\alpha^{P_{1k}}\eta^{P_{2k}})^{R_{i2}},\qquad
 B_P=-6P^{-1}e_1e_1^tP,\\
 (x,y)^t=P^{-1}e_2,\quad
 w=(x,y,-x-y),\quad d=(-x-2y,-2x-y,x-y),\\
 l=(B_{22}-B_{21},B_{11}-B_{12},B_{22}),\\
 I_i=-3w_i+d_i(l_i+3),\qquad J_i=-3w_i-d_i(l_i+3).
 \end{gathered}
\]
The complete vertical ordinary seed is
\[
 \begin{gathered}
 V=\begin{pmatrix}
 18j(j-1)&6-12j&12j-6\\
 6-12j&0&0\\12j-6&0&0
 \end{pmatrix},\\
 \delta=(-12-18j(j-1),12j-12,-12j),\qquad
 s=-1-3j(j-1).
 \end{gathered}
\]
Use the actual identical-fan column swaps and triangle translations
to reach \(T_j\), or the positive representative
\(P=\left(\begin{smallmatrix}p&-r\\q&s\end{smallmatrix}\right)\)
with \(p,r>0\), \(ps+rq=1\). Its parent is \(PQ_0^{-1}\)
for \(r\ge p\), and \(PQ_1^{-1}\) for \(p>r\), where
\[
 Q_0=\begin{pmatrix}1&-1\\0&1\end{pmatrix},\qquad
 Q_1=\begin{pmatrix}1&0\\-1&1\end{pmatrix}.
\]
The positive sum decreases, so this reaches a vertical seed.
Reverse the finite word. At each forward step through old edge
\(i=0\) or \(1\), write \(a=V_{ii},b=\delta_i\); then
\[
 s'=s+(b-a)/2,\qquad V'_{22}=a,\qquad\delta'_2=2a-b .
\]
The full target prism is
\[
 V'(w_ie_i-w_0e_0)=I_ie_i-I_0e_0,\qquad
 w_i\delta'_i-w_0\delta'_0=J_i-J_0,\qquad i=1,2.
\]
Its nonzero-weight target solution is uniquely
\[
 \begin{gathered}
 V'_{rs}=\frac{I_r}{w_r}\delta_{rs}+\frac{C}{w_rw_s},
 \qquad\delta'_r=\frac{J_r+E}{w_r},\\
 C=w_2^2a-w_2I_2,\qquad E=w_2(2a-b)-J_2 .
 \end{gathered}
\]
The final \(s\bmod28\) is the physical return. All divisions
occur in the actual integral prism and native coordinates before
reduction modulo 28. Both determinant signs are carried by the
explicit identical absolute fan maps.

The proof includes the original height and full inverse source
coordinates, every point/O and tensor Cartier factor, both
Chern caps, the actual physical inverse-clutch map, the signed
vertical relative-line and tangent comparisons, and all complete
coefficient and inverse-height maps at each fan step. The scope
is this ordered loop with the stated original cap and reference.
\end{theorem}
\begin{proof}[Supplied source argument; continuation review in progress]

% Expanded verified cp36 input: notes/part3_os45_selected_cusp_proof.tex
\par\bigskip\noindent\textbf{Full supporting proof 1: Original source, Miller maps and reflection.}\par\nobreak

% Expanded verified cp36 input: notes/part3_os45_selected_cusp_local_proof.tex
This separate calculation uses the literal source and point-line maps
of O45.1--5. It records the selected-cusp source collar. The compact
trace, both normal-cap comparison maps, and its physical return have
not yet been completed here. The unselected-II disc theorem is the
separate completed K116 calculation.

\par

\medskip\noindent\textbf{SC45.1. The original parameter and the exact source height}

\par\nobreak

Fix \(u_0^3=-3/4\). On a disc about \(v=0\), let \(u(v)\) be the
unique branch of \(4u^3+3=v^2\) with value \(u_0\). Retain
\[
 a=-\frac u3-\frac{u^4}{9},\qquad
 t_c=\frac1{2u}-\frac{u^2}{3},\qquad
 \xi=\frac1{4u^2}+\frac u3,\qquad
 c=\frac v{2u}.
 \tag{SC45.1}
\]
Then \(3\xi=c^2\). The map from \(v\) to \(c\) has derivative
\(1/(2u_0)\ne0\). Its inverse is the unique branch
\[
 u(c)^3-u(c)^2c^2+\frac34=0,\qquad
 u(0)=u_0,\qquad v=2u(c)c.
 \tag{SC45.2}
\]
Thus \(u,a,t_c,\xi\) are even functions of either ordered parameter,
with these explicit reciprocal parameter maps. The original source
parameter has the full fourth-order expression
\[
 a-a_0
 =-\frac{v^4\{u^2+2u_0u+3u_0^2\}}
 {144\{u^2+u_0u+u_0^2\}^2},\qquad
 a_0=-u_0/4.
 \tag{SC45.3}
\]
Indeed \(a-a_0=-(u-u_0)^2(u^2+2u_0u+3u_0^2)/9\) and
\(u-u_0=v^2/[4(u^2+u_0u+u_0^2)]\). The leading coefficient in
(SC45.3) is \(-1/(216u_0^2)\). A positive ordered-parameter meridian
therefore makes four positive turns in the original \(a\)-parameter.
This is a statement about this displayed parameter map.

Put \(q=t-t_c\), \(X=x-\xi\), \(R_c=\xi-u\), and \(d=2t_c+q\).
The entire original cubic becomes
\[
 y^2=X^2(X+c^2)+q\,H(q,X,c),
 \tag{SC45.4}
\]
where the actual polynomial in \(q,X\), with the full analytic
coefficients of (SC45.1), is
\[
\begin{split}
H={}&d\{-6R_c(u+X)+3t_c\}
   +3qd^2(2R_c-\xi-X)+2q^2d^3\\
 &\hspace{12mm}-3(u+X)+3qd .
\end{split}
\tag{SC45.5}
\]
There is no omitted source coefficient in (SC45.5):
\(R(t_c+q)=R_c+qd\), and substitution in the original \(f,g\)
gives (SC45.4) directly. At \(q=X=0\),
\[
 H(0,0,c)=-2u(c)^4,\qquad H(0,0,0)=3u_0/2\ne0.
 \tag{SC45.6}
\]
The implicit function theorem therefore gives the unique original
height \(q=q(X,y,c)\) from (SC45.4). Conversely \(x=X+\xi(c)\),
\(t=t_c(c)+q(X,y,c)\), \(a=a(c)\), and the unchanged \(y\)
give its full inverse source chart. These maps retain the original
height even at the selected cusp.

\par

\medskip\noindent\textbf{SC45.2. Ordered normalization, source translation and normal arrows}

\par\nobreak

The selected normalization, including both marked branches, is
\[
 X=S^2-c^2,\qquad y=S(S^2-c^2),\qquad
 S_G=-\frac{\sqrt3}{2u(c)},\qquad
 Z=\frac{S-c}{S+c},\qquad
 P=\frac{\sqrt3+2u(c)c}{\sqrt3-2u(c)c}.
 \tag{SC45.7}
\]
The actual minus-\(G\) source arrow on this normalization is
\[
 S\longmapsto S_T=\frac{S_GS-c^2}{S_G-S},
 \qquad
 S=\frac{S_GS_T+c^2}{S_G+S_T}.
 \tag{SC45.8}
\]
Substitution gives \(Z(S_T)=P^{-1}Z(S)\). Both denominators
are units near \(S=c=0\); (SC45.8) fixes the branches \(S=\pm c\).
The full off-normalization arrow is the original chord map
(O45.17), evaluated in the inverse source chart of SC45.1.
Its plus-\(G\) inverse uses the actual target source point.

Write \(x_-=S-c\), \(x_+=S+c\), and
\[
 n_f=(3S^2-c^2)dX-2S\,dy+2c(S^2-c^2)\,dc .
 \tag{SC45.9}
\]
Differentiating the full original height (SC45.4), along its
normalization, gives
\[
 dq=-\frac{x_-x_+}{H(0,x_-x_+,c)}n_f .
 \tag{SC45.10}
\]
On the actual native chart \(q=ZAB\), and its opposite chart
\(Z'=Z^{-1}\), \(A'=ZA\), \(B'=ZB\), this becomes
\[
 d(AB)=-\frac{x_+^2}{H(0,x_-x_+,c)}n_f,\qquad
 d(A'B')=-\frac{x_-^2}{H(0,x_-x_+,c)}n_f .
 \tag{SC45.11}
\]
The common denominator is the extending original source unit from
(SC45.5). Thus these are the full normal coefficients of this source
chart. A later comparison must carry that coefficient and its inverse.

\par

\medskip\noindent\textbf{SC45.3. The complete nodal point-line coefficient}

\par\nobreak

For \(m=2,4,8\), put
\[
 S_m=c\,\frac{1+P^m}{1-P^m}.
 \tag{SC45.12}
\]
The apparent quotient at \(c=0\) has the analytic limit
\(S_m(0)=S_G(0)/m\). The original tangent-and-doubling functions
of (O45.18), restricted to the nodal normalization, are exactly
\[
 f_m(S)=\frac{(S-S_G)^m}{S-S_m}.
 \tag{SC45.13}
\]
The coefficient is fixed by the original tangent quotient at \(O\):
its leading term in \(S\) is \(S^{m-1}\). For the initial tangent
the slope is \((3S_G^2-c^2)/(2S_G)\), the double point has
parameter \((S_G^2+c^2)/(2S_G)\), and its literal quotient is
(SC45.13) for \(m=2\). The two following doublings give \(m=4,8\)
with the same leading coefficient. In particular,
\[
 f_2(c)=f_2(-c)=-2S_G=\frac{\sqrt3}{u},\qquad
 B_{8,1}|_{\text{either node}}=\frac{81}{u^8}.
 \tag{SC45.14}
\]
For the second equality the actual translation (SC45.8) fixes each
branch and \(f_8\) is nonzero there, so the ratio
\(f_8(S_T)/f_8(S)\) is literally one at the branch.

At \(Q=0\), the complete theta point function restricts to
\[
 f_{P^m}(Z)=\frac{S-S_m}{c-S_m}.
 \tag{SC45.15}
\]
For the genuine point coefficient and the tensor coefficient this
specifies the direction
\[
 V_{\mathrm{pt}}=f_{P^8}(Z)v_M,\qquad
 V_{\mathrm{tensor}}=f_P(Z)^8 f_8(S)^{-1}v_M
       =C_8(c)V_{\mathrm{pt}},\qquad
 C_8(c)=\frac{c-S_8}{(c-S_G)^8}.
 \tag{SC45.16}
\]
Both \(c-S_8\) and \(c-S_G\) are analytic units at the selected cusp;
their central values are \(\sqrt3/(16u_0)\) and
\(\sqrt3/(2u_0)\), respectively. The inverse arrow multiplies by
\(C_8^{-1}\).

The meromorphic section \(\sigma_8\) of divisor \([8G]-[O]\)
is trivialized on the entire normalization by
\(\sigma_8\mapsto S-S_8\). Its zero and pole cancel those of
the actual point and \(O\) Cartier modules. This is a line
isomorphism through both sections. The normalization coefficient
in (SC45.16) and this genuine module map are retained separately.

\par

\medskip\noindent\textbf{SC45.4. Full height winding and the directed reflection}

\par\nobreak

For every nonzero ordered parameter, use the full inverse marked
Tate uniformization of O45.5 with the original height \(q\).
The original \(j\)-function gives the exact coefficient
\[
 \kappa(0,v)=\left.\frac{Q}{q}\right|_{q=0}
            =\frac{2u(v)^{10}}{v^6}.
 \tag{SC45.17}
\]
To derive it directly, use
\(g^2-f^3=D/64\),
\(D_t(t_c)=-2v^6/(27u^2)\), and
\(f(t_c)^3=\xi^6=v^{12}/(12^6u^{12})\) in
\(j=1728f^3/(f^3-g^2)\), with \(Qj(Q)\) tending to one.
The winding is therefore \(-6\).
This Tate chart is taken on the actual shrinking selected collar;
it does not assert that the Tate coordinates extend over a fixed
\((q,v)\)-bidisc through \(v=0\). The full source charts
(SC45.4)--(SC45.11) give the retained chart at that point.

The complete genuine point-period unit of O45.27 has central
node coefficient
\[
 \beta(0,v)=C_*\,81/u(v)^8 .
 \tag{SC45.18}
\]
Indeed the node is fixed by the source translation, and its actual
point/O module is disjoint from the branch. This formula retains
the original scalar \(C_*\).

Choose the exterior, ordinary cap, and original reference first
over the \(\zeta=v^2\) disc, then pull back by \(v^2=\zeta\).
The original coefficients, marked section, selected height,
and every original rational point-line factor are holomorphic
there. On the full smooth selected collar, reversal of the
ordered branches is consequently the actual identity of the
source in the two opposite uniformizations:
\[
 Q(-v)=Q(v),\qquad P(-v)=P(v)^{-1},\qquad
 f_{P^{-8}}(Z^{-1})=P^8f_{P^8}(Z).
 \tag{SC45.19}
\]
The last identity follows from the full product identity
\(\Theta_Q(z^{-1})=-z^{-1}\Theta_Q(z)\); it retains every
higher \(Q\)-coefficient. Thus the point-line reflection and
the subsequent actual target source-period map are
\[
 I:(Z,V_{\mathrm{pt}})\longmapsto(Z^{-1},P^8V_{\mathrm{pt}}),
 \qquad
 h_{\mathrm{target}}I:
 (Z,V_{\mathrm{pt}})\longmapsto(Q/Z,V_{\mathrm{pt}}).
 \tag{SC45.20}
\]
Here the target line multiplier is \(P^{-8}\), so its cancellation
is literal. The integral height-one source-period reindex remains
part of the second map.

The same parameter reversal fixes \(S_G,S_8\), while
\[
 (c-S_8)(-v)=P^{-8}(c-S_8)(v),\qquad
 (c-S_G)(-v)=P^{-1}(c-S_G)(v),\qquad C_8(-v)=C_8(v).
 \tag{SC45.21}
\]
These are the exact tensor and point-frame compatibility equations.
They establish the local collar and reflection maps recorded here.
The compact trace, its outer product, both normal caps and the
physical characteristic calculation remain the next proof steps.

% End expanded verified cp36 input

\par\bigskip\noindent\textbf{Full supporting proof 2: Complete source, orbit, cap, vertical and all-fan construction.}\par\nobreak

% Expanded verified cp36 input: notes/part3_os45_selected_cusp_completion_proof.tex
\par

\medskip\noindent\textbf{Completed proof and retained submission history.}

\par

The original pending-review notice is retained below. The complete C45.1-\/-8 proof has now passed both full independent geometric audits reproduced afterward. Its primary readiness and all preceding attempts remain preserved.

This completes the local SC45.1--4 calculation on the literal OS45
source. Fix one of the three points \(u_0^3=-3/4\) on the ordered
cover, retain \(M_8=\mathcal O([8G]-[O])\), translation by \(-G\),
the original order-one ordinary cap at \(t=0\), and the entire
prescribed source, line and annular reference. The parameter is one
positive small \(v\)-circle with \(4u^3+3=v^2\). The scalar is fixed.
The proof first uses a common strict-contraction scalar and then gives
the actual annular transport to every scalar admitted on the original
boundary family. It concerns this ordered loop, with its four turns
in the original \(a\)-plane. No shorter endpoint is assigned here.

All SC45 formulas are retained. In particular the original height is
\(q=t-t_c\), the source chart is
\(y^2=X^2(X+c^2)+qH(q,X,c)\), and
\(c=v/(2u)\). The full polynomial \(H\) is SC45.5, not only its
nonzero central value. The local source and point-line proof remains
in \path{notes/part3_os45_selected_cusp_working.md}; the present
text supplies the compact trace, its full normal arrows, both cap
framings, physical transfer and the fanwise products.

\par

\medskip\noindent\textbf{C45.1. The source disc, critical values and complete outside marking}

\par\nobreak

Put \(H_0=H(0,0,0)=3u_0/2\). A second critical point of the
original height, besides the selected \(X=y=q=0\), has
\(X=c^2R(c),y=0,q=c^6J(c)\), where the exact equations are
\[
\begin{split}
 R^2(R+1)+JH(c^6J,c^2R,c)&=0,\\
 3R^2+2R+c^2JH_X(c^6J,c^2R,c)&=0,\\
 R(0)=-2/3,\qquad J(0)&=-4/(27H_0).
\end{split}
\tag{C45.1}
\]
The Jacobian in \((R,J)\) at the last row is
\(\left(\begin{smallmatrix}0&H_0\\-2&0\end{smallmatrix}\right)\),
with determinant \(2H_0\ne0\). The implicit solutions are even
in \(c\), because every coefficient of the original source is even.
Their inverse is the displayed original height, with no base factor
removed. Hence the ordered critical-value difference makes six
positive turns on this loop. In the positive half-twist convention
the two-value braid is \(\sigma^{12}\).

The actual weighted source chart is
\[
 X=c^2\widetilde X,\quad y=c^3\widetilde Y,\quad q=c^6Q,
 \qquad
 \widetilde Y^2=\widetilde X^2(\widetilde X+1)
              +QH(c^6Q,c^2\widetilde X,c).
\tag{C45.2}
\]
Its exceptional pencil is
\([Y^2-X^2(X+R^2):H_0R^6]\) on \(\mathbf P(2,3,1)\).
The six successive base-ideal charts \(z=R^iz_i\),
\(1\le i\le6\), and the order-two and order-three quotient
charts are exactly the F56 pencil charts, with the displayed
constant \(H_0\). On the overlap (C45.2), the full \(H\) stays
in the equation. Its quotient by \(H_0\) is an analytic unit on
the protected local chart; its actual logarithm carries its
boundary phase throughout the overlap.

The adjacent vanishing arcs are the two adjacent root arcs of the
original cubic. Their double-cover lifts meet once. If their positive
twists are \((g_1,g_2)\), the actual right Hurwitz arrow is
\((g_1,g_2)\mapsto(g_2,g_2^{-1}g_1g_2)\). Six full twists
conjugate both curves by \((g_1g_2)^{-6}\). The branched boundary
angle is \(\vartheta=2\widetilde\vartheta\); one full turn of
the three roots lifts to a half boundary turn, and its square is
\((g_1g_2)^6=T_\partial\). The weighted chart (C45.2) has the
opposite inner-to-outer arrow \(T_\partial\). These are the two
directions of the same retained collar map. The full unit logarithm
in (C45.2) transports that map by its actual boundary vector field;
it introduces no extra integer turn.

Choose a finite \(t\)-disc containing this pair and no other
singular value. At its boundary the original elliptic map is a
submersion for the whole closed \(c\)-disc. The source total
family is smooth through the selected cusp by SC45.4--6. Its
projection to \(c\) is a submersion there since \(q\) is the
actual implicit function of \((X,y,c)\). A horizontal field
tangent to that fixed outside boundary gives a proper source-disc
trivialization, with its inverse obtained by reverse flow.

The rest of the original surface and the cap extend on the disc
\(\zeta=v^2\) before pullback. In particular
\[
 a_0=-u_0/4\ne0,\quad t_c(0)=3/(4u_0)\ne0,\quad
 D(0,a_0)=153/16,
\quad Q_\infty=z^8\kappa_\infty,\quad
 P_\infty=z^3\pi,
\tag{C45.3}
\]
with \(\kappa_\infty(0)=-a/2304\) and \(\pi(0)=-1/24\).
The full inverse Tate functions and the actual consecutive charts
\(X_i=Z/z^i,Y_i=z^{i+1}/Z\), \(X_iY_i=z\), are the O45.5
maps throughout this disc. Their source and line periods retain
\(f_{Q^3b_0}(Z)=(-Z/b_0)^3Q^{-6}f_{b_0}(Z)\),
\(b_0=P_\infty^8/Q_\infty^3\), and the full reciprocal words
\(H=h^{-1}B^{-3},K=h^3B^8\),
\(h=H^8K^3,B=H^{-3}K^{-1}\). The unit \(b_0(0)=-256/27\)
keeps the two infinity point modules separated.

On every fixed angular sector of the original ordinary cap retain
its actual extension-group period \(\gamma(t,c,\lambda)\). The
map, its inverse, and its branch-change square are
\[
 \Psi(\omega,t,c)=
 (\omega+\log(t)\gamma/(2\pi i),t,c),\qquad
 \Psi^{-1}(\omega,t,c)=
 (\omega-\log(t)\gamma/(2\pi i),t,c),\qquad
 \Psi_{\beta}=B^k\Psi_{\alpha}
\tag{C45.4}
\]
when \(\log_\beta t-\log_\alpha t=2\pi ik\).
Here \(\gamma\) is the full source/line/scalar period of
\(B_\lambda=(T,\lambda C_*qB_{8,1})\); the given annular
integer cochain is transported from the reference. Every component
of (C45.4) is holomorphic in \(\zeta\). The inverse subtracts
the full target period. Thus the positive ordered loop closes by
the literal identity in the original source, line and cap charts.
The height-one reindex in SC45.20 is retained only when making
that explicitly stated reflected comparison.

\par

\medskip\noindent\textbf{C45.2. The actual compact orbit body and the embedded normalization trace}

\par\nobreak

Choose the scalar below the common norm bound for the full
unscaled suspension lift on the compact protected exterior and
cap. Its strict contraction gives a compact annular fundamental
domain and excludes an additional physical circle stabilizer:
an equality between a unit-circle action and \(B^k\) forces
\(k=0\) by the norm inequality. The remaining selected charts
are the literal fixed-edge charts
\[
 q=ZAB=Z'A'B',\qquad Z'=Z^{-1},\quad A'=ZA,\quad B'=ZB,
 \qquad \operatorname{wt}_\delta(A,B)=(-1,+1).
\tag{C45.5}
\]
The other edges are primitive for this physical cocharacter.
These are the same actual character rings used by the original
K10 construction; all four units \(\kappa,P^{-1},P^8,\lambda\beta\)
are retained on their protected overlaps.

For completeness the orbit-pair hypotheses used below follow from
these maps. Carry the K10 sphere along a path in the connected
marked open cover to the present boundary parameter. The full
O45.1--6 equations, line arrows, infinity charts and cap square
give the proper interval family and inverse flow. A prescribed
fan is carried from the reference by the finite two-ruling maps
written in C45.8 below. For the current vertical fan the only
fixed set is \(F=S^2\). If \(E=X\setminus\operatorname{int}
(F\times D^4)\), Alexander duality gives the homology of \(S^3\)
and codimension four gives \(\pi_1E=0\). Conjugating the negative
normal line in (C45.5) gives \(\mathcal O(1)\oplus\mathcal O(-1)\).
Its SU(2) clutch bounds the explicit disc with first column
\((re^{i\phi},\sqrt{1-r^2})\). Thus the actual boundary Hopf
map is \(S^2_F\times S^3\to S^2_F\times S^2_{\rm mer}\).

The Mayer--Vietoris map \(H_3(\partial E)\to H_3(E)\) is
primitive and an isomorphism. Its Hopf fibre integration and
Gysin naturality give \(H^2(E/S^1)=\mathbb Ze\),
\(H^3(E/S^1)=H^4(E/S^1)=0\). On the boundary the classes are
the base and meridian classes. For the original fixed-tube filling
\(F\times D^3\), their difference matrix is
\(\left(\begin{smallmatrix}0&1\\-1&0\end{smallmatrix}\right)\).
It gives a simply connected homology five-sphere. The verified
F8 five-sphere and fixed-pair classification applies to precisely
this Hopf attachment and identifies \((Y,F)\) with the standard
\((S^5,S^2)\). The later surgery replaces this filling by
\(D^3_F\times S^2_{\rm mer}\), using the same Hopf line; these
two displayed attachments are kept separately.

The source translation has an actual logarithm near the coalescing
pair. At the central cuspidal smooth group, divide the value of
\(-G\) by a sufficiently large integer, continue that division
section by the inverse function theorem, and take its small-group
logarithm. Multiplication by the same integer gives the original
translation. The invertible lift \(B/q\) is treated in the
line-extension group with its full analytic scalar logarithm and
the given reference. Exponentiating a real multiple gives a
source-and-line isotopy with its literal inverse. Its endpoint is
the original lift, including \(C_*\), \(\lambda\), and every
Cartier coefficient.

The relative orbit construction of F9 is applied to that actual
isotopy. In its node charts its two coordinates satisfy
\[
 2b=r-|y_0|^2/r,\qquad
 2b'=|x_0|^2r-r^{-1},\qquad
 \frac{db'}{db}=
 \frac{|x_0|^2+r^{-2}}{1+|y_0|^2r^{-2}}>0.
\tag{C45.6}
\]
The positive derivative, full positive unit norms and protected
outside coordinate give a proper marked circle coordinate. The
ordinary source and line logarithms extend through the parameter
disc, and (C45.4) closes literally. A nonzero suspension deck
would change the strictly contracting line norm. It therefore
cannot occur in that closing. The fixed graph has a real lift in
one common short circle arc, with the branch \(S=-c\) at the
lower F9 endpoint and \(S=c\) at the upper endpoint. This
statement uses the outward F9 coordinate; the coefficient of the
contracting period has the opposite orientation, as in K63.

In the source product the projection of this fixed sphere is
\[
 f(S,c)=(S^2-c^2,S(S^2-c^2),c),
 \qquad N=\mathbf P^1_S\times D_c.
\tag{C45.7}
\]
The full fractional translation used above restricts to
\[
 T_t(S,c)=
 \frac{(1+g_t)S+(1-g_t)c}{((1-g_t)/c)S+(1+g_t)},
 \quad g_t=\exp(t\log P^{-1})=1+ch_t,
 \quad T_t-S=\frac{h_t(S^2-c^2)}{2+h_t(c-S)}.
\tag{C45.8}
\]
It fixes both branches and has derivative \(I+O(|S|+|c|)\),
uniformly in the bounded height and its first derivative. On a
common smaller ball the derivative error has norm below one half.
This gives the actual orientation-preserving germ isotopy removing
the fractional source shear. Construct the height interpolation
first in the real circle cover. Quotienting its short arc creates
no additional double points; (C45.8) then gives its literal source
reparametrization and inverse.

The derivative-minor ideal of (C45.7) is \((S,c^2)\); its only
distinct double pair is \((S,c)=(c,c),(-c,c)\). Put
\[
 h_f(S,c)=\frac{\operatorname{Re}(S\overline c)}
                    {\sqrt{|S|^2+|c|^2}}.
\tag{C45.9}
\]
A small positive multiple has the same strict branch order.
Interpolate to it on the parameter annulus by a convex combination
in the marked short arc, keeping the actual outer graph fixed.
Away from the displayed double pair the source map is injective;
at that pair the height difference stays positive. The result is
an embedded graph of
\(W=N\setminus\operatorname{int}B_f\), where
\(B_f=f^{-1}(B^6_\epsilon)\) is the actual star-shaped four-ball.
The norm of (C45.7) increases strictly on each positive ray, so its
boundary has the complex-domain orientation. The ambient inner
ball is the rounded product
\(B^6_\epsilon\times[-2\eta\epsilon,2\eta\epsilon]\).
Because \(|h_f|\le|c|\le\epsilon\), the graph meets only its
cylindrical boundary. Its link is literally
\((f,\eta h_f):f^{-1}(S^5_\epsilon)\to
S^5_\epsilon\times(-2\eta\epsilon,2\eta\epsilon)\).
Using the orientation-preserving F55 stereographic chart gives
the exact F55 lifted link, with the same ordered branches and
projection-normal plus increasing-height frame. Thus its proved
invariants are complex Smale integer \(-2\) and Haefliger integer
\(+1\); this application is through the displayed germ and
framed boundary map.

\par

\medskip\noindent\textbf{C45.3. Source determinant, genuine point line and the normal spinor}

\par\nobreak

The original finite-base form \(dt\wedge dx/y\) is nowhere zero
on the smooth minimal source surface. On the actual threefold
chart SC45.4 its wedge with \(dc\) is
\[
 \Omega=-\frac{2}{H+qH_q}\,dX\wedge dy\wedge dc.
\tag{C45.10}
\]
This follows by differentiating the original implicit equation:
\(q_y=2y/(H+qH_q)\); the \(c\)-dependent shifts of \(x,t\)
vanish after wedging with \(dc\). The coefficient is an extending
unit on the inner source chart. Write \(x_-=S-c,x_+=S+c\),
\(D_\pm=(x_\mp=0)\), with \(D_+=(x_-=0)\) and
\(D_-=(x_+=0)\). The actual cross product is
\[
 n_f=(3S^2-c^2)dX-2S\,dy+2c(S^2-c^2)dc.
\tag{C45.11}
\]
Its determinant isomorphism identifies the complex normal line
with
\(L=\mathcal O(-D_+-D_-)\). The precise residue module arrow is
\[
 \mathcal O(-D_+-D_-)\longrightarrow K_N,\qquad
 f\longmapsto f\frac{dS\wedge dc}{x_-x_+}.
\tag{C45.12}
\]
Its frame \(x_-x_+\) maps to \(dS\wedge dc\); at infinity
the displayed form is \(-dz\wedge dc/(1-c^2z^2)\).
Thus the full determinant type is fixed also at \(S=\infty\).
The graph-normal comparison adds only the actual triangular
height shear; its contraction is a positive triangular homotopy.

The genuine source line is carried by its literal Cartier map
\[
 \sigma_8\longmapsto S-S_8(c),\qquad
 W=(S-S_8)v_M,\qquad
 V_{\rm pt}=\frac{W}{c-S_8},\qquad
 V_{\rm tensor}=\frac{W}{(c-S_G)^8}.
\tag{C45.13}
\]
The zero at \(8G\) and the pole at \(O\) cancel the matching
Cartier generators, so this is a nonzero line-module isomorphism
on the entire normalization, with its reciprocal. Both remaining
base factors are actual analytic units. This also includes the
central point value and the full reference coefficient from
SC45.13--18. For any one coefficient arrow \(W=gV\), the native
normal arrow is \((A,B)\mapsto(g^{-1}A,gB)\), and the ordered
positive spinor changes by \((\overline g^{-1},g)\). The two
phases agree. All positive norms stay in the corresponding
equivariant normal rescaling.

The determinant-one spinor on \(W\) is
\[
 E=\mathcal O(D_+)\oplus\mathcal O(-D_-),\qquad
 \det E=\mathcal O(D_+-D_-),\qquad
 \det E\xrightarrow{\;Z=x_-/x_+\;}\mathcal O.
\tag{C45.14}
\]
Its finite Cartier frames are \((1/x_-,x_+)\), and its adjoint
is \(\mathbb R_\gamma\oplus L_{\mathbb R}\). For the map
\(r=[x_-:x_+]:W\to\mathbf P^1\), the coefficient arrow to
\(r^*(\mathcal O(1)\oplus\mathcal O(-1))\) is
\(\operatorname{diag}(1/x_+,x_+)\) on \(x_+\ne0\), with
the literal reciprocal arrow on the other chart. The full original
height gives, without substituting its Tate multiple,
\[
 dq=-\frac{x_-x_+}{H}n_f,\qquad
 d(AB)=-\frac{x_+^2}{H}n_f,\qquad
 d(A'B')=-\frac{x_-^2}{H}n_f,
 \quad H=H(0,x_-x_+,c).
\tag{C45.15}
\]
The same \(-1/H\) multiplies both Hom arrows. It has an analytic
logarithm and square root on the inner ball; its full spinor
adjustment and inverse extend over that ball. The residue factor
in (C45.10) gives the identical determinant comparison up to the
explicit nonzero constant. It inserts no undetermined normal clutch.

After this product arrow is matched, the remaining weight-preserving
normal ratio has the exact form \((A,B)\mapsto(gA,g^{-1}B)\).
Removing its positive norm by positive rescaling makes its action
on the right-complex quaternion the physical phase
\(q_H\mapsto q_Hg/|g|\). Its adjoint normal action is the identity.
Its possible parameter winding is retained: the restriction
\(H^1(D^3\times S^1)\to H^1(S^2\times S^1)\) is an
isomorphism, so the full phase extends on the outer product cap.
The inner logarithms just constructed extend on the inner ball.

The actual positive-spinor to right-Cox map, with physical action
\(q_H\mapsto q_He^{-i\theta}\), is
\[
 q_H=(1+jZ)(A+j\overline B),\qquad
 (a,b)\longmapsto
 \frac1{\sqrt{|x_-|^2+|x_+|^2}}
 \begin{pmatrix}x_+&-\overline{x_-}\\x_-&\overline{x_+}\end{pmatrix}
 \overline{(a,b)}.
\tag{C45.16}
\]
The Hermitian dual map on the associated line reverses that circle
character. With \(J=\left(\begin{smallmatrix}0&-1\\1&0\end{smallmatrix}\right)\),
complex conjugation of every SU(2) transition is conjugation by
\(J\). Hence the global anti-linear map
\((a,b)\mapsto(\overline b,-\overline a)\) commutes with the
clutches and extends both cap discs. Its real four-dimensional
determinant is positive; its adjoint normal map is the antipodal
three-plane map. In the literal F9 chart at \(A=1,B=0\), the
projective coordinate \(B/\overline A\), the outward normal
orientation and the circle-first convention give degree \(-1\)
for the physical nonzero-vector line \(M\). Its Hermitian dual
\(U=M^*\) has meridian degree \(+1\), with these exact maps
on both caps.

\par

\medskip\noindent\textbf{C45.4. The outer product, both relative Chern caps and the physical pair}

\par\nobreak

For each outer parameter put
\[
 A_c=\begin{pmatrix}1&-c\\1&c\end{pmatrix},\qquad
 \mathcal B_c(P)=\frac{A_cPA_c^*}{\operatorname{tr}(A_cPA_c^*)}.
\tag{C45.17}
\]
Positive semidefinite Hermitian matrices of trace one form the
closed three-ball; its rank-one boundary is \(\mathbf P^1\).
The inverse uses \(A_c^{-1}\). These are periodic ball
diffeomorphisms and their boundary is exactly \(Z=(S-c)/(S+c)\).
The actual product-cap arrow is
\[
 (P,q_H,c)\longmapsto(\mathcal B_c(P),q_H,c).
\tag{C45.18}
\]
It holds the transported Cox quaternion fixed. Its quotient keeps
the Hopf normal sphere fixed as well. In particular the complete
determinant factor stays present:
\[
 x_+^2dZ=2c\,dS-2S\,dc,\qquad
 x_+^2d_{\rm rel}Z=2c\,d_{\rm rel}S,\qquad \det A_c=2c.
\tag{C45.19}
\]
Here \(d_{\rm rel}\) is the restriction to fixed \(c\).
The full \(-2S\,dc\) column is retained in the triangular
parameter-to-source derivative; wedging with \(dc\) gives
the displayed determinant factor. The parameter winding is one.
Equations (C45.15) and (C45.18)
use the original \(q\); the winding \(-6\) of \(Q/q\) remains
in the adjacent source-period chart. No derivative of
\(\mathcal B_c\) is substituted for the held normal quaternion.
The fractional source, Cartier and line isotopies extend by their
actual boundary velocities with a collar cutoff and physical-circle
averaging. This carries their full units through (C45.18).

The complement of the small inner seven-ball in \(Y\times D_c\)
has \(H^2=H^3=0\). Its codimension-three Thom sequence for the
connected embedded \(W\) therefore identifies degree-two line
classes on the complement with the integral normal-meridian value.
On the outer complement the Wang sequence for the unknotted
\(S^2\subset S^5\) has precisely that same degree-two generator
and no \(H^1\) term. The actual dual line \(U\), with its
meridian value one, consequently extends through this complement.
On \(S(\nu W)\) its outer restriction determines both the base
and meridian summands. The Cox map (C45.16) fixes that restriction
and the determinant trivialization. Any remaining outer boundary
isomorphism is a circle gauge; its full parameter winding extends
over the stated outer cap by the \(H^1\) isomorphism above.
The inner \(S^3\) has no \(H^1\) obstruction. These are actual
line isomorphisms, obtained by a nonzero section of the trivial
ratio line and the extending boundary gauge.

Cap \(W\) with the inner projection-normal four-ball and the
outer \(D^3\times S^1\) carrying (C45.18). The base is
\[
 \overline W=(S^2\times D^2)\cup(D^3\times S^1)
            =\partial(D^3\times D^2)=S^4.
\tag{C45.20}
\]
The first piece has its complex orientation and the second has
the negative displayed product orientation. The spinor extends
over the inner ball in its actual Cartier frame. Its determinant
trivialization then extends over \(N\simeq S^2\), so its SU(2)
bundle there is trivial. Choose a global frame agreeing with the
inner frame: the difference of two frames on that ball is an
extending gauge, and its contraction on the ball extends with a
collar cutoff. After the constant
anti-linear conjugacy in (C45.16), the frame arrow is the displayed
SU(2) matrix \(g(x_-,x_+)\). Its first column is \((x_+,x_-)\).
The complex-linear map \((S,c)\mapsto(x_+,x_-)\) has determinant
\(-2\), hence its complex-oriented boundary degree is one.
This comparison extends on \(W\), so its outer degree equals
that inner positive-domain degree. In the first-column orientation
the matrices
\(I=\operatorname{diag}(i,-i)\),
\(J=\left(\begin{smallmatrix}0&-1\\1&0\end{smallmatrix}\right)\),
\(K=\left(\begin{smallmatrix}0&i\\i&0\end{smallmatrix}\right)\)
satisfy \(IJ=-K\), \(\operatorname{tr}(IJK)=2\). Thus
\(\int\operatorname{tr}(dg\,g^{-1})^3=24\pi^2\).
A zero connection in the first frame is \(-dg\,g^{-1}\) in
the cap frame. Its Chern--Simons form is one third of this cubic
form; the cap boundary has the negative orientation. Therefore
\[
 c_2(\overline E)[S^4]=-1,\qquad
 p_1(\operatorname{ad}\overline E)[S^4]=4.
\tag{C45.21}
\]
This is the actual relative clutch, including the determinant
winding in (C45.19).

On the sphere boundary \(B=S(\operatorname{ad}\overline E)\),
the degree-one line class \(x_U\) has
\(e(T_vB)=2x_U\), and
\(4x_U^2=\pi^*p_1(\operatorname{ad}\overline E)\).
With \(\pi_*x_U=1\) this gives the individual pair \((1,4)\).
For the inner lifted-link surgery use its actual projection-normal
and increasing-height frame from C45.2. The F55/K42 Seifert
identities are \(2(-2)=3\sigma+e^2\) and
\(8=e^2-\sigma\), hence \(\sigma=-3,e^2=5\).
The truncated Seifert surface caps by the constant \(-\nu_1\)
section of \(D^4\times S^2\). Its oriented normal plane is the
positive meridian plane, so it is the Poincare dual of \(x_U\).
Its own pair is \((e^2,3\sigma+e^2)=(5,-4)\).

The outward boundary of the actual quotient bordism is
\(-P_f\sqcup-N_e\sqcup-B\). The line \(U\) extends with
the two just-computed restrictions. Stokes gives the positive
family pair \((-6,0)\) for \(U\). Jin\textquotesingle s first-to-second
endpoint convention reverses this parameter. The duality in
(C45.16) then gives
\[
 (k_U,\ell_U)_{\rm Jin}=(6,0),\qquad
 (k_M,\ell_M)_{\rm Jin}=(-6,0),\qquad A_M=-1.
\tag{C45.22}
\]
Here the physical Hopf transfer retains its inverse: the actual
exterior is \(D^3_y\times S^3_q\), the sphere is
\(-\partial(D^3_y\times D^4_q)\), and the physical marking
uses \(u_M=q,y_M=-y\). Its two pieces are the negative of the
second and first pieces of
\(\Sigma_F=(D^4_y\times S^3_q)\cup_F(S^3_y\times D^4_q)\).
Reordering the actual first-to-second relation inverts \(F\).
Thus \(\lambda_M(F)=-\Sigma_{F^{-1}}=\Sigma_F\), while
\(\lambda_U(F)=\Sigma_{F^{-1}}\). The K68 cut-out ball has
the same source/Hopf coordinates and line gauges. This identifies
the physical return class \(-1\bmod28\) and its actual native
fixed-star pair \((\mu,\nu)=(0,-6)\), through
\(k_M=\mu+\nu,\ell_M=4\mu\) in this held product.

\par

\medskip\noindent\textbf{C45.5. Every vertical fan: the full relative source and line maps}

\par\nobreak

For each integer \(j\), retain the actual vertical fan
\(T_j=T(e_1+je_2,e_2)\). On the same shrinking selected collar
write the complete periods as
\[
 h=(\kappa qZ,\alpha V),\quad B=(gZ,\eta qV),\quad
 g=P^{-1},\quad\alpha=P^8,\quad\eta=\lambda\beta.
\tag{C45.23}
\]
Their ordered-parameter windings are \((-6,0,0,0)\), by SC45.17--18.
Each of the last three full functions has its actual periodic
logarithm on the protected \(q\)-disc times parameter annulus.
The first has that logarithm after its explicit winding is retained.
The coordinate/period arrow and inverse are
\[
 \begin{gathered}
 V_j=Z^{-j}V,\quad H_j=hB^j,\quad Q_j=\kappa qg^j,\\
 \alpha_j=\alpha\eta^j\kappa^{-j}g^{-j^2},\quad
 \eta_j=\eta g^{-j};\qquad V=Z^jV_j,\ h=H_jB^{-j}.
 \end{gathered}\tag{C45.24}
\]
Every exponent follows by substitution into the full commuting
period maps; the height stays \(q\).

The exact source interpolation is \(Q_t=Q\exp(tj\log g)\).
In a node chart \(xy=Q_t\), its horizontal field is
\(\partial_t+j\log g\,y\partial_y\). On the protected
outside collar prescribe the real marked-period map
\(u+\tau_tv\mapsto u+\tau_0v\). These fields project to
the same \(\partial_t\); their difference is vertical.
A partition of unity, fixed on the two smaller protected collars,
gives a smooth horizontal field on the compact source with
boundary. Its forward and reverse flows give the actual map,
periodic on this parameter annulus. The full F9 coordinate is
formed from the real period coefficients: its source deck has
circle lift zero and its suspension deck lift one. The two fixed
endpoint lifts stay in one short arc. Thus the transported graph
isotopy is the convex interpolation used in C45.2, now with the
actual \(j\)-dependent source flow.

On the spinor the full Cartier arrow is multiplication by \(Z^j\),
with the precise source type
\[
 E_j=E\otimes\mathcal O(j(D_+-D_-)),\qquad
 (x_-^{-j-1}x_+^j,\ x_-^{-j}x_+^{j+1})
       \xrightarrow{\ Z^j\ }(1/x_-,x_+).
\tag{C45.25}
\]
Its determinant trivializations agree. The normal source-flow
derivative and triangular graph shear extend through their actual
isotopies. Extend the transported sphere velocity over the
three-ball cap with a collar cutoff and retain the adjoint of
(C45.25) on its Hopf factor. This gives the actual relative
completed-orbit-piece diffeomorphism, with the physical line
carried separately by the following full boundary arrow.

Use real coefficients \((u,b,d,\gamma)\) of the vanishing,
source, physical-delta and contracting periods. In full logarithmic
coordinates
\[
 (z,\ell)=u(1,0)+b(\tau,\mu)+d(0,1)+\gamma(G,\mathcal B),
\quad
 (\tau,\mu,G,\mathcal B)=
 \frac1{2\pi i}(\log(\kappa q),\log\alpha,\log g,\log(\eta q)).
\tag{C45.26}
\]
In particular \(d=\operatorname{Re}\ell-\operatorname{Re}\mu\,b
-\operatorname{Re}\mathcal B\,\gamma\); differentiating retains
every full unit coefficient and derivative. With \(\alpha_q\)
the \(q\)-angle and \(\theta\) the positive parameter angle,
the actual deck changes and character arrow are
\[
 (u,d)\xrightarrow{q\text{-deck}}(u-b,d-\gamma),\quad
 (u,d)\xrightarrow{c\text{-deck}}(u+6b,d),\qquad
 (u_j,b_j,d_j,\gamma_j)=(u,b,d-ju,\gamma-jb).
\tag{C45.27}
\]
The physical principal connections and their complete difference
can therefore be chosen as
\[
 A_0=dd+\alpha_q\,d\gamma,\quad
 A_j=dd_j+\alpha_q\,d\gamma_j+6j\theta\,db,\quad
 \beta_j=A_j-A_0=-j\{du+(\alpha_q-6\theta)db\}.
\tag{C45.28}
\]
These are formulas in the actual coefficient frame (C45.26), so
the zero-winding unit derivatives are included rather than set to
zero. For the associated lines \(c_1(M)=-dA\) and
\(c_1(U)=dA\), with the exact Hermitian duality of C45.3.

Let the source graph map \(b\) be extended on the completed
piece and made constant near its fixed tube. The theta-cut line
\(\mathcal T\) has coefficient arrow \(e^{2\pi ib}\),
connection \(d-2\pi i\theta\,db\), and curvature
\(\Xi=d\theta\wedge db\), \(\Xi^2=0\).
The full deck square in (C45.27) gives the relative line map
\[
 U_j\longrightarrow U_0\otimes\mathcal T^{6j}.
\tag{C45.29}
\]
Here is its boundary-integer verification. The \(q\)-commutator
kills \(u\) primitively; the parameter commutator is \(6u\).
Thus the outer \(H_1\) is freely generated by
\(\alpha_q,\theta,b,\gamma\), while the last three extend to
the completed piece by the source graph, marked circle and cap.
The only possible nonextending gauge is the \(q\)-angle.
Choose \(Z_*\) near \(-1\), away from the original point/O
sections and both node branches, and the free lift \(\gamma=0\).
The actual section disc and its inverse image are
\[
 (q/Z_*,Z_*,1)\longmapsto
 (qZ_*^{j-1},Z_*,Z_*^{-j}).
\tag{C45.30}
\]
The source flow can be prescribed to fix this disc. All full
point/O modules are nonzero there. The line arrow is the displayed
constant nonzero power and extends across \(q=0\), so its
\(q\)-winding is zero. At another free circle lift, the full
factors \(Z_*\exp(\gamma\log g)\) and
\(\eta_j^\gamma=\eta^\gamma\exp(-j\gamma\log g)\)
cancel in the ratio; each is retained and extending on the disc.

To obtain (C45.29) with the prescribed boundary value, form the
actual ratio line with its full transitions. Remove its integral
Cech two-cocycle by the integral cochain given by the deck square
and meridian/base evaluations. For remaining real logarithms
\(r_{ij}\) of vector-frame transitions, a partition of unity
gives \(a_i=\sum_k\rho_kr_{ki}\), \(a_j-a_i=r_{ij}\).
The frames multiplied by \(e^{-2\pi ia_i}\) give a global
nonzero section of the ratio line; take its inverse in direction
(C45.29). Remove the remaining extending \(\theta,b,\gamma\)
characters from its boundary ratio, and extend the final real
boundary logarithm by a collar cutoff. This constructs precisely
the specified boundary arrow and its inverse, including the cap.

\par

\medskip\noindent\textbf{C45.6. Relative tangent pairing and the complete vertical integers}

\par\nobreak

On the boundary let
\(\omega=du+\operatorname{Re}\tau\,db+
\operatorname{Re}G\,d\gamma\). The full new period is
\(\tau_j=\tau+jG\), and \(d\gamma_j=d\gamma-j\,db\),
so its exact pullback is
\[
 \omega_j=du+\operatorname{Re}(\tau+jG)\,db+
       \operatorname{Re}G(d\gamma-j\,db)=\omega,
 \qquad
 \begin{pmatrix}\omega_j\\d\gamma_j\\db_j\end{pmatrix}
 =\begin{pmatrix}1&0&0\\0&1&-j\\0&0&1\end{pmatrix}
  \begin{pmatrix}\omega\\d\gamma\\db\end{pmatrix}.
\tag{C45.31a}
\]
The base coordinates are fixed on this marked boundary; their
full coefficient derivatives give the retained base-to-source
off-diagonal blocks. The source/graph derivative is therefore
triangular in this actual order, with diagonal entries one.
Its strictly triangular source connection
contributes zero to \(\operatorname{tr}(F^2)\). The base block
on the boundary depends on the two base angles; its interpolation
adds one collar variable, so its curvature-square is zero by
dimension. The entire block-triangular tangent comparison therefore
has zero Pontryagin transgression. All full unit logarithms occur
in its off-diagonal entries and the line connection (C45.28).

For its remaining relative pairing choose nested neighbourhoods
\(U_1\Subset U_2\Subset U_3\) of the fixed tube, containing all
surgery and connection changes in \(U_1\). On \(U_3\) take a
real lift \(\widetilde b\) of the original source-product circle
map. With \(\chi=1\) on \(U_2\), supported in \(U_3\), set
\[
 b_t=b\exp(2\pi it\chi(b_* -\widetilde b)),\quad0\le t\le1.
\tag{C45.31}
\]
The homotopy is literally the original map outside \(U_3\),
and is constant on \(U_2\) at the final endpoint. A common regular
value different from \(e^{2\pi ib_*}\) gives a relative level
cobordism fixed at the outer boundary. Its terminal level misses
all surgery changes and is the same cycle on the two pieces.
The original level is the actual source three-level times the
contracting circle; the source-pulled Pontryagin form integrates
to zero there. Stokes and the just-proved zero boundary
transgression give zero on the terminal level. This uses the
displayed relative cobordism; it does not replace that terminal
level by an unproved product identification. Hence
\(\ell_{M,j}=\ell_{M,0}=0\).

Use (C45.29) to choose curvatures \(F_j=F_0+6j\Xi\) on the
fixed completed piece. Their cubic interior difference is affine
in \(j\), since \(\Xi^2=0\). On the boundary, written in the
new piece coefficient \(\gamma=\gamma_j\), the full exterior
pullback and the line difference are
\[
 F_{{\rm ext},j}=d\alpha_q\wedge d\gamma+j\,d\alpha_q\wedge db,
\quad d\beta_j=-j\,d\alpha_q\wedge db+6j\,d\theta\wedge db.
\tag{C45.32}
\]
Orient the positive family by
\(dr_q\,d\alpha_q\,du\,db\,d\gamma\,d\theta\).
Its outward five-form is
\(\Omega_\partial=d\alpha_q\wedge du\wedge db\wedge d\gamma
\wedge d\theta\); its primitive marked cell has integral one.
The exact wedge identity is
\[
 3\int_0^1\beta_j\wedge(F_{{\rm ext},j}+s\,d\beta_j)^2ds
                 =18j^2\Omega_\partial.
\tag{C45.33}
\]
The outward gluing collar runs from the piece to the exterior,
so \(s\) runs from one to zero and contributes \(-18j^2\).
The physical/Jin convention in C45.4 gives
\(k_{M,j}=k_{M,0}+Lj-18j^2\) with one constant \(L\).
This follows from the complete relative connection, including its
boundary cochain and every full coefficient.

The second endpoint is fixed by the actual full reflection, not
by its residue. SC45.19--21 give
\((Z,V)\mapsto(Z^{-1},P^8V)\), followed by the actual target
source period to give \((Q/Z,V)\). The entire outside and cap
reference was constructed on \(\zeta=v^2\); it is carried by
this exact map. Its integral height translation is one, so its
image fan is \(T_1\), with both physical line and product caps
preserved. It gives \(k_{M,1}=k_{M,0}=-6\), hence \(L=18\).
For every integer \(j\) the full answer is consequently
\[
\begin{split}
 (k_M,\ell_M)_j&=(-6-18j(j-1),0),\\
 A_{M,j}&=-1-3j(j-1),\\
 (\mu_j,\nu_j)_{
 \rm native}&=(0,-6-18j(j-1)).
\end{split}
\tag{C45.34}
\]
The last row uses the actual held Cox/K68 product from C45.4,
with \(k_M=\mu+\nu\), \(\ell_M=4\mu\). The ordinary-to-native
comparison on the plumbing is the full pure-left cylinder given
next; its fixed ordinary entries are therefore
\(V_{00}=18j(j-1)\), \(\delta_0=-12-18j(j-1)\).

\par

\medskip\noindent\textbf{C45.7. Full ordinary products and the circle/derivative prism}

\par\nobreak

Let \(P_{\rm bas}\in GL(2,\mathbb Z)\), put \(R=P_{\rm bas}^{-1}\),
and abbreviate its entries by \(P_{ij}\). The actual character
rows and period columns give
\[
 A_{ik}=(\kappa^{P_{1k}}g^{P_{2k}})^{R_{i1}}
         (\alpha^{P_{1k}}\eta^{P_{2k}})^{R_{i2}},\qquad
 B_P=-6P_{\rm bas}^{-1}e_1e_1^tP_{\rm bas}.
\tag{C45.35}
\]
The original \(q\)-valuation is \(RP=I\). The full periodic
logarithm of each coefficient after its displayed winding is
\(R_{i1}(P_{1k}f_\kappa+P_{2k}f_g)+
R_{i2}(P_{1k}f_\alpha+P_{2k}f_\eta)\). Here the four \(f\)\textquotesingle s
are the actual based logarithms from (C45.23), with only the
specified \(-6\) winding separated in \(f_\kappa\).

At the two vertices the three exact edge rows are
\[
\begin{array}{c|ccc}
 C_0&A_{11}A_{21}bc&A_{21}^{-1}ab&A_{11}^{-1}b^{-1}\\
 C_1&A_{12}A_{22}bc&A_{22}^{-1}c^{-1}&A_{12}^{-1}ac\\
 C_2&a^{-1}&ab&ac.
\end{array}
\tag{C45.36}
\]
Each triple product is \(q\). The actual native second charts are respectively
\[
\begin{gathered}
(A_{11}w,A_{21}v,(A_{11}A_{21})^{-1}u),\\
(A_{22}v,A_{12}w,(A_{12}A_{22})^{-1}u),\qquad (u,v,w);
\end{gathered}
\]
their first normal pairs are \((a,c),(a,b),(b,c)\).
They give \(z'=z^{-1},n'=zn\) literally. For the ordinary product,
keep the first vertex and use the full second-vertex arrow
\[
 (u,v,w)\longmapsto
 (u,v',w')=(u,A_{22}(s,q)v,A_{11}(s,q)w),\qquad
 h_s(q)=qA_{11}(s,q)A_{22}(s,q).
\tag{C45.36a}
\]
Here \(s\) denotes the parameter on the retained compact annulus.
The new coordinate product is \(uv'w'=h_s(q)\). The inverse first
recovers \(q=h_s^{-1}(uv'w')\), then divides \(v'\) by
\(A_{22}(s,q)\) and \(w'\) by \(A_{11}(s,q)\), evaluated at that
recovered original height. The derivative of \(h_s\) at zero is
the nonzero product of the two actual units. The parameter annulus
is compact, so a uniform smaller selected disc gives this inverse.
In the ordered normal pairs \((a,c),(a,b),(b,c)\), the ordinary
transitions are \(zD_i\), where
\[
 D_0=\operatorname{diag}(A_{22}/A_{21},A_{11}A_{21}),\quad
 D_1=\operatorname{diag}(A_{11}/A_{12},A_{12}A_{22}),\quad
 D_2=\operatorname{diag}(A_{22},A_{11}).
\tag{C45.36b}
\]
Thus the first normal windings are
\(l=(B_{22}-B_{21},B_{11}-B_{12},B_{22})\), each with
determinant winding \(-6\).

Write
\[
 H_d(s,\theta)=
 (\sin\vartheta(s)e^{id\theta}+j\cos\vartheta(s))
 (\sin\vartheta(s)+j\cos\vartheta(s))^{-1}.
\tag{C45.37}
\]
The fixed smooth radial function has \(H_d=1\) on the central
collar and \(H_d=e^{id\theta}\) on the outer collar; it is
based at \(\theta=0\) and periodic. On edge \(i\), use
\(Q\mapsto H_{-3}QH_{l_i+3}\), together with the actual
periodic, zero-winding diagonal logarithms \(k_{i0},k_{i1}\)
obtained from (C45.35) after separating these two winding factors.
The corresponding unit arrow is
\[
 (z,n_0,n_1)\longmapsto
 (z,e^{\rho(|z|)k_{i0}(s,q)}n_0,
    e^{\rho(|z|)k_{i1}(s,q)}n_1),\qquad
 H_{s,z}(q)=q\exp\{\rho(|z|)(k_{i0}+k_{i1})(s,q)\}.
\tag{C45.37a}
\]
The fixed cutoff is zero on the first vertex collar and one on the
second. Given the target normals, first recover
\(q=H_{s,z}^{-1}(z n'_0n'_1)\) and then set
\(n_j=e^{-\rho(|z|)k_{ij}(s,q)}n'_j\).
After division by the nonzero derivative at zero, the derivative
of \(H_{s,z}\) differs from one by less than one half on a
uniform smaller disc. The full parameter and radial collars are
compact. Integrating along a straight segment proves injectivity;
the inverse function theorem supplies the displayed inverse with
smooth radial dependence. This is OF.2\textquotesingle s exact height comparison,
applied to the present full periodic coefficients. It retains every
\(q\)-derivative and the original height in the target; the
coefficients in the inverse are evaluated at the recovered \(q\).
At the boundary its two complex winding factors are precisely
\(e^{il_i\theta},e^{i(-6-l_i)\theta}\). The identity
\(zD_{i,*}K_i=zD_i\) glues both vertices. Rescaling the full
normal germ and integrating its cutoff velocity carries every
higher normal jet and its inverse, fixed on the protected collars.
The Hopf product map is
\((UH_{-3}U^{-1})q_HH_{l_i+3}\). Its left degree is \(-6\)
and its right degree is zero, giving the actual inverse cylinder
\(L^6\), with the retained odd Spin endpoints. This is the
GF.9--12 product construction applied to these exact coefficients.

Put
\[
 \begin{gathered}
 (x,y)^t=P_{\rm bas}^{-1}e_2,\quad w=(x,y,-x-y),\\
 d=(-x-2y,-2x-y,x-y),\\
 I_i=-3w_i+d_i(l_i+3),\qquad J_i=-3w_i-d_i(l_i+3).
 \end{gathered}\tag{C45.38}
\]
The physical orbit moves the native normal quaternion by left
angle \(-w_i\phi/2\) and right angle \(d_i\phi/2\), conjugated
by the two actual factors (C45.37). Their axis degrees are
\(3,l_i+3\). At \(\phi=2\pi\) the base and normal Spin
endpoints both equal \((-1)^{w_i}\), hence their total is one.
The full GF.15 circle map and its derivative therefore give the
two integers (C45.38), with all unit conjugations and normal jets
retained. The three ordinary variations \(V\) and derivatives
\(\delta\) satisfy the actual integral prism equations
\[
 V(w_ie_i-w_0e_0)=I_ie_i-I_0e_0,\qquad
 w_i\delta_i-w_0\delta_0=J_i-J_0,\qquad i=1,2.
\tag{C45.39}
\]
On \(T_j\), (C45.34) and the full ordinary \(L^6\) comparison
give the complete seed
\[
 \begin{gathered}
 V=\begin{pmatrix}
 18j(j-1)&6-12j&12j-6\\6-12j&0&0\\12j-6&0&0
 \end{pmatrix},\\
 \delta=(-12-18j(j-1),12j-12,-12j),\qquad s=-1-3j(j-1).
 \end{gathered}\tag{C45.40}
\]

\par

\medskip\noindent\textbf{C45.8. Actual all-fan continuation and the exact terminating value algorithm}

\par\nobreak

For columns \(v,w\), the two triangles of \(T(P)\) are
\(\{n,n+v,n+w\}\) and \(\{n+v,n+w,n+v+w\}\).
Column swap and right multiplication by
\(R_0=\left(\begin{smallmatrix}-1&-1\\1&0\end{smallmatrix}\right)\),
\(R_0^2,-I,-R_0,-R_0^2\) give the identical absolute fan by
the explicit triangle translations of AF.1. The identity on its
absolute character semigroups is the inverse holomorphic chart
map, commutes with every full coefficient (C45.35), and is the
identity on the original source/Cartier/cap overlap. It carries
the prescribed closing and physical circle without a new phase.

After these actual identical-fan maps there is one representative
\(T_j\), or \(P=\left(\begin{smallmatrix}p&-r\\q&s\end{smallmatrix}\right)\)
with \(p,r>0\), \(ps+rq=1\). In the latter case its parent is
\(PQ_0^{-1}\) for \(r\ge p\), and \(PQ_1^{-1}\) for
\(p>r\), where
\[
 Q_0=\begin{pmatrix}1&-1\\0&1\end{pmatrix},\qquad
 Q_1=\begin{pmatrix}1&0\\-1&1\end{pmatrix}.
\tag{C45.41}
\]
The positive pair becomes \((p,r-p)\) or \((p-r,r)\).
Its sum decreases and its gcd is one, so equality is \((1,1)\)
and the next parent is a unique \(T_j\). Reversing first uses
\(Q_0\) and then gives a finite word of actual forward maps with all new weights
\((r,p,-r-p)\) nonzero.

Their dense character arrows are \((ZV,V)\) and \((Z,ZV)\),
their first triples are \((a/c,bc,c)\) and \((a/b,b,bc)\),
and their other-ruling native charts are
\((n_0/n_1,zn_1,n_1)\), \((n_1/n_0,zn_0,n_0)\).
They have common rank-one matrix
\(\left(\begin{smallmatrix}n_0&n_1\\zn_0&zn_1\end{smallmatrix}\right)\)
and unchanged height \(q=zn_0n_1\). For the second generator
the actual normal interchange is joined to the identity by
\(\frac12\left(\begin{smallmatrix}1+e^{i\pi t}&1-e^{i\pi t}\\
1-e^{i\pi t}&1+e^{i\pi t}\end{smallmatrix}\right)\); it commutes
with \(zI\) and preserves the oriented real four-plane. The
full coefficients transform by the literal monomial
\[
 A'_{ik}=\prod_{h,l}A_{hl}^{(Q^{-1})_{ih}Q_{lk}}.
\tag{C45.42}
\]
Thus the whole family flop, full normal jets, original point line,
outside annular maps and ordinary cap carry through each step.
The compact common blowup and its two projective ruling
projections, with their integral Gysin/van Kampen maps, preserve
the reference sphere topology as in the explicit F35/K23
construction. This proves the geometric existence of every
specified fan object before evaluating its return.

For the selected old edge \(i=0\) or \(1\), let
\(a=V_{ii},b=\delta_i\). The ordinary/native cylinder above
shows that its actual native right integer is \((b-a)/2\).
On the common link the full Cox arrow is
\((h,q_H)\mapsto(q_H^{-1}hq_H,q_H)\). With its carried
complex reset it sends
\((\mu,\nu,s)\mapsto(\mu+2\nu,-\nu,s+\nu)\).
The K68 physical commutator orientation is positive. Carrying the
same \(L^6\) cylinder on both sides gives
\[
 V'_{22}=a,\qquad \delta'_2=2a-b,\qquad
 s'=s+(b-a)/2\pmod {28}.
\tag{C45.43}
\]
The full target solution of (C45.39) is unique. Explicitly,
\[
 V'_{rs}=\frac{I_r}{w_r}\delta_{rs}+\frac{C}{w_rw_s},
 \quad\delta'_r=\frac{J_r+E}{w_r},\qquad
 C=w_2^2a-w_2I_2,\quad E=w_2(2a-b)-J_2.
\tag{C45.44}
\]
The difference of two symmetric solutions annihilates the plane
spanned by \(w_ie_i-w_0e_0\), so it is a multiple of
\((1/w_r)(1/w_s)\). The derivative difference has the analogous
single constant. The two selected values fix these constants.
Existence is furnished by the actual target family and full
products just constructed; its integral invariants equal the
displayed rational solution. Thus \((b-a)/2\) is the actual
integer native coordinate. No class is divided in a finite group.

Starting from (C45.40), these steps are exactly the proved K71
algorithm, now obtained from the original OS45 source/line/cap
maps. They determine the physical return for every integral
unimodular selected fan and both determinant signs. In particular
the initial fan has pair \((-6,0)\), native pair \((0,-6)\),
and physical class \(-1\); the vertical values are (C45.34).
Every retained full scalar and parameter unit travels through
the displayed arrows. Decreasing an admitted boundary scalar
along its constant-argument interval, carrying its actual
\(\log\lambda\) in (C45.4) and (C45.23), gives the relative
annular family to the common small scalar used in C45.2.
Its reverse flow returns to the original scalar with no added
scalar turn. This extends the same results to every admitted
fixed scalar in the stated original ordinary-cap family.

The exact local source checker SC45 and the accompanying completion
checker retain the source height, ordered braid, determinant/module
arrows, Chern signs, full connection wedge, actual coefficient
composition and fresh full prism solves. The complete geometric
maps above are required in addition to those executions. This
working completion is awaiting full independent geometric audit
before its reserved numbered export P3K.117.

The executed completion certificate passes 4,514 exact checks.
It freshly eliminates the complete nine-variable rational systems
on 217 actual canonical paths from \(j=-3,\ldots,3\), through
word length five, and compares the resulting full variations and
derivatives with (C45.44). Each path first enters the nonvertical
region by \(Q_0\); the other adjacent vertical fan is already
covered by its independently checked seed. The same execution
checks all full multiplier exponents and inverse coefficient maps,
the retained boundary coframes, oriented fifth form, source
critical-point Jacobian, residue/Cartier arrows and actual cap
signs. It imports no preexisting fan calculator. The unrestricted
all-fan conclusion is proved by the exact terminating construction
in C45.8, with this finite execution recorded as verification.

% End expanded verified cp36 input

\par\bigskip\noindent\textbf{Full supporting proof 3: Complete independent local source and point-line proof.}\par\nobreak

% Expanded verified cp36 input: notes/part3_os45_selected_cusp_local_blind_peer.tex
This proof reads SC45.1-\/-4 of
notes/part3\_os45\_selected\_cusp\_working.md in full. Its scope is the
original selected source, both normal coefficients, the complete nodal
Miller/point map and the smooth selected-collar reflection. The compact
trace and its return are treated in the author\textquotesingle s separate continuation.
No numerical return is assigned by this local audit.

\par

\medskip\noindent\textbf{SCP45.1. The entire original height and its inverse}

\par\nobreak

Keep the actual source
\[
 y^2=x^3-3\{(t^2+a)^2+t\}x+
       2\{(t^2+a)^3+\tfrac32t(t^2+a)+\tfrac38\}.
 \tag{SCP45.1}
\]
At the ordered selected cusp use exactly
\[
 a=-u/3-u^4/9,\quad t_c=1/(2u)-u^2/3,\quad
 \xi=1/(4u^2)+u/3,\quad v^2=4u^3+3,\quad c=v/(2u).
 \tag{SCP45.2}
\]
Thus \(3\xi=c^2\). The map \(v\mapsto c\) has derivative
\(1/(2u_0)\ne0\), and its actual inverse is the branch
\(u^3-u^2c^2+3/4=0\), \(v=2uc\), with \(u(0)=u_0\).
All original coefficients are therefore even in the ordered parameter.
The exact factorization
\[
 a-a_0=-\frac{(u-u_0)^2(u^2+2u_0u+3u_0^2)}9,\qquad
 u-u_0=\frac{v^2}{4(u^2+u_0u+u_0^2)}
 \tag{SCP45.3}
\]
has leading coefficient \(-1/(216u_0^2)\) in \(v^4\).
It supplies the asserted positive degree-four original parameter map.

Put \(q=t-t_c\), \(X=x-\xi\), \(R_c=\xi-u\) and \(d=2t_c+q\).
Direct substitution in SCP45.1 gives, without deleting a coefficient,
\[
 y^2=X^2(X+c^2)+qH,\qquad
 H=d[-6R_c(u+X)+3t_c]+3qd^2(2R_c-\xi-X)
                   +2q^2d^3-3(u+X)+3qd .
 \tag{SCP45.4}
\]
The coefficient \(H(0,0,c)=-2u^4\) is a unit at the cusp.
The inverse source map is consequently the unique solution of this
equation for \(q\), followed by
\(x=X+\xi(c)\), \(t=t_c(c)+q\), \(a=a(c)\), with \(y\) unchanged.
This proves the full source chart and its reciprocal chart, including
the original height rather than just its first nonzero term.

\par

\medskip\noindent\textbf{SCP45.2. Both original normal products}

\par\nobreak

On \(q=0\), use
\[
 X=S^2-c^2,\quad y=S(S^2-c^2),\quad
 S_G=-\sqrt3/(2u),\quad Z=(S-c)/(S+c).
 \tag{SCP45.5}
\]
The literal source translation and inverse are
\[
 S_T=\frac{S_GS-c^2}{S_G-S},\qquad
 S=\frac{S_GS_T+c^2}{S_G+S_T}.
 \tag{SCP45.6}
\]
Both denominators are units near the source cusp. With
\(P=(S_G-c)/(S_G+c)\), substitution gives
\(Z(S_T)=P^{-1}Z(S)\); \(S=\pm c\) are fixed separately.
Away from the normalization, the original chord arrow and its
plus-\(G\) inverse are evaluated at the source and target points of
SCP45.4. No normalization formula is substituted for that whole arrow.

Let
\[
 n_f=(3S^2-c^2)dX-2S\,dy+2c(S^2-c^2)dc.
 \tag{SCP45.7}
\]
For \(F=y^2-X^2(X+c^2)-qH\), restriction of its full differential
to the normalization gives
\[
 dF=-(S^2-c^2)n_f-H(0,S^2-c^2,c)dq.
 \tag{SCP45.8}
\]
Terms \(q\,dH\) vanish there because the actual parameter is \(q=0\).
Thus
\[
 dq=-\frac{(S-c)(S+c)}{H(0,S^2-c^2,c)}n_f.
 \tag{SCP45.9}
\]
In the native chart \(q=ZAB\), its opposite coordinates are
\(Z'=Z^{-1}\), \(A'=ZA\), \(B'=ZB\), with inverse obtained by
the same reciprocal operations. At \(AB=0\), \(dq=Z\,d(AB)\).
Consequently the two full normal products are
\[
 d(AB)=-\frac{(S+c)^2}{H(0,S^2-c^2,c)}n_f,\qquad
 d(A'B')=-\frac{(S-c)^2}{H(0,S^2-c^2,c)}n_f .
 \tag{SCP45.10}
\]
The same original unit and its reciprocal occur in both maps.

\par

\medskip\noindent\textbf{SCP45.3. Literal tangent functions and the two line frames}

\par\nobreak

For a node with arbitrary nonzero \(c\), write
\[
 S_m=c\frac{1+P^m}{1-P^m}.
 \tag{SCP45.11}
\]
The tangent at \(S_G\) has slope
\((3S_G^2-c^2)/(2S_G)\), and its double has parameter
\((S_G^2+c^2)/(2S_G)\). In the original \(x,y\) coordinates
the tangent quotient is therefore
\((S-S_G)^2/(S-S_2)\), with principal coefficient \(S\) at \(O\).
The actual chronological doublings give
\[
 f_4=f_2^2f_{2,2G}=\frac{(S-S_G)^4}{S-S_4},\qquad
 f_8=f_4^2f_{2,4G}=\frac{(S-S_G)^8}{S-S_8}.
 \tag{SCP45.12}
\]
These retain their original principal coefficient at \(O\).
In particular \(f_2(\pm c)=-2S_G\). In the full coefficient
\(B_{8,1}=f_2^8 f_8(S_T)/f_8(S)\), the source fixes each node
branch; all factors are nonzero there. Its node value is exactly
\((-2S_G)^8=81/u^8\).

The nodal theta function is
\[
 f_{P^m}(Z)=\frac{S-S_m}{c-S_m}.
 \tag{SCP45.13}
\]
Hence the actual directed point-to-tensor map and inverse are
\[
 V_{\rm tensor}=f_P(Z)^8f_8(S)^{-1}v_M
              =C_8 V_{\rm pt},\qquad
 V_{\rm pt}=C_8^{-1}V_{\rm tensor},\qquad
 C_8=\frac{c-S_8}{(c-S_G)^8}.
 \tag{SCP45.14}
\]
Here \(V_{\rm pt}=f_{P^8}(Z)v_M\). The analytic limits are
\(S_m(0)=S_G(0)/m\); thus both factors defining \(C_8\) are
units. The separate genuine point-line trivialization sends the
meromorphic section \(\sigma_8\) to \(S-S_8\), carrying its
point and \(O\) Cartier modules. That is the exact module cancellation
through the zero and pole; it does not replace the coefficient
arrow SCP45.14.

\par

\medskip\noindent\textbf{SCP45.4. Complete collar and reflection}

\par\nobreak

The original discriminant calculation has
\(D_t(t_c)=-2v^6/(27u^2)\) and
\(f(t_c)^3=\xi^6=v^{12}/(12^6u^{12})\). Using
\(Qj(Q)\to1\) in the original \(j\)-function gives
\[
 \left.Q/q\right|_{q=0}=2u^{10}/v^6,\qquad
 \beta(0,v)=C_*81/u^8 .
 \tag{SCP45.15}
\]
The height has winding \(-6\) on the actual shrinking selected
collar. These equations do not claim a Tate chart on a fixed
\((q,v)\)-bidisc at the cusp; SCP45.4 is the source chart there.

The exterior and reference are first chosen over \(\zeta=v^2\).
The marked source identity reverses the two ordered branches in
the complete smooth-collar Tate uniformization. It fixes \(Q\),
sends \(Z\) to \(Z^{-1}\), and sends \(P\) to \(P^{-1}\).
The origin \(O\) and the original marked section determine these
representatives, and the subsequent height-one period reindex
is retained explicitly below. The complete product relation
\(\Theta_Q(z^{-1})=-z^{-1}\Theta_Q(z)\) gives
\[
 f_{P^{-8}}(Z^{-1})=P^8f_{P^8}(Z).
 \tag{SCP45.16}
\]
This retains every higher \(Q\)-coefficient. Thus
\[
 I:(Z,V_{\rm pt})\longmapsto(Z^{-1},P^8V_{\rm pt}),\qquad
 h_{\rm target}I:(Z,V_{\rm pt})\longmapsto(Q/Z,V_{\rm pt}),
 \tag{SCP45.17}
\]
because the actual target line period is \(P^{-8}\).
Each map has its literal reciprocal inverse at the target source
point. On the nodal frame, reversing \(c\) gives
\[
 (-c-S_m)=P^{-m}(c-S_m),\qquad
 C_8(-c)=C_8(c).
 \tag{SCP45.18}
\]
This verifies both point and tensor frame compatibility. The full
source, normal, line, height and reflection arrows SC45.1-\/-4 pass
this independent read. The accompanying executable derives these
identities directly from SCP45.1 and the original tangent functions;
it imports no author\textquotesingle s checker.

% End expanded verified cp36 input

\par\bigskip\noindent\textbf{Full supporting proof 4: Complete independent root source, cap and fan proof.}\par\nobreak

% Expanded verified cp36 input: notes/part3_os45_selected_cusp_root_peer.tex
This proof checks the complete SC45 and C45 construction on its original
source. The fixed line is \(\mathcal M_8=\mathcal O(8G-O)\), the
suspension source map is \(-G\), and the ordinary cap has order one
and norm \(+1\). Every original source, line and annular reference is
retained. The parameter is the positive ordered \(v\)-circle at one
of the three points \(u_0^3=-3/4\). The proof gives the initial
physical pair, the full vertical comparison, and the actual algorithm
on every selected basis fan. The separate source loop for a shorter
original-\(a\) meridian is not included in this scope.

\par

\medskip\noindent\textbf{CR45.1. Original coordinates, the source disc and the closing}

\par\nobreak

Start with the actual cubic
\[
 y^2=x^3-3fx+2g,\quad
 R=t^2+a,\quad f=R^2+t,\quad
 g=R^3+\frac32tR+\frac38,\quad G=(R,\sqrt3/2).
 \tag{CR45.1}
\]
The ordered selected-root map and its inverse parameter chart are
\[
 \begin{gathered}
 a=-u/3-u^4/9,\quad t_c=1/(2u)-u^2/3,\\
 \xi=1/(4u^2)+u/3,\quad c=v/(2u),\quad 4u^3+3=v^2,\\
 u^3-u^2c^2+3/4=0,\qquad v=2uc.
 \end{gathered}\tag{CR45.2}
\]
The last equation has derivative \(3u_0^2\ne0\) in \(u\).
It specifies a holomorphic inverse with \(u(0)=u_0\).
Consequently \(u,a,t_c\) are even in \(c\), and \(3\xi=c^2\).
Substitution of \(x=X+\xi,t=q+t_c\) gives, with no base
coefficient removed,
\[
 y^2=X^2(X+c^2)+qH,\qquad
 H=d\{-6R_c(u+X)+3t_c\}
 +3qd^2(2R_c-\xi-X)+2q^2d^3-3(u+X)+3qd,
 \tag{CR45.3}
\]
where \(d=2t_c+q,R_c=\xi-u\).
Its value at \(q=X=0\) is \(-2u^4\); at the central parameter
it is \(3u_0/2\ne0\). The implicit function \(q(X,y,c)\)
and \(x=X+\xi(c),t=q+t_c(c)\) are actual reciprocal source
charts. In particular the total source and projection to \(c\)
are smooth there.

The other height critical point is obtained by
\(X=c^2r,y=0,q=c^6j\). Its equations are
\[
 r^2(r+1)+jH(c^6j,c^2r,c)=0,\qquad
 3r^2+2r+c^2jH_X(c^6j,c^2r,c)=0.
 \tag{CR45.4}
\]
At \(c=0\), \(r=-2/3,j=-4/(27H_0)\).
The derivative matrix in \((r,j)\) is
\(\left(\begin{smallmatrix}0&H_0\\-2&0\end{smallmatrix}\right)\).
Thus the difference of the two actual critical values has
order six in the ordered parameter. The weighted chart
\((X,y,q)=(c^2\widetilde X,c^3\widetilde Y,c^6Q)\)
keeps the full \(H\) in its equation and has the exact reciprocal
coordinates for \(c\ne0\). The adjacent double-cover vanishing
arcs meet once. The twelve positive half-twists conjugate their
positive Dehn twists by \((g_1g_2)^{-6}\).
The weighted-pencil collar supplies the opposite inner-to-outer
boundary twist \((g_1g_2)^6=T_\partial\). These actual two
directions agree on their common collar. The logarithm of \(H/H_0\)
is a full extending unit there and transports that collar by its
boundary velocity.

The original parameter map has
\[
 a-a_0=
 -\frac{v^4(u^2+2u_0u+3u_0^2)}
 {144(u^2+u_0u+u_0^2)^2},\qquad a_0=-u_0/4.
 \tag{CR45.5}
\]
Its nonzero leading coefficient is \(-1/(216u_0^2)\).
Thus the ordered loop has degree four in the original \(a\)-plane.
We keep that degree throughout.

The exterior and ordinary cap are defined on the entire
\(\zeta=v^2\) disc. The ordinary fibre at \(t=0\) has
discriminant coefficient \(153/16\) and is separated from
\(t_c(0)=3/(4u_0)\).
At infinity the full functions are \(Q_\infty=z^8\kappa_\infty\),
\(P_\infty=z^3\pi\), with central coefficients
\(-a/2304,-1/24\). The point-line factor remains
\[
 f_{Q^3b_0}(Z)=(-Z/b_0)^3Q^{-6}f_{b_0}(Z),\qquad
 b_0=P_\infty^8/Q_\infty^3,\quad b_0(0)=-256/27.
 \tag{CR45.6}
\]
The period change has actual words
\(H=h^{-1}B^{-3},K=h^3B^8\), with inverse
\(h=H^8K^3,B=H^{-3}K^{-1}\).
The ordinary cap arrows are
\[
 \Psi_\alpha(\omega,t,c)=
 \omega+\frac{\log_\alpha t}{2\pi i}\gamma(t,c,\lambda),
 \qquad
 \Psi_\alpha^{-1}(\omega,t,c)=
 \omega-\frac{\log_\alpha t}{2\pi i}\gamma(t,c,\lambda),
 \quad \Psi_\beta=B^k\Psi_\alpha .
 \tag{CR45.7}
\]
The full period \(\gamma\) includes \(-G\), its Miller coefficient,
the physical scalar and the prescribed integer reference.
These arrows and their target-point inverses are holomorphic in
\(\zeta\). Hence the positive ordered loop closes by the actual
identity in the original exterior and cap.

\par

\medskip\noindent\textbf{CR45.2. The complete orbit trace and the lifted link}

\par\nobreak

The common small scalar gives a strictly contracting full lift
on the compact protected exterior and cap. Its norm shell is
a compact fundamental region. A physical unit circle element
equal to \(B^n\) has \(n=0\), by that same strict norm inequality;
there are no additional stabilizers. In the selected chart the
actual fixed edge is
\[
 q=ZAB,\quad Z'=Z^{-1},\quad A'=ZA,\quad B'=ZB,
 \qquad \operatorname{wt}_\delta(A,B)=(-1,+1).
 \tag{CR45.8}
\]
The original K10 sphere is transported to this source along
the connected marked open cover by the full equations and
cap maps CR45.1-\/-7. Forward and reverse proper interval flow
give its diffeomorphism. Common blowup and ruling maps carry
the topology to the chosen fan.

The fixed sphere has complex normal
\(\mathcal O(-1)\oplus\mathcal O(-1)\) with the displayed
opposite circle weights. Conjugating the negative summand
gives the determinant-one spinor
\(\mathcal O(1)\oplus\mathcal O(-1)\); its SU(2) clutch
contracts through the first columns
\((r e^{i\phi},\sqrt{1-r^2})\).
The actual normal boundary is therefore the product Hopf map
\(S^2\times S^3\to S^2\times S^2\).
Alexander duality and codimension four give
\(\pi_1E=0,H_*(E)=H_*(S^3)\) for its exterior.
The primitive boundary \(H_3\) map and the Hopf Gysin sequence
give the orbit exterior cohomology
\(H^2=\mathbb Z,H^3=H^4=0\).
Its original fixed-tube filling has degree-two matrix
\(\left(\begin{smallmatrix}0&1\\-1&0\end{smallmatrix}\right)\).
These are precisely the proved F8 fixed-pair hypotheses,
so its orbit pair is \((S^5,S^2)\) in the stated Hopf marking.
The surgery filling used below is separately
\(D^3\times S^2\), with this same line.

Near the coalescing pair the source translation has a
holomorphic logarithm in the smooth group: divide its central
additive value by a sufficiently large integer, continue the
small division section, take the local logarithm and multiply
back by that integer. The lift \(B/q\) has its full
line-extension logarithm on the same protected source.
Exponentiating real multiples, with opposite multiples as
inverse, carries the original source and line to its F9
relative orbit model. In its two real coordinates
\[
 2b=r-|y_0|^2/r,\qquad
 2b'=|x_0|^2r-r^{-1},\qquad
 \frac{db'}{db}=
 \frac{|x_0|^2+r^{-2}}{1+|y_0|^2r^{-2}}>0.
 \tag{CR45.9}
\]
F9\textquotesingle s proper increasing fibre coordinate, constructed by a
positive complete vertical field and fixed on the outside
collar, is an actual product coordinate with inverse flow.
It contracts the fixed interval into an arbitrarily short
marked circle arc. Its derivative is triangular with positive
last entry; contraction of that triangular factor preserves
the relative stable framing.

The fractional source translation on the normalization has
\[
 T_t(S,c)=
 \frac{(1+g_t)S+(1-g_t)c}
 {((1-g_t)/c)S+1+g_t},\qquad
 g_t=1+ch_t,\qquad
 T_t-S=\frac{h_t(S^2-c^2)}{2+h_t(c-S)} .
 \tag{CR45.10}
\]
Its inverse replaces \(g_t\) by \(g_t^{-1}\).
The determinant of its Möbius matrix is \(4g_t\).
It fixes \(S=\pm c\) and has derivative tending uniformly
to identity as \((S,c)\to0\), including its bounded height
derivative. On a smaller ball the error has norm less than
one half, giving the actual relative source germ isotopy.

The projected fixed trace is now exactly
\[
 f(S,c)=(S^2-c^2,S(S^2-c^2),c),\qquad
 h_f(S,c)=\frac{\operatorname{Re}(S\bar c)}
 {\sqrt{|S|^2+|c|^2}}.
 \tag{CR45.11}
\]
The derivative minors generate \((S,c^2)\). Equal source
images at distinct points are exactly \(S=c,-c\).
Their heights have opposite signs and difference
\(\sqrt2|c|\). A convex height interpolation in the actual
short circle lift therefore remains embedded. The three
summands of \(\|f(rS,rc)\|^2\) scale by \(r^4,r^6,r^2\),
and are not all zero away from the origin. Thus its inner
preimage is an actual star-shaped four-ball with the
complex-domain boundary orientation.
The graph meets only the cylindrical face of the rounded
ambient seven-ball. Its boundary in the F55 oriented
stereographic chart is literally the F55 link and ordered
height, with projection-normal plus increasing-height frame.
The proved link values consequently apply by this framed
map: complex Smale integer \(-2\), Haefliger integer \(+1\).
The source, cap and normal maps below supply their
transfer to this specific OS45 family.

\par

\medskip\noindent\textbf{CR45.3. Original determinant, point modules and physical duality}

\par\nobreak

The full original source form is
\[
 \frac{dt\wedge dx}{y}\wedge dc
 =-\frac{2}{H+qH_q}\,dX\wedge dy\wedge dc.
 \tag{CR45.12}
\]
Indeed \(q_y=2y/(H+qH_q)\); the parameter shifts in \(x,t\)
vanish only after the displayed wedge with \(dc\).
On the normalization let \(x_-=S-c,x_+=S+c\).
The cross product of the two source derivatives is
\[
 n=(3S^2-c^2)dX-2S\,dy+2c(S^2-c^2)dc.
 \tag{CR45.13}
\]
The complete normal and determinant arrows are
\[
 \begin{gathered}
 dq=-x_-x_+\,n/H,\quad
 d(AB)=-x_+^2n/H,\quad d(A'B')=-x_-^2n/H,\\
 \mathcal O(-D_+-D_-)\longrightarrow K_N,\qquad
 x_-x_+\longmapsto dS\wedge dc.
 \end{gathered}\tag{CR45.14}
\]
At \(z=1/S\) the residue is
\(-dz\wedge dc/(1-c^2z^2)\).
This gives the exact normal type on the entire
\(\mathbb P^1\)-normalization. The common multiplier
\(-1/H\) is an extending unit with a specified logarithm
and square root on the inner ball. Its inverse and the
triangular graph shear give the complete normal-frame
comparison there.

For the source point set \(d=S_G=-\sqrt3/(2u)\) and
\(S_m=c(1+P^m)/(1-P^m)\).
The actual doubling functions and negative source arrow are
\[
 f_m(S)=\frac{(S-d)^m}{S-S_m}\ (m=2,4,8),\quad
 T(S)=\frac{dS-c^2}{d-S},\quad
 T^{-1}(S')=\frac{dS'+c^2}{d+S'}.
 \tag{CR45.15}
\]
The original leading coefficient at \(O\) is one.
Direct cancellation gives
\(f_2(S)(T(S)-d)/(S-d)=-2d\). Consequently
\[
 B_{8,1}(S)=(-2d)^8\frac{S-S_8}{T(S)-S_8},\qquad
 W=(S-S_8)v_M,\qquad W\longmapsto(-2d)^8W .
 \tag{CR45.16}
\]
Here \((-2d)^8=81/u^8\). The coefficient \(S-S_8\)
cancels the actual zero and pole point modules at
\(8G,O\), including those points; it and its reciprocal
are line-module maps, rather than nonzero functions on
unrelated trivial bundles.
The entire point and tensor coefficients are
\[
 V_{\rm pt}=W/(c-S_8),\qquad
 V_{\rm tensor}=W/(c-S_G)^8,\qquad
 V_{\rm tensor}=\frac{c-S_8}{(c-S_G)^8}V_{\rm pt}.
 \tag{CR45.17}
\]
Both displayed base factors are extending units.
On a line-arrow coefficient \(g\), the normal map is
\((A,B)\mapsto(g^{-1}A,gB)\) with reciprocal inverse.
Its positive-spinor phases coincide.

The determinant-one spinor is
\[
 E=\mathcal O(D_+)\oplus\mathcal O(-D_-),\qquad
 (1/x_-,x_+),\qquad
 \det E\xrightarrow{\;x_-/x_+\;}\mathcal O.
 \tag{CR45.18}
\]
Its adjoint is the increasing real height plus the
complex normal from CR45.14.
The exact coefficient arrow to the pulled-back
dual-tautological and tautological frames is
\(\operatorname{diag}(1/x_+,x_+)\), with the reciprocal
formula on the other chart.
Thus the actual right-Cox map after positive norm
rescaling is
\[
 q_H=(1+jZ)(A+j\bar B),\qquad
 (a,b)\longmapsto
 \frac1{\sqrt{|x_-|^2+|x_+|^2}}
 \begin{pmatrix}x_+&-\bar x_-\\x_-&\bar x_+\end{pmatrix}
 \overline{(a,b)}.
 \tag{CR45.19}
\]
Its anti-linear map reverses the scalar character.
For \(J=\left(\begin{smallmatrix}0&-1\\1&0\end{smallmatrix}\right)\),
\(\bar g=JgJ^{-1}\) on every SU(2) transition.
The global anti-linear arrow
\((a,b)\mapsto(\bar b,-\bar a)\) commutes with every
clutch and extends both cap discs. Its real
four-dimensional determinant is positive; its adjoint
normal action is the negative identity on three
dimensions. The literal F9 circle-first meridian
calculation gives degree \(-1\) for the nonzero-vector
physical line \(M\). Its Hermitian dual \(U\) has
degree \(+1\), with precisely this extending duality.

Once the normal product is matched, a residual
weight-preserving gauge is
\((A,B)\mapsto(gA,g^{-1}B)\).
Its positive norm contracts by positive rescaling.
Its phase acts as \(q_H\mapsto q_Hg/|g|\); the Hopf
normal point is fixed because this phase commutes
with the chosen quaternionic axis.
Any parameter winding extends over \(D^3\times S^1\)
by the actual isomorphism on \(H^1\) of its boundary.
The inner unit logarithms extend on the inner ball.
There is no unassigned normal clutch in these arrows.

\par

\medskip\noindent\textbf{CR45.4. Both actual caps and the initial physical integer}

\par\nobreak

For each outer \(c\ne0\), use
\[
 A_c=\begin{pmatrix}1&-c\\1&c\end{pmatrix},\qquad
 \mathcal B_c(P)=
 \frac{A_cPA_c^*}{\operatorname{tr}(A_cPA_c^*)},\qquad
 (P,q_H,c)\longmapsto(\mathcal B_c(P),q_H,c).
 \tag{CR45.20}
\]
Positive semidefinite trace-one Hermitian matrices form
the closed three-ball, and its inverse uses \(A_c^{-1}\).
The boundary map is \(Z=(S-c)/(S+c)\), while the
actual Cox quaternion and its Hopf line are held fixed.
The full differential is
\[
 x_+^2dZ=2c\,dS-2S\,dc,\qquad
 x_+^2d_{\rm rel}Z=2c\,d_{\rm rel}S,\qquad
 \det A_c=2c .
 \tag{CR45.21}
\]
The first formula retains the parameter column, whose
triangular homotopy is explicit. The last factor has
winding one and remains in the coefficient matrix
CR45.19. The source-period factor
\[
 (Q/q)|_{q=0}=\frac{2u^{10}}{v^6}
 \tag{CR45.22}
\]
has winding \(-6\) on its shrinking Tate collar and
remains in the adjacent period chart. The normal
coordinate in CR45.14 and the held cap CR45.20
use the original \(q\) throughout.

The complement of the removed inner ball in
\(S^5\times D^2\) has \(H^2=H^3=0\).
The codimension-three Thom sequence of the connected
trace \(W\) makes the complement line class exactly
its normal-meridian integer. The outer complement
has the same generator, by the Wang sequence for
the proved unknotted sphere pair.
Therefore the actual \(U\) with meridian value
one extends over this bordism complement.
On its normal boundary both the base and meridian
summands are fixed by the outer restriction and
CR45.19. A nonzero section of the trivial ratio
line supplies the actual line isomorphism; its
outer parameter gauge extends by the \(H^1\)
isomorphism just stated, and its inner gauge
extends since \(H^1(S^3)=0\).

The inner projection-normal ball and the outer
product cap close the base to
\[
 (S^2\times D^2)\cup(D^3\times S^1)
 =\partial(D^3\times D^2)=S^4.
 \tag{CR45.23}
\]
The first piece has its complex orientation,
the second the negative product orientation.
The spinor is trivial over \(N\simeq S^2\);
choose its global frame to agree with the extending
Cartier frame on the inner ball.
The first column of the comparison CR45.19 is
\((x_+,x_-)\). Its linear map from \((S,c)\)
has complex determinant \(-2\) and hence
complex-oriented boundary degree one.
The comparison extends on \(W\), so the outer
degree is that same inner positive-domain degree.
For the first-column orientation, the tangent
matrices \(I=\operatorname{diag}(i,-i),J\) above,
and \(K=\left(\begin{smallmatrix}0&i\\i&0\end{smallmatrix}\right)\)
have \(IJ=-K,\operatorname{tr}(IJK)=2\).
It follows that the full outer integral of
\(\operatorname{tr}(dg\,g^{-1})^3\) is \(24\pi^2\).
A zero connection in the first frame is
\(-dg\,g^{-1}\) in the cap frame. Its Chern-\/-Simons
form is one third of that cubic form, while the
cap boundary orientation is negative. Thus
\[
 c_2(\bar E)[S^4]=-1,\qquad
 p_1(\operatorname{ad}\bar E)[S^4]=4.
 \tag{CR45.24}
\]
The latter follows also from the actual circle
splitting with spinor roots \(a,-a\) and adjoint
Euler root \(2a\).
For its sphere bundle the degree-one line \(x_U\)
satisfies \(4x_U^2=\pi^*p_1\) and \(\pi_*x_U=1\).
This cap has the individual pair \((1,4)\).

For the inner lifted-link surgery, the actual
projection-normal plus height frame of CR45.11
is the F55/K42 frame. Its two signed Seifert
identities are \(3\sigma+e^2=-4\) and
\(e^2-\sigma=8\), giving \(\sigma=-3,e^2=5\).
The truncated surface caps by the constant
\(-\nu_1\) section of \(D^4\times S^2\).
Its oriented normal plane is the positive
meridian plane, so it is the dual four-cycle
of \(x_U\). Its pair is \((5,-4)\).
The actual outward bordism boundary is
\(-P_f\sqcup-N_e\sqcup-B\); Stokes gives
the positive-family \(U\) pair \((-6,0)\).
The first-to-second Jin parameter reverses it,
and the exact Hermitian duality CR45.19 reverses
the line. Hence
\[
 (k_U,\ell_U)_{\rm Jin}=(6,0),\qquad
 (k_M,\ell_M)_{\rm Jin}=(-6,0).
 \tag{CR45.25}
\]
The physical Hopf marking uses the actual exterior
\(D^3_y\times S^3_q\), the sphere
\(-\partial(D^3_y\times D^4_q)\), and
\((u_M,y_M)=(q,-y)\). Its two pieces are
the negatives of the second and first pieces
of \(\Sigma_F\). Reordering the relation
inverts \(F\). Thus
\(\lambda_M(F)=-\Sigma_{F^{-1}}=\Sigma_F\);
the dual line has \(\Sigma_{F^{-1}}\).
The K68 free and fixed faces have the same
source, Hopf and line maps, so the physical
return is \(-1\pmod{28}\) and the native
pair is \((0,-6)\). Every sign has been
carried by its displayed source or cap map.

\par

\medskip\noindent\textbf{CR45.5. Full vertical comparison}

\par\nobreak

Retain the entire periods
\[
 h=(\kappa qZ,\alpha V),\quad B=(gZ,\eta qV),\quad
 g=P^{-1},\quad\alpha=P^8,\quad\eta=\lambda\beta .
 \tag{CR45.26}
\]
Their ordered-parameter windings are \((-6,0,0,0)\).
The character and period maps at \(T_j\) are
\[
 V_j=Z^{-j}V,\quad H_j=hB^j,\quad
 Q_j=\kappa qg^j,\quad
 \alpha_j=\alpha\eta^j\kappa^{-j}g^{-j^2},\quad
 \eta_j=\eta g^{-j}.
 \tag{CR45.27}
\]
Their inverse has \(V=Z^jV_j,h=H_jB^{-j}\).
The source flow \(Q_t=Q\exp(tj\log g)\) uses
the full periodic logarithm. On a node
\(xy=Q_t\), its field is
\(\partial_t+j\log g\,y\partial_y\);
on the outside collar use the prescribed
real period map. Their difference is vertical,
so a partition of unity gives a relative
horizontal field. Its reverse flow supplies
the inverse and retains every source unit.

The complete Cartier map is
\[
 E_j=E\otimes\mathcal O(j(D_+-D_-)),\qquad
 (x_-^{-j-1}x_+^j,x_-^{-j}x_+^{j+1})
 \xrightarrow{\;Z^j\;}(1/x_-,x_+).
 \tag{CR45.28}
\]
It preserves the determinant trivializations
and carries the adjoint normal over the
transported ball cap.
Use full real period coefficients
\[
 (z,\ell)=u(1,0)+b(\tau,\mu)+d(0,1)
                         +\gamma(G,\mathcal B),
 \quad(\tau,\mu,G,\mathcal B)
 =(\log(\kappa q),\log\alpha,\log g,\log(\eta q))/(2\pi i).
 \tag{CR45.29}
\]
They give exactly
\[
 \begin{gathered}
 (u,d)\mapsto(u-b,d-\gamma)\quad(q\text{ deck}),\\
 (u,d)\mapsto(u+6b,d)\quad(c\text{ deck}),\\
 (u_j,b_j,d_j,\gamma_j)=(u,b,d-ju,\gamma-jb).
 \end{gathered}\tag{CR45.30}
\]
Differentiation retains all coefficient derivatives.
The invariant physical connections are
\[
 A_0=dd+\alpha_qd\gamma,\quad
 A_j=dd_j+\alpha_qd\gamma_j+6j\theta\,db,\quad
 \beta_j=-j\{du+(\alpha_q-6\theta)db\}.
 \tag{CR45.31}
\]
The theta-cut line \(\mathcal T\) has coefficient
\(e^{2\pi ib}\) and curvature
\(\Xi=d\theta\wedge db,\Xi^2=0\).
The deck square gives the full relative line
\(U_j\cong U_0\otimes\mathcal T^{6j}\).
Its potentially nonextending \(q\)-character
is fixed by the actual free section disc
\[
 (q/Z_*,Z_*,1)\longmapsto
 (qZ_*^{j-1},Z_*,Z_*^{-j}).
 \tag{CR45.32}
\]
Choose \(Z_*\) near \(-1\), away from the
original point and zero modules. This map,
its inverse and the actual line coefficient
extend across \(q=0\), giving zero
\(q\)-winding. At another free circle lift
the full factors involving \(\exp(\gamma\log g)\)
cancel between the two coefficients and
remain extending units.
The \(q\)-commutator kills \(u\) primitively,
and the parameter commutator is \(6u\);
the remaining outer \(H_1\) generators
\(\theta,b,\gamma\) extend over the completed
piece. Removing their characters leaves a
real boundary logarithm. The explicit
partition-of-unity solution of its transition
cocycle and a collar extension of that
logarithm supply the desired nonzero
ratio-line section and its inverse.
This constructs the full relative line
arrow with its prescribed boundary value.

For the tangent comparison put
\(\omega=du+\operatorname{Re}\tau\,db+
\operatorname{Re}G\,d\gamma\).
Then \(\tau_j=\tau+jG\) and
\(\gamma_j=\gamma-jb\) give \(\omega_j=\omega\).
The coframe change is
\(\left(\begin{smallmatrix}1&0&0\\0&1&-j\\0&0&1\end{smallmatrix}\right)\).
All parameter-to-source derivatives are in
its triangular off-diagonal blocks. Their
curvature-square trace is zero; the remaining
base block has two angles and one collar
variable, so its four-form is zero.
This is the complete relative zero
Pontryagin transgression.
Make the source circle map \(b\) constant
near the surgery tube by
\(b_t=b\exp(2\pi it\chi(b_*-\widetilde b))\),
where the real lift is on the indicated
neighbourhood and the cutoff is fixed outside.
A common regular level gives a relative
cobordism to a level missing all surgery
changes. The initial source three-level
times circle has zero source-pulled
Pontryagin integral. Stokes with the
just-computed boundary transgression
therefore gives \(\ell_{M,j}=0\).

On the boundary the complete forms are
\[
 F_{{\rm ext},j}=d\alpha_q\wedge d\gamma
                     +j\,d\alpha_q\wedge db,\qquad
 d\beta_j=-j\,d\alpha_q\wedge db
                     +6j\,d\theta\wedge db.
 \tag{CR45.33}
\]
In the positive outward order
\((\alpha_q,u,b,\gamma,\theta)\), direct
exterior multiplication gives
\[
 3\int_0^1\beta_j(F_{{\rm ext},j}+s\,d\beta_j)^2\,ds
 =18j^2\,d\alpha_q\,du\,db\,d\gamma\,d\theta .
 \tag{CR45.34}
\]
The actual outward collar runs from the
piece to the exterior and reverses that
integral. The interior cube change is
affine in \(j\), since \(\Xi^2=0\).
Thus \(k_{M,j}=-6+Lj-18j^2\).
The full theta identity
\(f_{P^{-8}}(Z^{-1})=P^8 f_{P^8}(Z)\)
gives the reflected map
\((Z,V)\mapsto(Z^{-1},P^8V)\).
Following it by the actual target source
period gives \((Q/Z,V)\).
It preserves the even exterior and cap
reference and translates height by one,
carrying \(T_0\) to \(T_1\) with the held
product caps. Therefore \(k_{M,1}=k_{M,0}\)
and \(L=18\). We obtain
\[
 (k_M,\ell_M)_j=(-6-18j(j-1),0),\quad
 A_j=-1-3j(j-1),\quad
 (\mu_j,\nu_j)=(0,-6-18j(j-1)).
 \tag{CR45.35}
\]

\par

\medskip\noindent\textbf{CR45.6. Every fan with the full height inverse}

\par\nobreak

Let \(P\in GL_2(\mathbb Z)\), \(R=P^{-1}\).
The complete multiplier matrix and its winding are
\[
 A_{ik}=(\kappa^{P_{1k}}g^{P_{2k}})^{R_{i1}}
        (\alpha^{P_{1k}}\eta^{P_{2k}})^{R_{i2}},
 \qquad B_P=-6R e_1e_1^tP .
 \tag{CR45.36}
\]
The original \(q\)-valuation is \(RP=I\).
The three edge rows are
\[
 (A_{11}A_{21}bc,A_{21}^{-1}ab,A_{11}^{-1}b^{-1}),\quad
 (A_{12}A_{22}bc,A_{22}^{-1}c^{-1},A_{12}^{-1}ac),\quad
 (a^{-1},ab,ac).
 \tag{CR45.37}
\]
Every row has product \(q\). Their actual native
coordinates give \(z'=z^{-1},n'=zn\).
The ordinary second-vertex map retains the full
height function
\[
 (u,v,w)\mapsto
 (u,A_{22}(s,q)v,A_{11}(s,q)w),\qquad
 h_s(q)=qA_{11}(s,q)A_{22}(s,q).
 \tag{CR45.38}
\]
Its inverse first solves \(q=h_s^{-1}(uv'w')\)
and then divides by the actual coefficients at
that recovered original \(q\).
Its zero derivative is the nonzero product of
the two units. On the compact parameter annulus
a uniform smaller disc makes the normalized
derivative differ from one by less than one half;
integration on a straight segment proves
injectivity there. The inverse function theorem
then gives the stated uniform inverse.
The ordinary normal blocks are
\[
 D_0=\operatorname{diag}(A_{22}/A_{21},A_{11}A_{21}),
 \quad D_1=\operatorname{diag}(A_{11}/A_{12},A_{12}A_{22}),
 \quad D_2=\operatorname{diag}(A_{22},A_{11}).
 \tag{CR45.39}
\]
Their first windings are
\(l=(B_{22}-B_{21},B_{11}-B_{12},B_{22})\);
every determinant winding is \(-6\).

The zero-winding diagonal logarithms \(k_{i0},k_{i1}\)
are carried by the cutoff map
\[
 (z,n_0,n_1)\mapsto
 (z,e^{\rho k_{i0}(s,q)}n_0,e^{\rho k_{i1}(s,q)}n_1),
 \quad H_{s,z}(q)=q e^{\rho(k_{i0}+k_{i1})(s,q)}.
 \tag{CR45.40}
\]
Its inverse first recovers \(q\) from
\(H_{s,z}(q)=zn'_0n'_1\), then multiplies
by the reciprocal coefficients evaluated at
that \(q\). Compactness and the same derivative
bound give its smooth radial inverse.
Thus every original height derivative is
retained in the ordinary products.
The winding cylinder is the explicit quaternion
\[
 H_d(s,\theta)=
 (\sin\vartheta(s)e^{id\theta}+j\cos\vartheta(s))
 (\sin\vartheta(s)+j\cos\vartheta(s))^{-1}.
 \tag{CR45.41}
\]
On each edge its map is \(Q\mapsto
H_{-3}QH_{l_i+3}\). It is based, periodic
and has the exact two boundary normal
phases. Its Hopf arrow has left degree
\(-6\), right degree zero, hence inverse
ordinary/native cylinder \(L^6\).
Normal rescaling followed by the cutoff
velocity flow carries the full higher jets
and inverse, relative to both protected collars.
The half-angle Spin endpoints on both factors
are \((-1)^{w_i}\), so their total is one.

For \((x,y)^t=P^{-1}e_2\), put
\[
 w=(x,y,-x-y),\quad d=(-x-2y,-2x-y,x-y),\quad
 I_i=-3w_i+d_i(l_i+3),\quad
 J_i=-3w_i-d_i(l_i+3).
 \tag{CR45.42}
\]
The actual circle and derivative prisms obey
\[
 V(w_ie_i-w_0e_0)=I_ie_i-I_0e_0,\qquad
 w_iD_i-w_0D_0=J_i-J_0,\quad i=1,2.
 \tag{CR45.43}
\]
The complete vertical ordinary seed, through
CR45.35 and the literal \(L^6\) cylinder, is
\[
 \begin{gathered}
 V=\begin{pmatrix}
 18j(j-1)&6-12j&12j-6\\
 6-12j&0&0\\12j-6&0&0
 \end{pmatrix},\\
 D=(-12-18j(j-1),12j-12,-12j),\qquad S=-1-3j(j-1).
 \end{gathered}\tag{CR45.44}
\]
The identical-fan maps act by the identity
on absolute character semigroups, including
their full coefficients and original cap
overlaps. Euclidean subtraction gives a
unique seed and a finite word in
\(Q_0=\left(\begin{smallmatrix}1&-1\\0&1\end{smallmatrix}\right)\),
\(Q_1=\left(\begin{smallmatrix}1&0\\-1&1\end{smallmatrix}\right)\).
Their actual dense character arrows are
\((ZV,V),(Z,ZV)\), and their native first
triples are \((a/c,bc,c),(a/b,b,bc)\).
The other-ruling charts come from the common
rank-one matrix
\(\left(\begin{smallmatrix}n_0&n_1\\zn_0&zn_1\end{smallmatrix}\right)\).
Their inverse maps are the reciprocal charts
of the common projective blowup. The full
coefficients transform by
\[
 A'_{ik}=\prod_{h,l}A_{hl}^{(Q^{-1})_{ih}Q_{lk}}.
 \tag{CR45.45}
\]
The common proper blowup and the two ruling
projections carry the integral F35 maps and
the sphere topology before any return
number is evaluated.

At the selected old edge set \(a=V_{ii},b=D_i\).
The full Cox common-link map is
\((h,q)\mapsto(q^{-1}hq,q)\); its carried
complex reset sends
\((\mu,\nu,S)\mapsto(\mu+2\nu,-\nu,S+\nu)\).
The actual native right integer is
\((b-a)/2\), with the same \(L^6\) cylinder
on both sides. Thus the target conditions are
\[
 V'_{22}=a,\qquad D'_2=2a-b,\qquad
 S'=S+(b-a)/2.
 \tag{CR45.46}
\]
The six symmetric \(V'\) coefficients and
three \(D'\) coefficients are uniquely
determined by the two full prism columns
and these target conditions.
A homogeneous symmetric difference
annihilates the two-dimensional column
plane and is a multiple of
\((1/w_r)(1/w_s)\); the last diagonal
condition kills it. The derivative
difference is a multiple of \((1/w_r)\)
and its last entry kills it.
Existence and integrality are supplied
by the actual target maps and products
above. Fresh rational elimination in
the independent certificate checks
the complete equations on each tested
path; no target formula is imported.
The terminating chart construction
proves the conclusion for every fan.

\par

\medskip\noindent\textbf{CR45.7. Scope and audit result}

\par\nobreak

The full positive ordered OS45 source
loop has the initial physical pair
\((-6,0)\), native pair \((0,-6)\)
and class \(-1\), with vertical
values CR45.35 and all-fan values
CR45.36-\/-46.
The full source maps hold at each
of the three \(u_0\) points.
For a fixed scalar in the original
admitted boundary family, decrease
it along its constant-argument
interval to the common strict bound.
The actual logarithm in CR45.7,26
is carried along that interval;
proper forward and reverse flow
gives the complete annular
comparison with no scalar winding.
This proves the same result for
every admitted fixed scalar.

This peer verifies the complete
original source, point line,
normalization trace, two caps,
physical duality and inverse
clutch, full vertical relative
line, tangent pairing, ordinary
height inverses and all-fan
products. The independent exact
execution supports these maps
and their signs. It assigns no
shorter source endpoint or
additional source-moduli return.

% End expanded verified cp36 input

\par\bigskip\noindent\textbf{Full supporting proof 5: Complete independent source, relative-line, cap and fan proof.}\par\nobreak

% Expanded verified cp36 input: notes/part3_os45_selected_cusp_completion_blind_peer.tex
This proof audits C45.1--8 at the source

\par

\noindent{\small\texttt{a32c16121016f2a255a7269aa91b77c6a40f16761340ebc78acd2a3194f2f6a9}}.

\par

\noindent
It retains the original OS45 Weierstrass equation, the actual line
\(M_8=\mathcal O([8G]-[O])\), translation by \(-G\), the original
order-one ordinary cap at \(t=0\), its complete annular reference,
and a fixed admissible scalar. The ordered parameter is a positive
\(v\)-circle at one of the three points \(u_0^3=-3/4\), with
\(v^2=4u^3+3\). It covers four positive turns in the original
\(a\)-plane. The argument computes this actual ordered loop.

The preceding independent local proof SCP45.1--4, in
\path{notes/part3_os45_selected_cusp_local_blind_peer.md}, gives the
full original polynomial, chronological Miller maps, point and tensor
frames, inverse source translation and smooth-height theta
reflection. Its 68 independent identities remain a separate locked
input. The present proof supplies the complete source-disc,
orbit, normal, cap, relative-line and fan comparisons. Every
coefficient below is a function of the original height; its
displayed value at the node is used only after its extending
module map has been established.

\par

\medskip\noindent\textbf{CP45.1. Original source and the outside closing square}

\par\nobreak

Write
\[
\begin{gathered}
 R=t^2+a,\qquad f=R^2+t,\qquad
 g=R^3+\tfrac32tR+\tfrac38,\\
 y^2=x^3-3fx+2g,\qquad G=(R,\sqrt3/2),\\
 a=-u/3-u^4/9,\quad t_c=1/(2u)-u^2/3,\quad
 c=v/(2u),\quad \xi=c^2/3,\quad q=t-t_c,\quad X=x-\xi.
\end{gathered}
\tag{CP45.1}
\]
Put \(R_c=\xi-u\) and \(d=2t_c+q\). Direct substitution gives
\[
\begin{split}
y^2&=X^2(X+c^2)+qH,\\
H&=d[-6R_c(u+X)+3t_c]+3qd^2(2R_c-\xi-X)\\
 &\qquad+2q^2d^3-3(u+X)+3qd,\\
H(0,0,c)&=-2u^4,\qquad H(0,0,0)=H_0=3u_0/2.
\end{split}
\tag{CP45.2}
\]
The derivative in \(q\) of the right side is \(H+qH_q\), a
unit near the inner origin. Thus \(q\) is the actual implicit
holomorphic function of \((X,y,c)\), and the source projection
to \(c\) is a submersion there. No change of its height is used.
The local cover and its inverse have the exact factorization
\[
a-a_0=-\frac{(u-u_0)^2(u^2+2uu_0+3u_0^2)}9,\qquad
u-u_0=\frac{v^2}{4(u^2+uu_0+u_0^2)}.
\tag{CP45.3}
\]
Its leading coefficient is \(-1/(216u_0^2)\) times \(v^4\).
The holomorphic unit in this formula fixes the positive degree
four of the ordered loop.

For the other nearby critical point set \(X=c^2R_1(c)\),
\(q=c^6J_1(c)\), \(y=0\). Its exact equations are
\[
R_1^2(R_1+1)+J_1H(c^6J_1,c^2R_1,c)=0,\quad
3R_1^2+2R_1+c^2J_1H_X(c^6J_1,c^2R_1,c)=0.
\tag{CP45.4}
\]
At \(c=0\) the solution is \((-2/3,-4/(27H_0))\);
the Jacobian is
\(\left(\begin{smallmatrix}0&H_0\\-2&0\end{smallmatrix}\right)\).
It is invertible, and the unique solutions are even. Hence the
ordered critical-value difference has six positive turns, giving
the braid \(\sigma^{12}\) in the stated half-twist convention.
In the literal weighted chart
\[
X=c^2\widetilde X,\quad y=c^3\widetilde Y,\quad q=c^6Q,
\qquad
\widetilde Y^2=\widetilde X^2(\widetilde X+1)
 +QH(c^6Q,c^2\widetilde X,c),
\tag{CP45.5}
\]
the exceptional pencil is
\([Y^2-X^2(X+R^2):H_0R^6]\). Its six actual base-ideal
charts and its order-two and order-three quotient charts are the
F56 charts with this constant. The analytic logarithm of the
entire quotient \(H/H_0\) carries its boundary phase. The two
adjacent vanishing arcs of the original cubic lift to curves
meeting once. The right Hurwitz arrow is
\((g_1,g_2)\mapsto(g_2,g_2^{-1}g_1g_2)\); six full twists
conjugate by \((g_1g_2)^{-6}\). The branched boundary coordinate
has \(\vartheta=2\widetilde\vartheta\), and the actual weighted
inner-to-outer map is the reciprocal boundary twist
\((g_1g_2)^6\). Their composite fixes the retained outside
marking. The extending logarithm changes this composite through
the boundary vector field of the actual source map, without an
integer reindex.

Choose a fixed finite \(t\)-disc around the two critical values.
Its boundary is a source submersion for the whole closed small
\(c\)-disc. The horizontal source field tangent to that boundary
has a proper flow and the inverse reverse flow. Outside it all
original coefficients are constructed first over \(\zeta=v^2\):
\[
a_0=-u_0/4,\quad t_c(0)=3/(4u_0),\quad
D(0,a_0)=153/16,\quad
Q_\infty=z^8\kappa_\infty,\quad P_\infty=z^3\pi,
\tag{CP45.6}
\]
where \(\kappa_\infty(0)=-a/2304\), \(\pi(0)=-1/24\).
The full infinity point map retains
\[
f_{Q^3b_0}(Z)=(-Z/b_0)^3Q^{-6}f_{b_0}(Z),\quad
b_0=P_\infty^8/Q_\infty^3,\quad b_0(0)=-256/27.
\tag{CP45.7}
\]
It also retains the actual reciprocal period words
\(H=h^{-1}B^{-3},K=h^3B^8\), \(h=H^8K^3,B=H^{-3}K^{-1}\).
The original cap is smooth on this parameter disc. On each of
its fixed sectors its full period \(\gamma\) gives
\[
\Psi_\alpha(\omega,t,c)=
 \omega+\frac{\log_\alpha t}{2\pi i}\gamma(t,c,\lambda),\quad
\Psi_\alpha^{-1}(\omega,t,c)=
 \omega-\frac{\log_\alpha t}{2\pi i}\gamma(t,c,\lambda),\quad
\Psi_\beta=B^k\Psi_\alpha .
\tag{CP45.8}
\]
Here \(\gamma\) is the entire source, physical line and scalar
period of \(B_\lambda=(T,\lambda C_*qB_{8,1})\), including the
prescribed integer cochain; the inverse uses the actual target
period. Each map is defined over \(\zeta\). Therefore the
ordered \(v\)-loop has the literal original identity closing on
this exterior and cap. The height-one reflection used below is
a separate, explicitly written comparison.

\par

\medskip\noindent\textbf{CP45.2. The actual trace and its orbit pair}

\par\nobreak

Choose the scalar below the common contraction bound on the
compact protected source, line and cap pieces. A physical
unit-circle equality with \(B^k\) forces \(k=0\) by that strict
norm inequality. At a vertical fan the only fixed edge has the
actual two charts
\[
q=ZAB=Z'A'B',\quad
Z'=Z^{-1},\quad A'=ZA,\quad B'=ZB,\qquad
\operatorname{wt}_\delta(A,B)=(-1,+1).
\tag{CP45.9}
\]
The other edge cocharacters are primitive for this circle.
The original K10 sphere is transported to the chosen boundary
point through the full O45 source, point-line, infinity and cap
maps. For every prescribed fan, the common blowup, both ruling
projections, and the primitive integral Gysin and van Kampen
maps of F35 give the actual finite fan path from that reference.
This supplies the sphere antecedent before the following orbit
calculation.

For a current vertical fan let \(E\) be the complement of a
fixed \(S^2\) tube in that sphere. Alexander duality gives
\(H_*(E)=H_*(S^3)\), and codimension four gives
\(\pi_1(E)=0\). Conjugating the negative native normal line in
(CP45.9) gives the positive spinor
\(\mathcal O(1)\oplus\mathcal O(-1)\). Its determinant-one
clutch bounds the explicit disc whose first column is
\((re^{i\phi},\sqrt{1-r^2})\). Thus the actual boundary is
the product Hopf map \(S^2\times S^3\to S^2\times S^2\).
The map \(H_3(\partial E)\to H_3(E)\) is a primitive
isomorphism. Fibre integration and Gysin naturality give
\(H^2(E/S^1)=\mathbb Z e\), \(H^3=H^4=0\). Gluing the
original \(S^2\times D^3\) has difference matrix
\(\left(\begin{smallmatrix}0&1\\-1&0\end{smallmatrix}\right)\)
on the boundary base and meridian. Van Kampen and this matrix
give a simply connected homology five-sphere with the stated
embedded fixed sphere. The F8 fixed-pair result applies to
this precise Hopf attachment and gives \((S^5,S^2)\).
The later \(D^3\times S^2\) surgery uses the same Hopf line;
the two fillings have the displayed different attaching maps.

Near the coalescing pair, divide \(-G\) in the cuspidal smooth
group by a large integer, continue the small division section,
take its actual group logarithm and multiply by that integer.
The full invertible lift \(B/q\), including \(C_*,\lambda\)
and its original Cartier modules, has the corresponding
logarithm in the line extension. Exponentiation of a real
multiple and its negative give the source-and-line isotopy and
its inverse. The F9 node coordinates satisfy
\[
2b=r-|y_0|^2/r,\quad
2b'=|x_0|^2r-r^{-1},\qquad
\frac{db'}{db}=
\frac{|x_0|^2+r^{-2}}{1+|y_0|^2r^{-2}}>0 .
\tag{CP45.10}
\]
The positive unit norms are retained. With the original closing
(CP45.8), a nonzero suspension deck would change the strict
contracting norm and is excluded. The graph has a common real
short-arc lift, with \(S=-c\) below \(S=c\) in the outward
F9 coordinate. Its opposite is the coefficient of the contracting
period.

The source projection is the full germ
\[
f(S,c)=(S^2-c^2,S(S^2-c^2),c).
\tag{CP45.11}
\]
For \(P=(S_G-c)/(S_G+c)\), \(S_G=-\sqrt3/(2u)\), its
actual fractional source isotopy is
\[
T_t(S,c)=
\frac{(1+g_t)S+(1-g_t)c}{((1-g_t)/c)S+1+g_t},\quad
g_t=e^{t\log P^{-1}}=1+ch_t,\quad
T_t-S=\frac{h_t(S^2-c^2)}{2+h_t(c-S)}.
\tag{CP45.12}
\]
Both branches are fixed. On a smaller ball the derivative error
from the identity is less than \(1/2\), uniformly in the
bounded F9 coordinate. The source reparametrization and its
inverse therefore remove this shear with its orientation.
The derivative-minor ideal of \(f\) is \((S,c^2)\), and its
only distinct double pair is \((c,c),(-c,c)\). The actual
height
\[
h_f=\frac{\operatorname{Re}(S\overline c)}
             {\sqrt{|S|^2+|c|^2}}
\tag{CP45.13}
\]
has the same strict branch order. Convex interpolation in the
marked short arc preserves the embedded graph away from the
origin and stays fixed on the actual outside graph. The function
\(\lVert f\rVert\) is strictly increasing on every positive
ray: its squared terms have positive powers of the ray
parameter, with at least one nonzero term. Thus
\(B_f=f^{-1}(B_\epsilon^6)\) is the actual star-shaped
four-ball. The graph of a small multiple of \(h_f\) meets
only the cylindrical boundary of the rounded product
\(B_\epsilon^6\times[-2\eta\epsilon,2\eta\epsilon]\).
The F55 stereographic map identifies its boundary with the
literal lifted link of (CP45.11), with projection-normal and
increasing-height frame. Its complex Smale integer is \(-2\)
and its Haefliger integer is \(+1\) by that framed germ.

\par

\medskip\noindent\textbf{CP45.3. Entire point modules and the normal determinant}

\par\nobreak

Put \(x_-=S-c,x_+=S+c\), \(D_+=(x_-=0)\),
\(D_-=(x_+=0)\). The actual normalization has
\[
S_m=c\frac{1+P^m}{1-P^m},\qquad
f_m=\frac{(S-S_G)^m}{S-S_m}\quad(m=2,4,8),\qquad
S_m(0)=S_G(0)/m .
\tag{CP45.14}
\]
These are the chronological original tangent quotients and
their doublings, not separately chosen rational functions.
The genuine point-line coefficient and both theta coefficients
are
\[
W=(S-S_8)v_M,\qquad
V_{\rm pt}=\frac W{c-S_8},\qquad
V_{\rm tensor}=\frac W{(c-S_G)^8}.
\tag{CP45.15}
\]
The zero at \(8G\) and the pole at \(O\) cancel the defining
Cartier generators. Their reciprocals are maps of those modules
also at the two marked points. Both remaining base factors are
analytic units. In particular the full tensor-to-point ratio
\((c-S_8)/(c-S_G)^8\) is retained. A coefficient arrow
\(W=gV\) has native normal arrow
\((A,B)\mapsto(g^{-1}A,gB)\), and ordered positive-spinor
arrow \((\overline g^{-1},g)\). Its common phase is retained
on the physical right Hopf factor, while its positive norm
has the literal reciprocal positive rescaling.

Differentiating the original equation and wedging with \(dc\)
gives
\[
dt\wedge dx/y\wedge dc
 =-\frac2{H+qH_q}\,dX\wedge dy\wedge dc .
\tag{CP45.16}
\]
The unit includes all original height derivatives. The
normal covector and its two native products are
\[
\begin{gathered}
n_f=(3S^2-c^2)dX-2S\,dy+2c(S^2-c^2)dc,\\
dq=-\frac{x_-x_+}{H}n_f,\quad
d(AB)=-\frac{x_+^2}{H}n_f,\quad
d(A'B')=-\frac{x_-^2}{H}n_f,\quad
H=H(0,x_-x_+,c).
\end{gathered}
\tag{CP45.17}
\]
Their common \(-1/H\) has an extending logarithm and square
root on the inner source ball. Its full normal adjustment and
inverse agree with (CP45.16), with its displayed constant.
The normal line is
\(L=\mathcal O(-D_+-D_-)\), through the actual residue arrow
\[
f\longmapsto f\,\frac{dS\wedge dc}{x_-x_+}.
\tag{CP45.18}
\]
At \(z=1/S\) the displayed meromorphic form is
\(-dz\wedge dc/(1-c^2z^2)\); hence the comparison includes
the infinity chart. The triangular graph-height derivative is
carried with its positive triangular contraction.

The determinant-one positive spinor and its specified
trivialization are
\[
E=\mathcal O(D_+)\oplus\mathcal O(-D_-),\quad
\det E=\mathcal O(D_+-D_-)
\xrightarrow{\ Z=x_-/x_+\ }\mathcal O .
\tag{CP45.19}
\]
In frames \((1/x_-,x_+)\), the coefficient map to
\(r^*(\mathcal O(1)\oplus\mathcal O(-1))\), for
\(r=[x_-:x_+]\), is \(\operatorname{diag}(1/x_+,x_+)\);
the other chart carries its reciprocal. Its adjoint is
\(\mathbb R_\gamma\oplus L_{\mathbb R}\).
The full positive-spinor/right-Cox map is
\[
q_H=(1+jZ)(A+j\overline B),\qquad
(a,b)\longmapsto
\frac1{\sqrt{|x_-|^2+|x_+|^2}}
\begin{pmatrix}x_+&-\overline{x_-}\\x_-&\overline{x_+}\end{pmatrix}
\overline{(a,b)} .
\tag{CP45.20}
\]
It intertwines the physical action with
\(q_H\mapsto q_He^{-i\theta}\). The Hermitian dual reverses
this character. The anti-linear map
\((a,b)\mapsto(\overline b,-\overline a)\) commutes with
the SU(2) clutches, has positive real determinant and gives
the antipodal adjoint three-plane map. In the F9 chart
\(A=1,B=0\), the coordinate \(B/\overline A\), the outward
normal and the circle-first orientation give physical
meridian degree \(-1\); the dual \(U=M^*\) has degree \(+1\).
These are maps on both cap discs, not a sign change on a
number after integration.

\par

\medskip\noindent\textbf{CP45.4. Held-Cox cap and both Chern evaluations}

\par\nobreak

For \(c\ne0\) on the outer circle let
\[
A_c=\begin{pmatrix}1&-c\\1&c\end{pmatrix},\qquad
\mathcal B_c(P)=\frac{A_cPA_c^*}{\operatorname{tr}(A_cPA_c^*)}.
\tag{CP45.21}
\]
On the trace-one positive Hermitian ball its inverse is the
same expression with \(A_c^{-1}\). Its rank-one boundary
is \(Z=(S-c)/(S+c)\). The actual cap arrow holds the normal
quaternion:
\[
(P,q_H,c)\longmapsto(\mathcal B_c(P),q_H,c),\qquad
x_+^2dZ=2c\,dS-2S\,dc,\quad\det A_c=2c.
\tag{CP45.22}
\]
At fixed \(c\) the differential restricts to
\(x_+^2d_{\rm rel}Z=2c\,d_{\rm rel}S\). The full
parameter column \(-2S\,dc\) is triangular and remains in
the tangent comparison. The determinant has winding one;
the separate source ratio \(Q/q\) retains winding \(-6\).
No derivative of \(\mathcal B_c\) replaces the held normal
quaternion. The full boundary source and line isotopies
extend by their cutoff velocities and physical-circle
averaging, with inverse reverse flow.

The complement of the inner seven-ball in \(Y\times D_c\)
has \(H^2=H^3=0\). The codimension-three Thom sequence for
the connected embedded trace therefore classifies the
extending line by its integral normal-meridian degree.
The outside complement of the unknotted sphere has the same
degree-two generator. Its Wang sequence has no extra
\(H^1\) summand. Consequently the actual \(U\) of degree
one extends with the Cox restriction on \(S(\nu W)\).
Its residual boundary ratio is a circle gauge.
\(H^1(D^3\times S^1)\to H^1(S^2\times S^1)\) is an
isomorphism, so its entire parameter winding extends over
the outer cap. The inner \(S^3\) has \(H^1=0\).
A nonzero section of the ratio line followed by these
extending gauges constructs the required isomorphism and
its inverse, with the full point coefficients.

After the inner projection-normal ball and the held outer
cap are attached, the trace base is
\[
\overline W=(S^2\times D^2)\cup(D^3\times S^1)
 =\partial(D^3\times D^2)=S^4 .
\tag{CP45.23}
\]
The first piece has complex orientation; the second has
negative displayed product orientation. The inner Cartier
frame and determinant extend over \(N\simeq S^2\), where
the SU(2) bundle is trivial. A global frame can be chosen
to agree with that inner frame using its extending
ball gauge. The frame comparison is the matrix in
(CP45.20), with first column \((x_+,x_-)\).
The complex-linear map \((S,c)\mapsto(x_+,x_-)\) has
determinant \(-2\), and positive real determinant \(4\);
its boundary degree is one. The comparison extends on
the intervening trace, so the outer degree equals this
inner positive-domain degree.

For \(I=\operatorname{diag}(i,-i)\),
\(J=\left(\begin{smallmatrix}0&-1\\1&0\end{smallmatrix}\right)\),
\(K=\left(\begin{smallmatrix}0&i\\i&0\end{smallmatrix}\right)\),
one has \(IJ=-K\), \(\operatorname{tr}(IJK)=2\).
Thus in the first-column orientation the cubic integral
is \(24\pi^2\). The zero connection becomes
\(-dg\,g^{-1}\) in the cap frame. Its Chern--Simons
cubic has factor \(1/3\), and the cap boundary has
the negative orientation just specified. It follows that
\[
c_2(\overline E)[S^4]=-1,\qquad
p_1(\operatorname{ad}\overline E)[S^4]=4 .
\tag{CP45.24}
\]
On \(B=S(\operatorname{ad}\overline E)\) the dual line
has \(\pi_*x_U=1\), \(e(T_vB)=2x_U\), and
\(4x_U^2=\pi^*p_1(\operatorname{ad}\overline E)\).
Its individual pair is therefore \((1,4)\).

The inner lifted link has the exact projection-normal
and increasing-height framing of CP45.2. Its Seifert
identities are
\[
-4=3\sigma+e^2,\qquad 8=e^2-\sigma,
\quad\text{so}\quad \sigma=-3,\ e^2=5.
\tag{CP45.25}
\]
Its constant negative-normal cap has positive meridian
plane and is dual to \(x_U\). Its pair is \((5,-4)\).
The actual quotient bordism has outward boundary
\(-P_f\sqcup-N_e\sqcup-B\). Stokes, with the extending
line, gives positive-family \(U\)-pair \((-6,0)\).
Jin reversal gives \((6,0)\), and the full physical
Hermitian dual in (CP45.20) gives
\[
(k_M,\ell_M)=(-6,0),\qquad
(\mu,\nu)=(0,-6),\qquad A_M=-1.
\tag{CP45.26}
\]
The native coordinates here are typed by
\(k_M=\mu+\nu,\ell_M=4\mu\) in this held product.
For the final sphere identification the actual physical
exterior is \(D_y^3\times S_q^3\), with sphere
\(-\partial(D_y^3\times D_q^4)\). The physical marking
is \(u_M=q,y_M=-y\). Its pieces are the negatives of
the second and first pieces of the K68 clutch.
Reordering the first-to-second seam inverts \(F\), so
\(\lambda_M(F)=-\Sigma_{F^{-1}}=\Sigma_F\);
the dual gives \(\lambda_U(F)=\Sigma_{F^{-1}}\).
This literal piece map fixes the class in (CP45.26).

\par

\medskip\noindent\textbf{CP45.5. Vertical source map and the full relative line}

\par\nobreak

Write the full periods on the protected collar as
\[
h=(\kappa qZ,\alpha V),\quad B=(gZ,\eta qV),\qquad
g=P^{-1},\quad\alpha=P^8,\quad\eta=\lambda\beta .
\tag{CP45.27}
\]
Their parameter windings are \((-6,0,0,0)\).
Each remaining coefficient has its actual periodic
logarithm on the height disc and parameter annulus.
For every integer \(j\), the exact vertical arrow is
\[
\begin{gathered}
V_j=Z^{-j}V,\quad H_j=hB^j,\quad
Q_j=\kappa qg^j,\quad
\alpha_j=\alpha\eta^j\kappa^{-j}g^{-j^2},\quad
\eta_j=\eta g^{-j},\\
V=Z^jV_j,\qquad h=H_jB^{-j}.
\end{gathered}
\tag{CP45.28}
\]
Its source interpolation is \(Q_t=Qe^{tj\log g}\).
In the node \(xy=Q_t\), the exact horizontal field is
\(\partial_t+j\log g\,y\partial_y\), with \(q\) fixed.
On the protected outside prescribe the marked real-period
map \(u+\tau_t b\mapsto u+\tau_0b\). Their difference
is vertical. A partition fixed on smaller collars glues
the two fields; compact forward and reverse flow gives
the actual periodic source map. The F9 circle has source
deck zero and suspension deck one. Its graph interpolation
therefore keeps the same short-arc endpoint order.

The exact Cartier spinor arrow is
\[
E_j=E\otimes\mathcal O(j(D_+-D_-)),\qquad
(x_-^{-j-1}x_+^j,x_-^{-j}x_+^{j+1})
\xrightarrow{\ Z^j\ }(1/x_-,x_+).
\tag{CP45.29}
\]
It respects the determinant trivializations, including
their reciprocal charts. Extend the transported fixed-sphere
velocity over the outer ball; retain its adjoint on the
Hopf factor. This constructs the relative completed
orbit-piece diffeomorphism. The physical line is carried
by the following separate, explicit coefficient map.

Use the actual logarithmic coefficients
\[
(z,\ell)=u(1,0)+b(\tau,\mu)+d(0,1)+\gamma(G,\mathcal B),
\quad
(\tau,\mu,G,\mathcal B)=
\frac1{2\pi i}(\log\kappa q,\log\alpha,\log g,\log\eta q).
\tag{CP45.30}
\]
In particular \(d=\operatorname{Re}\ell-
\operatorname{Re}\mu\,b-\operatorname{Re}\mathcal B\,\gamma\);
all coefficient derivatives stay present. The full deck
and shear arrows are
\[
q:\ (u,d)\mapsto(u-b,d-\gamma),\qquad
c:\ (u,d)\mapsto(u+6b,d),\qquad
(u_j,b_j,d_j,\gamma_j)=(u,b,d-ju,\gamma-jb).
\tag{CP45.31}
\]
For the \(q\)-angle \(\alpha_q\) and parameter angle
\(\theta\), choose
\[
A_0=dd+\alpha_q\,d\gamma,\quad
A_j=dd_j+\alpha_q\,d\gamma_j+6j\theta\,db,\quad
\beta_j=-j\{du+(\alpha_q-6\theta)db\}.
\tag{CP45.32}
\]
These are expressions in the full coefficient frame.
The physical and dual line conventions are
\(c_1(M)=-dA,c_1(U)=dA\).

The theta-cut line \(\mathcal T\) has the actual coefficient
seam \(e^{2\pi ib}\), connection \(d-2\pi i\theta\,db\),
and curvature \(\Xi=d\theta\wedge db\), with \(\Xi^2=0\).
The integral deck square gives
\[
U_j\longrightarrow U_0\otimes\mathcal T^{6j}.
\tag{CP45.33}
\]
Here the degree-two and boundary integers can be checked
in the actual K63 completed piece. Before surgery the
source-region product has \(H_1=\mathbb Z b\oplus
\mathbb Z\gamma\) and \(H_2=\mathbb Z F\oplus
\mathbb Z(b\gamma)\). Removing the fixed fibre adds its
primitive meridian: the fibre-times-\(\gamma\) three-cycle
can be represented disjointly, so the Thom map to the
intersection number is zero. The displayed Hopf surgery
kills the primitive fibre class and retains the meridian.
Its integral Mayer--Vietoris maps therefore give
\(H_1=\mathbb Z b\oplus\mathbb Z\gamma\) and torsion-free
\(H_2=\mathbb Z\,{\rm mer}\oplus\mathbb Z(b\gamma)\).
The actual source and line arrows fix both classes;
the Euler line has degree one on the meridian and zero
on the regular torus, which excludes an added meridian
in the transported torus class. The parameter Wang
sequence adds \(b\theta,\gamma\theta\). The full deck
square (CP45.31) evaluates the ratio as zero on the
meridian, \(b\gamma\), and \(\gamma\theta\), and as
\(6j\) on \(b\theta\). These are precisely the integral
degrees of \(\mathcal T^{6j}\).

The boundary \(q\)-commutator kills \(u\) primitively;
the parameter commutator is \(6u\). Thus its \(H_1\)
has free generators \(\alpha_q,\theta,b,\gamma\).
The last three extend through the completed piece.
For the possible remaining \(q\)-angle gauge choose
\(Z_*\) near \(-1\), away from all actual point/O
sections and node branches, with \(\gamma=0\).
The full section arrow is
\[
(q/Z_*,Z_*,1)\longmapsto
(qZ_*^{j-1},Z_*,Z_*^{-j}).
\tag{CP45.34}
\]
The source field fixes this disc; the point modules are
nonzero on it. The displayed coefficient is constant
and extends at \(q=0\), so the boundary winding is zero.
At another free lift the factors
\(Z_*e^{\gamma\log g}\) and
\(\eta_j^\gamma=\eta^\gamma e^{-j\gamma\log g}\)
cancel literally. None of their full logarithms is
discarded.

Remove the integral Cech two-cocycle of the ratio line
using the primitive cochain just computed. For its
remaining real logarithms \(r_{ij}\), the partition
formula \(a_i=\sum_k\rho_k r_{ki}\) has
\(a_j-a_i=r_{ij}\). Multiplying the vector frames
by \(e^{-2\pi ia_i}\) gives a nonzero global section.
Its inverse gives the direction (CP45.33). Remove
the extending \(\theta,b,\gamma\) characters on the
boundary; the remaining real boundary logarithm extends
by a collar cutoff. This constructs the prescribed
line map with inverse, including its cap value.

\par

\medskip\noindent\textbf{CP45.6. Tangent transgression, cubic collar and reflection}

\par\nobreak

The full source coframe satisfies
\[
\omega=du+\operatorname{Re}\tau\,db+
\operatorname{Re}G\,d\gamma,\qquad
\omega_j=du+\operatorname{Re}(\tau+jG)\,db+
\operatorname{Re}G(d\gamma-j\,db)=\omega .
\tag{CP45.35}
\]
In the ordered frame \((\omega,d\gamma,db)\) its matrix
is
\(\left(\begin{smallmatrix}1&0&0\\0&1&-j\\0&0&1\end{smallmatrix}\right)\).
The base variables are fixed. Differentiating the entire
coefficients gives the actual base-to-source upper
blocks. The strictly triangular source block has zero
\(\operatorname{tr}(F^2)\); its full curvature remains
strictly triangular. The base block depends on two
base angles and one interpolation variable, so its
four-form curvature square vanishes. The complete
block-triangular tangent arrow has zero Pontryagin
transgression, retaining all the coefficient derivatives.

For the residual pairing take
\(U_1\Subset U_2\Subset U_3\) containing the tube and
all surgery changes. With a real lift \(\widetilde b\)
on \(U_3\), a cutoff equal to one on \(U_2\), and a
constant \(b_*\), use the actual homotopy
\[
b_t=b\exp(2\pi it\chi(b_*-\widetilde b)).
\tag{CP45.36}
\]
It is fixed outside \(U_3\) and constant near the surgery
at the endpoint. A common regular value different from
that constant gives the relative level cobordism.
Its original representative is the source three-level
times the contracting circle, where the source-pulled
Pontryagin form has integral zero. The endpoint avoids
the surgery. Stokes and the zero boundary transgression
transport the zero value through that actual cobordism.
This proves \(\ell_{M,j}=0\).

On the common completed piece (CP45.33) gives
\(F_j=F_0+6j\Xi\). Since \(\Xi^2=0\), its interior
cubic difference is affine in \(j\). On the boundary
use the new coefficient \(\gamma=\gamma_j\); then
\[
F_{{\rm ext},j}=d\alpha_q\wedge d\gamma+
j\,d\alpha_q\wedge db,\qquad
d\beta_j=-j\,d\alpha_q\wedge db+6j\,d\theta\wedge db.
\tag{CP45.37}
\]
The outward primitive marked five-form is
\(\Omega_\partial=d\alpha_q\wedge du\wedge db\wedge
d\gamma\wedge d\theta\). Direct exterior multiplication gives
\[
3\int_0^1\beta_j\wedge
 (F_{{\rm ext},j}+s\,d\beta_j)^2\,ds
=18j^2\Omega_\partial.
\tag{CP45.38}
\]
The outward collar runs from the piece to the exterior,
hence from \(s=1\) to \(s=0\). It contributes
\(-18j^2\), giving \(k_{M,j}=-6+Lj-18j^2\).
The exact second endpoint determines the remaining
linear term. Complete theta inversion, including the
original point and tensor modules, gives
\[
(Z,V)\longmapsto(Z^{-1},P^8V)
\xrightarrow{\ \text{actual target }h\ }(Q/Z,V).
\tag{CP45.39}
\]
Both representatives of \(P\) are extending units at
\(q=0\). Their quotient cannot contain \(Q^h\) for
nonzero integer \(h\), by its valuation. Thus this is
the stated height-one reflection with no extra period.
The exterior and cap were constructed over \(v^2\).
The physical line and both held product caps are
transported by (CP45.39), which sends \(T_0\) to
\(T_1\). Therefore \(k_{M,1}=k_{M,0}\), giving \(L=18\).
For every signed integer \(j\),
\[
(k_M,\ell_M)_j=(-6-18j(j-1),0),\quad
(\mu_j,\nu_j)=(0,-6-18j(j-1)),\quad
A_j=-1-3j(j-1).
\tag{CP45.40}
\]

\par

\medskip\noindent\textbf{CP45.7. Ordinary products with their full height inverses}

\par\nobreak

For a unimodular basis \(P\), write \(R=P^{-1}\).
The four actual atom coefficients give
\[
A_{ik}=(\kappa^{P_{1k}}g^{P_{2k}})^{R_{i1}}
       (\alpha^{P_{1k}}\eta^{P_{2k}})^{R_{i2}},
\qquad B_P=-6P^{-1}e_1e_1^tP.
\tag{CP45.41}
\]
The height-valuation matrix is \(RP=I\). All four
based logarithms remain in these coefficients, with
the single explicit \(-6\) winding separated.
The exact three edge maps are
\[
\begin{array}{c|ccc}
C_0&A_{11}A_{21}bc&A_{21}^{-1}ab&A_{11}^{-1}b^{-1}\\
C_1&A_{12}A_{22}bc&A_{22}^{-1}c^{-1}&A_{12}^{-1}ac\\
C_2&a^{-1}&ab&ac .
\end{array}
\tag{CP45.42}
\]
Their products are the original \(q=abc\). Their native
second charts are respectively
\((A_{11}w,A_{21}v,(A_{11}A_{21})^{-1}u)\),
\((A_{22}v,A_{12}w,(A_{12}A_{22})^{-1}u)\), and
\((u,v,w)\). These give \(z'=z^{-1},n'=zn\) with
the first normal pairs \((a,c),(a,b),(b,c)\).

For the ordinary product keep the first vertex and use
\[
(u,v,w)\mapsto(u,A_{22}(s,q)v,A_{11}(s,q)w),\qquad
h_s(q)=qA_{11}(s,q)A_{22}(s,q).
\tag{CP45.43}
\]
The inverse first recovers \(q=h_s^{-1}(uv'w')\),
then divides by the two units evaluated at that
recovered height. Its derivative at zero is nonzero
uniformly on the compact parameter annulus. The
three full normal blocks are
\[
D_0=\operatorname{diag}(A_{22}/A_{21},A_{11}A_{21}),\quad
D_1=\operatorname{diag}(A_{11}/A_{12},A_{12}A_{22}),\quad
D_2=\operatorname{diag}(A_{22},A_{11}).
\tag{CP45.44}
\]
Their first windings are
\(l=(B_{22}-B_{21},B_{11}-B_{12},B_{22})\), and
their determinant winding is \(-6\).

The radial quaternion maps
\[
H_d=
(\sin\vartheta\,e^{id\theta}+j\cos\vartheta)
(\sin\vartheta+j\cos\vartheta)^{-1}
\tag{CP45.45}
\]
are one on the inner collar, \(e^{id\theta}\) on
the outer collar, and based periodic maps. On edge
\(i\) use \(Q\mapsto H_{-3}QH_{l_i+3}\).
The complete periodic zero-winding logarithms
\(k_{i0},k_{i1}\) have their actual unit map
\[
(z,n_0,n_1)\mapsto
(z,e^{\rho k_{i0}(s,q)}n_0,e^{\rho k_{i1}(s,q)}n_1),
\quad H_{s,z}(q)=q e^{\rho(k_{i0}+k_{i1})(s,q)} .
\tag{CP45.46}
\]
Its inverse is
\(q=H_{s,z}^{-1}(zn'_0n'_1)\), followed by the
two reciprocal coefficients at this \(q\).
Divide by the nonzero derivative at zero and
shrink uniformly so that the derivative differs
from one by less than \(1/2\). Integration on a
straight segment proves injectivity; the inverse
function theorem gives the displayed inverse with
smooth radial dependence. This keeps the full
height dependence in both directions.
The two-vertex identity is \(zD_{i,*}K_i=zD_i\).
Normal rescaling of the complete germ, followed by
the cutoff-velocity flow, retains all higher jets
and their inverse, fixed on the protected collars.
The Hopf map is
\((UH_{-3}U^{-1})q_HH_{l_i+3}\); its left degree
is \(-6\), its right degree zero. Its actual inverse
cylinder is \(L^6\), with the paired Spin endpoints.

For \((x,y)^t=P^{-1}e_2\), define
\[
w=(x,y,-x-y),\quad d=(-x-2y,-2x-y,x-y),\quad
I_i=-3w_i+d_i(l_i+3),\quad
J_i=-3w_i-d_i(l_i+3).
\tag{CP45.47}
\]
The physical orbit has left angle \(-w_i\phi/2\)
and right angle \(d_i\phi/2\) in the native normal.
Conjugation by (CP45.45) has axis degrees
\(3,l_i+3\). Since \(d_i\equiv w_i\pmod2\),
the two Spin endpoints at \(2\pi\) multiply to one.
The full circle map and derivative give \(I,J\),
and the actual integral prism equations are
\[
V(w_ie_i-w_0e_0)=I_ie_i-I_0e_0,\qquad
w_i\delta_i-w_0\delta_0=J_i-J_0,\quad i=1,2.
\tag{CP45.48}
\]
The vertical complete ordinary seed, obtained through
the same \(L^6\) cylinder from (CP45.40), is
\[
V=
\begin{pmatrix}
18j(j-1)&6-12j&12j-6\\
6-12j&0&0\\12j-6&0&0
\end{pmatrix},
\quad
\delta=(-12-18j(j-1),12j-12,-12j).
\tag{CP45.49}
\]
All six vector equations and both derivative
equations hold in this seed, including the fixed
zero-weight edge.

\par

\medskip\noindent\textbf{CP45.8. Complete fan paths, native reset and uniqueness}

\par\nobreak

The triangles of \(T(P)\) are
\(\{n,n+v,n+w\}\) and
\(\{n+v,n+w,n+v+w\}\).
Column swap and right multiplication by
\[
R_0=\begin{pmatrix}-1&-1\\1&0\end{pmatrix},
\quad R_0^2,\quad -I,\quad -R_0,\quad -R_0^2
\tag{CP45.50}
\]
give the same absolute cones by their integral
triangle translations. The identity on each absolute
character semigroup and its inverse commute with
the full coefficient maps, and are the identity
on the original source, point-line and cap
overlap. Thus these are actual marked same-fan
maps for both determinant signs.

The resulting representative is a vertical \(T_j\),
or \(P=\left(\begin{smallmatrix}p&-r\\q&s\end{smallmatrix}\right)\)
with \(p,r>0,ps+rq=1\). The exact parent matrices
are
\[
Q_0=\begin{pmatrix}1&-1\\0&1\end{pmatrix},\quad
Q_1=\begin{pmatrix}1&0\\-1&1\end{pmatrix};
\quad
P\mapsto PQ_0^{-1}\ (r\ge p),\quad
P\mapsto PQ_1^{-1}\ (p>r).
\tag{CP45.51}
\]
They change \((p,r)\) to \((p,r-p)\) or \((p-r,r)\).
The sum strictly decreases and the gcd stays one.
Equality is \((1,1)\); its next parent is the
unique vertical basis. Reversal first uses \(Q_0\)
and then stays in the nonzero-weight region.

The dense character arrows are \((ZV,V)\) and
\((Z,ZV)\); the first triples are
\((a/c,bc,c)\) and \((a/b,b,bc)\).
The other-ruling charts are
\((n_0/n_1,zn_1,n_1)\) and
\((n_1/n_0,zn_0,n_0)\), with common rank-one
matrix
\(\left(\begin{smallmatrix}n_0&n_1\\zn_0&zn_1\end{smallmatrix}\right)\)
and unchanged \(q=zn_0n_1\).
For \(Q_1\), the actual normal interchange is
joined to the identity by
\[
\frac12
\begin{pmatrix}1+e^{i\pi t}&1-e^{i\pi t}\\
1-e^{i\pi t}&1+e^{i\pi t}\end{pmatrix};
\tag{CP45.52}
\]
it commutes with \(zI\) and has positive real
determinant. Every atom coefficient transforms by
\[
A'_{ik}=\prod_{h,l}A_{hl}^{(Q^{-1})_{ih}Q_{lk}}.
\tag{CP45.53}
\]
The common blowup and both ruling projections
carry this entire source, line, normal, exterior
and cap family. Their primitive integral
Gysin and van Kampen maps preserve the reference
sphere. No source or period marking is changed
to obtain a numerical seed.

On the selected old edge write \(a=V_{ii}\),
\(b=\delta_i\). The explicit \(L^6\) comparison
types the ordinary and native coordinates as
\[
\mu=-6-\frac{a+b}2,\qquad
\nu=\frac{b-a}2.
\tag{CP45.54}
\]
The actual Cox link map
\((h,q_H)\mapsto(q_H^{-1}hq_H,q_H)\), with its
carried complex reset, gives
\((\mu,\nu,s)\mapsto(\mu+2\nu,-\nu,s+\nu)\).
Applying (CP45.54) on the target with the same
cylinder gives
\[
V'_{22}=a,\qquad\delta'_2=2a-b,\qquad
s'=s+(b-a)/2 .
\tag{CP45.55}
\]
This is an integer relation in the actual native
coordinates before reduction modulo 28.

All target weights are nonzero. The full target
prism has the unique solution
\[
V'_{rs}=\frac{I_r}{w_r}\delta_{rs}
       +\frac C{w_rw_s},\quad
\delta'_r=\frac{J_r+E}{w_r},\qquad
C=w_2^2a-w_2I_2,\quad E=w_2(2a-b)-J_2.
\tag{CP45.56}
\]
Indeed a symmetric difference annihilating the
two prism vectors is a multiple of the outer
product \((1/w_r)(1/w_s)\); the derivative
difference has one analogous constant. The
two target edge values fix both constants.
Existence is the actual target family above,
and its integer variation equals this unique
rationally written solution. The independent
certificate solves the entire nine-variable
linear system afresh and checks this expression,
the integer reset and all full atom exponents
on finite canonical paths. The Euclidean
decrease in (CP45.51) proves termination for
every integral basis without a boundedness
assumption.

These comparisons prove the stated initial
\((-6,0)\), vertical (CP45.40), and full
all-fan algorithm for the original OS45 loop.
The full scalar is held fixed. An admitted
scalar is transported to the common strict
bound by its constant-argument interval;
the actual logarithm in (CP45.8), (CP45.27)
and all reference sectors is retained.
The reciprocal interval returns to its
original value with no scalar turn. This
gives the same conclusion for every admitted
fixed scalar in the stated order-one cap
family. The proof assigns no shorter
parameter endpoint or additional cap order.

\par

\medskip\noindent\textbf{CP45.9. Independent execution and provenance}

\par\nobreak

The independent checker is
\path{code/part3/os45_selected_cusp_completion_blind_peer.py}.
It uses the source equations and the displayed full
character arrows directly; it does not import the
author\textquotesingle s completion routine. It verifies the
critical Jacobian, both columns of the Möbius
differential, quaternion signs, the two relative
Chern systems, exact period/deck changes, the
five-form collar identity, both native overlaps,
the full recovered-height inverse construction,
all four atom exponents and their reciprocal
transformations, and fresh nine-variable Gaussian
prism solutions on 217 complete canonical paths
from seven signed vertical seeds through word
length five. The source and all proof dependencies
are locked before and after execution.

The local 68-check proof is retained separately.
The full geometric audit includes the protected
source and cap, the actual orbit pair, the
embedded link, all point/O modules, both
determinant and Hopf arrows, the relative line
and tangent maps, and the entire fan path.
The two located source refinements were the
full \(-2S\,dc\) column in (CP45.22) and the
recovered-height inverses (CP45.43), (CP45.46).
Both are present in the locked final author
source. Earlier author versions and their
certificates are preserved by the author.
The subsequent 4,514-count wording refresh
changed no map or characteristic number.

% End expanded verified cp36 input


% End expanded verified cp36 input

\par\medskip\noindent\textbf{Complete execution and provenance.}\par
The original local certificate passes 55 checks, and its independent local proof passes 68. The complete primary execution passes 4,514 checks with 217 fresh nine-variable target solves. The full root peer passes 1,500 checks with 90 fresh target solves; the full blind peer passes 17,598 checks with 217 fresh target solves. All five supporting proofs are reproduced without omitted mathematical sections or displays. Exact input, code, audit, source and page locks are retained in \path{records/k117_export_manifest.json}.
\end{proof}


\subsection{The original E6 source with both additive caps: all scalar endpoints and fan returns}
\begin{theorem}[Source result P3K.118 under continuation review]\label{night:P3K.118}
\StatementRecord{night:P3K.118}{open}{Historical events: GPT-6 Astra (OpenAI), ChatGPT Pro mode, effort Pro, S6\_06\_BLIND, 16--18 November 2026: targeted proof check. Also checked in this writing. Continuation mathematical and source checks: GPT-6 Astra (OpenAI), ChatGPT Pro mode, effort Pro, S6\_10\_PAPERS continuation, 19 November 2026: computation or formal certificate; used and rechecked in this run. Continuation metadata events: GPT-6 Astra (OpenAI), ChatGPT Pro mode, effort Pro, S6\_10\_PAPERS continuation, 19 November 2026: used and rechecked in this run; receiving source placement; metadata reconciliation only; no new mathematical check. GPT-6 Astra (OpenAI), ChatGPT Pro mode, effort Pro, S6\_10\_PAPERS continuation, 19 November 2026: used and rechecked in this run; explicitly scoped metadata reconciliation; metadata reconciliation only; no new mathematical check. GPT-6 Astra (OpenAI), ChatGPT Pro mode, effort Pro, S6\_10\_PAPERS continuation, 19 November 2026: used and rechecked in this run; source-row intake; metadata reconciliation only; no new mathematical check. \emph{Pre-continuation recorded proof scope.} Historical event: GPT-6 Astra (OpenAI), ChatGPT Pro mode, effort Pro, S6\_06\_BLIND, 16--18 November 2026: targeted proof check, recorded producing-run re-derivation. Also checked in this writing by GPT-6 Astra (OpenAI), ChatGPT Pro mode, effort Pro, S6\_10\_PAPERS continuation, 19 November 2026: computation or formal certificate for the explicitly listed algebraic substeps in the event ledger; full continuation rederivation and renewed credit comparison remain pending.}{the original source, point line, roots, characters, reference, selected fan and scalar range stated here}
Retain the original marked source, selected whole nodal fibre and point line
\[
y^2=x^3+tx+t^2,\qquad G=(0,t),\qquad
q=t+4/27,\qquad \mathcal M_a=\mathcal O([aG]-[O]),\qquad ab=6,
\]
with source translation by \(-bG\). Its two original additive caps,
including their full Cartier maps, finite characters and prescribed
annular references, have orders in the following four domains:
\[
\begin{array}{c|c|c}
(a,b)&m=m_\infty&h=m_0\\ \hline
(1,6)&3l,\ l\ge1&4k,\ k\ge1\\
(2,3)&3l,\ l\ge1&4k+1,\ k\ge0\\
(3,2)&3l+2,\ l\ge0&4k,\ k\ge1\\
(6,1)&3l+2,\ l\ge0&4k+1,\ k\ge0.
\end{array}
\]
The prescribed norms satisfy \(hn_\infty+mn_0=\varepsilon\), with
\(\varepsilon=\pm1\). The additional ordinary total norm is zero.
Retain every prescribed selected integral unimodular fan and the
original strictly contracting scalar domain. These objects have the
original six-sphere topology through the full F46 source/cap and
common-blowup maps.

The holomorphic automorphism group of each such marked object is
exactly its faithful physical \(\mathbb C^*\). Every holomorphic
scalar endpoint fixes the original base \(t\), has positive affine
linear part on the regular two-tori, and has equal original scalar
values at its two ends. Its possible integer scalar windings are
exactly
\[
J_0\mathbb Z,\qquad J_0=\frac{mh}{b}.
\]
At a fixed winding the full holomorphic endpoints differ precisely
by the extending physical scalar. This classification includes both
root caps and their complete reference words.

Retain the rational positive-\(G\) line operator \(U\), the actual
reference constants \(C_a,c_0\), and the full original identities
\[
B_0=C_a\frac q{t^3}U^{-b},\qquad
h_0=c_0q^\varepsilon t^{-mn_0}.
\]
On \(\lambda=r^{mh}\), the primitive forward endpoint is
\[
E_r=\gamma_0r^{-\varepsilon J_0}t^{d_0}U^\varepsilon,
\qquad
d_0=\frac{3\varepsilon-mn_0}{b},\qquad
\gamma_0^b=c_0C_a^{-\varepsilon},
\qquad r\longmapsto\zeta_b r.
\]
Its inverse is evaluated at the actual inverse source point. Its
\(b\) chronological copies give exactly
\(h_0B_\lambda^{-\varepsilon}\), with all original constants retained.
The complete K67 parameter-family comparison is the physical map
\[
F_r=r^{\varepsilon J_0(b-1)/2},\qquad
F_{\zeta_b r}E_r=
\exp\!\left(\frac{\pi i\varepsilon J_0(b-1)}b\right)E_rF_r.
\]
The exponent is integral on every displayed domain. The original
family is the literal degree-\(b\) pullback of the completed shorter
family over \(w=r^b\), \(\lambda=w^{J_0}\), including both caps,
the selected Cox maps, physical line, stable tangent and all seams.

For one positive primitive scalar circle, the initial physical pair
\((k_{\rm phys},\ell_{\rm phys})\) and index \(S_0\) are
\[
\begin{array}{c|c|c|c}
(a,b)&k_{\rm phys}&\ell_{\rm phys}&S_0\\ \hline
(1,6)&kl&4kl&0\\
(2,3)&2l(4k+1)&8l(k+1)&kl\\
(3,2)&9k(3l+2)&12k(l+2)&k(4l+2)\\
(6,1)&18(3l+2)(4k+1)&24(kl+2k+l+1)&35kl+22k+8l+5.
\end{array}
\]
The actual held-quaternion native pair is
\[
(\mu,\nu)=(\ell_{\rm phys}/4,\,
k_{\rm phys}-\ell_{\rm phys}/4).
\]
For the vertical fan \(T_j\), the full primitive data are
\[
S_j=S_0-\frac{J_0}{2}j(j-1),\qquad
k_j=k_{\rm phys}-3J_0j(j-1),\qquad
\ell_j=\ell_{\rm phys}.
\]
For every selected lattice-basis fan \(P\), the complete terminating
algorithm RJC.78--80, equivalently the full native maps and
nine-variable target systems JFP.41--43, gives its actual integer
return \(S(P)\). This includes every full source coefficient, original
height, ordinary product, physical normal and cap reset.

Every holomorphic scalar path of winding \(sJ_0\) and every
holomorphic closing at that winding has return
\(sS(P)\pmod{28}\) on its prescribed fan. Conversely every
integer \(s\) is attained by the explicit shorter-family annular
pullback. The exact images over all selected fans and all such
paths are
\[
\begin{array}{c|c|c}
(a,b)&\text{holomorphic winding group}&
\text{complete all-fan return image}\\ \hline
(1,6)&2kl\mathbb Z&kl\mathbb Z/28\\
(2,3)&l(4k+1)\mathbb Z&l\mathbb Z/28\\
(3,2)&2k(3l+2)\mathbb Z&k\mathbb Z/28\\
(6,1)&(3l+2)(4k+1)\mathbb Z&\mathbb Z/28.
\end{array}
\]
For the positive seam \((x,1)\sim(\Phi(x),0)\), the oriented
mapping torus is
\[
(S^1\times S^6)\mathbin{\#}\Sigma_{sS(P)}.
\]
The inverse seam gives \(\Sigma_{-sS(P)}\). For each of the two
spin structures the corresponding Eells--Kuiper values are
\(sS(P)/28\) and \(-sS(P)/28\) in \(\mathbb Q/\mathbb Z\).
The two signed sphere equations retain the same positive scalar
parameter orientation throughout these statements.

All assertions retain the four stated root progressions, their
original norm representatives, every finite character, full annular
reference and selected fan. The scope is the fixed original source
and zero additional ordinary total norm.
\end{theorem}
\begin{proof}[Supplied source argument; continuation review in progress]

% Expanded verified cp36 input: notes/part3_e6_iii_joint_endpoints_proof.tex
\par\bigskip\noindent\textbf{Complete supporting proof 1: RJC.}\par


% Expanded verified cp36 input: notes/part3_e6_iii_joint_endpoints_rjc_proof.tex
\emph{Retained preparation history.} The primary source\textquotesingle s introductory and final preparation-stage audit sentences are reproduced verbatim. The complete independent derivation and geometric audit follow as JFP.1-\/-12. The exact current gates and their source locks are recorded in \path{records/k118_export_manifest.json}. The completed reciprocal full-source review is included in that gate.

\par

\medskip\noindent\textbf{The original E6 source with both additive caps: endpoint continuation}

\par\nobreak

This continuation retains K113-\/-115 and the full F46 joint-cap construction.
The original source is
\[
 E:\quad y^2=x^3+tx+t^2,\qquad G=(0,t),\qquad
 t_c=-4/27,\quad q=t+4/27.                         \tag{RJC.1}
\]
This is the full primary continuation on the four actual two-cap
domains below. It retains every point-line module, root comparison,
finite character and prescribed annular reference. It constructs the
primitive endpoint, classifies all holomorphic scalar endpoints, and
carries their parameter families to every lattice-basis fan. The
independent full audit and the manuscript statement remain pending.
The earlier derivation, certificate and output are preserved exactly in
\path{records/attempts/e6_iii_joint_before_complete_proof/}.

\par

\medskip\noindent\textbf{RJC.1. The actual four root domains and the signed sphere equation}

\par\nobreak

Write \(m=m_\infty\), \(h=m_0\), \(M=mh\). The actual root domains of
F46, including the finite root rather than its order-one special case, are
\[
\begin{array}{c|c|c|c}
 (a,b)&m&h&\text{minimal root fibres}\\ \hline
 (1,6)&3l,\ l\ge1&4k,\ k\ge1&\text{smooth, smooth}\\
 (2,3)&3l,\ l\ge1&4k+1,\ k\ge0&\text{smooth, III}\\
 (3,2)&3l+2,\ l\ge0&4k,\ k\ge1&\text{IV, smooth}\\
 (6,1)&3l+2,\ l\ge0&4k+1,\ k\ge0&\text{IV, III}.
\end{array}                                                        \tag{RJC.2}
\]
The two retained norms satisfy
\[
 hn_\infty+mn_0=\varepsilon,\qquad \varepsilon\in\{1,-1\}.             \tag{RJC.3}
\]
This is the actual two-cap sphere equation proved by all integral maps
F46.J.7-\/-12. It forces \(\gcd(m,h)=1\) and each local primitive norm.
The ordinary order-one framing in this first continuation has zero total
additional norm. Its full annular reference and both finite characters
are retained. Additional ordinary clutch slots have their separate
K32 exponents and will be accounted for explicitly before that scope is
included.

The minimal root substitutions are, at infinity,
\[
\begin{array}{c|c|c}
m&(t,x,y)&\text{root equation}\\ \hline
3l&(z^{-3l},z^{-2l}X,z^{-3l}Y)&Y^2=X^3+z^lX+1\\
3l+2&(z^{-(3l+2)},z^{-2(l+1)}X,z^{-3(l+1)}Y)
 &Y^2=X^3+z^{l+2}X+z^2,
\end{array}                                                        \tag{RJC.4}
\]
and at zero,
\[
\begin{array}{c|c|c}
h&(t,x,y)&\text{root equation}\\ \hline
4k&(z^{4k},z^{2k}X,z^{3k}Y)&Y^2=X^3+X+z^{2k}\\
4k+1&(z^{4k+1},z^{2k}X,z^{3k}Y)
 &Y^2=X^3+zX+z^{2k+2}.
\end{array}                                                        \tag{RJC.5}
\]
The good III central point is \(T=(0,0)\) on \(Y^2=X^3+X\).
The good IV-star central point is \((0,1)\) on \(Y^2=X^3+1\).
Their orders are respectively two and three, by the tangent divisors.
The inverse substitutions over \(z\ne0\) are multiplication by the
reciprocal displayed powers. The deck acts by multiplying \(z\) by
\(\zeta_{m_i}\) and multiplying \(X,Y\) by the inverse original
coordinate weights. Thus each original coordinate is literally
invariant. The reduced root maps include the original A1 or A2
resolution.

Here are the full point-line correspondences used below. At a reduced
root the point \(aG'\) has identity component. Both the original line
and \(N=\mathcal O(aG'-O)\) have zero degrees on every component. On
the resolved fibre product the first line has zero degree on every
curve contracted by relative minimalization. The blowup Picard
decomposition therefore descends it. Its ratio to \(N\) has vertical
divisor with zero component intersections; the III or IV intersection
matrix has primitive integral kernel the whole-fibre vector. The
actual base equation, with its integer exponent and deck character,
is the resulting Cartier comparison. Its reciprocal is the inverse.

For the good IV-star root the minimal cubic quotient has tangent
weights \((2,1)\) at its three fixed points. On the common root
refinement use rays
\[
 (1,0),(2,1),(1,1),(1,2),(0,1).
 \tag{RJC.19}
\]
For a primitive ray \(v=(v_\eta,v_z)\), the difference between the
root section and the pulled-back resolved section is
\(\min(v_\eta,2v_z)\). Its three coefficients are \(2,1,1\).
Subtracting the \(O\) contribution from the \(aG'\) contribution
gives the complete F46.J.1 Cartier equality. At order \(3l\), pull
the entire diagram back by \(w=z^l\) and normalize its common
blowups. The coefficients are now
\[
 D_P(v)=\min(v_\eta,2l v_z),\qquad
 \beta^*N=\pi^*\mathcal M+D_{aG'}-D_O.
 \tag{RJC.20}
\]
Every toric chart map sends its two monomials to the original
\((\eta,z)\); inversion of the unimodular cone matrix gives its
inverse on the common torus. The strict-section function remains
in the Cartier numerator. This specifies the original point module
as well as every exceptional factor.

For the minimal good III root \(w\), the action is
\[
 (w,X,Y)\longmapsto(iw,-X,iY).
 \tag{RJC.21}
\]
It fixes \(O,T=(0,0)\) and has the order-two orbit
\((i,0),(-i,0)\). Blow up those four points and quotient. The
exceptional images have squares \(-4,-4,-2\); the central image
has square \(-1\), and the fibre is \(4C+D_0+D_1+2D_2\).
Contract \(C\), then the resulting \((-1)\)-curve \(D_2\).
The remaining curves have squares \(-2\) and intersection two.
The invariant substitution in RJC.5 gives the actual III
tangency, with its length-two intersection scheme. The marked
sections avoid both contractions. Consequently
\[
 \beta^*N=\pi^*\alpha^*\mathcal M+E_{aG'}-E_O.
 \tag{RJC.22}
\]
For order \(4k\) the pullback \(w=z^k\) resolves the ideals
\((\eta,z^k)\) at \(O,T\). A complete regular fan has rays
\((j,1)\), \(0\le j\le k\), and \((1,0)\). On successive cones its
literal maps are
\[
 \eta=u^jv^{j+1},\quad z=uv\quad(0\le j<k),\qquad
 \eta=uv^k,\quad z=v
 \tag{RJC.23}
\]
on the last cone. The cone determinants are \(-1\), so their inverse
monomials are integral on the common torus. Pulling the ideal through
these maps gives exactly \(\min(v_\eta,k v_z)\); retaining these
generators in RJC.22 gives its higher-order comparison. The two
remaining fixed-point charts retain their source blowups and carry
no marked point divisor.

For a reduced III root \(h=4k+1\), the original formal parameter
\(\tau=-x/y\) pulls back as \(\tau'=z^k\tau\). The actual
\(aG'\) graph, for \(a=2,6\), is
\(\tau'=z^k A_a(z^h)\), where \(A_a(0)=a\).
The normalized blowup of \((\tau',z^k)\) carries the exact two
Cartier ratios
\[
 \frac{p-A_a(z^h)}p,\qquad 1-vA_a(z^h),\qquad
 \tau'=z^kp,\quad z^k=\tau'v.
 \tag{RJC.24}
\]
The intermediate regular charts are the monomial charts RJC.23.
The source A1 chart is
\[
 X=z\alpha,\quad Y=z\beta,\quad
 \beta^2=z\alpha^3+\alpha+z^{2k}.
 \tag{RJC.25}
\]
At \(k=0\), this includes the original equation with the constant
term \(1\). The inverse is division by \(z\) off the exceptional
curve, and the equation supplies its regular extension. Translation
by \(\pm G'\) uses the actual chord slopes
\((Y\mp z^{k+1})/X\). Resolving their graph gives the mutually
inverse translations of the minimal regular surface. Their
point-module maps carry RJC.24 to the translated point graphs.

The reduced IV root uses the A2 charts of R0E.60-\/-63: in the first
chart \(X=zX_1,Y=zY_1\), the two nonidentity labels are \(Y_1=\pm1\).
The original coordinate substitutions in RJC.4 identify these
labels and the whole exceptional chain. The point graph is the
full ratio \(\tau'/(\tau'-z^{2l+1}L_a(z))\). For the actual
\(3G\) point in RJC.26 its original-coordinate substitution gives
\[
 L_3(z)=-\frac{64+8z^m}{512+96z^m+3z^{2m}},
 \qquad L_3(0)=-1/8.
\]
For the actual doubled point \(6G\) in the same formulae,
\(L_6(0)=-1/4\). The full rational function
\(L_a=-z^{-(2l+1)}X_{aG'}/Y_{aG'}\)
is retained in its own point graph.
Its reciprocal and the inverse chord map give the reverse line
and source arrows. These are the complete F46 and R0E diagrams
for the specific four rows, with their full Cartier functions.

\par

\medskip\noindent\textbf{RJC.2. The ordinary lift, exponents and exact root characters}

\par\nobreak

Retain the rational point-line lift \(U\) of \(+G\) obtained from the
actual source Miller maps. The chosen original forward Hom section is
\[
 B_0=C_a\,\frac{q}{t^3}\,U^{-b},\qquad C_a\in\mathbb C^*.             \tag{RJC.6}
\]
Here \(C_a\) is the coefficient of that specific section in these
specific Miller frames, fixed by evaluation at the same ordinary
reference point; it remains in every map. We verify RJC.6 on the
original regular surface. Put \(T(p)=p-G\). The tangent and chord
formulae give
\[
 \begin{split}
 f_1&=1,\qquad
 f_2=\frac{y-t-x/2}{x-1/4},\\
 f_3&=\frac{(x/4)(y-t-x/2)}
 {(y+t+x/2)/4-(1/4+2t)x},\\
 (x_3,y_3)&=(8t+64t^2,\,3t+96t^2+512t^3),\\
 L_3&=(3x_3^2+t)/(2y_3),\\
 x_6&=L_3^2-2x_3,\quad y_6=L_3(x_3-x_6)-y_3,\\
 f_6&=f_3^2\frac{y-y_3-L_3(x-x_3)}{x-x_6}.
 \end{split}
 \tag{RJC.26}
\]
Their divisors are \(a(G)-(aG)-(a-1)(O)\). Therefore
\[
 B_{a,1}(p)=f_2(p)^a\frac{f_a(Tp)}{f_a(p)},\qquad
 U(p,v)=\bigl(p+G,B_{a,1}(p+G)^{-1}v\bigr)
 \tag{RJC.27}
\]
are the actual negative one-step point-line map and its positive
inverse. The coefficient of the inverse is evaluated at \(p+G\).
All zero and pole factors in RJC.26 belong to the corresponding
point Cartier modules, so these identities are Hom identities
before any specialization.

On the original III components the orders of \(f_2,f_3,f_6\) are
\((0,1),(0,1),(0,3)\). To see these in the original coordinates,
use \(x=p^{-2},y=-p^{-3}\) on the identity component and
\(x=t\alpha,y=t\beta\),
\(\alpha+t\alpha^3=\beta^2-1\), on the other.
The translation \(T\) has directed component formulae
\(\beta'=p-1,p'=\beta-1\).
Substitution into the entire functions RJC.26 gives the orders
just stated and preserves the point factors of R0R.11.
It follows that the two component orders of \(B_{a,1}\) are
\[
 \begin{array}{c|cccc}
 a&1&2&3&6\\ \hline
 \operatorname{ord}_{(C_0,C_1)}B_{a,1}
 &(0,1)&(1,1)&(1,2)&(3,3).
 \end{array}
 \tag{RJC.28}
\]
The \(b\) chronological factors of \(U^{-b}\) thus have order
three on each original component in all four rows. There is no
other finite whole-fibre divisor: the two substitutions in
\(f_a(Tp)/f_a(p)\) have equal generic orders on an irreducible
fibre, and \(f_2\) has zero whole-fibre order at every finite
nonzero \(t\), including the selected node.
The exact integral source translation in F46.J.1,J.6 gives the
Hom class \(F\). Its only remaining whole-fibre coefficient, at
infinity, is consequently \(-2\): \(3F-2F=F\).
Multiplication by \(q/t^3\) leaves exactly the chosen simple
whole-fibre zero at \(t_c\). Its ratio to the original specified
section \(B_0\) has zero divisor on the compact regular source.
Both it and its reciprocal are regular and hence constant.
Reference evaluation defines precisely the nonzero \(C_a\) in
RJC.6. This proves that identity for all four actual point lines.
Its literal rational norm identity is
\[
 \rho=\frac{B_0^{-1}}{U^b}=\frac{t^3}{C_aq}.                         \tag{RJC.7}
\]
For completeness, the new odd-line good-root coefficients can be
read directly before taking central values. At the finite good root
put \(w=z^k\). Then \(f_{2,\rm old}=w^2\phi\),
\(f_{3,\rm old}=w^2\psi\), where
\[
 \phi=\frac{wY-w^2-X/2}{w^2X-1/4},\qquad
 \psi=\frac{X(2wY-2w^2-X)}
 {2wY+2w^2-X-16w^4X}.
 \tag{RJC.29}
\]
The central negative source translation is
\((X,Y)\mapsto(1/X,-Y/X^2)\).
Thus the normalized negative coefficients for \(a=1,3\) are
\(2X,8X\). Their positive inverse coefficients, evaluated at the
inverse source point, are \(X/2,X/8\). The negative twofold
norms are \(4,64\), and the \(a=1\) sixfold norm is \(64\).
These are values of the full functions RJC.29 in the full
Cartier comparison RJC.22-\/-23.

At the good infinity root put \(v=z^l\). The full coefficient is
\[
 f_{2,\rm old}=v^{-1}
 \frac{Y-1-vX/2}{X-v^2/4}.
 \tag{RJC.30}
\]
On the central cubic the positive source map is
\[
 X_+=2(1-Y)/X^2,\qquad
 Y_+=(Y-1)(X-X_+)/X-Y.
 \tag{RJC.31}
\]
Its \(a=1\) inverse-source line coefficient is \(X/2\);
its three chronological factors multiply to \(-1/2\) and
its six factors to \(1/4\). The opposite sixfold norm is \(4\).
Together with RJC.20 these formulae carry the odd point line
through its full original infinity comparison.

Put
\[
 J_0=\frac{mh}{b},\qquad
 d_0=\frac{3\varepsilon-mn_0}{b}
     =\frac{2\varepsilon+hn_\infty}{b}.                              \tag{RJC.8}
\]
Both are integers on every domain RJC.2. The equality follows by
subtracting RJC.3. Substituting that same actual equation gives the
explicit integral expressions
\[
\begin{array}{c|c|c}
(a,b)&J_0&d_0\\ \hline
(1,6)&2kl&2kn_\infty+ln_0\\
(2,3)&l(4k+1)&(4k+1)n_\infty+2ln_0\\
(3,2)&2k(3l+2)&6kn_\infty+(3l+2)n_0\\
(6,1)&(3l+2)(4k+1)&3(4k+1)n_\infty+2(3l+2)n_0.
\end{array}                                                        \tag{RJC.9}
\]

The actual extending normalized root lift \(L_i\) in K67.SE.3 satisfies
\(L_i^b=B_0^{-1}\). The line correspondences and that norm identity
retain the following source comparison:
\[
\begin{array}{c|c|c}
i&e_i&U=z^{e_i}u_i(z)L_i,\quad u_i^b\\ \hline
0&-3h/b&C_aq\\
\infty&2m/b&C_a h_\infty,\quad h_\infty=1+4z^m/27.
\end{array}                                                        \tag{RJC.10}
\]
The constants of the two analytic unit roots are fixed by the actual
source/Cartier correspondence. Here is the extension and normalization
that produces them. The root chord lift of \(+G'\) is a meromorphic
Hom section. Its horizontal divisor is the prescribed difference of
point modules. After the exact point corrections RJC.20-\/-25 its
remaining vertical divisor has zero component degrees. The primitive
kernel of the resolved intersection matrix makes it an integer
whole fibre. Divide by that actual power of \(z\).
The resulting section and its inverse have zero orders at every
codimension-one point. Normal local rings express each as the
intersection of their height-one valuation rings; hence both are
regular, including at intersections and section points. They give
the extending positive lift with its extending inverse.

Its deck commutator is a nonzero base unit. Choose its logarithm
with the prescribed value at the reference. For its norm-zero
nonconstant part \(g\), the finite averaging
\(-m_i^{-1}\sum_{j=0}^{m_i-1}jD_i^jg\) supplies the correcting
logarithm. An invariant analytic \(b\)-th root of the remaining
norm unit then gives \(L_i^b=B_0^{-1}\). Both corrections are
actual scalar maps on the root line, with their reciprocal maps.
The finite value of the commutator is retained as \(\chi_i\).
Comparing the norm to RJC.7 fixes exactly the integer powers
and analytic units in RJC.10. In particular these units are
retained, and their reciprocals give the inverse comparisons.
Their exponent identity is \(\rho=z^{-be_i}u_i^{-b}\).

The deck characters are the actual fibre-weight ratios:
\[
\begin{array}{c|c|c}
(a,b)&\chi_\infty&\chi_0\\ \hline
(1,6)&\zeta_3&-1\\
(2,3)&\zeta_3^2&1\\
(3,2)&1&-1\\
(6,1)&1&1.
\end{array}                                                        \tag{RJC.11}
\]
They give, with the complete cap generator \(A_{i,r}\),
\[
 L_i^\varepsilon A_{i,r}L_i^{-\varepsilon}
 =\chi_i^{-\varepsilon}A_{i,r}
 =\zeta_b^{Mn_i/m_i}A_{i,r}
 =A_{i,\zeta_b r}.                                                  \tag{RJC.12}
\]
The middle equality is exactly RJC.3 modulo two and three. For example,
in the first row \(kn_\infty\equiv\varepsilon\pmod3\) and
\(ln_0\equiv\varepsilon\pmod2\). The first congruence gives
\(\zeta_6^{4kn_\infty}=\zeta_3^{-\varepsilon}\), and the second gives
\(\zeta_6^{3ln_0}=(-1)^{-\varepsilon}\). The other rows have the same
displayed reductions with the trivial partner character retained.
Every prescribed common finite line character cancels in this
conjugation.

\par

\medskip\noindent\textbf{RJC.3. The full ordinary endpoint and its cyclic norm}

\par\nobreak

Write the actual one-gamma forward closing as
\[
 h_0(t)=c_0 q^\varepsilon t^{-mn_0},\qquad c_0\in\mathbb C^*,          \tag{RJC.13}
\]
with its original annular-reference constant \(c_0\).
Choose the prescribed scalar \(\gamma_0\) with
\(\gamma_0^b=c_0C_a^{-\varepsilon}\).
The rational root required by the full K67 comparison is literally
\[
 g_0=\gamma_0t^{d_0},\qquad
 g_0^b=h_0\rho^\varepsilon.                                         \tag{RJC.14}
\]
On \(\lambda=r^M\) the ordinary positive endpoint is therefore
\[
 E_r=\gamma_0 r^{-\varepsilon J_0}t^{d_0}U^\varepsilon.               \tag{RJC.15}
\]
It covers \(+\varepsilon G\). Its inverse evaluates \(U^{-\varepsilon}\)
at the actual inverse source point and uses the reciprocal displayed
base and scalar factors. At the finite root its total \(z\)-power is
\[
 hd_0-3h\varepsilon/b=-J_0n_0,
\]
and at infinity it is
\[
 -md_0+2m\varepsilon/b=-J_0n_\infty.                                 \tag{RJC.16}
\]
The full forward ordinary-to-cap cochain for the path winding \(J_0\)
contributes \(z^{J_0n_i}\). Thus the remaining cap arrow is
\(\gamma_0 r^{-\varepsilon J_0}u_i^\varepsilon L_i^\varepsilon\),
including the full invariant unit and finite character in RJC.10-\/-12.
The full annular terms and all original Cartier generators in
this square are given in RJC.4 below.

The \(b\) chronological source points of the parameter are
\(r,\zeta_b r,\ldots,\zeta_b^{b-1}r\). Their scalar product is
\[
 r^{-\varepsilon bJ_0}
 \exp\!\left(-\pi i\varepsilon J_0(b-1)\right).
\]
For all four rows \(J_0(b-1)\) is even: the even \(b\) rows have even
\(J_0\), and the odd \(b\) rows have even \(b-1\). Keeping the
reference constants, RJC.6,9,13 therefore give the exact ordinary
composite
\[
 E_{\zeta_b^{b-1}r}\cdots E_{\zeta_b r}E_r
       =h_0(t)B_\lambda^{-\varepsilon}.                             \tag{RJC.17}
\]
The K67 bounded origin has the additional constant
\[
 c=\exp\!\left(\frac{2\pi i\varepsilon M(b-1)}{2b^2}\right).
\]
Its exact family comparison is the physical scalar
\[
 F_r=r^d,\qquad d=\frac{\varepsilon J_0(b-1)}2\in\mathbb Z,
 \qquad F_{\zeta_b r}E_r=(cE_r)F_r.                                 \tag{RJC.18}
\]
The inverse is \(r^{-d}\). This keeps the finite phase even when \(c\)
is a nontrivial root of unity; its \(b\)-fold product is already the
same exact closing in RJC.17.

The next sections prove that this is a family comparison through
the complete roots, references and fans and that its powers
exhaust all holomorphic endpoints on this domain.

\par

\medskip\noindent\textbf{RJC.4. Both complete annular references and their scalar variation}

\par\nobreak

At site \(i\in\{0,\infty\}\), write \(q_0=t\),
\(q_\infty=1/t\), \(m_0=h,m_\infty=m\).
The actual root division section \(R_i\) satisfies
\(m_iR_i=bG'\). At a reduced root multiplication by \(m_i\)
is invertible on both the additive identity group and the finite
component group. At a good root use its analytic inverse branch
at \(O\). These give the unique marked analytic division section,
and the deck preserves it. Its negative translation is the
source part of the cap lift.

For \(P_i=aG'\), the chord-over-vertical function
\[
 g_{P_i,R_i}(p)=
 \frac{\ell_{P_i,R_i}(p)}{x(p)-x(P_i+R_i)}
 \tag{RJC.32}
\]
has divisor \((P_i)+(R_i)-(P_i+R_i)-(O)\).
It is precisely the point-line Hom coefficient over \(p\mapsto
p-R_i\). Its whole-fibre power is removed by the actual graph
generators RJC.20-\/-25; the full strict-section factors remain.
Finite averaging as in RJC.2 supplies the deck-equivariant lift.
For this division lift the remaining deck character is one.
At a good root \(R_i(0)=O\), and its central line map is a
constant scalar; its commutator with the deck is one there.
At a reduced root of order greater than one, the explicit
graphs RJC.24 and the reduced-IV graph show \(aG'(0)=O\).
The central line is trivial with its actual point/O
identification, so the central translation lift is again
a constant in that identification. At order one the deck
is identity. In each case the constant character left by
finite averaging equals its central value one.
The ratio of the \(m_i\)-fold norm to \(B_0\) is therefore an
invariant nonzero analytic base function. Taking the root fixed
by the prescribed reference gives the full commuting lift
\(\widetilde F_i^{\,m_i}=B_0\). Thus the
actual marked cap generator, including its prescribed character,
is
\[
 A_{i,\lambda}=R_{\lambda,i}^{\,n_i}
 \zeta_{m_i}^{d_i}\widetilde F_i^{\,n_i}D_i,\qquad
 R_{\lambda,i}^{\,m_i}=\lambda,\qquad
 A_{i,\lambda}^{\,m_i}=B_\lambda^{\,n_i}.
 \tag{RJC.33}
\]
The finite character \(d_i\), division-section branch and scalar
root are those of the original cap. The source and line maps
commute with \(B_\lambda\). In constant-period coordinates they
are translations in the commutative extension
\((\mathbb C\times\mathbb C^*)/\Lambda\); the Cartier extensions
above carry this identity through the whole root.

On the reference sectors use
\[
 s_{i\alpha}=\log_\alpha q_i/(2\pi i),\quad
 z_{i\alpha}=e^{2\pi i s_{i\alpha}/m_i},\quad
 c_{\lambda,i}=
 n_i\log R_{\lambda,i}/(2\pi i)+d_i/m_i,\quad
 f_{i\alpha}=\operatorname{Log}_\alpha\widetilde F_i.
 \tag{RJC.34}
\]
With the full original-coordinate and Cartier map \(I_{z_{i\alpha}}\)
on the annulus, the gluing arrow is
\[
 \Phi_{i\alpha}=
 \operatorname{Exp}
 \bigl(n_i s_{i\alpha}f_{i\alpha}
       +s_{i\alpha}c_{\lambda,i}\delta+h_{i\alpha}\bigr)
 I_{z_{i\alpha}}.
 \tag{RJC.35}
\]
It includes the negative source displacement
\(-n_i s_{i\alpha}R_i\), every point-line factor, and the full
physical scalar. On a labelled overlap
\(s_{i\beta}=s_{i\alpha}+k\),
\(f_{i\beta}=D_{i*}^kf_{i\alpha}+\Pi_{i\alpha\beta}\).
The complete cochain equation is
\[
 h_{i\beta}-D_{i*}^k h_{i\alpha}
 =-n_i s_{i\alpha}\Pi_{i\alpha\beta}
  +j_{i\alpha\beta}\operatorname{Log}_\beta B_\lambda
  +\Lambda_{i\alpha\beta}.
 \tag{RJC.36}
\]
The source periods \(\Pi,\Lambda\) and integral suspension
columns \(j\) remain those of the specified reference.
An explicit solution in its ordinary Lie frame starts with its
fixed smooth extension \(h^0\), its residual cocycle \(E\) and
a partition \(\chi_\beta\):
\[
 k_\alpha=-\sum_\beta\chi_\beta E_{\alpha\beta},\qquad
 v(q)=\frac1\pi\int
 \frac{\bar\partial k(w)}{q-w}\,dA(w),\qquad
 h_\alpha=h_\alpha^0+k_\alpha-v.
 \tag{RJC.37}
\]
The residual is supported on the prescribed compact subannulus.
The common holomorphic constant of \(v\) is fixed by the original
reference value; equivalently subtract \(v(q_{\rm ref})\) and
apply the corresponding common constant to the rigidification.
This gives exactly that value and leaves every integral period
in RJC.36 unchanged. Subtracting the target-sheet logarithm
\(kn_i f_{i\beta}+kc_{\lambda,i}\delta\) proves
\[
 \Phi_{i\beta}=
 A_{i,\lambda}^{\,k}B_\lambda^{\,j_{i\alpha\beta}}\Phi_{i\alpha},
 \qquad
 \Phi_{i\alpha}B_\lambda=B_\lambda\Phi_{i\alpha}.
 \tag{RJC.38}
\]
The inverse is the inverse root map followed by the negative
of the displayed Lie translation at the inverse source point,
with the reciprocal point and scalar coefficients.

For scalar variation put
\(\lambda_\theta=e^{2\pi i\theta}\lambda_*\), and keep the
actual integer representatives \(J_{i\alpha}\) with
\(J_{i\beta}-J_{i\alpha}=j_{i\alpha\beta}\).
The full reference restoration is
\[
 \Phi_{\theta,i\alpha}
 =e^{2\pi i n_i\theta s_{i\alpha}/m_i}
 B_\theta^{\,J_{i\alpha}}
 \bigl(B_*^{-J_{i\alpha}}\Phi_{*,i\alpha}\bigr).
 \tag{RJC.39}
\]
On an integer winding \(J\), the last two factors have ratio
\(e^{2\pi iJ J_{i\alpha}}=1\), and the first factor gives
\[
 \Phi_{\theta+J,i\alpha}
 =z_{i\alpha}^{\,n_iJ}\Phi_{\theta,i\alpha}.
 \tag{RJC.40}
\]
This establishes the root factors used in RJC.16 with every
source period and line reference still present. The remaining
cap map is the complete \(u_i^\varepsilon L_i^\varepsilon\)
map of RJC.10, multiplied by the specified scalar.
RJC.12 supplies its finite generator square.

The integral form of that same square is the K67 cochain
\[
 \frac{P_i\delta-\varepsilon(A_i-I)e_\gamma}{b}
 =\bigl(0,-\varepsilon k_{i,2},\varepsilon k_{i,1},
                (P_i+\varepsilon z_i)/b\bigr)^t,\qquad
 P_i=Mn_i/m_i.
 \tag{RJC.41}
\]
Here \(z_i=2\) at infinity and \(-3\) at zero are representatives
of the actual characters \(\chi_i=\zeta_b^{z_i}\).
RJC.8-\/-12 make the last coordinate integral, and the source
coordinates are the original integral Abel-period columns.
This exhibits the entire ordinary/root lattice arrow.

\par

\medskip\noindent\textbf{RJC.5. The intrinsic base and full affine endpoint map}

\par\nobreak

The integral chain maps F46.J.7-\/-12 give the actual sphere
topology under RJC.3: their degree-two map is an integral
isomorphism and both degree-one and degree-three determinants
are \(-\varepsilon\). The cap fundamental-group presentation
reduces to the cyclic group of order \(|\varepsilon|=1\).
The common blowups and both ruling maps of F35 carry these
integral maps through each prescribed fan. In particular every
object here has \(H^2=H^4=0\).

Let \(f\) be a meromorphic function on one of these compact
objects. A horizontal component of its polar divisor restricts
to a nonempty effective divisor on a generic compact two-torus
and has positive pairing with a flat Kähler form there.
Its cohomology class is the restriction of a class in
\(H^2(X)=0\), a contradiction. All poles are vertical, and
restriction to each regular compact torus is constant.
The original transverse maps at special components are
\(u\mapsto t_0+u^e\). Their invariant Laurent expansions
contain precisely multiples of \(e\); hence the constants
descend meromorphically across each special value. We obtain
the actual field isomorphism
\[
 \pi^*:\mathbb C(t)\xrightarrow{\;\cong\;}\mathbb C(X).
 \tag{RJC.42}
\]
Its inverse is restriction to a regular fibre value followed
by the just-described Laurent descent. A holomorphic endpoint
and its inverse consequently induce inverse Möbius maps of
the original base.

The complete generic point-line period module is, for all four
rows,
\[
 \gamma=(-bg,\ell),\quad e_1=(1,0),\quad
 e_2=(\tau,ag),\quad\delta=(0,1).
 \tag{RJC.43}
\]
The full theta/Miller reference constant occurs in \(\ell\);
both point Cartier modules are included in these coordinates.
For a source loop \(\alpha\) its monodromy is
\[
 M_\alpha=
 \begin{pmatrix}
 1&0&0\\
 b v_\alpha&A_\alpha&0\\
 c_\alpha&a v_\alpha^t J A_\alpha&1
 \end{pmatrix},\quad
 J=\begin{pmatrix}0&-1\\1&0\end{pmatrix},
 \tag{RJC.44}
\]
with the \(1,2,1\) block order. Direct multiplication gives
\(v_{\alpha\beta}=v_\alpha+A_\alpha v_\beta\) and
\[
 c_{\alpha\beta}=c_\alpha+c_\beta+
 ab\,v_\alpha^tJ A_\alpha v_\beta.
 \tag{RJC.45}
\]
Thus no central cochain has been removed.

The Abel cocycle \(v_\alpha\) is nonzero over \(\mathbb Q\).
A rational coboundary, after clearing its denominator and the
finite component orders, would give a global holomorphic Lie
logarithm of a nonzero multiple of \(G\). On the minimal good
III chart the logarithm of \(2G'\) divided by its actual
factor \(w\) extends; the order-four elliptic action has no
invariant period. On the minimal good infinity chart
\(Y^2=X^3+wX+1\), the logarithm of \(3G'\) has leading term
\(-w^2/8\); division by \(w^2\) gives the original extending
Lie coefficient. Its order-three elliptic action has no
invariant period. At the selected node the invariant annular
period extends through \(xy=q\).
The original Lie transition is
\[
 u_\infty(q_\infty)
 =q_\infty^{-1}u_0(q_\infty^{-1}).
 \tag{RJC.46}
\]
An entire \(u_0\) satisfying this transition tends to zero
at infinity and is zero. But \(3G\) has original infinity
value \((64,512)\) and additive parameter \(-1/8\);
each of its nonzero multiples has nonzero additive parameter.
This contradiction proves the asserted Abel obstruction.
The higher-root diagrams RJC.20-\/-25 are the pullbacks of
these actual minimal diagrams, so the argument applies to
every cap order in RJC.2.

The vector and dual equations of RJC.44 now give
\[
 \Lambda^{\pi_1}=\mathbb Z\delta,\qquad
 (\Lambda^*)^{\pi_1}=\mathbb Z\psi,\quad
 \psi(\gamma)=1,\quad\psi(e_1)=\psi(e_2)=\psi(\delta)=0.
 \tag{RJC.47}
\]
A nonzero \(\gamma\) coefficient in an invariant vector makes
the Abel cocycle a coboundary. With that coefficient zero,
the order-three infinity matrix kills the remaining invariant
elliptic vector. The dual equation has the same nonzero Abel
obstruction. These are primitive integral subgroups, so the
extra deck of the kernel cover is the actual primitive
\(\gamma\). The kernel cover is the nonzero point line
with its actual Cartier modules.

An endpoint preserves these two subgroups with their signs.
The derivative of its universal-cover lift is periodic for
the full compact torus lattice and therefore constant.
Complex linearity and \(\delta\mapsto e\delta\) give the
full affine map and its inverse
\[
 \begin{split}
 (z,w)&\longmapsto(Az+z_0,Cz+e w+w_0),\\
 (z',w')&\longmapsto
 \bigl((z'-z_0)/A,\,
 e(w'-w_0-C(z'-z_0)/A)\bigr),
 \quad A\ne0,\quad e=\pm1.
 \end{split}
 \tag{RJC.48}
\]
On the kernel lattice its complete integral matrix and inverse
are
\[
 K=\begin{pmatrix}p&q_1&0\\r&s&0\\u&v&e\end{pmatrix},
 \quad
 K^{-1}=
 \begin{pmatrix}
 s&-q_1&0\\-r&p&0\\e(vr-us)&e(uq_1-vp)&e
 \end{pmatrix},\qquad ps-q_1r=1.
 \tag{RJC.49}
\]
The positive determinant follows from the oriented area of
the nonzero complex-linear first coordinate. The first
coordinates of RJC.48 are actual inverse elliptic maps.
For this source
\[
 j_0(t)=\frac{6912}{27t+4}.
 \tag{RJC.50}
\]
The equality \(j_0(h(t))=j_0(t)\) forces \(h(t)=t\).
The endpoint therefore fixes the original base everywhere,
including its two original cap centres and selected height.

The full linear centralizer of RJC.44 is \(eI+c\delta\psi\).
For the elliptic block, the infinity matrix generates
\(\mathbb Q(\zeta_3)\); the selected unipotent has a nonzero
nilpotent part. A nonscalar element of that quadratic field
generates the field, whose nonzero elements are invertible.
Its simultaneous centralizer with the unipotent is thus
\(\mathbb QI\). The two off-diagonal block equations use
the nonzero Abel class to identify the adjoining scalars
with this one and to set their elliptic vectors to zero.
The central cochains remain in the calculation and leave
exactly \(c\delta\psi\). Integral unimodularity gives
\(e=\pm1\).
The two complex-independent period vectors \(e_1,\delta\)
force the derivative to be exactly \(eI\). The \(\gamma\)
equation then says that the two suspension logarithms
differ by an integer. Thus an endpoint has exactly equal
original scalar values at its two ends.

\par

\medskip\noindent\textbf{RJC.6. Extending the source displacement and the full section lattice}

\par\nobreak

Consider first a positive endpoint. Its ordinary lift has the
form \(H(t)U^j\) after extending its source displacement.
We give the extension rather than assume it. Pull the cap
map and its inverse to the normalized fibre product
\(z_i^{m_i}=q_i\). On the actual point-line cover, a
nonzero vector over the \(O\)-section kills the root meridian;
the remaining vanishing cycles lie in \(\ker\psi\).
These are the Cartier charts RJC.20-\/-25, so the lift is
single-valued on the full root source.
Its commutator with the cap generator is an integer power
\(B^s\). Taking the \(m_i\)-fold relation gives
\(B^{m_i s}=1\). The original strict contraction has no
nonzero finite-order deck, hence \(s=0\).
The source displacement is therefore an equivariant
holomorphic root section. The original-coordinate maps
RJC.4-\/-5 descend it, and the proper minimal regular
source resolves the indicated graph.

At the selected fibre a pole of order \(-K\) in the
physical coefficient is cancelled by the actual \(B^K\)
reindex. Its remaining source displacement is shifted by
\(-bKG\), and its scalar coefficient is a nonzero
holomorphic torus unit. In the two selected vertex
charts these maps are the unit maps displayed in
RJC.65 below, with reciprocal inverses.
They extend through each vertex and open curve axis.
The \(B\)-labels of two ordinary local lifts differ by
a locally constant integer. After the selected unit
has been joined, their joining torsor is over the
two-point root sphere with presentation
\[
 \langle x,y\mid x^m=y^h=xy=1\rangle.
 \tag{RJC.51}
\]
This group is cyclic of order \(\gcd(m,h)=1\).
Its integer joining torsor is trivial, with the
constant join fixed at the original reference.
The source displacement extends to an actual global
section of the original elliptic surface.

The section group here is exactly \(\mathbb ZG\).
To verify the integral index, use the birational chart
\[
 u=\frac{y-t-x/2}{x},\quad
 v=\frac{y+t+x/2}{x},\quad
 x=uv+\frac14,\quad
 t=\frac{x(v-u-1)}2,\quad
 y=\frac{x(u+v)}2.
 \tag{RJC.52}
\]
The displayed reciprocal functions are inverse on their
common open. The source is a rational elliptic surface
with Euler number \(8+3+1=12\).
Its ten classes
\[
 (O,F,C_0,\ldots,C_5,A_1,G)
 \tag{RJC.53}
\]
have \(O^2=G^2=-1\), \(OF=GF=1\),
the negative \(E_6\) Cartan block with edges
\((0,1),(1,2),(2,3),(3,4),(2,5)\), and \(A_1^2=-2\).
The additional nonzero intersections are
\(G C_0=G A_1=1\); all unspecified ones are zero.
Their determinant is \(-1\).
They span the entire rank-ten unimodular source lattice.
Quotienting by its first nine classes is free cyclic,
with generator the image of \(G\). This is the full
section quotient, with no remaining torsion or index.
Its height is
\(2-4/3-1/2=1/6\).
Thus the extended displacement above is \(jG\) for
an integer \(j\), and its full ordinary point-line
map is precisely \(H(t)U^j\), up to its physical
scalar already included in \(H\).

\par

\medskip\noindent\textbf{RJC.7. All positive endpoints and the physical automorphism group}

\par\nobreak

For one object, the exact zero and infinity comparisons
RJC.10 and the selected unit reindex give
\[
 \operatorname{ord}_0H=3j/b,\quad
 \operatorname{ord}_\infty H=-2j/b,\quad
 \operatorname{ord}_{t_c}H=-K.
 \tag{RJC.54}
\]
There are no other zeros or poles: the map and inverse
are nonzero Hom arrows on every other fibre.
The degree of this principal divisor gives \(j=bK\),
and \(H=c(t^3/q)^K\).
Using the exact \(C_a\) in RJC.6, the quotient map is
the physical scalar
\[
 H U^{bK}=cC_a^K B_0^{-K}
        =cC_a^K\lambda^K B_\lambda^{-K}
        \quad\text{on the covering line.}
 \tag{RJC.55}
\]
All physical scalars extend by the displayed root,
Cartier and selected unit maps. A scalar acting as an
identity deck would translate the source by a multiple
of \(bG\); the nonzero Abel class shows that \(G\) has
infinite order, so that multiple is zero and the scalar
is one. The physical action is faithful.

A negative automorphism of this one object has
sixteen transverse fixed points on each regular
two-torus, because its affine linear part is \(-I\).
Its compact analytic equalizer contains a horizontal
curve of positive base degree. The degree equals
pairing that curve with \(\pi^*H^2(\mathbb P^1)\),
whose image in \(H^2(X)\) is zero. The contradiction
excludes every negative automorphism. We have the
actual group isomorphism, with its explicit inverse
from RJC.55,
\[
 \mathbb C^*_{\rm physical}
 \xrightarrow{\;\cong\;}\operatorname{Aut}_{\rm hol}(X).
 \tag{RJC.56}
\]

For a positive endpoint of scalar winding \(J\), RJC.40
contributes both root factors. Its full divisor equations
are
\[
 \begin{split}
 \operatorname{ord}_0H&=\frac{3j}{b}-\frac{n_0J}{h},\\
 \operatorname{ord}_\infty H&=-\frac{2j}{b}
                                  -\frac{n_\infty J}{m},\\
 \operatorname{ord}_{t_c}H&=-K.
 \end{split}
 \tag{RJC.57}
\]
These are integer orders of the single-valued ordinary
base function; all corresponding local units were retained
in RJC.10,35-\/-40. Their degree and RJC.3 give
\[
 \frac jb-K-\frac{\varepsilon J}{mh}=0,\qquad
 j=bK+\frac{b\varepsilon J}{mh}.
 \tag{RJC.58}
\]
It follows that \(M\mid bJ\), or \(J\in J_0\mathbb Z\).
The endpoint RJC.15 attains \(J_0\). Its chronological
powers and inverse-source inverses attain every integer
multiple. At a fixed winding the ratio of two endpoints
is an automorphism RJC.56, hence exactly a physical
scalar with its full extending map.

\par

\medskip\noindent\textbf{RJC.8. Single-valued negative endpoints}

\par\nobreak

The selected inversion
\[
 (Z,V,q)\longmapsto(Z^{-1},V^{-1},q)
 \tag{RJC.59}
\]
reverses both full periods. Negating a lattice-basis
triangle gives its opposite triangle, translated by
\(-p-u-v\). The actual vertex exchange in R0E.90
preserves \(q\), the factors \(\kappa g,\alpha\eta\),
and their reciprocal factors. It therefore extends
on every selected fan, with itself as inverse.

A negative endpoint has single-valued nonzero-line
lifts on the ordinary and root pieces: the point/O
vector and root-cover argument of RJC.6 applies to
its inverse as well. Composing its selected arrow
with RJC.59 gives the extending positive selected
unit arrow, so its continuation has no additional
selected integer monodromy. The constant integral
joins use RJC.51. Both lifts reverse the full
suspension,
\(\widetilde F B=B^{-1}\widetilde F\).

On the original infinity annulus the complete fixed
reference is \(A_\infty B^\nu\), with the actual
retained integer \(\nu\). One root circle consequently
has reference \(B^N\), where
\[
 N=n_\infty+m\nu.
 \tag{RJC.60}
\]
The two scalar endpoints have the same original
scalar and the same reference operators by RJC.39-\/-40.
Choose the one constant deck making their two lifts
equal on a cut annulus. Continuation around the
root circle gives
\[
 B^N\widetilde F B^{-N}=B^{2N}\widetilde F.
 \tag{RJC.61}
\]
Single-valuedness implies \(B^{2N}=1\), and strict
contraction gives \(N=0\).
But \(\gcd(m,n_\infty)=1\) and every allowed \(m\)
is at least two, so \(n_\infty+m\nu\ne0\) for
every integer \(\nu\). This excludes all negative
endpoints, including the quadratic reduced infinity
root. The complete holomorphic winding group on
each stated joint domain is therefore exactly
\(J_0\mathbb Z\).

\par

\medskip\noindent\textbf{RJC.9. The entire parameter-family comparison on every fan}

\par\nobreak

We now carry RJC.18 through the actual selected charts.
Let \(v=\log r/(2\pi i)\), let \(\Gamma\) be the bounded
Tate logarithm of \(+bG\), and retain the full ordinary
closing logarithm
\(H_0=\varepsilon s+U_0(q)\).
Here \(s=\log q/(2\pi i)\); the zero-cap factor
\(-mn_0\log t/(2\pi i)\) and the fixed constant
\(\log c_0/(2\pi i)\) are included in \(U_0\), which is
holomorphic near the selected \(t_c\ne0\).
The full outward suspension period there is
\(\gamma=(\Gamma,C-s-Mv)\).
The bounded endpoint displacement is
\[
 \tau(v)=\frac{H_0\delta+\varepsilon\gamma}{b}
 =\left(\frac{\varepsilon\Gamma}{b},
 \frac{U_0+\varepsilon C-\varepsilon Mv}{b}\right).
 \tag{RJC.62}
\]
Every entry is holomorphic on the whole selected disc.
Put \(\alpha=\varepsilon M(b-1)/(2b^2)\) and
\[
 \sigma(v)=
 \left(\varepsilon\Gamma v,\,
 v(U_0+\varepsilon C)
       -\frac{\varepsilon M}{2}(v^2-v)\right).
 \tag{RJC.63}
\]
Direct subtraction gives
\(\sigma(v+1/b)-\sigma(v)=\tau(v)+\alpha\delta\).
Thus \(\xi=z-\sigma(v)\) is the closed native coordinate
for the forward endpoint \(cE_r\).
For every cut interval,
\[
 C_s(v)=\exp(\sigma(v)-\sigma(v+s)),\quad
 P_j(v)=\exp(\sigma(v)-\sigma(v+j/b)),
 \tag{RJC.64}
\]
and \(C_s(v+j/b)P_j(v)=C_{j/b+s}(v)\) literally.
Their inverse maps have the negative exponents.
At \(j=b\) the product is
\(\exp(-H_0\delta-\varepsilon\gamma)\), the full
one-gamma closing. This equality holds on parameter
neighborhoods, so the first and last collars agree
with all derivatives.

For a lattice basis \(P\in GL_2(\mathbb Z)\), conjugate
the full period matrix by \(P\); all factors
\(\kappa,\alpha_0,g,\eta\) stay in that matrix.
The log characters are \(P^{-1}(z_1,z_2)^t\).
The physical scalar \(F_r=r^d\) of RJC.18 acts on
their multiplicative coordinates by
\[
 A=r^{d(P^{-1})_{12}},\qquad
 D=r^{d(P^{-1})_{22}}.
\]
For these factors, or for the full two factors
of any \(C_s,P_j\), the actual two vertex maps are
\[
 \begin{split}
 (a,b_1,c_1)&\longmapsto
 ((AD)^{-1}a,A b_1,D c_1),\\
 (u,v_1,w_1)&\longmapsto
 (ADu,D^{-1}v_1,A^{-1}w_1).
 \end{split}
 \tag{RJC.65}
\]
Their inverses use \(A^{-1},D^{-1}\). They preserve
\(q=ab_1c_1\) and intertwine every full transition,
because each transition is the specified monomial
with its unchanged source coefficient.
For example the vertical fan \(T_j\) retains
\[
 \kappa_j=\kappa g^j,\quad
 \alpha_j=\alpha_0\eta^j\kappa^{-j}g^{-j^2},\quad
 g_j=g,\quad\eta_j=\eta g^{-j},
 \tag{RJC.66}
\]
and its three full transition rows are
\[
 \begin{array}{c|ccc}
 &u&v_1&w_1\\ \hline
 C_0&\kappa_j\alpha_j b_1c_1&
       \alpha_j^{-1}ab_1&\kappa_j^{-1}b_1^{-1}\\
 C_1&g_j\eta_j b_1c_1&
       \eta_j^{-1}c_1^{-1}&g_j^{-1}ac_1\\
 C_2&a^{-1}&ab_1&ac_1.
 \end{array}
 \tag{RJC.67}
\]
Substitution of RJC.65 gives the same character
factor on each side, with every source unit present.

The corresponding native Cox map is
\[
 (a_0,a_1,c_0,c_1)\longmapsto
 (a_0,Aa_1,(AD)^{-1}c_0,Dc_1).
 \tag{RJC.68}
\]
It preserves the Cox relation and descends to both
native charts. Its reciprocal map proves the
inverse on the exceptional curve and the normal
bundle. Normal rescaling retains the complete
\(q\)-jets: \(q\) has order two in the normal
coordinates of an open curve axis and order three
at a vertex, so the remaining source dependence
is precisely the higher normal jet carried by
that rescaling.
RJC.64 is an equality of both full characters,
and hence of RJC.65 and RJC.68 on every sheet.
The ordinary, cap and Cartier maps are
RJC.15-\/-18 and RJC.35-\/-41.
They prove that \(F_r\), with inverse \(r^{-d}\),
is an isomorphism of the complete parameter
families, including every endpoint and intermediate
seam.

The action \(r\mapsto\zeta_b r\) on the parameter
punctured disc is free. Its quotient is \(w=r^b\),
with \(\lambda=w^{J_0}\).
The complete \(r\)-family is the literal degree-\(b\)
pullback of this shorter family.
This is a covering of the completed orbit pieces:
the physical line, cap identification and stable
tangent bundle pull back by the displayed maps.
Integration of their integral characteristic
classes multiplies by \(b\). This supplies the
integer characteristic comparison used below,
before reduction modulo \(28\).

\par

\medskip\noindent\textbf{RJC.10. The original page, its line and the primitive pair}

\par\nobreak

The actual canonical page over an ordinary binding
value \(\tau\) has source-base curve
\[
 \eta^h-\tau\xi^m=1,\qquad
 t=\tau+\xi^{-m}=\eta^h\xi^{-m}.
 \tag{RJC.69}
\]
The reciprocal functions on its ordinary open are
\(\xi^m=(t-\tau)^{-1}\) and
\(\eta^h=t/(t-\tau)\), with all chosen roots
given by the original cap labels.
The degree is \(M=mh\); the three branch indices
at \(0,\infty,\tau\) are \(h,m,M\).
Riemann-\/-Hurwitz gives genus
\((m-1)(h-1)/2\).
The original source page is the pullback by
these actual maps, followed by the minimal root
models RJC.4-\/-5 and their stated resolutions.

The source translation comparison of the page
uses the rigidified relative Picard evaluation
of O4-\/-O5 in the proved order-one-page construction.
It is worth specifying its map on this source.
On each translation-parameter chart, form the
ratio of the original pulled-back point line
and the translated degree-zero root point line.
On every smooth fibre it is trivial; direct
image and evaluation supply its scalar base line.
The evaluation divisor has zero intersection
with every component by RJC.20-\/-25.
The connected Kodaira intersection matrix
therefore makes it an integral sum of whole
fibres. Division by their actual base equations
gives an invertible evaluation map, with
reciprocal inverse, at singular as well as
smooth fibres.
Rigidifying at \(O\) and pulling back by the
smooth translating section gives the exact
page-line comparison.
Every root character and full delta transition
remains as base-line holonomy.
The only resulting change of the line class
is \(c_1(N)\mapsto c_1(N)+zF\), where
\[
 c_1(N)F=F^2=0.
 \tag{RJC.70}
\]
Thus its square is unchanged. On the binding
collar this is the specified comparison; the
base line extends over the ordinary disc and
has zero boundary square ambiguity by the
full K5 relative gauge formula. The stable
tangent comparison is the same source map.

Let \(e_\infty,e_0\) be the Euler numbers of
the actual minimal root fibres:
\((0,0),(0,3),(4,0),(4,3)\) in the four
rows of RJC.2.
The single original I1 away from the cap
centres has \(M\) preimages counted with
degree, infinity has \(h\) root fibres,
and zero has \(m\) root fibres.
The page height and holomorphic Euler number
are therefore
\[
 H=\frac{Ma^2}{6},\qquad
 \chi=\frac{M+h e_\infty+m e_0}{12}.
 \tag{RJC.71}
\]
The height \(1/6\) is the exact original
section height from RJC.53; the line square
is \(-H\).
All singular line corrections have been
carried before this intersection calculation.

For the signed equation RJC.3, the actual
binding coordinate is
\[
 q_b=(t-\tau)(t/\tau)^{-\varepsilon m n_0}.
 \tag{RJC.72}
\]
It has derivative one at \(\tau\), and its
entire unit factor remains in the graph.
After retaining its constant phase, the
normal angles satisfy
\(\alpha_{\rm normal}=\phi-M\beta\), with
\(\beta=\varepsilon\gamma\).
The source-preserving boundary graph has
\(\beta=\theta+g(p)\).
K30\textquotesingle s triangular tangent shear carries
the derivatives of \(g\) and the full unit
in RJC.72; its inverse subtracts them.
It gives the physical Hopf normal sphere
with the same positive parameter orientation.
Equivalently the complete origin obeys
\[
 F_\varepsilon(v+\theta)-F_\varepsilon(v)
 =\varepsilon\theta\gamma
   -\frac{\varepsilon M}{2}\theta(\theta-1)\delta.
 \tag{RJC.73}
\]
Both the source graph and this quadratic
term are retained for each sign.

The full physical dual-line and inverse-clutch
maps of K68-\/-K69 now give, on the literal
shorter cover RJC.64-\/-68,
\[
 (k_{\rm phys},\ell_{\rm phys})
   =(3H/b,24\chi/b),\qquad
 (m_{\rm native},n_{\rm native})
   =(6\chi/b,\,3H/b-6\chi/b).
 \tag{RJC.74}
\]
These follow by the oriented degree-\(b\)
covering of integral characteristic pairs,
so no division in the finite group is used.
Substitution of the actual root Euler
numbers gives the integer table
\[
 \begin{array}{c|c|c|c}
 (a,b)&k_{\rm phys}&\ell_{\rm phys}&
 S_0=(4k_{\rm phys}-\ell_{\rm phys})/24\\ \hline
 (1,6)&kl&4kl&0\\
 (2,3)&2l(4k+1)&8l(k+1)&kl\\
 (3,2)&9k(3l+2)&12k(l+2)&k(4l+2)\\
 (6,1)&18(3l+2)(4k+1)&
 24\bigl(k(l+2)+l+1\bigr)&k(35l+22)+8l+5.
 \end{array}
 \tag{RJC.75}
\]
For each row the native pair is also exactly
\((\ell_{\rm phys}/4,
k_{\rm phys}-\ell_{\rm phys}/4)\).
The proof retains both signs of RJC.3;
the positive scalar-parameter orientation
in RJC.72-\/-73 gives these displayed pairs
for both.

\par

\medskip\noindent\textbf{RJC.11. The full fan algorithm and its exact image}

\par\nobreak

On the vertical fan \(T_j\), the primitive
return and pair are
\[
 S_j=S_0-\frac{J_0}{2}j(j-1),\quad
 k_j=k_{\rm phys}-3J_0j(j-1),\quad
 \ell_j=\ell_{\rm phys}.
 \tag{RJC.76}
\]
For every other lattice-basis fan the exact
K73 inputs, now carried by RJC.62-\/-74, are
\[
 J=J_0,\qquad \widehat H=H/b,\qquad
 \widehat\chi=\chi/b.
 \tag{RJC.77}
\]
To make the computation fully specified,
take the canonical finite word from the
actual identical-fan maps and terminating
Euclidean subtraction of K71.AF.1-\/-2.
At its vertical seed \(T_j\), start with
\[
\begin{gathered}
 V=\begin{pmatrix}
 2J-6\widehat H+6Jj(j-1)&2J(1-2j)&2J(2j-1)\\
 2J(1-2j)&2J&0\\
 2J(2j-1)&0&2J
 \end{pmatrix},\\
 D=\begin{pmatrix}
 2J+6\widehat H-24\widehat\chi-6Jj(j-1)\\
 2J(2j-1)\\2J(1-2j)
 \end{pmatrix},\\
 S=\widehat H/2-\widehat\chi-\frac J2j(j-1).
\end{gathered}
\tag{RJC.78}
\]
The matrix and vector are the actual double-cover
variation and stable-derivative invariants.
At a directed step \(P'=PQ_i\), let
\(a_*=V_{ii},b_*=D_i\), and set
\[
 \begin{split}
 B&=2J(P')^{-1}e_2e_2^tP'=(b_{uv}),\\
 L&=(b_{22}-b_{21},\,b_{11}-b_{12},\,b_{22}),\\
 (x,y)^t&=(P')^{-1}e_2,\quad
 w=(x,y,-x-y),\\
 d&=(-x-2y,-2x-y,x-y),\\
 I_r&=Jw_r+d_r(L_r-J),\quad
 W_r=Jw_r-d_r(L_r-J).
 \end{split}
 \tag{RJC.79}
\]
For the actual nonvertical step all three
target weights are nonzero. The unique
target prism solution is
\[
 \begin{split}
 C_*&=w_2^2a_*-w_2I_2,\quad
 E_*=w_2(2a_*-b_*)-W_2,\\
 V'_{rs}&=(I_r/w_r)\delta_{rs}+C_* /(w_rw_s),\\
 D'_r&=(W_r+E_*)/w_r,\qquad
 S'=S+(b_*-a_*)/4.
 \end{split}
 \tag{RJC.80}
\]
These are the full target systems of K73,
whose exterior, physical line, tangent
and two cap arrows have now been matched
to this joint source. In particular the
native resets are the maps RJC.65-\/-68,
including each full source coefficient.
Every intermediate free invariant and
increment is integral. The actual fan
return is \(S(P)\bmod28\).
The Euclidean word terminates for every
integral unimodular \(P\), and its
reciprocal chart maps give the identical
fan inverses. This proves coverage of
all basis fans, beyond the finite tests.

We compute the exact image over all fans
and all integer scalar windings for
each fixed \(k,l\).
The prism equations and their unique
solutions are linear in
\((J,\widehat H,\widehat\chi)\).
Thus scaling those three inputs scales
the full \(V,D,S\) at every step.
All four comparisons below use an actual
admissible unscaled source profile, so
its return is an integer for every fan.

For \((a,b)=(1,6)\), the inputs are
\(kl\) times the actual \(k=l=1\)
profile. Therefore every \(S(P)\)
is divisible by \(kl\).
The explicit fan
\(P=\left(\begin{smallmatrix}1&-2\\0&1\end{smallmatrix}\right)\)
has word \(Q_0Q_0\) from \(T_0\).
For the unscaled profile,
\((J,\widehat H,\widehat\chi)=(2,1/3,1/6)\);
the two selected pairs \((V_{00},D_0)\)
are \((2,2)\), then \((-4,8)\).
The seed is zero and the increments
are \(0,3\), so \(S=3kl\).
Since \(3\) is a unit modulo \(28\),
the exact image is \(kl\mathbb Z/28\).

For \((a,b)=(2,3)\), all inputs are
\(l\) times those of the actual
\((k,1)\) profile. Its admissibility
follows from
\(\gcd(3l,4k+1)=1\).
Consequently every return is divisible
by \(l\). The normalized vertical
integer is
\[
 S_j/l=k-(4k+1)j(j-1)/2.
 \tag{RJC.81}
\]
For odd \(k\), use \(j=0\) when
\(k\not\equiv0\pmod7\) and \(j=4\)
when \(k\equiv0\pmod7\).
The coefficients are \(k\) and
\(-23k-6\), both odd and nonzero
modulo \(7\) in their stated cases.
For even \(k\), use \(j=2\) when
\(k\not\equiv2\pmod7\) and \(j=3\)
when \(k\equiv2\pmod7\).
The coefficients are \(-3k-1\)
and \(-11k-3\), again odd and
nonzero modulo \(7\).
Each supplies a unit modulo \(28\);
the exact image is \(l\mathbb Z/28\).

For \((a,b)=(3,2)\), admissibility
\(\gcd(3l+2,4k)=1\) forces \(l\)
odd. All inputs are \(k\) times
those of the actual \((1,l)\)
profile; every fan return is
divisible by \(k\).
Use
\(P=\left(\begin{smallmatrix}1&-3\\0&1\end{smallmatrix}\right)\),
with word \(Q_0Q_0Q_0\) from \(T_0\).
Its full selected target-system data,
divided by \(k\), are
\[
 \begin{array}{c|c|c|c}
 \text{stage}&V_{00}/k&D_0/k&(D_0-V_{00})/(4k)\\ \hline
 T_0&-42l-28&54l+20&24l+12\\
 Q_0&-204l-136&312l+176&129l+78\\
 Q_0Q_0&-474l-316&1110l+692&396l+252.
 \end{array}
 \tag{RJC.82}
\]
Adding these increments to the seed
\(4l+2\) gives \(553l+344\).
This is odd for odd \(l\), and
\(553l+344=7(79l+49)+1\).
It is a unit modulo \(28\).
The exact image is \(k\mathbb Z/28\).

For \((a,b)=(6,1)\), put
\(A=k(35l+22)+8l+5\) and
\(J=(3l+2)(4k+1)\).
The four explicit vertical integers
needed are
\[
 \begin{split}
 S_0&=A,\\
 S_4&=-k(37l+26)-10l-7,\\
 S_2&=k(23l+14)+5l+3,\\
 S_3&=-k(l+2)-l-1.
 \end{split}
 \tag{RJC.83}
\]
When \(kl\) is even, \(S_0,S_4\)
are both odd. Their residues modulo
\(7\) are \(k+l+5\) and
\(-k(2l+5)-3l\).
Common zero would give \(k=2-l\)
and \(2(l-4)^2=0\pmod7\), hence
\((l,k)=(4,5)\pmod7\).
When \(kl\) is odd, \(S_2,S_3\)
are both odd. Their common zero
has \(l\ne5\pmod7\), since
\(S_3\) is then nonzero.
Eliminating \(k\) gives
\(3l^2+11l+6=3(l-4)^2=0\pmod7\),
and again \((l,k)=(4,5)\pmod7\).
But that residue makes both
\(3l+2\) and \(4k+1\) divisible
by \(7\), contradicting the
actual coprimality in RJC.3.
Thus one of the displayed two
odd integers in the relevant
parity case is a unit modulo
\(28\). The exact image is all
of \(\mathbb Z/28\).

The result is the full table
\[
 \begin{array}{c|c|c}
 (a,b)&\text{holomorphic winding group}&
 \text{all-fan, all-loop return image}\\ \hline
 (1,6)&2kl\,\mathbb Z&kl\,\mathbb Z/28\\
 (2,3)&l(4k+1)\,\mathbb Z&l\,\mathbb Z/28\\
 (3,2)&2k(3l+2)\,\mathbb Z&k\,\mathbb Z/28\\
 (6,1)&(3l+2)(4k+1)\,\mathbb Z&\mathbb Z/28 .
 \end{array}
 \tag{RJC.84}
\]
The displayed nonvertical witnesses use
the entire prism solves, not a polynomial
extrapolation in the fan entries.

\par

\medskip\noindent\textbf{RJC.12. Every holomorphic scalar path and both mapping-torus orientations}

\par\nobreak

Let a holomorphic scalar path on its
parameter annulus have a holomorphic
endpoint identification on these fixed
marked objects. RJC.5-\/-8 forces equal
scalar values at its ends and winding
\(J=sJ_0\), \(s\in\mathbb Z\).
On the universal logarithmic strip
choose its full logarithm \(L\), with
\(L(\theta+1)-L(\theta)=2\pi iJ\).
The explicit deformation
\[
 L_u(\theta)=(1-u)L(\theta)
       +u\bigl(L(0)+2\pi iJ\theta\bigr),
 \qquad 0\le u\le1
 \tag{RJC.85}
\]
is holomorphic in the annular parameter
and preserves the integral winding.
The contraction domain is the original
punctured sufficiently small scalar
disc; convexity of its logarithmic
real bound keeps this deformation
inside it. On the compact deformation
interval the original source and line
norm bounds give common strictly
contracting cap tubes.

Use \(L_u\) in the full equations
RJC.33-\/-41, keeping every integer
period, finite character and reference
value. Use RJC.64-\/-68 on every selected
chart. Their inverse arrows have the
same continuous dependence. These
maps form the actual smooth proper
family over the deformation square.
Local lifts of its parameter vector
field, patched by a partition of
unity and chosen to match the given
end collars, integrate for the
whole compact interval. This gives
the relative endpoint isotopy and
preserves the specified source,
line and cap comparisons.

The ratio of its closing to the
chronological \(s\)-th power of
RJC.15 is a physical scalar by
RJC.56-\/-58. Choosing any logarithm
of that scalar, the actual path
of multiplying by \(\exp(u\log c)\)
extends on every point line,
cap and fan and gives its isotopy
to identity. Thus all holomorphic
closings at that winding have
the same smooth return.
Every integer \(s\) is attained
on an annulus about \(w_*\ne0\)
by the explicit scalar parameter
map
\[
 w\longmapsto
 w_*\left(\frac{w}{w_*}\right)^s.
 \tag{RJC.86}
\]
Pulling back the actual shorter
family gives its closing. Negative
\(s\) uses the actual inverse
maps; \(s=0\) is the constant
family. At a fixed fan \(P\) the
complete set of returns is
\[
 [\Phi]=s\,S(P)\pmod{28}.
 \tag{RJC.87}
\]
Combining this with the unit
witnesses proves the images
in RJC.84 for all paths and
all holomorphic closings.

Finally use the complete seam
and disk-transfer maps of the
proved mapping-torus construction.
For the positive seam
\((x,1)\sim(\Phi(x),0)\),
\[
 T_+(\Phi)\cong
 (S^1\times S^6)\mathbin{\#}\Sigma_{sS(P)}.
 \tag{RJC.88}
\]
The inverse seam has
\(\Sigma_{-sS(P)}\).
The physical inverse-clutch
identification is
\(\lambda_{\rm physical}(F)
=-\Sigma_{F^{-1}}=\Sigma_F\);
it is the one transported in
RJC.72-\/-74.
For each of the two spin
structures the corresponding
Eells-\/-Kuiper values are
\[
 \mu(T_+)=\frac{sS(P)}{28},
 \qquad
 \mu(T_-)=-\frac{sS(P)}{28}
 \quad\text{in }\mathbb Q/\mathbb Z.
 \tag{RJC.89}
\]
Every nonzero residue obstructs
an orientation-preserving product
identification in that orientation;
reversing the seam changes its
sign by the explicit inverse.

\par

\medskip\noindent\textbf{RJC.13. Exact checks and remaining scope}

\par\nobreak

The primary executable reconstructs
the original odd-line Miller
pullbacks, inverse-source
coefficients, root characters,
signed norm equations, both
root exponents and the full
cyclic scalar identity.
It calculates the four page
pairs and performs fresh
nine-unknown target-prism solves,
then checks the actual unit-fan
witnesses and the unbounded
modulo-seven identities.
The geometric statements above
use the complete source, line,
root, reference and parameter
maps before their numerical
characteristic evaluations.
The bounded tests support those
proofs; the terminating fan
construction and the exact
residue arguments give their
unbounded conclusions.

This proof has the two actual
additive caps of RJC.2 and
signed equation RJC.3. The
additional ordinary norm is
zero throughout. The prescribed
ordinary references, all finite
characters and every displayed
root order are included.
Further ordinary clutch slots
with nonzero total norm require
their actual redistributed
cochain comparison before they
enter this statement. Other
marked source-moduli loops are
also outside this scalar-path
classification. Part3.4 remains
incomplete. Independent full
geometric audit and export are
still pending for this
continuation.

% End expanded verified cp36 input


\par\bigskip\noindent\textbf{Complete supporting proof 2: JFP.}\par


% Expanded verified cp36 input: notes/part3_e6_iii_joint_endpoints_jfp_proof.tex
\par

\medskip\noindent\textbf{Independent full comparison for the original E6 source with both additive caps}

\par\nobreak

This proof starts from the actual F46 source and Cartier construction,
the K67 shorter parameter family, the K73 native products, and the
completed K114--115 endpoint and scalar-family maps. The root exponents,
cyclic endpoint, characteristic inputs and fresh nine-variable prism
systems were derived and executed before the primary RJC manuscript
was read. The present proof also records the subsequent complete
RJC.1--13 geometric review. Every point below uses the original source
and its specified references. The ordinary extra norm is zero.

\par

\medskip\noindent\textbf{JFP.1. Source, both root maps, and all point divisors}

\par\nobreak

Keep
\[
 y^2=x^3+tx+t^2,\quad G=(0,t),\quad q=t+4/27,
 \quad \mathcal M_a=\mathcal O([aG]-[O]),\quad ab=6.
\tag{JFP.1}
\]
The four domains, with their original positive orders, are
\[
\begin{array}{c|c|c|c}
(a,b)&m=m_\infty&h=m_0&\text{root fibres at infinity and zero}\\\hline
(1,6)&3l, l\ge1&4k, k\ge1&\text{smooth, smooth}\\
(2,3)&3l, l\ge1&4k+1, k\ge0&\text{smooth, III}\\
(3,2)&3l+2, l\ge0&4k, k\ge1&\text{IV, smooth}\\
(6,1)&3l+2, l\ge0&4k+1, k\ge0&\text{IV, III}.
\end{array}
\tag{JFP.2}
\]
The prescribed signed norms obey
\[
 h n_\infty+m n_0=\varepsilon,\qquad \varepsilon=\pm1.
\tag{JFP.3}
\]
In particular \(\gcd(m,h)=1\) and both local norms are primitive.
This is F46\textquotesingle s integral sphere equation. Its actual source/cap
chain maps have degree-two map an integral isomorphism and
degree-one and degree-three determinants \(-\varepsilon\).
Their fundamental-group presentation is cyclic of order one.
The common blowup and its two ruling projections carry these maps
through each prescribed lattice-basis fan, preserving the primitive
normal circle and the physical line. Thus the F46 sphere and its
vanishing \(H^2,H^4\) apply on every fan used here.

At infinity the full substitutions are
\[
\begin{array}{c|c|c}
m&(t,x,y)&\text{equation}\\\hline
3l&(z^{-3l},z^{-2l}X,z^{-3l}Y)&Y^2=X^3+z^lX+1\\
3l+2&(z^{-m},z^{-2(l+1)}X,z^{-3(l+1)}Y)&
Y^2=X^3+z^{l+2}X+z^2.
\end{array}
\tag{JFP.4}
\]
At zero they are
\[
\begin{array}{c|c|c}
h&(t,x,y)&\text{equation}\\\hline
4k&(z^{4k},z^{2k}X,z^{3k}Y)&Y^2=X^3+X+z^{2k}\\
4k+1&(z^h,z^{2k}X,z^{3k}Y)&Y^2=X^3+zX+z^{2k+2}.
\end{array}
\tag{JFP.5}
\]
The deck multiplies \(z\) by its specified root of unity and
\(X,Y\) by the inverse displayed weights. Each original coordinate
is therefore invariant. The reciprocal powers give the inverse
on the original annulus. The reduced equations retain their A1 or
A2 resolution charts.

The full good-cubic point-line graph uses the quotient weights
\((2,1)\). On its common regular refinement the rays
\((1,0),(2,1),(1,1),(1,2),(0,1)\) give the discrepancy
\(\min(v_\eta,2v_z)\). After the actual pullback \(w=z^l\),
the correspondence is
\[
 \beta^*\mathcal O(aG'-O)=\pi^*\mathcal M_a+D_{aG'}-D_O,
 \qquad D_P(v)=\min(v_\eta,2l v_z).
\tag{JFP.6}
\]
Each cone map is its two original monomials in \(\eta,z\), with
the integral inverse on the common torus. The strict point-section
factor is retained in that generator. This carries both the original
point module and all its exceptional factors.

On the good quartic root, the actual action
\((w,X,Y)\mapsto(iw,-X,iY)\) fixes \(O,T=(0,0)\) and exchanges
\((i,0),(-i,0)\). Blow up those four points and take the quotient.
The exceptional images have squares \(-4,-4,-2\), the central
curve has square \(-1\), and the fibre is
\(4C+D_O+D_T+2D_2\). Contracting \(C\), then the resulting
\((-1)\)-curve \(D_2\), gives the original two \((-2)\)-curves
with intersection scheme of length two. The invariant maps (JFP.5)
give their actual tangency. The marked sections avoid these two
contractions. The line correspondence is
\[
 \beta^*\mathcal O(aG'-O)
 =\pi^*\alpha^*\mathcal M_a+E_{aG'}-E_O.
\tag{JFP.7}
\]
For \(w=z^k\), the section graphs resolve the ideals
\((\eta,z^k)\). The complete successive cone maps are
\[
 \eta=u^jv^{j+1},\quad z=uv\quad(0\le j<k),
 \qquad \eta=uv^k,\quad z=v
\tag{JFP.8}
\]
on the last cone. Their determinants are \(-1\). They pull the
ideal to the indicated generator with valuation
\(\min(v_\eta,k v_z)\); the corresponding reciprocal module maps
are the inverse. For odd \(a\), translation exchanges the two
marked exceptional contributions. Its actual Hom coefficient
carries that exchange, as computed in JFP.2 below.

For the reduced III root the original formal parameter is
\(\tau=-x/y\) and \(\tau'=z^k\tau\). For \(a=2,6\),
\(\tau'_{aG'}=z^kA_a(z^h)\), where \(A_a(0)=a\).
The two actual point/O ratios are
\[
 \frac{p-A_a(z^h)}p,\quad 1-vA_a(z^h),
 \qquad \tau'=z^kp,\quad z^k=\tau'v.
\tag{JFP.9}
\]
The intermediate charts are (JFP.8). The A1 chart is
\(X=z\alpha,Y=z\beta\),
\(\beta^2=z\alpha^3+\alpha+z^{2k}\). The slopes
\((Y\mp z^{k+1})/X\) give the two inverse source translations
on their resolved graphs. For the reduced IV root, use the full
A2 charts, beginning with \(X=zX_1,Y=zY_1\), whose two outer
component labels are \(Y_1=\pm1\). Its point graph is
\(\tau'/(\tau'-z^{2l+1}L_a(z))\), with the actual functions
\[
 L_3(z)=-\frac{64+8z^m}{512+96z^m+3z^{2m}},
 \quad L_3(0)=-\frac18,
 \qquad L_6(0)=-\frac14.
\tag{JFP.10}
\]
Here \(L_a=-z^{-2l-1}X_{aG'}/Y_{aG'}\), so its complete
rational function, and its reciprocal module map, stay in the
correspondence. The distinct constants are obtained by the original
\(3G\) and doubled \(6G\) formulae. This was the one source-graph
correction found in the full primary read; the preceding primary
and independent derivation are preserved in their attempt records.

At either reduced root the point \(aG'\) lies in the identity
component. The original pulled-back line and the displayed root
line have zero degrees on all components. The blowup Picard
decomposition descends the original line through every contracted
curve. Their remaining vertical divisor is an integer whole fibre,
because the resolved Kodaira intersection matrix has primitive
kernel its fibre vector. Division by the actual base equation to
that integer power gives the line map; multiplication gives its
inverse. The power and its deck character remain specified.

\par

\medskip\noindent\textbf{JFP.2. Original Miller maps and the odd-line coefficients}

\par\nobreak

Let \(T(p)=p-G\). On the original source retain
\[
\begin{split}
 f_1&=1,\qquad f_2=\frac{y-t-x/2}{x-1/4},\\
 f_3&=\frac{(x/4)(y-t-x/2)}{(y+t+x/2)/4-(1/4+2t)x},\\
 (x_3,y_3)&=(8t+64t^2,3t+96t^2+512t^3),\\
 L_3^{\rm tan}&=(3x_3^2+t)/(2y_3),\quad
 x_6=(L_3^{\rm tan})^2-2x_3,\\
 y_6&=L_3^{\rm tan}(x_3-x_6)-y_3,\quad
 f_6=f_3^2\frac{y-y_3-L_3^{\rm tan}(x-x_3)}{x-x_6}.
\end{split}
\tag{JFP.11}
\]
Their divisors are \(a(G)-(aG)-(a-1)(O)\). The full one-step
negative lift and its positive inverse therefore are
\[
 B_{a,1}(p)=f_2(p)^a\frac{f_a(Tp)}{f_a(p)},\qquad
 U(p,v)=\bigl(p+G,B_{a,1}(p+G)^{-1}v\bigr).
\tag{JFP.12}
\]
The inverse coefficient is evaluated at the inverse source point.
These are Hom-module identities even where a displayed rational
coefficient has a zero or pole; the point/O Cartier generators
carry those factors.

On the two original III components use
\(x=p^{-2},y=-p^{-3}\) and
\(x=t\alpha,y=t\beta\),
\(\alpha+t\alpha^3=\beta^2-1\). The directed source maps are
\(\beta'=p-1,p'=\beta-1\). Substitution into (JFP.11) gives
\(f_2,f_3,f_6\) component orders \((0,1),(0,1),(0,3)\).
Consequently the one-step Hom orders for \(a=1,2,3,6\) are
\((0,1),(1,1),(1,2),(3,3)\). The chronological \(b\) factors
have order three on both components. At every other finite fibre
the two translated point factors have the same generic order;
\(f_2\) has no whole-fibre zero. The actual original translation
Hom has class \(F\), so its infinity coefficient is \(-2\).
The specified suspension section is exactly
\[
 B_0=C_a\frac q{t^3}U^{-b},\qquad
 \rho=\frac{B_0^{-1}}{U^b}=\frac{t^3}{C_aq}.
\tag{JFP.13}
\]
The ratio of the two sections has zero divisor on the compact
regular source and is an invertible constant. Its evaluation at
the prescribed reference fixes the actual \(C_a\in\mathbb C^*\).

The new odd-line good-root coefficients can also be computed
without relying on a divisor alone. At the finite good root,
\(w=z^k\), direct substitution gives
\[
 f_{2,\rm old}=w^2\frac{wY-w^2-X/2}{w^2X-1/4},\quad
 f_{3,\rm old}=w^2\frac{X(2wY-2w^2-X)}
 {2wY+2w^2-X-16w^4X}.
\tag{JFP.14}
\]
The central source map is \((X,Y)\mapsto(1/X,-Y/X^2)\).
The full root functions specialize to \(f_1=1,f_2=2X,
f_3=X,f_6=6X^3\). Hence the one-step negative coefficients
are \(2^aX\) for odd \(a\), and \(2^a\) for even \(a\).
The positive inverse coefficients are \(X/2^a\) and \(1/2^a\),
respectively. For the two odd rows their chronological \(b\)-fold
positive norm is \(1/64\). The exchanged \(E_{aG'}-E_O\)
generators in (JFP.7) are carried by these exact coefficients.

At the good infinity root, \(w=z^l\),
\[
 f_{2,\rm old}=w^{-1}\frac{Y-1-wX/2}{X-w^2/4}.
\tag{JFP.15}
\]
On \(Y^2=X^3+1\), the mutually inverse source maps are
\[
 (X,Y)\mapsto\left(-\frac{2X}{Y+1},\frac{Y-3}{Y+1}\right),
 \qquad
 (X,Y)\mapsto\left(\frac{2X}{Y-1},\frac{-Y-3}{Y-1}\right).
\tag{JFP.16}
\]
The positive \(M_1\) coefficient is \(X/2\), whose three
chronological factors give \(-1/2\) and whose six give \(1/4\).
The positive \(M_2\) coefficient is \((Y+1)/(2X)\), whose
three factors give \(1/4\). The source third iterate is the
identity. These full source and line calculations carry the
good-root point corrections (JFP.6); their inverse maps retain
the full off-central functions (JFP.14--15).

\par

\medskip\noindent\textbf{JFP.3. Both root lifts, division lifts, and finite characters}

\par\nobreak

The extending positive root lift \(L_i\) is chosen with its
specified rigidification and \(L_i^b=B_0^{-1}\). The entire
ordinary-to-root comparison is
\[
 U=z^{e_i}u_i L_i,\qquad
 \begin{array}{c|c|c}
 i&e_i&u_i^b\\\hline
 0&-3h/b&C_aq\\
 \infty&2m/b&C_a(1+4z^m/27).
 \end{array}
\tag{JFP.17}
\]
To construct it, apply the exact point modules of JFP.1 to the
root chord lift. Its horizontal divisor is the prescribed Hom
divisor. Its remaining vertical orders have zero component degrees;
the primitive fibre kernel makes them one integer base power.
Remove that displayed power. The section and its inverse then
have zero orders at every codimension-one point. Normal local
rings are intersections of their height-one valuation rings, so
both extend through component and section intersections. The
nonconstant logarithmic deck commutator is removed by finite
averaging. For a norm-zero logarithm \(g\), the explicit average
\(-m_i^{-1}\sum_{j=0}^{m_i-1}jD_i^jg\) gives the required
coefficient correction. The invariant analytic \(b\)-th root
of the norm ratio, with the prescribed reference value, fixes
\(L_i^b=B_0^{-1}\). Its reciprocal gives the inverse correction.
The retained constant character is obtained from (JFP.17).

The coefficient convention matters. Conjugating multiplication
by \(z^e\) by \(D:z\mapsto\zeta z\) evaluates that function
at \(z/\zeta\), giving \(\zeta^{-e}z^e\). Original \(U\)
is invariant, so
\[
 D_iL_iD_i^{-1}=\zeta_{m_i}^{e_i}L_i,\qquad
 \chi_\infty=\zeta_b^2,\quad\chi_0=\zeta_b^{-3}.
\tag{JFP.18}
\]
Thus the four character pairs are
\((\zeta_3,-1),(\zeta_3^2,1),(1,-1),(1,1)\).
The norm equation (JFP.3) gives
\(\chi_i^{-\varepsilon}=\zeta_b^{Mn_i/m_i}\),
with \(M=mh\). Every prescribed common finite line character
cancels in this conjugation, while it remains in the cap operator.

For completeness the cap division lift has character one.
The marked division section \(R_i\) satisfies \(m_iR_i=bG'\).
At a good root it is the unique analytic inverse branch with
\(R_i(0)=O\). At a reduced root multiplication by \(m_i\)
is invertible on the additive identity group and on its finite
component group, so it gives the unique marked analytic section.
The deck preserves it. For \(P_i=aG'\), the full chord ratio
\[
 g_{P_i,R_i}(p)=
 \frac{\ell_{P_i,R_i}(p)}{x(p)-x(P_i+R_i)}
\tag{JFP.19}
\]
has divisor \((P_i)+(R_i)-(P_i+R_i)-(O)\) and is the Hom
coefficient for \(p\mapsto p-R_i\). Its whole-fibre power is
removed through the actual generators (JFP.6--10). At a good root
the central source translation is the identity and the central
line map is a constant. At a reduced root of order greater than one, the actual \(aG'\)
section limits to \(O\), so its point/O line is trivial with
the displayed identification; the central translation lift is
again a constant in that identification. At the finite order-one
root, the full graph is \(\tau'=A_a(t)\) with
\(A_a(0)=a\), as in (JFP.9), and the deck is the identity.
Its commutator is therefore one directly on that actual point line.
The commutator\textquotesingle s
constant character is therefore one. Finite averaging and the
invariant norm root give the complete commuting operator
\(\widetilde F_i^{m_i}=B_0\).

With its retained character \(d_i\), the actual cap generator is
\[
 A_{i,\lambda}=R_{\lambda,i}^{n_i}\zeta_{m_i}^{d_i}
 \widetilde F_i^{n_i}D_i,\qquad
 R_{\lambda,i}^{m_i}=\lambda,\qquad
 A_{i,\lambda}^{m_i}=B_\lambda^{n_i}.
\tag{JFP.20}
\]
All factors commute with the suspension. For integers \(u_i,v_i\)
with \(n_iu_i+m_iv_i=1\), the primitive generator
\(A_i^{u_i}B^{v_i}\) is the strict contraction with its actual
scalar root. This proves the cap shell and its physical-circle
freeness on the stated domain for either sign, with the same
Cartier and finite-character factors.

\par

\medskip\noindent\textbf{JFP.4. Complete annular cochains and their integer restoration}

\par\nobreak

Set \(q_0=t,q_\infty=1/t\), \(m_0=h,m_\infty=m\), and retain
\[
 s_{i\alpha}=\log_\alpha q_i/(2\pi i),\quad
 z_{i\alpha}=e^{2\pi i s_{i\alpha}/m_i},\quad
 c_{\lambda,i}=n_i\log R_{\lambda,i}/(2\pi i)+d_i/m_i,
 \quad f_{i\alpha}=\operatorname{Log}_\alpha\widetilde F_i.
\tag{JFP.21}
\]
The original-coordinate and point-module root map is \(I_z\).
The complete gluing is
\[
 \Phi_{i\alpha}=\operatorname{Exp}
 (n_i s_{i\alpha}f_{i\alpha}+s_{i\alpha}c_{\lambda,i}\delta
       +h_{i\alpha})I_{z_{i\alpha}}.
\tag{JFP.22}
\]
Its source translation is \(-n_i s_{i\alpha}R_i\). On an
overlap labelled \(k\),
\(f_{i\beta}=D_{i*}^kf_{i\alpha}+\Pi_{i\alpha\beta}\), and
the retained reference cochain satisfies
\[
 h_{i\beta}-D_{i*}^kh_{i\alpha}
 =-n_i s_{i\alpha}\Pi_{i\alpha\beta}
   +j_{i\alpha\beta}\operatorname{Log}_\beta B_\lambda
   +\Lambda_{i\alpha\beta}.
\tag{JFP.23}
\]
The periods \(\Pi,\Lambda\), suspension integers \(j\), and
reference value are unchanged. Starting from its prescribed smooth
extension \(h^0\), subtract the residual with
\(k_\alpha=-\sum_\beta\chi_\beta E_{\alpha\beta}\), then
subtract the common Cauchy integral
\(\pi^{-1}\int\bar\partial k(w)/(q-w)\,dA(w)\).
The residual has compact subannular support. Fixing the common
holomorphic constant at the original reference gives exactly that
reference value. Subtracting the target-sheet logarithm proves
\[
 \Phi_{i\beta}=A_{i,\lambda}^{k}B_\lambda^{j_{i\alpha\beta}}
 \Phi_{i\alpha},\qquad \Phi_{i\alpha}B_\lambda=B_\lambda\Phi_{i\alpha}.
\tag{JFP.24}
\]
The inverse uses the negative Lie translation at the inverse
source point, followed by the inverse root map and reciprocal
Cartier factors.

For \(\lambda_\theta=e^{2\pi i\theta}\lambda_*\), retain the
actual integers \(J_{i\beta}-J_{i\alpha}=j_{i\alpha\beta}\).
The full restoration formula is
\[
 \Phi_{\theta,i\alpha}
 =e^{2\pi i n_i\theta s_{i\alpha}/m_i}
 B_\theta^{J_{i\alpha}}
 (B_*^{-J_{i\alpha}}\Phi_{*,i\alpha}).
\tag{JFP.25}
\]
At integer winding \(J\), the last two factors have scalar ratio
\(e^{2\pi iJ J_{i\alpha}}=1\), and the first gives
\(\Phi_{\theta+J,i\alpha}=z_{i\alpha}^{n_iJ}\Phi_{\theta,i\alpha}\).
This establishes both root exponents used below through the full
reference, with every integral period still present.

\par

\medskip\noindent\textbf{JFP.5. Primitive endpoint and its exact cyclic phase}

\par\nobreak

Put \(J_0=M/b\) and retain the actual one-gamma forward closing
\(h_0=c_0q^\varepsilon t^{-mn_0}\). The single-valued integer
\[
 d_0=(3\varepsilon-mn_0)/b=(2\varepsilon+hn_\infty)/b
\tag{JFP.26}
\]
has respectively the integral expressions
\(2kn_\infty+ln_0\),
\((4k+1)n_\infty+2ln_0\),
\(6kn_\infty+(3l+2)n_0\), and
\(3(4k+1)n_\infty+2(3l+2)n_0\).
Choose the specified root \(\gamma_0^b=c_0C_a^{-\varepsilon}\).
The full positive endpoint over \(\lambda=r^M\) is
\[
 E_r=\gamma_0r^{-\varepsilon J_0}t^{d_0}U^\varepsilon,
 \qquad r\longmapsto\zeta_b r.
\tag{JFP.27}
\]
At zero the full base exponent is
\(hd_0-3h\varepsilon/b=-J_0n_0\); at infinity it is
\(-md_0+2m\varepsilon/b=-J_0n_\infty\).
The reference factors of (JFP.25) cancel exactly these powers.
The remaining cap coefficient is
\(\gamma_0r^{-\varepsilon J_0}u_i^\varepsilon L_i^\varepsilon\),
and (JFP.18) gives its exact finite-generator square. The inverse
is the reciprocal coefficient at the inverse source point.

The chronological product keeps the scalar phase
\(\exp[-\pi i\varepsilon J_0(b-1)]\). In all four rows
\(J_0(b-1)\) is even. With the full constants in (JFP.13),
\[
 E_{\zeta_b^{b-1}r}\cdots E_{\zeta_b r}E_r
 =h_0 B_\lambda^{-\varepsilon}.
\tag{JFP.28}
\]
The K67 bounded origin uses
\(c=\exp[\pi i\varepsilon J_0(b-1)/b]\). The exact family
intertwiner is physical multiplication by
\[
 F_r=r^N,\qquad N=\varepsilon J_0(b-1)/2\in\mathbb Z,
 \qquad F_{\zeta_b r}E_r=(cE_r)F_r.
\tag{JFP.29}
\]
The four exponents are \(5\varepsilon kl\),
\(\varepsilon l(4k+1)\), \(\varepsilon k(3l+2)\), and zero.
The phase has this explicit map and reciprocal inverse; it remains
in the complete parameter family.

\par

\medskip\noindent\textbf{JFP.6. Intrinsic elliptic quotient and base rigidity}

\par\nobreak

The F46 topology in JFP.1 gives \(H^2=H^4=0\). A horizontal
polar divisor of a meromorphic function restricts to a nonempty
effective divisor on a compact regular two-torus. Pairing with a
flat Kähler form is positive, whereas its restricted cohomology
class is zero. Thus all poles are vertical. The function is
constant on each regular compact torus. At a special component,
the full transverse map \(u\mapsto t_0+u^e\) and invariance
force Laurent powers divisible by \(e\). Consequently pullback
gives the precise field isomorphism \(\mathbb C(t)\cong\mathbb C(X)\),
whose inverse is fibre restriction followed by this Laurent descent.

The full generic period module is
\[
\begin{gathered}
 \gamma=(-bg,\ell),\quad e_1=(1,0),\quad e_2=(\tau,ag),
 \quad\delta=(0,1),\\
 M_\alpha=
 \begin{pmatrix}1&0&0\\bv_\alpha&A_\alpha&0\\
 c_\alpha&a v_\alpha^tJ A_\alpha&1\end{pmatrix},
 \quad J=\begin{pmatrix}0&-1\\1&0\end{pmatrix}.
\end{gathered}
\tag{JFP.30}
\]
The full theta and Miller constant is in \(\ell\). Multiplication
gives the retained central cochain
\(c_{\alpha\beta}=c_\alpha+c_\beta+abv_\alpha^tJ A_\alpha v_\beta\).
The Abel cocycle is nonzero over \(\mathbb Q\). A rational
coboundary, after clearing denominators and component orders,
would supply a global Lie logarithm of a nonzero multiple of
\(G\). At the minimal quartic root divide the logarithm of
\(2G'\) by its actual factor \(w\); at the cubic root divide
the logarithm of \(3G'\) by \(w^2\). The latter has leading
coefficient \(-1/8\). Both finite elliptic actions have no
invariant period; at the I1 the invariant annular period extends
through \(xy=q\). The original Lie transition is
\(u_\infty(q_\infty)=q_\infty^{-1}u_0(q_\infty^{-1})\).
Its entire global coefficient is zero. But original \(3G\)
has infinity value \((64,512)\), with nonzero additive parameter
\(-1/8\), and every nonzero multiple remains nonzero. The
contradiction proves the Abel assertion. The higher maps JFP.1
are the actual pullbacks of these minimal diagrams.

The vector and dual equations of (JFP.30) now give exactly
\(\Lambda^{\pi_1}=\mathbb Z\delta\) and
\((\Lambda^*)^{\pi_1}=\mathbb Z\psi\), where \(\psi(\gamma)=1\)
and it vanishes on the other three generators. Both are primitive.
The kernel cover is the actual nonzero point line with its Cartier
modules. The derivative of a lifted endpoint is periodic for the
compact torus lattice and hence constant. The full affine map and
inverse are
\[
\begin{gathered}
 (z,w)\mapsto(Az+z_0,Cz+e w+w_0),\\
 (z',w')\mapsto((z'-z_0)/A,\,
 e[w'-w_0-C(z'-z_0)/A]),\\
 A\ne0,\qquad e=\pm1.
\end{gathered}
\tag{JFP.31}
\]
On the kernel lattice its matrix is
\(K=\left(\begin{smallmatrix}p&q_1&0\\r&s&0\\u&v&e\end{smallmatrix}\right)\),
with \(ps-q_1r=1\). Its inverse has last row
\((e(vr-us),e(uq_1-vp),e)\) and upper block
\(\left(\begin{smallmatrix}s&-q_1\\-r&p\end{smallmatrix}\right)\).
The positive determinant is the complex area identity. The first
coordinates therefore descend to inverse elliptic source maps.
The original cubic gives \(j_0(t)=6912/(27t+4)\). Equality under
that actual elliptic isomorphism forces the base Möbius map to be
the identity.

The order-three infinity block and the selected nonzero unipotent
have simultaneous rational commutant \(\mathbb QI\): a nonscalar
element of \(\mathbb Q(\zeta_3)\) generates a field, which has no
nonzero nilpotent. The off-diagonal equations and nonzero Abel
class leave the full lattice commutant \(eI+c\delta\psi\).
Complex linearity on the independent vectors \(e_1,\delta\)
sets the derivative to \(eI\). The remaining suspension equation
makes the scalar logarithm difference an integer. The two ends
therefore have equal original scalar values.

\par

\medskip\noindent\textbf{JFP.7. Global section extension and all holomorphic endpoints}

\par\nobreak

Pull an endpoint and its inverse to the normalized root fibre
product \(z_i^{m_i}=q_i\). A nonzero line vector over the
O-section kills the root meridian; the remaining vanishing cycles
lie in \(\ker\psi\). The actual graphs in JFP.1 supply this
lift on the whole root source. Its commutator with the cap
generator is \(B^s\). Iterating the relation \(m_i\) times gives
\(B^{m_is}=1\); strict contraction implies \(s=0\).
The source displacement is an equivariant holomorphic root
section. Descend by (JFP.4--5) and resolve its graph in the proper
minimal original source.

At the selected source a coefficient pole of order \(-K\)
is cancelled by the actual \(B^K\) reindex. It shifts the source
displacement by \(-bKG\). The remaining two character factors
are holomorphic units; their complete vertex arrows are written
in JFP.8. Their inverses extend over both vertices and the open
axes. After this selected reindex the joining torsor has
presentation
\[
 \langle x,y\mid x^m=y^h=xy=1\rangle.
\tag{JFP.32}
\]
It is cyclic of order \(\gcd(m,h)=1\). Thus the integer joins
are globally trivial, with the one constant join fixed at the
prescribed reference. This proves extension to a global source
section.

The source is rational by the inverse birational chart
\[
 u=(y-t-x/2)/x,\quad v=(y+t+x/2)/x,\quad
 x=uv+1/4,\quad t=x(v-u-1)/2,\quad y=x(u+v)/2.
\tag{JFP.33}
\]
Its ten classes \((O,F,C_0,\ldots,C_5,A_1,G)\) have the
negative E6 block with edges \((0,1),\allowbreak(1,2),\allowbreak(2,3),\allowbreak(3,4),\allowbreak(2,5)\),
\(A_1^2=-2\), \(O^2=G^2=-1\), \(OF=GF=1\), and
\(GC_0=GA_1=1\). All other unspecified intersections are zero.
The complete determinant is \(-1\). These classes span the
rank-ten unimodular lattice; quotienting the first nine is free
cyclic generated by \(G\). Hence the section group is exactly
\(\mathbb ZG\), with height \(2-4/3-1/2=1/6\).

A positive endpoint is therefore the full map \(H(t)U^j\).
For winding \(J\), the two root comparisons and selected reindex
give the actual divisor orders
\[
 \operatorname{ord}_0H=3j/b-n_0J/h,\quad
 \operatorname{ord}_\infty H=-2j/b-n_\infty J/m,\quad
 \operatorname{ord}_{t_c}H=-K.
\tag{JFP.34}
\]
Elsewhere the Hom and its inverse have no zero. Summing the
principal divisor gives
\(j=bK+b\varepsilon J/M\), so \(J\in J_0\mathbb Z\).
The endpoint (JFP.27), its chronological powers and its inverse
attain every such integer.

For one object \(J=0\), so \(j=bK\),
\(H=c(t^3/q)^K\), and the whole map is
\(cC_a^K\lambda^KB_\lambda^{-K}\). Its quotient is the
displayed physical scalar. Faithfulness follows because a
nonzero suspension deck translates by a nonzero multiple of
the non-torsion \(bG\). A negative automorphism has sixteen
transverse fixed points on each regular two-torus. Its compact
analytic equalizer has a horizontal curve, whose positive base
degree contradicts pairing with \(\pi^*H^2(\mathbb P^1)=0\).
Thus \(\operatorname{Aut}_{\rm hol}(X)=\mathbb C^*_{\rm physical}\),
with the inverse identification just given. Two positive
endpoints at the same winding differ by precisely that scalar.

The full selected negative arrow is
\((Z,V,q)\mapsto(Z^{-1},V^{-1},q)\). Opposite triangles are
translated by \(-p-u-v\); the vertex exchange retains all
\(\kappa g,\alpha\eta\) factors and their reciprocals, and
extends on every fan. The source/root lifts of a negative
endpoint exist by the same point-vector and inverse-map
argument. Compose at the selected charts with this inversion
to remove its selected integer join. Both lifts satisfy
\(\widetilde F B=B^{-1}\widetilde F\). On the infinity annulus
the full reference word is \(A_\infty B^\nu\). One root circle
has \(B^N\), \(N=n_\infty+m\nu\). After fixing one constant
deck on a cut annulus, continuation yields
\[
 B^N\widetilde F B^{-N}=B^{2N}\widetilde F.
\tag{JFP.35}
\]
Single-valuedness and contraction force \(N=0\). But
\(\gcd(m,n_\infty)=1\) and every allowed \(m\ge2\), so
\(n_\infty+m\nu\ne0\). This proves the negative exclusion
through the actual root word, including the quadratic reduced
case, and completes the exact winding group classification.

\par

\medskip\noindent\textbf{JFP.8. Complete parameter comparison and the finite cover}

\par\nobreak

Let \(v=\log r/(2\pi i)\), and keep the bounded source period
\(\Gamma\), ordinary logarithm \(H_0=\varepsilon s+U_0(q)\),
and outward suspension \(\gamma=(\Gamma,C-s-Mv)\).
The constant \(c_0\) and whole zero-cap term
\(-mn_0\log t/(2\pi i)\) are in \(U_0\). Put
\[
 \tau(v)=\left(\frac{\varepsilon\Gamma}{b},
 \frac{U_0+\varepsilon C-\varepsilon Mv}{b}\right),
 \quad
 \sigma(v)=\left(\varepsilon\Gamma v,
 v(U_0+\varepsilon C)-\frac{\varepsilon M}{2}(v^2-v)\right).
\tag{JFP.36}
\]
Then \(\sigma(v+1/b)-\sigma(v)=\tau(v)+
\varepsilon M(b-1)\delta/(2b^2)\). The coordinate
\(\xi=z-\sigma(v)\) closes for \(cE_r\). The entire
intermediate products
\(C_s(v)=\exp(\sigma(v)-\sigma(v+s))\) and
\(P_j(v)=\exp(\sigma(v)-\sigma(v+j/b))\) satisfy
\(C_s(v+j/b)P_j(v)=C_{j/b+s}(v)\).
Their negative exponents are the inverse. At \(j=b\) this
is the actual inverse one-gamma closing, with every derivative
matched on the end collars.

For a basis \(P\), retain all four full coefficient units in
the conjugated period matrix. If its two multiplicative factors
are \(A,D\), the actual vertex and Cox arrows are
\[
\begin{split}
 (a,b_1,c_1)&\mapsto((AD)^{-1}a,Ab_1,Dc_1),\\
 (u,v_1,w_1)&\mapsto(ADu,D^{-1}v_1,A^{-1}w_1),\\
 (a_0,a_1,c_0,c_1)&\mapsto(a_0,Aa_1,(AD)^{-1}c_0,Dc_1).
\end{split}
\tag{JFP.37}
\]
They preserve \(q\), the Cox relation and both ruling charts;
reciprocal factors give the inverse on the exceptional curve
as well. Substitution into the three complete transition rows
\((\kappa\alpha b_1c_1,\alpha^{-1}ab_1,\kappa^{-1}b_1^{-1})\),
\((g\eta b_1c_1,\eta^{-1}c_1^{-1},g^{-1}ac_1)\), and
\((a^{-1},ab_1,ac_1)\) gives the same factor on both sides.
For \(F_r\) these factors are
\(r^{N(P^{-1})_{12}},r^{N(P^{-1})_{22}}\). Every full
\(q\)-dependent unit and higher normal jet is retained by the
normal rescaling; its cutoff flow fixes the protected collars.
The ordinary and cap maps are exactly JFP.3--5. Thus (JFP.29)
is an isomorphism of the entire parameter families and every
intermediate seam.

The free parameter action has quotient \(w=r^b\),
\(\lambda=w^{J_0}\). Its complete \(r\)-family is the literal
degree-\(b\) pullback. The physical line, cap identification and
stable tangent pull back by the maps just written. Integration
of their integral characteristic classes therefore multiplies
by \(b\). This gives the primitive integral comparison before
any reduction modulo 28.

\par

\medskip\noindent\textbf{JFP.9. Page, physical pair and native product}

\par\nobreak

The actual original canonical page over ordinary binding
\(\tau\) is
\[
 \eta^h-\tau\xi^m=1,\qquad
 t=\tau+\xi^{-m}=\eta^h\xi^{-m}.
\tag{JFP.38}
\]
The reciprocal roots are \(\xi^m=(t-\tau)^{-1}\),
\(\eta^h=t/(t-\tau)\), with the original cap labels.
The degree is \(M\), with branch indices \(h,m,M\) at
\(0,\infty,\tau\); its genus is \((m-1)(h-1)/2\).
Pull back the original source by these maps and apply the
actual root resolutions in JFP.1. The relative Picard
evaluation compares its pulled-back point line with the
translated degree-zero line on each fibre. Its component
degrees vanish by (JFP.6--10), so the remaining evaluation
divisor is an integral sum of whole fibres. Division by their
base equations gives an invertible map at every singular
component. Rigidification at O fixes its scalar. The residual
base line changes the class by \(zF\), with \(NF=F^2=0\).
Its boundary value is the specified one and extends over the
ordinary binding disc; the full relative gauge has zero square
ambiguity. The tangent comparison uses the same source map.

The original section height is \(1/6\). The I1 has \(M\)
preimages, infinity has \(h\) root fibres, and zero has \(m\).
Thus
\[
 H=Ma^2/6,\quad N^2=-H,\quad
 \chi=(M+h e_\infty+m e_0)/12,
 \quad(e_\infty,e_0)=(0,0),(0,3),(4,0),(4,3).
\tag{JFP.39}
\]
The full signed binding graph is
\(q_b=(t-\tau)(t/\tau)^{-\varepsilon mn_0}\), whose
derivative at \(\tau\) is one. Retain its constant phase
and all unit derivatives. The normal angles obey
\(\alpha_{\rm normal}=\phi-M\beta\), \(\beta=\varepsilon\gamma\),
and the source graph has \(\beta=\theta+g(p)\).
The K30 triangular source/tangent shear subtracts the full
derivatives of \(g\) and that binding unit. Equivalently the
whole origin difference is
\(\varepsilon\theta\gamma-\varepsilon M\theta(\theta-1)\delta/2\).
Its positive parameter orientation is retained for both signs.
The actual K68 dual-line and inverse-clutch Hopf map consequently
gives the primitive pair and its held-quaternion native product
\[
 (k,\ell)=(3H/b,24\chi/b),\qquad
 (\mu,\nu)=(6\chi/b,3H/b-6\chi/b).
\tag{JFP.40}
\]
The cover in JFP.8 proves this as an equality of integral
characteristic pairs. The native left/right relation is
\(k=\mu+\nu,\ell=4\mu\) under that actual product map.
The four primitive indices are
\(0,kl,k(4l+2),35kl+22k+8l+5\). In full pairs they are
\((kl,4kl)\), \((2l(4k+1),8l(k+1))\),
\((9k(3l+2),12k(l+2))\), and
\((18(3l+2)(4k+1),24(kl+2k+l+1))\).

\par

\medskip\noindent\textbf{JFP.10. Actual fan maps and independent full target systems}

\par\nobreak

Use the literal K71 identical-fan maps (column swap and the
six right triangle rotations) on the absolute semigroups.
They are inverse holomorphic maps, identity on the original
source/Cartier/cap overlaps. Euclidean subtraction of the
positive pair terminates at one vertical \(T_j\). Reversing
uses the actual two flops
\(Q_0=\left(\begin{smallmatrix}1&-1\\0&1\end{smallmatrix}\right)\),
\(Q_1=\left(\begin{smallmatrix}1&0\\-1&1\end{smallmatrix}\right)\).
Their dense characters are \((ZV,V)\), \((Z,ZV)\), their
first triples \((a/c,bc,c)\), \((a/b,b,bc)\), and their
two ruling charts
\((n_0/n_1,zn_1,n_1)\), \((n_1/n_0,zn_0,n_0)\).
The common rank-one matrix and \(q=zn_0n_1\) are retained.
Every full coefficient is transformed by
\(A'_{ik}=\prod_{hl}A_{hl}^{(Q^{-1})_{ih}Q_{lk}}\), with
inverse obtained from \(Q^{-1}\). The source point line and
two original caps therefore carry through each finite word.

For the scalar family of primitive winding \(J=J_0\), the
complete ordinary product has determinant winding \(2J\).
Its native quaternion factors are \(H_J QH_{L_i-J}\), with
the periodic coefficient logarithms and all normal jets retained.
At the second vertex the full ordinary height is
\(qA_{11}(q)A_{22}(q)\); its inverse first recovers the original
\(q\) and divides at that recovered value. The cutoff unit
product similarly has height \(q\exp[\rho(k_{i0}+k_{i1})(q)]\),
whose uniform inverse follows from its unit derivative on a
smaller protected disc. The held Hopf map is
\((UH_JU^{-1})q_HH_{L_i-J}\), of left degree \(2J\) and
right degree zero. Its inverse cylinder is \(L^{-2J}\).
The full circle and derivative prisms therefore have exactly
the K73 data below, with their base and normal Spin endpoints
paired. This is the ordinary-to-native identification carried
by (JFP.37), including the full source-dependent coefficients.

Set \(A=S_0=H/(2b)-\chi/b\), \(\mu=6\chi/b\).
At \(T_j\), \(\nu_j=6A-3Jj(j-1)\),
\[
\begin{split}
 V&=\begin{pmatrix}
 2J-2\mu-2\nu_j&2J(1-2j)&2J(2j-1)\\
 2J(1-2j)&2J&0\\2J(2j-1)&0&2J
 \end{pmatrix},\\
 D&=(2J-2\mu+2\nu_j,2J(2j-1),2J(1-2j)),
 \qquad S=\nu_j/6.
\end{split}
\tag{JFP.41}
\]
At a directed edge select \(a_*=V_{ii},b_*=D_i\) and
\(P'=PQ_i\). Put
\[
\begin{split}
 B&=2J(P')^{-1}e_2e_2^tP',\quad
 L=(B_{22}-B_{21},B_{11}-B_{12},B_{22}),\\
 (x,y)^t&=(P')^{-1}e_2,\quad
 w=(x,y,-x-y),\quad d=(-x-2y,-2x-y,x-y),\\
 I_r&=Jw_r+d_r(L_r-J),\quad W_r=Jw_r-d_r(L_r-J).
\end{split}
\tag{JFP.42}
\]
The actual common-link map conjugates the first quaternion by
the held second quaternion. The selected native reset gives
\(V'_{22}=a_*\), \(D'_2=2a_*-b_*\),
\(S'=S+(b_*-a_*)/4\). Together with all equations
\[
 V'(w_ie_i-w_0e_0)=I_ie_i-I_0e_0,\quad
 w_iD'_i-w_0D'_0=W_i-W_0,\qquad i=1,2,
\tag{JFP.43}
\]
these determine all six symmetric entries and all three
derivatives. The independent executable solves these complete
systems by Fraction row elimination, without importing the
primary or another fan program. Each target has rank nine,
and every redundant original equation is checked. The resulting
entries agree with
\(V'_{rs}=(I_r/w_r)\delta_{rs}+C/(w_rw_s)\),
\(D'_r=(W_r+E)/w_r\), where
\(C=w_2^2a_*-w_2I_2\), \(E=w_2(2a_*-b_*)-W_2\).
All target weights are nonzero in the finite canonical word.
The inverse reset gives the negative of its forward right
increment. Thus these are the actual all-fan returns rather
than only solutions of a numerical recursion.

\par

\medskip\noindent\textbf{JFP.11. Exact unbounded images and all scalar paths}

\par\nobreak

Linearity of the complete target systems in \((J,\mu,A)\)
gives the divisibility upper bounds from actual admissible
unscaled profiles. For \(a=1\), all inputs are \(kl\) times
the \((k,l)=(1,1)\) source. The two full selected pairs along
\(Q_0Q_0\) are \((2,2)\), then \((-4,8)\), and the
return is \(3kl\). Since 3 is a unit modulo 28, the image is
\(kl\mathbb Z/28\).

For \(a=2\), all inputs are \(l\) times the actual \((k,1)\)
profile. Its vertical coefficients are
\(S_j/l=k-(4k+1)j(j-1)/2\). For odd \(k\), choose
\(j=0\) when \(k\ne0\pmod7\), and \(j=4\) otherwise;
the coefficients are \(k\), \(-23k-6\). For even \(k\),
choose \(j=2\) when \(k\ne2\pmod7\), and \(j=3\)
otherwise; the coefficients are \(-3k-1\), \(-11k-3\).
Each is odd and nonzero modulo 7. The image is
\(l\mathbb Z/28\).

For \(a=3\), admissibility forces \(l\) odd. All inputs
are \(k\) times the actual \((1,l)\) profile. The independent
three full target solves along \(Q_0^3\), applied to
\(\mu/k=3(l+2)\), \(A/k=4l+2\), give
\[
 S(Q_0^3)/k=15\cdot3(l+2)+127(4l+2)
 =553l+344=7(79l+49)+1.
\tag{JFP.44}
\]
It is odd and nonzero modulo 7. The image is
\(k\mathbb Z/28\).

For \(a=6\), write \(A=35kl+22k+8l+5\) and
\(J=(3l+2)(4k+1)\). When \(kl\) is even, the two odd
integers \(A,A-6J\) have a common zero modulo 7 only at
\((l,k)=(4,5)\pmod7\): substitution yields
\(2(l-4)^2=0\). When \(kl\) is odd, use the odd integers
\(A-J,A-3J\); their common zero again forces that pair,
by \(3l^2+11l+6=3(l-4)^2\pmod7\).
It makes both \(m,h\) divisible by 7, contradicting their
coprimality. Hence one displayed vertical fan has a return
unit modulo 28, giving the whole group. These finite residue
identities are polynomial identities in the unrestricted
orders, rather than bounded tests of the orders.

The complete table is
\[
\begin{array}{c|c|c}
(a,b)&\text{all holomorphic winding values}&
\text{all-fan, all-loop return image}\\\hline
(1,6)&2kl\mathbb Z&kl\mathbb Z/28\\
(2,3)&l(4k+1)\mathbb Z&l\mathbb Z/28\\
(3,2)&2k(3l+2)\mathbb Z&k\mathbb Z/28\\
(6,1)&(3l+2)(4k+1)\mathbb Z&\mathbb Z/28.
\end{array}
\tag{JFP.45}
\]
For a full holomorphic scalar path, its logarithm on the
universal strip obeys \(L(\theta+1)-L(\theta)=2\pi iJ\).
Interpolate to \(L(0)+2\pi iJ\theta\). Convexity of the
logarithmic modulus bound keeps it in the original contraction
disc. Substitute this full logarithm into (JFP.20--25), retaining
the original finite characters, integer periods and references,
and use (JFP.36--37) on every selected chart. Common strict
contraction tubes exist on the compact interpolation interval.
Local parameter lifts, patched relative to the specified end
collars, integrate to a proper relative isotopy. The endpoint
ratio is a physical scalar by JFP.7; multiplication by
\(\exp(u\log c)\) gives its extending isotopy. Winding
\(sJ_0\) is attained by the literal annular pullback
\(w\mapsto w_*(w/w_*)^s\), with inverse maps for negative
\(s\). At a fixed fan the complete return set is therefore
\(sS(P)\bmod28\), and (JFP.45) follows for all paths and
all holomorphic closings.

\par

\medskip\noindent\textbf{JFP.12. Mapping-torus orientations and audit scope}

\par\nobreak

The full disk-transfer and seam maps retained in K114--115 use
the physical inverse-clutch convention
\(\lambda_{\rm physical}(F)=-\Sigma_{F^{-1}}=\Sigma_F\).
For \((x,1)\sim(\Phi(x),0)\), the oriented mapping torus
is \((S^1\times S^6)\#\Sigma_{sS(P)}\). Reversing the
seam uses the actual inverse endpoint and gives
\(\Sigma_{-sS(P)}\). For each of the two spin structures
the corresponding Eells--Kuiper values are
\(sS(P)/28\) and \(-sS(P)/28\) in \(\mathbb Q/\mathbb Z\).
These signs are those of the binding graph and Hopf map in
JFP.9; no change of the signed sphere equation changes the
chosen positive scalar-parameter orientation.

The independent certificate includes the four symbolic root
and page rows, all signed norm and local character identities,
the whole finite good-root Miller functions, the actual
reduced-IV marked graphs, the integral source lattice, the
bounded-origin and vertex squares, and 867 fresh complete
nine-variable systems. The first executed anchor-only result
and its script are preserved before this full-source refresh.
The complete primary RJC.1--13 and all named stable anchors
have been read. The a=3 graph constant correction is retained
in (JFP.10); no further geometric discrepancy was found.
The full source/cap/reference and negative-lift arguments,
both mapping-torus signs, and the unbounded image proof are
part of this peer, beyond the finite arithmetic checks.

The scope is precisely (JFP.1--3), the original two additive
caps, zero additional ordinary total norm, every retained
finite character and annular reference, and all selected
lattice-basis fans. Other source-moduli paths and redistributed
nonzero ordinary clutch totals have their own full comparison
problems and are not asserted by this proof.

% End expanded verified cp36 input



\par\medskip\noindent\textbf{Exact construction and audit dependencies.}\par
The complete original F46 joint-cap construction is
\path{notes/part3_iv_star_iii_joint_working.md}.
Its actual point/O graph, normalized blowups, primitive cap generators
and integral sphere maps are retained in both supporting proofs.
The K67 shorter-parameter family is
\path{notes/part3_shorter_scalar_vertical_fans_working.md};
the exact K73 all-fan product and prism construction is
\path{notes/part3_scalar_all_fans.tex}.
The completed K114--115 source and endpoint comparisons are
\path{notes/part3_e6_iii_endpoint_fourfold_peer.md} and
\path{notes/part3_e6_complete_scalar_returns_working.md}.
The canonical two-cap page and its signed source graph are
\path{notes/part3_joint_additive_returns_working.md}.

The primary executable is
\path{code/part3_root/e6_iii_joint_endpoint_derivation.py}.
The separate independent executable is
\path{code/part3_fourfold/e6_iii_joint_endpoint_peer.py}.
The independent source and endpoint derivation, fresh full target
systems, original-coordinate marked graphs and complete residue
witnesses are locked with their entire supporting proofs.
Every source, executable, exact certificate and full geometric audit
dependency is listed in
\path{records/k118_export_manifest.json}.
The bounded executions support the displayed morphisms; the source
extension, finite quotient, terminating fan maps and exact modular
identities prove the full stated domains.


% End expanded verified cp36 input

\end{proof}


\subsection{OS45 ordinary-cap collision meridians with the full original reference}
\begin{theorem}[Source result P3K.119 under continuation review]\label{night:P3K.119}
\StatementRecord{night:P3K.119}{open}{Historical events: GPT-6 Astra (OpenAI), ChatGPT Pro mode, effort Pro, S6\_06\_BLIND, 16--18 November 2026: targeted proof check. Also checked in this writing. Continuation mathematical and source checks: GPT-6 Astra (OpenAI), ChatGPT Pro mode, effort Pro, S6\_10\_PAPERS continuation, 19 November 2026: computation or formal certificate; used and rechecked in this run. Continuation metadata events: GPT-6 Astra (OpenAI), ChatGPT Pro mode, effort Pro, S6\_10\_PAPERS continuation, 19 November 2026: used and rechecked in this run; explicitly scoped metadata reconciliation; metadata reconciliation only; no new mathematical check. GPT-6 Astra (OpenAI), ChatGPT Pro mode, effort Pro, S6\_10\_PAPERS continuation, 19 November 2026: used and rechecked in this run; receiving source placement; metadata reconciliation only; no new mathematical check. GPT-6 Astra (OpenAI), ChatGPT Pro mode, effort Pro, S6\_10\_PAPERS continuation, 19 November 2026: used and rechecked in this run; source-row intake; metadata reconciliation only; no new mathematical check. \emph{Pre-continuation recorded proof scope.} Historical event: GPT-6 Astra (OpenAI), ChatGPT Pro mode, effort Pro, S6\_06\_BLIND, 16--18 November 2026: targeted proof check, recorded producing-run re-derivation. Also checked in this writing by GPT-6 Astra (OpenAI), ChatGPT Pro mode, effort Pro, S6\_10\_PAPERS continuation, 19 November 2026: computation or formal certificate for the explicitly listed algebraic substeps in the event ledger; full continuation rederivation and renewed credit comparison remain pending.}{the original source, point line, roots, characters, reference, selected fan and scalar range stated here}
Retain the original OS45 marked source and point line
\[
\begin{gathered}
R=t^2+a,\qquad f=R^2+t,\qquad
g=R^3+\frac32tR+\frac38,\\
y^2=x^3-3fx+2g,\qquad G=(R,\sqrt3/2),\qquad
\mathcal M=\mathcal O([8G]-[O]),
\end{gathered}
\]
with source translation by \(-G\), selected whole nodal fibre at
\(t=t_c\), and original height \(q_{\rm sel}=t-t_c\). Retain the
original order-one cap at \(t=0\), its positive norm
\(\psi(\widehat\gamma)=1\), every prescribed source and physical
integer in its full period reference, every prescribed integral
lattice-basis selected fan, and an original admitted fixed nonzero
scalar \(\lambda\).

On the original ordered selected-source cover
\[
a=-u/3-u^4/9,\qquad t_c=1/(2u)-u^2/3,\qquad
v^2=4u^3+3,
\]
the ordinary-cap collision locus consists of exactly twenty-four
marked points:
\[
\begin{gathered}
16a^3+3=0,\\
2187-16u^3(u^3+3)^3
=-(2u^3-3)(8u^9+84u^6+342u^3+729)=0.
\end{gathered}
\]
Each of the twelve distinct \(u\)-values has two distinct ordered
\(v\)-values. The six points with \(u^3=3/2\) have \(t_c=0\);
the other eighteen have \(t_c\ne0\). These two loci and the other
excluded source boundaries have the disjointness and full local
coordinate inverses proved in OC45.1--2 and ORC.1--2.

Every positive meridian at any of these twenty-four marked points,
closed by its actual original source, line and cap-reference
endpoint, has
\[
[\Phi]=0\qquad\text{in }\Theta_7=\mathbb Z/28.
\]
Every inverse meridian has the same zero class through the
reciprocal source, Cartier, cap, endpoint and transport maps.

The comparison giving this return is an explicit family map.
At each collision let \(s=t_n(a)\) be the actual moving nodal
value, with \(s(a_*)=0\) and \(ds/da=-2a_*\ne0\). Its full
source chart has \(XY=t-s\), the inverse
\[
t=s+XY,\qquad y=(X+h_nY)/2,\qquad
2(x-\xi)\sqrt{x+2\xi}=h_nY-X,
\]
and the complete functions \(\xi,h_n\) of OC45.8--11.
In the actual point coefficient \(V=f_{P^8}(Z;Q)v_M\), the
original periods are
\[
h:(Z,V)\longmapsto(QZ,P^8V),\qquad
B:(Z,V)\longmapsto(P^{-1}Z,CV),\qquad
C=\lambda(t-t_c)\beta .
\]
Here \(Q=(t-s)\kappa\), \(P\), \(\beta\), the fixed Hom
normalization, and all higher coefficients remain the original
functions. At the central node,
\[
P_*=-2\pm\sqrt3,\qquad
\kappa(0,0)=\frac1{288a_*^2},\qquad
\beta(0,0)=C_*\frac9{a_*^2}.
\]
Their full Cartier maps and their target-point inverses occur in
the construction.

Move the reference cap to the protected ordinary point \(t=d\)
by the full based integral and holomorphic comparison
OC45.39--56. Its gluing is
\[
\mathcal F_w(Z,V)=(e^{-w\log P}Z,e^{w\log C}V),\qquad
\Psi_d=\mathcal F_{\log(t-d)/(2\pi i)},\qquad
\Psi_d^{-1}=\mathcal F_{-\log(t-d)/(2\pi i)}.
\]
These maps extend on the original source, point/O modules,
infinity components and every prescribed selected fan by
OC45.29--35. The resulting family \(Y\) is smooth and proper
over the full \(s\)-disc with the original oriented sphere fibre.

On the original cap, the complete retained period is
\[
\gamma_{\rm ref}
=\frac{(-\log P,\log C)}{2\pi i}
+k\,\frac{(\log Q,8\log P)}{2\pi i}
+r\alpha+d_\delta\delta .
\]
Its positive parameter jump is \(k\alpha\) at the eighteen
unselected-cap points and \(k\alpha+\delta\) at the six
selected-cap points. The actual integral preimages, based-cell
inverse and Lie-cochain solution carry each of those full jumps.
With \(F_\theta:X_\theta\to Y_\theta\) the resulting comparison
and \(E_X,E_Y\) the forward marked endpoints, they give
\[
F_1E_XF_0^{-1}=CE_Y,\qquad
C\in\mathbb C^*_{\rm physical}.
\]
For the actual transport \(P_Y:Y_0\to Y_1\), the full return
identity is
\[
F_0\bigl(E_X^{-1}F_1^{-1}P_YF_0\bigr)F_0^{-1}
=E_Y^{-1}C^{-1}P_Y .
\]
The explicit physical isotopy \(e^{\theta\log C}\) and the
disc transport prove the stated class with its prescribed
endpoint. The argument-preserving radial scalar comparison
of OC45.10 carries these same maps to every original admitted
fixed scalar. The complete scope is these twenty-four
ordinary-cap collision meridians and their inverses.
\end{theorem}
\begin{proof}[Supplied source argument; continuation review in progress]

% Expanded verified cp36 input: notes/part3_os45_cap_collisions_proof.tex
\par\bigskip\noindent\textbf{Complete supporting proof 1: OC45.}\par


% Expanded verified cp36 input: notes/part3_os45_cap_collisions_oc45_proof.tex
\emph{Completed proof and retained preparation history.} The original opening and OC45.1-\/-5 retain the local-stage scope; OC45.6-\/-11 supplies the full original-reference comparison and all twenty-four points. The original closing preparation-stage audit notice is reproduced verbatim. The complete OCP, OCGC and ORC supporting proofs follow below. Their completed geometric audits, exact certificates and source locks are recorded in \path{records/k119_export_manifest.json}.

\par

\medskip\noindent\textbf{OS45 ordinary-cap collisions: exact local source and cap maps}

\par\nobreak

This separate continuation preserves K116 and its complete unselected-II
disc proof. It concerns the original cap at \(t=0\) in the original
OS45 source. The full prescribed cap-period comparison is recorded
explicitly at the end. No return for that prescribed datum is assigned
before its integral reference comparison is constructed.

\par

\medskip\noindent\textbf{OC45.1. The twenty-four marked points and their two actual strata}

\par\nobreak

Retain
\[
\begin{gathered}
 R=t^2+a,\quad f=R^2+t,\quad
 g=R^3+\frac32tR+\frac38,\\
 y^2=x^3-3fx+2g,\quad G=(R,\sqrt3/2),
 \qquad q_{\rm sel}=t-t_c.
\end{gathered}
 \tag{OC45.1}
\]
The selected source cover is
\[
 a=-u/3-u^4/9,\quad t_c=1/(2u)-u^2/3,\quad
 v^2=4u^3+3.
 \tag{OC45.2}
\]
The ordinary cap becomes singular exactly at
\[
 16a^3+3=0,\qquad
 2187-16u^3(u^3+3)^3=0.
 \tag{OC45.3}
\]
Writing \(w=u^3\), its complete factorization is
\[
 2187-16w(w+3)^3
       =-(2w-3)(8w^3+84w^2+342w+729).
 \tag{OC45.4}
\]
The twelve distinct \(u\)-values each have two ordered values of \(v\).
At \(u^3=3/2\), the actual selected height satisfies \(t_c=0\).
This gives six selected-cap collision points.
The other three \(w\)-roots give eighteen points with \(t_c\ne0\);
the node at the original cap is then an unselected source node.
The two sets are disjoint by their polynomial gcd. Both are also
disjoint from the selected-cusp, infinity-\(I_9\), and unselected-II
loci of O45.2.

The current construction below treats the latter eighteen points.
The six selected-cap points retain their own rank-two selected/cap
interaction; their actual polynomial and source label are kept in
OC45.2-\/-4.

\par

\medskip\noindent\textbf{OC45.2. Original nodal source, moving height and complete point line}

\par\nobreak

Fix one of the eighteen points and write \(a_*\) for its source value.
Then \(a_*^3=-3/16\), the original node is
\((x,y,t)=(-a_*,0,0)\), and the nodal cubic is
\[
 y^2=(x+a_*)^2(x-2a_*).
 \tag{OC45.5}
\]
The original finite discriminant factor is
\(D=48a(t^2+a)^2+8t^3+72at+9\). At this node
\[
 D_t=72a_*,\quad D_a=144a_*^2.
 \tag{OC45.6}
\]
Thus its actual critical base value is a single-valued holomorphic
function \(s=t_n(a)\), with
\[
 s(a_*)=0,\qquad \frac{ds}{da}(a_*)=-2a_*\ne0.
 \tag{OC45.7}
\]
The inverse \(a=a(s)\) is fixed by this original implicit equation.
Set \(q_n=t-s\), keeping the original cap coordinate \(t\).

Choose \(\xi(q_n,s)=\sqrt{f(s+q_n,a(s))}\) with value
\(\xi(0,0)=-a_*\), and put
\[
 h_n(q_n,s)=
  \frac{2(g(s+q_n,a(s))-\xi(q_n,s)^3)}{q_n}.
 \tag{OC45.8}
\]
Its numerator vanishes at \(q_n=0\).
At the collision
\[
 f_t=1,\quad g_t=3a_*/2,\qquad h_n(0,0)=6a_*\ne0.
 \tag{OC45.9}
\]
The actual source chart and inverse are
\[
 X=y-(x-\xi)\sqrt{x+2\xi},\qquad
 Y=h_n^{-1}\bigl(y+(x-\xi)\sqrt{x+2\xi}\bigr),\qquad XY=q_n,
 \tag{OC45.10}
\]
\[
 q_n=XY,\quad t=s+XY,\quad y=(X+h_nY)/2,\quad
 2(x-\xi)\sqrt{x+2\xi}=h_nY-X.
 \tag{OC45.11}
\]
The last implicit derivative is \(2\sqrt{3\xi}\ne0\).
Consequently these are inverse holomorphic source maps over the
joint \(s\)-disc, not merely formulas on its punctured boundary.

On the central nodal normalization put
\[
 x=S^2+2a_*,\quad y=S(S^2+3a_*),\quad
 c^2=-3a_*,\quad Z=(S-c)/(S+c).
 \tag{OC45.12}
\]
The original section has \(S_G=\sqrt3/(4a_*)\).
Its square is \(-a_*\), so \(c/S_G=\pm\sqrt3\), and therefore
\[
 P_*=Z(G)=-2+\sqrt3\quad\hbox{or}\quad -2-\sqrt3
 \tag{OC45.13}
\]
according to the chosen ordered branch.
The two displayed values are reciprocal. Neither is zero,
one, or a root of unity.

The original Miller functions give a useful full nodal identity.
For an arbitrary nonzero node parameter \(c\) and marked
\(P=Z(G)\), the original tangent normalization at \(O\) gives
\[
 f_m(Z)=(-2c)^{m-1}\frac{1-P^m}{(1-P)^m}\,
       \frac{(Z-P)^m}{(Z-P^m)(Z-1)^{m-1}},
       \qquad m=2,4,8.
 \tag{OC45.14}
\]
Indeed the horizontal divisor is the required one, the values
at the two node branches agree, and its original principal
coefficient at \(O\) is \(S^{m-1}\), with
\(S=c(1+Z)/(1-Z)\). Alternatively the three original tangent
quotients successively give the same constants.
The whole chronological coefficient is
\[
 B_{8,1}(Z)=
  \left(-2c\frac{1+P}{1-P}\right)^8
  \frac{(Z-P)(Z-P^8)}{(Z-P^9)(Z-1)}.
 \tag{OC45.15}
\]
For the actual point coordinate
\(V_{\rm pt}=f_{P^8}(Z)v_M\),
where \(f_b=(1-Z/b)/(1-Z)\) on the nodal normalization,
the directed coefficient is
\[
 B_{8,1}(Z)\frac{f_{P^8}(Z/P)}{f_{P^8}(Z)}
       =(-2S_G)^8.
 \tag{OC45.16}
\]
All four point/O factors are retained before this cancellation.
At the cap node this value is
\[
 (-2S_G)^8=\frac{81}{256a_*^8}=\frac9{a_*^2}.
 \tag{OC45.17}
\]
This also verifies the independent selected-node formula:
at \(S_G=-\sqrt3/(2u)\), the value in OC45.16 is \(81/u^8\).

Let \(Q=q_n\kappa(q_n,s)\) and \(P=P(q_n,s)\) be the full
inverse-Tate source functions with this ordered branch and
the original reference scale. Their first coefficient is
\[
 \kappa(0,0)
   =-\frac{h_n(0,0)}{1728\xi(0,0)^3}
   =\frac1{288a_*^2}.
 \tag{OC45.18}
\]
The full physical point-line coefficient is defined, with
its reciprocal, by evaluation of the actual Miller/theta
module ratio at a protected regular point. Write it as
\(\beta(q_n,s)\), including the original normalization
\(C_*\). Then \(\beta(0,0)=C_*9/a_*^2\); its full function
is retained below. Since \(q_{\rm sel}(0,0)=-t_c(a_*)\ne0\),
the complete suspension multiplier
\[
 C(q_n,s)=\lambda q_{\rm sel}\beta(q_n,s)
 \tag{OC45.19}
\]
is an extending unit. The original period arrows are
\[
 h:(Z,V_{\rm pt})\mapsto(QZ,P^8V_{\rm pt}),\qquad
 B:(Z,V_{\rm pt})\mapsto(P^{-1}Z,CV_{\rm pt}).
 \tag{OC45.20}
\]
No central value replaces \(P,\kappa,\beta\), or \(C\) in
these arrows.

\par

\medskip\noindent\textbf{OC45.3. A complete extending local suspension flow}

\par\nobreak

Choose the actual analytic logarithms of the nonzero
functions \(P\) and \(C\) on the small closed joint
\((q_n,s)\)-bidisc. Their values at its reference point
are prescribed. For a complex number \(w\), define
\[
 \mathcal F_w(Z,V_{\rm pt},q_n,s)
   =\left(e^{-w\log P}Z,\ e^{w\log C}V_{\rm pt},q_n,s\right).
 \tag{OC45.21}
\]
Its inverse is \(\mathcal F_{-w}\). The two maps commute
with \(h\) and \(B\), retain the original \(q_n\), and
\[
 \mathcal F_{w+1}=B\mathcal F_w,\qquad
 \mathcal F_{w+w'}=\mathcal F_w\mathcal F_{w'}.
 \tag{OC45.22}
\]
These are identities of the full coefficients.

They extend through the whole nodal source and point
line. On the Tate chain
\(X_i=Z/Q^i,\ Y_i=Q^{i+1}/Z\), the source arrow is
\[
 (X_i,Y_i)\mapsto
       (e^{-w\log P}X_i,\ e^{w\log P}Y_i).
 \tag{OC45.23}
\]
It preserves \(X_iY_i=Q\) and has the written reciprocal
inverse. The exact \(Q=q_n\kappa\) coordinate change
then carries it to the original source chart.
In the original point frame its line coefficient is
\[
 e^{w\log C}\,
       \frac{f_{P^8}(Z)}{f_{P^8}(e^{-w\log P}Z)}.
 \tag{OC45.24}
\]
The defining numerator and denominator modules are
the source and target point/O modules. Their graphs
are carried by the displayed source arrow, so this
is an isomorphism at their zeros and poles as well.
At each node both point sections lie in its smooth
identity component, and the Tate-chain coefficient
in OC45.21 is an actual unit. Thus the complete
map extends across the node, with its inverse.
This proves the local flow as a source/line
isomorphism, not just an equality of a central
coefficient.

In additive logarithmic coordinates its period at
time one is
\[
 \gamma_0=\frac1{2\pi i}
               \binom{-\log P}{\log C}.
 \tag{OC45.25}
\]
The subscript names this particular extending lift.
The prescribed original cap-period lift has not
yet been identified with it.

\par

\medskip\noindent\textbf{OC45.4. Exact enlargement of the cap with this period}

\par\nobreak

For the period \(\gamma_0\), the cap-to-ordinary
map is the original formula, now with its full
extending flow:
\[
 \Psi_0(\zeta,t,s)
       =\mathcal F_{\log t/(2\pi i)}(\zeta,t,s),
 \qquad
 \Psi_0^{-1}
       =\mathcal F_{-\log t/(2\pi i)}.
 \tag{OC45.26}
\]
A logarithm branch change multiplies by \(B\)
by OC45.22. Hence both maps are single-valued
on the actual suspension quotient.
For \(t\ne0\) this includes the moving source
node \(t=s\); the extension follows from
OC45.23-\/-24. Thus a fixed larger cap disc may
contain the moving node.

Here is the precise comparison with the usual
small ordinary cap for \(s\ne0\). On the
small cap use the identity. On the intervening
ordinary region use \(\Psi_0^{-1}\) into the
large cap, and outside the large cap use the
identity on the ordinary region. On the small
cap overlap the original gluing is \(\Psi_0\),
so the two maps compose to the identity.
On the large-cap boundary its gluing is
again \(\Psi_0\), giving the same equality.
These piecewise maps therefore glue to a
holomorphic isomorphism. The reciprocal
piecewise maps give its inverse. Their full
source, point-module and scalar coefficients
are precisely OC45.21-\/-26.

The large cap copy is the original nodal
source/line quotient, with \(q_n=t-s\).
At \((t,s)=(0,0)\) its source total space is
smooth by OC45.10-\/-11; its projection to \(s\)
is a submersion. The lift \(B\) is uniformly
contracting for a common sufficiently small
fixed scalar, since OC45.19 is nonzero and
the line-module map is holomorphic on the
compact local source. Its line-complement
quotient is therefore smooth and proper over
the closed cap disc. Gluing its fixed outer
annulus by OC45.26 gives an actual cap
family over the full \(s\)-disc. All other
source, selected-fan and infinity pieces
retain the O45 maps on this disc.

For reference, the exact base blowup on the
punctured comparison is
\[
 t=s\tau,\qquad
 Q=s(\tau-1)\kappa(s(\tau-1),s),\qquad
 \log t=\log s+\log\tau.
 \tag{OC45.27}
\]
The additive coordinate
\(\bar\zeta=\zeta_{\rm cap}
+(\log s)/(2\pi i)\,\gamma_0\)
then has ordinary transition
\(\zeta_{\rm ord}=\bar\zeta+
(\log\tau)/(2\pi i)\,\gamma_0\).
Both inverse translations retain the complete
period vector. Around the \(s\)-loop this
coordinate changes by precisely \(B\);
it is already the period of the quotient.
No fractional source-power coefficient is
declared to extend as a single-valued
function on the original source chart.

\par

\medskip\noindent\textbf{OC45.5. The exact remaining prescribed-reference equation}

\par\nobreak

On a cut parameter annulus, compare the
prescribed original cap period with OC45.25.
They are lifts of the same original
suspension map on the torus. In the complete
original source/physical lattice their
difference is
\[
 \gamma_{\rm ref}-\gamma_0
        =k\,h+r\,\alpha+d\,\delta,\qquad k,r,d\in\mathbb Z,
 \tag{OC45.28}
\]
where \(h\) is the full source period in
OC45.20, \(\alpha\) is its source circle
period and \(\delta\) the physical circle.
Every one of these integers is retained;
the difference has suspension coordinate
\(\psi=0\). Replacing the cap map by
OC45.26 without the actual comparison
for OC45.28 would change the prescribed
cap datum.

The original five source matrices of K10
give a primitive image lattice
\(\sum_i\operatorname{im}(T_i^{-t}-I)
=\ker\psi\). This exact integral equation
is rechecked in the O45P certificate.
The next step is its full groupoid
realization for OC45.28: explicit integral
preimages on the ordinary source collars,
their extending source/point-line arrows
at all five fillings, a parameter-dependent
Lie cochain with its full logarithms,
and the resulting directed endpoint
square. This is the remaining comparison
before assigning a return for the
prescribed original cap meridian.
The construction OC45.1-\/-27 keeps the
local source and actual cap model ready
for that comparison.

\par

\medskip\noindent\textbf{OC45.6. The original extending translation sheaf on all five fillings}

\par\nobreak

The opening status and OC45.1-\/-5 above are retained as the completed
local-stage record. The following construction supplies the original
reference comparison left open there. It also moves the cap to a
protected ordinary point, using explicit maps, and therefore treats
the six selected-cap points as well as the eighteen unselected-cap
points. The cap retains its original order one and positive norm
\(\psi(\widehat\gamma)=1\). No additional root or finite character is
introduced.

Fix the original five source discs and the prescribed selected fan.
All source, point-line and scalar functions in this section are the
ones in O45.1-\/-5, with the full functions of OC45.18-\/-20 on a nodal
disc. A translation on a regular source torus is written in its
actual point coordinates as
\[
 \mathcal T_{A,D}(Z,V,q,s)
       =(e^A Z,e^D V,q,s),\qquad
 \mathcal T_{A,D}^{-1}=\mathcal T_{-A,-D}.
 \tag{OC45.29}
\]
Here \(A,D\) are the two coefficients of the actual Lie vector;
they can depend on the original base and parameter. The scalar
\(\lambda\), the source height and every point-line function remain
fixed in this arrow. In the original meromorphic point section the
line coefficient is exactly
\[
 v_M\longmapsto e^D
       \frac{f_{P^8}(Z;Q)}{f_{P^8}(e^A Z;Q)}v_M.
 \tag{OC45.30}
\]
The inverse evaluates its reciprocal at \(e^A Z\). The numerator and
denominator are the actual source and target point/O modules, so
this is their module isomorphism also at their zeros and poles.
This same formula gives the ordinary torus exponential, using the
complete inverse marked uniformization.

At an unselected node the source-chain coordinates in OC45.23
transform as \((X_i,Y_i)\mapsto(e^A X_i,e^{-A}Y_i)\).
Thus both the flow and its inverse extend on the whole source
chain and its genuine point line. At the selected zero the original
periods have the full form
\[
 h=(\kappa q,P^8),\qquad
 B=(P^{-1},\lambda q\beta)
 \tag{OC45.31}
\]
as multiplicative source/line columns. The two valuation vectors
are \((1,0)\) and \((0,1)\). For every actual character
\(Z^rV^d q^e\) on every cone of the prescribed periodic fan,
OC45.29 multiplies that coordinate by \(e^{rA+dD}\).
These are nonzero analytic units at its unchanged original height,
with inverse \(e^{-rA-dD}\). They agree on all toric overlaps and
commute with the full periods, including \(\kappa,P,\lambda,\beta\).
This proves the extension on each prescribed fan by its actual
coordinate map.

At the original infinity put \(z=1/t\). O45.4 gives the full
\[
 P_\infty=z^3\pi(z,s),\qquad Q_\infty=z^8\kappa_\infty(z,s),
 \qquad b_0=P_\infty^8/Q_\infty^3,
 \qquad b_0(0,s)=-1/(9a(s)^3).
 \tag{OC45.32}
\]
All displayed base coefficients here remain analytic units near
the cap-collision disc. The exact product identity
\[
 f_{Q_\infty^3b_0}(Z;Q_\infty)
       =(-Z/b_0)^3Q_\infty^{-6}f_{b_0}(Z;Q_\infty)
 \tag{OC45.33}
\]
identifies the original point coordinate with the multidegree-zero
coordinate \(V_0=f_{b_0}(Z;Q_\infty)v_M\). Consequently OC45.29
acts there by
\[
 Z\longmapsto e^A Z,\qquad V_0\longmapsto e^{D-3A}V_0.
 \tag{OC45.34}
\]
The factors \(b_0^3Q_\infty^6\) in the reciprocal coordinate
comparison are retained; they cancel in this displayed action
because the original base is fixed. In that frame the actual
period reindex is
\[
 H=h^{-1}B^{-3},\quad K=h^3B^8,\qquad
 h=H^8K^3,\quad B=H^{-3}K^{-1}.
 \tag{OC45.35}
\]
It has determinant one, source valuations one and zero, and the
complete line multiplier of \(K\) is the original contracting
unit. Formula OC45.34 commutes with both full columns. On the
primitive rank-one chain it and its inverse multiply every
coordinate by its actual unit character. Thus it extends through
the entire infinity filling. No constant term is substituted for
either period or the point-to-tensor factor.

Let \(\mathscr E\) denote this holomorphic Lie bundle on the
original base, and \(\mathscr H\) the kernel of its exponential
in the sheaf \(\mathscr A\) of the extending translations just
constructed. The period sheaf is \(j_*\Lambda\), with local stalk
\(\ker(H_i-I)\). Here is its actual local identification. The
logarithmic radial coefficients of the full periods have exactly
the primitive valuation maps of OC45.31 and OC45.35. A period
extends in the Lie bundle precisely when these radial coefficients
vanish. The remaining angular periods are the logarithms of
the extending unit actions in OC45.29, with the full point frame
OC45.30. This is exactly the kernel of the original monodromy.
On two slit sectors and the centre the marked circle complex
\(\Lambda\xrightarrow{H_i-I}\Lambda\) has its degree-one
coinvariant term cancelled by the identical annular term.
Thus the invariant extension has no higher cohomology on a
source disc. This is an integral cancellation in the stated
circle coordinates.

Logarithms of the two nonzero unit multipliers give the local
inverse to the exponential. The primitive period reindexes above
remove the integral radial orders before these logarithms are
taken. The resulting exact sequence is therefore the sequence
of the actual source and point-line maps:
\[
 0\longrightarrow\mathscr H\longrightarrow\mathscr E
       \xrightarrow{\operatorname{Exp}}\mathscr A\longrightarrow0.
 \tag{OC45.36}
\]
Only the existing source fillings and their same fan occur here.

The Lie source line can be calculated directly on OS45. In the
original minimal infinity chart \(x=z^{-2}X\), \(y=z^{-3}Y\),
the relative invariant differentials satisfy
\[
 \omega_t=\frac{dx}{2y}
           =z\frac{dX}{2Y}=z\omega_\infty .
 \tag{OC45.37}
\]
Their dual vectors therefore have
\(e_\infty=t^{-1}e_0\). At a nodal point the relative dualizing
differential is a unit times \(dX_i/X_i=-dY_i/Y_i\), and its
dual is the extending field \(X_i\partial_{X_i}-Y_i\partial_{Y_i}\).
At the selected fan this is the source unit action already
displayed; at infinity it is carried by OC45.33-\/-35. Hence the
source Lie line is literally \(\mathcal O(-1)\) with these frames.
The physical line is the global field \(\delta\) whose flow is
multiplication by \(e^{2\pi iz}\). The complete Lie extension is
\[
 0\longrightarrow\mathcal O\delta\longrightarrow\mathscr E
       \longrightarrow\mathcal O(-1)\longrightarrow0.
 \tag{OC45.38}
\]
In point-module charts its off-diagonal entry is the source
invariant vector applied to the logarithm of the entire actual
line transition, including OC45.30 and OC45.33, divided by the
chosen logarithmic normalization. No central line constant
replaces this derivative.

\par

\medskip\noindent\textbf{OC45.7. The complete integral comparison in the actual based cells}

\par\nobreak

Use exactly K10\textquotesingle s five matrices \(T_i\), in order
\((\psi,u,w,\delta)\), and set
\[
 H_i=T_i^{-t},\quad N_i=H_i-I,\quad
 P_0=I,\quad P_i=H_0\cdots H_{i-1},\quad
 M_i=P_iN_i=P_{i+1}-P_i.
 \tag{OC45.39}
\]
The product is \(H_0\cdots H_4=I\). The complete source basis
and transported cells are those of K10, not an untransported
sum of abstract image lattices. For selected zero at site
\(\sigma\in\{1,2,3,4\}\) put
\[
 d_1=e_\psi,\quad d_2=e_\psi+2e_w,\quad
 d_3=e_\psi+e_w,\quad d_4=e_\psi .
 \tag{OC45.40}
\]
Direct inversion of each of the four actual central-entry choices
gives
\[
 N_\sigma d_\sigma=e_\delta,\qquad
 N_1(-e_u)=e_w,\qquad N_3e_w=e_u+8e_\delta.
 \tag{OC45.41}
\]
For an arbitrary retained difference
\(v=Ae_u+Be_w+Ce_\delta\), start with all \(r_i=0\), set
\[
 r_1=-Be_u,\qquad r_3=Ae_w,\qquad
 r_\sigma\mathrel{+}=(C-8A)d_\sigma .
 \tag{OC45.42}
\]
Repeated indices receive the indicated sum. Then
\(\sum_iN_ir_i=v\), with no division. Its actual based-cell
preimages and inverse are
\[
 t_j=r_j-\sum_{i>j}N_i r_i,\qquad
 r_j=t_j+\sum_{i>j}N_i r_i\quad(j=4,3,2,1,0).
 \tag{OC45.43}
\]
The identity \(\sum_{j<i}M_j=P_i-I\) proves
\(\sum_jM_jt_j=\sum_iN_ir_i=v\).
This is an upper triangular integral coordinate isomorphism
with identity diagonal.

For clarity the full cone sign is retained too. With exterior
one-cochains \(p_0,\ldots,p_3\) and annular zero-cochains \(z_i\),
the degree-two boundary is
\[
 b_i=p_i-N_iz_i\ (i<4),\qquad
 b_4=-H_4\sum_{i<4}P_ip_i-N_4z_4.
 \tag{OC45.44}
\]
The coordinate \(v_*=\sum_iP_i b_i\) cancels every exterior
\(p_i\), and its boundary is \(-\sum_iM_i z_i\).
Thus the positive difference \(v\) in OC45.42 is produced
by exactly \(z_i=-t_i\). These are the restrictions along
the actual based cuts. Smooth cap discs add their transported
period vector to \(v_*\); their monodromy is the identity.

We also retain the entire degree-one calculation. The primitive
local image basis is
\[
 g_0=e_w+3e_\delta,\quad g_1=e_w,\quad
 g_2=e_u+3e_w+16e_\delta,\quad
 g_3=e_u+8e_\delta,\quad g_4=e_u-3e_w,
 \tag{OC45.45}
\]
with one additional \(e_\delta\) at the selected site.
Record these six coefficients in that order, placing the
additional coefficient last. In these actual integral bases
the parabolic complex has matrices
\[
 D=\begin{pmatrix}
 -3&-8&0&0\\0&-1&0&0\\-2&-3&1&0\\
 -1&0&1&0\\0&3&1&0\\1&0&0&0
 \end{pmatrix},\qquad
 B=\begin{pmatrix}
 0&0&0&0&0&0\\
 0&0&1&-2&1&0\\
 1&1&-6&9&-3&0\\
 3&0&-8&8&0&1
 \end{pmatrix}.
 \tag{OC45.46}
\]
Here \(D(v)=(N_iv)_i\); the first five columns of \(B\) are
\(P_i g_i\), and its last is \(P_\sigma e_\delta=e_\delta\).
These matrices are independently obtained for every selected
site after the displayed placement of the extra generator.
For every integral cocycle,
\[
 Bz=0\quad\Longrightarrow\quad
 z=D(z_5,-z_1,z_3+z_5,0)^t.
 \tag{OC45.47}
\]
Let \(D'\) be the first three columns of \(D\), \(B'\) the
last three rows of \(B\), and define
\[
 Lz=(z_5,-z_1,z_3+z_5)^t,\qquad
 J=\begin{pmatrix}
 -24&-8&3\\0&0&0\\-9&-3&1\\0&0&0\\
 10&3&-1\\0&0&0
 \end{pmatrix}.
 \tag{OC45.48}
\]
Their full integer identities are
\[
 LD'=I_3,\quad B'J=I_3,\quad D'L+JB'=I_6,\quad B'D'=0.
 \tag{OC45.49}
\]
They give explicit chain contraction and splitting maps,
including all torsion information. With the actual local
invariant-stalk cancellation of OC45.36, they prove
\[
 H^0(\mathscr H)=\mathbb Z\delta,\qquad
 H^1(\mathscr H)=0,\qquad
 H^2(\mathscr H)\xrightarrow{\ \psi\ }\mathbb Z.
 \tag{OC45.50}
\]
The complete independent derivation, including the twenty
original inverse matrices and the full twenty-by-twenty map
OC45.43, is retained in OCP.1-\/-6 and its executable certificate.

\par

\medskip\noindent\textbf{OC45.8. The holomorphic comparison with every original reference integer}

\par\nobreak

Compare two actual order-one cap constructions on this same
source, line, scalar and selected fan. Each cap has the original
positive norm. The first is the prescribed cap at \(t=0\).
The second will be placed at a protected ordinary value \(t=d\)
in OC45.10. On a cut parameter disc join their based cap collars
by a specified path in the original regular source base.
Transport the complete original lattice along this path.
The difference of their two cap cocycles in the coordinate
of OC45.44 is an integral vector \(v\) with \(\psi(v)=1-1=0\).
Every other component of \(v\) is retained as \(A,B,C\) in
OC45.42. The two base logarithms are \(\log t\) and
\(\log(t-d)\), respectively; neither is replaced by the other.
The actual collar maps are their full plus-log translations.

Take logarithms of the actual difference of these translation
gluing maps on slit overlaps. Their triple-overlap discrepancy
is their full period matrix \(\Pi(t,s)\) times the integral
two-cocycle just described. Subtract the periods of the precise
total one-cochain OC45.42-\/-44. The actual \(e_\psi\) terms
in \(d_\sigma\) are the original suspension deck, with its
whole source and line multiplier OC45.31 or OC45.35.
They are inserted at both ends of the corresponding overlap.
The corrected Lie cocycle is single valued on each annulus:
its two slit branches differ by the cancelled \(N_i z_i\).
The full source periods in \(\Pi\), including \(\log Q\),
\(\log P\), the physical period and the suspension logarithm,
remain in these subtractions. No period is declared constant
as a parameter function.

Here is a reproducible analytic solution of this remaining
Lie cocycle. Use the actual two based Lie frames of OC45.38.
Their transition has the form
\[
 \mathcal L(t,s)=
     \begin{pmatrix}t^{-1}&0\\a(t,s)&1\end{pmatrix}.
 \tag{OC45.51}
\]
The entire coefficient \(a(t,s)\) is the point-module derivative
specified after OC45.38. Holomorphic lifts of the source frame
on the two affine base charts are obtained by the scalar Cauchy
operator on those charts applied to their point-frame differences;
this retains the same coefficient on the overlap.
For a holomorphic overlap vector \((f,g)\), expand in Laurent
series on a fixed protected circle \(|t|=\rho\). Its coefficients
are the actual integrals
\[
 f_j(s)=\frac1{2\pi i}\int_{|t|=\rho}f(t,s)t^{-j-1}\,dt .
 \tag{OC45.52}
\]
The equation
\(\mathcal L(u_\infty,v_\infty)^t-(u_0,v_0)^t=(f,g)^t\)
has the explicit solution
\[
 u_0=-\sum_{j\ge0}f_jt^j,\qquad
 u_\infty=\sum_{j<0}f_jt^{j+1},\qquad
 G=g-a(t,s)u_\infty,
 \tag{OC45.53}
\]
\[
 v_0=-\sum_{j\ge0}G_jt^j,\qquad
 v_\infty=\sum_{j<0}G_jt^j.
 \tag{OC45.54}
\]
The negative powers in the infinity components are holomorphic
in \(1/t\), including the \(j=-1\) constant in \(u_\infty\).
Direct substitution proves both equations and keeps the whole
off-diagonal term.

For the original many-disc cover, choose a fixed smooth partition
of unity \(\rho_k\) on its refinement. If the corrected cocycle is
\(\theta_{ij}\) and its Lie transition is \(\mathcal L_{ij}\),
the smooth cochain
\(s_i^{\rm sm}=\sum_k\rho_k\theta_{ik}\), in frame \(i\),
satisfies \(s_i^{\rm sm}-\mathcal L_{ij}s_j^{\rm sm}=\theta_{ij}\).
Its \(\bar\partial\) is a global \(\mathscr E\)-valued form.
On the two affine charts take its componentwise Cauchy primitives,
using fixed compactly supported extensions and the kernel
\((\pi(t-w))^{-1}\). Their difference on the overlap is holomorphic.
Apply OC45.53-\/-54 to this difference to obtain matching primitives.
Subtracting the resulting global primitive from \(s_i^{\rm sm}\)
produces holomorphic \(s_i\) with exactly the same cocycle equation.
All partitions and integration kernels are fixed in the parameter;
every resulting coefficient is holomorphic in \(s\) on the
chosen cut. This is also an explicit proof of \(H^1(\mathscr E)=0\).

Write the original and reference overlaps in the directed convention
\[
 g^X_{ij}=t_{\theta_{ij}}g^Y_{ij},\qquad
 \theta_{ij}=s_i-\mathcal L_{ij}s_j.
 \tag{OC45.55}
\]
Then the full family map and inverse are
\[
 F_i=\operatorname{Exp}(-s_i),\qquad F_i^{-1}=\operatorname{Exp}(s_i),
 \qquad F_i g^X_{ij}=g^Y_{ij}F_j.
 \tag{OC45.56}
\]
The last identity is an equality of the whole original source
and point-line coefficients. Formulae OC45.29-\/-35 carry it to
every ordinary, nodal, infinity and prescribed selected-fan
chart, with reciprocal coefficients in the inverse. On a
smooth cap it is the actual torus exponential before its
plus-log gluing. It therefore includes the cap circle and
the complete annular reference, rather than just the integral
class of that reference.

\par

\medskip\noindent\textbf{OC45.9. The full parameter closing square}

\par\nobreak

The parameter monodromy of the prescribed original period is
visible in the local coordinates; it is not suppressed by
OC45.56. At the original cap \(t=0\) put \(q_n=-s\). On a cut
disc write its complete period as
\[
 \gamma_{\rm ref}=\gamma_0+k h+r\alpha+d_\delta\delta,
 \qquad \gamma_0=(-\log P,\log C)/(2\pi i),
 \tag{OC45.57}
\]
where \(h=(\log Q,8\log P)/(2\pi i)\) and \(\alpha,\delta\)
are the original source and physical circle periods.
All integers are constant on this cut in the actual transported
marking. A positive \(s\)-turn has \(h\mapsto h+\alpha\).
At an unselected-cap collision \(C(0,s)\) is an extending unit,
so \(\gamma_0\) returns; at a selected-cap collision it is
exactly \(-\lambda s\beta(-s,s)\), so
\(\gamma_0\mapsto\gamma_0+\delta\). Therefore the full jumps are
\[
\begin{gathered}
 \Delta\gamma_{\rm ref}=k\alpha
       \quad\hbox{at each of the eighteen unselected-cap points},\\
 \Delta\gamma_{\rm ref}=k\alpha+\delta
       \quad\hbox{at each of the six selected-cap points}.
\end{gathered}
 \tag{OC45.58}
\]
The analytic logarithms of \(P,\kappa,\beta\) and every other unit
have their actual continued values; their integer changes are
already included in the chosen \(k,r,d_\delta\) marking.
After its full based transport, each vector in OC45.58 is
inserted into OC45.42-\/-44. It has zero \(\psi\)-coordinate.
This explicitly carries the parameter endpoint of the original
cap, including the additional physical period at the six points.

The group of differences of the local parameter comparisons
can be calculated from the same exact maps. A global section
of OC45.38 projects to a global section of \(\mathcal O(-1)\).
In the two charts its equation is \(u_0=t^{-1}u_\infty\);
the regular Laurent powers on the two sides have no common
nonzero term. Hence it is zero. The remaining component is a
constant physical vector. By OC45.36 and OC45.50,
\[
 0\longrightarrow\mathbb Z\delta\longrightarrow\mathbb C\delta
 \longrightarrow H^0(\mathscr A)\longrightarrow0,\qquad
 [z\delta]\longmapsto(v_M\mapsto e^{2\pi iz}v_M).
 \tag{OC45.59}
\]
The last arrow is onto because the actual integral \(H^1\)
is zero by OC45.47-\/-49. Its kernel is exactly the primitive
physical period. Thus every ratio of the constructed global
comparisons is the indicated physical \(\mathbb C^*\) action,
including its literal action and inverse on every cap and fan.

For completeness, these local comparisons have an explicit
closed holomorphic choice on a thin parameter annulus.
Write \(F_\alpha=c_{\alpha\beta}F_\beta\) on parameter overlaps.
The coefficients \(c_{\alpha\beta}\) are the physical scalars
of OC45.59, retaining the original endpoint substitutions.
For a partition of unity \(\rho_\beta\) put
\[
 \mathfrak a_\alpha=\sum_\beta\rho_\beta\,d\log c_{\alpha\beta},
 \qquad
 \mathfrak a_\alpha-\mathfrak a_\gamma=d\log c_{\alpha\gamma}.
 \tag{OC45.60}
\]
Angular parallel transport gives a nonzero comparison with
holonomy \(c\in\mathbb C^*\). Choose a logarithm of this one
number and a smooth function \(\chi(\theta+1)=\chi(\theta)+1\)
which is constant on the endpoint collars. Multiplication by
\(\exp[-\chi(\theta)\log c]\) makes it periodic with matching
derivatives. Its extension to a thin radial annulus is the
actual connection equation
\[
 \partial_\rho z_\alpha
       =-\mathfrak a_\alpha(\partial_\rho)z_\alpha,\qquad
 z_\alpha(\rho,\theta)=z_\alpha(\rho_0,\theta)
       \exp\left[-\int_{\rho_0}^{\rho}
                 \mathfrak a_\alpha(\partial_r)\,dr\right].
 \tag{OC45.61}
\]
Here the comparison is \(z_\alpha F_\alpha\). The transition
rule in OC45.60 proves that these solutions agree. Its
antiholomorphic derivative is a single scalar form
\(\bar\partial F^{\rm sm}=\eta F^{\rm sm}\). Extend its
coefficient smoothly with compact support in the parameter
plane and set
\[
 u(s)=\frac1\pi\int_{\mathbb C}
                    \frac{\eta_{\bar w}(w)}{s-w}\,dA(w).
 \tag{OC45.62}
\]
Then \(\bar\partial u=\eta\) on the smaller annulus, and
\(e^{-u}F^{\rm sm}\) is an actual nonzero holomorphic family
comparison. This multiplies only by the already proved
physical action, with inverse multiplication. It does not
discard a source period, a cap integer or a parameter seam.

Equivalently, the endpoint can be retained directly on a cut.
Let \(F_\theta:X_\theta\to Y_\theta\) be OC45.56 with all
the continued periods in OC45.58. Let \(E_X:X_0\to X_1\)
and \(E_Y:Y_0\to Y_1\) be the forward marked endpoints.
The constructed cap endpoint is \(F_1^{-1}E_YF_0\); its
local coefficients and reciprocal are exactly OC45.56,
with the integer correction OC45.58 at the final cut.
The supplied reference endpoint is carried with those same
period decks. Its difference is therefore the actual scalar
\[
 C=F_1E_XF_0^{-1}E_Y^{-1}\in\mathbb C^*_{\rm physical},
 \qquad F_1E_XF_0^{-1}=CE_Y .
 \tag{OC45.63}
\]
This scalar is retained, and its explicit isotopy
\(e^{\theta\log C}\) extends on every original line, fan and cap.
For a smooth parallel transport \(P_Y:Y_0\to Y_1\), choose
the transport on \(X\) through \(F_\theta\). Its inverse-endpoint
return satisfies the exact conjugacy equation
\[
 F_0\bigl(E_X^{-1}F_1^{-1}P_YF_0\bigr)F_0^{-1}
           =E_Y^{-1}C^{-1}P_Y .
 \tag{OC45.64}
\]
Thus the closed-annulus construction and the cut construction
give the same original marked return, with its full endpoint.

\par

\medskip\noindent\textbf{OC45.10. A protected ordinary cap across all twenty-four points}

\par\nobreak

Fix any one of the twenty-four marked points OC45.3. The
source critical fibre at the original cap is the ordinary
node OC45.5, with the exact moving height \(s\) of OC45.7.
This holds at the six points too: there the selected nodal
value is precisely \(t_c=s\). The two branch values in
OC45.13 are nonzero and distinct. Choose constants
\(0<\epsilon\ll |d|\) inside the same source-coordinate
neighbourhood, and a closed disc about \(t=d\) of radius
less than \(|d|/4\). Shrink \(|s|<\epsilon\) so that this
disc is disjoint from every source critical value and from
the selected value. Its original coordinate is \(t-d\);
its source functions are still the full functions of
OC45.1 with \(a=a(s)\).

On this protected disc, \(Q=(t-s)\kappa(t-s,s)\),
\(P=P(t-s,s)\), and \(C=\lambda(t-t_c(s))\beta(t-s,s)\)
are nonzero. Choose their original analytic logarithms
on this joint simply connected disc, fixed at one reference
point. The source/line flow is precisely OC45.21 there.
The cap-to-ordinary gluing and its inverse are
\[
 \Psi_d=\mathcal F_{\log(t-d)/(2\pi i)},\qquad
 \Psi_d^{-1}=\mathcal F_{-\log(t-d)/(2\pi i)}.
 \tag{OC45.65}
\]
A logarithm branch changes this by the entire original
suspension \(B\). The maps therefore descend literally
to the compact torus quotient. Its cap period has norm
one and all its source/physical components are retained
by the displayed logarithms.

Call the resulting family \(Y\). At the original point
\(t=0\) it has no extra logarithmic clutch: it uses the
actual unselected source quotient for the eighteen points,
and the prescribed selected periodic fan for the other
six points. Their charts extend over the parameter centre.
For the unselected source this is OC45.10-\/-11 and
OC45.29-\/-30. For the selected zero the exact full
\(q=t-s\) columns OC45.31 have analytic nonzero
\(\kappa,P,\beta\), with \(P_*\) given in OC45.13.
Each fan chart and its reciprocal unit transition hence
extend with the unchanged \(q\). At original infinity
\(a_*^3=-3/16\), so OC45.32 gives
\(b_0(0,0)=16/27\), a nonzero nonidentity point. The full
source/line maps OC45.32-\/-35 extend there as well.
Every other source critical point is simple and distinct
on this parameter disc, by the exact polynomial gcds
in OC45.1 and the nonzero original discriminant derivative.
The remaining O45 source and Cartier-module charts
therefore extend on the same disc.

This is a smooth proper family, as follows from those
actual local maps. The source total space near the moving
node is \(XY=t-s\), so its coordinates \((X,Y,s)\)
exhibit a submersion to \(s\). At the selected fan the
smooth unimodular cone coordinates together with \(s\)
give the same submersion. The original infinity and
ordinary source graphs retain their smooth local
coordinates and extending analytic units. The cap at
\(d\) is the actual smooth torus disc before OC45.65.
Choose a common strictly contracting scalar for these
compact source pieces. Their contracting fundamental
annuli are compact, and OC45.35 gives the primitive
rank-one period coordinates at infinity. The selected
valuation basis OC45.31 and a finite set of cones
modulo its two periods give a compact fundamental
set there; all remaining multipliers are bounded
nonzero units on the closed parameter disc.
The finitely many indicated compact pieces are glued
along the displayed collars and their inverses.
They give a compact Hausdorff total inverse image of
the closed parameter disc, and the local submersions
agree. This proves properness and smoothness of \(Y\)
without an unlisted piece at the cap collision.

The original cap family \(X\) over the punctured parameter
disc and \(Y\) have the same degree-two class by
OC45.42-\/-50. Their full comparison is OC45.51-\/-65;
in particular the source/physical periods of the original
cap are carried by an actual map and inverse. On a chosen
reference parameter, the member \(X\) is the original
K10 sphere transported through the O45 source, line,
scalar, cap and selected-fan charts. For every prescribed
fan its reference sphere is obtained by the finite
canonical word and the actual common-blowup/two-ruling
maps, primitive Gysin and van Kampen maps of F35,
as retained in O45P.6. Thus \(Y\) has a sphere member.
Its proper smooth disc family transports that member,
with its complex orientation, to every fibre.

The comparison is made at the original fixed nonzero
scalar. The usual contraction domain is a punctured disc.
For a larger admitted scalar on a compact source loop,
the explicit radial path
\(\lambda_\tau=[(1-\tau)+\tau\varepsilon]\lambda_0\),
with \(0<\varepsilon<1\), remains contracting because
the entire suspension multiplier is linear in \(\lambda\).
Apply the same source, point-line and cap maps throughout
this compact interval. It joins the given loop to one
in the common disc domain above, with the same endpoint
periods and no scalar turn. Its smooth cylinder transport
conjugates the two marked returns. This carries the
disc-filling conclusion to that original admitted scalar.

\par

\medskip\noindent\textbf{OC45.11. The marked return and the exact completed scope}

\par\nobreak

For \(Y\), the marked endpoint around the parameter disc
is its identity, because its full source, period, line,
cap and fan maps were constructed over that disc. Smooth
parallel transport of this proper submersion gives a
return isotopic to the identity. The comparison
OC45.63-\/-64 carries the original cap reference and
its actual endpoint to this disc family. The remaining
factor \(C^{-1}\) has the literal global physical
isotopy in OC45.63. Hence each of the twenty-four
ordinary-cap collision meridians in the ordered OS45
source cover has
\[
 [\Phi]=0\quad\hbox{in }\Theta_7=\mathbb Z/28
 \tag{OC45.66}
\]
for the original order-one positive-norm cap, its
complete prescribed period reference and every
prescribed lattice-basis selected fan. The inverse
meridian uses the reciprocal maps and has the same
zero class. The statement concerns this original
source, fixed scalar and cap-reference closing;
the full source and linear markings in the comparison
have been kept throughout.

The local enlargement OC45.21-\/-27 is retained as a
second explicit filling for the eighteen unselected
points and its particular extending period. The
protected-point comparison proves the original
reference assertion and the additional six points;
it does not identify the two periods by notation.
All previous local-stage bytes and the PASS64
certificate are preserved under
\path{records/attempts/os45_cap_before_global_reference/}.
The full proof remains a working manuscript until
the independent global source/line/cochain and
endpoint audits are recorded.

% End expanded verified cp36 input


\par\bigskip\noindent\textbf{Complete supporting proof 2: OCP.}\par


% Expanded verified cp36 input: notes/part3_os45_cap_collisions_ocp_proof.tex
\par

\medskip\noindent\textbf{OS45 ordinary-cap period lattice: exact based-cell comparison}

\par\nobreak

This independent calculation uses the five actual K10 matrices in
\path{notes/part3_os45.md}, in the retained order \((\psi,u,w,\delta)\).
It supplies integral preimages and the directed analytic cochain map.
The parameter endpoint of OC45\textquotesingle s prescribed reference remains in that
map; this note assigns no collision return.

\par

\medskip\noindent\textbf{OCP.1. Primitive integral images in the original marking}

\par\nobreak

Let \(T_i\), \(0\leq i\leq4\), be exactly K10\textquotesingle s displayed cohomology
matrices and put \(H_i=T_i^{-t}\), \(N_i=H_i-I\). The selected zero is
first at site4. Direct inversion gives
\[
\begin{aligned}
N_0v&=(-3v_\psi-8v_u)(e_w+3e_\delta),\\
N_1v&=-v_u e_w,\\
N_2v&=(-2v_\psi-3v_u+v_w)(e_u+3e_w+16e_\delta),\\
N_3v&=(-v_\psi+v_w)(e_u+8e_\delta),\\
N_4v&=(3v_u+v_w)(e_u-3e_w)+v_\psi e_\delta.
\end{aligned}
\tag{OCP.1}
\]
Every displayed rank-one coefficient row has gcd1. At site4 the two
coefficient rows separately take the values \((1,0)\) and \((0,1)\).
These are the primitive image lattices of the actual local fillings.

For selected site \(s\in\{1,2,3,4\}\), the source K10 rule adds one to
the central entry of \(T_4\) and subtracts one at \(T_s\). Set
\[
d_1=e_\psi,\quad d_2=e_\psi+2e_w,\quad
d_3=e_\psi+e_w,\quad d_4=e_\psi.
\tag{OCP.2}
\]
For every one of these four actual choices the full inverse-transpose
matrices satisfy
\[
N_s d_s=e_\delta,\qquad N_1(-e_u)=e_w,
\qquad N_3e_w=e_u+8e_\delta.
\tag{OCP.3}
\]
Thus for the retained difference
\(v=Ae_u+Be_w+Ce_\delta\), define the five source cochains by
\[
r_1=-Be_u,\qquad r_3=Ae_w,\qquad
r_s\mathrel{+}=(C-8A)d_s,
\tag{OCP.4}
\]
starting with all other entries zero and adding terms at repeated
indices. Their full image is exactly \(\sum_iN_ir_i=v\).
The first row of every \(H_i\) is \((1,0,0,0)\). Consequently
\[
\sum_i\operatorname{im}N_i=\ker\psi,\qquad
\Lambda/\sum_i\operatorname{im}N_i\xrightarrow{\ \psi\ }\mathbb Z
\tag{OCP.5}
\]
is an integral isomorphism, with section \(1\mapsto[e_\psi]\).
All source-cycle and physical-circle coefficients in \(v\) are kept.

\par

\medskip\noindent\textbf{OCP.2. The actual based polygon and its inverse coordinate change}

\par\nobreak

Use the oriented based boundary word \(H_0H_1H_2H_3H_4=I\) of K10.
Write
\[
P_0=I,\qquad P_i=H_0\cdots H_{i-1},\qquad
M_i=P_i(H_i-I)=P_{i+1}-P_i.
\tag{OCP.6}
\]
These are the transported cellular boundary columns. The following
explicit cochain map carries the raw preimages into these columns:
\[
t_j=r_j-\sum_{i>j}N_i r_i,\qquad 0\leq j\leq4.
\tag{OCP.7}
\]
Indeed \(P_i-I=\sum_{j<i}M_j\), and therefore
\[
\sum_jM_jt_j
=\sum_i\{P_iN_i-(P_i-I)N_i\}r_i
=\sum_iN_ir_i=v.
\tag{OCP.8}
\]
The map in OCP.7 is upper triangular over \(\mathbb Z\), with identity
diagonal blocks. Its actual inverse is the descending recursion
\[
r_j=t_j+\sum_{i>j}N_ir_i,\qquad j=4,3,2,1,0.
\tag{OCP.9}
\]
Here is the full cone direction. On the exterior take four one-cochains
\(p_0,\ldots,p_3\); the last restriction is
\(-H_4\sum_{i<4}P_ip_i\). With local annular zero-cochains \(z_i\),
the degree-two boundaries are
\[
b_i=p_i-N_i z_i\quad(i<4),\qquad
b_4=-H_4\sum_{i<4}P_ip_i-N_4z_4.
\tag{OCP.9a}
\]
The integral coordinate \(v_* = \sum_iP_i b_i\) eliminates every
\(p_i\) and sends this boundary to \(-\sum_iM_i z_i\). To produce
the positive difference \(v\) of OCP.8 use exactly \(z_i=-t_i\).
These are the actual exterior/local restrictions and their signs.

Thus the transported cell quotient is literally \(\mathbb Z\) by
\(\psi\). K10\textquotesingle s primitive local specialization maps and the exact
slit-sector invariant-stalk contraction of K78, OCS.1, identify this
cone with the no-finite-root extending period sheaf. Its
degree-two obstruction group is \(H^2(\mathbf P^1,\mathscr H)=\mathbb Z\)
with the specified generator. The ordinary cap vector has
\(\psi(\gamma)=1\), so its actual clutch represents the displayed
generator. Two norm-one cap vectors have the same degree-two class;
OCP.7 and the sign in OCP.9a supply the full integral cochain for their
difference.

\par

\medskip\noindent\textbf{OCP.3. Directed local logarithms and the full line transition}

\par\nobreak

For OC45\textquotesingle s local flow retain its full period
\[
\gamma_0=( -\log P,\log C)/(2\pi i),\qquad
C=\lambda q_{\rm sel}\beta.
\tag{OCP.10}
\]
Express the prescribed difference \(\gamma_{\rm ref}-\gamma_0\)
in the original K10 common period marking before using OCP.4.
The integer \(\psi\)-component is zero. OCP.4-\/-9 subtract precisely
that integral total coboundary from the actual atlas logarithms.
An occurrence of \(e_\psi\) in \(d_s\) is the actual suspension
deck, including its original source and line multiplier. Its full
deck map is retained at both ends of the affected overlap.

After this integral adjustment the difference of the actual annular
translation cocycles has a vector-sheaf lift. The source translation
line is the original \(\mathcal O(-1)\), and the vertical physical
line is \(\mathcal O\). These are the invariant differential and
physical-circle frames in K41/K44; their exponentials are the extending
translations on the primitive source and selected Cox charts. The
off-diagonal entry is retained from differentiating the full original
point-line transition. In two based source charts its coefficient
matrix has the form
\[
L(t)=\begin{pmatrix}t^{-1}&0\\a(t)&1\end{pmatrix}.
\tag{OCP.11}
\]
Here \(a(t)\) is that actual holomorphic overlap coefficient; no value
is substituted for it. For a vector cocycle \((f,g)\) the directed
equation \(L(t)(u_\infty,v_\infty)^t-(u_0,v_0)^t=(f,g)^t\) is solved by
the following fixed-contour Laurent operators:
\[
u_0=-\sum_{n\geq0}f_nt^n,\qquad
u_\infty=\sum_{n<0}f_nt^{n+1},\qquad
G=g-a(t)u_\infty,
\tag{OCP.12}
\]
\[
v_0=-\sum_{n\geq0}G_nt^n,\qquad
v_\infty=\sum_{n<0}G_nt^n.
\tag{OCP.13}
\]
The equations follow by direct substitution, including the entire
\(a(t)u_\infty\) term. The Laurent coefficients are the Cauchy
integrals on the same protected base circle. Thus the solution retains
the full line module. On the protected parameter sector the actual
fixed-contour Cauchy integrals are holomorphic parameter functions.

With gluing convention \(g_{0i}^{\rm ref}=g_{0i}^0\exp(c_{0i})\),
the solved cochain is \(c_{0i}=b_0-b_i\). The directed comparison is
\[
\Psi_i=\exp(-b_i),\qquad
\Psi_0g_{0i}^{\rm ref}=g_{0i}^0\Psi_i,
\qquad \Psi_i^{-1}=\exp(b_i).
\tag{OCP.14}
\]
K44/K78\textquotesingle s actual translation and point-line maps realize this exponential
on the compactified charts; the integral deck terms in OCP.7 are inserted
before exponentiation. This is the full source/line cochain arrow.

\par

\medskip\noindent\textbf{OCP.4. The retained parameter endpoint}

\par\nobreak

OCP.14 preserves the endpoint carried by the supplied logarithms. On a
parameter interval its endpoint comparison is the actual composition
\(\Psi_i(1)\Psi_i(0)^{-1}\), with the supplied original closing maps
on the two sides. No condition making this composition the identity
has been imposed. OC45\textquotesingle s joint source/cap flow supplies the local
disc maps; its original prescribed parameter reference still supplies
the logarithm jumps to be inserted into OCP.7 and OCP.14. The lattice
calculation here removes the integral degree-two difference and gives
the displayed analytic comparison without deleting those endpoint data.

The independent certificate re-inverts all twenty selected matrices,
checks every complete product, all primitive generators, the complete
20-by-20 integral map and its inverse, and all directed cochain identities.

\par

\medskip\noindent\textbf{OCP.5. The full integral parabolic complex and its contraction}

\par\nobreak

The same actual slit-sector fillings give the based cellular complex
\[
\Lambda\xrightarrow{d^0}\bigoplus_{i=0}^4\operatorname{im}N_i
\xrightarrow{d^1}\Lambda,
\qquad d^0(v)=(N_i v)_i,\qquad
d^1((z_i)_i)=\sum_iP_i z_i.
\tag{OCP.15}
\]
The local images in this formula are the integral image lattices,
including their displayed primitive generators. The identity
\(d^1d^0=\sum_i(P_{i+1}-P_i)=0\) retains the based-cell transport.
For all four selected sites use
\[
g_0=e_w+3e_\delta,\quad g_1=e_w,\quad
g_2=e_u+3e_w+16e_\delta,\quad
g_3=e_u+8e_\delta,\quad g_4=e_u-3e_w.
\tag{OCP.16}
\]
At the selected site add \(e_\delta\) to the basis. Each unselected
image has the indicated single generator; the selected image is
exactly \(\mathbb Zg_s+\mathbb Ze_\delta\). The coefficient rows
and the primitive preimages in OCP.2-\/-3 verify these assertions over
\(\mathbb Z\). Record the five \(g_i\) coefficients first and the
selected \(e_\delta\) coefficient last, giving coordinates
\(z=(z_0,z_1,z_2,z_3,z_4,z_5)\). In these actual bases the two
matrices, independently rederived at each selected site, are
\[
D=\begin{pmatrix}
-3&-8&0&0\\
0&-1&0&0\\
-2&-3&1&0\\
-1&0&1&0\\
0&3&1&0\\
1&0&0&0
\end{pmatrix},\qquad
B=\begin{pmatrix}
0&0&0&0&0&0\\
0&0&1&-2&1&0\\
1&1&-6&9&-3&0\\
3&0&-8&8&0&1
\end{pmatrix}.
\tag{OCP.17}
\]
Thus \(D\) represents \(d^0\), and the columns of \(B\) are
\(P_i g_i\) followed by \(P_s e_\delta=e_\delta\).
The placement of the extra local generator changes with the selected
site; these matrices do not change after the stated integral basis map.

Here is a primitive preimage for every integral cocycle:
\[
Bz=0\quad\Longrightarrow\quad
z=D\begin{pmatrix}z_5\\-z_1\\z_3+z_5\\0\end{pmatrix}.
\tag{OCP.18}
\]
Indeed the three nonzero equations of \(Bz=0\) successively give
\(z_4=2z_3-z_2\),
\(z_0=3z_2-3z_3-z_1\), and
\(z_2=z_3+3z_1-z_5\). These are precisely the six rows of OCP.18.
No division occurs. For a full chain identity put
\[
Lz=(z_5,-z_1,z_3+z_5)^t,
\qquad
J=\begin{pmatrix}
-24&-8&3\\
0&0&0\\
-9&-3&1\\
0&0&0\\
10&3&-1\\
0&0&0
\end{pmatrix}.
\tag{OCP.19}
\]
Let \(D'\) be the first three columns of \(D\), and let \(B'\)
be the last three rows of \(B\), in the original \((u,w,\delta)\)
order. Direct integer multiplication gives
\[
LD'=I_3,\qquad B'J=I_3,\qquad
D'L+JB'=I_6,\qquad B'D'=0.
\tag{OCP.20}
\]
These identities are explicit splitting and contraction maps for the
actual integral complex. They prove, including all torsion information,
\[
H^0(\mathbf P^1,\mathscr H)=\mathbb Ze_\delta,
\qquad H^1(\mathbf P^1,\mathscr H)=0,
\qquad H^2(\mathbf P^1,\mathscr H)
\xrightarrow{\ \psi\ }\mathbb Z.
\tag{OCP.21}
\]
The degree-zero identification follows from the fourth zero column and
the left inverse of \(D'\). The degree-two identification follows from
the zero first row of \(B\) and its right inverse \(J\) onto the last
three coordinates. All of these are the maps for the primitive K10
fillings used in OCP.9a and OCP.15.

\par

\medskip\noindent\textbf{OCP.6. Global translating automorphisms with the retained period kernel}

\par\nobreak

For the actual Lie sheaf in OCP.11, a global section projects to a
global section of the original source-translation line
\(\mathcal O(-1)\). Its two-chart Laurent equation has only the zero
solution. The remaining physical component is a constant, so
\(H^0(\mathscr E)=\mathbb Ce_\delta\). The full off-diagonal
coefficient is still present in OCP.11-\/-13; the source component is
zero on this global-section calculation.

The actual translation exponential sequence
\(0\to\mathscr H\to\mathscr E\to\mathscr A\to0\), with the
same primitive local period kernels as OCP.15, now has global part
\[
0\longrightarrow\mathbb Ze_\delta
\longrightarrow\mathbb Ce_\delta
\longrightarrow H^0(\mathbf P^1,\mathscr A)
\longrightarrow0.
\tag{OCP.22}
\]
The last surjectivity is exactly OCP.21, with its explicit integral
contraction. In the original physical line the resulting isomorphism
is the literal map
\[
\begin{gathered}
\mathbb C/\mathbb Z\longrightarrow\mathbb C^*,
\qquad [z]\longmapsto e^{2\pi i z},\\
v_M\longmapsto e^{2\pi i z}v_M.
\end{gathered}
\tag{OCP.23}
\]
Its local inverse is \((2\pi i)^{-1}\log c\), whose two choices differ
by the retained primitive physical period. Thus every globally
extending translation is exactly such a physical scalar. The endpoint
in OCP.14 is evaluated through this same map, retaining its scalar
value and its logarithm jump; the original parameter comparison is
not replaced by a selected endpoint.

% End expanded verified cp36 input


\par\bigskip\noindent\textbf{Complete supporting proof 3: OCGC.}\par


% Expanded verified cp36 input: notes/part3_os45_cap_collisions_ocgc_proof.tex
\par

\medskip\noindent\textbf{OS45 prescribed-cap comparison: independent full global audit}

\par\nobreak

This is the complete geometric read of OC45.6-\/-11 in
\path{notes/part3_os45_ordinary_cap_collision_working.md}, source SHA
\texttt{c202e568\allowbreak{}332d9d05\allowbreak{}54d30d3c\allowbreak{}439ff2a1\allowbreak{}fad74c78\allowbreak{}3b7e0a53\allowbreak{}34316eee\allowbreak{}391aaa3b}.
The earlier local source/cap maps and the complete integral calculation
OCP.1-\/-6 are retained as their respective supporting proofs. The following
checks concern the actual global comparison and its parameter endpoint.

\par

\medskip\noindent\textbf{OCGC.1. Original line, infinity frame, and extending translations}

\par\nobreak

Write the actual point coefficient as \(V=f_{P^8}(Z;Q)v_M\).
The translation \((Z,V)\mapsto(e^A Z,e^D V)\) acts in the original
Cartier module by
\[
v_M\longmapsto e^D\frac{f_{P^8}(Z;Q)}{f_{P^8}(e^A Z;Q)}v_M.
\tag{OCGC.1}
\]
Its inverse evaluates the reciprocal at the translated source point.
Composition cancels the intermediate actual point function. Thus its
source and target modules agree also at the point and zero-section
divisors. This is the coefficient map used in OC45.30; it is retained
before passing to a nonzero point coefficient.

The three theta reindexes at infinity give, with the full base factors,
\[
f_{Q^3b}(Z;Q)=(-Z/b)^3Q^{-6}f_b(Z;Q).
\tag{OCGC.2}
\]
Consequently the actual coefficient \(V_0=f_b(Z;Q)v_M\) changes by
\(e^{D-3A}\), while the source changes by \(e^A\). The factors
\(b^3Q^6\) remain in the forward and inverse point-coordinate maps;
they cancel in this action because the base is fixed. The period
matrix with columns \((-1,-3)^t,(3,8)^t\) has determinant one and
inverse with columns \((8,3)^t,(-3,-1)^t\). Its source valuation
row \((8,-3)\) becomes \((1,0)\). These are precisely the full
\(H=h^{-1}B^{-3},K=h^3B^8\) charts and inverse words.

At every source node the two coordinates receive reciprocal source
units; their product and the original height are fixed. At a selected
cone, a character \(Z^rV^d q^e\) receives \(e^{rA+dD}\), with its
literal reciprocal in the inverse. This respects every character
relation and every full period column. The primitive infinity chain
has the same property after the preceding determinant-one period
change. Thus the exponential used by the global comparison is the
actual extending map on the ordinary, nodal, selected, infinity, and
point/O charts, with no replacement of its units by central values.

\par

\medskip\noindent\textbf{OCGC.2. The original Lie bundle and integral comparison}

\par\nobreak

The original infinity change \(x=z^{-2}X,y=z^{-3}Y\) gives the relative
differential equality \(dx/(2y)=z\,dX/(2Y)\). Its dual frame equality
is \(e_\infty=t^{-1}e_0\). On a nodal chart the extending source
field is \(X\partial_X-Y\partial_Y\). This gives the source line
\(\mathcal O(-1)\) in the actual OS45 coordinates. The physical
vertical field is the globally defined \(\delta\). The off-diagonal
Lie coefficient is the invariant source derivative of the logarithm
of the full original point-module transition in OCGC.1-\/-2. It remains
the full function \(a(t,s)\) in
\[
\mathcal L(t,s)=\begin{pmatrix}t^{-1}&0\\a(t,s)&1\end{pmatrix}.
\tag{OCGC.3}
\]

OCP.1-\/-6 independently invert the actual twenty K10 matrices for all
four selected-zero choices. Their based matrices are
\(M_i=P_i(H_i-I)=P_{i+1}-P_i\); the explicit integral map
\(t_j=r_j-\sum_{i>j}(H_i-I)r_i\) carries the primitive raw preimages
to these transported columns. The full exterior cone sends the local
cochain \(z_i=-t_i\) to the positive retained cap-vector difference.
The parabolic complex has the displayed integral contraction
\(D'L+JB'=I_6\), \(LD'=B'J=I_3\). Thus its degree-zero kernel is
\(\mathbb Z\delta\), its degree-one group is zero integrally, and its
degree-two quotient is the original \(\mathbb Z\) measured by \(\psi\).
Both original order-one caps have \(\psi=1\). Their actual difference
therefore has the precise integral preimages in OCP.4-\/-9, including
all source-circle, physical-circle, and suspension components.

After inserting those complete period decks, the remaining Lie
cocycle is solved by OC45.53-\/-54. Substitution in OCGC.3 gives
\(t^{-1}u_\infty-u_0=f\) and
\(a u_\infty+v_\infty-v_0=g\). In particular the full term
\(a u_\infty\) is subtracted before the scalar Laurent split.
The many-disc construction first forms the actual partition cochain,
then subtracts its global Cauchy primitive. The fixed kernels give
holomorphic parameter dependence on each cut. The directed comparison
is exactly
\[
F_i=\operatorname{Exp}(-s_i),\qquad
F_i g^X_{ij}=g^Y_{ij}F_j,\qquad F_i^{-1}=\operatorname{Exp}(s_i).
\tag{OCGC.4}
\]
These are whole source and line maps through all of the preceding
charts. The base logarithms \(\log t\) and \(\log(t-d)\) retain their
individual prescribed values throughout this construction.

\par

\medskip\noindent\textbf{OCGC.3. Full parameter seam and physical scalar comparison}

\par\nobreak

On the original cap \(q_n=-s\), a positive parameter turn changes
\(h\) by \(\alpha\). For the eighteen unselected collisions the full
\(C(0,s)\) is an extending nonzero unit. For the six selected
collisions it is \(-\lambda s\beta(-s,s)\). The literal reference
\(\gamma_0+k h+r\alpha+d_\delta\delta\) therefore changes by
\(k\alpha\) and \(k\alpha+\delta\), respectively. These full vectors
are inserted into the same based integral cochain before the analytic
solve. No source or physical period has been set to zero.

A global Lie section has source component zero: the equation
\(u_0=t^{-1}u_\infty\) has no common nonzero regular Laurent power.
Its remaining component is a constant multiple of \(\delta\).
The integral degree-one vanishing just proved then identifies all
global extending translations with the exact physical action
\[
\mathbb C\delta/\mathbb Z\delta\longrightarrow\mathbb C^*,
\qquad [z\delta]\longmapsto e^{2\pi i z}.
\tag{OCGC.5}
\]
In particular ratios of the constructed comparisons are this action
on every original cap and selected fan.

For the local parameter comparisons \(F_\alpha=c_{\alpha\beta}F_\beta\),
the connection satisfies
\(\mathfrak a_\alpha-\mathfrak a_\beta=d\log c_{\alpha\beta}\).
The scalar coefficient transforms as
\(z_\alpha=c_{\alpha\beta}^{-1}z_\beta\), so its actual horizontal
equation is \(dz_\alpha=-\mathfrak a_\alpha z_\alpha\). This checks
the sign in OC45.61. The retained angular holonomy is flattened by
its explicit exponential, and the radial equation respects these
same transitions. Its antiholomorphic derivative is therefore a
single scalar form; the stated Cauchy kernel gives the global
primitive on a smaller annulus. Multiplication by its negative
exponential gives the closed holomorphic family and its reciprocal.

On the cut, the complete endpoint equality is
\[
F_1E_XF_0^{-1}=CE_Y,
\qquad C=F_1E_XF_0^{-1}E_Y^{-1}.
\tag{OCGC.6}
\]
All source/period/Cartier comparisons above prove that this particular
ratio is the physical scalar of OCGC.5. Taking the actual inverse
gives \(F_0E_X^{-1}F_1^{-1}=E_Y^{-1}C^{-1}\), without commuting any
of the other maps. Inserting the full parallel transport \(P_Y\)
then gives precisely OC45.64. The literal physical isotopy
\(e^{\theta\log C}\) extends through the same cap and cone charts.
This proves the return comparison with its supplied endpoint.

\par

\medskip\noindent\textbf{OCGC.4. Protected cap, proper disc, and the stated return}

\par\nobreak

The point \(d\) and its cap disc are chosen inside the actual nodal
coordinate neighbourhood and outside all moving critical values.
There the full \(Q,P,C\) are nonzero, and the fixed joint logarithms
give the exact plus-log map \(\mathcal F_{\log(t-d)/(2\pi i)}\)
with inverse \(\mathcal F_{-\log(t-d)/(2\pi i)}\). A base turn is the
whole original suspension \(B\), so its norm is the same positive
one and its source/line components remain those of the original flow.

At the old cap location the reference family uses the actual moving
node \(XY=t-s\), or the prescribed selected fan with this same
height. These are submersions over \(s\). The original infinity
point is \(b_0=16/27\) at the parameter centre, so its point/O module
and both full period words extend there. The protected ordinary cap,
all other simple source nodes, and all selected cone charts have the
displayed reciprocal overlap maps. Uniform contraction on the finitely
many compact pieces and the finite period representatives give the
stated compact total preimage of a closed parameter disc. Thus this
is the actual smooth proper disc family used in OC45.10.

The integral and analytic comparison carries the original positive
norm cap to that family. Its reference sphere and each prescribed
fan are carried by the retained actual ruling/Gysin/van Kampen maps.
The disc family has its identity marked endpoint; a partition of
unity applied to its local submersion lifts gives an oriented smooth
transport with inverse. OCGC.6 identifies the original return with
that transport followed by the explicitly isotopic physical scalar.
Consequently OC45.66 holds for all twenty-four stated meridians,
their reciprocals, the original order-one positive-norm reference,
and every prescribed selected fan. The radial scalar comparison
retains its argument and introduces no scalar turn. It transports
this same conclusion to every admitted original fixed scalar.

This audit assigns exactly the scope of OC45.6-\/-11. The primitive
integral calculation, original point-module maps, complete parameter
jumps, and physical endpoint morphism all precede the zero-return
conclusion.

% End expanded verified cp36 input


\par\bigskip\noindent\textbf{Complete supporting proof 4: ORC.}\par


% Expanded verified cp36 input: notes/part3_os45_cap_collisions_orc_proof.tex
\par

\medskip\noindent\textbf{OS45 ordinary-cap collisions: complete root geometric peer}

\par\nobreak

This independent root read covers all OC45.1-\/-11, the complete original
O45 source/line/collar construction, OCP.1-\/-6, and OCGC.1-\/-4. The actual
source, its order-one positive-norm cap, the prescribed original
reference, the nonzero scalar, and the selected integral fan are
retained throughout. The exact executable accompanying this read
rederives the algebra from the original five K10 matrices and the
displayed source equation. The geometric arrows below establish the
extension and endpoint assertions.

\par

\medskip\noindent\textbf{ORC.1. Original parameter and all twenty-four marked points}

\par\nobreak

Use the original source and original selected-node map
\[
\begin{gathered}
R=t^2+a,\quad f=R^2+t,\quad
g=R^3+\tfrac32tR+\tfrac38,\qquad
y^2=x^3-3fx+2g,\quad G=(R,\sqrt3/2),\\
a=-u/3-u^4/9,\qquad t_c=1/(2u)-u^2/3,\qquad
v^2=4u^3+3 .
\end{gathered}
\tag{ORC.1}
\]
The inverse to the marked-node substitution is the original
\(u=\xi-(t_c^2+a)\), with \(\xi\) the selected double cubic root;
the ordered normalization has \(c=v/(2u)\) and
\(Z=(S-c)/(S+c)\), \(S=c(1+Z)/(1-Z)\).
Thus the \(u,v\) labels identify the actual source branches.

The ordinary cap at \(t=0\) meets a source node exactly at
\[
16a^3+3=0,\qquad
2187-16u^3(u^3+3)^3
=-(2u^3-3)(8u^9+84u^6+342u^3+729)=0.
\tag{ORC.2}
\]
The degree-twelve polynomial is square-free, is nonzero at \(u=0\),
and has no common root with \(4u^3+3\), \(u^3+3\), or
\(16u^6+120u^3+243\). Each of its twelve roots consequently has
two distinct ordered \(v\)-values. The factor \(2u^3-3\) makes
\(t_c=0\), and gives six selected-cap points; the other factor
gives the eighteen unselected-cap points. The reciprocal branch
coordinate \(Z^{-1}\) is retained for the other \(v\)-value.

At every such point write \(a_*\) for the actual source parameter.
The finite discriminant factor and its derivatives give
\[
D=48a(t^2+a)^2+8t^3+72at+9,\quad
D_t(0,a_*)=72a_*,\quad D_a(0,a_*)=144a_*^2 .
\tag{ORC.3}
\]
The moving critical value \(s=t_n(a)\) therefore has
\(ds/da=-2a_*\). Moreover \(da/du=-(4u^3+3)/9=-v^2/9\)
is nonzero on this locus. The actual maps
\((u,v)\mapsto a\mapsto s\) are holomorphic local coordinates
with their uniquely fixed implicit inverses. A positive meridian
in the ordered source curve makes one positive \(s\)-turn.
The selected value is exactly \(s\) at the six selected-cap points.

\par

\medskip\noindent\textbf{ORC.2. Original node and complete point-line coefficient}

\par\nobreak

At \(s=0\) the source node is \((x,y,t)=(-a_*,0,0)\) and
\[
y^2=(x+a_*)^2(x-2a_*).
\tag{ORC.4}
\]
Set \(q=t-s\), choose the original branch
\(\xi(q,s)=\sqrt{f(s+q,a(s))}\) with \(\xi(0,0)=-a_*\),
and use the complete analytic unit
\[
h(q,s)=\frac{2(g(s+q,a(s))-\xi(q,s)^3)}q,\qquad h(0,0)=6a_*.
\tag{ORC.5}
\]
The numerator vanishes along \(q=0\), since it is the moving
double-root equation. Direct multiplication gives the actual maps
\[
X=y-(x-\xi)\sqrt{x+2\xi},\qquad
Y=\frac{y+(x-\xi)\sqrt{x+2\xi}}h,\qquad XY=q,
\tag{ORC.6}
\]
\[
q=XY,\quad t=s+XY,\quad y=(X+hY)/2,\qquad
2(x-\xi)\sqrt{x+2\xi}=hY-X.
\tag{ORC.7}
\]
The last equation has derivative \(2\sqrt{3\xi}\) at the node.
These maps and their inverse extend on the joint source/parameter
neighbourhood. In particular \((X,Y,s)\) are smooth total-space
coordinates with submersive projection to \(s\).

The central normalization and section are
\[
x=S^2+2a_*,\quad y=S(S^2+3a_*),\quad
c^2=-3a_*,\quad S_G=\frac{\sqrt3}{4a_*},\quad S_G^2=-a_*.
\tag{ORC.8}
\]
Thus \(c/S_G=\pm\sqrt3\) and
\(P_*=Z(G)=-2\pm\sqrt3\). Their product is one and their
absolute values differ from one. Every point \(P_*^j\) used in
the three original Miller doublings is consequently defined
and different from the zero-section coordinate one.

The original tangent normalization gives, for \(m=2,4,8\),
\[
f_m(Z)=(-2c)^{m-1}\frac{1-P^m}{(1-P)^m}
 \frac{(Z-P)^m}{(Z-P^m)(Z-1)^{m-1}}.
\tag{ORC.9}
\]
Its divisor is \(m[G]-[mG]-(m-1)[O]\); substituting
\(S=c(1+Z)/(1-Z)\) shows the principal coefficient at \(O\)
is exactly \(S^{m-1}\). These two checks agree with the
successive original tangent quotients. For the actual direction
\(T=-G\), the full chronological coefficient is
\[
B_{8,1}=f_2^8\frac{f_8(Z/P)}{f_8(Z)}
 =(-2S_G)^8\frac{(Z-P)(Z-P^8)}{(Z-P^9)(Z-1)}.
\tag{ORC.10}
\]
With the actual point function \(f_b=(1-Z/b)/(1-Z)\),
\[
B_{8,1}(Z)\frac{f_{P^8}(Z/P)}{f_{P^8}(Z)}
=(-2S_G)^8=\frac{81}{256a_*^8}=\frac9{a_*^2}.
\tag{ORC.11}
\]
The equality cancels the four displayed source and target Cartier
factors; their point/O modules remain the modules of the map.
The inverse evaluates the reciprocal at the inverse source point.

Retain the complete inverse-Tate functions \(Q=q\kappa(q,s)\),
\(P=P(q,s)\), and the actual Miller/theta coefficient \(\beta(q,s)\),
fixed at the original regular reference. Their central values are
\[
\kappa(0,0)=-\frac{h(0,0)}{1728\xi(0,0)^3}
 =\frac1{288a_*^2},\qquad
\beta(0,0)=C_*\frac9{a_*^2}.
\tag{ORC.12}
\]
Here \(C_*\) is the original fixed Hom normalization. The full
period arrows, before any local logarithm, are
\[
h_{\mathrm{per}}:(Z,V)\longmapsto(QZ,P^8V),\qquad
B:(Z,V)\longmapsto(P^{-1}Z,CV),\quad
C=\lambda(t-t_c)\beta .
\tag{ORC.13}
\]
All higher coefficients of these functions occur in the maps below.

\par

\medskip\noindent\textbf{ORC.3. Extending translations and the infinity component map}

\par\nobreak

A Lie translation with complete analytic coefficients \(A,D\) is
\[
\mathcal T_{A,D}(Z,V,t,s)=(e^AZ,e^DV,t,s),
\quad
\mathcal T_{A,D}^{-1}=\mathcal T_{-A,-D}.
\tag{ORC.14}
\]
The base is fixed. On the original point module its coefficient is
\[
e^D\frac{f_{P^8}(Z;Q)}{f_{P^8}(e^AZ;Q)}.
\tag{ORC.15}
\]
Composition cancels the middle point function at the actual
intermediate source point. The reciprocal there proves the
inverse identity. At a zero or pole the numerator and denominator
are the source and target Cartier generators, so the same formula
is the module isomorphism through their graphs.

On a source chain, \(X_i=Z/Q^i,\ Y_i=Q^{i+1}/Z\) receive the
units \(e^A,e^{-A}\). Their product \(Q\), and hence the original
height \(q\) under \(Q=q\kappa\), is fixed. At a selected fan cone,
the actual character \(Z^rV^d q^e\) receives \(e^{rA+dD}\).
This proves the map, its inverse, and every monomial overlap on
each prescribed unimodular fan.

At infinity keep \(z=1/t\), the original minimal source coordinates,
and the full functions
\[
Q_\infty=z^8\kappa_\infty,\qquad
P_\infty=z^3\pi,\qquad
b_0=P_\infty^8/Q_\infty^3,\qquad
b_0(0,s)=-1/(9a(s)^3).
\tag{ORC.16}
\]
The complete theta identity is
\[
f_{Q_\infty^3b_0}(Z;Q_\infty)
 =(-Z/b_0)^3Q_\infty^{-6}f_{b_0}(Z;Q_\infty).
\tag{ORC.17}
\]
Writing \(V_0=f_{b_0}v_M\) therefore gives
\((Z,V_0)\mapsto(e^AZ,e^{D-3A}V_0)\).
Both base factors \(b_0^3Q_\infty^6\) occur in the reciprocal
coordinate map; fixing the base cancels them in this action.

The original source valuations of \(h_{\mathrm{per}},B\) are
\((8,-3)\). Their actual period words and inverse are
\[
H=h_{\mathrm{per}}^{-1}B^{-3},\quad
K=h_{\mathrm{per}}^3B^8,\qquad
h_{\mathrm{per}}=H^8K^3,\quad B=H^{-3}K^{-1}.
\tag{ORC.18}
\]
The integer matrix
\(\left(\begin{smallmatrix}-1&3\\-3&8\end{smallmatrix}\right)\)
has inverse
\(\left(\begin{smallmatrix}8&-3\\3&-1\end{smallmatrix}\right)\);
the row \((8,-3)\) becomes \((1,0)\).
The line multiplier of \(K\) is the full unit \(\lambda^8\)
times its original analytic coefficient. In particular every
component shift has its explicit inverse period word: a source
order \(n\) is removed by \(H^{-n}\). These words are evaluated
on the original source and line, including their unit factors.
The rank-one chain action of (ORC.14) consequently extends
through every original infinity component and descends to the
same compact quotient. This proves the component transfer used
by the exponential, with its inverse.

\par

\medskip\noindent\textbf{ORC.4. The actual local exponential and original Lie bundle}

\par\nobreak

Let \(\mathscr A\) be the sheaf of extending fibre translations
in the just displayed charts, \(\mathscr E\) their Lie bundle,
and \(\mathscr H\) the period kernel. The local inverse to
\(\operatorname{Exp}\) is obtained as follows. On an unselected
node, remove the integral source order by the period
\(h_{\mathrm{per}}\), whose source valuation is one. At the
selected filling the two period valuations are exactly
\((1,0),(0,1)\); their integral inverse removes the source
and line orders. At infinity use (ORC.18) to remove the source
order. The remaining line multiplier has order zero because
the map and its inverse preserve the nonzero line quotient.
In each chart the remaining multipliers are nonvanishing
analytic units; their analytic logarithms give \(A,D\).
The differences of two logarithms are precisely the retained
angular periods and the extending zero-radial-order period
words. All removals are period maps on the quotient, and
their inverse words were specified above.

The radial coefficients of a local full period vanish exactly
when the original monodromy fixes that period. Thus
\(\mathscr H=j_*\Lambda\), with stalk \(\ker(H_i-I)\) at the
five source fillings, and the explicit maps give
\[
0\longrightarrow\mathscr H\longrightarrow\mathscr E
 \xrightarrow{\operatorname{Exp}}\mathscr A\longrightarrow0.
\tag{ORC.19}
\]
On a slit source disc the annular complex is
\(\Lambda\xrightarrow{H_i-I}\Lambda\). Attaching its centre
keeps its invariant zero-cochains; the same annular
coinvariant term occurs with the opposite boundary in the
disc/annulus mapping cone and cancels by its identity map.
This gives the invariant stalk and zero positive local
cohomology used in the global based complex below.

The original coordinate equality
\[
\frac{dx}{2y}=z\,\frac{dX}{2Y},\qquad
e_\infty=t^{-1}e_0
\tag{ORC.20}
\]
identifies the source Lie line with \(\mathcal O(-1)\).
At a source node its vector is
\(X_i\partial_{X_i}-Y_i\partial_{Y_i}\).
The physical vector \(\delta\) has literal flow
\(v_M\mapsto e^{2\pi iw}v_M\). Differentiating the entire
point-module transitions (ORC.15),(ORC.17) gives
\[
0\longrightarrow\mathcal O\delta\longrightarrow\mathscr E
 \longrightarrow\mathcal O(-1)\longrightarrow0,\qquad
\mathcal L(t,s)=
\begin{pmatrix}t^{-1}&0\\ a(t,s)&1\end{pmatrix}.
\tag{ORC.21}
\]
The coefficient \(a(t,s)\) is this full logarithmic derivative.
No truncation of it occurs in the cochain solution.

\par

\medskip\noindent\textbf{ORC.5. Literal integral contraction of the five source fillings}

\par\nobreak

Take the five original K10 matrices in the order
\((\psi,u,w,\delta)\) and set
\[
H_i=T_i^{-t},\quad N_i=H_i-I,\quad
P_0=I,\quad P_i=H_0\cdots H_{i-1},\quad
M_i=P_iN_i=P_{i+1}-P_i.
\tag{ORC.22}
\]
The original product is \(H_0\cdots H_4=I\).
For the selected site \(\sigma=1,2,3,4\), respectively, use
\[
d_1=e_\psi,\quad d_2=e_\psi+2e_w,\quad
d_3=e_\psi+e_w,\quad d_4=e_\psi .
\tag{ORC.23}
\]
Literal inversion of the original matrices, with their selected
central entries, gives
\[
N_\sigma d_\sigma=e_\delta,\quad
N_1(-e_u)=e_w,\quad
N_3e_w=e_u+8e_\delta.
\tag{ORC.24}
\]
For every retained vector \(v=Ae_u+Be_w+Ce_\delta\), set
\[
r_1=-Be_u,\quad r_3=Ae_w,\quad
r_\sigma\mathrel{+}=(C-8A)d_\sigma,\qquad
t_j=r_j-\sum_{i>j}N_ir_i.
\tag{ORC.25}
\]
Repeated indices receive both terms. The inverse is the
descending recursion \(r_j=t_j+\sum_{i>j}N_ir_i\).
The telescoping equality \(\sum_{j<i}M_j=P_i-I\) gives
\(\sum_jM_jt_j=\sum_iN_ir_i=v\).
Thus these are inverse integral based-cell coordinate maps.

The oriented cone boundary is
\[
b_i=p_i-N_iz_i\quad(i<4),\qquad
b_4=-H_4\sum_{i<4}P_ip_i-N_4z_4.
\tag{ORC.26}
\]
Its coordinate \(v_*=\sum_iP_ib_i\) cancels every \(p_i\)
and equals \(-\sum_iM_iz_i\); hence the positive \(v\)
uses precisely \(z_i=-t_i\). This keeps the original
suspension terms in \(d_\sigma\) with their whole period
arrows, including at both ends of the corresponding cut.

After the local disc/annulus cancellation in (ORC.19), the
remaining integral complex is
\[
\Lambda\ \xrightarrow{D}\ \mathbb Z^6\
 \xrightarrow{B}\ \Lambda .
\tag{ORC.27}
\]
The first five local generators are
\[
e_w+3e_\delta,\ e_w,\ e_u+3e_w+16e_\delta,\
e_u+8e_\delta,\ e_u-3e_w;
\]
the sixth is \(e_\delta\) at the selected site. Their primitive
preimages in (ORC.24) give the full image lattices. In those
actual bases, for each of the four selected sites,
\[
D=\begin{pmatrix}
-3&-8&0&0\\0&-1&0&0\\-2&-3&1&0\\
-1&0&1&0\\0&3&1&0\\1&0&0&0
\end{pmatrix},\qquad
B=\begin{pmatrix}
0&0&0&0&0&0\\0&0&1&-2&1&0\\
1&1&-6&9&-3&0\\3&0&-8&8&0&1
\end{pmatrix}.
\tag{ORC.28}
\]
Write \(D'\) for the first three columns and \(B'\) for the
last three rows. The exact splitting maps are
\[
Lz=(z_5,-z_1,z_3+z_5)^t,\qquad
J=\begin{pmatrix}
-24&-8&3\\0&0&0\\-9&-3&1\\0&0&0\\
10&3&-1\\0&0&0
\end{pmatrix},
\tag{ORC.29}
\]
\[
LD'=I_3,\quad B'J=I_3,\quad
D'L+JB'=I_6,\quad B'D'=0.
\tag{ORC.30}
\]
Every map is integral. These identities supply the inverse
on each image and the contraction of degree one, including
all torsion information. They prove
\[
H^0(\mathscr H)=\mathbb Z\delta,\qquad
H^1(\mathscr H)=0,\qquad
H^2(\mathscr H)\xrightarrow{\psi}\mathbb Z ,
\tag{ORC.31}
\]
where the last map has section \(1\mapsto[e_\psi]\).
Smooth cap discs add their transported full period vector
to \(v_*\). The two order-one positive-norm caps considered
below therefore differ by exactly the vector \(v\) in
(ORC.25), retaining all three of its integer coordinates.

\par

\medskip\noindent\textbf{ORC.6. Full analytic map between the two cap constructions}

\par\nobreak

Take the original cap at \(t=0\) and a protected ordinary cap
at \(t=d\). Use a specified based path on a cut parameter
disc to transport their full lattice markings. Their actual
clutch logarithms use \(\log t\) and \(\log(t-d)\).
Their difference has a period-valued two-cocycle with
\(\psi=1-1=0\). Subtract the explicit cochain
(ORC.25)-\/-(ORC.26), evaluating every inserted period using
the original source, line and scalar functions. The two
branches on a slit annulus differ by the cancelled
\(N_iz_i\), so the remaining Lie cocycle is single valued.

Here is its analytic solution with the full off-diagonal
coefficient. For an overlap vector \((f,g)\) in the frames
(ORC.21), let its Laurent coefficients be the fixed-contour
Cauchy integrals and put
\[
u_0=-\sum_{j\geq0}f_jt^j,\qquad
u_\infty=\sum_{j<0}f_jt^{j+1},\qquad
G=g-a(t,s)u_\infty ,
\tag{ORC.32}
\]
\[
v_0=-\sum_{j\geq0}G_jt^j,\qquad
v_\infty=\sum_{j<0}G_jt^j.
\tag{ORC.33}
\]
Direct substitution gives
\(\mathcal L(u_\infty,v_\infty)^t-(u_0,v_0)^t=(f,g)^t\).
Both infinity components are holomorphic in \(1/t\);
the \(j=-1\) source coefficient is an allowed constant.

On the full many-disc cover first form the partition cochain
\(s_i^{\mathrm{sm}}=\sum_k\rho_k\theta_{ik}\).
The cocycle equation gives
\(s_i^{\mathrm{sm}}-\mathcal L_{ij}s_j^{\mathrm{sm}}=\theta_{ij}\).
Its antiholomorphic derivative is a global \(\mathscr E\)-valued
form. Componentwise Cauchy primitives on the two affine
charts have holomorphic difference; (ORC.32)-\/-(ORC.33)
remove that difference. Subtracting this global primitive
gives holomorphic \(s_i\) with the same cocycle.
Fixed partitions and kernels retain holomorphic dependence
on the cut parameter. This construction also proves
\(H^1(\mathscr E)=0\) with its actual transition.

For directed overlaps \(g^X_{ij}=t_{\theta_{ij}}g^Y_{ij}\),
the resulting full maps are
\[
F_i=\operatorname{Exp}(-s_i),\qquad
F_i^{-1}=\operatorname{Exp}(s_i),\qquad
F_i g^X_{ij}=g^Y_{ij}F_j .
\tag{ORC.34}
\]
Equations (ORC.14)-\/-(ORC.18) give their literal coefficients
and reciprocals in the original point/O, source-chain,
infinity and selected-fan coordinates. On each smooth cap
they are the same exponential before its plus-log gluing.
Thus the comparison carries the complete original cap
reference and its source/physical integers.

\par

\medskip\noindent\textbf{ORC.7. Both parameter jumps and the global comparison group}

\par\nobreak

At the original cap \(q=-s\), write the actual period on a cut as
\[
\gamma_{\mathrm{ref}}=
\frac{(-\log P,\log C)}{2\pi i}
+k\,\frac{(\log Q,8\log P)}{2\pi i}
+r\alpha+d_\delta\delta .
\tag{ORC.35}
\]
All coefficients \(k,r,d_\delta\) are the original integers
in the based marking. A positive \(s\)-turn adds \(\alpha\)
to the second displayed full period. At an unselected-cap
point \(C(0,s)\) is an extending unit. At a selected-cap
point it is exactly \(-\lambda s\beta(-s,s)\). Therefore
\[
\Delta\gamma_{\mathrm{ref}}=k\alpha
\quad(18\text{ points}),\qquad
\Delta\gamma_{\mathrm{ref}}=k\alpha+\delta
\quad(6\text{ points}).
\tag{ORC.36}
\]
The logarithms of every extending unit have their full
continued values. These vectors have zero \(\psi\)-coordinate;
their based transports are inserted into (ORC.25)-\/-(ORC.34)
at the final parameter cut. This retains the physical jump
at the six selected-cap points.

A global source Lie section satisfies \(u_0=t^{-1}u_\infty\).
Regular powers at zero and infinity force it to vanish.
The remaining global Lie section is a constant physical
vector. The exact exponential (ORC.19), together with the
integral degree-one contraction (ORC.30), gives
\[
H^0(\mathscr A)=\mathbb C\delta/\mathbb Z\delta
 \xrightarrow{\ \sim\ }\mathbb C^*,\qquad
[w\delta]\longmapsto(v_M\mapsto e^{2\pi iw}v_M).
\tag{ORC.37}
\]
The local inverse is \((2\pi i)^{-1}\log c\), and two
choices differ by the retained primitive period.
Consequently every ratio of the constructed global
comparisons is this literal physical scalar on every
original chart.

\par

\medskip\noindent\textbf{ORC.8. Closed parameter comparison and directed endpoint}

\par\nobreak

For local parameter comparisons \(F_\alpha=c_{\alpha\beta}F_\beta\),
all \(c_{\alpha\beta}\) are the scalars in (ORC.37). Choose a
partition and form
\(\mathfrak a_\alpha=\sum_\beta\rho_\beta\,d\log c_{\alpha\beta}\).
Then
\(\mathfrak a_\alpha-\mathfrak a_\beta=d\log c_{\alpha\beta}\).
The scalar coefficients of a comparison transform by
\(z_\alpha=c_{\alpha\beta}^{-1}z_\beta\), so their horizontal
equation is exactly
\[
dz_\alpha=-\mathfrak a_\alpha z_\alpha,\qquad
z_\alpha(\rho,\theta)=z_\alpha(\rho_0,\theta)
\exp\!\left(-\int_{\rho_0}^{\rho}
\mathfrak a_\alpha(\partial_r)\,dr\right).
\tag{ORC.38}
\]
Angular transport has a nonzero scalar holonomy \(c\).
Multiplication by \(e^{-\chi(\theta)\log c}\), with
\(\chi(\theta+1)=\chi(\theta)+1\) and constant endpoint
collars, makes the comparison periodic with all derivatives.
The radial equation extends it to a thin annulus.
Its antiholomorphic derivative is a single scalar form
\(\bar\partial F^{\mathrm{sm}}=\eta F^{\mathrm{sm}}\).
After a compactly supported extension, the Cauchy integral
\[
u(s)=\frac1\pi\int_{\mathbb C}\frac{\eta_{\bar w}(w)}{s-w}\,dA(w)
\tag{ORC.39}
\]
satisfies \(\bar\partial u=\eta\) on a smaller annulus.
Thus \(e^{-u}F^{\mathrm{sm}}\) is a nonzero closed
holomorphic comparison, with its explicit reciprocal.

The endpoint direction is also visible without choosing
this closed section. Let \(F_\theta:X_\theta\to Y_\theta\)
be (ORC.34), and let \(E_X:X_0\to X_1\),
\(E_Y:Y_0\to Y_1\) be the full forward marked endpoints.
All continued periods in (ORC.36) are included. The exact
ratio and its inverse equation are
\[
C=F_1E_XF_0^{-1}E_Y^{-1}\in\mathbb C^*_{\mathrm{physical}},
\qquad F_0E_X^{-1}F_1^{-1}=E_Y^{-1}C^{-1}.
\tag{ORC.40}
\]
For actual parallel transport \(P_Y:Y_0\to Y_1\),
\[
F_0\bigl(E_X^{-1}F_1^{-1}P_YF_0\bigr)F_0^{-1}
=E_Y^{-1}C^{-1}P_Y.
\tag{ORC.41}
\]
This follows by composition in the written order. No
commutation with a source endpoint or transport is used.
The scalar factor has the physical isotopy
\(e^{\theta\log C}\), whose maps and inverses extend on
every original cap, line, infinity and selected fan.

\par

\medskip\noindent\textbf{ORC.9. Actual protected-cap family across the parameter centre}

\par\nobreak

Choose a fixed \(d\ne0\) in the original nodal coordinate
neighbourhood and a cap disc of radius less than \(|d|/4\).
A sufficiently small closed \(s\)-disc keeps it disjoint
from all moving critical values and the selected value.
There the complete \(Q,P,C\) are nonzero, and their
analytic logarithms are fixed at the same reference.
The actual flow and cap maps are
\[
\mathcal F_w(Z,V)=
(e^{-w\log P}Z,e^{w\log C}V),\quad
\mathcal F_w^{-1}=\mathcal F_{-w},\qquad
\Psi_d=\mathcal F_{\log(t-d)/(2\pi i)}.
\tag{ORC.42}
\]
A branch change gives the whole \(B\)-period. The maps
therefore descend to the compact quotient, and their
reciprocal is \(\mathcal F_{-\log(t-d)/(2\pi i)}\).
This is the same positive norm one, with all other
period components recorded by the full logarithms.

At the old cap location the family \(Y\) uses the actual
source node (ORC.6)-\/-(ORC.7), or the prescribed selected
fan with \(q=t-s\). The source coordinates \((X,Y,s)\)
and the unimodular cone coordinates with \(s\) give
submersions. At infinity (ORC.16) has
\(b_0(0,0)=16/27\), so its point/O modules and full
period words extend. The other source critical values
remain distinct and simple by the original discriminant
and the disjoint polynomial loci in ORC.1. Their original
charts and full unit overlaps extend on the same disc.

Choose a common strictly contracting scalar for these
compact source pieces. The contracting line fundamental
annuli are compact. At infinity the actual \(H,K\) words
give a compact fundamental set for the primitive chain
quotient; the original eight components are carried by
their inverse words. At the selected filling the two
primitive valuation columns and finitely many cones
modulo their periods give a compact fundamental set.
The protected smooth cap and the remaining finite
source pieces are compact over their closed base discs.
All unit multipliers are bounded away from zero on
the closed parameter disc. Gluing these finitely many
compact pieces by the stated collar maps and their
inverses gives a Hausdorff compact preimage of that disc.
The local submersions agree on the same maps. This
establishes a smooth proper family \(Y\) across \(s=0\).

On a reference parameter, (ORC.34) identifies \(Y\)
with the original K10 sphere with its full cap reference.
Each prescribed fan is carried from that sphere by the
retained finite canonical fan word, the actual common
blowup and two-ruling maps, and their primitive Gysin
and van Kampen maps in F35/O45P.6. Thus the disc
family has the required oriented sphere fibre.
Partitions of unity on the displayed submersion charts
give a global horizontal lift; compactness gives its
flow and inverse for every radial and angular path.
A radial contraction of the parameter circle supplies
the isotopy between its return and the identity.

For an original admitted scalar outside the common
small disc, use the actual path
\(\lambda_\tau=((1-\tau)+\tau\varepsilon)\lambda_0\)
with \(0<\varepsilon<1\). The complete suspension
coefficient is linear in \(\lambda\), so this path
remains contracting and has zero argument change.
The same source, point-line, cap and fan arrows give
the compact cylinder and its inverse transports.
This conjugates the marked returns to the common-disc
calculation while preserving the original endpoints.

\par

\medskip\noindent\textbf{ORC.10. Original-reference return and precise scope}

\par\nobreak

The full maps of \(Y\) are defined on the parameter
disc, so its forward endpoint around the circle is
the identity. Its actual return is isotopic to the
identity by the flow just constructed.
Equation (ORC.41) transfers this return to the original
family \(X\), and (ORC.40) removes the remaining
physical scalar by its explicit isotopy. Therefore
\[
[\Phi]=0\quad\text{in }\Theta_7=\mathbb Z/28
\tag{ORC.43}
\]
for each of the twenty-four ordinary-cap collision
meridians of the original ordered OS45 source cover,
with its original order-one positive-norm cap, complete
prescribed reference, every prescribed lattice-basis
selected fan, and original admitted fixed scalar.
Each inverse meridian uses the reciprocal source,
line, cap, endpoint and transport maps and has the
same zero class.

For the eighteen unselected-cap points the original
particular period also gives the local enlarged-cap
construction: the middle ordinary region maps by
\(\Psi_0^{-1}\), the small cap and exterior map by
their identities, and both collar equations are
\(\Psi_0^{-1}\Psi_0=1\). Its inverse is obtained
piecewise with the same full flow. The original
reference assertion above additionally carries every
retained integer and the six selected-cap jumps.
Part3.4 remains active on the separately labelled
original-source generators outside this completed
ordinary-cap domain.

% End expanded verified cp36 input



\par\medskip\noindent\textbf{Exact construction and audit dependencies.}\par
The original K10 source, point line, five marked monodromy matrices,
primitive local fillings and sphere construction are retained in
\path{notes/part3_os45.md}. The complete K116 original marked-source,
infinity and cap construction is
\path{notes/part3_os45_marked_source_working.md};
its independent all-fan reference-sphere antecedent is
\path{notes/part3_os45_unselected_ii_blind_peer.md}.

The primary local and global-reference executables are
\path{code/part3/os45_ordinary_cap_collision.py} and
\path{code/part3/os45_ordinary_cap_reference.py}.
The independent integral and global-reference executables are
\path{code/part3/os45_cap_period_checks.py} and
\path{code/part3/os45_cap_reference_checks_peer.py}.
The full root execution is
\path{code/part3_root/os45_cap_collision_full_peer.py}.
Their exact results and full geometric audit records retain their
separate execution scopes in
\path{records/k119_export_manifest.json}.
The original preparation-stage readiness record remains preserved.
The completed source, Cartier, integral, analytic, cap, properness
and endpoint proofs appear above in full.


% End expanded verified cp36 input

\end{proof}


% END CONTINUATION_NIGHT36_RECEIVING
